% !TeX program = xelatex
% 数学题面为原生 LaTeX。图题为内嵌 TikZ 路径。
\documentclass[UTF8,zihao=-4,fontset=fandol]{ctexart}
\ifdefined\QuestionBankAfour
\usepackage[a4paper,margin=17mm]{geometry}
\else
\usepackage[paperwidth=800bp,paperheight=600bp,left=52bp,right=52bp,top=42bp,bottom=47bp]{geometry}
\fi
\usepackage{amsmath,amssymb,xcolor,fancyhdr,enumitem,tikz}
\usetikzlibrary{svg.path}
\DeclareMathOperator{\arcsec}{arcsec}
\renewcommand{\lim}{\operatorname*{lim}\limits}
\usepackage[hidelinks,bookmarksopen=true]{hyperref}
\definecolor{ink}{HTML}{192635}
\definecolor{muted}{HTML}{45566A}
\definecolor{accent}{HTML}{175C57}
\setlength{\parindent}{0pt}
\setlength{\headheight}{15bp}
\pagestyle{fancy}
\fancyhf{}
\renewcommand{\headrulewidth}{0pt}
\fancyfoot[L]{\small\textcolor{muted}{\hyperlink{contents}{返回目录}}}
\fancyfoot[R]{\small\textcolor{muted}{\thepage}}
\raggedbottom
\newcommand{\SectionHeading}[4]{%
  \clearpage\hypertarget{sec-#1}{}%
  {\small\textcolor{accent}{#2\hfill 第 #1 节}}\par\medskip
  {\LARGE\bfseries #3}\par\smallskip
  {\small\textcolor{muted}{#4}}\par\bigskip}
\newcommand{\Theory}[4]{%
  {\small\textcolor{accent}{#1}\quad\textbf{#3}\hfill\textcolor{muted}{教材 PDF 第 #2 页}}\par
  {\small #4}\par\smallskip}
\newcommand{\blank}{\underline{\hspace{5em}}}
\newcommand{\pickout}{\hfill（\hspace{1.4em}）}
\newcommand{\circnum}[1]{\textcircled{\scriptsize #1}}
\setlist[enumerate,1]{label=\textup{(\arabic*)}}
\newenvironment{abcd}{\begin{enumerate}[label=\textup{(\Alph*)},leftmargin=2.3em,itemsep=1pt,topsep=4pt,parsep=1pt]}{\end{enumerate}}
\newcommand{\choices}[4]{%
  \begin{abcd}
    \item #1
    \item #2
    \item #3
    \item #4
  \end{abcd}}
\newcommand{\graphchoice}[2]{\begin{minipage}[t]{.235\linewidth}\centering\textup{(#1)}\par #2\end{minipage}}
\newcommand{\fourchoices}[4]{%
  \par\smallskip\noindent
  \graphchoice{A}{#1}\hfill\graphchoice{B}{#2}\hfill\graphchoice{C}{#3}\hfill\graphchoice{D}{#4}\par}
\newcommand{\widechoice}[2]{\begin{minipage}[t]{.48\linewidth}\raggedright\textup{(#1)}\enspace #2\end{minipage}}
\newcommand{\twochoices}[4]{%
  \par\medskip\noindent\widechoice{A}{#1}\hfill\widechoice{B}{#2}\par
  \medskip\noindent\widechoice{C}{#3}\hfill\widechoice{D}{#4}\par}
\newcounter{BankQuestion}
\ifdefined\QuestionBankAfour
\newcommand{\DiagramScale}{.85}
\newcommand{\BankFormat}{A4 纵向}
\newcommand{\BankSubtitle}{分节练习版}
\newcommand{\QuestionPage}[7]{%
  \stepcounter{BankQuestion}%
  \par\noindent\begin{minipage}{\linewidth}
  \raggedright
  {\small\textcolor{black!60}{#3\quad #5\quad 难度 #7}\hfill\textcolor{accent}{#2}}\par\smallskip}
\newcommand{\QuestionEnd}{\end{minipage}\par\medskip}
\else
\newcommand{\DiagramScale}{1}
\newcommand{\BankFormat}{4:3 横屏}
\newcommand{\BankSubtitle}{iPad 手写练习版}
\newcommand{\QuestionPage}[7]{%
  \clearpage
  \stepcounter{BankQuestion}%
  {\small\textcolor{accent}{#1\hfill #2}}\par\medskip
  {\small\textcolor{black!60}{#3\quad #5\quad 难度 #7\hfill \arabic{BankQuestion} / 234}}\par\vspace{9bp}}
\newcommand{\QuestionEnd}{}
\fi
\newenvironment{QuestionSolution}[1]{%
  \clearpage
  {\small\textcolor{black!60}{本题解析\quad #1}}\par\medskip
}{\par\clearpage}
\begin{document}
\hypertarget{contents}{}
{\small\textcolor{accent}{302\quad 全国硕士研究生招生考试\hfill \BankFormat}}\par\bigskip
{\huge\bfseries 数学二分节真题集}\par\smallskip
{\large \BankSubtitle}\par\smallskip
{\small 2009--2026 年\quad 234 题\quad 上册 19 节}\par\smallskip
{\small 点击目录进入对应小节。每道题的解析在下一页。}\par\bigskip
{\small 难度用 1--5 个 * 表示：* 直接应用，** 常规推算，*** 综合分析，**** 多步推理，***** 较难综合。各节由易到难，同难度保留原顺序。评级供学习安排参考。}\par\medskip

\begin{minipage}[t]{0.48\linewidth}
\medskip\textbf{函数、极限、连续}\par
{\small\hyperlink{sec-1.1}{1.1\quad 集合、映射与函数}\hfill 1 题}\par
{\small\hyperlink{sec-1.2}{1.2\quad 数列的极限}\hfill 10 题}\par
{\small\hyperlink{sec-1.3}{1.3\quad 函数的极限}\hfill 5 题}\par
{\small\hyperlink{sec-1.4}{1.4\quad 无穷小量与无穷大量}\hfill 23 题}\par
{\small\hyperlink{sec-1.5}{1.5\quad 连续函数}\hfill 8 题}\par
\medskip\textbf{一元函数微分学及其应用}\par
{\small\hyperlink{sec-2.1}{2.1\quad 导数的概念}\hfill 8 题}\par
{\small\hyperlink{sec-2.2}{2.2\quad 求导的基本法则}\hfill 28 题}\par
{\small\hyperlink{sec-2.3}{2.3\quad 微分}\hfill 0 题}\par
{\small\hyperlink{sec-2.4}{2.4\quad 微分中值定理及其应用}\hfill 12 题}\par
{\small\hyperlink{sec-2.5}{2.5\quad Taylor定理及其应用}\hfill 9 题}\par
{\small\hyperlink{sec-2.6}{2.6\quad 函数性态的研究}\hfill 25 题}\par
\end{minipage}%
\hfill
\begin{minipage}[t]{0.48\linewidth}
\medskip\textbf{一元函数积分学及其应用}\par
{\small\hyperlink{sec-3.1}{3.1\quad 定积分的概念、存在条件与性质}\hfill 11 题}\par
{\small\hyperlink{sec-3.2}{3.2\quad 微积分基本公式与基本定理}\hfill 11 题}\par
{\small\hyperlink{sec-3.3}{3.3\quad 两种基本积分法}\hfill 12 题}\par
{\small\hyperlink{sec-3.4}{3.4\quad 定积分的应用}\hfill 25 题}\par
{\small\hyperlink{sec-3.5}{3.5\quad 反常积分}\hfill 13 题}\par
\medskip\textbf{常微分方程}\par
{\small\hyperlink{sec-4.1}{4.1\quad 几类简单的微分方程}\hfill 17 题}\par
{\small\hyperlink{sec-4.2}{4.2\quad 高阶线性微分方程}\hfill 16 题}\par
{\small\hyperlink{sec-4.3}{4.3\quad 线性微分方程组（教材星号节）}\hfill 0 题}\par
\end{minipage}%

% 原卷取自本项目 302真题 所对应的 PDF。参校 https://lvjr.bitbucket.io/kysx.html
% 2020--2026 年参校 https://www.cnblogs.com/zhaokaifeng/category/2511758.html

\SectionHeading{1.1}{函数、极限、连续}{集合、映射与函数}{定理与要点\quad 教材 PDF 第 24 页起\hfill 本节 1 题}
\pdfbookmark[1]{函数、极限、连续}{chapter-1.1}
\pdfbookmark[2]{1.1 集合、映射与函数}{section-1.1}
\Theory{定理 1.1}{29}{确界存在定理}{非空实数集若有上界，就有上确界。若有下界，就有下确界。}
\Theory{定理 1.2}{36}{逆映射存在定理}{双射 \(f\colon A\to B\) 有唯一的逆映射 \(f^{-1}\colon B\to A\)。}
\Theory{要点}{33--36}{函数性质与复合}{奇偶性先核对定义域关于原点对称。求复合函数时，同时限制内层的定义域和外层的输入。}
% Q 2024 3
\QuestionPage{1.1 集合、映射与函数}{选择题}{2024 年第 3 题}{1}{函数的奇偶性}{001 / 234}{**}
已知
\[f(x)=\int_0^{\sin x}\sin t^3\,\mathrm{d}t,\qquad
g(x)=\int_0^x f(t)\,\mathrm{d}t,
\]
则（ ）
\choices{\(f(x)\) 为奇函数，\(g(x)\) 为奇函数}{\(f(x)\) 为奇函数，\(g(x)\) 为偶函数}{\(f(x)\) 为偶函数，\(g(x)\) 为偶函数}{\(f(x)\) 为偶函数，\(g(x)\) 为奇函数}
\QuestionEnd

\begin{QuestionSolution}{2024 年第 3 题}
\textbf{已知与目标}\quad \(f\) 是以 \(\sin x\) 为上限的积分，\(g\) 又是 \(f\) 的积分。要分别判断两者在自变量变号时是否保持原值或改变符号。

\textbf{方法与条件}\quad 偶函数满足 \(h(-x)=h(x)\)，奇函数满足 \(h(-x)=-h(x)\)，两者都要求定义域关于原点对称。这里被积函数连续，\(f,g\) 在全体实数上有定义，可以用积分换元比较。

\textbf{推导}\quad 设 \(F(u)=\int_0^u\sin(t^3)\,\mathrm{d}t\)。令 \(t=-s\)，则 \(\mathrm{d}t=-\mathrm{d}s\)，上下限由 \(0,-u\) 变为 \(0,u\)，因此
\[
F(-u)=-\int_0^u\sin((-s)^3)\,\mathrm{d}s
=\int_0^u\sin(s^3)\,\mathrm{d}s=F(u).
\]
用到 \(\sin(-s^3)=-\sin(s^3)\)。于是 \(f(-x)=F(-\sin x)=F(\sin x)=f(x)\)。对 \(g\) 再作同样换元：
\[
g(-x)=-\int_0^x f(-s)\,\mathrm{d}s=-\int_0^x f(s)\,\mathrm{d}s=-g(x).
\]
\textbf{结论与检查}\quad \(f\) 为偶函数，\(g\) 为奇函数，选 \(\boxed{\mathrm D}\)。其中 \(g(0)=0\)，与奇函数在原点的必要条件一致。
\end{QuestionSolution}

\SectionHeading{1.2}{函数、极限、连续}{数列的极限}{定理与要点\quad 教材 PDF 第 41 页起\hfill 本节 10 题}
\pdfbookmark[2]{1.2 数列的极限}{section-1.2}
\Theory{定理 2.1}{46}{唯一性}{收敛数列只有一个极限。}
\Theory{定理 2.2}{47}{有界性}{收敛数列必有界。}
\Theory{定理 2.3}{47}{有理运算法则}{两收敛数列可逐项作和、差、积、商再取极限，取商时分母极限须非零。}
\Theory{定理 2.4}{49}{保号性}{数列收敛到非零值时，从某一项起与极限同号，并且绝对值离零有正的下界。}
\Theory{定理 2.5}{49}{夹逼性}{两端数列收敛到同一值时，最终被夹在其间的数列也收敛到该值。}
\Theory{定理 2.6}{51}{单调有界准则}{单调增且有上界，或单调减且有下界的数列必收敛。}
\Theory{定理 2.7}{54}{子列与极限}{数列收敛于 \(a\) 的充要条件是它的每个子列都收敛于 \(a\)。}
\Theory{定理 2.8}{54}{闭区间套定理}{逐级包含且长度趋于零的非空闭区间列，恰有一个公共点。}
\Theory{定理 2.9}{55}{Weierstrass 定理}{有界实数列至少有一个收敛子列。}
\Theory{定理 2.10}{56}{Cauchy 收敛原理}{实数列收敛，当且仅当任给正误差，充分靠后的任意两项之差的绝对值都小于该误差。}
% Q 2012 3
\QuestionPage{1.2 数列的极限}{选择题}{2012 年第 3 题}{1}{数列与部分和的收敛}{002 / 234}{**}
设$a_n > 0$($n = 1,2, \cdots$)，$S_n = a_1 + a_2 + \cdots + a_n$，
则数列$\{s_n\}$有界是数列$\{a_n\}$收敛的\pickout{}
\begin{abcd}
\item 充分必要条件．
\item 充分非必要条件．
\item 必要非充分条件．
\item 即非充分也非必要条件．
\end{abcd}
\QuestionEnd

\begin{QuestionSolution}{2012 年第 3 题}
\textbf{已知与目标}\quad \(a_n>0\)，\(S_n\) 是前 \(n\) 项的和，题中的 \(s_n\) 按已定义的部分和理解。要判断“部分和有界”与“通项数列收敛”的两个蕴含方向。

\textbf{方法与条件}\quad 充分条件表示前者能够推出后者；必要条件表示后者能够推出前者。正项保证部分和递增，递增且有上界的实数列必收敛。

\textbf{推导}\quad 若 \(\{S_n\}\) 有界，由 \(S_{n+1}-S_n=a_{n+1}>0\) 可设 \(S_n\to S\)。相邻两部分和之差为 \(a_n\)，所以
\[
\lim_{n\to\infty}a_n=\lim_{n\to\infty}(S_n-S_{n-1})=S-S=0.
\]
这证明了充分性。检验必要性时取 \(a_n=1\)：每项都为正，且数列收敛于 \(1\)，部分和却为 \(S_n=n\)，没有上界。

\textbf{结论与检查}\quad 前者能推出后者，反向不成立，故为充分非必要条件，选 \(\boxed{\mathrm B}\)。反例只需满足“通项收敛”，无须额外要求通项趋于零。
\end{QuestionSolution}

% Q 2017 3
\QuestionPage{1.2 数列的极限}{选择题}{2017 年第 3 题}{1}{收敛数列的性质}{003 / 234}{**}
设数列$\left\{x_n \right\}$收敛，则\pickout{}
\begin{abcd}
\item 当$\lim_{n \to \infty} \sin x_n = 0$时，$\lim_{n \to \infty} x_n = 0$．
\item 当$\lim_{n \to \infty} (x_n + \sqrt{\left| x_n \right|}) = 0$时，$\lim_{n \to \infty} x_n = 0$．
\item 当$\lim_{n \to \infty} (x_n + x_n^2) = 0$时，$\lim_{n \to \infty} x_n = 0$．
\item 当$\lim_{n \to \infty} (x_n + \sin x_n) = 0$时，$\lim_{n \to \infty} x_n = 0$．
\end{abcd}
\QuestionEnd

\begin{QuestionSolution}{2017 年第 3 题}
\textbf{已知与目标}\quad 已知 \(x_n\to L\)，但 \(L\) 未知。要找出哪一个变换后的极限为零能够强制 \(L=0\)。

\textbf{方法与条件}\quad 各选项中的函数均连续，因此可将 \(x_n\) 替换为其极限 \(L\)。如果所得方程还有非零解，对应的常数数列就能作为反例。

\textbf{推导}\quad A 取 \(x_n=\pi\)，则 \(\sin x_n=0\)，但极限为 \(\pi\)。B、C 均取 \(x_n=-1\)，分别有 \(-1+\sqrt{|-1|}=0\)、\(-1+(-1)^2=0\)，但极限为 \(-1\)。

对 D，连续性给出 \(L+\sin L=0\)。当 \(L>0\) 时，\(0<L\le1\) 有 \(\sin L>0\)，而 \(L>1\) 有 \(L+\sin L\ge L-1>0\)。当 \(L<0\) 时，由 \(t+\sin t\) 为奇函数知其值为负。因此等于零只能发生在 \(L=0\)。

\textbf{结论与检查}\quad 只有 D 的前提能推出 \(x_n\to0\)，选 \(\boxed{\mathrm D}\)。三个反例均为收敛数列，符合题目的共同条件。
\end{QuestionSolution}

% Q 2022 6
\QuestionPage{1.2 数列的极限}{选择题}{2022 年第 6 题}{1}{递推数列的性质}{004 / 234}{**}
设数列 \(\left\{ x_{n} \right\}\) 满足 \(- \frac{\pi}{2} \le x_{n} \le \frac{\pi}{2}\)，则 \(\left( \ \right)\)
\choices{若 \(\lim_{n \to \infty} \cos\left( \sin x_{n} \right)\) 存在，则 \(\lim_{n \to \infty} x_{n}\) 存在}{若 \(\lim_{n \to \infty} \sin\left( \cos x_{n} \right)\) 存在，则 \(\lim_{n \to \infty} x_{n}\) 存在}{若 \(\lim_{n \to \infty} \cos\left( \sin x_{n} \right)\) 存在，则 \(\lim_{n \to \infty} \sin x_{n}\) 存在，但 \(\lim_{n \to \infty} x_{n}\) 不一定存在}{若 \(\lim_{n \to \infty} \sin\left( \cos x_{n} \right)\) 存在，则 \(\lim_{n \to \infty} \cos x_{n}\) 存在，但 \(\lim_{n \to \infty} x_{n}\) 不一定存在}
\QuestionEnd

\begin{QuestionSolution}{2022 年第 6 题}
\textbf{已知与目标}\quad 每个 \(x_n\) 都在 \([-\frac{\pi}{2},\frac{\pi}{2}]\) 中。题目要求辨认复合函数值收敛时，哪些内层数列一定收敛。

\textbf{方法与条件}\quad 反函数用于从函数值恢复自变量，必须限定在函数一一对应的区间。这里 \(\cos x_n\in[0,1]\)，正弦在 \([0,1]\) 严格递增，反函数 \(\arcsin\) 在对应值域连续。

\textbf{推导}\quad 若 \(b_n=\sin(\cos x_n)\to L\)，由于 \(0\le b_n\le\sin1\)，有 \(L\in[0,\sin1]\)。因此
\[
\cos x_n=\arcsin b_n\longrightarrow\arcsin L,
\]
证明 D 中前一结论必然成立。再取 \(x_n=\frac{(-1)^n\pi}{4}\)。其奇偶项分别为两个不同常数，故不收敛；\(\sin x_n\) 也在 \(-\frac{\sqrt2}{2}\) 与 \(\frac{\sqrt2}{2}\) 间交替。可是 \(\cos x_n=\frac{\sqrt2}{2}\)，并且 \(\cos(\sin x_n)=\cos(\frac{\sqrt2}{2})\)，均为常数。

\textbf{结论与检查}\quad 该反例同时排除 A、B、C，并说明 D 中 \(x_n\) 确实可以不收敛，选 \(\boxed{\mathrm D}\)。反例的所有项均满足给定区间限制。
\end{QuestionSolution}

% Q 2024 4
\QuestionPage{1.2 数列的极限}{选择题}{2024 年第 4 题}{1}{数列极限的敛散性}{006 / 234}{**}
已知数列 \(\{a_n\}\)（\(a_n\ne 0\)），若 \(\{a_n\}\) 发散，则（ ）
\choices{\(\left\{a_n+\frac{1}{a_n}\right\}\) 发散}{\(\left\{a_n-\frac{1}{a_n}\right\}\) 发散}{\(\left\{\mathrm{e}^{a_n}+\frac{1}{\mathrm{e}^{a_n}}\right\}\) 发散}{\(\left\{\mathrm{e}^{a_n}-\frac{1}{\mathrm{e}^{a_n}}\right\}\) 发散}
\QuestionEnd

\begin{QuestionSolution}{2024 年第 4 题}
\textbf{已知与目标}\quad \(a_n\ne0\) 且数列发散，即没有有限的实数极限。要找出一定保持发散的变换。

\textbf{方法与条件}\quad 不同输入若能产生同一输出，交替取这两个输入就可能使变换后的数列收敛。对能连续反解的变换，可用反证法：输出收敛将迫使输入收敛。

\textbf{推导}\quad A 取 \(a_n\) 在 \(2\) 和 \(\frac{1}{2}\) 间交替，输出恒为 \(2+\frac{1}{2}=\frac{5}{2}\)。B、C 取 \(a_n=(-1)^n\)，分别得到恒为零的 \(a_n-\frac{1}{a_n}\) 和恒为 \(\mathrm e+\mathrm e^{-1}\) 的 \(\mathrm e^{a_n}+\mathrm e^{-a_n}\)。这些输入均发散，因此三项都不能保证。

对 D，假设 \(b_n=\mathrm e^{a_n}-\mathrm e^{-a_n}\to b\)。令 \(z_n=\mathrm e^{a_n}>0\)，则 \(z_n^2-b_nz_n-1=0\)。舍去负根，得
\[
a_n=\ln\frac{b_n+\sqrt{b_n^2+4}}2.
\]
因为 \(\sqrt{b_n^2+4}>|b_n|\)，对数中的量恒正；该式关于 \(b_n\) 连续，所以 \(a_n\to\ln\frac{b+\sqrt{b^2+4}}{2}\)，与发散矛盾。

\textbf{结论与检查}\quad D 的变换一定发散，选 \(\boxed{\mathrm D}\)。反证覆盖交替振荡及趋于无穷等所有没有有限极限的情形。
\end{QuestionSolution}

% Q 2012 10
\QuestionPage{1.2 数列的极限}{填空题}{2012 年第 10 题}{2}{和式型数列极限与夹逼准则}{007 / 234}{**}
计算$\lim_{n \to \infty} n\left(\frac{1}{1 + n^2} + \frac{1}{2^2 + n^2}
+ \cdots + \frac{1}{n^2 + n^2} \right) =$ \blank{}．
\QuestionEnd

\begin{QuestionSolution}{2012 年第 10 题}
\textbf{已知与目标}\quad 求含 \(n\) 项的和乘以 \(n\) 后的极限。项数和每项分母都随 \(n\) 改变，不能逐项取极限后相加。

\textbf{方法与条件}\quad Riemann 和是“函数值乘小区间长度再求和”。连续函数在闭区间上的等分和趋于定积分，因此先把式子整理成区间 \([0,1]\) 上的等分和。

\textbf{推导}\quad 每项分母提出 \(n^2\)，外面的 \(n\) 与 \(n^2\) 约去一个因子，得到
\[
n\sum_{k=1}^n\frac{1}{n^2+k^2}
=\frac1n\sum_{k=1}^n\frac{1}{1+(\frac{k}{n})^2}.
\]
这里小区间长为 \(\frac{1}{n}\)，取样点是 \(t_k=\frac{k}{n}\)，函数 \(\frac{1}{1+t^2}\) 在 \([0,1]\) 上连续。因此极限为
\[
\int_0^1\frac{1}{1+t^2}\,\mathrm{d}t
=[\arctan t]_0^1=\frac{\pi}{4}-0=\frac{\pi}{4}.
\]

\textbf{结论与检查}\quad 所求值为 \(\boxed{\frac{\pi}{4}}\)。各取样函数值介于 \(\frac{1}{2}\) 与 \(1\) 之间，平均值也在该范围，答案与此相符。
\end{QuestionSolution}

% Q 2023 3
\QuestionPage{1.2 数列的极限}{选择题}{2023 年第 3 题}{1}{递推数列与无穷小比较}{005 / 234}{***}
已知数列 \(\{x_n\}\)，\(\{y_n\}\) 满足
\[x_1=y_1=\frac12,\qquad x_{n+1}=\sin x_n,\qquad y_{n+1}=y_n^2\quad (n=1,2,\cdots),
\]
则当 \(n\to\infty\) 时，（ ）。
\choices{\(x_n\) 是 \(y_n\) 的高阶无穷小}{\(y_n\) 是 \(x_n\) 的高阶无穷小}{\(x_n\) 与 \(y_n\) 是等价无穷小}{\(x_n\) 与 \(y_n\) 是同阶但不等价的无穷小}
\QuestionEnd

\begin{QuestionSolution}{2023 年第 3 题}
\textbf{已知与目标}\quad 两个数列从 \(\frac{1}{2}\) 开始，分别反复取正弦和平方。要先确认两者都趋于零，再比较它们趋于零的速度。

\textbf{方法与条件}\quad 若 \(\frac{y_n}{x_n}\to0\)，称 \(y_n\) 为 \(x_n\) 的高阶无穷小。直接求 \(x_n\) 的通项较困难，可以用下界控制这个比值。

\textbf{推导}\quad 对 \(0<t\le\frac{1}{2}\)，由 \(\sin t=\int_0^t\cos u\,\mathrm{d}u\) 以及 \(\frac{1}{2}<\cos u\le1\)，得 \(\frac{t}{2}<\sin t<t\)。归纳可知 \(0<x_{n+1}<x_n\le\frac{1}{2}\)。故 \(x_n\) 递减且有下界，极限 \(L\in[0,\frac{1}{2}]\) 存在。递推式取极限得到 \(L=\sin L\)，而正数 \(L\le\frac{1}{2}\) 满足 \(\sin L<L\)，所以 \(L=0\)。

由 \(x_{n+1}>\frac{x_n}{2}\) 和 \(x_1=\frac{1}{2}\)，逐步得到 \(x_n\ge2^{-n}\)。另一边，连续平方给出 \(y_n=2^{-2^{n-1}}\to0\)。于是
\[
0<\frac{y_n}{x_n}\le 2^{\,n-2^{n-1}}\longrightarrow0.
\]
上界趋于零，因为 \(2^{n-1}-n\to+\infty\)。夹逼后得到 \(\frac{y_n}{x_n}\to0\)。

\textbf{结论与检查}\quad \(y_n\) 是 \(x_n\) 的高阶无穷小，选 \(\boxed{\mathrm B}\)。两数列始终为正，比较时分母不会为零。
\end{QuestionSolution}

% Q 2011 19
\QuestionPage{1.2 数列的极限}{解答题}{2011 年第 19 题}{3}{调和和式与数列收敛}{008 / 234}{***}
\begin{enumerate}
\item
证明：对任意的正整数$n$，都有$\frac{1}{n + 1} < \ln \left(1 + \frac{1}{n} \right) < \frac{1}{n}$成立．
\item
设$a_n = 1 + \frac{1}{2} + \cdots + \frac{1}{n} - \ln n$（$n = 1,2, \cdots$），
证明数列$\left\{a_n \right\}$收敛．
\end{enumerate}
\QuestionEnd

\begin{QuestionSolution}{2011 年第 19 题}
\textbf{已知与目标}\quad 第一问证明对数增量的双边估计；第二问利用它证明 \(a_n\) 收敛，不要求求出极限的具体数值。

\textbf{方法与条件}\quad \(\ln(n+1)-\ln n\) 是 \(\frac{1}{t}\) 在 \([n,n+1]\) 上的积分。先估计积分，再证明 \(a_n\) 单调且有下界，即可应用单调有界准则。

\textbf{推导}\quad （1）由于 \(n\ge1\)，函数 \(\frac{1}{t}\) 在 \([n,n+1]\) 连续且严格递减。在内部各点都有 \(\frac{1}{n+1}<\frac{1}{t}<\frac{1}{n}\)，区间长度为 \(1\)，故积分满足严格不等式
\[
\frac1{n+1}<\int_n^{n+1}\frac{\mathrm{d}t}{t}
=\ln(1+\frac{1}{n})<\frac1n.
\]
（2）相邻两项相减，前 \(n\) 项的调和和抵消，得到
\[
a_{n+1}-a_n=\frac{1}{n+1}-\ln\left(1+\frac{1}{n}\right)<0,
\]
这里使用第一问的左边估计。再将右边估计对 \(k=1,\ldots,n-1\) 求和，对数可合并为 \(\ln2+\ln\frac{3}{2}+\cdots+\ln\frac{n}{n-1}=\ln n\)，于是 \(\ln n\le\sum_{k=1}^{n-1}\frac{1}{k}\)。因此 \(a_n\ge\frac{1}{n}>0\)；当 \(n=1\) 时直接有 \(a_1=1\)，结论同样成立。

\textbf{结论与检查}\quad 第一问的两个严格不等式均成立；第二问的数列递减且以 \(0\) 为下界，故收敛。仅证明相邻差趋零不足以证明收敛，这里已经另外给出了下界。
\end{QuestionSolution}

% Q 2013 20
\QuestionPage{1.2 数列的极限}{解答题}{2013 年第 20 题}{3}{不等式约束下的数列极限}{010 / 234}{***}
设函数$f(x) = \ln x + \frac{1}{x}$．
\begin{enumerate}
\item
求$f(x)$的最小值；
\item
设数列$\{x_n\}$满足$\ln x_n + \frac{1}{x_{n + 1}} < 1$，证明$\lim_{n \to \infty} x_n$存在，并求此极限．
\end{enumerate}
\QuestionEnd

\begin{QuestionSolution}{2013 年第 20 题}
\textbf{已知与目标}\quad 第一问求 \(\ln x+\frac{1}{x}\) 的全局最小值；第二问中相邻两项满足不等式，要由此推出收敛并确定极限。

\textbf{方法与条件}\quad 先利用导数符号求出对所有正数有效的下界，再与题设不等式比较。对数有定义保证各项为正，因而比较倒数时可以改变不等号方向。

\textbf{推导}\quad （1）定义域为 \(x>0\)，
\[
f'(x)=\frac{1}{x}-\frac{1}{x^2}=\frac{x-1}{x^2}.
\]
分母恒正，因此 \(0<x<1\) 时导数为负，\(x>1\) 时为正。函数先减后增，故全局最小值为 \(f(1)=\ln1+1=1\)，且仅在 \(x=1\) 取到。

（2）每个 \(x_n>0\)。将第一问代入 \(x_n\)，并结合题设，有
\[
\ln x_n+\frac1{x_{n+1}}<1\le\ln x_n+\frac1{x_n}
\]
消去 \(\ln x_n\)，得 \(\frac{1}{x_{n+1}}<\frac{1}{x_n}\)，所以 \(x_{n+1}>x_n\)。原不等式中倒数为正，故 \(\ln x_n<1\)，即 \(x_n<\mathrm e\)。数列递增且有上界，因而 \(x_n\to L\)，并且 \(L\ge x_1>0\)。这保证 \(\ln L\) 与 \(\frac{1}{L}\) 都有意义。取极限时严格不等式只能保留为
\[
\ln L+\frac{1}{L}\le1.
\]
第一问又给出左端至少为 \(1\)，等号只能在 \(L=1\) 取得。

\textbf{结论与检查}\quad 最小值为 \(\boxed{1}\)，数列极限为 \(\boxed{1}\)。取极限前已证明收敛且极限为正，使用连续性和倒数运算的条件齐备。
\end{QuestionSolution}

% Q 2012 21
\QuestionPage{1.2 数列的极限}{解答题}{2012 年第 21 题}{4}{实根构成的数列极限}{009 / 234}{****}
\begin{enumerate}
\item
证明方程$x^n + x^{n - 1} + \cdots + x = 1$（$n$为大于$1$的整数）
在区间$\left(\frac{1}{2},1 \right)$内有且仅有一个实根；
\item
记第（1）问中的实根为$x_n$，证明$\lim_{n \to \infty} x_n$存在，并求此极限．
\end{enumerate}
\QuestionEnd

\begin{QuestionSolution}{2012 年第 21 题}
\textbf{已知与目标}\quad 对每个整数 \(n>1\)，先证明方程在 \((\frac{1}{2},1)\) 内恰有一根，再证明这些根组成的数列收敛并求极限。

\textbf{方法与条件}\quad 零点存在定理要求函数在闭区间上连续且端点值异号；严格单调可保证根唯一。比较相邻方程在同一点的值，可以判断根的大小，继而使用单调有界准则。

\textbf{推导}\quad （1）记 \(P_n(t)=\sum_{k=1}^n t^k\)。多项式连续，且在正半轴上 \(P_n'(t)=\sum_{k=1}^n kt^{k-1}>0\)，故严格递增。端点代入得
\[
P_n\left(\frac{1}{2}\right)=\frac{\frac{1}{2}(1-2^{-n})}{1-\frac{1}{2}}=1-2^{-n}<1,
\qquad P_n(1)=n>1.
\]
因此 \(P_n(t)-1\) 在两端异号，内部有零点，且严格递增使零点至多一个，记为 \(x_n\)。

（2）将 \(x_n\) 代入下一方程，得到 \(P_{n+1}(x_n)=P_n(x_n)+x_n^{n+1}=1+x_n^{n+1}>1\)。而 \(P_{n+1}(x_{n+1})=1\)，由严格递增得 \(x_{n+1}<x_n\)。又 \(x_n>\frac{1}{2}\)，故极限存在。

由 \(x_n\ne1\) 可用等比求和：\(\frac{x_n(1-x_n^n)}{1-x_n}=1\)，整理为 \(2x_n-1=x_n^{n+1}\)。不能仅由 \(x_n<1\) 就判定右侧趋零，须使用统一上界：当 \(n\ge2\) 时 \(0<x_n\le x_2<1\)，所以 \(0<x_n^{n+1}\le x_2^{n+1}\to0\)。

\textbf{结论与检查}\quad 对每个 \(n>1\) 均有唯一根，且 \(\boxed{\lim_{n\to\infty}x_n=\frac{1}{2}}\)。每项虽严格大于 \(\frac{1}{2}\)，极限仍可等于这个区间端点。
\end{QuestionSolution}

% Q 2018 21
\QuestionPage{1.2 数列的极限}{解答题}{2018 年第 21 题}{4}{递推数列的收敛与极限}{011 / 234}{****}
设数列$\left\{x_n \right\}$满足：$x_1 > 0$，$x_n \mathrm{e}^{x_{n + 1}} = \mathrm{e}^{x_n} - 1$（$n = 1,2, \cdots$）．
证明$\left\{x_n \right\}$收敛，并求$\lim_{n \to \infty} x_n$．
\QuestionEnd

\begin{QuestionSolution}{2018 年第 21 题}
\textbf{已知与目标}\quad 仅已知 \(x_1>0\) 及相邻项的指数关系。需要证明所有后续项保持正值、数列确实收敛，再求极限。

\textbf{方法与条件}\quad 对正数 \(x\)，\(\frac{\mathrm e^x-1}{x}\) 是 \(\mathrm e^t\) 在 \([0,x]\) 上的平均值，即积分除以区间长度。它严格介于端点函数值之间，可以同时给出正性和递减性。

\textbf{推导}\quad 假设 \(x_n>0\)。因为 \(1<\mathrm e^t<\mathrm e^{x_n}\) 对 \(0<t<x_n\) 成立，积分后除以正数 \(x_n\)，得到
\[
1<\mathrm{e}^{x_{n+1}}=\frac1{x_n}\int_0^{x_n}\mathrm{e}^t\,\mathrm{d}t<\mathrm{e}^{x_n}.
\]
指数函数严格递增，比较指数便得 \(0<x_{n+1}<x_n\)。从 \(x_1>0\) 出发归纳，所有项都为正并严格递减，因此以 \(0\) 为下界，必收敛。设极限为 \(L\ge0\)。

若 \(L>0\)，分母的极限非零，可在递推式中取极限，得到
\[
\mathrm e^L=\frac{\mathrm e^L-1}{L}
=\frac{1}{L}\int_0^L\mathrm e^t\,\mathrm{d}t<\mathrm e^L,
\]
最后的严格不等式仍来自正长度区间上的平均值估计，形成矛盾。故 \(L>0\) 不可能。

\textbf{结论与检查}\quad 数列收敛且 \(\boxed{\lim_{n\to\infty}x_n=0}\)。将 \(L=0\) 代入原递推式的极限关系 \(L\mathrm e^L=\mathrm e^L-1\)，两边均为零，结果相容。
\end{QuestionSolution}

\SectionHeading{1.3}{函数、极限、连续}{函数的极限}{定理与要点\quad 教材 PDF 第 60 页起\hfill 本节 5 题}
\pdfbookmark[2]{1.3 函数的极限}{section-1.3}
\Theory{定理 3.1}{65}{Heine 定理}{当 \(x\to x_0\) 时，函数 \(f\) 的极限为 \(a\) 的充要条件是，对定义域内任意满足 \(x_n\ne x_0\) 且 \(x_n\to x_0\) 的数列，都有 \(f(x_n)\to a\)。}
\Theory{定理 3.2}{66}{唯一性与局部有界性}{函数的有限极限唯一，且函数在相应去心邻域内有界。}
\Theory{定理 3.3}{67}{保号、保序和夹逼}{非零极限在局部保号。局部大小关系传给极限，两端极限相同即可夹逼。}
\Theory{定理 3.4}{67}{有理运算法则}{和、差、积的极限由各自极限运算得到。商的分母极限须非零。}
\Theory{定理 3.5}{69}{复合函数的极限}{内层函数趋于 \(u_0\) 且在去心邻域内避开 \(u_0\) 时，可应用外层函数在 \(u_0\) 处的极限。}
\Theory{定理 3.6}{73}{单调有界准则}{在右半无限区间单调增且有上界或单调减且有下界时，右端极限存在。单调函数在区间内点有左右极限。}
\Theory{定理 3.7}{74}{Cauchy 收敛原理}{有限极限存在，当且仅当充分接近目标点的任意两处函数值都可彼此任意接近。}
\Theory{公式}{70--72}{两个重要极限}{\(\lim_{x\to0}\frac{\sin x}{x}=1\)，\(\lim_{u\to0}(1+u)^{\frac{1}{u}}=\mathrm{e}\)。}
% Q 2011 9
\QuestionPage{1.3 函数的极限}{填空题}{2011 年第 9 题}{1}{函数极限}{012 / 234}{**}
$\lim_{x \to 0} \left(\frac{1 + 2^x}{2} \right)^{\frac{1}{x}} =$ \blank{}．
\QuestionEnd

\begin{QuestionSolution}{2011 年第 9 题}
\textbf{已知与目标}\quad 底数趋于 \(1\)，指数 \(\frac{1}{x}\) 无界，不能直接代入确定幂的极限。

\textbf{方法与条件}\quad 正底数的变指数幂可写成指数函数：\(A(x)^{B(x)}=\mathrm e^{B(x)\ln A(x)}\)。这里底数始终为正。取对数后用 \(\frac{\ln(1+u)}{u}\to1\) 和 \(\frac{\mathrm e^v-1}{v}\to1\) 两个基本极限。

\textbf{推导}\quad 令 \(u=\frac{2^x-1}{2}\)，则 \(u\to0\)，底数为 \(1+u\)。原幂的对数满足
\[
\frac{\ln(1+u)}{x}
=\frac{\ln(1+u)}u\frac{2^x-1}{2x}
\longrightarrow 1\cdot\frac{\ln 2}{2}.
\]
其中 \(\frac{2^x-1}{x}=\ln2\cdot\frac{\mathrm e^{x\ln2}-1}{x\ln2}\to\ln2\)。

\textbf{结论与检查}\quad 由指数函数连续，原极限为 \(\boxed{\mathrm e^{\frac{\ln2}{2}}=\sqrt2}\)。上述基本极限对左右两侧均成立，所以得到的是双侧极限。
\end{QuestionSolution}

% Q 2013 9
\QuestionPage{1.3 函数的极限}{填空题}{2013 年第 9 题}{2}{指数型函数极限}{013 / 234}{**}
$\lim_{x \to 0} \left(2 - \frac{\ln(1 + x)}{x} \right)^{\frac{1}{x}} =$ \blank{}．
\QuestionEnd

\begin{QuestionSolution}{2013 年第 9 题}
\textbf{已知与目标}\quad 底数 \(2-\frac{\ln(1+x)}{x}\to1\)，需要求其偏离 \(1\) 的速度，才能处理指数 \(\frac{1}{x}\)。

\textbf{方法与条件}\quad Taylor 展开用低次幂表示函数在一点附近的变化。记号 \(o(x^2)\) 表示余项除以 \(x^2\) 趋于零。分子将除以 \(x\)，所以对数必须展开到二次。

\textbf{推导}\quad 在 \(x=0\) 附近，\(\ln(1+x)=x-\frac{x^2}{2}+o(x^2)\)。除以 \(x\ne0\) 并相减，得
\[
2-\frac{\ln(1+x)}{x}=2-\left(1-\frac{x}{2}+o(x)\right)
=1+\frac{x}{2}+o(x).
\]
底数趋于正数 \(1\)，在足够小的去心邻域中为正，因而可取对数。利用 \(\ln(1+v)\sim v\)，原幂的对数为
\[
\frac{1}{x}\ln\bigl(1+\frac{x}{2}+o(x)\bigr)\longrightarrow\frac{1}{2}.
\]

\textbf{结论与检查}\quad 原极限为 \(\boxed{\mathrm e^{\frac{1}{2}}=\sqrt{\mathrm e}}\)。对数的二次项除以 \(x\) 后产生底数的一次修正，符号为正。
\end{QuestionSolution}

% Q 2019 9
\QuestionPage{1.3 函数的极限}{填空题}{2019 年第 9 题}{1}{指数型函数极限}{014 / 234}{**}
$\lim_{x \to 0} \left(x + 2^x \right)^{\frac{2}{x}} =$ \blank{}．
\QuestionEnd

\begin{QuestionSolution}{2019 年第 9 题}
\textbf{已知与目标}\quad 求 \((x+2^x)^{\frac{2}{x}}\) 的极限。底数趋于 \(1\)，需要同时计入 \(x\) 和 \(2^x-1\) 的变化。

\textbf{方法与条件}\quad 底数在零点附近为正，可以取对数；用 \(2^x=1+x\ln2+o(x)\)，其中 \(o(x)\) 表示除以 \(x\) 后趋于零的余项，再应用 \(\ln(1+u)\sim u\)。

\textbf{推导}\quad 将底数合并为 \(x+2^x=1+(1+\ln2)x+o(x)\)。令 \(u=(1+\ln2)x+o(x)\)，则 \(\frac{u}{x}\to1+\ln2\)，故对数极限为
\[
\frac{2}{x}\ln(x+2^x)=2\frac{\ln(1+u)}{u}\frac{u}{x}
\longrightarrow2(1+\ln2).
\]

\textbf{结论与检查}\quad 原极限为 \(\boxed{\mathrm e^{2(1+\ln2)}=4\mathrm e^2}\)，其中 \(\mathrm e^{2\ln2}=2^2=4\)。两项的一次变化均已计入对数系数。
\end{QuestionSolution}

% Q 2022 11
\QuestionPage{1.3 函数的极限}{填空题}{2022 年第 11 题}{2}{指数型函数极限}{015 / 234}{**}
\(\lim_{x \to 0} \left( \frac{1 + \mathrm{e}^{x}}{2} \right)^{\cot x} = \blank{}\)
\QuestionEnd

\begin{QuestionSolution}{2022 年第 11 题}
\textbf{已知与目标}\quad 底数 \(\frac{1+\mathrm e^x}{2}\to1\)，指数 \(\cot x\) 在零点附近无界，需要求乘上底数对数后的极限。

\textbf{方法与条件}\quad \(\cot x=\frac{\cos x}{\sin x}\)。用 \(\frac{\sin x}{x}\to1\) 得 \(x\cot x\to1\)，再将对数除以 \(x\) 单独求极限。底数始终为正，取对数合法。

\textbf{推导}\quad 由 \(\mathrm e^x=1+x+o(x)\)，有 \(\frac{1+\mathrm e^x}{2}=1+\frac{x}{2}+o(x)\)。这里 \(o(x)\) 的含义是余项除以 \(x\) 后趋零。于是
\[
\cot x\ln\frac{1+\mathrm{e}^x}{2}
=(x\cot x)\frac{\ln(1+\frac{x}{2}+o(x))}{x}\longrightarrow1\cdot\frac{1}{2}.
\]

\textbf{结论与检查}\quad 指数化得 \(\boxed{\mathrm e^{\frac{1}{2}}}\)。虽然 \(\cot x\) 在两侧符号不同，对数也随 \(x\) 变号，乘积的两侧极限相同。
\end{QuestionSolution}

% Q 2011 15
\QuestionPage{1.3 函数的极限}{解答题}{2011 年第 15 题}{2}{两处函数极限与参数}{016 / 234}{***}
已知函数$F(x) = \frac{\int_0^x \ln(1 + t^2)\,\mathrm{d} t}{x^a}$，
设$\!\!\lim_{x \to +\infty}\!\!\! F(x) = \!\lim_{x \to 0^+}\! F(x) = 0$，试求$a$的取值范围．
\QuestionEnd

\begin{QuestionSolution}{2011 年第 15 题}
\textbf{已知与目标}\quad 同一个 \(F(x)\) 在 \(x\to0^+\) 与 \(x\to+\infty\) 时都要趋于零，参数 \(a\) 必须同时满足两端的条件。

\textbf{方法与条件}\quad 先找出积分分子在两端的等价量，即比值趋于 \(1\) 的较简单函数。L'Hospital 法则可用于 \(\frac{0}{0}\) 或 \(\frac{\infty}{\infty}\) 型，要求附近可导、分母导数非零且导数比极限存在。

\textbf{推导}\quad 记分子为 \(I(x)\)。当 \(x\to0^+\) 时，\(I(x),x^3\) 都趋零；被积函数连续保证 \(I'(x)=\ln(1+x^2)\)，且 \(3x^2>0\)。因此
\[
\lim_{x\to0^+}\frac{\int_0^x\ln(1+t^2)\,\mathrm{d}t}{x^3}
=\lim_{x\to0^+}\frac{\ln(1+x^2)}{3x^2}=\frac{1}{3}.
\]
故 \(F(x)\sim\frac{1}{3}x^{3-a}\)。右极限为零恰要求 \(a<3\)；\(a=3\) 时极限为 \(\frac{1}{3}\)，\(a>3\) 时趋于正无穷。

当 \(x\to+\infty\) 时，对 \(t\ge1\) 有 \(\ln(1+t^2)\ge\ln2\)，故 \(I(x)\to+\infty\)。与 \(x\ln x\) 比较属于 \(\frac{\infty}{\infty}\) 型，分母导数 \(1+\ln x\) 最终为正，且
\[
\frac{I'(x)}{(x\ln x)'}=\frac{\ln(1+x^2)}{1+\ln x}
=\frac{2\ln x+\ln(1+x^{-2})}{1+\ln x}\longrightarrow2.
\]
所以 \(F(x)\sim2x^{1-a}\ln x\)。若 \(a>1\)，令 \(p=a-1>0\)，由 \(\frac{\ln x}{x^p}\to0\) 得极限为零；此基本比较也可由 L'Hospital 法则求得导数比 \(\frac{1}{p x^p}\to0\)。若 \(a\le1\)，该等价量趋于正无穷。

\textbf{结论与检查}\quad 两端条件取交集，得到 \(\boxed{1<a<3}\)。两个端点 \(a=1,3\) 都不能使两个极限同时为零。
\end{QuestionSolution}

\SectionHeading{1.4}{函数、极限、连续}{无穷小量与无穷大量}{定理与要点\quad 教材 PDF 第 77 页起\hfill 本节 23 题}
\pdfbookmark[2]{1.4 无穷小量与无穷大量}{section-1.4}
\Theory{定理 4.1}{78}{极限与无穷小}{函数 \(f\) 趋于有限值 \(a\) 的充要条件是 \(f-a\) 为无穷小量。}
\Theory{定理 4.2}{78}{无穷小的和与积}{在同一变化趋势下，有限个无穷小的代数和及乘积仍为无穷小。}
\Theory{定理 4.3}{78}{有界量乘无穷小}{局部有界函数乘无穷小，结果仍为无穷小。}
\Theory{定理 4.4}{81}{等价无穷小}{两非零无穷小等价，表示两者之差相对于其中任一量为高阶无穷小。}
\Theory{定理 4.5}{82}{无穷小等价代换}{求商的极限时，分子和分母的无穷小因子可分别换成各自的等价量。}
\Theory{定理 4.6}{84}{倒数关系}{非零无穷小的倒数为无穷大量，无穷大量的倒数为无穷小。}
\Theory{定理 4.7}{84}{无穷大量的运算}{有限个无穷大量的乘积仍是无穷大量。无穷大量与有界量的和仍是无穷大量。}
\Theory{公式}{82--85}{常用等价式与渐近线}{当 \(x\to0\) 时，\(\sin x\sim x\)，\(1-\cos x\sim\frac{x^2}{2}\)。求斜渐近线 \(y=kx+b\) 时，可依次求 \(\frac{f(x)}{x}\) 和 \(f(x)-kx\) 的极限。}
% Q 2011 1
\QuestionPage{1.4 无穷小量与无穷大量}{选择题}{2011 年第 1 题}{1}{三角函数的等价无穷小}{018 / 234}{*}
已知当$x \to 0$时，$f(x) = 3\sin x - \sin 3x$与$cx^k$是等价无穷小，则\pickout{}
\begin{abcd}
\item $k = 1,c = 4$．
\item $k = 1,c = - 4$．
\item $k = 3,c = 4$．
\item $k = 3,c = - 4$．
\end{abcd}
\QuestionEnd

\begin{QuestionSolution}{2011 年第 1 题}
\textbf{已知与目标}\quad 要把 \(3\sin x-\sin3x\) 写成 \(cx^k\) 的等价形式，即求使两者比值趋于 \(1\) 的系数和幂次。

\textbf{方法与条件}\quad 三倍角恒等式 \(\sin3x=3\sin x-4\sin^3x\) 能直接化简差式，再使用弧度制下的基本极限 \(\frac{\sin x}{x}\to1\)。

\textbf{推导}\quad 代入恒等式得 \(f(x)=4\sin^3x\)，所以
\[
\frac{f(x)}{x^3}=4\left(\frac{\sin x}{x}\right)^3\longrightarrow4.
\]
因此最低非零幂次为三次，系数为 \(4\)，也就是 \(f(x)\sim4x^3\)。

\textbf{结论与检查}\quad \(\boxed{k=3,\ c=4}\)，选 C。除以 \(4x^3\) 后极限为 \(1\)，符合等价的定义。
\end{QuestionSolution}

% Q 2019 1
\QuestionPage{1.4 无穷小量与无穷大量}{选择题}{2019 年第 1 题}{1}{正切函数的同阶无穷小}{024 / 234}{*}
当$x \to 0$时，若$x - \tan x$与$x^k$是同阶无穷小，则$k =$\pickout{}
\begin{abcd}
\item $1$．
\item $2$．
\item $3$．
\item $4$．
\end{abcd}
\QuestionEnd

\begin{QuestionSolution}{2019 年第 1 题}
\textbf{已知与目标}\quad 求使 \(x-\tan x\) 与 \(x^k\) 同阶的幂次 \(k\)。同阶指比值趋于有限非零常数，常数可以为负。

\textbf{方法与条件}\quad 对 \(\frac{x-\tan x}{x^3}\) 使用 L'Hospital 法则：分子、分母均趋零，在足够小的去心邻域中可导，分母导数 \(3x^2\ne0\)。

\textbf{推导}\quad 因为 \((\tan x)'=\sec^2x\) 且 \(1-\sec^2x=-\tan^2x\)，所以
\[
\lim_{x\to0}\frac{x-\tan x}{x^3}
=-\frac{1}{3}\lim_{x\to0}\left(\frac{\tan x}{x}\right)^2=-\frac{1}{3}.
\]
这里 \(\frac{\tan x}{x}=\frac{\sin x}{x}\frac{1}{\cos x}\to1\)，导数比的极限存在，法则适用。

\textbf{结论与检查}\quad \(\boxed{k=3}\)，选 C。\(x-\tan x\sim-\frac{x^3}{3}\)，所以它与 \(x^3\) 同阶，但二者并不等价。
\end{QuestionSolution}

% Q 2010 10
\QuestionPage{1.4 无穷小量与无穷大量}{填空题}{2010 年第 10 题}{2}{曲线的渐近线}{031 / 234}{*}
曲线$y = \frac{2x^3}{x^2 + 1}$的渐近线方程为 \blank{}．
\QuestionEnd

\begin{QuestionSolution}{2010 年第 10 题}
\textbf{已知与目标}\quad 求有理曲线在有限点或无穷远处的直线渐近线。

\textbf{方法与条件}\quad 若 \(y-(kx+b)\to0\)，则 \(y=kx+b\) 是无穷远处的渐近线。分子次数比分母高一次，直接作多项式除法即可提取线性部分。

\textbf{推导}\quad 由 \(2x^3=2x(x^2+1)-2x\)，得到
\[
y=\frac{2x^3}{x^2+1}=2x-\frac{2x}{x^2+1}.
\]
余项 \(-\frac{2x}{x^2+1}=-\frac{\frac{2}{x}}{1+\frac{1}{x^2}}\to0\)，故候选直线为 \(y=2x\)。又因 \(x^2+1>0\)，函数在每个有限点连续，附近有界，没有竖直渐近线。

\textbf{结论与检查}\quad 渐近线为 \(\boxed{y=2x}\)。曲线与直线的差在正负两个无穷远处都趋零，无常数截距修正。
\end{QuestionSolution}

% Q 2009 2
\QuestionPage{1.4 无穷小量与无穷大量}{选择题}{2009 年第 2 题}{1}{等价无穷小}{017 / 234}{**}
当$x \to 0$时，$f(x) = x - \sin ax$与$g(x) = x^2\ln (1 - bx)$是等价无穷小，则\pickout{}
\begin{abcd}
\item $a = 1,b = - \frac{1}{6}$．
\item $a = 1,b = \frac{1}{6}$．
\item $a = - 1,b = - \frac{1}{6}$．
\item $a = - 1,b = \frac{1}{6}$．
\end{abcd}
\QuestionEnd

\begin{QuestionSolution}{2009 年第 2 题}
\textbf{已知与目标}\quad 已知 \(f\) 与 \(g\) 等价，要求确定参数 \(a,b\)。等价的含义是 \(\frac{f(x)}{g(x)}\to1\)，因此二者最低非零幂次及其系数都必须相同。

\textbf{方法与条件}\quad 用 Taylor 展开找出最低非零项。\(o(x^3)\) 表示除以 \(x^3\) 后趋零的余项。若 \(b=0\)，则 \(g\) 恒为零，无法作为等价关系中的分母，故 \(b\ne0\)。固定 \(b\) 后，\(1-bx\) 在零点附近为正。

\textbf{推导}\quad 由 \(\ln(1+u)=u+o(u)\) 和 \(\sin u=u-\frac{u^3}{6}+o(u^3)\)，得
\[
g(x)=-bx^3+o(x^3),\qquad
f(x)=(1-a)x+\frac{a^3x^3}{6}+o(x^3).
\]
右侧 \(g\) 为三阶。若 \(1-a\ne0\)，\(f\) 含有一次项，与三阶量的比值不可能趋于 \(1\)，故 \(a=1\)。此时再比较三次系数，\(\frac{1}{6}=-b\)，得到 \(b=-\frac{1}{6}\)。

\textbf{结论与检查}\quad \(\boxed{a=1,\ b=-\frac{1}{6}}\)，选 A。代回后 \(f,g\) 都等价于 \(\frac{x^3}{6}\)，其比值确实趋于 \(1\)。
\end{QuestionSolution}

% Q 2012 1
\QuestionPage{1.4 无穷小量与无穷大量}{选择题}{2012 年第 1 题}{1}{渐近线条数}{019 / 234}{**}
曲线$y = \frac{x^2 + x}{x^2 - 1}$渐近线的条数为\pickout{}
\begin{abcd}
\item $0$．
\item $1$．
\item $2$．
\item $3$．
\end{abcd}
\QuestionEnd

\begin{QuestionSolution}{2012 年第 1 题}
\textbf{已知与目标}\quad 曲线在 \(x=\pm1\) 无定义，要计算全部直线渐近线的条数。

\textbf{方法与条件}\quad 某有限点至少一侧函数值趋于无穷时，该竖直线是渐近线；当 \(x\to\pm\infty\) 时函数趋于常数，该常数给出水平渐近线。分母为零的点需先检查能否约去因子。

\textbf{推导}\quad 保留原定义域限制 \(x\ne\pm1\)，有
\[
y=\frac{x(x+1)}{(x-1)(x+1)}=\frac{x}{x-1}=1+\frac{1}{x-1}.
\]
当 \(x\to-1\) 时，\(y\to1-\frac{1}{2}=\frac{1}{2}\)，没有无穷趋势；当 \(x\to1^\pm\) 时，\(\frac{1}{x-1}\to\pm\infty\)，故 \(x=1\) 是竖直渐近线。无穷远处 \(\frac{1}{x-1}\to0\)，所以 \(y=1\) 是水平渐近线。此时 \(\frac{y}{x}\to0\)，没有斜率非零的直线渐近线。

\textbf{结论与检查}\quad 共有 \(\boxed{2}\) 条：\(x=1\)、\(y=1\)，选 C。\(x=-1\) 处只是可以补值的缺口，不增加渐近线。
\end{QuestionSolution}

% Q 2013 1
\QuestionPage{1.4 无穷小量与无穷大量}{选择题}{2013 年第 1 题}{1}{三角无穷小的比较}{020 / 234}{**}
设$\cos x - 1 = x\sin \alpha(x)$，其中$\left| \alpha(x) \right| < \frac{\pi}{2}$，
则当$x \to 0$时，$\alpha (x)$是\pickout{}
\begin{abcd}
\item 比$x$高阶的无穷小．
\item 比$x$低阶的无穷小．
\item 与$x$同阶但不等价的无穷小．
\item 与$x$等价的无穷小．
\end{abcd}
\QuestionEnd

\begin{QuestionSolution}{2013 年第 1 题}
\textbf{已知与目标}\quad 等式给出 \(\sin\alpha(x)\)，而问题比较的是角度 \(\alpha(x)\) 与 \(x\)。需要先证明角度本身趋零。

\textbf{方法与条件}\quad 正弦在 \((-\frac{\pi}{2},\frac{\pi}{2})\) 内可由连续反函数 \(\arcsin\) 恢复角度。两无穷小的比值趋于有限非零常数称为同阶，其中常数等于 \(1\) 才称等价。

\textbf{推导}\quad 对 \(x\ne0\)，由题设和 \(1-\cos x\sim\frac{x^2}{2}\)，得 \(\sin\alpha(x)=\frac{\cos x-1}{x}\to0\)。由给定角度范围，\(\alpha(x)=\arcsin(\sin\alpha(x))\to\arcsin0=0\)，因而可使用小角度极限。于是
\[
\frac{\alpha(x)}{x}
=\frac{\alpha(x)}{\sin\alpha(x)}\frac{\cos x-1}{x^2}
\longrightarrow1\cdot\left(-\frac{1}{2}\right)=-\frac{1}{2}.
\]

\textbf{结论与检查}\quad 比值有限非零且不为 \(1\)，所以同阶但不等价，选 \(\boxed{\mathrm C}\)。负号与 \(x>0\) 时 \(\cos x-1<0\)、\(\alpha(x)<0\) 一致。
\end{QuestionSolution}

% Q 2014 1
\QuestionPage{1.4 无穷小量与无穷大量}{选择题}{2014 年第 1 题}{1}{幂次与无穷小阶}{021 / 234}{**}
当$x \to 0^+$时，若$\ln^{\alpha}(1 + 2x)$，$(1 - \cos x)^{\frac{1}{\alpha}}$
均是比$x$高阶的无穷小，则$\alpha$的可能取值范围是\pickout{}
\begin{abcd}
\item $(2, +\infty)$．
\item $(1,2)$．
\item $\left(\frac{1}{2},1 \right)$．
\item $\left(0,\frac{1}{2} \right)$．
\end{abcd}
\QuestionEnd

\begin{QuestionSolution}{2014 年第 1 题}
\textbf{已知与目标}\quad 两个量都必须比 \(x\) 更快趋于零，即分别除以 \(x\) 后都趋于零。求同时满足这两条要求的 \(\alpha\)。

\textbf{方法与条件}\quad 当 \(x\to0^+\) 时，\(x^p\to0\) 当且仅当 \(p>0\)。这里 \(\ln(1+2x)>0\)、\(1-\cos x>0\)，固定实数幂的运算合法；第二个式子还要求 \(\alpha\ne0\)。

\textbf{推导}\quad 由 \(\ln(1+2x)\sim2x\) 及 \(1-\cos x\sim\frac{x^2}{2}\)，得
\[
\frac{\ln^\alpha(1+2x)}x\sim 2^\alpha x^{\alpha-1},
\qquad
\frac{(1-\cos x)^{\frac{1}{\alpha}}}x\sim 2^{-\frac{1}{\alpha}}x^{\frac{2}{\alpha}-1}.
\]
第一个比值趋零要求 \(\alpha-1>0\)，即 \(\alpha>1\)。在此条件下 \(\alpha\) 为正，第二个比值趋零要求 \(\frac{2}{\alpha}-1>0\)，两边乘以 \(\alpha\) 得 \(\alpha<2\)。

\textbf{结论与检查}\quad 交集为 \(\boxed{1<\alpha<2}\)，选 B。\(\alpha=1\) 时第一个比值趋于 \(2\)，\(\alpha=2\) 时第二个比值趋于 \(\frac{1}{\sqrt2}\)，端点都需排除。
\end{QuestionSolution}

% Q 2014 2
\QuestionPage{1.4 无穷小量与无穷大量}{选择题}{2014 年第 2 题}{1}{渐近线的判定}{022 / 234}{**}
下列曲线有渐近线的是\pickout{}
\begin{abcd}
\item $y = x + \sin x$．
\item $y = x^2 + \sin x$．
\item $y = x + \sin \frac{1}{x}$．
\item $y = x^2 + \sin \frac{1}{x}$．
\end{abcd}
\QuestionEnd

\begin{QuestionSolution}{2014 年第 2 题}
\textbf{已知与目标}\quad 四个函数都由多项式加上有界的正弦项组成，要判断哪条曲线接近某条直线。

\textbf{方法与条件}\quad 无穷远处若有直线渐近线 \(y=kx+b\)，必须满足 \(k=\lim\frac{y}{x}\) 为有限值，且 \(y-kx\to b\)。有限点若要产生竖直渐近线，函数至少一侧必须无界。

\textbf{推导}\quad 正弦项绝对值不超过 \(1\)，因此四个函数在有限点附近都有界，C、D 在 \(0\) 附近也如此，均无竖直渐近线。B、D 中 \(\frac{y}{x}=x+\frac{\text{有界量}}{x}\) 在两端均没有有限极限，故无直线渐近线。

A 中 \(\frac{y}{x}\to1\)，但 \(y-x=\sin x\) 在 \(-1\) 与 \(1\) 之间持续振荡，不趋于固定截距。C 中同样有斜率 \(1\)，而 \(y-x=\sin(\frac{1}{x})\to0\)。

\textbf{结论与检查}\quad C 有渐近线 \(\boxed{y=x}\)，选 C。判断依据是曲线与这条直线之差趋零，仅有差值有界还不够。
\end{QuestionSolution}

% Q 2016 1
\QuestionPage{1.4 无穷小量与无穷大量}{选择题}{2016 年第 1 题}{1}{无穷小量的阶}{023 / 234}{**}
设$a_1 = x(\cos \sqrt{x} - 1)$，$a_2 = \sqrt{x} \ln(1 + \sqrt[3]{x})$，$a_3 = \sqrt[3]{x + 1} - 1$．
当$x \to 0^+$时，以上$3$个无穷小量按照从低阶到高阶的排序是\pickout{}
\begin{abcd}
\item $a_1,a_2,a_3$．
\item $a_2,a_3,a_1$．
\item $a_2,a_1,a_3$．
\item $a_3,a_2,a_1$．
\end{abcd}
\QuestionEnd

\begin{QuestionSolution}{2016 年第 1 题}
\textbf{已知与目标}\quad 三个无穷小分别含平方根、立方根和对数，要将它们按趋零速度从慢到快排列。

\textbf{方法与条件}\quad 对正数 \(x\to0^+\)，若 \(q>p\)，则 \(\frac{x^q}{x^p}=x^{q-p}\to0\)，因此幂指数越大，阶越高。非零常数系数的正负不影响阶。

\textbf{推导}\quad 对 \(a_1\)，令 \(u=\sqrt x\)，有 \(\cos\sqrt x-1\sim-\frac{x}{2}\)，乘外面的 \(x\) 得 \(-\frac{x^2}{2}\)。对 \(a_2\)，\(\ln(1+\sqrt[3]x)\sim x^{\frac{1}{3}}\)，再乘 \(x^{\frac{1}{2}}\) 得 \(x^{\frac{1}{2}+\frac{1}{3}}=x^{\frac{5}{6}}\)。对 \(a_3\)，用 \((1+u)^{\frac{1}{3}}-1\sim\frac{u}{3}\)。综上
\[
a_1\sim-\frac{1}{2}x^2,\qquad a_2\sim x^{\frac{5}{6}},\qquad a_3\sim\frac{1}{3}x.
\]

\textbf{结论与检查}\quad \(\frac{5}{6}<1<2\)，故顺序为 \(\boxed{a_2,a_3,a_1}\)，选 B。相邻比值 \(\frac{a_3}{a_2}\sim\frac{1}{3}x^{\frac{1}{6}}\to0\)、\(\frac{a_1}{a_3}\sim-\frac{3}{2}x\to0\)，与该顺序一致。
\end{QuestionSolution}

% Q 2021 1
\QuestionPage{1.4 无穷小量与无穷大量}{选择题}{2021 年第 1 题}{1}{变上限积分的无穷小阶}{026 / 234}{**}
当 \(x \to 0\) 时，\(\int_{0}^{x^{2}} \left(\mathrm{e}^{t^3}-1\right) \mathrm{d} t\) 是 \(x^{7}\) 的
\choices{低阶无穷小}{等价无穷小}{高阶无穷小}{同阶但非等价无穷小}
\QuestionEnd

\begin{QuestionSolution}{2021 年第 1 题}
\textbf{已知与目标}\quad 要比较 \(I(x)=\int_0^{x^2}(\mathrm e^{t^3}-1)\,\mathrm{d}t\) 与 \(x^7\)，因此计算 \(\frac{I(x)}{x^7}\) 的极限即可。

\textbf{方法与条件}\quad 当 \(t\to0\) 时，指数展开为 \(\mathrm e^{t^3}-1=t^3+O(t^6)\)。\(O(t^6)\) 表示余项的绝对值不超过某个固定常数乘 \(|t|^6\)；积分后也可以控制余项。

\textbf{推导}\quad 上限 \(x^2\ge0\)，对两侧趋零的 \(x\) 都能在 \([0,x^2]\) 上积分。主项积分为 \(\int_0^{x^2}t^3\,\mathrm{d}t=\frac{(x^2)^4}{4}=\frac{x^8}{4}\)，余项积分绝对值至多为常数乘 \(\int_0^{x^2}t^6\,\mathrm{d}t=\frac{x^{14}}{7}\)。所以
\[
\int_0^{x^2}(\mathrm{e}^{t^3}-1)\,\mathrm{d}t
=\frac{x^8}{4}+O(x^{14}).
\]
对 \(x\ne0\) 除以 \(x^7\)，得 \(\frac{I(x)}{x^7}=\frac{x}{4}+O(|x|^7)\to0\)。

\textbf{结论与检查}\quad 原积分是 \(x^7\) 的高阶无穷小，选 \(\boxed{\mathrm C}\)。主项为八次，幂次高于七次，与比值计算一致。
\end{QuestionSolution}

% Q 2023 1
\QuestionPage{1.4 无穷小量与无穷大量}{选择题}{2023 年第 1 题}{1}{斜渐近线}{028 / 234}{**}
曲线
\[y=x\ln\left(\mathrm{e}+\frac{1}{x-1}\right)
\]
的斜渐近线方程为（ ）。
\choices{\(y=x+\mathrm{e}\)}{\(y=x+\dfrac{1}{\mathrm{e}}\)}{\(y=x\)}{\(y=x-\dfrac{1}{\mathrm{e}}\)}
\QuestionEnd

\begin{QuestionSolution}{2023 年第 1 题}
\textbf{已知与目标}\quad 求曲线无穷远处的斜渐近线 \(y=kx+b\)，需要分别求斜率 \(k\) 和截距 \(b\)。

\textbf{方法与条件}\quad 若 \(k=\lim\frac{y}{x}\) 有限非零，再求 \(b=\lim(y-kx)\)。当 \(|x|\) 足够大时对数真数接近 \(\mathrm e>0\)，两端的运算都有定义。

\textbf{推导}\quad 首先 \(\frac{y}{x}=\ln(\mathrm e+\frac{1}{x-1})\to1\)，故 \(k=1\)。为求截距，将对数中的 \(\mathrm e\) 提出，利用 \(\ln(\mathrm e A)=1+\ln A\)，得到
\[
y-x=x\ln\left(1+\frac{1}{\mathrm{e}(x-1)}\right)
\sim\frac{x}{\mathrm{e}(x-1)}\longrightarrow\frac{1}{\mathrm{e}}.
\]
其中使用 \(\ln(1+u)\sim u\)，这里 \(u=\frac{1}{\mathrm e(x-1)}\to0\)，且 \(\frac{x}{x-1}\to1\)。

\textbf{结论与检查}\quad 斜渐近线为 \(\boxed{y=x+\frac{1}{\mathrm e}}\)，选 B。上述计算对 \(x\to+\infty\) 和 \(x\to-\infty\) 得到相同截距，曲线与该直线之差均趋于零。
\end{QuestionSolution}

% Q 2026 1
\QuestionPage{1.4 无穷小量与无穷大量}{选择题}{2026 年第 1 题}{1}{等价无穷小求系数}{030 / 234}{**}
已知当 \(x \to 0\) 时，\(a x^{2} + b x + \arcsin x\) 与 \(\sqrt[3]{1 + x^{2}} - 1\) 是等价无穷小，则
\choices{\(a = \frac{1}{3},\ b = -1\)}{\(a = \frac{1}{3},\ b = 1\)}{\(a = \frac{2}{3},\ b = -1\)}{\(a = \frac{2}{3},\ b = 1\)}
\QuestionEnd

\begin{QuestionSolution}{2026 年第 1 题}
\textbf{已知与目标}\quad 左式含有可调的一次、二次项，右式为立方根增量。等价要求两式的比值趋于 \(1\)，据此求 \(a,b\)。

\textbf{方法与条件}\quad Taylor 展开保留到二次，\(o(x^2)\) 表示比 \(x^2\) 更快趋零的余项。\(\arcsin\) 在零点附近光滑，其二次系数为零；立方根在 \(1\) 附近也可展开。

\textbf{推导}\quad 由 \(\arcsin x=x+o(x^2)\) 及 \((1+u)^{\frac{1}{3}}=1+\frac{u}{3}+o(u)\)，取 \(u=x^2\)，得到
\[
ax^2+bx+\arcsin x=(b+1)x+ax^2+o(x^2),
\qquad \sqrt[3]{1+x^2}-1=\frac{x^2}{3}+o(x^2).
\]
右式为二阶，所以左式的一次项必须消失，\(b+1=0\)。此时两式之比趋于 \(\frac{a}{\frac{1}{3}}=3a\)，令其等于 \(1\)，得 \(a=\frac{1}{3}\)。

\textbf{结论与检查}\quad \(\boxed{a=\frac{1}{3},\ b=-1}\)，选 A。代回后二者主项同为 \(\frac{x^2}{3}\)，满足等价关系。
\end{QuestionSolution}

% Q 2016 9
\QuestionPage{1.4 无穷小量与无穷大量}{填空题}{2016 年第 9 题}{2}{斜渐近线}{032 / 234}{**}
曲线$y = \frac{x^3}{1 + x^2} + \arctan(1 + x^2)$的斜渐近线方程为 \blank{}．
\QuestionEnd

\begin{QuestionSolution}{2016 年第 9 题}
\textbf{已知与目标}\quad 求 \(\frac{x^3}{1+x^2}+\arctan(1+x^2)\) 的斜渐近线，要计入两个加项对斜率和截距的贡献。

\textbf{方法与条件}\quad 有理项用多项式除法分出 \(x\)，再计算剩余部分的极限。反正切函数 \(\arctan u\) 当 \(u\to+\infty\) 时趋于 \(\frac{\pi}{2}\)。

\textbf{推导}\quad 先写 \(\frac{x^3}{1+x^2}=x-\frac{x}{1+x^2}\)。有理余项趋零，反正切项有界，所以 \(\frac{y}{x}\to1\)，斜率为 \(1\)。截距为
\[
y-x=-\frac{x}{1+x^2}+\arctan(1+x^2)\longrightarrow0+\frac{\pi}{2}
\quad(x\to\pm\infty).
\]

\textbf{结论与检查}\quad 斜渐近线为 \(\boxed{y=x+\frac{\pi}{2}}\)。正负无穷远处都有 \(1+x^2\to+\infty\)，所以两端的截距均为正的 \(\frac{\pi}{2}\)。
\end{QuestionSolution}

% Q 2017 9
\QuestionPage{1.4 无穷小量与无穷大量}{填空题}{2017 年第 9 题}{2}{斜渐近线}{033 / 234}{**}
曲线$y = x\left(1 + \arcsin \frac{2}{x} \right)$的斜渐近线方程为 \blank{}．
\QuestionEnd

\begin{QuestionSolution}{2017 年第 9 题}
\textbf{已知与目标}\quad 曲线含 \(x\arcsin(\frac{2}{x})\) 这一“无界量乘无穷小”，要求其最终给斜渐近线留下的截距。

\textbf{方法与条件}\quad 反正弦有定义要求 \(|\frac{2}{x}|\le1\)，即 \(|x|\ge2\)，因此正负无穷远均可考虑。用 \(\frac{\arcsin u}{u}\to1\) 计算乘积。

\textbf{推导}\quad 因 \(\frac{2}{x}\to0\)，\(\frac{y}{x}=1+\arcsin(\frac{2}{x})\to1\)，故斜率为 \(1\)。截距满足
\[
y-x=x\arcsin(\frac{2}{x})=2\,\frac{\arcsin(\frac{2}{x})}{\frac{2}{x}}\longrightarrow2.
\]

\textbf{结论与检查}\quad 斜渐近线为 \(\boxed{y=x+2}\)。两端的乘积极限都是 \(2\)；当 \(x<0\) 时两个因子同时为负，截距符号仍为正。
\end{QuestionSolution}

% Q 2023 11
\QuestionPage{1.4 无穷小量与无穷大量}{填空题}{2023 年第 11 题}{2}{等价无穷小求系数}{034 / 234}{**}
当 \(x\to 0\) 时，函数
\[f(x)=ax+bx^2+\ln(1+x)
\]
与
\[g(x)=\mathrm{e}^{x^2}-\cos x
\]
是等价无穷小，则
\[ab=\blank{}.
\]
\QuestionEnd

\begin{QuestionSolution}{2023 年第 11 题}
\textbf{已知与目标}\quad 由 \(f\sim g\) 求参数乘积 \(ab\)，需先确定两边最低非零项。

\textbf{方法与条件}\quad 等价表示比值趋于 \(1\)。将指数、余弦、对数分别 Taylor 展开至二次；\(o(x^2)\) 表示余项除以 \(x^2\) 后趋零。

\textbf{推导}\quad 使用 \(\mathrm e^{x^2}=1+x^2+o(x^2)\)、\(\cos x=1-\frac{x^2}{2}+o(x^2)\)，得到
\[
g(x)=(1+x^2)-(1-\frac{x^2}{2})+o(x^2)=\frac{3}{2}x^2+o(x^2).
\]
对左式代入 \(\ln(1+x)=x-\frac{x^2}{2}+o(x^2)\)，得
\[
f(x)=(a+1)x+\left(b-\frac{1}{2}\right)x^2+o(x^2).
\]
右式是二阶，故必须先令一次项系数 \(a+1=0\)，再令二次系数 \(b-\frac{1}{2}=\frac{3}{2}\)。解得 \(a=-1,b=2\)。

\textbf{结论与检查}\quad \(\boxed{ab=-2}\)。代回后 \(f,g\) 都以 \(\frac{3}{2}x^2\) 为主项，其比值趋于 \(1\)。
\end{QuestionSolution}

% Q 2025 12
\QuestionPage{1.4 无穷小量与无穷大量}{填空题}{2025 年第 12 题}{2}{斜渐近线}{035 / 234}{**}
曲线 \(y = \sqrt[3]{x^{3} - 3x^{2} + 1}\) 的渐近线方程为 \blank{}.
\QuestionEnd

\begin{QuestionSolution}{2025 年第 12 题}
\textbf{已知与目标}\quad 求立方根曲线的渐近线，需要保留提出 \(x\) 后的首个修正项来确定截距。

\textbf{方法与条件}\quad 实数立方根满足 \(\sqrt[3]{x^3A}=x\sqrt[3]A\)，对负 \(x\) 也成立。再用 \(\frac{(1+u)^{\frac{1}{3}}-1}{u}\to\frac{1}{3}\) 求剩余乘积。

\textbf{推导}\quad 对 \(x\ne0\)，写成 \(y=x(1-\frac{3}{x}+\frac{1}{x^3})^{\frac{1}{3}}\)，得到 \(\frac{y}{x}\to1\)。令 \(u=-\frac{3}{x}+\frac{1}{x^3}\to0\)，则 \(xu=-3+\frac{1}{x^2}\to-3\)，所以
\[
y-x=xu\,\frac{(1+u)^{\frac{1}{3}}-1}{u}\longrightarrow-3\cdot\frac{1}{3}=-1
\quad(x\to\pm\infty).
\]

\textbf{结论与检查}\quad 渐近线为 \(\boxed{y=x-1}\)。实数立方根在所有有限点连续，不会有函数值趋于无穷的有限点；两端的斜率、截距也相同。
\end{QuestionSolution}

% Q 2012 15
\QuestionPage{1.4 无穷小量与无穷大量}{解答题}{2012 年第 15 题}{2}{同阶无穷小与参数}{036 / 234}{**}
已知函数$f(x) = \frac{1 + x}{\sin x} - \frac{1}{x}$，记$a = \lim_{x \to 0} f(x)$．\par
\begin{enumerate}
\item 求$a$的值；
\item 若当$x \to 0$时，$f(x) - a$是$x^k$的同阶无穷小，求$k$．
\end{enumerate}
\QuestionEnd

\begin{QuestionSolution}{2012 年第 15 题}
\textbf{已知与目标}\quad 第一问求 \(f(x)\) 的极限 \(a\)；第二问比较去掉极限后的余量 \(f(x)-a\) 与 \(x^k\) 的阶。

\textbf{方法与条件}\quad 原式是两个分别无界的量相减，先通分或拆成可求极限的两项。使用 \(\sin x=x-\frac{x^3}{6}+o(x^3)\)，其中余项除以 \(x^3\) 趋零；小去心邻域内 \(x\sin x\ne0\)。

\textbf{推导}\quad （1）把分子 \(1+x\) 分开，并将 \(\frac{1}{\sin x}-\frac{1}{x}\) 通分，得
\[
f(x)=\frac{x}{\sin x}+\frac{x-\sin x}{x\sin x}.
\]
第一项 \(\frac{x}{\sin x}\to1\)，第二项的分子为 \(\frac{x^3}{6}+o(x^3)\)，分母等价于 \(x^2\)，因此第二项等价于 \(\frac{x}{6}\to0\)，得到 \(a=1\)。

（2）扣去 \(1\) 后，第一项余量为 \(\frac{x-\sin x}{\sin x}\sim\frac{x^2}{6}\)，它除以 \(x\) 趋零。第二项除以 \(x\) 则趋于 \(\frac{1}{6}\)。故
\[
\lim_{x\to0}\frac{f(x)-1}{x}=0+\frac{1}{6}=\frac{1}{6}.
\]

\textbf{结论与检查}\quad \(\boxed{a=1,\ k=1}\)。比值为有限非零常数，符合“同阶”；常数不为 \(1\)，所以 \(f(x)-a\) 与 \(x\) 并不等价。
\end{QuestionSolution}

% Q 2013 15
\QuestionPage{1.4 无穷小量与无穷大量}{解答题}{2013 年第 15 题}{2}{三角乘积的等价无穷小}{037 / 234}{**}
当$x \to 0$时，$1 - \cos x \cdot \cos 2x \cdot \cos 3x$与$ax^n$为等价无穷小，求$n$与$a$的值．
\QuestionEnd

\begin{QuestionSolution}{2013 年第 15 题}
\textbf{已知与目标}\quad 找出 \(1-\cos x\cos2x\cos3x\) 的最低非零幂次和系数，使它除以 \(ax^n\) 后趋于 \(1\)。

\textbf{方法与条件}\quad 每个余弦都趋于 \(1\)，需要保留二次项才能看到乘积与 \(1\) 的差。使用 Taylor 展开，\(o(x^2)\) 表示比 \(x^2\) 更高阶的余项。

\textbf{推导}\quad 对 \(j=1,2,3\)，有 \(\cos(jx)=1-\frac{j^2x^2}{2}+o(x^2)\)。乘积中的二次项来自分别选一个因子的二次项，其余两个选常数 \(1\)。选两个二次项相乘已为四次，对二次主项无贡献。因此
\[
1-\cos x\cos 2x\cos 3x
=\frac{1+4+9}{2}x^2+o(x^2)=7x^2+o(x^2).
\]

\textbf{结论与检查}\quad \(\boxed{n=2,\ a=7}\)。原式除以 \(7x^2\) 后为 \(1+o(1)\)，极限为 \(1\)，与等价要求一致。
\end{QuestionSolution}

% Q 2020 1
\QuestionPage{1.4 无穷小量与无穷大量}{选择题}{2020 年第 1 题}{1}{积分形式的无穷小比阶}{025 / 234}{***}
当 \(x \to 0^{+}\) 时，下列无穷小量中阶最高的是
\choices{\(\int_{0}^{x}{\left( \mathrm{e}^{t^{2}} - 1 \right) \,\mathrm{d} t}\)}{\(\int_{0}^{x}{\ln \left( 1 + \sqrt{t^{3}} \right) \,\mathrm{d} t}\)}{\(\int_{0}^{\sin x}{\sin t^{2} \,\mathrm{d} t}\)}{\(\int_{0}^{1-\cos x}{\sqrt{\sin^{3} t} \,\mathrm{d} t}\)}
\QuestionEnd

\begin{QuestionSolution}{2020 年第 1 题}
\textbf{已知与目标}\quad 比较四个积分作为 \(x\to0^+\) 时的无穷小阶。既要考虑被积函数的幂次，也要考虑积分上限的幂次。

\textbf{方法与条件}\quad 若连续的 \(\,h(t)\sim ct^p\)、\(c\ne0\)、\(p>-1\)，则 \(\int_0^u h(t)\,\mathrm{d}t\sim\frac{c}{p+1}u^{p+1}\)。这里分子和比较量均趋零，L'Hospital 法则将二者比值化为 \(\frac{h(u)}{cu^p}\to1\)，说明这一积分等价式成立。

\textbf{推导}\quad A 中 \(\mathrm e^{t^2}-1\sim t^2\)，积分后为 \(\frac{x^3}{3}\)。B 中 \(\sqrt{t^3}=t^{\frac{3}{2}}\)，故对数等价于 \(t^{\frac{3}{2}}\)，积分后为 \(\frac{2}{5}x^{\frac{5}{2}}\)。C 中 \(\sin(t^2)\sim t^2\)，先得上限 \(u\) 的三次量，再代入 \(u=\sin x\sim x\)。D 中 \(t\ge0\) 且接近零时 \(\sin t\ge0\)，所以 \(\sqrt{\sin^3t}\sim t^{\frac{3}{2}}\)，再代入 \(u=1-\cos x\sim\frac{x^2}{2}\)。四项依次等价于
\[
\frac{x^3}{3},\qquad \frac{2}{5}x^{\frac{5}{2}},\qquad
\frac{x^3}{3},\qquad \frac{2}{5}\left(\frac{x^2}{2}\right)^{\frac{5}{2}}.
\]

\textbf{结论与检查}\quad 幂次分别为 \(3,\frac{5}{2},3,5\)，所以 D 阶最高，选 \(\boxed{\mathrm D}\)。四个最低项系数都为正且非零，比较幂次足以判定。
\end{QuestionSolution}

% Q 2022 1
\QuestionPage{1.4 无穷小量与无穷大量}{选择题}{2022 年第 1 题}{1}{等价无穷小的性质}{027 / 234}{***}
若当 \(x \to 0\) 时，\(\alpha\left( x \right), \beta\left( x \right)\) 是非零无穷小量，则以下的命题中，
\par\circnum{1} 若 \(\alpha\left( x \right) \sim \beta\left( x \right)\)，则 \(\alpha^{2}\left( x \right) \sim \beta^{2}\left( x \right)\)；
\par\circnum{2} 若 \(\alpha^{2}\left( x \right) \sim \beta^{2}\left( x \right)\)，则 \(\alpha\left( x \right) \sim \beta\left( x \right)\)；
\par\circnum{3} 若 \(\alpha\left( x \right) \sim \beta\left( x \right)\)，则 \(\alpha\left( x \right) - \beta\left( x \right) = o\left( \alpha\left( x \right) \right)\)；
\par\circnum{4} 若 \(\alpha\left( x \right) - \beta\left( x \right) = o\left( \alpha\left( x \right) \right)\)，则 \(\alpha\left( x \right) \sim \beta\left( x \right)\)，
真命题的序号为 \(\left( \ \right)\)
\choices{\circnum{1}\circnum{3}}{\circnum{1}\circnum{4}}{\circnum{1}\circnum{3}\circnum{4}}{\circnum{2}\circnum{3}\circnum{4}}
\QuestionEnd

\begin{QuestionSolution}{2022 年第 1 题}
\textbf{已知与目标}\quad \(\alpha,\beta\) 都趋于零且在去心邻域非零，要逐条判断四个关系是否必然成立。

\textbf{方法与条件}\quad \(\alpha\sim\beta\) 表示 \(\frac{\alpha}{\beta}\to1\)；\(r=o(\alpha)\) 表示 \(\frac{r}{\alpha}\to0\)。非零条件保证可以作商，而极限为 \(1\) 时也可以取倒数。

\textbf{推导}\quad 命题 1：若 \(\frac{\alpha}{\beta}\to1\)，则 \(\frac{\alpha^2}{\beta^2}=(\frac{\alpha}{\beta})^2\to1\)，故成立。

命题 2：取 \(\alpha(x)=x,\beta(x)=-x\)，二者在 \(x\ne0\) 时非零且都趋零。平方比恒为 \(1\)，但原比恒为 \(-1\)，所以不成立。

命题 3：等价关系给出 \(\frac{\beta}{\alpha}\to1\)，故 \(\frac{\alpha-\beta}{\alpha}=1-\frac{\beta}{\alpha}\to0\)，成立。命题 4：若这个差商趋零，同一等式反推 \(\frac{\beta}{\alpha}\to1\)，再取倒数得 \(\frac{\alpha}{\beta}\to1\)，也成立。

\textbf{结论与检查}\quad 真命题的序号为 \(\boxed{1,3,4}\)，选 C。平方会丢失正负号，故命题 2 不能从平方等价反推原量等价。
\end{QuestionSolution}

% Q 2025 4
\QuestionPage{1.4 无穷小量与无穷大量}{选择题}{2025 年第 4 题}{1}{无穷小比阶}{029 / 234}{***}
设函数 \(f\left( x \right)\)，\(g\left( x \right)\) 在 \(x = 0\) 的某去心邻域内有定义且恒不为零.若当 \(x \to 0\) 时，\(f\left( x \right)\) 是 \(g\left( x \right)\) 的高阶无穷小，则当 \(x \to 0\) 时，
\choices{\(f\left( x \right) + g\left( x \right) = o\left( g\left( x \right) \right)\).}{\(f\left( x \right) g\left( x \right) = o\left( f^{2}\left( x \right) \right)\).}{\(f\left( x \right) = o\left( \mathrm{e}^{g\left( x \right)} - 1 \right)\).}{\(f\left( x \right) = o\left( g^{2}\left( x \right) \right)\).}
\QuestionEnd

\begin{QuestionSolution}{2025 年第 4 题}
\textbf{已知与目标}\quad 题设中的高阶无穷小表示 \(f\to0\)、\(g\to0\) 且 \(\frac{f}{g}\to0\)。要判断这个条件能保证四个小 \(o\) 关系中的哪一个。

\textbf{方法与条件}\quad \(u=o(v)\) 的含义是 \(\frac{u}{v}\to0\)，逐项作商即可。题设保证 \(f,g\) 在去心邻域非零；又因实数 \(g\ne0\) 时 \(\mathrm e^g-1\ne0\)，各商均有定义。

\textbf{推导}\quad A 中 \(\frac{f+g}{g}=1+\frac{f}{g}\to1\)，故不成立。B 中 \(\frac{fg}{f^2}=\frac{g}{f}\)，其绝对值是 \(|\frac{f}{g}|\) 的倒数，趋于正无穷，不会趋零。

C 中，利用 \(g\to0\) 和基本等价式 \(\mathrm e^g-1\sim g\)，有
\[
\frac{f}{\mathrm{e}^g-1}=\frac{f}{g}\frac{g}{\mathrm{e}^g-1}\longrightarrow0\cdot1=0.
\]
D 则不能保证：取 \(g(x)=x,f(x)=x^2\)，在 \(x\ne0\) 时二者非零，且 \(\frac{f}{g}=x\to0\)，但 \(\frac{f}{g^2}=1\)。

\textbf{结论与检查}\quad 只有 C 必然成立，选 \(\boxed{\mathrm C}\)。高于 \(g\) 的阶，并不能保证也高于 \(g^2\) 的阶。
\end{QuestionSolution}

% Q 2015 15
\QuestionPage{1.4 无穷小量与无穷大量}{解答题}{2015 年第 15 题}{2}{等价无穷小求参数}{038 / 234}{***}
设函数$f(x) = x + a\ln(1 + x) + bx\sin x$，$g(x) = kx^3$，
若$f(x)$与$g(x)$在$x \to 0$是等价无穷小，求$a,b,k$的值．
\QuestionEnd

\begin{QuestionSolution}{2015 年第 15 题}
\textbf{已知与目标}\quad \(f(x)\) 与 \(kx^3\) 等价，需要确定 \(a,b,k\)，并保证三次主项的系数非零。

\textbf{方法与条件}\quad 等价表示 \(\frac{f(x)}{kx^3}\to1\)，所以低于三次的项都必须消失。对数在零点附近有定义，可以 Taylor 展开至三次；\(o(x^3)\) 表示除以 \(x^3\) 后趋零的余项。

\textbf{推导}\quad 代入 \(\ln(1+x)=x-\frac{x^2}{2}+\frac{x^3}{3}+o(x^3)\)。由 \(\sin x=x-\frac{x^3}{6}+o(x^3)\)，又得 \(x\sin x=x^2+o(x^3)\)，故
\[
f(x)=(1+a)x+(b-\frac{a}{2})x^2+(\frac{a}{3})x^3+o(x^3).
\]
一次系数为零给出 \(1+a=0\)，所以 \(a=-1\)。二次系数为零给出 \(b-\frac{a}{2}=0\)，代入得 \(b=-\frac{1}{2}\)。此时 \(f(x)=-\frac{x^3}{3}+o(x^3)\)，因此 \(k=-\frac{1}{3}\)。

\textbf{结论与检查}\quad \(\boxed{a=-1,\ b=-\frac{1}{2},\ k=-\frac{1}{3}}\)。所求 \(k\ne0\)，代回后与 \(kx^3\) 的比值趋于 \(1\)，三个参数均符合条件。
\end{QuestionSolution}

% Q 2020 15
\QuestionPage{1.4 无穷小量与无穷大量}{解答题}{2020 年第 15 题}{2}{斜渐近线}{039 / 234}{***}
求曲线 \(y = \dfrac{x^{1+x}}{\left( 1+x \right)^{x}} \left( x>0 \right)\) 的斜渐近线方程。
\QuestionEnd

\begin{QuestionSolution}{2020 年第 15 题}
\textbf{已知与目标}\quad 定义域为 \(x>0\)，只需在 \(x\to+\infty\) 求斜渐近线 \(y=kx+b\)。斜率由 \(\frac{y}{x}\) 决定，截距还需要更精确的一阶修正。

\textbf{方法与条件}\quad 将正底数的幂写成指数函数：\(y=x(1+\frac{1}{x})^{-x}\)。取对数后对 \(\ln(1+v)\) 展开至二次，因为后续会乘上 \(x\)。记号 \(o(\frac{1}{x})\) 表示余项乘以 \(x\) 后趋零。

\textbf{推导}\quad 令 \(u(x)=-x\ln(1+\frac{1}{x})\)，则 \(y=x\mathrm e^{u(x)}\)。代入 \(\ln(1+v)=v-\frac{v^2}{2}+o(v^2)\)，得
\[
u(x)=-1+\frac{1}{2x}+o(\frac{1}{x}),
\qquad
\mathrm{e}^{u(x)}=\mathrm{e}^{-1}\left(1+\frac{1}{2x}+o(\frac{1}{x})\right).
\]
第二个等式用了 \(\mathrm e^w=1+w+o(w)\)，其中 \(w=u+1=\frac{1}{2x}+o(\frac{1}{x})\)。因此
\[
k=\lim_{x\to+\infty}\frac{y}{x}=\frac{1}{\mathrm e},\qquad
y-\frac{x}{\mathrm e}=\frac{1}{2\mathrm e}+o(1)\longrightarrow\frac{1}{2\mathrm e}.
\]

\textbf{结论与检查}\quad 斜渐近线为 \(\boxed{y=\frac{x}{\mathrm e}+\frac{1}{2\mathrm e}}\)。仅使用 \((1+\frac{1}{x})^x\to\mathrm e\) 只能求斜率，保留二次对数项才确定了截距。
\end{QuestionSolution}

\SectionHeading{1.5}{函数、极限、连续}{连续函数}{定理与要点\quad 教材 PDF 第 86 页起\hfill 本节 8 题}
\pdfbookmark[2]{1.5 连续函数}{section-1.5}
\Theory{定理 5.1}{90}{连续函数的四则运算}{在同一点连续的函数，其和、差、积、商仍连续，取商时分母在该点非零。}
\Theory{定理 5.2}{90}{复合函数的连续性}{内层函数在 \(x_0\) 处连续，且外层函数在对应的内层函数值处连续，则复合函数在 \(x_0\) 处连续。}
\Theory{定理 5.3}{91}{反函数的连续性}{区间上的连续严格单调函数有反函数，且反函数在值域上连续并保持单调方向。}
\Theory{定理 5.4}{94}{有界性}{闭区间上的连续函数有界。}
\Theory{定理 5.5}{94}{最大最小值定理}{闭区间上的连续函数能取到最大值和最小值。}
\Theory{定理 5.6}{95}{零点存在定理}{闭区间上连续且两端值异号，区间内部至少有一处零点。}
\Theory{定理 5.7}{97}{介值定理}{闭区间上的连续函数取到介于端点函数值之间的每一个值。}
\Theory{定理 5.8}{100}{一致连续性}{闭区间上的连续函数一致连续。}
\Theory{定理 5.9}{101}{压缩映射原理（星号内容）}{实直线上的压缩映射有唯一不动点，从任意实数开始迭代均收敛到该点。}
\Theory{要点}{86--89}{间断点分类}{先检查两侧极限是否有限且相等，再将共同极限与函数值比较。}
% Q 2009 1
\QuestionPage{1.5 连续函数}{选择题}{2009 年第 1 题}{1}{可去间断点}{040 / 234}{**}
函数$f(x) = \frac{x - x^3}{\sin\pi x}$的可去间断点的个数为\pickout{}
\begin{abcd}
\item $1$．
\item $2$．
\item $3$．
\item 无穷多个．
\end{abcd}
\QuestionEnd

\begin{QuestionSolution}{2009 年第 1 题}
\textbf{已知与目标}\quad 分母 \(\sin\pi x\) 在所有整数处为零，需要在这些无定义点中找出极限有限的可去间断点。

\textbf{方法与条件}\quad 可去间断点指双侧极限存在且有限，但该点无定义或函数值不等于极限。分子也为零的点可用 L'Hospital 法则求 \(\frac{0}{0}\) 型极限；分母导数须在附近非零。

\textbf{推导}\quad \(\sin\pi x=0\) 恰在整数 \(k\) 处成立。分子 \(x-x^3=x(1-x)(1+x)\) 仅在 \(k=-1,0,1\) 为零。对这三点，分母导数在点上为 \(\pi\cos\pi k=\pi(-1)^k\ne0\)，附近也非零，因此
\[
\lim_{x\to k}\frac{x-x^3}{\sin\pi x}=\frac{1-3k^2}{\pi(-1)^k}
=\begin{cases}\frac{1}{\pi},&k=0,\\[2pt]\frac{2}{\pi},&k=\pm1.\end{cases}
\]
其余整数处分子趋于非零值，分母趋零，商的绝对值趋于无穷，不属于可去间断点。非整数点处分母不为零，函数连续。

\textbf{结论与检查}\quad 可去间断点为 \(-1,0,1\)，共 \(\boxed{3}\) 个，选 C。所有可能间断的点均已按整数与非整数排查。
\end{QuestionSolution}

% Q 2010 1
\QuestionPage{1.5 连续函数}{选择题}{2010 年第 1 题}{1}{无穷间断点}{041 / 234}{**}
函数$f(x) = \frac{x^2 - x}{x^2 - 1}\sqrt{1 + \frac{1}{x^2}}$的无穷间断点的个数为\pickout{}
\begin{abcd}
\item $0$．
\item $1$．
\item $2$．
\item $3$．
\end{abcd}
\QuestionEnd

\begin{QuestionSolution}{2010 年第 1 题}
\textbf{已知与目标}\quad 原式在 \(-1,0,1\) 处无定义。要找的是至少有一侧函数值趋于无穷的无穷间断点。

\textbf{方法与条件}\quad 约分仅在原定义域内成立。化简根号必须用 \(\sqrt{x^2}=|x|\)，否则会丢失左右两侧的符号差异。

\textbf{推导}\quad 当 \(x\ne-1,0,1\) 时，
\[
f(x)=\frac{x}{x+1}\cdot\frac{\sqrt{x^2+1}}{|x|}
=\frac{x}{|x|}\frac{\sqrt{x^2+1}}{x+1}.
\]
在 \(-1\) 附近 \(\frac{x}{|x|}=-1\)，分子趋于 \(-\sqrt2\)，而 \(x+1\to0\)，故左极限为 \(+\infty\)，右极限为 \(-\infty\)。在 \(0\) 附近，\(\frac{\sqrt{x^2+1}}{x+1}\to1\)，但 \(\frac{x}{|x|}\) 左侧为 \(-1\)、右侧为 \(1\)，所以两个单侧极限有限且不相等，为跳跃间断点。在 \(1\) 处，化简式的极限为 \(\frac{\sqrt2}{2}\)，是可去间断点。

\textbf{结论与检查}\quad 无穷间断点只有 \(x=-1\)，共 \(\boxed{1}\) 个，选 B。其余定义域内各因子连续，没有额外间断点。
\end{QuestionSolution}

% Q 2015 2
\QuestionPage{1.5 连续函数}{选择题}{2015 年第 2 题}{1}{由极限定义的函数及间断点}{042 / 234}{**}
函数$f(x) = \lim_{t \to 0} \left(1 + \frac{\sin t}{x} \right)^{\frac{x^2}{t}}$
在$(-\infty , +\infty)$内\pickout{}
\begin{abcd}
\item 连续．
\item 有可去间断点．
\item 有跳跃间断点．
\item 有无穷间断点．
\end{abcd}
\QuestionEnd

\begin{QuestionSolution}{2015 年第 2 题}
\textbf{已知与目标}\quad 函数值由关于 \(t\) 的极限定义，最终要判断它作为 \(x\) 的函数在实轴上的连续性。

\textbf{方法与条件}\quad 先固定 \(x\ne0\)，此时 \(\frac{\sin t}{x}\to0\)，底数在 \(t=0\) 附近为正，可以取对数。求完 \(t\) 极限后再检查 \(x=0\)，两次极限的变量不能混淆。

\textbf{推导}\quad 固定 \(x\ne0\)，原幂的对数满足
\[
\frac{x^2}{t}\ln\left(1+\frac{\sin t}{x}\right)
=x\,\frac{\sin t}{t}
\frac{\ln\!\left(1+\frac{\sin t}{x}\right)}{\frac{\sin t}{x}}\longrightarrow x.
\]
其中两个比值分别趋于 \(1\)，所以指数化得 \(f(x)=\mathrm e^x\)（\(x\ne0\)）。该函数在正负半轴均连续；原定义因含 \(\frac{\sin t}{x}\)，在 \(x=0\) 无定义。可是 \(\lim_{x\to0^-}f(x)=\lim_{x\to0^+}f(x)=1\)。

\textbf{结论与检查}\quad \(x=0\) 为可去间断点，选 \(\boxed{\mathrm B}\)。补上 \(f(0)=1\) 可以使之连续，但题目没有给出这个补定义，不能直接判全轴连续。
\end{QuestionSolution}

% Q 2017 1
\QuestionPage{1.5 连续函数}{选择题}{2017 年第 1 题}{1}{分段函数的连续条件}{044 / 234}{**}
若函数$f(x) = \begin{cases}
  \frac{1 - \cos \sqrt{x}}{ax}, & x > 0 \\
                             b, & x \le 0
\end{cases}$在$x = 0$处连续，则\pickout{}
\begin{abcd}
\item $ab = \frac{1}{2}$．
\item $ab = - \frac{1}{2}$．
\item $ab = 0$．
\item $ab = 2$．
\end{abcd}
\QuestionEnd

\begin{QuestionSolution}{2017 年第 1 题}
\textbf{已知与目标}\quad 左段及零点的函数值都为 \(b\)，要求右段也在 \(x\to0^+\) 时趋于 \(b\)，由此确定 \(ab\)。

\textbf{方法与条件}\quad 零点连续要求左极限、右极限、函数值三者相等。右段作为函数有定义需 \(a\ne0\)，并可用 \(1-\cos u\sim\frac{u^2}{2}\)。

\textbf{推导}\quad 令 \(u=\sqrt x\to0^+\)，则 \(u^2=x\)，所以
\[
\lim_{x\to0^+}\frac{1-\cos\sqrt{x}}{ax}
=\frac{1}{a}\lim_{u\to0^+}\frac{1-\cos u}{u^2}=\frac{1}{2a}.
\]
左极限为 \(b\)，函数值也为 \(b\)，故连续条件是 \(b=\frac{1}{2a}\)。两边乘以非零的 \(a\)，得 \(ab=\frac{1}{2}\)。

\textbf{结论与检查}\quad \(\boxed{ab=\frac{1}{2}}\)，选 A。该关系同时保证 \(a,b\) 非零，并使两个单侧极限与函数值相等。
\end{QuestionSolution}

% Q 2018 3
\QuestionPage{1.5 连续函数}{选择题}{2018 年第 3 题}{1}{分段函数和的连续性}{045 / 234}{**}
设函数$f(x) = \begin{cases}
  -1, & x < 0, \\
   1, & x \ge 0;
\end{cases}$，$g(x) = \begin{cases}
  2 - ax, & x \le - 1, \\
       x, & -1 < x < 0, \\
   x - b, & x \ge 0.
\end{cases}$若$f(x) + g(x)$在$R$上连续，则\pickout{}
\begin{abcd}
\item $a = 3\;,\;b = 1$．
\item $a = 3\;,\;b = 2$．
\item $a = - 3\;,\;b = 1$．
\item $a = - 3\;,\;b = 2$．
\end{abcd}
\QuestionEnd

\begin{QuestionSolution}{2018 年第 3 题}
\textbf{已知与目标}\quad 要求和函数 \(h=f+g\) 在全体实数上连续，求参数 \(a,b\)。各分段内都是一次式，唯一需检查的拼接点是 \(-1,0\)。

\textbf{方法与条件}\quad 先将两个分段函数相加，再在每个拼接点要求左右极限等于该点函数值。两函数单独有跳跃时，和仍可能通过相互抵消而连续。

\textbf{推导}\quad 相加得到
\[
h(x)=\begin{cases}1-ax,&x\le-1,\\ x-1,&-1<x<0,\\x+1-b,&x\ge0.\end{cases}
\]
在 \(-1\) 处，左极限和函数值为 \(1+a\)，右极限为 \(-1-1=-2\)，故 \(1+a=-2\)，得 \(a=-3\)。在 \(0\) 处，左极限为 \(-1\)，右极限及函数值为 \(1-b\)，故 \(1-b=-1\)，得 \(b=2\)。

\textbf{结论与检查}\quad \(\boxed{a=-3,\ b=2}\)，选 D。代回后在 \(-1\) 处两段都取 \(-2\)，在 \(0\) 处两段都趋于 \(-1\)，其余点各段本来就连续。
\end{QuestionSolution}

% Q 2015 3
\QuestionPage{1.5 连续函数}{选择题}{2015 年第 3 题}{1}{分段函数连续性与参数}{043 / 234}{***}
设函数$f(x) = \begin{cases}
  x^{\alpha}\cos \frac{1}{x^{\beta}}, & x > 0 \\
                                   0, & x \le 0
\end{cases}$（$\alpha > 0,\beta > 0$），若$f'(x)$在$x = 0$处连续，则\pickout{}
\begin{abcd}
\item $\alpha - \beta > 1$．
\item $0 < \alpha - \beta \le 1$．
\item $\alpha - \beta > 2$．
\item $0 < \alpha - \beta \le 2$．
\end{abcd}
\QuestionEnd

\begin{QuestionSolution}{2015 年第 3 题}
\textbf{已知与目标}\quad \(\alpha,\beta>0\)，左段恒为零。要求导函数 \(f'\) 在零点连续，既须保证 \(f'(0)\) 存在，也须使右侧导数趋于同一个值。

\textbf{方法与条件}\quad 零点导数用差商定义，不能把只对 \(x>0\) 成立的求导式直接代入零。对振荡项的必要条件，可选正弦为 \(1\)、余弦为 \(0\) 的点，避免两项抵消。

\textbf{推导}\quad 在 \(x>0\) 上，用乘积法则和复合求导法则，\((x^{-\beta})'=-\beta x^{-\beta-1}\)，得到
\[
f'(x)=\alpha x^{\alpha-1}\cos(x^{-\beta})
+\beta x^{\alpha-\beta-1}\sin(x^{-\beta}).
\]
必要性：若 \(f'\) 连续，由左段导数恒为零知 \(f'(0)=0\)。取 \(x_k=(\frac{\pi}{2}+2k\pi)^{-\frac{1}{\beta}}\to0^+\)，此时余弦为零、正弦为 \(1\)，所以 \(f'(x_k)=\beta x_k^{\alpha-\beta-1}\)。其趋零要求指数 \(\alpha-\beta-1>0\)；指数为零时值为 \(\beta\ne0\)，指数为负时无界。

充分性：若 \(\alpha-\beta>1\)，则 \(\alpha>1\)。零点右侧差商的绝对值为 \(|x^{\alpha-1}\cos(x^{-\beta})|\le x^{\alpha-1}\to0\)，与左侧差商零相等，故 \(f'(0)=0\)。上面的右侧导数满足
\[
|f'(x)|\le\alpha x^{\alpha-1}+\beta x^{\alpha-\beta-1}\longrightarrow0,
\]
所以 \(f'\) 在零点连续。

\textbf{结论与检查}\quad 条件为 \(\boxed{\alpha-\beta>1}\)，选 A。必要性与充分性都已证明，等号情形被选取的数列明确排除。
\end{QuestionSolution}

% Q 2020 2
\QuestionPage{1.5 连续函数}{选择题}{2020 年第 2 题}{1}{第二类间断点}{046 / 234}{***}
函数 \(f\left( x \right) = \dfrac{\mathrm{e}^{\frac{1}{x-1}} \ln \left| 1+x \right|}{\left( \mathrm{e}^{x} - 1 \right)\left( x-2 \right)}\) 的第二类间断点的个数为
\choices{\(1\) 个}{\(2\) 个}{\(3\) 个}{\(4\) 个}
\QuestionEnd

\begin{QuestionSolution}{2020 年第 2 题}
\textbf{已知与目标}\quad 指数中的分母、对数真数和总分母共同限制定义域，要统计第二类间断点的个数。

\textbf{方法与条件}\quad 第一类间断点要求左右极限都存在且有限；至少有一侧没有有限极限的为第二类。先列出全部可疑点，再逐一判断，不能仅按分母为零计数。

\textbf{推导}\quad 原式恰在 \(-1,0,1,2\) 处被排除，其余各点各因子连续且分母非零。

在 \(-1\) 处，\(\ln|1+x|\to-\infty\)，系数 \(\frac{\mathrm e^{\frac{1}{x-1}}}{(\mathrm e^x-1)(x-2)}\) 趋于正的非零有限值，故 \(f\to-\infty\)，为第二类。在 \(0\) 处，\(\frac{\ln(1+x)}{\mathrm e^x-1}\to1\)，其他因子趋于 \(\frac{\mathrm e^{-1}}{-2}\)，所以双侧极限为 \(-\frac{1}{2\mathrm e}\)，为可去间断点。

在 \(1\) 处，剩余因子 \(\frac{\ln|1+x|}{(\mathrm e^x-1)(x-2)}\to-\frac{\ln2}{\mathrm e-1}\ne0\)。指数因子左侧趋零、右侧趋正无穷，故 \(f(1^-)=0\)、\(f(1^+)=-\infty\)，为第二类。在 \(2\) 处，\((x-2)f(x)\to\frac{\mathrm e\ln3}{\mathrm e^2-1}>0\)，所以左右分别趋于 \(-\infty,+\infty\)，也为第二类。

\textbf{结论与检查}\quad 第二类间断点为 \(-1,1,2\)，共 \(\boxed{3}\) 个，选 C。\(x=0\) 虽不在定义域，却有有限双侧极限，已从第二类中排除。
\end{QuestionSolution}

% Q 2024 1
\QuestionPage{1.5 连续函数}{选择题}{2024 年第 1 题}{1}{间断点个数}{047 / 234}{***}
函数
\[f(x)=|x|^{\frac{1}{(1-x)(x-2)}}
\]
的第一类间断点的个数为（ ）
\choices{\(3\)}{\(2\)}{\(1\)}{\(0\)}
\QuestionEnd

\begin{QuestionSolution}{2024 年第 1 题}
\textbf{已知与目标}\quad 幂函数的底数与指数都随 \(x\) 改变，要统计左右极限都有限的第一类间断点。

\textbf{方法与条件}\quad 对 \(x\ne0,1,2\) 有 \(|x|>0\)，可写成 \(f(x)=\mathrm e^{H(x)}\)，其中 \(H(x)=\frac{\ln|x|}{(1-x)(x-2)}\)。只有这三个点可能间断。先求 \(H\) 的单侧或双侧极限，再由指数函数确定 \(f\) 的极限。

\textbf{推导}\quad 当 \(x\to0\) 时，\(\ln|x|\to-\infty\)，分母趋于 \(-2\)，所以 \(H\to+\infty\)，\(f\to+\infty\)，属于第二类。

当 \(x\to1\) 时，附近 \(x>0\)，将两个负号抵消，得到
\[
H(x)=\frac{\ln x}{x-1}\frac{1}{2-x}\longrightarrow1,
\]
这里使用 \(\frac{\ln x}{x-1}\to1\)。故 \(f(x)\to\mathrm e\)，但原式在 \(1\) 无定义，属于可去间断点，也是第一类。

当 \(x\to2^-\) 时，分子趋于正数 \(\ln2\)，分母 \((1-x)(x-2)\to0^+\)，所以 \(H\to+\infty\)、\(f\to+\infty\)；当 \(x\to2^+\) 时，分母趋 \(0^-\)，故 \(H\to-\infty\)、\(f\to0\)。有一侧无穷，仍是第二类。

\textbf{结论与检查}\quad 第一类间断点只有 \(x=1\)，共 \(\boxed{1}\) 个，选 C。指数无界并不单独决定幂的极限，\(x=1\) 处底数趋 \(1\) 的抵消已通过对数计算处理。
\end{QuestionSolution}

\SectionHeading{2.1}{一元函数微分学及其应用}{导数的概念}{定理与要点\quad 教材 PDF 第 108 页起\hfill 本节 8 题}
\pdfbookmark[1]{一元函数微分学及其应用}{chapter-2.1}
\pdfbookmark[2]{2.1 导数的概念}{section-2.1}
\Theory{定理 1.1}{115}{可导推出连续}{函数在某点可导，则在该点连续。}
\Theory{要点}{108--115}{导数定义与几何意义}{差商的左右极限须存在且相等。有限导数给出切线斜率，非零时法线斜率为其负倒数。}
% Q 2010 3
\QuestionPage{2.1 导数的概念}{选择题}{2010 年第 3 题}{1}{曲线相切}{048 / 234}{**}
曲线$y = x^2$与曲线$y = a\ln x$（$a \ne 0$）相切，则$a =$\pickout{}
\begin{abcd}
\item $4\mathrm{e}$．
\item $3\mathrm{e}$．
\item $2\mathrm{e}$．
\item $\mathrm{e}$．
\end{abcd}
\QuestionEnd

\begin{QuestionSolution}{2010 年第 3 题}
\textbf{已知与目标}\quad 两条曲线相切，要求参数 \(a\)。对数要求切点横坐标 \(x_0>0\)。

\textbf{方法与条件}\quad 相切表示两条曲线经过同一点，且在该点有同一条切线。两函数在 \(x_0>0\) 时都可导，切线斜率就是导数，所以要同时满足函数值相等、导数相等。

\textbf{推导}\quad 分别求导得 \((x^2)'=2x\)、\((a\ln x)'=\frac{a}{x}\)，因此
\[
x_0^2=a\ln x_0,\qquad 2x_0=\frac{a}{x_0}.
\]
由第二式得 \(a=2x_0^2\)。代入第一式并除以非零的 \(x_0^2\)，得 \(1=2\ln x_0\)，故 \(x_0=\sqrt{\mathrm{e}}\)，\(a=2\mathrm{e}\)。

\textbf{结论与检查}\quad \(\boxed{a=2\mathrm{e}}\)，选 C。代入时两曲线的纵坐标都是 \(\mathrm{e}\)，斜率都是 \(2\sqrt{\mathrm{e}}\)，且 \(a\ne0\)，符合条件。
\end{QuestionSolution}

% Q 2011 2
\QuestionPage{2.1 导数的概念}{选择题}{2011 年第 2 题}{1}{导数定义与极限}{049 / 234}{**}
已知$f(x)$在$x = 0$处可导，且$f(0) = 0$，
则$\lim_{x \to 0} \frac{x^2 f(x) - 2f\left(x^3 \right)}{x^3} =$\pickout{}
\begin{abcd}
\item $- 2f'(0)$．
\item $- f'(0)$．
\item $f'(0)$．
\item $0$．
\end{abcd}
\QuestionEnd

\begin{QuestionSolution}{2011 年第 2 题}
\textbf{已知与目标}\quad 已知 \(f(0)=0\)，且 \(f'(0)\) 存在，要求含 \(f(x)\) 与 \(f(x^3)\) 的极限。

\textbf{方法与条件}\quad 差商是函数值的改变量除以自变量的改变量。导数定义给出 \(\lim_{h\to0}\frac{f(h)-f(0)}{h}=f'(0)\)；只需把两项分别整理成这种形式。

\textbf{推导}\quad 对 \(x\ne0\)，利用 \(f(0)=0\)，原商为
\[
\frac{f(x)-f(0)}x-2\frac{f(x^3)-f(0)}{x^3}.
\]
当 \(x\to0\) 时，\(x\to0\)、\(x^3\to0\)，两项差商分别趋于 \(f'(0)\)，所以原极限等于 \(f'(0)-2f'(0)=-f'(0)\)。

\textbf{结论与检查}\quad 所求为 \(\boxed{-f'(0)}\)，选 B。这里没有对 \(f\) 求更高阶导数，只用了题目给出的零点可导性。
\end{QuestionSolution}

% Q 2021 2
\QuestionPage{2.1 导数的概念}{选择题}{2021 年第 2 题}{1}{分段函数的可导性与极值}{051 / 234}{**}
函数
\[f(x)=
\begin{cases}
\dfrac{\mathrm{e}^x-1}{x}, & x\ne 0,\\
1, & x=0
\end{cases}
\]
在 \(x=0\) 处
\choices{连续且取极大值}{连续且取极小值}{可导且导数为 \(0\)}{可导且导数不为 \(0\)}
\QuestionEnd

\begin{QuestionSolution}{2021 年第 2 题}
\textbf{已知与目标}\quad 已知零点的函数值为 \(1\)，要求判断其连续性、可导性及是否有极值。

\textbf{方法与条件}\quad 可导性用差商定义判断。Taylor 展开用多项式描述函数在零点附近的变化；\(o(x^2)\) 表示除以 \(x^2\) 后趋零的余项。本题差商还要再除以 \(x\)，所以指数展开须保留二次项。

\textbf{推导}\quad 由 \(\mathrm{e}^x-1=x+\frac{x^2}{2}+o(x^2)\)，对 \(x\ne0\) 有
\[
f(x)=1+\frac{x}{2}+o(x),\qquad
\frac{f(x)-f(0)}{x}=\frac12+o(1)\longrightarrow\frac12.
\]
因此 \(f(x)\to1=f(0)\)，并且 \(f'(0)=\frac12\)。

\textbf{结论与检查}\quad 选 \(\boxed{\mathrm D}\)。导数非零，可导函数在内部极值点的导数必须为零，所以零点也不是极大值点或极小值点。
\end{QuestionSolution}

% Q 2009 9
\QuestionPage{2.1 导数的概念}{填空题}{2009 年第 9 题}{2}{切线方程}{053 / 234}{**}
曲线$\begin{cases}
  x = \int_0^{1- t} \mathrm{e}^{- u^2} \,\mathrm{d} u \\
  y = t^2\ln(2 - t^2)
\end{cases}$在$(0,0)$处的切线方程为 \blank{}．
\QuestionEnd

\begin{QuestionSolution}{2009 年第 9 题}
\textbf{已知与目标}\quad 曲线以参数 \(t\) 表示，要求经过原点的切线。

\textbf{方法与条件}\quad 先找原点对应的参数，再用 \(\frac{\mathrm dy}{\mathrm dx}=\frac{y'(t)}{x'(t)}\) 求斜率；该公式要求对应参数处 \(x'(t)\ne0\)。变上限积分的导数等于上限处被积函数乘以上限的导数。

\textbf{推导}\quad 被积函数 \(\mathrm e^{-u^2}\) 恒正，积分只有上下限相等时才为零。因此 \(x=0\) 要求 \(1-t=0\)，即 \(t=1\)；此时 \(y=\ln1=0\)。对参数求导得
\[
x'=-\mathrm{e}^{-(1-t)^2},\qquad
y'=2t\ln(2-t^2)-\frac{2t^3}{2-t^2}.
\]
在 \(t=1\) 时，\(x'=-\mathrm e^0=-1\)，\(y'=2\ln1-\frac{2}{1}=-2\)，故切线斜率为 \(\frac{-2}{-1}=2\)。

\textbf{结论与检查}\quad 由点斜式得 \(y-0=2(x-0)\)，即 \(\boxed{y=2x}\)。参数处 \(2-t^2=1>0\)，对数有意义，且 \(x'\ne0\)，求导条件满足。
\end{QuestionSolution}

% Q 2018 2
\QuestionPage{2.1 导数的概念}{选择题}{2018 年第 2 题}{1}{原点处的可导性}{050 / 234}{***}
下列函数中，在$x = 0$处不可导的是\pickout{}
\begin{abcd}
\item $f(x) = |x|\sin |x|$．
\item $f(x) = |x|\sin \sqrt{|x|}$．
\item $f(x) = \cos |x|$．
\item $f(x) = \cos \sqrt{|x|}$．
\end{abcd}
\QuestionEnd

\begin{QuestionSolution}{2018 年第 2 题}
\textbf{已知与目标}\quad 要在四个含绝对值或根号的函数中找出零点不可导者。

\textbf{方法与条件}\quad 零点可导要求 \(\frac{f(x)-f(0)}{x}\) 的左、右极限存在且相等。分段表达式在零点不能直接套用通常的求导公式，需回到这个定义。

\textbf{推导}\quad A、B 的 \(f(0)=0\)，差商分别为 \(\frac{|x|}{x}\sin|x|\)、\(\frac{|x|}{x}\sin\sqrt{|x|}\)。由于 \(\left|\frac{|x|}{x}\right|=1\)，而两个正弦因子趋零，差商都趋于零。C 中 \(\cos|x|=\cos x\)，且 \(f(0)=1\)，故差商为 \(\frac{\cos x-1}{x}\sim-\frac{x}{2}\to0\)。D 中 \(f(0)=1\)，利用 \(\cos u-1=-\frac{u^2}{2}+o(u^2)\)，其中 \(o(u^2)\) 表示除以 \(u^2\) 后趋零的余项，得
\[
\frac{\cos\sqrt{|x|}-1}{x}=-\frac{|x|}{2x}+o(1).
\]
当 \(x<0\) 时 \(\frac{|x|}{x}=-1\)，当 \(x>0\) 时该比值为 \(1\)，所以 D 的左、右导数分别为 \(\frac{1}{2}\)、\(-\frac{1}{2}\)。

\textbf{结论与检查}\quad 选 \(\boxed{\mathrm D}\)。A、B、C 在零点的导数均为零；D 的两个单侧导数不相等，因此不能组成导数。
\end{QuestionSolution}

% Q 2025 7
\QuestionPage{2.1 导数的概念}{选择题}{2025 年第 7 题}{2}{可导的等价条件}{052 / 234}{***}
设函数 \(f\left( x \right)\) 连续，给出下列四个条件：
\par\circnum{1} \(\lim_{x \to 0} \frac{\left| f\left( x \right) \right| - f\left( 0 \right)}{x}\) 存在；
\par\circnum{2} \(\lim_{x \to 0} \frac{f\left( x \right) - \left| f\left( 0 \right) \right|}{x}\) 存在；
\par\circnum{3} \(\lim_{x \to 0} \frac{\left| f\left( x \right) \right|}{x}\) 存在；
\par\circnum{4} \(\lim_{x \to 0} \frac{\left| f\left( x \right) \right| - \left| f\left( 0 \right) \right|}{x}\) 存在.
其中能得到 “\(f\left( x \right)\) 在 \(x = 0\) 处可导” 的条件的个数是
\choices{1}{2}{3}{4}
\QuestionEnd

\begin{QuestionSolution}{2025 年第 7 题}
\textbf{已知与目标}\quad 已知 \(f\) 连续，要求数出四个条件中能保证零点可导的条件。

\textbf{方法与条件}\quad 可导要求有限差商极限存在。绝对值可能改变符号，因此按 \(f(0)=0\) 和 \(f(0)\ne0\) 分类。连续性保证：若 \(f(0)\ne0\)，则零点附近 \(f(x)\) 与 \(f(0)\) 同号。

\textbf{推导}\quad 若 \(f(0)=0\)，第 1、3、4 个条件都成为 \(\frac{|f(x)|}{x}\) 的有限极限存在。该商在右侧非负、左侧非正，故共同极限只能为零。由 \(\left|\frac{f(x)}{x}\right|=\left|\frac{|f(x)|}{x}\right|\to0\)，可得 \(f'(0)=0\)。第 2 个条件此时直接是导数定义。

若 \(f(0)\ne0\)，有限商极限要求分子趋零，因而第 1、2 个条件都要求 \(|f(0)|-f(0)=0\)，即 \(f(0)>0\)。零点附近此时有 \(f(x)>0\)，两个差商都等于 \(\frac{f(x)-f(0)}{x}\)，从而可导。第 3 个条件要求 \(|f(x)|\to0\)，由连续性必有 \(f(0)=0\)，已经在上一种情况处理。第 4 个条件在 \(f(0)>0\) 时等于原函数差商，在 \(f(0)<0\) 时等于其相反数；极限存在都能推出原函数可导。

\textbf{结论与检查}\quad 四个条件各自都充分，故共有 \(\boxed{4}\) 个，选 D。对第 3 个条件已说明它本身会强制 \(f(0)=0\)，没有遗漏非零函数值的情形。
\end{QuestionSolution}

% Q 2022 17
\QuestionPage{2.1 导数的概念}{解答题}{2022 年第 17 题}{2}{由极限求导数}{054 / 234}{***}
已知函数 \(f\left( x \right)\) 在 \(x = 1\) 处可导，且
\(\lim_{x \to 0} \frac{f\left( \mathrm{e}^{x^{2}} \right) - 3f\left( 1 + \sin^{2}x \right)}{x^{2}} = 2\)
求 \(f^{\prime}\left( 1 \right)\).
\QuestionEnd

\begin{QuestionSolution}{2022 年第 17 题}
\textbf{已知与目标}\quad 已知 \(f\) 在 \(1\) 处可导，以及一个商的极限等于 \(2\)，要求 \(f'(1)\)。

\textbf{方法与条件}\quad 有限极限乘上趋零的分母仍趋零，故先用分子极限求 \(f(1)\)。再用可导的等价表示 \(f(1+h)=f(1)+f'(1)h+o(h)\)，其中余项除以 \(h\) 趋零。

\textbf{推导}\quad 可导推出连续，分子趋于 \(f(1)-3f(1)=-2f(1)\)，而分母 \(x^2\to0\)，因此 \(f(1)=0\)。设 \(u=\mathrm{e}^{x^2}-1\)、\(v=\sin^2x\)。由基本极限，
\[
\frac{u}{x^2}=\frac{\mathrm e^{x^2}-1}{x^2}\to1,
\qquad \frac{v}{x^2}=\left(\frac{\sin x}{x}\right)^2\to1.
\]
因而 \(o(u)\)、\(o(v)\) 除以 \(x^2\) 也都趋零，代入可导表示得
\[
2=\lim_{x\to0}\frac{f(1+u)-3f(1+v)}{x^2}
=(1-3)f'(1).
\]
解得 \(f'(1)=-1\)。

\textbf{结论与检查}\quad \(\boxed{f'(1)=-1}\)。代回主要项得 \((-1)\times(1-3)=2\)，与给定极限相符；全过程只要求 \(f\) 在 \(1\) 处可导。
\end{QuestionSolution}

% Q 2025 18
\QuestionPage{2.1 导数的概念}{解答题}{2025 年第 18 题}{3}{由极限判断可导性}{055 / 234}{***}
（本题满分 12 分）
设函数 \(f\left( x \right)\) 在 \(x = 0\) 处连续，且 \(\lim_{x \to 0} \frac{x f\left( x \right) - \mathrm{e}^{2 \sin x} + 1}{\ln \left( 1 + x \right) + \ln \left( 1 - x \right)} = -3\)，证明 \(f\left( x \right)\) 在 \(x = 0\) 处可导，并求 \(f^{\prime}\left( 0 \right)\).
\QuestionEnd

\begin{QuestionSolution}{2025 年第 18 题}
\textbf{已知与目标}\quad 只知道 \(f\) 在零点连续，要从题给极限证明可导并求 \(f'(0)\)。

\textbf{方法与条件}\quad 不能预先假定 \(f\) 可导。先把极限转成函数增量的等式，再得到 \(f(x)=f(0)+Lx+o(x)\)，即可由差商定义证明导数为 \(L\)。这里 \(o(x^2)\) 表示除以 \(x^2\) 后趋零的余项。

\textbf{推导}\quad 在 \(|x|<1\) 时两个对数有意义，且
\(\ln(1+x)+\ln(1-x)=\ln(1-x^2)=-x^2+o(x^2)\)。题给商趋于 \(-3\)，故其分子等于 \((-3+o(1))(-x^2+o(x^2))\)，即
\[
xf(x)-\mathrm{e}^{2\sin x}+1=3x^2+o(x^2).
\]
由 \(\sin x=x+o(x^2)\)，得
\[
\mathrm{e}^{2\sin x}=1+2\sin x+\frac{(2\sin x)^2}{2}+o(x^2)
=1+2x+2x^2+o(x^2).
\]
代入上式，得到 \(xf(x)=2x+5x^2+o(x^2)\)。对 \(x\ne0\) 除以 \(x\)，得 \(f(x)=2+5x+o(x)\)。令 \(x\to0\)，结合连续性可知 \(f(0)=2\)，所以
\[\frac{f(x)-f(0)}{x}=5+o(1)\to5.\]

\textbf{结论与检查}\quad 零点的双侧差商极限存在，故 \(f\) 在零点可导，且 \(\boxed{f'(0)=5}\)。将 \(f(x)=2+5x+o(x)\) 代回，分子的二次项为 \(3x^2\)，与分母的 \(-x^2\) 相除确实趋于 \(-3\)。
\end{QuestionSolution}

\SectionHeading{2.2}{一元函数微分学及其应用}{求导的基本法则}{定理与要点\quad 教材 PDF 第 121 页起\hfill 本节 28 题}
\pdfbookmark[2]{2.2 求导的基本法则}{section-2.2}
\Theory{定理 2.1}{121}{导数的有理运算法则}{可导函数的和、差、积、商可导。乘积用乘积法则，商须分母非零。}
\Theory{定理 2.2}{123}{链式法则}{内外两层在相应点可导时，复合函数的导数等于外层导数乘内层导数。}
\Theory{定理 2.3}{125}{反函数求导法则}{严格单调连续函数在对应点可导且导数非零时，反函数导数等于原导数的倒数。}
\Theory{定理 2.4}{130}{高阶导数的运算}{\(n\) 阶导数满足线性法则。乘积的 \(n\) 阶导数按 Leibniz 公式作二项式系数加权求和。}
\Theory{定理 2.5}{132}{参数方程求导法则}{当 \(x'(t)\ne0\) 时，\(\frac{\mathrm{d}y}{\mathrm{d}x}=\frac{y'(t)}{x'(t)}\)。将所得一阶导数对 \(t\) 求导，再除以 \(x'(t)\) 即得二阶导数。}
\Theory{方法}{130--135}{隐函数与相关变化率}{先按变量依赖关系求导，再代入给定点或时刻。隐函数项需乘上相应的导数。}
% Q 2012 2
\QuestionPage{2.2 求导的基本法则}{选择题}{2012 年第 2 题}{1}{乘积求导}{056 / 234}{**}
设函数$f(x) = (\mathrm{e}^x - 1)(\mathrm{e}^{2x} - 2) \cdots(\mathrm{e}^{nx} - n)$，其中$n$为正整数，则$f'(0) =$\pickout{}
\begin{abcd}
\item $(- 1)^{n - 1}(n - 1)!$．
\item $(- 1)^n(n - 1)!$．
\item $(- 1)^{n - 1} n!$．
\item $(- 1)^n n!$．
\end{abcd}
\QuestionEnd

\begin{QuestionSolution}{2012 年第 2 题}
\textbf{已知与目标}\quad \(f\) 是 \(n\) 个可导因子的乘积，要计算 \(f'(0)\)。

\textbf{方法与条件}\quad 乘积法则为 \((uv)'=u'v+uv'\)。第一因子 \(\mathrm e^x-1\) 在零点为零，因此将它单独分出，可使第二项在代入时消失。

\textbf{推导}\quad 写 \(f(x)=(\mathrm{e}^x-1)P(x)\)，其中 \(P(x)=\prod_{k=2}^n(\mathrm{e}^{kx}-k)\)；乘积符号表示将各因子相乘，\(n=1\) 时没有剩余因子，取 \(P=1\)。于是 \(f'=\mathrm e^xP+(\mathrm e^x-1)P'\)，从而
\[
f'(0)=P(0)=\prod_{k=2}^n(1-k)=(-1)^{n-1}(n-1)!.
\]
其中各因子依次为 \(-1,-2,\ldots,-(n-1)\)，故共有 \(n-1\) 个负号。

\textbf{结论与检查}\quad \(\boxed{f'(0)=(-1)^{n-1}(n-1)!}\)，选 A。\(n=1\) 时结果为 \(0!=1\)，也与 \((\mathrm e^x-1)'_{x=0}=1\) 一致。
\end{QuestionSolution}

% Q 2013 2
\QuestionPage{2.2 求导的基本法则}{选择题}{2013 年第 2 题}{1}{隐函数求导与极限}{057 / 234}{**}
设函数$y = f(x)$由方程$\cos (xy) + \ln y - x = 1$确定，
则$\lim_{n \to \infty} n\left[f\left(\frac{2}{n} \right) - 1 \right] =$\pickout{}
\begin{abcd}
\item $2$．
\item $1$．
\item $- 1$．
\item $- 2$．
\end{abcd}
\QuestionEnd

\begin{QuestionSolution}{2013 年第 2 题}
\textbf{已知与目标}\quad \(y=f(x)\) 由含 \(x,y\) 的方程确定，要计算 \(n[f(\frac{2}{n})-1]\) 的极限。

\textbf{方法与条件}\quad 隐函数求导是对方程两侧求导，并将 \(y\) 视为 \(x\) 的函数。先求 \(f(0)\) 和 \(f'(0)\)，再将原式改写为导数定义。方程左侧对 \(y\) 的偏导在 \((0,1)\) 为 \(1\ne0\)，保证这里存在可导的隐函数。

\textbf{推导}\quad 令 \(x=0\)，得 \(1+\ln y=1\)，故 \(f(0)=1\)。因为 \((xy)'=y+xy'\)、\((\ln y)'=\frac{y'}{y}\)，对方程求导得
\[
-\sin(xy)(y+xy')+\frac{y'}y-1=0.
\]
在 \((0,1)\) 代入得 \(0+y'-1=0\)，即 \(f'(0)=1\)。又 \(\frac{2}{n}\to0\)，所以
\[
n\left[f\left(\frac{2}{n}\right)-1\right]
=2\frac{f(\frac{2}{n})-f(0)}{\frac{2}{n}}\longrightarrow2f'(0)=2.
\]

\textbf{结论与检查}\quad 极限为 \(\boxed{2}\)，选 A。改写时增量为 \(\frac{2}{n}\)，因而需要保留前面的系数 \(2\)。
\end{QuestionSolution}

% Q 2021 3
\QuestionPage{2.2 求导的基本法则}{选择题}{2021 年第 3 题}{1}{圆柱体的相关变化率}{058 / 234}{**}
有一圆柱体底面半径与高随时间变化的速率分别为 \(2\,\mathrm{cm}\,\mathrm{s}^{-1}\)、\(-3\,\mathrm{cm}\,\mathrm{s}^{-1}\)。当底面半径为 \(10\,\mathrm{cm}\)，高为 \(5\,\mathrm{cm}\) 时，圆柱体的体积与表面积随时间变化的速率分别为
\choices{\(125\pi\,\mathrm{cm}^{3}\,\mathrm{s}^{-1}\)，\(40\pi\,\mathrm{cm}^{2}\,\mathrm{s}^{-1}\)}{\(125\pi\,\mathrm{cm}^{3}\,\mathrm{s}^{-1}\)，\(-40\pi\,\mathrm{cm}^{2}\,\mathrm{s}^{-1}\)}{\(-100\pi\,\mathrm{cm}^{3}\,\mathrm{s}^{-1}\)，\(40\pi\,\mathrm{cm}^{2}\,\mathrm{s}^{-1}\)}{\(-100\pi\,\mathrm{cm}^{3}\,\mathrm{s}^{-1}\)，\(-40\pi\,\mathrm{cm}^{2}\,\mathrm{s}^{-1}\)}
\QuestionEnd

\begin{QuestionSolution}{2021 年第 3 题}
\textbf{已知与目标}\quad 在所问时刻 \(r=10,h=5\)，半径变化率 \(r'=2\)，高度变化率 \(h'=-3\)，要求体积和表面积对时间的导数。

\textbf{方法与条件}\quad 相关变化率指几个随时间变化的量通过几何关系联系起来的导数。圆柱体的 \(V=\pi r^2h\)、\(S=2\pi r^2+2\pi rh\)，表面积含两个底面；求导时 \(r,h\) 都依赖时间。

\textbf{推导}\quad 由乘积法则和链式法则，
\[
\begin{aligned}
V'&=\pi(2rr'h+r^2h'),\\
S'&=2\pi(2rr'+r'h+rh').
\end{aligned}
\]
逐项代入可得
\[
V'=\pi(2\cdot10\cdot2\cdot5+10^2\cdot(-3))=-100\pi,
\qquad S'=2\pi(2\cdot10\cdot2+2\cdot5+10\cdot(-3))=40\pi.
\]

\textbf{结论与检查}\quad 速率依次为 \(\boxed{-100\pi\,\mathrm{cm}^3\,\mathrm{s}^{-1}}\)、\(\boxed{40\pi\,\mathrm{cm}^2\,\mathrm{s}^{-1}}\)，选 C。体积变化率为负表示此时体积减小；两个答案的单位分别为体积和面积除以时间。
\end{QuestionSolution}

% Q 2024 2
\QuestionPage{2.2 求导的基本法则}{选择题}{2024 年第 2 题}{1}{参数方程求导与导数定义}{060 / 234}{**}
已知
\[\left\{
\begin{aligned}
x &= 1+t^3,\\
y &= \mathrm{e}^{t^2},
\end{aligned}
\right.
\]
则
\[\lim_{x\to+\infty}x\!\left[f\!\left(2+\frac{2}{x}\right)-f(2)\right]
\]
等于（ ）
\choices{\(2\mathrm{e}\)}{\(\frac{4}{3}\mathrm{e}\)}{\(\frac{2}{3}\mathrm{e}\)}{\(\frac{\mathrm{e}}{3}\)}
\QuestionEnd

\begin{QuestionSolution}{2024 年第 2 题}
\textbf{已知与目标}\quad 参数方程定义 \(y=f(x)\)，要计算 \(f\) 在自变量 \(2\) 附近的增量极限。

\textbf{方法与条件}\quad 令 \(h=\frac{2}{x}\)，则极限中的 \(x\to+\infty\) 对应 \(h\to0^+\)，原式成为 \(2\) 倍右差商。先用参数导数比求 \(f'(2)\)，需检查 \(\frac{\mathrm dx}{\mathrm dt}\ne0\)。

\textbf{推导}\quad 参数坐标 \(x=2\) 要求 \(1+t^3=2\)，即 \(t=1\)。此时 \(\frac{\mathrm dx}{\mathrm dt}=3t^2=3\ne0\)，所以
\[
\left.\frac{\mathrm{d}y}{\mathrm{d}x}\right|_{t=1}
=\left.\frac{2t \mathrm{e}^{t^2}}{3t^2}\right|_{t=1}=\frac{2\mathrm{e}}{3}.
\]
又 \(x=\frac{2}{h}\)，原式写为 \(2\frac{f(2+h)-f(2)}h\)，故极限为 \(2\cdot\frac{2\mathrm e}{3}=\frac{4\mathrm e}{3}\)。

\textbf{结论与检查}\quad \(\boxed{\frac{4\mathrm{e}}{3}}\)，选 B。由于 \(f'(2)\) 存在，右差商极限等于该导数；系数 \(2\) 来源于代换 \(h=\frac2x\)。
\end{QuestionSolution}

% Q 2026 6
\QuestionPage{2.2 求导的基本法则}{选择题}{2026 年第 6 题}{1}{反函数与变上限积分求导}{061 / 234}{**}
已知函数
\[f(x)=\int_1^{x^3}\frac{\mathrm{e}^t}{1+t^2}\,\mathrm{d}t,
\]
\(f\) 的反函数为 \(g\)，则
\choices{\(g(0)=1,\ g'(0)=\dfrac{3}{2\mathrm{e}}\)}{\(g(0)=1,\ g'(0)=\dfrac{2}{3\mathrm{e}}\)}{\(g(1)=0,\ g'(1)=\dfrac{3}{2\mathrm{e}}\)}{\(g(1)=0,\ g'(1)=\dfrac{2}{3\mathrm{e}}\)}
\QuestionEnd

\begin{QuestionSolution}{2026 年第 6 题}
\textbf{已知与目标}\quad \(g\) 是积分函数 \(f\) 的反函数，要求相应函数值和导数值。

\textbf{方法与条件}\quad 反函数把原函数的输入与输出互换：\(f(a)=b\) 等价于 \(g(b)=a\)。若原函数在该点可导且 \(f'(a)\ne0\)，则 \(g'(b)=\frac1{f'(a)}\)。积分求导还要乘上上限 \(x^3\) 的导数。

\textbf{推导}\quad 被积函数恒正，积分随上限严格增加，且 \(x^3\) 严格递增，故 \(f\) 严格递增，有反函数。\(f(1)=\int_1^1\frac{\mathrm e^t}{1+t^2}\,\mathrm dt=0\)，故 \(g(0)=1\)。由微积分基本定理及链式法则，
\[
f'(x)=3x^2\frac{\mathrm{e}^{x^3}}{1+x^6},\qquad f'(1)=\frac{3\mathrm{e}}{2}.
\]
\textbf{结论与检查}\quad \(\boxed{g(0)=1,\ g'(0)=\frac{2}{3\mathrm e}}\)，选 B。因为 \(f'(1)=\frac{3\mathrm e}{2}\ne0\)，反函数求导条件满足，两导数的乘积等于 \(1\)。
\end{QuestionSolution}

% Q 2009 12
\QuestionPage{2.2 求导的基本法则}{填空题}{2009 年第 12 题}{2}{隐函数求导}{062 / 234}{**}
设$y = y(x)$是由方程$xy + \mathrm{e}^y = x + 1$确定的隐函数，
则$\frac{\partial{}^2 y}{\partial{} x^2}\Big|_{x = 0} =$ \blank{}．
\QuestionEnd

\begin{QuestionSolution}{2009 年第 12 题}
\textbf{已知与目标}\quad \(y\) 由 \(xy+\mathrm e^y=x+1\) 确定，要求 \(y''(0)\)。

\textbf{方法与条件}\quad 隐函数求导时 \(y\) 随 \(x\) 改变，所以 \((xy)'=y+xy'\)、\((\mathrm e^y)'=\mathrm e^yy'\)。二次求导要保留 \(y'\) 的乘积项，最后才代入点值。方程对 \(y\) 的偏导在 \((0,0)\) 为 \(1\)，可在此求导。

\textbf{推导}\quad \(x=0\) 时 \(\mathrm{e}^y=1\)，故 \(y(0)=0\)。一次求导得 \(y+(x+\mathrm{e}^y)y'=1\)，代入得 \(0+(0+1)y'(0)=1\)，即 \(y'(0)=1\)。再次求导得
\[
2y'+\mathrm{e}^y(y')^2+(x+\mathrm{e}^y)y''=0.
\]
代入 \(x=0,y=0,y'=1\)，得 \(2+1+y''(0)=0\)。

\textbf{结论与检查}\quad \(\boxed{y''(0)=-3}\)。二次求导式中的两个 \(y'\) 分别来自 \(y\) 和 \(xy'\)，而 \(\mathrm e^y(y')^2\) 来自指数复合函数，均不能遗漏。
\end{QuestionSolution}

% Q 2010 11
\QuestionPage{2.2 求导的基本法则}{填空题}{2010 年第 11 题}{2}{对数函数的高阶导数}{063 / 234}{**}
函数$y = \ln(1 - 2x)$在$x = 0$处的$n$阶导数$y^{(n)}(0) =$ \blank{}．
\QuestionEnd

\begin{QuestionSolution}{2010 年第 11 题}
\textbf{已知与目标}\quad 对 \(y=\ln(1-2x)\)，要求正整数 \(n\) 阶导数在零点的值。

\textbf{方法与条件}\quad 高阶导数是反复求导。对数要求 \(1-2x>0\)，零点附近满足此条件。先求前几阶，再用链式法则验证一般式。

\textbf{推导}\quad 直接得到
\[
y'=-2(1-2x)^{-1},\qquad y''=-4(1-2x)^{-2},\qquad y'''=-16(1-2x)^{-3}.
\]
一般地，\(y^{(n)}=-2^n(n-1)!(1-2x)^{-n}\)。若该式在 \(n\) 阶成立，再求导时额外乘 \((-n)(-2)=2n\)，恰得到 \(-2^{n+1}n!(1-2x)^{-(n+1)}\)，所以公式对所有 \(n\ge1\) 成立。

\textbf{结论与检查}\quad 代入 \(x=0\)，\(\boxed{y^{(n)}(0)=-2^n(n-1)!}\)。\(n=1,2\) 时分别为 \(-2,-4\)，与直接求导一致。
\end{QuestionSolution}

% Q 2010 13
\QuestionPage{2.2 求导的基本法则}{填空题}{2010 年第 13 题}{2}{相关变化率}{064 / 234}{**}
已知一个长方形的长$l$以$2\,\mathrm{cm}\,\mathrm{s}^{-1}$的速率增加，宽$w$以$3\,\mathrm{cm}\,\mathrm{s}^{-1}$的速率增加．
则当$l = 12\text{cm}$，$w = 5\text{cm}$时，它的对角线增加的速率为 \blank{}．
\QuestionEnd

\begin{QuestionSolution}{2010 年第 13 题}
\textbf{已知与目标}\quad 长、宽对时间的导数分别为 \(l'=2,w'=3\)，要求 \(l=12,w=5\) 时对角线的增长速率。

\textbf{方法与条件}\quad 用勾股定理联系三条边，再对时间求导。各边都是时间函数，所以平方求导还须乘上对应变化率。

\textbf{推导}\quad 设对角线为 \(d>0\)，则 \(d^2=l^2+w^2\)，求导得 \(2dd'=2ll'+2ww'\)，所以 \(d'=\frac{ll'+ww'}d\)。在所问时刻 \(d=\sqrt{12^2+5^2}=13\)，代入得
\[
d'=\frac{12\cdot2+5\cdot3}{13}=\boxed{3\,\mathrm{cm}\,\mathrm{s}^{-1}}.
\]
\textbf{结论与检查}\quad 对角线以 \(\boxed{3\,\mathrm{cm}\,\mathrm{s}^{-1}}\) 的速率增加。结果为正，符合长、宽均增加的情形；除以的对角线长 \(13\) 非零。
\end{QuestionSolution}

% Q 2012 9
\QuestionPage{2.2 求导的基本法则}{填空题}{2012 年第 9 题}{2}{隐函数求导}{065 / 234}{**}
设$y = y(x)$是由方程$x^2 - y + 1 = \mathrm{e}^y$所确定的隐函数，
则$\left. \frac{\mathrm{d}^2y}{\mathrm{d} x^2} \right|_{x = 0} =$ \blank{}．
\QuestionEnd

\begin{QuestionSolution}{2012 年第 9 题}
\textbf{已知与目标}\quad 已知 \(y+\mathrm e^y=x^2+1\)，要求 \(y''(0)\)。

\textbf{方法与条件}\quad 对原方程连续求两次导数。因 \(y\) 随 \(x\) 变化，指数求导须乘 \(y'\)；二次求导前先确定 \(y(0)\)、\(y'(0)\)。方程关于 \(y\) 的导数 \(1+\mathrm e^y>0\)，保证隐函数可导。

\textbf{推导}\quad \(x=0\) 时 \(y+\mathrm e^y=1\)。左侧严格递增，在 \(y=0\) 取值 \(1\)，所以唯一解为 \(y(0)=0\)。一次求导得 \((1+\mathrm{e}^y)y'=2x\)，代入零点得 \(2y'(0)=0\)，即 \(y'(0)=0\)。再求导得
\[
(1+\mathrm{e}^y)y''+\mathrm{e}^y(y')^2=2.
\]
代入 \(y=0,y'=0\)，得 \(2y''(0)+0=2\)。

\textbf{结论与检查}\quad \(\boxed{y''(0)=1}\)。二次导数式中的平方项因 \(y'(0)=0\) 消失，系数 \(1+\mathrm e^0=2\) 非零，可合法除去。
\end{QuestionSolution}

% Q 2013 10
\QuestionPage{2.2 求导的基本法则}{填空题}{2013 年第 10 题}{2}{反函数求导}{066 / 234}{**}
设函数$f(x) = \int_{-1}^x \sqrt{1 - \mathrm{e}^t} \,\mathrm{d} t$，则$y = f(x)$的反函数$x = f^{-1}(y)$
在$y = 0$处的导数$\left. \frac{\mathrm{d} x}{\mathrm{d} y} \right|_{y = 0} =$ \blank{}．
\QuestionEnd

\begin{QuestionSolution}{2013 年第 10 题}
\textbf{已知与目标}\quad 原函数是变上限积分，要计算反函数在输出值 \(y=0\) 处的导数。

\textbf{方法与条件}\quad 先确定 \(f(a)=0\) 的输入 \(a\)，再用 \(\frac{\mathrm dx}{\mathrm dy}=\frac1{f'(a)}\)。该公式要求 \(f\) 在对应点可导且导数非零。

\textbf{推导}\quad 根号要求 \(x\le0\)。在 \(x<0\) 时被积函数为正，所以 \(f\) 严格递增；当上下限相等时 \(f(-1)=0\)，故 \(y=0\) 对应 \(x=-1\)。由微积分基本定理，\(f'(-1)=\sqrt{1-\mathrm{e}^{-1}}>0\)，从而
\[
\left.\frac{\mathrm{d}x}{\mathrm{d}y}\right|_{y=0}
=\frac{1}{f'(-1)}=\boxed{\frac{1}{\sqrt{1-\mathrm{e}^{-1}}}}.
\]
\textbf{结论与检查}\quad 答案为 \(\boxed{\frac{1}{\sqrt{1-\mathrm e^{-1}}}}\)。对应点 \(-1\) 在定义域内部，根号内为正，故倒数有限且为正，与反函数严格递增一致。
\end{QuestionSolution}

% Q 2013 12
\QuestionPage{2.2 求导的基本法则}{填空题}{2013 年第 12 题}{2}{参数曲线的法线}{067 / 234}{**}
曲线$\begin{cases}
  x = \arctan t, \\
  y = \ln \sqrt{1 + t^2}
\end{cases}$上对应于$t = 1$的点处的法线方程为 \blank{}．
\QuestionEnd

\begin{QuestionSolution}{2013 年第 12 题}
\textbf{已知与目标}\quad 曲线用参数 \(t\) 表示，要求 \(t=1\) 处的法线方程。

\textbf{方法与条件}\quad 法线是经过曲线上该点且垂直于切线的直线。切线斜率 \(k=\frac{y'(t)}{x'(t)}\) 有限且非零时，法线斜率为 \(-\frac1k\)。

\textbf{推导}\quad \(t=1\) 时点坐标为 \((\arctan1,\ln\sqrt2)=(\frac{\pi}{4},\frac{1}{2}\ln2)\)。将 \(y=\frac12\ln(1+t^2)\) 求导，得到
\[x'=\frac{1}{1+t^2},\qquad y'=\frac{t}{1+t^2},\qquad k=\frac{y'}{x'}=t.\]
因此该点切线斜率为 \(1\)，法线斜率为 \(-1\)，代入直线点斜式得
\[
\boxed{y-\frac{1}{2}\ln2=-(x-\frac{\pi}{4})}.
\]
\textbf{结论与检查}\quad 上式即所求法线。它经过指定点，且斜率与切线斜率的乘积为 \(-1\)，满足垂直条件。
\end{QuestionSolution}

% Q 2014 10
\QuestionPage{2.2 求导的基本法则}{填空题}{2014 年第 10 题}{2}{奇偶性、周期性与求导}{068 / 234}{**}
设$f(x)$是周期为$4$的可导奇函数，且$f'(x)= 2(x - 1)$, $x \in[0,2]$，则$f(7) =$ \blank{}．
\QuestionEnd

\begin{QuestionSolution}{2014 年第 10 题}
\textbf{已知与目标}\quad 已知 \(f(x+4)=f(x)\)、\(f(-x)=-f(x)\)，以及 \([0,2]\) 上的导数，要求 \(f(7)\)。

\textbf{方法与条件}\quad 周期性允许自变量减去 \(4\) 的整数倍，奇性允许把负数化为正数。奇函数在零点满足 \(f(0)=-f(0)\)，所以 \(f(0)=0\)，可用于确定原函数的常数。

\textbf{推导}\quad 由周期性、奇性，\(f(7)=f(7-8)=f(-1)=-f(1)\)。在 \([0,1]\) 上导数 \(2(x-1)\) 连续，故
\[
f(1)=f(0)+\int_0^1 2(x-1)\,\mathrm{d}x
=0+\left[x^2-2x\right]_0^1=-1.
\]
\textbf{结论与检查}\quad \(\boxed{f(7)=1}\)。也可在已知区间写成 \(f(x)=x^2-2x\)，代入 \(x=1\) 得 \(-1\)，再用奇性取相反数。
\end{QuestionSolution}

% Q 2014 12
\QuestionPage{2.2 求导的基本法则}{填空题}{2014 年第 12 题}{2}{极坐标曲线的切线}{069 / 234}{**}
曲线$L$的极坐标方程为$r = \theta$，
则$L$在点$(r,\theta) = \left(\frac{\pi}{2},\frac{\pi}{2}\right)$处的切线方程为 \blank{}．
\QuestionEnd

\begin{QuestionSolution}{2014 年第 12 题}
\textbf{已知与目标}\quad 极坐标曲线满足 \(r=\theta\)，要求指定点处的直角坐标切线方程。

\textbf{方法与条件}\quad 极坐标中的 \(r\) 是点到原点的距离，\(\theta\) 是方向角。用 \(x=r\cos\theta,y=r\sin\theta\) 转为以 \(\theta\) 为参数的方程，然后求导数比；需检查 \(x'(\theta)\ne0\)。

\textbf{推导}\quad 有 \(x=\theta\cos\theta,y=\theta\sin\theta\)。当 \(\theta=\frac{\pi}{2}\) 时，\(\cos\theta=0,\sin\theta=1\)，故点为 \((0,\frac{\pi}{2})\)。对参数求导并代入，
\[
x'=\cos\theta-\theta\sin\theta=-\frac{\pi}{2},\quad
y'=\sin\theta+\theta\cos\theta=1.
\]
因此斜率为 \(\frac{1}{-\frac{\pi}{2}}=-\frac2\pi\)，点斜式为 \(y-\frac\pi2=-\frac2\pi(x-0)\)。

\textbf{结论与检查}\quad \(\boxed{y=\frac{\pi}{2}-\frac{2x}{\pi}}\)。它经过指定点，且参数处 \(x'=-\frac\pi2\ne0\)，所以导数比确实给出切线斜率。
\end{QuestionSolution}

% Q 2015 9
\QuestionPage{2.2 求导的基本法则}{填空题}{2015 年第 9 题}{2}{隐函数求导}{070 / 234}{**}
设$\begin{cases}
  x = \arctan t, \\
  y = 3t + t^3,
\end{cases}$则$\frac{\mathrm{d}^2y}{\mathrm{d} x^2}\big|_{t = 1} =$ \blank{}．
\QuestionEnd

\begin{QuestionSolution}{2015 年第 9 题}
\textbf{已知与目标}\quad 已知 \(x=\arctan t,y=3t+t^3\)，要求 \(t=1\) 处关于 \(x\) 的二阶导数。

\textbf{方法与条件}\quad 参数曲线的一阶导数为 \(\frac{y_t'}{x_t'}\)。二阶导数需先对这个比值关于 \(t\) 求导，再除以 \(x_t'\)，因为对 \(x\) 求导等于 \(\frac1{x_t'}\frac{\mathrm d}{\mathrm dt}\)。这里 \(x_t'=\frac1{1+t^2}>0\)，公式可用。

\textbf{推导}\quad \(\frac{\mathrm{d}x}{\mathrm{d}t}=(1+t^2)^{-1}\)，\(\frac{\mathrm{d}y}{\mathrm{d}t}=3(1+t^2)\)，故 \(\frac{\mathrm{d}y}{\mathrm{d}x}=3(1+t^2)^2\)。再对 \(t\) 求导得到 \(12t(1+t^2)\)，从而
\[
\frac{\mathrm{d}^2y}{\mathrm{d}x^2}
=\frac{12t(1+t^2)}{\frac{1}{1+t^2}}
=12t(1+t^2)^2.
\]
在 \(t=1\) 时为 \(12\cdot1\cdot(1+1)^2=48\)。

\textbf{结论与检查}\quad \(\boxed{48}\)。对参数求导得到的 \(24\) 还须除以当时的 \(x_t'=\frac12\)，才能得到所求的 \(x\) 二阶导数。
\end{QuestionSolution}

% Q 2015 10
\QuestionPage{2.2 求导的基本法则}{填空题}{2015 年第 10 题}{2}{高阶导数}{071 / 234}{**}
函数$f(x) = x^2 \cdot 2^x$在$x = 0$处的$n$阶导数$f^{(n)}(0) =$ \blank{}．
\QuestionEnd

\begin{QuestionSolution}{2015 年第 10 题}
\textbf{已知与目标}\quad \(f=x^2 2^x\)，要求零点的 \(n\) 阶导数。

\textbf{方法与条件}\quad Leibniz 公式是乘积的高阶求导公式：分配给两个因子的求导次数之和为 \(n\)，并乘相应二项式系数。在零点，\(x^2\) 只有求导两次才非零，所以 \(n\ge2\) 时只剩一项。

\textbf{推导}\quad 因 \((2^x)^{(k)}=(\ln2)^k2^x\)，当 \(n\ge2\) 时有
\[
f^{(n)}(0)=\binom n2\cdot2\cdot(\ln2)^{n-2}\cdot2^0
=n(n-1)(\ln2)^{n-2}.
\]
因此
\[
\boxed{f^{(n)}(0)=n(n-1)(\ln2)^{n-2}\quad(n\ge2)}.
\]
\textbf{结论与检查}\quad \(n\ge2\) 时答案如上；\(n=1\) 时直接求导得 \(f'(0)=0\)。\(n=2\) 时公式给出 \(2\)，与 \(f(x)=x^2+o(x^2)\) 的二次项一致。
\end{QuestionSolution}

% Q 2016 13
\QuestionPage{2.2 求导的基本法则}{填空题}{2016 年第 13 题}{2}{动点距离的相关变化率}{072 / 234}{**}
已知动点$P$在曲线$y = x^3$上运动，记坐标原点与点$P$间的距离为$l$．
若点$P$的横坐标对时间的变化率为常数$v_0$，则当点$P$运动到点$(1,1)$时，$l$对时间的变化率是 \blank{}．
\QuestionEnd

\begin{QuestionSolution}{2016 年第 13 题}
\textbf{已知与目标}\quad 点在 \(y=x^3\) 上运动，且 \(\frac{\mathrm dx}{\mathrm dt}=v_0\)，要求经过 \((1,1)\) 时距离 \(l\) 的变化率。

\textbf{方法与条件}\quad 距离满足 \(l^2=x^2+y^2\)。代入曲线关系可只保留一个随时间变化的坐标，再用链式法则求导；在所问点 \(l>0\)，允许除以 \(l\)。

\textbf{推导}\quad 代入 \(y=x^3\) 得 \(l^2=x^2+x^6\)。对时间求导得
\[
2l\frac{\mathrm{d}l}{\mathrm{d}t}=(2x+6x^5)\frac{\mathrm{d}x}{\mathrm{d}t}.
\]
在 \((1,1)\) 有 \(l=\sqrt2\)，所以 \(2\sqrt2\frac{\mathrm dl}{\mathrm dt}=(2+6)v_0\)，解得 \(\frac{\mathrm dl}{\mathrm dt}=\frac{8v_0}{2\sqrt2}=2\sqrt2v_0\)。

\textbf{结论与检查}\quad \(\boxed{\frac{\mathrm{d}l}{\mathrm{d}t}=2\sqrt2\,v_0}\)。结果与 \(v_0\) 同号：在该点向右运动时离原点更远，向左运动时则更近。
\end{QuestionSolution}

% Q 2017 10
\QuestionPage{2.2 求导的基本法则}{填空题}{2017 年第 10 题}{2}{参数方程求导}{073 / 234}{**}
设函数$y = y(x)$由参数方程$\begin{cases}
  x = t + \mathrm{e}^t \\
  y = \sin t
\end{cases}$确定，则$\left.\frac{\mathrm{d}^2y}{\mathrm{d} x^2}\right|_{t=0} =$ \blank{}．
\QuestionEnd

\begin{QuestionSolution}{2017 年第 10 题}
\textbf{已知与目标}\quad 参数曲线满足 \(x=t+\mathrm e^t,y=\sin t\)，要求 \(t=0\) 时的二阶导数。

\textbf{方法与条件}\quad 当 \(x'(t)\ne0\) 时，\(\frac{\mathrm dy}{\mathrm dx}=\frac{y'(t)}{x'(t)}\)，再对这个商求 \(t\) 导数并除以 \(x'(t)\)，得到 \(\frac{\mathrm d^2y}{\mathrm dx^2}\)。

\textbf{推导}\quad \(x'=1+\mathrm{e}^t>0\)、\(y'=\cos t\)，故 \(\frac{\mathrm{d}y}{\mathrm{d}x}=\frac{\cos t}{1+\mathrm{e}^t}\)。商法则给出的分母为 \((1+\mathrm e^t)^2\)，再除以 \(x'\) 后有
\[
\frac{\mathrm{d}^2y}{\mathrm{d}x^2}
=\frac{-\sin t(1+\mathrm{e}^t)-\mathrm{e}^t\cos t}{(1+\mathrm{e}^t)^3}.
\]
在 \(t=0\) 时，分子为 \(-0\cdot2-1\cdot1=-1\)，分母为 \(2^3=8\)。

\textbf{结论与检查}\quad \(\boxed{-\frac{1}{8}}\)。分母非零，且三次幂来自商求导的平方分母再乘一次 \(x'\)。
\end{QuestionSolution}

% Q 2019 10
\QuestionPage{2.2 求导的基本法则}{填空题}{2019 年第 10 题}{1}{参数曲线切线的截距}{074 / 234}{**}
曲线$\begin{cases}
  x = t - \sin t \\
  y = 1 - \cos t
\end{cases}$在$t = \frac{3\pi}{2}$对应点处的切线在$y$轴的截距为 \blank{}．
\QuestionEnd

\begin{QuestionSolution}{2019 年第 10 题}
\textbf{已知与目标}\quad 要求参数曲线在 \(t=\frac{3\pi}{2}\) 处的切线与 \(y\) 轴交点的纵坐标，即纵截距。

\textbf{方法与条件}\quad 先求切点和切线斜率，再写直线方程，令 \(x=0\) 得纵截距。斜率公式 \(\frac{y'(t)}{x'(t)}\) 需 \(x'(t)\ne0\)。

\textbf{推导}\quad 此时 \(\sin t=-1,\cos t=0\)，切点为 \((\frac{3\pi}{2}+1,1)\)。又 \(x'=1-\cos t=1\)、\(y'=\sin t=-1\)，故斜率为 \(-1\)。切线为
\[y-1=-\left(x-\frac{3\pi}{2}-1\right),\qquad y=-x+2+\frac{3\pi}{2}.\]

\textbf{结论与检查}\quad 令 \(x=0\)，纵截距为 \(\boxed{2+\frac{3\pi}{2}}\)。代入切点横坐标得到 \(y=1\)，直线确实经过该点。
\end{QuestionSolution}

% Q 2020 9
\QuestionPage{2.2 求导的基本法则}{填空题}{2020 年第 9 题}{2}{参数方程的导数}{076 / 234}{**}
设
\[\begin{cases}
x=\sqrt{t^2+1},\\
y=\ln\!\left(t+\sqrt{t^2+1}\right),
\end{cases}
\]
则 \(\left.\dfrac{\mathrm{d}^2y}{\mathrm{d}x^2}\right|_{t=1}=\)
\blank{}
\QuestionEnd

\begin{QuestionSolution}{2020 年第 9 题}
\textbf{已知与目标}\quad 要计算参数曲线在 \(t=1\) 处的 \(\frac{\mathrm d^2y}{\mathrm dx^2}\)。

\textbf{方法与条件}\quad 对参数求导后，\(\frac{\mathrm dy}{\mathrm dx}=\frac{y'}{x'}\)，二阶导数再对 \(t\) 求导并除以 \(x'\)。在 \(t=1\) 附近 \(x'\ne0\)，可以将曲线局部视为 \(y(x)\)。

\textbf{推导}\quad 记 \(s=\sqrt{1+t^2}\)。因 \(s>|t|\)，对数真数 \(t+s>0\)。由链式法则，\(x'=\frac ts\)，\(y'=\frac{1+\frac ts}{t+s}=\frac1s\)。于是
\[
\frac{\mathrm{d}y}{\mathrm{d}x}=\frac{1}{t},\qquad
\frac{\mathrm{d}^2y}{\mathrm{d}x^2}=\frac{-\frac{1}{t^2}}{\frac{t}{s}}=-\frac{s}{t^3}.
\]
\textbf{结论与检查}\quad 代入 \(t=1,s=\sqrt2\)，得 \(\boxed{-\sqrt2}\)。此时 \(x'=\frac1{\sqrt2}\ne0\)，所用除法合法；第二次求导仍须除以 \(x'\)，不能只把 \(\frac1t\) 对 \(t\) 求导。
\end{QuestionSolution}

% Q 2021 12
\QuestionPage{2.2 求导的基本法则}{填空题}{2021 年第 12 题}{2}{参数方程求导}{077 / 234}{**}
设函数 \(y=y(x)\) 由参数方程
\[\begin{cases}
x=2\mathrm{e}^t+t+1,\\
y=4(t-1)\mathrm{e}^t+t^2
\end{cases}
\]
确定，则
\[\left.\frac{\mathrm{d}^2 y}{\mathrm{d} x^2}\right|_{t=0}= \blank{}
\]
\QuestionEnd

\begin{QuestionSolution}{2021 年第 12 题}
\textbf{已知与目标}\quad 已知 \(x,y\) 关于 \(t\) 的表达式，要计算 \(t=0\) 时 \(y\) 关于 \(x\) 的二阶导数。

\textbf{方法与条件}\quad 参数式中撇号先表示对 \(t\) 求导。由 \(\frac{\mathrm dy}{\mathrm dx}=\frac{y'}{x'}\)，对商求 \(t\) 导数再除以 \(x'\)，得到 \(\frac{\mathrm d^2y}{\mathrm dx^2}=\frac{y''x'-y'x''}{(x')^3}\)。这里 \(x'=2\mathrm e^t+1>0\)，公式可用。

\textbf{推导}\quad 逐项求导为
\[
x'=2\mathrm{e}^t+1,\qquad y'=4t\mathrm{e}^t+2t,\qquad y''=4(1+t)\mathrm{e}^t+2,
\]
其中 \(x''=2\mathrm e^t\)。在 \(t=0\) 有 \(x'=3,y'=0,y''=6,x''=2\)，故
\[\left.\frac{\mathrm d^2y}{\mathrm dx^2}\right|_{t=0}=\frac{6\cdot3-0\cdot2}{3^3}=\frac{18}{27}=\frac23.\]

\textbf{结论与检查}\quad \(\boxed{\frac{2}{3}}\)。也可先化简 \(y'=2t(2\mathrm e^t+1)\)，得到一阶斜率 \(2t\)，再求二阶导数 \(\frac2{2\mathrm e^t+1}\)，零点同样为 \(\frac23\)。
\end{QuestionSolution}

% Q 2022 12
\QuestionPage{2.2 求导的基本法则}{填空题}{2022 年第 12 题}{2}{隐函数的二阶导数}{078 / 234}{**}
已知函数 \(y = y\left( x \right)\) 由方程
\(x^{2} + xy + y^{3} = 3\)
确定，则 \(y^{\prime\prime}\left( 1 \right) = \blank{}\)
\QuestionEnd

\begin{QuestionSolution}{2022 年第 12 题}
\textbf{已知与目标}\quad 已知 \(x^2+xy+y^3=3\)，要求隐函数的 \(y''(1)\)。

\textbf{方法与条件}\quad 先确定所问点，再进行两次隐函数求导。方程对 \(y\) 的偏导为 \(x+3y^2\)，在所问点不为零时可以求出 \(y'\)、\(y''\)。

\textbf{推导}\quad \(x=1\) 时 \(y^3+y=2\)，因左侧导数 \(3y^2+1>0\)，唯一解为 \(y=1\)，且 \(x+3y^2=4\ne0\)。一次求导得
\[
2x+y+(x+3y^2)y'=0,
\]
所以 \(y'(1)=-\frac{3}{4}\)。二次求导得
\[
2+2y'+6y(y')^2+(x+3y^2)y''=0.
\]
代入 \((x,y,y')=(1,1,-\frac{3}{4})\)，得
\[4y''=-2-2\left(-\frac34\right)-6\left(-\frac34\right)^2
=-2+\frac32-\frac{27}{8}=-\frac{31}{8}.\]

\textbf{结论与检查}\quad \(\boxed{y''(1)=-\frac{31}{32}}\)。代回二阶关系，\(2-\frac32+\frac{27}{8}-\frac{31}{8}=0\)，符合求导后的方程。
\end{QuestionSolution}

% Q 2023 14
\QuestionPage{2.2 求导的基本法则}{填空题}{2023 年第 14 题}{2}{隐函数曲线的法线}{079 / 234}{**}
曲线
\[3x^3=y^5+2y^3
\]
在 \(x=1\) 对应点处的法线斜率为
\[\blank{}.
\]
\QuestionEnd

\begin{QuestionSolution}{2023 年第 14 题}
\textbf{已知与目标}\quad 曲线由 \(3x^3=y^5+2y^3\) 确定，要求 \(x=1\) 对应点处的法线斜率。

\textbf{方法与条件}\quad 法线与切线垂直。切线斜率 \(y'\) 有限且非零时，法线斜率为 \(-\frac1{y'}\)。先对隐方程求导，不必求出 \(y(x)\) 的显式表达式。

\textbf{推导}\quad \(y^5+2y^3\) 是严格递增函数（两个奇次幂均严格递增），在 \(y=1\) 时为 \(3\)，故所问点为 \((1,1)\)。对原式求导得
\[
9x^2=(5y^4+6y^2)y',
\]
在 \((1,1)\) 代入得 \(9=(5+6)y'\)，故 \(y'=\frac9{11}\)。

\textbf{结论与检查}\quad 法线斜率为 \(\boxed{-\frac{11}{9}}\)。所问点的系数 \(5y^4+6y^2=11\ne0\)，切线斜率有限且不为零，负倒数公式适用。
\end{QuestionSolution}

% Q 2024 14
\QuestionPage{2.2 求导的基本法则}{填空题}{2024 年第 14 题}{2}{高阶导数}{080 / 234}{**}
已知函数
\[f(x)=x^2(\mathrm{e}^x-1),
\]
则
\[f^{(5)}(1)=\ \blank{}
\]
\QuestionEnd

\begin{QuestionSolution}{2024 年第 14 题}
\textbf{已知与目标}\quad 要计算 \(f=x^2(\mathrm e^x-1)\) 在 \(1\) 处的五阶导数。

\textbf{方法与条件}\quad Leibniz 公式将乘积的五次求导分配给两个因子。\(x^2\) 求导三次后为零，因此只需考虑它求导零次、一次和两次的三项；指数函数任意阶导数均为自身。

\textbf{推导}\quad 写 \(f=x^2\mathrm e^x-x^2\)。单独的 \(-x^2\) 的五阶导数为零，所以
\[
f^{(5)}(x)=\mathrm{e}^x\left[x^2+5(2x)+\binom52\cdot2\right].
\]
其中 \(5=\binom51\)，\(\binom52=10\)，分别表示五次求导中挑出一次或两次落在多项式上的组合数。代入 \(x=1\)，得 \(\mathrm e(1+10+20)=31\mathrm e\)。

\textbf{结论与检查}\quad \(\boxed{31\mathrm{e}}\)。多项式的三阶以上导数全部为零，所列三项已经包含全部非零项。
\end{QuestionSolution}

% Q 2025 14
\QuestionPage{2.2 求导的基本法则}{填空题}{2025 年第 14 题}{2}{隐函数与参数方程求导}{081 / 234}{**}
已知函数 \(y = y\left( x \right)\) 由 \(\left\{ \begin{aligned} x &= \ln \left( 1 + 2t \right), \\ 2t - \int_{1}^{y + t^{2}} \mathrm{e}^{-u^{2}} \,\mathrm{d} u &= 0 \end{aligned} \right.\)

确定，则 \(\left. \frac{\mathrm{d} y}{\mathrm{d} x} \right|_{x = 0} =\) \blank{}.
\QuestionEnd

\begin{QuestionSolution}{2025 年第 14 题}
\textbf{已知与目标}\quad \(x\) 由参数 \(t\) 给出，\(y\) 还满足积分关系，要求 \(x=0\) 时 \(\frac{\mathrm dy}{\mathrm dx}\)。

\textbf{方法与条件}\quad 先求对应的 \(t,y\)，对两式分别关于 \(t\) 求导，最后求导数比。积分上限 \(y+t^2\) 随 \(t\) 变化，求导时须乘 \(\frac{\mathrm dy}{\mathrm dt}+2t\)。

\textbf{推导}\quad \(x=\ln(1+2t)=0\) 给出 \(t=0\)。此时 \(\int_1^y\mathrm e^{-u^2}\,\mathrm du=0\)，因被积函数恒正，只有 \(y=1\) 时积分为零。对 \(2t=\int_1^{y+t^2}\mathrm e^{-u^2}\,\mathrm du\) 关于 \(t\) 求导，得到
\[
2=\mathrm{e}^{-(y+t^2)^2}\left(\frac{\mathrm{d}y}{\mathrm{d}t}+2t\right).
\]
在 \((t,y)=(0,1)\) 代入得 \(2=\mathrm e^{-1}\frac{\mathrm dy}{\mathrm dt}\)，所以 \(\frac{\mathrm dy}{\mathrm dt}=2\mathrm e\)。又 \(\frac{\mathrm dx}{\mathrm dt}=\frac2{1+2t}\)，在 \(t=0\) 为 \(2\)，故导数比为 \(\frac{2\mathrm e}{2}=\mathrm e\)。

\textbf{结论与检查}\quad \(\boxed{\left.\frac{\mathrm{d}y}{\mathrm{d}x}\right|_{x=0}=\mathrm{e}}\)。该点对数真数为 \(1>0\)，积分关系对 \(y\) 的导数为 \(\mathrm e^{-1}\ne0\)，且 \(x_t'=2\ne0\)，两次求导所需条件均满足。
\end{QuestionSolution}

% Q 2023 5
\QuestionPage{2.2 求导的基本法则}{选择题}{2023 年第 5 题}{1}{参数方程与可导性}{059 / 234}{***}
设函数 \(y=f(x)\) 由
\[\begin{cases}
x=2t+|t|,\\
y=|t|\sin t
\end{cases}
\]
确定，则（ ）。
\choices{\(f(x)\) 连续，\(f'(0)\) 不存在}{\(f'(0)\) 存在，\(f'(x)\) 在 \(x=0\) 处不连续}{\(f'(x)\) 连续，\(f''(0)\) 不存在}{\(f''(0)\) 存在，\(f''(x)\) 在 \(x=0\) 处不连续}
\QuestionEnd

\begin{QuestionSolution}{2023 年第 5 题}
\textbf{已知与目标}\quad 参数式含绝对值，要求判断消去参数后函数在零点的连续性、一阶导数及二阶导数是否存在。

\textbf{方法与条件}\quad 先按 \(t\) 的正负消去参数，再用导数定义比较左右差商。二阶导数是对一阶导数再次求导，不能只用 \(f\) 连续或 \(f'\) 连续代替判断。

\textbf{推导}\quad 当 \(t\ge0\) 时，\(x=3t\ge0\)；当 \(t<0\) 时，\(x=t<0\)。两段都能唯一反解，故
\[
f(x)=\begin{cases}
\frac{x}{3}\sin\frac{x}{3},&x\ge0,\\
-x\sin x,&x<0.
\end{cases}
\]
两侧函数值均趋于 \(0=f(0)\)，所以连续。零点处差商在右侧为 \(\frac{1}{3}\sin\frac{x}{3}\)，在左侧为 \(-\sin x\)，均趋于零，因此 \(f'(0)=0\)。在非零点由乘积及复合函数求导，得
\[
f'(x)=\begin{cases}
\frac{1}{3}\sin\frac{x}{3}+\frac{x}{9}\cos\frac{x}{3},&x>0,\\
-\sin x-x\cos x,&x<0.
\end{cases}
\]
两侧都趋于零，故 \(f'\) 在零点连续。再计算 \(\frac{f'(x)-f'(0)}{x}\)：右极限为 \(\frac{1}{9}\lim_{x\to0^+}\frac{\sin\left(\frac{x}{3}\right)}{\frac{x}{3}}+\frac{1}{9}\lim_{x\to0^+}\cos\frac{x}{3}=\frac{2}{9}\)；左极限为 \(-\lim_{x\to0^-}\frac{\sin x}{x}-\lim_{x\to0^-}\cos x=-2\)。两者不相等。

\textbf{结论与检查}\quad \(f'\) 连续而 \(f''(0)\) 不存在，选 \(\boxed{\mathrm C}\)。两侧参数与横坐标同号，分段区间没有颠倒；二阶导数的判断已经使用一阶导数的差商。
\end{QuestionSolution}

% Q 2019 13
\QuestionPage{2.2 求导的基本法则}{填空题}{2019 年第 13 题}{2}{高阶导数}{075 / 234}{***}
已知函数$f(x) = x\int_1^x \frac{\sin t^2}{t} \,\mathrm{d} t$，则$\int_0^1 f(x)\,\mathrm{d} x =$ \blank{}．
\QuestionEnd

\begin{QuestionSolution}{2019 年第 13 题}
\textbf{已知与目标}\quad 函数由外部因子 \(x\) 与变上限积分相乘构成，要求它在 \([0,1]\) 上的积分。

\textbf{方法与条件}\quad 令内部积分为 \(H(x)\)，分部积分时对它求导、对外面的 \(x\) 积分，可以把复杂的积分函数化成普通三角函数。零点处先检查连续延拓，以保证积分及边界项有意义。

\textbf{推导}\quad 记 \(H(x)=\int_1^x\frac{\sin(t^2)}{t}\,\mathrm{d}t\)。由于 \(\frac{\sin(t^2)}{t}=t\frac{\sin(t^2)}{t^2}\to0\)，被积函数在零点可补值为零，使其连续。因此 \(H\) 在 \([0,1]\) 有界，且对 \(x>0\) 有 \(H'(x)=\frac{\sin(x^2)}{x}\)。分部积分取 \(u=H(x)\)、\(\mathrm{d}v=x\,\mathrm{d}x\)，得到
\[
\int_0^1xH(x)\,\mathrm{d}x
=\left[\frac{x^2}{2}H(x)\right]_0^1
-\frac{1}{2}\int_0^1x\sin(x^2)\,\mathrm{d}x.
\]
上端因 \(H(1)=0\) 为零，下端因 \(H\) 有界、\(x^2\to0\) 也为零。余项令 \(u=x^2\)，\(\mathrm{d}u=2x\,\mathrm{d}x\)，上下限仍为 \(0,1\)，故
\[
\int_0^1 f(x)\,\mathrm{d}x
=-\frac{1}{4}\int_0^1\sin u\,\mathrm{d}u
=\frac{1}{4}[\cos u]_0^1
=\boxed{\frac{\cos1-1}{4}}.
\]

\textbf{结论与检查}\quad 在 \(0<x<1\) 上，内部积分从 \(1\) 积到较小的 \(x\)，被积函数为正，所以 \(f(x)<0\)。所得 \(\cos1-1<0\)，符号一致；零点的边界项也已检查。
\end{QuestionSolution}

% Q 2010 17
\QuestionPage{2.2 求导的基本法则}{解答题}{2010 年第 17 题}{3}{参数方程的二阶导数}{082 / 234}{***}
设函数$y = f(x)$由参数方程$\begin{cases}
  x = 2t + t^2 \\
  y = \psi(t)
\end{cases}(t > - 1)$所确定，其中$\psi (t)$具有$2$阶导数，且$\psi (1) = \frac{5}{2}$，$\psi'(1) = 6$．
已知$\frac{\mathrm{d}^2y}{\mathrm{d} x^2} = \frac{3}{4(1 + t)}$，求函数$\psi (t)$．
\QuestionEnd

\begin{QuestionSolution}{2010 年第 17 题}
\textbf{已知与目标}\quad 已知 \(x=2t+t^2\)、\(y=\psi(t)\) 的二阶导数和 \(t=1\) 时的两个初值，要求 \(\psi(t)\)。

\textbf{方法与条件}\quad 将一阶斜率记作新函数 \(p(t)\)。二阶导数给出 \(\frac{\mathrm dp}{\mathrm dx}\)，乘上 \(\frac{\mathrm dx}{\mathrm dt}\) 后就是 \(p'(t)\)，可通过积分恢复 \(p\)，再恢复 \(\psi\)。因为 \(t>-1\)，有 \(\frac{\mathrm dx}{\mathrm dt}=2(1+t)>0\)，各次转换均合法。

\textbf{推导}\quad 令 \(p(t)=\frac{\mathrm dy}{\mathrm dx}=\frac{\psi'(t)}{2(1+t)}\)，由链式法则，
\[
p'(t)=\frac{3}{4(1+t)}\cdot2(1+t)=\frac{3}{2}.
\]
积分得 \(p(t)=\frac32t+C_1\)。由 \(p(1)=\frac{\psi'(1)}4=\frac64=\frac32\)，得 \(C_1=0\)。所以 \(\psi'(t)=2(1+t)p(t)=3t+3t^2\)，再次积分得 \(\psi(t)=\frac32t^2+t^3+C_2\)。由 \(\psi(1)=\frac52\)，得 \(C_2=0\)。

\textbf{结论与检查}\quad \(\boxed{\psi(t)=t^3+\frac{3}{2}t^2\quad(t>-1)}\)。其初值为 \(\psi(1)=\frac52\)、\(\psi'(1)=6\)，一阶斜率为 \(\frac32t\)，再对 \(x\) 求导得 \(\frac{\frac32}{2(1+t)}=\frac3{4(1+t)}\)，全部符合题设。
\end{QuestionSolution}

% Q 2018 20
\QuestionPage{2.2 求导的基本法则}{解答题}{2018 年第 20 题}{3}{面积的相关变化率}{083 / 234}{***}
已知曲线$L:y = \frac{4}{9}x^2\left(x \ge 0 \right)$，点$O\left(0,0 \right)$，点$A\left(0,1 \right)$．
设$P$是$L$上的动点，$S$是直线$OA$与直线$AP$及曲线$L$所围图形的面积．
若$P$运动到$\left(3,4 \right)$时沿$x$轴正向的速度是$4$，求此时$S$关于时间$t$的变化率．
\QuestionEnd

\begin{QuestionSolution}{2018 年第 20 题}
\textbf{已知与目标}\quad \(P\) 在抛物线上运动，所问时刻 \(P=(3,4)\)、横向速度为 \(4\)，要求围成面积 \(S\) 对时间的导数。

\textbf{方法与条件}\quad 用横坐标 \(a\) 表示动点，再将面积写为 \(S(a)\)。链式法则给出 \(\frac{\mathrm dS}{\mathrm dt}=\frac{\mathrm dS}{\mathrm da}\frac{\mathrm da}{\mathrm dt}\)。用上下曲线之差积分求面积前，需要确定哪条曲线在上方。

\textbf{推导}\quad 设 \(P=(a,\frac{4a^2}{9})\)，在所问时刻附近 \(a>0\)。直线 \(AP\) 经过 \((0,1)\)，所以其方程为 \(y=1+\frac{(\frac{4a^2}{9}-1)x}{a}\)。它减去抛物线的纵坐标，得到
\[1+\frac{(\frac{4a^2}{9}-1)x}{a}-\frac49x^2=(a-x)\left(\frac{4x}{9}+\frac1a\right)\ge0\quad(0\le x\le a).\]
因此该区间上直线在上、抛物线在下；直线下方是梯形，可得
\[
S=\frac{a}{2}\left(1+\frac{4a^2}{9}\right)-\int_0^a\frac{4}{9}x^2\,\mathrm{d}x
=\frac{a}{2}+\frac{2a^3}{9}-\frac{4a^3}{27}=\frac{a}{2}+\frac{2a^3}{27}.
\]
因此 \(\frac{\mathrm dS}{\mathrm dt}=(\frac12+\frac{2a^2}{9})\frac{\mathrm da}{\mathrm dt}\)，代入 \(a=3,\frac{\mathrm da}{\mathrm dt}=4\)，得 \((\frac12+2)\cdot4=10\)。

\textbf{结论与检查}\quad \(\boxed{\frac{\mathrm{d}S}{\mathrm{d}t}=10}\)。面积导数 \(\frac12+\frac{2a^2}{9}>0\)，横坐标正向变化时面积增加，所得正号符合几何关系。
\end{QuestionSolution}

\SectionHeading{2.3}{一元函数微分学及其应用}{微分}{定理与要点\quad 教材 PDF 第 140 页起\hfill 本节 0 题}
\pdfbookmark[2]{2.3 微分}{section-2.3}
\Theory{定理 3.1}{141}{可微与可导的等价性}{一元函数 \(f\) 在一点可微的充要条件是它在该点可导，此时 \(\mathrm{d}f=f'(x)\,\mathrm{d}x\)。}
\Theory{公式}{143--144}{微分及线性近似}{一阶微分形式适用于复合变量。增量充分小时，可用 \(f'(x)\,\mathrm{d}x\) 近似函数增量。}
本节在所收原卷中没有以该节为主考点的题目。
\SectionHeading{2.4}{一元函数微分学及其应用}{微分中值定理及其应用}{定理与要点\quad 教材 PDF 第 147 页起\hfill 本节 12 题}
\pdfbookmark[2]{2.4 微分中值定理及其应用}{section-2.4}
\Theory{定理 4.1}{147}{Fermat 定理}{内点取得局部极值且在该点可导，则导数为零。}
\Theory{定理 4.2}{148}{Rolle 定理}{闭区间上连续、开区间内可导且端点值相等，内部存在导数为零的点。}
\Theory{定理 4.3}{150}{Lagrange 中值定理}{闭区间上连续、开区间内可导，内部存在一点，其导数等于全区间的平均变化率。}
\Theory{定理 4.4}{154}{Cauchy 中值定理}{两函数在闭区间上连续、开区间内可导，且 \(g'\) 在开区间内恒不为零，则区间内某点的导数比等于端点增量比。}
\Theory{定理 4.5}{154--155}{L'Hospital 法则（\(\frac{0}{0}\) 型）}{两函数趋零，邻域内可导且分母导数非零。若导数之比有极限，则原商有同一极限。}
\Theory{定理 4.6}{156}{L'Hospital 法则（\(\frac{\infty}{\infty}\) 型）}{两函数在同一变化趋势下为无穷大量，其余条件沿用定理 4.5，可由导数比求原商极限。}
\Theory{推论 4.3}{152}{端点导数}{函数在端点连续、在相邻开区间可导且导数的单侧极限有限时，该极限就是端点的单侧导数。}
% Q 2014 5
\QuestionPage{2.4 微分中值定理及其应用}{选择题}{2014 年第 5 题}{1}{L'Hospital 法则与函数极限}{084 / 234}{**}
设函数$f(x) = \arctan x$，若$f(x) = xf'(\xi)$，则$\lim_{x \to 0} \frac{\xi^2}{x^2} =$\pickout{}
\begin{abcd}
\item $1$．
\item $\frac{2}{3}$．
\item $\frac{1}{2}$．
\item $\frac{1}{3}$．
\end{abcd}
\QuestionEnd

\begin{QuestionSolution}{2014 年第 5 题}
\textbf{已知与目标}\quad \(f(x)=\arctan x\)，且 \(f(x)=xf'(\xi)\)，要求 \(\frac{\xi^2}{x^2}\) 的极限。

\textbf{方法与条件}\quad 先代入 \(f'(\xi)=\frac1{1+\xi^2}\) 消去导数，避免直接求 \(\xi\)。化成只含 \(x\) 的极限后，对分子分母都趋零的部分用 L'Hospital 法则；需检查分母导数在去掉零点的邻域内非零。

\textbf{推导}\quad 由 \(\arctan x=\frac{x}{1+\xi^2}\) 得 \(\xi^2=\frac{x-\arctan x}{\arctan x}\)，所以
\[
\frac{\xi^2}{x^2}
=\frac{x-\arctan x}{x^3}\frac{x}{\arctan x}.
\]
由于 \(\frac{\arctan x}{x}\to1\)，第二个因子趋于 \(1\)。第一个因子的分子分母都趋零，分母导数 \(3x^2\) 在 \(x\ne0\) 时非零，且
\[\frac{(x-\arctan x)'}{(x^3)'}=\frac{1-\frac1{1+x^2}}{3x^2}=\frac1{3(1+x^2)}\to\frac13.\]

\textbf{结论与检查}\quad 原极限为 \(\boxed{\frac{1}{3}}\)，选 D。所得正数也符合 \(\frac{\xi^2}{x^2}\ge0\)；计算只涉及 \(\xi^2\)，无需选择 \(\xi\) 的正负。
\end{QuestionSolution}

% Q 2018 9
\QuestionPage{2.4 微分中值定理及其应用}{填空题}{2018 年第 9 题}{2}{反正切差的极限与中值定理}{086 / 234}{**}
$\lim_{x \to +\infty} x^2[\arctan(x + 1) - \arctan x] =$ \blank{}．
\QuestionEnd

\begin{QuestionSolution}{2018 年第 9 题}
\textbf{已知与目标}\quad 两个反正切函数的自变量相差 \(1\)，要求其差与 \(x^2\) 的乘积在无穷远处的极限。

\textbf{方法与条件}\quad Lagrange 中值定理把函数增量写成区间内某点的导数乘区间长度。\(\arctan\) 在 \([x,x+1]\) 上连续、内部可导，满足定理条件。

\textbf{推导}\quad 存在 \(\xi_x\in(x,x+1)\) 使
\[\arctan(x+1)-\arctan x=(1)\frac1{1+\xi_x^2}.\]
对充分大的正 \(x\)，有 \(1<\frac{\xi_x}{x}<1+\frac1x\)，由夹逼可知 \(\frac{\xi_x}{x}\to1\)。将分子分母都除以 \(x^2\)，得
\[
\frac{x^2}{1+\xi_x^2}=\frac{1}{\frac{1}{x^2}+(\frac{\xi_x}{x})^2}\longrightarrow\boxed{1}.
\]
\textbf{结论与检查}\quad 极限为 \(\boxed{1}\)。反正切严格递增，所以差为正，结果的符号符合原式；\(\xi_x\) 随 \(x\) 变化，但上述夹逼对所有可能的中值点均成立。
\end{QuestionSolution}

% Q 2017 15
\QuestionPage{2.4 微分中值定理及其应用}{解答题}{2017 年第 15 题}{2}{变上限积分的未定式极限}{091 / 234}{**}
求极限$\lim_{x \to 0^+} \frac{\int_0^x \sqrt{x - t} \mathrm{e}^t \,\mathrm{d} t}{\sqrt{x^3}}$．
\QuestionEnd

\begin{QuestionSolution}{2017 年第 15 题}
\textbf{已知与目标}\quad \(x\to0^+\)，要计算带参数积分除以 \(\sqrt{x^3}=x^{\frac32}\) 的极限。

\textbf{方法与条件}\quad 夹逼准则是把待求量夹在趋于同一值的上下界之间。积分区间上 \(1\le\mathrm e^t\le\mathrm e^x\)，根式非负，故乘上根式并积分都保持不等号方向。

\textbf{推导}\quad 对 \(0\le t\le x\)，有 \(\sqrt{x-t}\le\sqrt{x-t}\mathrm e^t\le\mathrm e^x\sqrt{x-t}\)。令 \(u=x-t\)，直接积分得
\[\int_0^x\sqrt{x-t}\,\mathrm dt=\int_0^x u^{\frac12}\,\mathrm du=\left[\frac23u^{\frac32}\right]_0^x=\frac23x^{\frac32}.\]
将不等式积分，再除以正数 \(x^{\frac32}\)，得到
\[
\frac23\le\frac{\int_0^x\sqrt{x-t}\mathrm{e}^t\,\mathrm{d}t}{x^{\frac{3}{2}}}\le\frac23 \mathrm{e}^x.
\]
\textbf{结论与检查}\quad 两端分别趋于 \(\frac23\) 和 \(\frac23\mathrm e^0=\frac23\)，故原极限为 \(\boxed{\frac{2}{3}}\)。所用根式、正分母和积分区间均要求 \(x>0\)，与右极限一致。
\end{QuestionSolution}

% Q 2017 2
\QuestionPage{2.4 微分中值定理及其应用}{选择题}{2017 年第 2 题}{1}{由可导条件推断中值点}{085 / 234}{***}
设二阶可导函数$f(x)$满足$f(1) = f(- 1) = 1$，$f(0) = - 1$，且$f''(x) > 0$，则\pickout{}
\begin{abcd}
\item $\int_{-1}^1 f(x)\,\mathrm{d} x > 0$．
\item $\int_{-1}^1 f(x)\,\mathrm{d} x < 0$．
\item $\int_{-1}^0 f(x)\,\mathrm{d} x > \int_0^1 f(x)\,\mathrm{d} x$．
\item $\int_{-1}^0 f(x)\,\mathrm{d} x < \int_0^1 f(x)\,\mathrm{d} x$．
\end{abcd}
\QuestionEnd

\begin{QuestionSolution}{2017 年第 2 题}
\textbf{已知与目标}\quad 已知三个函数值和 \(f''>0\)，要判断整个区间积分的符号及两半区间的关系。

\textbf{方法与条件}\quad \(f''>0\) 表示斜率严格增加，曲线在任意两点连成的弦线下方。分别在两个半区间作弦线，便能用容易积分的直线估计原积分。

\textbf{推导}\quad 左半区间端点为 \((-1,1),(0,-1)\)，弦线为 \(y=-2x-1\)，所以 \(f(x)<-2x-1\)（\(-1<x<0\)）。右半区间弦线为 \(y=2x-1\)，所以 \(f(x)<2x-1\)（\(0<x<1\)）。由连续性和严格不等式，
\[
\int_{-1}^0 f(x)\,\mathrm dx<[-x^2-x]_{-1}^0=0,
\qquad \int_0^1 f(x)\,\mathrm dx<[x^2-x]_0^1=0.
\]
相加得到总积分小于零。若要判断 C、D，取符合题设的 \(f(x)=2x^2-1\)，因它是偶函数，两半积分相等，不能得到严格大小关系。

\textbf{结论与检查}\quad 选 \(\boxed{\mathrm B}\)。弦线估计分别利用了 \(f(0)=-1\)，只连接最外侧两端点不足以判断总积分为负。
\end{QuestionSolution}

% Q 2009 21
\QuestionPage{2.4 微分中值定理及其应用}{解答题}{2009 年第 21 题}{4}{Lagrange 中值定理}{087 / 234}{***}
\begin{enumerate}
\item 证明 Lagrange 中值定理：若函数$f(x)$在$\left[a,b \right]$上连续，在$(a,b)$可导，
      则存在$\xi \in \left(a,b \right)$，使得$f(b) - f(a) = f'\left(\xi \right)(b - a)$．
\item 证明：若函数$f(x)$在$x = 0$处连续，在$\left(0,\delta \right)\left(\delta > 0 \right)$内可导，
      且$\lim_{x \to 0^+} f'(x) = A$，则$f'_+(0)$存在，且$f'_+(0) = A$．
\end{enumerate}
\QuestionEnd

\begin{QuestionSolution}{2009 年第 21 题}
\textbf{已知与目标}\quad 第一问从连续、可导条件证明 Lagrange 中值定理；第二问用导数的右极限证明端点的右导数存在。

\textbf{方法与条件}\quad Rolle 定理要求闭区间连续、开区间可导、端点值相等。为满足第三项条件，从 \(f\) 减去两端点的连线；第二问再把第一问用于 \([0,x]\)。

\textbf{推导}\quad （1）因 \(a<b\)，可定义
\[
F(x)=f(x)-f(a)-\frac{f(b)-f(a)}{b-a}(x-a).
\]
减去的部分是直线，故 \(F\) 仍在 \([a,b]\) 连续、\((a,b)\) 可导，且代入端点得到 \(F(a)=F(b)=0\)。由 Rolle 定理，存在 \(\xi\in(a,b)\) 使 \(F'(\xi)=0\)。由于 \(F'(x)=f'(x)-\frac{f(b)-f(a)}{b-a}\)，整理得到第一问的等式。

（2）固定任意 \(x\in(0,\delta)\)。在 \(0\) 处连续由题设给出，在其余点可导所以连续，故 \(f\) 在 \([0,x]\) 连续、\((0,x)\) 可导。应用（1），存在 \(\xi_x\in(0,x)\) 使
\[
\frac{f(x)-f(0)}x=f'(\xi_x).
\]
由于 \(0<\xi_x<x\)，当 \(x\to0^+\) 时 \(\xi_x\to0^+\)，故 \(f'(\xi_x)\to A\)。

\textbf{结论与检查}\quad 第一问得到
\[\boxed{f(b)-f(a)=f'(\xi)(b-a)}.\]
第二问由右差商极限定义得到 \(\boxed{f'_+(0)=A}\)。第二问只使用右侧区间，没有假定左导数存在。
\end{QuestionSolution}

% Q 2010 21
\QuestionPage{2.4 微分中值定理及其应用}{解答题}{2010 年第 21 题}{4}{两段区间上的中值定理}{088 / 234}{***}
设函数$f(x)$在闭区间$[0,1]$上连续，在开区间$(0,1)$内可导，
且$f(0) = 0$, $f(1) = \frac{1}{3}$．证明：存在$\xi \in \left(0,\frac{1}{2}\right)$,
$\eta \in \left(\frac{1}{2},1\right)$，使得$f'(\xi) + f'(\eta) = \xi^2 + \eta^2$．
\QuestionEnd

\begin{QuestionSolution}{2010 年第 21 题}
\textbf{已知与目标}\quad 已知 \(f(0)=0,f(1)=\frac13\)，要在左右半区间各找到一点，使两个导数之和等于两点的平方和。

\textbf{方法与条件}\quad 因 \((\frac{x^3}{3})'=x^2\)，令 \(F=f-\frac{x^3}{3}\)，目标即为 \(F'(\xi)+F'(\eta)=0\)。两个半区间长度相等，可分别用 Lagrange 中值定理，使中点的函数值在相加时消失。

\textbf{推导}\quad \(F\) 保留题设的连续性、可导性，且 \(F(0)=0,F(1)=\frac13-\frac13=0\)。在 \([0,\frac12]\)、\([\frac12,1]\) 上分别满足中值定理条件，故存在 \(\xi\in(0,\frac12)\)、\(\eta\in(\frac12,1)\)，满足
\[
F'(\xi)=2[F(\frac{1}{2})-F(0)],\qquad
F'(\eta)=2[F(1)-F(\frac{1}{2})].
\]
两式相加得 \(F'(\xi)+F'(\eta)=2[F(1)-F(0)]=0\)。再用 \(F'(u)=f'(u)-u^2\)，得到
\[f'(\xi)-\xi^2+f'(\eta)-\eta^2=0.\]

\textbf{结论与检查}\quad \(\boxed{f'(\xi)+f'(\eta)=\xi^2+\eta^2}\)。两点分别由两个半区间内的中值定理取得，严格属于题目要求的开区间，不会取到中点或端点。
\end{QuestionSolution}

% Q 2015 21
\QuestionPage{2.4 微分中值定理及其应用}{解答题}{2015 年第 21 题}{4}{切线与中值定理证明}{090 / 234}{***}
已知函数$f(x)$在区间$[a,+\infty)$上具有$2$阶导数，
$f(a) = 0$，$f'(x) > 0$，$f''(x) > 0$．设$b > a$，
曲线$y = f(x)$在点$(b,f(b))$处的切线与$x$轴的交点是$(x_0,0)$，
证明$a < x_0 < b$．
\QuestionEnd

\begin{QuestionSolution}{2015 年第 21 题}
\textbf{已知与目标}\quad 已知 \(f(a)=0\)，且一、二阶导数均为正，要证明 \(b\) 处切线与横轴的交点位于 \(a,b\) 之间。

\textbf{方法与条件}\quad 先写出交点横坐标。\(f'>0\) 保证函数增加，\(f''>0\) 保证 \(f'\) 严格增加。将 \([a,b]\) 的平均斜率与末端斜率比较，就能给交点定位。

\textbf{推导}\quad 切线方程为 \(y=f(b)+f'(b)(x-b)\)。令 \(y=0\)，因 \(f'(b)>0\)，得到
\[
x_0=b-\frac{f(b)}{f'(b)}.
\]
\(f\) 在 \([a,b]\) 连续、内部可导，由 Lagrange 中值定理，存在 \(\xi\in(a,b)\) 使 \(f(b)-f(a)=(b-a)f'(\xi)\)。由于 \(f(a)=0\)，并且 \(0<f'(\xi)<f'(b)\)，有
\[0<\frac{f(b)}{f'(b)}=(b-a)\frac{f'(\xi)}{f'(b)}<b-a.\]
将整个不等式从 \(b\) 中减去，得到 \(b-(b-a)<b-\frac{f(b)}{f'(b)}<b\)。

\textbf{结论与检查}\quad \(\boxed{a<x_0<b}\)。\(f(b)>0\) 使交点在 \(b\) 左侧，严格增加的斜率使它仍在 \(a\) 右侧；两侧严格性均由题设的严格正号保证。
\end{QuestionSolution}

% Q 2013 18
\QuestionPage{2.4 微分中值定理及其应用}{解答题}{2013 年第 18 题}{3}{奇函数与中值定理}{089 / 234}{****}
设奇函数$f(x)$在$[- 1,1]$上具有$2$阶导数，且$f(1) = 1$．证明：
\begin{enumerate}
\item 存在$\xi \in(0,1)$，使得$f'(\xi) = 1$；
\item 存在$\eta \in (-1,1)$，使得$f''(\eta) + f'(\eta) = 1$．
\end{enumerate}
\QuestionEnd

\begin{QuestionSolution}{2013 年第 18 题}
\textbf{已知与目标}\quad \(f\) 为二阶可导奇函数，\(f(1)=1\)。第一问找 \(f'=1\) 的点，第二问找 \(f''+f'=1\) 的点。

\textbf{方法与条件}\quad 第一问用 Lagrange 中值定理。第二问需要构造一个导数含 \(f''+f'-1\) 的函数；乘 \(\mathrm e^x\) 后，乘积求导恰好产生两项之和。再利用奇函数的导数为偶函数，找到 Rolle 定理所需的两个零点。

\textbf{推导}\quad （1）奇性给出 \(f(0)=-f(0)\)，故 \(f(0)=0\)。二阶可导保证 \(f\) 在 \([0,1]\) 连续、内部可导，故存在 \(\xi\in(0,1)\) 使 \(f'(\xi)=\frac{f(1)-f(0)}{1-0}=1\)。

（2）对 \(f(-x)=-f(x)\) 求导，得到 \(-f'(-x)=-f'(x)\)，所以 \(f'(-x)=f'(x)\)，特别有 \(f'(-\xi)=1\)。令 \(H(x)=\mathrm e^x(f'(x)-1)\)。因为 \(f'\) 可导，\(H\) 在 \([-\xi,\xi]\) 连续、内部可导，且 \(H(-\xi)=H(\xi)=0\)。由 Rolle 定理，存在 \(\eta\in(-\xi,\xi)\) 使
\[
0=H'(\eta)=\mathrm{e}^\eta(f''(\eta)+f'(\eta)-1).
\]
\textbf{结论与检查}\quad 第一问得到 \(\boxed{f'(\xi)=1}\)。第二问由于 \(\mathrm e^\eta>0\)，可约去指数因子，得到
\[\boxed{f''(\eta)+f'(\eta)=1}.\]
又 \((-\xi,\xi)\subset(-1,1)\)，点的位置符合题意。
\end{QuestionSolution}

% Q 2017 19
\QuestionPage{2.4 微分中值定理及其应用}{解答题}{2017 年第 19 题}{3}{Rolle 定理与方程的根}{092 / 234}{****}
设函数$f(x)$在区间$[0,1]$上具有二阶导数，且$f(1) > 0$，$\lim_{x \to 0^+} \frac{f(x)}{x} < 0$．证明：
\begin{enumerate}
\item 方程$f(x) = 0$在区间$(0,1)$内至少存在一个实根．
\item 方程$f(x)f''(x) + \left(f'(x) \right)^2 = 0$在区间$(0,1)$内至少存在两个不同实根．
\end{enumerate}
\QuestionEnd

\begin{QuestionSolution}{2017 年第 19 题}
\textbf{已知与目标}\quad 已知 \(f(1)>0\)，而 \(\frac{f(x)}x\) 在零点右侧有负极限。先证明 \(f\) 有一个内部零点，再证明 \(ff''+(f')^2\) 有两个不同内部零点。

\textbf{方法与条件}\quad 第一问用连续函数的零点定理：函数值异号的两个端点之间存在零点。第二问利用乘积求导 \((ff')'=ff''+(f')^2\)，先构造 \(ff'\) 的三个零点，再在两个互不相交的区间内分别用 Rolle 定理。

\textbf{推导}\quad （1）设给定极限为 \(L<0\)。充分小的 \(x>0\) 有 \(\frac{f(x)}x<0\)，所以可取 \(s\in(0,1)\) 使 \(f(s)<0\)。二阶可导保证连续，结合 \(f(1)>0\)，零点定理给出 \(c\in(s,1)\) 使 \(f(c)=0\)。

（2）有限极限还给出 \(f(x)=x\frac{f(x)}x\to0\)，由零点连续性得 \(f(0)=0\)。在 \([0,c]\) 上 \(f\) 连续、内部可导且两端为零，由 Rolle 定理得到 \(d\in(0,c)\) 使 \(f'(d)=0\)。令 \(H=ff'\)，则 \(H(0)=H(d)=H(c)=0\)。因为 \(f'\) 可导，所以 \(H\) 在这些闭区间连续、内部可导。分别在 \([0,d]\)、\([d,c]\) 应用 Rolle 定理，得到 \(r_1\in(0,d)\)、\(r_2\in(d,c)\) 使
\[
H'(x)=f(x)f''(x)+[f'(x)]^2=0.
\]
\textbf{结论与检查}\quad 第一问已有零点 \(c\in(0,1)\)；第二问的 \(r_1,r_2\) 均为所求方程的根。因 \(r_1<d<r_2<c<1\)，两根严格不同且均属于 \((0,1)\)。
\end{QuestionSolution}

% Q 2019 21
\QuestionPage{2.4 微分中值定理及其应用}{解答题}{2019 年第 21 题}{4}{中值定理与导数估计}{093 / 234}{****}
已知函数$f(x)$在$[0,1]$上具有二阶导数，且
\[
f(0) = 0,\quad f(1) = 1,\quad \int_0^1 f(x)\,\mathrm{d} x = 1.
\]
证明：
\begin{enumerate}
\item 存在$\xi \in(0,1)$，使得$f'(\xi) = 0$；
\item 存在$\eta \in(0,1)$，使得$f''(\eta) < - 2$．
\end{enumerate}
\QuestionEnd

\begin{QuestionSolution}{2019 年第 21 题}
\textbf{已知与目标}\quad 端点值为 \(0,1\)，总积分却为 \(1\)，要证明内部存在导数为零的点和二阶导数小于 \(-2\) 的点。

\textbf{方法与条件}\quad 第一问由积分条件找出内部最大值，再用 Fermat 定理（可导函数在内部极值点导数为零）。第二问用反证法，假设 \(f''\ge-2\)，加上 \(x^2\) 后二阶导数非负，即曲线位于弦线下方，由此给积分一个不可能的上界。

\textbf{推导}\quad （1）若整个区间上 \(f\le1\)，由 \(f(0)=0\) 和连续性，可取 \(\delta\in(0,1)\)，使 \([0,\delta]\) 上 \(f\le\frac12\)。则
\[\int_0^1 f\le\frac\delta2+(1-\delta)=1-\frac\delta2<1,\]
与题设矛盾。故某处 \(f>1\)。函数在闭区间连续，最大值一定取得，且两端函数值都不超过 \(1\)，所以最大值在某个内点 \(\xi\) 取得；可导性给出 \(f'(\xi)=0\)。

（2）若处处 \(f''(x)\ge-2\)，令 \(H(x)=f(x)+x^2\)。则 \(H''\ge0\)，且 \(H(0)=0,H(1)=2\)，由弦线估计有 \(H(x)\le2x\)。于是
\[
1=\int_0^1 f(x)\,\mathrm{d}x\le\int_0^1(2x-x^2)\,\mathrm{d}x=\frac23,
\]
此处右端为 \([x^2-\frac{x^3}{3}]_0^1=\frac23\)，与左侧等于 \(1\) 矛盾。

\textbf{结论与检查}\quad 因而 \(\boxed{f'(\xi)=0}\)，并且存在 \(\eta\in(0,1)\) 使 \(\boxed{f''(\eta)<-2}\)。第二问反设的是整个开区间上 \(f''\ge-2\)，否定后得到严格小于号，恰与目标相符。
\end{QuestionSolution}

% Q 2023 21
\QuestionPage{2.4 微分中值定理及其应用}{解答题}{2023 年第 21 题}{4}{极值条件与二阶中值结论}{094 / 234}{****}
设函数 \(f(x)\) 在 \([-a,a]\) 上具有二阶连续导数。证明：
\begin{enumerate}
\item 若 \(f(0)=0\)，则存在 \(\xi\in(-a,a)\)，使得
\[f''(\xi)=\frac{f(a)+f(-a)}{a^2};
\]
\item 若 \(f(x)\) 在 \((-a,a)\) 内取得极值，则存在 \(\eta\in(-a,a)\)，使得
\[|f''(\eta)|\ge \frac{|f(a)-f(-a)|}{2a^2}.
\]
\end{enumerate}
\QuestionEnd

\begin{QuestionSolution}{2023 年第 21 题}
\textbf{已知与目标}\quad 在 \(a>0\) 的对称区间上，\(f\) 二阶连续可导。第一问由 \(f(0)=0\) 找一个二阶导数中值点，第二问由内部极值推出二阶导数的大小估计。

\textbf{方法与条件}\quad 第一问构造经过 \(-a,0,a\) 三点的二次多项式，用三次 Rolle 定理把三个函数零点转成一个二阶导数零点。第二问在极值点使用带 Lagrange 余项的 Taylor 公式，即函数增量等于一次项与某个中间点的二阶导数项之和。二阶连续导数保证上述定理均可用。

\textbf{推导}\quad （1）令 \(P(x)=Ax^2+Bx\)，其中
\[
A=\frac{f(a)+f(-a)}{2a^2},\qquad
B=\frac{f(a)-f(-a)}{2a}.
\]
代入得 \(P(a)=f(a),P(-a)=f(-a),P(0)=0=f(0)\)。因此 \(F=f-P\) 在 \(-a,0,a\) 均为零。\(F\) 在两个半闭区间连续、内部可导，分别由 Rolle 定理得到 \(r\in(-a,0),s\in(0,a)\) 使 \(F'(r)=F'(s)=0\)。再因 \(F'\) 在 \([r,s]\) 连续、内部可导，得到 \(\xi\in(r,s)\) 使 \(F''(\xi)=0\)。于是 \(f''(\xi)=P''(\xi)=2A\)，即第一问结论。

（2）设内点 \(c\) 为极值点，则 \(f'(c)=0\)。将 \(f(a),f(-a)\) 分别在 \(c\) 处作带 Lagrange 余项的一阶 Taylor 展开，得
\[
\begin{aligned}
f(a)-f(c)&=\frac{1}{2}f''(\eta_1)(a-c)^2,\\
f(-a)-f(c)&=\frac{1}{2}f''(\eta_2)(a+c)^2,
\end{aligned}
\]
其中 \(\eta_1\in(c,a)\)、\(\eta_2\in(-a,c)\)。取 \(\eta\) 为其中 \(|f''|\) 较大的点。相减、用三角不等式，并用 \(\frac12[(a-c)^2+(a+c)^2]=a^2+c^2\)，得到
\[
|f(a)-f(-a)|\le |f''(\eta)|(a^2+c^2)
\le 2a^2|f''(\eta)|.
\]
\textbf{结论与检查}\quad 第一问为
\[\boxed{f''(\xi)=\frac{f(a)+f(-a)}{a^2}}.\]
第二问除以正数 \(2a^2\)，得
\[\boxed{|f''(\eta)|\ge\frac{|f(a)-f(-a)|}{2a^2}}.\]
\(|c|<a\) 保证 \(a^2+c^2\le2a^2\)，所取 \(\xi,\eta\) 都在开区间内。
\end{QuestionSolution}

% Q 2024 21
\QuestionPage{2.4 微分中值定理及其应用}{解答题}{2024 年第 21 题}{3}{中值定理与导数估计}{095 / 234}{****}
设函数 \(f(x)\) 具有二阶导数，且
\[f'(0)=f'(1),\qquad |f''(x)|\le 1.
\]
证明：
\begin{enumerate}
\item 当 \(x\in(0,1)\) 时，
\[\left|f(x)-f(0)(1-x)-f(1)x\right|\le \frac{x(1-x)}{2};
\]
\item \[\left|\int_0^1 f(x)\,\mathrm{d}x-\frac{f(0)+f(1)}{2}\right|\le \frac{1}{12}.
\]
\end{enumerate}
\QuestionEnd

\begin{QuestionSolution}{2024 年第 21 题}
\textbf{已知与目标}\quad 已知 \(|f''|\le1\)，要估计函数与两端点连线的误差，以及积分与梯形面积的误差。

\textbf{方法与条件}\quad 将端点连线从 \(f\) 中减去，所得误差两端为零。再加减二次函数 \(p=\frac{x(1-x)}2\)，使二阶导数分别非负、非正。二阶导数非负的函数在弦线下方，非正时在弦线上方，这给出误差上下界。该证明在题设条件下成立，仅用二阶导数界即可完成估计。

\textbf{推导}\quad （1）设 \(u(x)=f(x)-f(0)(1-x)-f(1)x\)，\(p(x)=\frac{x(1-x)}{2}\)。两者在端点均为零，且 \(u''=f''\)、\(p''=-1\)，故
\[
(u-p)''=f''+1\ge0,\qquad
(u+p)''=f''-1\le0.
\]
\(u-p\)、\(u+p\) 在闭区间连续且内部二阶可导，弦线均为零函数。由弦线不等式，\(u-p\le0\)，而 \(u+p\ge0\)，即 \(-p\le u\le p\)。因 \(p\ge0\)，得到第一问的 \(|u(x)|\le p(x)\)。

（2）端点弦线在 \([0,1]\) 上的积分为 \(\frac{f(0)+f(1)}{2}\)。因此
\[
\left|\int_0^1 f(x)\,\mathrm{d}x-\frac{f(0)+f(1)}{2}\right|
\le\int_0^1 |u(x)|\,\mathrm{d}x
\le\int_0^1\frac{x(1-x)}{2}\,\mathrm{d}x
=\frac12\left[\frac{x^2}{2}-\frac{x^3}{3}\right]_0^1=\frac{1}{12}.
\]
\textbf{结论与检查}\quad 分别得到 \(\boxed{|u(x)|\le\frac{x(1-x)}2}\) 和 \(\boxed{\left|\int_0^1f-\frac{f(0)+f(1)}2\right|\le\frac1{12}}\)。端点误差为零，符合界函数 \(p\) 的端点值；积分估计先取绝对值再用 \(\int|u|\)，无须假设误差同号。
\end{QuestionSolution}

\SectionHeading{2.5}{一元函数微分学及其应用}{Taylor定理及其应用}{定理与要点\quad 教材 PDF 第 161 页起\hfill 本节 9 题}
\pdfbookmark[2]{2.5 Taylor定理及其应用}{section-2.5}
\Theory{定理 5.1}{162}{带 Peano 余项的 Taylor 定理}{若 \(f\) 在 \(x_0\) 处有 \(n\) 阶导数，则 \(f(x)\) 等于在 \(x_0\) 处的 \(n\) 次 Taylor 多项式加上 \(o((x-x_0)^n)\)。}
\Theory{定理 5.2}{164}{带 Lagrange 余项的 Taylor 定理}{区间上 \(n+1\) 阶可导时，余项为 \(\frac{f^{(n+1)}(\xi)}{(n+1)!}(x-x_0)^{n+1}\)，其中 \(\xi\) 在 \(x_0\) 与 \(x\) 之间。}
\Theory{公式}{165--167}{Maclaurin 公式}{在 \(x=0\) 处展开时，\(n\) 次项的系数为 \(\frac{f^{(n)}(0)}{n!}\)。求极限时，展开到相消后出现首个非零项为止。}
% Q 2021 5
\QuestionPage{2.5 Taylor定理及其应用}{选择题}{2021 年第 5 题}{1}{二次 Taylor 多项式}{098 / 234}{*}
设函数 \(f(x)=\sec x\) 在 \(x=0\) 处的 2 次 Taylor 多项式为 \(1+ax+bx^2\)，则
\choices{\(a=1\)，\(b=-\dfrac{1}{2}\)}{\(a=1\)，\(b=\dfrac{1}{2}\)}{\(a=0\)，\(b=-\dfrac{1}{2}\)}{\(a=0\)，\(b=\dfrac{1}{2}\)}
\QuestionEnd

\begin{QuestionSolution}{2021 年第 5 题}
\textbf{已知与目标}\quad 正割函数 \(\sec x=\frac1{\cos x}\) 的二次 Taylor 多项式为 \(1+ax+bx^2\)，要求 \(a,b\)。

\textbf{方法与条件}\quad 二次 Taylor 多项式的公式为 \(f(0)+f'(0)x+\frac{f''(0)}2x^2\)。\(\cos0=1\ne0\)，所以正割在零点附近光滑，可以直接求导代入。

\textbf{推导}\quad \(f'=\sec x\tan x\)，\(f''=\sec x\tan^2x+\sec^3x\)。在零点 \(\sec0=1,\tan0=0\)，得 \(f(0)=1,f'(0)=0,f''(0)=1\)，故多项式为 \(1+\frac{x^2}{2}\)。

\textbf{结论与检查}\quad \(\boxed{a=0,b=\frac12}\)，选 D。\(\sec x\) 是偶函数，展开式没有一次项，与 \(a=0\) 一致。
\end{QuestionSolution}

% Q 2020 4
\QuestionPage{2.5 Taylor定理及其应用}{选择题}{2020 年第 4 题}{1}{由幂级展开求高阶导数值}{097 / 234}{**}
已知函数 \(f\left( x \right) = x^{2} \ln \left( 1-x \right)\)，当 \(n \geq 3\) 时，\(f^{\left( n \right)}\left( 0 \right) =\)
\choices{\(-\dfrac{n!}{n-2}\)}{\(\dfrac{n!}{n-2}\)}{\(-\dfrac{\left( n-2 \right)!}{n}\)}{\(\dfrac{\left( n-2 \right)!}{n}\)}
\QuestionEnd

\begin{QuestionSolution}{2020 年第 4 题}
\textbf{已知与目标}\quad 对 \(f=x^2\ln(1-x)\)，要求 \(n\ge3\) 时的 \(f^{(n)}(0)\)。

\textbf{方法与条件}\quad Maclaurin 展开是以零点为中心的 Taylor 展开，其 \(x^n\) 系数为 \(\frac{f^{(n)}(0)}{n!}\)。在 \(|x|<1\) 内，对数可展开为收敛幂级数；乘上 \(x^2\) 后，每项次数提高两次。

\textbf{推导}\quad 在 \(|x|<1\) 内，
\[
f(x)=x^2\ln(1-x)=-\sum_{k=1}^{\infty}\frac{x^{k+2}}{k}.
\]
要得到 \(x^n\)，必须 \(k+2=n\)，即 \(k=n-2\ge1\)。因此 \(x^n\) 的系数为 \(-\frac1{n-2}\)，有 \(\frac{f^{(n)}(0)}{n!}=-\frac1{n-2}\)。

\textbf{结论与检查}\quad \(\boxed{f^{(n)}(0)=-\frac{n!}{n-2}}\)，选 A。\(n=3\) 时最低次项为 \(-x^3\)，导数值应为 \(-3!=-6\)，与公式一致。
\end{QuestionSolution}

% Q 2026 12
\QuestionPage{2.5 Taylor定理及其应用}{填空题}{2026 年第 12 题}{2}{局部展开求函数极限}{100 / 234}{**}
\[\lim_{x\to 0}\left(\frac1x-\frac{\ln(1+x)}{x\sin x}\right)=\ \blank{}.
\]
\QuestionEnd

\begin{QuestionSolution}{2026 年第 12 题}
\textbf{已知与目标}\quad 要求两个均含零点奇异项的分式之差的极限。

\textbf{方法与条件}\quad 两个分式各自没有有限极限，先通分合并。分母为二阶小量，所以分子需展开到二次项。\(o(x^2)\) 是除以 \(x^2\) 后趋零的余项。

\textbf{推导}\quad 在零点附近取 \(x\ne0\)，对数与分母均有意义，通分得
\[
\frac1x-\frac{\ln(1+x)}{x\sin x}
=\frac{\sin x-\ln(1+x)}{x\sin x}.
\]
其中 \(\sin x=x+o(x^2)\)，\(\ln(1+x)=x-\frac{x^2}{2}+o(x^2)\)，故分子为 \(\frac{x^2}{2}+o(x^2)\)。分母 \(x\sin x=x^2(\frac{\sin x}{x})\)，除以 \(x^2\) 后趋于 \(1\)。因此原商趋于 \(\frac{\frac12}{1}=\frac12\)。

\textbf{结论与检查}\quad \(\boxed{\frac{1}{2}}\)。一次项先相消，留下的二次项为正，且两侧分母都等价于正的 \(x^2\)，结果符合双侧极限。
\end{QuestionSolution}

% Q 2009 15
\QuestionPage{2.5 Taylor定理及其应用}{解答题}{2009 年第 15 题}{3}{展开求极限}{101 / 234}{**}
求极限$\lim_{x \to 0} \frac{\left(1 - \cos x \right)\left[x - \ln(1 + \tan x) \right]}{\sin^4 x}$．
\QuestionEnd

\begin{QuestionSolution}{2009 年第 15 题}
\textbf{已知与目标}\quad 要计算 \((1-\cos x)[x-\ln(1+\tan x)]\) 与 \(\sin^4x\) 之比的极限。

\textbf{方法与条件}\quad 先判断次数：\(1-\cos x\) 与 \(x^2\) 同阶，分母与 \(x^4\) 同阶，因此另一因子只需保留二次项。Taylor 展开中的 \(o(x^2)\) 表示更高阶余项；零点附近 \(1+\tan x>0\)，对数有定义。

\textbf{推导}\quad 因 \(\tan x=x+o(x^2)\)，将 \(u=\tan x\) 代入 \(\ln(1+u)=u-\frac{u^2}{2}+o(u^2)\)，得
\[
\ln(1+\tan x)=\tan x-\frac{1}{2}\tan^2x+o(x^2)
=x-\frac{1}{2}x^2+o(x^2).
\]
因此 \(x-\ln(1+\tan x)=\frac{x^2}{2}+o(x^2)\)。又 \(1-\cos x=\frac{x^2}{2}+o(x^2)\)，\(\sin^4x=x^4(1+o(1))\)，故原商等于 \(\frac{\frac14x^4+o(x^4)}{x^4(1+o(1))}\to\frac14\)。

\textbf{结论与检查}\quad \(\boxed{\frac{1}{4}}\)。两个分子因子都是二阶小量，乘积恰为四阶，与分母同阶，所以结果有限且非零。
\end{QuestionSolution}

% Q 2018 1
\QuestionPage{2.5 Taylor定理及其应用}{选择题}{2018 年第 1 题}{1}{指数型极限与展开}{096 / 234}{***}
若$\lim_{x \to 0} \left(\mathrm{e}^x + ax^2 + bx \right)^{\frac{1}{x^2}} = 1$，则\pickout{}
\begin{abcd}
\item $a = \frac{1}{2}\;,\;b = - 1$．
\item $a = - \frac{1}{2}\;,\;b = - 1$．
\item $a = \frac{1}{2}\;,\;b = 1$．
\item $a = - \frac{1}{2}\;,\;b = 1$．
\end{abcd}
\QuestionEnd

\begin{QuestionSolution}{2018 年第 1 题}
\textbf{已知与目标}\quad 原幂的极限等于 \(1\)，要据此确定 \(a,b\)。

\textbf{方法与条件}\quad 底数趋于 \(1\)，在零点附近为正，可以取对数。原幂趋于 \(1\) 等价于其对数趋于 \(0\)。用 \(\ln(1+u)\sim u\) 后，要让底数减 \(1\) 的结果除以 \(x^2\) 趋零；\(o(x^2)\) 表示除以 \(x^2\) 后趋零的余项。

\textbf{推导}\quad 由指数展开 \(\mathrm e^x=1+x+\frac{x^2}{2}+o(x^2)\)，写底数为 \(1+u\)，其中
\[
u=(1+b)x+\left(\frac{1}{2}+a\right)x^2+o(x^2).
\]
因 \(\frac{\ln(1+u)}u\to1\)（\(u=0\) 处补值为 \(1\)），对数趋零等价于
\[\frac{u}{x^2}=\frac{1+b}{x}+\left(\frac12+a\right)+o(1)\to0.\]
若 \(1+b\ne0\)，第一项无有限双侧极限，所以先有 \(b=-1\)；此后常数项还须为零，得 \(a=-\frac12\)。

\textbf{结论与检查}\quad \(\boxed{a=-\frac{1}{2},\ b=-1}\)，选 B。代回底数为 \(1+o(x^2)\)，对数除以 \(x^2\) 趋零，原幂因此趋于 \(\mathrm e^0=1\)。
\end{QuestionSolution}

% Q 2014 15
\QuestionPage{2.5 Taylor定理及其应用}{解答题}{2014 年第 15 题}{2}{变上限积分与局部展开}{102 / 234}{***}
求极限$\lim_{x \to +\infty}
\frac{\int_1^x \left[t^2\left(\mathrm{e}^{\frac{1}{t}} - 1\right) - t \right]\,\mathrm{d} t}
     {x^2\ln \left(1 + \frac{1}{x}\right)}$．
\QuestionEnd

\begin{QuestionSolution}{2014 年第 15 题}
\textbf{已知与目标}\quad 分子为上限趋于无穷的积分，分母为 \(x^2\ln(1+\frac1x)\)，要求二者的比值极限。

\textbf{方法与条件}\quad 先用 \(\ln(1+u)\sim u\) 看出分母等价于 \(x\)。将被积函数在 \(\frac1t=0\) 处展开，可看出抵消后的常数极限；再对积分与 \(x\) 的比值用 L'Hospital 法则。

\textbf{推导}\quad 记被积函数为 \(h(t)\)。因 \(\mathrm e^u=1+u+\frac{u^2}{2}+o(u^2)\)，代入 \(u=\frac1t\)，当 \(t\to+\infty\) 时，
\[
t^2(\mathrm{e}^{\frac{1}{t}}-1)-t
=t^2\left(\frac{1}{t}+\frac{1}{2t^2}+o(t^{-2})\right)-t
=\frac{1}{2}+o(1).
\]
因 \(h(t)\to\frac12>0\)，充分大的 \(t\) 有 \(h(t)>\frac14\)，故积分趋于正无穷。\(h\) 连续，积分可导；分母 \(x\) 的导数为 \(1\ne0\)。对这个无穷比无穷型用 L'Hospital 法则，得
\[\lim_{x\to+\infty}\frac{\int_1^x h(t)\,\mathrm dt}{x}=\lim_{x\to+\infty}h(x)=\frac12.\]
同时 \(x\ln(1+\frac1x)\to1\)，所以原式等于上述商再除以一个趋于 \(1\) 的因子。

\textbf{结论与检查}\quad \(\boxed{\frac{1}{2}}\)。分子、分母均以 \(x\) 的一次量级增长，二者系数分别为 \(\frac12\)、\(1\)，与所得比值一致。
\end{QuestionSolution}

% Q 2016 15
\QuestionPage{2.5 Taylor定理及其应用}{解答题}{2016 年第 15 题}{2}{三角函数展开求极限}{103 / 234}{***}
求极限$\lim_{x \to 0} (\cos 2x + 2x\sin x)^{\frac{1}{x^4}}$．
\QuestionEnd

\begin{QuestionSolution}{2016 年第 15 题}
\textbf{已知与目标}\quad 底数趋于 \(1\)，指数 \(\frac1{x^4}\) 趋于正无穷，要求该幂的极限。

\textbf{方法与条件}\quad 幂的底数在零点附近为正，可先取对数。取对数后需将底数与 \(1\) 的差除以 \(x^4\)，故必须展开到四次项。\(o(x^4)\) 表示除以 \(x^4\) 后趋零的余项。

\textbf{推导}\quad 利用余弦展开到四次、正弦展开到三次，得
\[
\begin{aligned}
\cos2x&=1-\frac{(2x)^2}{2}+\frac{(2x)^4}{24}+o(x^4)=1-2x^2+\frac{2}{3}x^4+o(x^4),\\
2x\sin x&=2x\left(x-\frac{x^3}{6}+o(x^3)\right)=2x^2-\frac{1}{3}x^4+o(x^4).
\end{aligned}
\]
相加后两次项抵消，底数为 \(1+\frac{x^4}{3}+o(x^4)\)。由于 \(\ln(1+u)=u+o(u)\)，原幂的对数为
\[\frac{\ln(\cos2x+2x\sin x)}{x^4}=\frac{\frac{x^4}{3}+o(x^4)}{x^4}\to\frac13.\]

\textbf{结论与检查}\quad 用指数函数的连续性还原，原极限为 \(\boxed{\mathrm{e}^{\frac{1}{3}}}\)。底数趋于 \(1>0\)，取对数合法；最终四次项系数为正，所得值大于 \(1\)，与原幂的变化一致。
\end{QuestionSolution}

% Q 2021 17
\QuestionPage{2.5 Taylor定理及其应用}{解答题}{2021 年第 17 题}{2}{局部展开求极限}{104 / 234}{***}
求极限
\[\lim_{x \to 0} \left(\frac{1+\int_{0}^{x} \mathrm{e}^{t^{2}} \,\mathrm{d} t}{\mathrm{e}^{x}-1}-\frac{1}{\sin x}\right).
\]
\QuestionEnd

\begin{QuestionSolution}{2021 年第 17 题}
\textbf{已知与目标}\quad 要求两个各自发散的分式之差的极限，其中第一项的分子含变上限积分。

\textbf{方法与条件}\quad 先通分，再把分子展开到与分母相同的二次量级，避免分别处理发散项。\(O(x^3)\) 表示绝对值不超过常数倍 \(|x|^3\) 的余项。

\textbf{推导}\quad 记 \(I(x)=\int_0^x\mathrm e^{t^2}\,\mathrm dt\)。因 \(\mathrm e^{t^2}=1+O(t^2)\)，两侧积分都给出 \(I(x)=x+O(x^3)\)。通分得
\[
\frac{(1+I(x))\sin x-(\mathrm e^x-1)}{(\mathrm e^x-1)\sin x}.
\]
由 \(\sin x=x+O(x^3)\)、\(\mathrm e^x-1=x+\frac{x^2}{2}+O(x^3)\)，分子为
\[(1+x+O(x^3))(x+O(x^3))-\left(x+\frac{x^2}{2}+O(x^3)\right)=\frac{x^2}{2}+O(x^3).\]
分母为 \((x+O(x^2))(x+O(x^3))=x^2(1+O(x))\)，所以原商趋于 \(\frac12\)。

\textbf{结论与检查}\quad \(\boxed{\frac{1}{2}}\)。零点附近取 \(x\ne0\)，两个分母均非零；一次项相消，保留二次项已足够确定极限。
\end{QuestionSolution}

% Q 2022 3
\QuestionPage{2.5 Taylor定理及其应用}{选择题}{2022 年第 3 题}{1}{二阶展开与极限}{099 / 234}{****}
设 \(f\left( x \right)\) 在 \(x = x_{0}\) 处有二阶导数，则 \(\left( \ \right)\)
\choices{当 \(f\left( x \right)\) 在 \(x_{0}\) 的某邻域内单调增加时，\(f^{\prime}\left( x_{0} \right) > 0\)}{当 \(f^{\prime}\left( x_{0} \right) > 0\) 时，\(f\left( x \right)\) 在 \(x_{0}\) 的某邻域内单调增加}{当 \(f\left( x \right)\) 在 \(x_{0}\) 的某邻域内是凹函数时，\(f^{\prime\prime}\left( x_{0} \right) > 0\)}{当 \(f^{\prime\prime}\left( x_{0} \right) > 0\) 时，\(f\left( x \right)\) 在 \(x_{0}\) 的某邻域内是凹函数}
\QuestionEnd

\begin{QuestionSolution}{2022 年第 3 题}
\textbf{已知与目标}\quad 只知道 \(f\) 在 \(x_0\) 处有二阶导数，要辨别点处导数条件与邻域性态之间的推断。

\textbf{方法与条件}\quad \(f''(x_0)\) 存在，表示 \(f'\) 在 \(x_0\) 可导，因此 \(f'\) 在该点连续。正的一阶导数可由连续性扩展成邻域内恒正；题目未给 \(f''\) 连续，不能同样处理二阶导数。此处沿用题目约定：凹函数的图像在弦线下方。

\textbf{推导}\quad 对 B，因 \(f'\) 在 \(x_0\) 连续且 \(f'(x_0)>0\)，可取一个邻域，使其中 \(f'(x)>\frac12f'(x_0)>0\)。对邻域内任意 \(u<v\)，Lagrange 中值定理给出 \(f(v)-f(u)=f'(c)(v-u)>0\)，故严格增加。

对 A，取 \(x_0=0,f=x^3\)，函数严格增加但 \(f'(0)=0\)。对 C，取 \(f=x^4\)，其 \(f''=12x^2\ge0\)，图像满足凹函数的弦线不等式，但 \(f''(0)=0\)。对 D，取 \(f(x)=\frac{x^2}{2}+2x^4\sin(\frac1x)\)（\(x\ne0\)），\(f(0)=0\)。零点差商为 \(\frac x2+2x^3\sin(\frac1x)\to0\)，故 \(f'(0)=0\)；非零点处
\[f'(x)=x+8x^3\sin(\frac1x)-2x^2\cos(\frac1x),\qquad \frac{f'(x)-f'(0)}x\to1,\]
故 \(f''(0)=1>0\)。然而
\[
f''(x)=1+24x^2\sin(\frac{1}{x})-12x\cos(\frac{1}{x})-2\sin(\frac{1}{x}).
\]
取 \(x_k=\frac1{\frac\pi2+2k\pi}\to0\)，则 \(\sin(\frac1{x_k})=1\)、\(\cos(\frac1{x_k})=0\)，故充分大时 \(f''(x_k)=-1+24x_k^2<0\)。任意零点邻域内都有这样的点，D 不能成立。

\textbf{结论与检查}\quad 选 \(\boxed{\mathrm B}\)。三个反例分别满足对应前提；D 的反例在零点确实二阶可导，说明缺少的条件是邻域内二阶导数的控制。
\end{QuestionSolution}

\SectionHeading{2.6}{一元函数微分学及其应用}{函数性态的研究}{定理与要点\quad 教材 PDF 第 171 页起\hfill 本节 25 题}
\pdfbookmark[2]{2.6 函数性态的研究}{section-2.6}
\Theory{定理 6.1}{171}{单调性判别}{区间上连续且内部可导，单调增（减）等价于内部导数非负（非正）。严格正（负）可推出严格单调。}
\Theory{定理 6.2}{173}{极值的第一判别法}{驻点左侧导数非负、右侧非正时为极大值，反向为极小值。两侧导数严格同号时无极值。}
\Theory{定理 6.3}{174}{极值的第二判别法}{驻点处二阶导数为正则取极小值，为负则取极大值。}
\Theory{定理 6.4}{175}{高阶导数判别极值}{驻点的首个非零高阶导数若为偶数阶，正值对应极小、负值对应极大。奇数阶无极值。}
\Theory{定理 6.5}{179}{凹凸性判别}{按教材约定，区间内 \(f''(x)>0\) 时图像为凹（下凸），\(f''(x)<0\) 时图像为凸（上凸）。}
\Theory{要点}{175--182}{最值与拐点}{闭区间求最值要比较内部候选点和端点。判定拐点要核对曲线两侧凹凸性变化。}
% Q 2011 3
\QuestionPage{2.6 函数性态的研究}{选择题}{2011 年第 3 题}{1}{驻点个数}{105 / 234}{**}
函数$f(x) = \ln \left| (x - 1)(x - 2)(x - 3) \right|$的驻点个数为\pickout{}
\begin{abcd}
\item $0$．
\item $1$．
\item $2$．
\item $3$．
\end{abcd}
\QuestionEnd

\begin{QuestionSolution}{2011 年第 3 题}
\textbf{已知与目标}\quad 要求 \(f=\ln|(x-1)(x-2)(x-3)|\) 的驻点个数。驻点指函数可导且导数等于零的点。

\textbf{方法与条件}\quad 对非零的可导函数 \(P\)，有 \((\ln|P|)'=\frac{P'}P\)，所以在定义域内解 \(P'=0\) 即可；必须排除 \(P=0\) 的点。

\textbf{推导}\quad 定义域排除 \(x=1,2,3\)。令 \(u=x-2\)，则 \(P=(u+1)u(u-1)=u^3-u\)，且 \(\frac{\mathrm du}{\mathrm dx}=1\)，故
\[
f'(x)=\frac{3u^2-1}{u^3-u}.
\]
分子为零等价于 \(3u^2=1\)，有 \(u=\pm\frac1{\sqrt3}\)。它们都不等于 \(-1,0,1\)，分母非零。

\textbf{结论与检查}\quad 两个驻点为 \(x=2\pm\frac1{\sqrt3}\)，共有 \(\boxed{2}\) 个，选 C。三个被排除的点处原函数无定义，不能计为驻点。
\end{QuestionSolution}

% Q 2015 4
\QuestionPage{2.6 函数性态的研究}{选择题}{2015 年第 4 题}{1}{由二阶导数图像判定拐点}{107 / 234}{**}
设函数$f(x)$在$\left(-\infty , +\infty \right)$内连续，
其中二阶导数$f''(x)$的图形如图所示，则曲线$y = f(x)$的拐点的个数为\pickout{}
\par\centerline{\begin{tikzpicture}[scale=\DiagramScale,x=1pt,y=1pt,xscale=0.483697,yscale=-0.483697]
\path[use as bounding box] (0,0) rectangle (400,248);
\path[fill=black!100,even odd rule,shift={(158,16)}] svg {M0 0 C2.3 2.1 3.8 4.2 5.2 7 C5.5 7.7 5.9 8.4 6.2 9.1 C7 11 7 11 7 14 C5.3 14 3.7 14 2 14 C2 60.9 2 107.7 2 156 C8.3 155.7 14.5 155.3 21 155 C21.6 153.2 22.1 151.5 22.7 149.7 C31.8 121.2 41.5 99.1 69 84.8 C88.5 75.1 109.6 69.6 130.5 63.8 C135.4 62.4 140.4 61 145.4 59.6 C146 59.4 146.7 59.2 147.3 59 C150.4 58.1 153.5 57.2 156.6 56.3 C157.6 56 158.7 55.7 159.8 55.4 C160.7 55.1 161.7 54.9 162.6 54.6 C165 54 165 54 168 54 C168 54.7 168 55.3 168 56 C157.6 59.6 147.1 62.7 136.5 65.5 C72.2 82.7 72.2 82.7 53 99 C52.3 99.6 51.6 100.2 50.9 100.8 C35.5 114.6 29.8 137.1 24 156 C76.1 156 128.3 156 182 156 C182 154 182 152 182 150 C187 151.5 191.4 153.6 196 156 C192.3 159.1 186.9 163 182 163 C182 161.3 182 159.7 182 158 C129.5 158 77.1 158 23 158 C22.3 161.6 21.7 165.3 21 169 C20.2 172.7 19.5 176.4 18.7 180.1 C18.4 181.1 18.2 182.1 18 183.1 C17.5 185.2 17.1 187.3 16.6 189.4 C15.9 192.6 15.2 195.8 14.5 199 C14.1 201 13.7 203.1 13.2 205.1 C13 206.1 12.8 207.1 12.6 208 C12.4 208.9 12.2 209.8 12 210.7 C11.8 211.5 11.7 212.3 11.5 213.1 C11 215 11 215 10 217 C9.3 217 8.7 217 8 217 C9.7 205.7 12.2 194.6 14.6 183.5 C15 181.6 15.4 179.6 15.8 177.7 C16.2 175.8 16.6 174 17 172.2 C17.1 171.7 17.1 171.7 17.6 169.5 C17.7 168.7 17.9 167.9 18.1 167.1 C18.2 166.4 18.4 165.7 18.6 165 C19 163 19.5 161 20 159 C14.1 158.7 8.1 158.3 2 158 C2 180.1 2 202.2 2 225 C1.3 225 0.7 225 0 225 C0 202.9 0 180.8 0 158 C-44.2 158 -88.4 158 -134 158 C-134.3 157.3 -134.7 156.7 -135 156 C-110.9 155.7 -86.8 155.3 -62 155 C-64.6 152.4 -67.3 149.7 -70 147 C-71.7 144.7 -73.1 142.5 -74.5 140 C-74.9 139.4 -75.3 138.7 -75.6 138 C-79.2 131.7 -82.5 125.3 -85.6 118.8 C-87 115.9 -88.4 113.1 -89.8 110.3 C-91.2 107.5 -92.6 104.6 -94 101.8 C-94.6 100.4 -95.3 99 -96 97.6 C-97 95.5 -98 93.5 -99 91.5 C-99.6 90.3 -100.2 89.1 -100.8 87.9 C-102 85 -102 85 -102 82 C-97.9 85.5 -96 89.6 -93.7 94.4 C-93.3 95.2 -92.9 96 -92.5 96.9 C-91.3 99.5 -90 102.2 -88.8 104.8 C-68.7 146.8 -68.7 146.8 -58 154 C-51.5 154.9 -44.5 152.5 -39 149 C-29.8 141.5 -25.5 128.7 -22.5 117.6 C-22.3 116.8 -22 116 -21.8 115.2 C-19.8 107.8 -18.5 100.2 -17.3 92.5 C-17 90.2 -16.6 87.8 -16.2 85.5 C-15.7 82.2 -15.2 79 -14.7 75.7 C-13.7 69.1 -12.7 62.5 -11.6 56 C-11.3 54 -11 52.1 -10.6 50.2 C-10.5 49.1 -10.3 47.9 -10.1 46.8 C-9.9 45.8 -9.8 44.8 -9.6 43.7 C-9 41 -8.1 38.6 -7 36 C-6.3 36.3 -5.7 36.7 -5 37 C-5.1 37.3 -5.1 37.3 -5.3 38.7 C-7.3 49.2 -9.2 59.8 -10.7 70.4 C-10.8 70.9 -10.8 70.9 -11.1 73.3 C-11.7 77.2 -12.2 81 -12.8 84.9 C-14.3 95.7 -16.1 106.3 -19 116.9 C-19.2 117.8 -19.5 118.7 -19.7 119.6 C-23.8 134.1 -30.2 146.5 -43 155 C-28.8 155.3 -14.6 155.7 0 156 C0 109.1 0 62.3 0 14 C-1.7 14 -3.3 14 -5 14 C-5.3 13.3 -5.7 12.7 -6 12 C-4.8 9.2 -4.8 9.2 -3.1 5.9 C-2.5 4.8 -1.9 3.7 -1.3 2.6 C-0.9 1.7 -0.5 0.9 0 0 Z };
\path[fill=black!100,even odd rule,shift={(145.625,180.25)}] svg {M0 0 C0.6 0.1 0.6 0.1 3.4 0.8 C4.9 3.8 4.6 6.4 4.4 9.8 C2.8 13.8 1.4 16.9 -2.6 18.8 C-6.3 18.6 -6.3 18.6 -9.6 17.8 C-11.2 14.5 -11.1 11.2 -10.6 7.8 C-8 2.8 -6 -0.4 0 0 Z M-3.6 0.8 C-7.3 5.5 -8.4 8.7 -8.6 14.8 C-8.3 15.2 -8.3 15.2 -6.6 17.8 C-4.3 17.6 -4.3 17.6 -1.6 16.8 C1.9 12 1.7 7.4 1.4 1.8 C-0.3 1.4 -1.9 1.1 -3.6 0.8 Z };
\path[fill=black!100,even odd rule,shift={(135,19)}] svg {M0 0 C1.3 0 2.6 0 4 0 C4.3 2.6 4.7 5.3 5 8 C6.5 6.8 6.5 6.8 8 5 C8.2 2.3 8.2 2.3 8 0 C9.3 0 10.6 0 12 0 C10.7 4.6 8.9 8 6.1 11.9 C5.4 12.8 4.7 13.8 4 14.8 C3.3 15.5 2.7 16.3 2 17 C1 17 0 17 -1 17 C-0.5 16.6 0 16.3 0.5 15.9 C2.7 13.2 2.1 11.3 1.8 7.9 C1.6 6.9 1.5 5.8 1.4 4.7 C1 2 1 2 0 0 Z };
\path[fill=black!100,even odd rule,shift={(343,184)}] svg {M0 0 C1.9 0.2 1.9 0.2 4 1 C5.2 3.2 5.2 3.2 6.2 5.9 C6.5 6.8 6.8 7.7 7.2 8.7 C7.4 9.4 7.7 10.2 8 11 C8.3 11.7 8.7 12.3 9 13 C7.1 12.7 7.1 12.7 5 12 C4.3 11 3.7 10 3 9 C2.7 10.3 2.3 11.6 2 13 C0.7 13 -0.6 13 -2 13 C-0.8 10.5 0 9 2 7 C1.7 4.4 1.2 2.4 0 0 Z };
\path[fill=black!100,even odd rule,shift={(350,184)}] svg {M0 0 C1 0 2 0 3 0 C2.3 1.3 1.7 2.6 1 4 C0.7 2.7 0.3 1.4 0 0 Z };
\end{tikzpicture}}
\begin{abcd}
\item $0$．
\item $1$．
\item $2$．
\item $3$．
\end{abcd}
\QuestionEnd

\begin{QuestionSolution}{2015 年第 4 题}
\textbf{已知与目标}\quad 已知 \(f\) 连续，给出的是 \(f''\) 的图形，要数出 \(y=f(x)\) 的拐点。

\textbf{方法与条件}\quad 拐点是连续曲线上凹凸性发生变化的点。按本书约定，\(f''>0\) 时图像为凹（下凸），\(f''<0\) 时为凸（上凸）。候选位置包括二阶导数零点及不存在的点，关键是两侧是否变号。

\textbf{推导}\quad 负半轴的零点处，图中的 \(f''\) 只从上方触及横轴，两侧均为正，所以这里没有拐点。\(x=0\) 处二阶导数虽不存在，但左侧 \(f''>0\)、右侧 \(f''<0\)，凹凸性改变，且 \(f\) 连续，故原点对应一个拐点。正半轴的零点处，\(f''\) 从负变正，再产生一个拐点。

\textbf{结论与检查}\quad 共 \(\boxed{2}\) 个，选 C。单凭 \(f''=0\) 不能判定拐点；二阶导数不存在也不妨碍连续曲线在该处改变凹凸性。
\end{QuestionSolution}

% Q 2019 2
\QuestionPage{2.6 函数性态的研究}{选择题}{2019 年第 2 题}{1}{拐点坐标}{110 / 234}{**}
曲线$y = x\sin x + 2\cos x$（$- \frac{\pi}{2} < x < \frac{3\pi}{2}$）的拐点是\pickout{}
\begin{abcd}
\item $(0,2)$．
\item $(\pi , - 2)$．
\item $\left(\frac{\pi}{2},\frac{\pi}{2} \right)$．
\item $\left(\frac{3\pi}{2}, - \frac{3\pi}{2} \right)$．
\end{abcd}
\QuestionEnd

\begin{QuestionSolution}{2019 年第 2 题}
\textbf{已知与目标}\quad 在开区间 \((-\frac\pi2,\frac{3\pi}2)\) 上求 \(y=x\sin x+2\cos x\) 的拐点坐标。

\textbf{方法与条件}\quad 函数在区间内光滑，拐点只可能出现在二阶导数的零点，但还要检查两侧符号变化。拐点坐标需要把横坐标代回原函数。

\textbf{推导}\quad 用乘积法则，\(y'=\sin x+x\cos x-2\sin x=x\cos x-\sin x\)，再求导得 \(y''=-x\sin x\)。在所给区间内零点为 \(0,\pi\)。在 \((-\frac\pi2,0)\) 上 \(x,\sin x\) 都为负，\(y''<0\)；在 \((0,\pi)\) 上两者都为正，仍有 \(y''<0\)；在 \((\pi,\frac{3\pi}2)\) 上 \(x>0,\sin x<0\)，故 \(y''>0\)。只有 \(x=\pi\) 两侧变号。

\textbf{结论与检查}\quad 代入 \(y(\pi)=\pi\sin\pi+2\cos\pi=-2\)，得 \(\boxed{(\pi,-2)}\)，选 B。\(x=0\) 虽然二阶导数为零，但不改变凹凸性；\(\frac{3\pi}2\) 是被排除的端点。
\end{QuestionSolution}

% Q 2020 6
\QuestionPage{2.6 函数性态的研究}{选择题}{2020 年第 6 题}{1}{由导数不等式判断函数比值}{111 / 234}{**}
设函数 \(f\left( x \right)\) 在区间 \(\left[ -2,2 \right]\) 上可导，且 \(f^{\prime}\left( x \right) > f\left( x \right) > 0\)，则
\choices{\(\dfrac{f\left( -2 \right)}{f\left( -1 \right)} > 1\)}{\(\dfrac{f\left( 0 \right)}{f\left( -1 \right)} > \mathrm{e}\)}{\(\dfrac{f\left( 1 \right)}{f\left( -1 \right)} < \mathrm{e}^{2}\)}{\(\dfrac{f\left( 2 \right)}{f\left( -1 \right)} < \mathrm{e}^{3}\)}
\QuestionEnd

\begin{QuestionSolution}{2020 年第 6 题}
\textbf{已知与目标}\quad 已知 \(f'>f>0\)，要判断不同自变量处函数值的比与指数常数的关系。

\textbf{方法与条件}\quad 因 \(f>0\)，可除以 \(f\) 并取对数，\(\frac{f'}f\) 正好是 \(\ln f\) 的导数。将 \(\frac{f'}f>1\) 写成某个函数导数恒正，便能比较两点增量。

\textbf{推导}\quad 令 \(H=\ln f-x\)，则 \(H'=\frac{f'}f-1>0\)，由中值定理知 \(H\) 严格增加。对任意区间内的 \(a<b\)，有 \(H(b)>H(a)\)，整理、取指数得
\[
\ln f(b)-\ln f(a)>b-a,
\qquad \frac{f(b)}{f(a)}>\mathrm{e}^{b-a}.
\]
取 \(a=-1,b=0\)，得 \(\frac{f(0)}{f(-1)}>\mathrm e\)，即 B。取 \(a=-1,b=1,2\)，分别得到 \(\frac{f(1)}{f(-1)}>\mathrm e^2\)、\(\frac{f(2)}{f(-1)}>\mathrm e^3\)，排除 C、D。又 \(f'>0\) 给出 \(f(-2)<f(-1)\)，故 A 的比值小于 \(1\)。

\textbf{结论与检查}\quad 选 \(\boxed{\mathrm B}\)。所有分母均因 \(f>0\) 而非零，取对数、除法和指数还原均保持所需不等号方向。
\end{QuestionSolution}

% Q 2009 13
\QuestionPage{2.6 函数性态的研究}{填空题}{2009 年第 13 题}{2}{区间最值}{117 / 234}{**}
函数$y = x^{2x}$在区间$(0,1]$上的最小值为 \blank{}．
\QuestionEnd

\begin{QuestionSolution}{2009 年第 13 题}
\textbf{已知与目标}\quad 在半开区间 \((0,1]\) 上求 \(x^{2x}\) 的最小值。

\textbf{方法与条件}\quad 底数 \(x>0\)，可写成 \(x^{2x}=\mathrm e^{2x\ln x}\)。指数函数严格增加，所以只需求指数 \(h=2x\ln x\) 的最小值，再代入指数函数。

\textbf{推导}\quad \(h'=2(\ln x+1)\)，由 \(h'=0\) 得 \(x=\mathrm e^{-1}\)。当 \(0<x<\mathrm e^{-1}\) 时 \(h'<0\)，当 \(\mathrm e^{-1}<x\le1\) 时 \(h'>0\)，所以这里是全区间最小值点。代入 \(h(\mathrm e^{-1})=2\mathrm e^{-1}(-1)=-\frac2{\mathrm e}\)。

\textbf{结论与检查}\quad 最小值为 \(\boxed{\mathrm e^{-\frac2{\mathrm e}}}\)。右端点值为 \(1\)，左端虽不包括零点，但 \(x\to0^+\) 时 \(x\ln x\to0\)，函数也趋于 \(1\)，均大于内部所得值。
\end{QuestionSolution}

% Q 2018 10
\QuestionPage{2.6 函数性态的研究}{填空题}{2018 年第 10 题}{2}{拐点及其切线}{118 / 234}{**}
曲线$y = x^2 + 2\ln x$在其拐点处的切线方程是 \blank{}．
\QuestionEnd

\begin{QuestionSolution}{2018 年第 10 题}
\textbf{已知与目标}\quad 曲线为 \(y=x^2+2\ln x\)，要先找到拐点，再求该点的切线。

\textbf{方法与条件}\quad 对数要求 \(x>0\)。函数在定义域内光滑，二阶导数零点若两侧变号，则为拐点；切线斜率由一阶导数给出。

\textbf{推导}\quad \(y'=2x+\frac2x\)，\(y''=2-\frac2{x^2}=\frac{2(x^2-1)}{x^2}\)。因 \(x>0\)，唯一候选点为 \(x=1\)。在 \((0,1)\) 上 \(y''<0\)，在 \((1,+\infty)\) 上 \(y''>0\)，确有拐点。其纵坐标为 \(1+2\ln1=1\)，斜率为 \(2+2=4\)。

\textbf{结论与检查}\quad 切线 \(y-1=4(x-1)\)，即 \(\boxed{y=4x-3}\)。代入 \(x=1\) 得 \(y=1\)，且直线斜率 \(4\) 等于曲线在该点的一阶导数。
\end{QuestionSolution}

% Q 2017 18
\QuestionPage{2.6 函数性态的研究}{解答题}{2017 年第 18 题}{3}{隐函数的极值}{124 / 234}{**}
已知函数$y(x)$由方程$x^3 + y^3 - 3x + 3y - 2 = 0$确定，求$y(x)$的极值．
\QuestionEnd

\begin{QuestionSolution}{2017 年第 18 题}
\textbf{已知与目标}\quad \(y(x)\) 由三次方程隐式确定，要找出它的所有极值。

\textbf{方法与条件}\quad 将方程对 \(x\) 求导，\(y^3\) 的导数为 \(3y^2y'\)。极值由 \(y'\) 的变号确定，不必解出三次方程的一般表达式。

\textbf{推导}\quad 原式为 \(y^3+3y=2-x^3+3x\)。左侧关于 \(y\) 严格递增、值域为全实数，且导数 \(3(1+y^2)>0\)，所以每个 \(x\) 唯一确定一个可导的 \(y\)。求导得 \(3x^2+3y^2y'-3+3y'=0\)，即
\[
y'=\frac{1-x^2}{1+y^2}.
\]
分母恒正，因此 \(y'\) 在 \((-\infty,-1)\) 为负、\((-1,1)\) 为正、\((1,+\infty)\) 为负。故 \(x=-1\) 为极小值点，\(x=1\) 为极大值点。代入原式：\(x=-1\) 时 \(y^3+3y=0\)，唯一解 \(y=0\)；\(x=1\) 时 \(y^3+3y=4\)，唯一解 \(y=1\)。

\textbf{结论与检查}\quad \(\boxed{y(-1)=0}\) 为极小值，\(\boxed{y(1)=1}\) 为极大值。导函数只有 \(x=\pm1\) 两个零点且处处存在，故不存在其他极值候选点。
\end{QuestionSolution}

% Q 2014 3
\QuestionPage{2.6 函数性态的研究}{选择题}{2014 年第 3 题}{1}{凹凸性与弦线}{106 / 234}{***}
设函数$f(x)$具有二阶导数，$g(x) = f(0)(1 - x) + f(1)x$，则在区间$[0,1]$上\pickout{}
\begin{abcd}
\item 当$f'(x) \ge 0$时，$f(x) \ge g(x)$．
\item 当$f'(x) \ge 0$时，$f(x) \le g(x)$．
\item 当$f''(x) \ge 0$时，$f(x) \ge g(x)$．
\item 当$f''(x) \ge 0$时，$f(x) \le g(x)$．
\end{abcd}
\QuestionEnd

\begin{QuestionSolution}{2014 年第 3 题}
\textbf{已知与目标}\quad \(g(x)=f(0)(1-x)+f(1)x\) 是通过 \((0,f(0)),(1,f(1))\) 的直线，要求判断曲线与它的位置关系。

\textbf{方法与条件}\quad 弦线是曲线上两点连成的直线。减去弦线后误差两端为零，二阶导数不变。若二阶导数非负，一阶导数非减，可比较左右两段的平均斜率。

\textbf{推导}\quad 令 \(h=f-g\)，则 \(h(0)=h(1)=0,h''=f''\)。若 \(f''\ge0\)，则 \(h'\) 非减。对任意 \(0<x<1\)，\(h\) 在 \([0,x]\)、\([x,1]\) 连续且内部可导，由 Lagrange 中值定理得
\[
\frac{h(x)}x=h'(\xi)\le h'(\eta)=\frac{-h(x)}{1-x},
\qquad \xi<x<\eta.
\]
乘以正数 \(x(1-x)\)，得到 \((1-x)h(x)\le-xh(x)\)，即 \(h(x)\le0\)。端点本来取等号，所以整个闭区间都有 \(f\le g\)。对于只有 \(f'\ge0\) 的情形，\(f=x^2\) 在弦线 \(y=x\) 下方，而 \(f=2x-x^2\) 在其上方，分别排除 A、B；前者同时排除 C。

\textbf{结论与检查}\quad 选 \(\boxed{\mathrm D}\)。二阶导数控制弯曲方向，一阶导数非负本身只能控制增加方向，不能保证与弦线的固定位置关系。
\end{QuestionSolution}

% Q 2016 4
\QuestionPage{2.6 函数性态的研究}{选择题}{2016 年第 4 题}{1}{由导函数图像判断性态}{108 / 234}{***}
设函数$f(x)$在$(-\infty , +\infty)$内连续，其导函数的图形如图所示，则\pickout{}
\par\centerline{\begin{tikzpicture}[scale=\DiagramScale,x=1pt,y=1pt,xscale=0.439724,yscale=-0.439724]
\path[use as bounding box] (0,0) rectangle (440,348);
\path[fill=black!100,even odd rule,shift={(53,10)}] svg {M0 0 C2.6 3.1 4 6.1 5.5 9.9 C5.7 10.4 5.7 10.4 6.8 13.2 C7.3 14.4 7.7 15.6 8.2 16.8 C8.7 17.9 9.1 19.1 9.6 20.3 C10.7 23.2 11.9 26.1 13 29 C8.9 28.4 5.7 27 2 25 C1 58.6 1 58.6 8.1 66.4 C10 68.4 12 70.2 14 72 C56.6 124.7 64.4 207.8 63 273 C67.6 273 72.2 273 77 273 C77 272 77.1 271 77.1 270 C79.4 216.3 87.4 143.2 129 104 C142.1 92.8 155.2 89.2 172 90 C180.1 91.1 185.9 95.9 191 102 C205.3 123.3 207.2 151.9 211.1 176.7 C215.1 202.2 220 227.8 234.5 249.7 C236 252 236 252 236 254 C236.7 254 237.3 254 238 254 C239.4 255.5 239.4 255.5 241.1 257.6 C247.6 265.2 254.8 271 265 272.4 C273.3 271.2 279.2 266 284.2 259.5 C308.8 225.9 316.2 177.6 326 138 C327 138 328 138 329 138 C328.4 143.3 327.5 148.5 326.2 153.7 C326.1 154.1 326.1 154.1 325.7 155.9 C325.1 158.2 324.5 160.5 323.9 162.8 C323.7 163.6 323.5 164.4 323.3 165.2 C319.9 178.9 316.1 192.5 311.9 206 C311.3 208 310.7 209.9 310.2 211.8 C303.4 234.1 293.9 254.1 279 272 C295.5 272 312 272 329 272 C328.3 270.4 327.7 268.7 327 267 C326.3 265.3 325.7 263.7 325 262 C335.1 264.6 344.5 268.9 354 273 C354 273.7 354 274.3 354 275 C350.3 276.7 346.5 278.3 342.8 279.9 C341.8 280.4 340.7 280.9 339.6 281.4 C338.6 281.8 337.6 282.3 336.5 282.7 C335.6 283.1 334.6 283.5 333.7 284 C331.1 285 328.7 285.6 326 286 C327.1 282.5 328.3 279.3 330 276 C288.8 276.2 288.8 276.2 80 277 C79.7 286.9 79.3 296.8 79 307 C77.7 307 76.4 307 75 307 C75.3 297.1 75.7 287.2 76 277 C74.7 277 73.4 277 72 277 C71.7 277.7 71.3 278.3 71 279 C70 279 69 279 68 279 C67.3 278.3 66.7 277.7 66 277 C65.3 277 64.7 277 64 277 C64 277.3 64 277.3 63.9 278.7 C63.8 281.3 63.7 284 63.6 286.6 C63.5 287.5 63.5 288.4 63.4 289.3 C63.2 294.2 62.8 299.1 62 304 C61 304 60 304 59 304 C59 294.8 59 285.5 59 276 C56.3 276.2 53.5 276.3 50.7 276.5 C34.8 277.3 18.9 277.1 3 277 C2.7 293.5 2.3 310 2 327 C1 327 0 327 -1 327 C-1 310.5 -1 294 -1 277 C-6.3 277 -11.5 277 -16.9 277.1 C-28.7 277.1 -40.3 277 -52 276 C-52 275 -52 274 -52 273 C-35.5 273 -19 273 -2 273 C-1.9 245.3 -1.8 217.6 -1.8 189.9 C-1.8 177.1 -1.7 164.2 -1.7 151.4 C-1.7 140.2 -1.6 129 -1.6 117.8 C-1.6 111.8 -1.6 105.9 -1.6 100 C-1.5 94.4 -1.5 88.8 -1.5 83.2 C-1.5 81.2 -1.5 79.1 -1.5 77.1 C-1.5 74.3 -1.5 71.5 -1.5 68.7 C-1.5 68.3 -1.5 68.3 -1.5 66.2 C-1.5 61.2 -2.4 58.6 -6 55 C-7.6 53.5 -7.6 53.5 -9.1 52.4 C-9.6 52 -10.2 51.5 -10.8 51.1 C-11.1 50.9 -11.1 50.9 -12.5 49.8 C-13.1 49.4 -13.6 48.9 -14.2 48.5 C-17 46.4 -19.8 44.4 -23 43 C-22.7 42 -22.3 41 -22 40 C-18.1 40.6 -15.6 42.6 -12.4 44.9 C-9.7 47 -6.9 49.1 -4 51 C-4 51.7 -4 52.3 -4 53 C-3 53 -2 53 -1 53 C-1.2 52.3 -1.3 51.6 -1.5 50.8 C-2 48 -2.1 45.5 -2.1 42.7 C-2.1 41.7 -2.1 40.7 -2.1 39.7 C-2.1 39.2 -2.1 39.2 -2.1 36.7 C-2.1 35.7 -2.1 34.6 -2 33.6 C-2 31.1 -2 28.5 -2 26 C-5.5 27.1 -8.8 28.3 -12 30 C-10.2 23.1 -7.7 16.5 -4.9 9.9 C-4.5 9.1 -4.2 8.2 -3.8 7.3 C-1.1 1.1 -1.1 1.1 0 0 Z M2 64 C2 133 2 201.9 2 273 C20.8 273 39.6 273 59 273 C61.4 268.3 59.9 261.7 59.8 256.5 C59.7 255.8 59.7 255 59.7 254.3 C58.2 208.9 52.8 162.5 36 120 C35.7 119.3 35.4 118.5 35.1 117.8 C29.1 102.2 21.2 87.8 11.6 74.1 C11.2 73.5 10.8 72.9 10.4 72.4 C7.9 68.8 5.6 66.9 2 64 Z M137 102 C135.9 102.6 134.9 103.2 133.8 103.9 C129.9 106.8 127 110.2 124 114 C123.3 114.9 122.6 115.7 121.8 116.6 C99.1 145.4 91.2 185.6 85.8 220.9 C85.6 221.8 85.5 222.6 85.4 223.4 C84 232.1 82.8 240.9 81.8 249.7 C81.8 250 81.8 250 81.6 251.9 C80.9 258.5 80.6 265.3 81 272 C82 273 82 273 84.8 273.1 C86 273.1 87.3 273.1 88.5 273.1 C88.9 273.1 88.9 273.1 90.6 273.1 C93 273.1 95.3 273.1 97.6 273.1 C99.3 273.1 101 273.1 102.6 273.1 C107.2 273.1 111.7 273.1 116.2 273.1 C121 273.1 125.7 273.1 130.4 273.1 C139.4 273.1 148.4 273.1 157.3 273.1 C167.6 273.1 177.8 273.1 188 273.1 C209 273 230 273 251 273 C249.7 269 247.3 267.5 244 265 C243.3 265 242.7 265 242 265 C242 264.3 242 263.7 242 263 C240.6 261.4 240.6 261.4 239 260 C238.3 260 237.7 260 237 260 C236.8 259.5 236.5 258.9 236.2 258.4 C234.7 255.4 233 252.7 231 250 C230.3 250 229.7 250 229 250 C228.9 249.5 228.8 249 228.7 248.4 C227.8 245.2 226.3 242.4 224.8 239.4 C213.4 215.4 209.5 189.2 205.5 163.2 C198.1 114.8 198.1 114.8 186 101 C185.3 101 184.7 101 184 101 C184 100.3 184 99.7 184 99 C182.6 98.2 181.2 97.4 179.8 96.6 C179.4 96.3 179.4 96.3 177.4 95.2 C163.8 88.4 148.8 95 137 102 Z };
\path[fill=black!100,even odd rule,shift={(38,294)}] svg {M0 0 C3.3 2.8 4.8 4.7 5.4 9 C5.3 14.9 2.4 20.6 -1.8 24.8 C-3.6 26.2 -3.6 26.2 -7 28 C-8 27.7 -9 27.3 -10 27 C-9.3 26.3 -8.5 25.7 -7.8 25 C-2.5 19.7 -0.1 13.4 0 6 C-1 3.5 -1 3.5 -3 2 C-6.7 1.8 -6.7 1.8 -10 2 C-7.9 -2.2 -3.9 -1 0 0 Z };
\path[fill=black!100,even odd rule,shift={(28,295)}] svg {M0 0 C-1 2.1 -2.1 4.2 -3.3 6.3 C-6.3 11.8 -8.3 16.6 -7 23 C-4.6 25.4 -3.2 25.5 0 26 C0 26.7 0 27.3 0 28 C-7.3 27.5 -7.3 27.5 -10.4 25.3 C-12.8 21.8 -13.4 19.1 -12.6 14.9 C-11 9.2 -8.3 5 -4 1 C-2 0 -2 0 0 0 Z };
\path[fill=black!100,even odd rule,shift={(73,4)}] svg {M0 0 C2.3 0 4.6 0 7 0 C7.7 4.3 8.3 8.6 9 13 C10 12.3 11 11.7 12 11 C10.3 15.2 8.6 16.4 5 20 C4.3 14.4 3.7 8.8 3 3 C2 2.7 1 2.3 0 2 C0 1.3 0 0.7 0 0 Z };
\path[fill=black!100,even odd rule,shift={(421,276)}] svg {M0 0 C4.5 -0.4 4.5 -0.4 6.8 1.4 C8.8 4 9 4.9 8.9 8.3 C9 11.8 9.5 13.9 11 17 C9.7 17.3 8.4 17.7 7 18 C6.3 16.4 5.7 14.7 5 13 C4 13.3 3 13.7 2 14 C2.1 13.4 2.1 12.8 2.2 12.1 C2.3 11.7 2.3 11.7 2.5 9.7 C2.6 8.9 2.7 8.1 2.8 7.3 C3 5 3 5 3 2 C2 1.7 1 1.3 0 1 C0 0.7 0 0.3 0 0 Z };
\path[fill=black!100,even odd rule,shift={(87,4)}] svg {M0 0 C1 0 2 0 3 0 C2.5 6.4 2.5 6.4 -0.1 9.4 C-0.7 10 -1.3 10.5 -2 11 C-1.6 10.3 -1.3 9.6 -0.9 8.8 C0.1 5.7 -0.2 4.1 -1 1 C-0.7 0.7 -0.3 0.3 0 0 Z };
\path[fill=black!100,even odd rule,shift={(76,25)}] svg {M0 0 C0 1 0 2 0 3 C-2.5 5.4 -3.5 6 -7 6 C-6.7 4.7 -6.3 3.4 -6 2 C-5.6 1.9 -5.6 1.9 -3.6 1.6 C-1 1 -1 1 0 0 Z };
\path[fill=black!100,even odd rule,shift={(422,290)}] svg {M0 0 C0 1 0 2 0 3 C-1.6 3.3 -3.3 3.7 -5 4 C-5 3 -5 2 -5 1 C-3.4 0.7 -1.7 0.3 0 0 Z };
\path[fill=black!100,even odd rule,shift={(121,218)}] svg {M0 0 C1 0 2 0 3 0 C3 1.7 3 3.3 3 5 C2 5 1 5 0 5 C0 3.3 0 1.7 0 0 Z };
\path[fill=black!100,even odd rule,shift={(121,118)}] svg {M0 0 C1 0 2 0 3 0 C3 1.7 3 3.3 3 5 C2 5 1 5 0 5 C0 3.3 0 1.7 0 0 Z };
\path[fill=black!100,even odd rule,shift={(121,251)}] svg {M0 0 C1 0.3 2 0.7 3 1 C3 2.3 3 3.6 3 5 C2 5 1 5 0 5 C0 3.3 0 1.7 0 0 Z };
\path[fill=black!100,even odd rule,shift={(121,85)}] svg {M0 0 C1 0 2 0 3 0 C2.7 1.7 2.3 3.3 2 5 C1.3 5 0.7 5 0 5 C0 3.3 0 1.7 0 0 Z };
\path[fill=black!100,even odd rule,shift={(121,318)}] svg {M0 0 C1 0 2 0 3 0 C2.7 1.6 2.3 3.3 2 5 C1.3 4.7 0.7 4.3 0 4 C0 2.7 0 1.4 0 0 Z };
\path[fill=black!100,even odd rule,shift={(121,293)}] svg {M0 0 C1 0 2 0 3 0 C2.7 1.6 2.3 3.3 2 5 C1.3 4.7 0.7 4.3 0 4 C0 2.7 0 1.4 0 0 Z };
\path[fill=black!100,even odd rule,shift={(436,276)}] svg {M0 0 C-0.3 1 -0.7 2 -1 3 C-2.6 3 -4.3 3 -6 3 C-6 2.3 -6 1.7 -6 1 C-4 0 -4 0 0 0 Z };
\path[fill=black!100,even odd rule,shift={(121,210)}] svg {M0 0 C1 0 2 0 3 0 C2.7 1.7 2.3 3.3 2 5 C1.3 4.7 0.7 4.3 0 4 C0 2.7 0 1.4 0 0 Z };
\path[fill=black!100,even odd rule,shift={(121,193)}] svg {M0 0 C1 0.3 2 0.7 3 1 C3 2.3 3 3.6 3 5 C2 5 1 5 0 5 C0 3.3 0 1.7 0 0 Z };
\path[fill=black!100,even odd rule,shift={(121,135)}] svg {M0 0 C1 0 2 0 3 0 C2.7 1.7 2.3 3.3 2 5 C1.3 4.7 0.7 4.3 0 4 C0 2.7 0 1.4 0 0 Z };
\path[fill=black!100,even odd rule,shift={(121,310)}] svg {M0 0 C1 0 2 0 3 0 C3 1.3 3 2.6 3 4 C2 4 1 4 0 4 C0 2.7 0 1.4 0 0 Z };
\path[fill=black!100,even odd rule,shift={(121,302)}] svg {M0 0 C1 0 2 0 3 0 C3 1.3 3 2.6 3 4 C2 4 1 4 0 4 C0 2.7 0 1.4 0 0 Z };
\path[fill=black!100,even odd rule,shift={(121,277)}] svg {M0 0 C1 0 2 0 3 0 C3 1.3 3 2.6 3 4 C2 4 1 4 0 4 C0 2.7 0 1.4 0 0 Z };
\path[fill=black!100,even odd rule,shift={(121,268)}] svg {M0 0 C1 0 2 0 3 0 C3 1.3 3 2.6 3 4 C2 4 1 4 0 4 C0 2.7 0 1.4 0 0 Z };
\path[fill=black!100,even odd rule,shift={(121,260)}] svg {M0 0 C1 0 2 0 3 0 C3 1.3 3 2.6 3 4 C2 4 1 4 0 4 C0 2.7 0 1.4 0 0 Z };
\path[fill=black!100,even odd rule,shift={(122,243)}] svg {M0 0 C0.7 0 1.3 0 2 0 C1.7 1.7 1.3 3.3 1 5 C0.3 4.7 -0.3 4.3 -1 4 C-0.7 2.7 -0.3 1.4 0 0 Z };
\path[fill=black!100,even odd rule,shift={(121,235)}] svg {M0 0 C1 0 2 0 3 0 C3 1.3 3 2.6 3 4 C2 4 1 4 0 4 C0 2.7 0 1.4 0 0 Z };
\path[fill=black!100,even odd rule,shift={(121,227)}] svg {M0 0 C1 0 2 0 3 0 C3 1.3 3 2.6 3 4 C2 4 1 4 0 4 C0 2.7 0 1.4 0 0 Z };
\path[fill=black!100,even odd rule,shift={(121,202)}] svg {M0 0 C1 0 2 0 3 0 C3 1.3 3 2.6 3 4 C2 4 1 4 0 4 C0 2.7 0 1.4 0 0 Z };
\path[fill=black!100,even odd rule,shift={(121,185)}] svg {M0 0 C1 0 2 0 3 0 C3 1.3 3 2.6 3 4 C2 4 1 4 0 4 C0 2.7 0 1.4 0 0 Z };
\path[fill=black!100,even odd rule,shift={(121,177)}] svg {M0 0 C1 0 2 0 3 0 C3 1.3 3 2.6 3 4 C2 4 1 4 0 4 C0 2.7 0 1.4 0 0 Z };
\path[fill=black!100,even odd rule,shift={(121,169)}] svg {M0 0 C1 0 2 0 3 0 C3 1.3 3 2.6 3 4 C2 4 1 4 0 4 C0 2.7 0 1.4 0 0 Z };
\path[fill=black!100,even odd rule,shift={(122,160)}] svg {M0 0 C0.7 0 1.3 0 2 0 C1.7 1.7 1.3 3.3 1 5 C0.3 4.7 -0.3 4.3 -1 4 C-0.7 2.7 -0.3 1.4 0 0 Z };
\path[fill=black!100,even odd rule,shift={(121,152)}] svg {M0 0 C1 0 2 0 3 0 C3 1.3 3 2.6 3 4 C2 4 1 4 0 4 C0 2.7 0 1.4 0 0 Z };
\path[fill=black!100,even odd rule,shift={(121,144)}] svg {M0 0 C1 0 2 0 3 0 C3 1.3 3 2.6 3 4 C2 4 1 4 0 4 C0 2.7 0 1.4 0 0 Z };
\path[fill=black!100,even odd rule,shift={(121,127)}] svg {M0 0 C1 0 2 0 3 0 C3 1.3 3 2.6 3 4 C2 4 1 4 0 4 C0 2.7 0 1.4 0 0 Z };
\path[fill=black!100,even odd rule,shift={(121,110)}] svg {M0 0 C1 0 2 0 3 0 C3 1.3 3 2.6 3 4 C2 4 1 4 0 4 C0 2.7 0 1.4 0 0 Z };
\path[fill=black!100,even odd rule,shift={(121,102)}] svg {M0 0 C1 0 2 0 3 0 C3 1.3 3 2.6 3 4 C2 4 1 4 0 4 C0 2.7 0 1.4 0 0 Z };
\path[fill=black!100,even odd rule,shift={(121,94)}] svg {M0 0 C1 0 2 0 3 0 C3 1.3 3 2.6 3 4 C2 4 1 4 0 4 C0 2.7 0 1.4 0 0 Z };
\path[fill=black!100,even odd rule,shift={(121,77)}] svg {M0 0 C1 0 2 0 3 0 C3 1.3 3 2.6 3 4 C2 4 1 4 0 4 C0 2.7 0 1.4 0 0 Z };
\path[fill=black!100,even odd rule,shift={(121,69)}] svg {M0 0 C1 0 2 0 3 0 C3 1.3 3 2.6 3 4 C2 4 1 4 0 4 C0 2.7 0 1.4 0 0 Z };
\path[fill=black!100,even odd rule,shift={(122,60)}] svg {M0 0 C0.7 0.3 1.3 0.7 2 1 C1.7 2.3 1.3 3.6 1 5 C0.3 5 -0.3 5 -1 5 C-0.7 3.4 -0.3 1.7 0 0 Z };
\path[fill=black!100,even odd rule,shift={(121,52)}] svg {M0 0 C1 0 2 0 3 0 C3 1.3 3 2.6 3 4 C2 4 1 4 0 4 C0 2.7 0 1.4 0 0 Z };
\path[fill=black!100,even odd rule,shift={(121,44)}] svg {M0 0 C1 0 2 0 3 0 C3 1.3 3 2.6 3 4 C2 4 1 4 0 4 C0 2.7 0 1.4 0 0 Z };
\path[fill=black!100,even odd rule,shift={(432,290)}] svg {M0 0 C0.3 0.2 0.3 0.2 2 1 C1.7 1.2 1.7 1.2 0 2 C0 1.3 0 0.7 0 0 Z };
\end{tikzpicture}}
\begin{abcd}
\item 函数$f(x)$有$2$个极值点，曲线$y = f(x)$有$2$个拐点．
\item 函数$f(x)$有$2$个极值点，曲线$y = f(x)$有$3$个拐点．
\item 函数$f(x)$有$3$个极值点，曲线$y = f(x)$有$1$个拐点．
\item 函数$f(x)$有$3$个极值点，曲线$y = f(x)$有$2$个拐点．
\end{abcd}
\QuestionEnd

\begin{QuestionSolution}{2016 年第 4 题}
\textbf{已知与目标}\quad \(f\) 连续，图中给出其导函数。要分别统计原函数的极值点和拐点。

\textbf{方法与条件}\quad \(f'\) 的正负控制 \(f\) 的增减，正变负对应极大值，负变正对应极小值；这类变化也可发生在 \(f'\) 不存在但 \(f\) 连续的点。\(f'\) 的增减控制原曲线的凹凸性，增减方向改变的位置对应拐点。

\textbf{推导}\quad 从左向右，导函数先穿过横轴由正变负，给出一个极大值点。到图中的竖直虚线位置，\(f'\) 左侧为负、右侧为正；由于 \(f\) 连续，原函数由减转增，给出一个极小值点。最右的横轴接触点两侧 \(f'\) 都为正，原函数保持增加，没有极值。

导函数的增减方向依次为减、增、减、增，改变方向的位置共三处：竖直虚线处、右侧峰顶和右侧谷底。每处两侧的凹凸性均改变，而 \(f\) 连续，所以共有三个拐点。

\textbf{结论与检查}\quad \(f\) 有 \(\boxed{2}\) 个极值点，曲线有 \(\boxed{3}\) 个拐点，选 B。计数时分别检查导函数的正负和增减，横轴接触点虽无极值，仍可能因导函数由减转增而对应拐点。
\end{QuestionSolution}

% Q 2018 4
\QuestionPage{2.6 函数性态的研究}{选择题}{2018 年第 4 题}{1}{凹凸性与积分平均值}{109 / 234}{***}
设函数$f(x)$在$[0,1]$上二阶可导，且$\int_0^1 f(x)\,\mathrm{d} x = 0$，则\pickout{}
\begin{abcd}
\item 当$f'(x) < 0$时，$f\left(\frac{1}{2} \right) < 0$．
\item 当$f''(x) < 0$时，$f\left(\frac{1}{2} \right) < 0$．
\item 当$f'(x) > 0$时，$f\left(\frac{1}{2} \right) < 0$．
\item 当$f''(x) > 0$时，$f\left(\frac{1}{2} \right) < 0$．
\end{abcd}
\QuestionEnd

\begin{QuestionSolution}{2018 年第 4 题}
\textbf{已知与目标}\quad 已知 \(\int_0^1 f=0\)，要求哪一种导数条件能保证中点值 \(f(\frac12)<0\)。

\textbf{方法与条件}\quad \(f''>0\) 时曲线严格位于任意弦线下方。取关于 \(\frac12\) 对称的两点，它们连线中点的纵坐标为两函数值的平均值，因而可与 \(f(\frac12)\) 比较。

\textbf{推导}\quad 若 \(f''>0\)，对 \(0\le x<\frac12\) 的两个不同点 \(x,1-x\)，有
\[
f(x)+f(1-x)>2f(\frac{1}{2}).
\]
在 \([0,\frac12]\) 积分，并对 \(f(1-x)\) 换元，有
\[0=\int_0^{\frac12}[f(x)+f(1-x)]\,\mathrm dx>\int_0^{\frac12}2f(\frac12)\,\mathrm dx=f(\frac12).\]
故 D 正确。若只有一阶导数符号，取 \(f=\pm(x-\frac12)\)，满足相应增加或减少条件且积分为零，中点值却为零，排除 A、C。若 \(f''<0\)，弦线不等式反向，积分后反而有 \(f(\frac12)>0\)，排除 B。

\textbf{结论与检查}\quad 选 \(\boxed{\mathrm D}\)。严格不等式来自除中点外的每对不同点，连续性保证积分后仍严格，不会只得到 \(f(\frac12)\le0\)。
\end{QuestionSolution}

% Q 2021 4
\QuestionPage{2.6 函数性态的研究}{选择题}{2021 年第 4 题}{1}{单调性与零点个数}{112 / 234}{***}
设函数 \(f(x)=ax-b\ln x\)，其中 \(a>0\)。若 \(f(x)\) 有两个零点，则 \(\dfrac{b}{a}\) 的取值范围是
\choices{\(\left(\mathrm{e},+\infty\right)\)}{\(\left(0,\mathrm{e}\right)\)}{\(\left(0,\dfrac{1}{\mathrm{e}}\right)\)}{\(\left(\dfrac{1}{\mathrm{e}},+\infty\right)\)}
\QuestionEnd

\begin{QuestionSolution}{2021 年第 4 题}
\textbf{已知与目标}\quad \(f=ax-b\ln x\)，其中 \(a>0\)，要使它恰有两个不同零点，求 \(\frac ba\) 的范围。

\textbf{方法与条件}\quad 先按 \(b\) 的符号研究单调性。若函数两端都趋于正无穷且只有一个极小值点，则极小值为负时两侧各有一根；极小值为零时只有一个根。

\textbf{推导}\quad 定义域为 \(x>0\)，\(f'=a-\frac bx\)。若 \(b\le0\)，\(f'>0\)，严格增加，至多一个零点。若 \(b>0\)，当 \(0<x<\frac ba\) 时 \(f'<0\)，当 \(x>\frac ba\) 时 \(f'>0\)，故唯一极小值点为 \(\frac ba\)。这时 \(x\to0^+\) 时 \(-b\ln x\to+\infty\)；\(x\to+\infty\) 时因 \(\frac{\ln x}x\to0\)，也有 \(f\to+\infty\)。由连续性和两侧的严格单调性，两个零点的充要条件为
\[
f(\frac{b}{a})=b\bigl[1-\ln(\frac{b}{a})\bigr]<0.
\]
因 \(b>0\)，这等价于 \(\ln\frac ba>1\)，即 \(\frac ba>\mathrm e\)。

\textbf{结论与检查}\quad \(\boxed{\frac ba\in(\mathrm e,+\infty)}\)，选 A。边界 \(\frac ba=\mathrm e\) 时最小值恰为零，只产生一个零点，因此不能包含等号。
\end{QuestionSolution}

% Q 2023 6
\QuestionPage{2.6 函数性态的研究}{选择题}{2023 年第 6 题}{1}{积分函数的最小值}{113 / 234}{***}
若函数
\[f(a)=\int_2^{+\infty}\frac{1}{x(\ln x)^{a+1}}\,\mathrm{d}x
\]
在 \(a=a_0\) 处取得最小值，则 \(a_0=\)（ ）。
\choices{\(-\dfrac{1}{\ln(\ln 2)}\)}{\(-\ln(\ln 2)\)}{\(\dfrac{1}{\ln 2}\)}{\(\ln 2\)}
\QuestionEnd

\begin{QuestionSolution}{2023 年第 6 题}
\textbf{已知与目标}\quad \(f(a)\) 是一个以 \(a\) 为参数的无穷区间积分，要找出它取得最小值时的参数。

\textbf{方法与条件}\quad 反常积分先检查收敛，收敛参数范围才是函数的定义域。令 \(u=\ln x\)，用幂函数积分求出 \(f(a)\)，再对这个显式表达式求导。

\textbf{推导}\quad 因 \(\mathrm du=\frac{\mathrm dx}x\)，积分化为 \(\int_{\ln2}^{+\infty}u^{-a-1}\,\mathrm du\)。幂函数尾积分要求指数 \(-a-1<-1\)，即 \(a>0\)。此时
\[f(a)=\left[-\frac{u^{-a}}a\right]_{\ln2}^{+\infty}=\frac{(\ln2)^{-a}}a>0.\]
取对数 \(\ln f=-a\ln(\ln2)-\ln a\)，求导得
\[
\frac{f'(a)}{f(a)}=-\ln(\ln2)-\frac1a.
\]
令导数为零，得 \(a_0=-\frac1{\ln(\ln2)}\)。由于 \(0<\ln2<1\)，有 \(\ln(\ln2)<0\)，故 \(a_0>0\)。\(-\ln(\ln2)-\frac1a\) 在 \((0,a_0)\) 为负、\((a_0,+\infty)\) 为正，\(f\) 同样先减后增。

\textbf{结论与检查}\quad \(\boxed{a_0=-\frac{1}{\ln(\ln2)}}\)，选 A。这是整个收敛域上的最小值点；\(a\le0\) 时积分不收敛，不能作为最值候选参数。
\end{QuestionSolution}

% Q 2023 7
\QuestionPage{2.6 函数性态的研究}{选择题}{2023 年第 7 题}{2}{极值点与拐点的参数条件}{114 / 234}{***}
设函数
\[f(x)=(x^2+a)\mathrm{e}^x,
\]
若 \(f(x)\) 没有极值点，但曲线 \(y=f(x)\) 有拐点，则 \(a\) 的取值范围是（ ）。
\choices{\([0,1)\)}{\([1,+\infty)\)}{\([1,2)\)}{\([2,+\infty)\)}
\QuestionEnd

\begin{QuestionSolution}{2023 年第 7 题}
\textbf{已知与目标}\quad \(f=(x^2+a)\mathrm e^x\)，要求同时满足无极值点、有拐点的参数范围。

\textbf{方法与条件}\quad 函数处处光滑。极值要求一阶导数在候选点两侧改变符号，拐点要求二阶导数两侧改变符号。指数因子恒正，所以配方后的二次多项式决定符号；重根处通常不变号，需单独检查边界。

\textbf{推导}\quad 求导并配方，得
\[
f'(x)=\mathrm{e}^x\bigl[(x+1)^2+a-1\bigr],\qquad
f''(x)=\mathrm{e}^x\bigl[(x+2)^2+a-2\bigr].
\]
若 \(a<1\)，\((x+1)^2=1-a\) 有两个不同实根，一阶导数符号依次为正、负、正，产生两个极值点。若 \(a\ge1\)，一阶导数恒非负，最多在 \(a=1,x=-1\) 时为零且不变号，故没有极值点。因此第一条件是 \(a\ge1\)。

若 \(a<2\)，\((x+2)^2=2-a\) 有两个不同实根，二阶导数在两根处都变号，产生拐点；若 \(a\ge2\)，二阶导数非负，不改变凹凸性。因此第二条件是 \(a<2\)。

\textbf{结论与检查}\quad 取交集得 \(\boxed{1\le a<2}\)，选 C。\(a=1\) 的一阶重根不产生极值，故包含；\(a=2\) 的二阶重根不产生拐点，故排除。
\end{QuestionSolution}

% Q 2025 2
\QuestionPage{2.6 函数性态的研究}{选择题}{2025 年第 2 题}{1}{极值点和拐点的定义}{115 / 234}{***}
已知函数 \(f\left( x \right) = \int_{0}^{x} \mathrm{e}^{t^{2}} \sin t \,\mathrm{d} t\)，\(g\left( x \right) = \int_{0}^{x} \mathrm{e}^{t^{2}} \,\mathrm{d} t \cdot \sin^{2} x\)，则
\choices{\(x = 0\) 是 \(f\left( x \right)\) 的极值点，也是 \(g\left( x \right)\) 的极值点.}{\(x = 0\) 是 \(f\left( x \right)\) 的极值点，\(\left( 0, 0 \right)\) 是曲线 \(y = g\left( x \right)\) 的拐点.}{\(x = 0\) 是 \(f\left( x \right)\) 的极值点，\(\left( 0, 0 \right)\) 是曲线 \(y = f\left( x \right)\) 的拐点.}{\(\left( 0, 0 \right)\) 是曲线 \(y = f\left( x \right)\) 的拐点，也是曲线 \(y = g\left( x \right)\) 的拐点.}
\QuestionEnd

\begin{QuestionSolution}{2025 年第 2 题}
\textbf{已知与目标}\quad 要判断两个积分函数在零点是否取得极值，以及原点是否为拐点。

\textbf{方法与条件}\quad 积分的被积函数光滑，故 \(f,g\) 均可连续多次求导。极值检查一阶导数变号，拐点检查二阶导数变号。若 \(g''(0)=0,g'''(0)>0\)，由导数定义可得 \(g''(x)\) 在零点附近与 \(x\) 同号。

\textbf{推导}\quad 由微积分基本定理，\(f'=\mathrm e^{x^2}\sin x\)，其符号在零点左负右正，故零点为严格极小值点。再求导，\(f''=\mathrm e^{x^2}(2x\sin x+\cos x)\)，其零点值为 \(1\)，连续性保证零点附近恒正，所以原点不是 \(f\) 的拐点。

令 \(H=\int_0^x\mathrm e^{t^2}\,\mathrm dt\)、\(p=\sin^2x\)，则 \(g=Hp\)。有 \(H(0)=0,H'(0)=1\)，\(p(0)=p'(0)=0,p''(0)=2\)。所以
\[g''=H''p+2H'p'+Hp'',\qquad g'''=H'''p+3H''p'+3H'p''+Hp'''.\]
代入零点得 \(g''(0)=0,g'''(0)=3\cdot1\cdot2=6\)，故 \(g''(x)=6x+o(x)\)，在零点两侧变号。又 \(g(0)=0\)，所以原点是 \(g\) 的拐点。\(H\) 与 \(x\) 同号，而 \(p>0\) 在充分小的非零点成立，所以 \(g\) 左负右正，零点也不是 \(g\) 的极值点。

\textbf{结论与检查}\quad 选 \(\boxed{\mathrm B}\)：零点是 \(f\) 的极小值点，\((0,0)\) 是 \(g\) 的拐点。极值与拐点分别通过不同阶导数检查，未将两种性质混同。
\end{QuestionSolution}

% Q 2010 15
\QuestionPage{2.6 函数性态的研究}{解答题}{2010 年第 15 题}{2}{变上限积分的单调性与极值}{119 / 234}{***}
求函数$f(x) = \int_1^{x^2} \left(x^2 - t \right) \mathrm{e}^{- t^2} \,\mathrm{d} t$的单调区间与极值．
\QuestionEnd

\begin{QuestionSolution}{2010 年第 15 题}
\textbf{已知与目标}\quad 积分的上限和被积函数都含 \(x\)，要研究整个实数域上的单调区间和极值。

\textbf{方法与条件}\quad 可先拆成 \(f=x^2\int_1^{x^2}\mathrm e^{-t^2}\,\mathrm dt-\int_1^{x^2}t\mathrm e^{-t^2}\,\mathrm dt\)，再用乘积法则与变上限积分求导。被积函数连续，求导合法。极值由导数在候选点两侧的符号变化决定。

\textbf{推导}\quad 求导时，由上限变化产生的两项为 \(x^2\mathrm e^{-x^4}2x\) 和 \(-x^2\mathrm e^{-x^4}2x\)，相互抵消，故
\[
f'(x)=2x\int_1^{x^2}\mathrm{e}^{-t^2}\,\mathrm{d}t.
\]
积分的符号与 \(x^2-1\) 相同。故 \(f'\) 在 \((-\infty,-1)\)、\((0,1)\) 上为负，在 \((-1,0)\)、\((1,+\infty)\) 上为正，这四个区间依次为递减、递增、递减、递增区间。在 \(x=\pm1\) 处导数由负变正，极小值均为 \(0\)。在 \(x=0\) 处导数由正变负，极大值为
\[
f(0)=\int_0^1 t\mathrm e^{-t^2}\,\mathrm dt
=\left[-\frac12\mathrm e^{-t^2}\right]_0^1=\frac{1-\mathrm e^{-1}}{2}.
\]
\textbf{结论与检查}\quad 递减区间为 \(\boxed{(-\infty,-1),(0,1)}\)，递增区间为 \(\boxed{(-1,0),(1,+\infty)}\)。极小值 \(\boxed{f(\pm1)=0}\)，极大值 \(\boxed{f(0)=\frac{1-\mathrm e^{-1}}2}\)。原函数只含 \(x^2\)，是偶函数，所得单调性与极值关于零点对称。
\end{QuestionSolution}

% Q 2011 16
\QuestionPage{2.6 函数性态的研究}{解答题}{2011 年第 16 题}{2}{参数曲线的极值和拐点}{120 / 234}{***}
设函数$y = y(x)$由参数方程$\begin{cases}
  x = \frac{1}{3}t^3 + t + \frac{1}{3}, \\
  y = \frac{1}{3}t^3 - t + \frac{1}{3},
\end{cases}$确定，求$y = y(x)$的极值和曲线$y = y(x)$的凹凸区间及拐点．
\QuestionEnd

\begin{QuestionSolution}{2011 年第 16 题}
\textbf{已知与目标}\quad 曲线用参数 \(t\) 表示，要求 \(y(x)\) 的极值、凹凸区间及拐点，最终要换回 \(x\) 的区间。

\textbf{方法与条件}\quad \(x'(t)=1+t^2>0\)，且 \(x(t)\) 的值域为整个实数轴，所以 \(t\) 的增加顺序等于 \(x\) 的增加顺序。参数二阶导数为 \(\frac1{x'}\frac{\mathrm d}{\mathrm dt}(\frac{y'}{x'})\)。本题采用 \(y''>0\) 为凹（下凸）、\(y''<0\) 为凸（上凸）的约定。

\textbf{推导}\quad 由 \(y'(t)=t^2-1\)，并对一阶导数比求导再除以 \(1+t^2\)，得
\[
\frac{\mathrm{d}y}{\mathrm{d}x}=\frac{t^2-1}{t^2+1},\qquad
\frac{\mathrm{d}^2y}{\mathrm{d}x^2}=\frac{4t}{(1+t^2)^3}.
\]
一阶导数的分母恒正，分子 \(t^2-1\) 在 \(t=-1\) 处由正变负，在 \(t=1\) 处由负变正。代入参数方程：\(t=-1\) 对应 \((x,y)=(-1,1)\)，\(t=1\) 对应 \((\frac53,-\frac13)\)，故分别取得极大值和极小值。

二阶导数与 \(t\) 同号。\(t=0\) 对应 \(x=y=\frac13\)，其左右二阶导数由负变正。因 \(x(t)\) 严格增加，\(t<0\) 对应 \(x<\frac13\)，\(t>0\) 对应 \(x>\frac13\)。

\textbf{结论与检查}\quad 极大值 \(\boxed{y(-1)=1}\)，极小值 \(\boxed{y(\frac53)=-\frac13}\)；凸区间 \(\boxed{(-\infty,\frac13)}\)，凹区间 \(\boxed{(\frac13,+\infty)}\)，拐点 \(\boxed{(\frac13,\frac13)}\)。所有参数导数分母均为正，未遗漏竖直切线或定义域断点。
\end{QuestionSolution}

% Q 2012 20
\QuestionPage{2.6 函数性态的研究}{解答题}{2012 年第 20 题}{3}{用函数性态证明不等式}{121 / 234}{***}
证明：$x\ln \frac{1 + x}{1 - x} + \cos x \ge 1 + \frac{x^2}{2}$（$- 1 < x < 1$）．
\QuestionEnd

\begin{QuestionSolution}{2012 年第 20 题}
\textbf{已知与目标}\quad 要在 \(-1<x<1\) 内证明对数项与余弦项之和的下界，并确定何时取等号。

\textbf{方法与条件}\quad \(\ln\frac{1+x}{1-x}\) 是奇函数，乘上 \(x\) 后为偶函数，余弦和右侧也是偶函数，所以只需证明 \(0\le x<1\)。用导数或积分分别给出两个部分的下界，再相加。

\textbf{推导}\quad 对 \(0\le t<1\)，\(\frac2{1-t^2}\ge2\)。将对数写成导数的积分，得到
\[
\ln\frac{1+x}{1-x}=\int_0^x\frac{2}{1-t^2}\,\mathrm{d}t\ge 2x .
\]
另一方面，\((t-\sin t)'=1-\cos t\ge0\)，且零点值为零，所以 \(\sin t\le t\)。于是 \(1-\cos x=\int_0^x\sin t\,\mathrm dt\le\frac{x^2}{2}\)，即 \(\cos x\ge1-\frac{x^2}{2}\)。因 \(x\ge0\)，第一项的不等式可乘 \(x\)，相加得到
\[
x\ln\frac{1+x}{1-x}+\cos x\ge 2x^2+1-\frac{x^2}{2}
\ge 1+\frac{x^2}{2}.
\]
\textbf{结论与检查}\quad 由偶性，结论对整个 \((-1,1)\) 成立。\(x=0\) 时两侧均为 \(1\)，取得等号；\(x\ne0\) 时上面的 \(1+\frac32x^2\) 严格大于目标 \(1+\frac12x^2\)，故只有 \(\boxed{x=0}\) 取等号。区间限制也保证两个对数真数均为正。
\end{QuestionSolution}

% Q 2015 19
\QuestionPage{2.6 函数性态的研究}{解答题}{2015 年第 19 题}{3}{函数零点的个数}{122 / 234}{***}
已知函数$f(x) = \int_x^1 \sqrt{1 + t^2} \,\mathrm{d} t + \int_1^{x^2} \sqrt{1 + t} \,\mathrm{d} t$，求$f(x)$的零点个数．
\QuestionEnd

\begin{QuestionSolution}{2015 年第 19 题}
\textbf{已知与目标}\quad 函数为两个变限积分之和，要确定整个定义域内零点的个数。

\textbf{方法与条件}\quad 被积函数在各积分路径上连续，故可用微积分基本定理求导。导数的符号可给出每段零点数的上限，再找函数值异号的点保证根存在；存在性和唯一性都需说明。

\textbf{推导}\quad 因 \(x^2\ge0\)，第二积分始终位于 \(t\ge0\) 的范围，故函数对所有实数 \(x\) 有定义。第一个积分的变化端点是下限，求导带负号；第二个上限为 \(x^2\)，需乘 \(2x\)。因此
\[
f'(x)=-\sqrt{1+x^2}+2x\sqrt{1+x^2}=(2x-1)\sqrt{1+x^2}.
\]
根式恒正，所以函数在 \((-\infty,\frac12)\) 严格递减，在 \((\frac12,+\infty)\) 严格递增，每段至多一根。\(f(1)=0\)，故右段已有且只有一个根；由严格增加知 \(f(\frac12)<f(1)=0\)。而 \(f(-1)=\int_{-1}^1\sqrt{1+t^2}\,\mathrm dt>0\)，连续函数的零点定理保证 \((-1,\frac12)\) 也有一根，严格递减保证该根唯一。

\textbf{结论与检查}\quad 总计 \(\boxed{2}\) 个零点。唯一的单调分界点 \(\frac12\) 的函数值严格为负，本身不额外产生零点；两段已覆盖全实数域。
\end{QuestionSolution}

% Q 2016 16
\QuestionPage{2.6 函数性态的研究}{解答题}{2016 年第 16 题}{2}{含绝对值的积分与最小值}{123 / 234}{***}
设函数$f(x) = \int_0^1 \left| t^2 - x^2 \right|\,\mathrm{d} t$（$x > 0$），求$f'(x)$并求$f(x)$的最小值．
\QuestionEnd

\begin{QuestionSolution}{2016 年第 16 题}
\textbf{已知与目标}\quad 对 \(x>0\)，要求含绝对值积分的导函数及全局最小值。

\textbf{方法与条件}\quad 积分变量是 \(t\)，\(x\) 是固定参数。因 \(t,x\ge0\)，\(t^2-x^2\) 在 \(t=x\) 处变号；该点位于积分区间内与否，决定两种表达式。分段求导后还要检查接合点 \(x=1\)。

\textbf{推导}\quad 当 \(0<x<1\) 时，\(0\le t<x\) 取 \(x^2-t^2\)，\(x<t\le1\) 取 \(t^2-x^2\)，所以
\[
\begin{aligned}
f(x)&=\int_0^x(x^2-t^2)\,\mathrm{d}t+\int_x^1(t^2-x^2)\,\mathrm{d}t\\
&=\left[x^2t-\frac{t^3}{3}\right]_0^x+\left[\frac{t^3}{3}-x^2t\right]_x^1
=\frac{4}{3}x^3-x^2+\frac{1}{3}.
\end{aligned}
\]
对 \(x\ge1\)，\(f(x)=\int_0^1(x^2-t^2)\,\mathrm{d}t=x^2-\frac{1}{3}\)。两段在 \(x=1\) 接合，左右导数均为 \(2\)，故
\[
f'(x)=\begin{cases}
2x(2x-1),&0<x<1,\\
2x,&x\ge1.
\end{cases}
\]
导数在 \((0,\frac12)\) 为负、在 \((\frac12,+\infty)\) 为正，故 \(\frac12\) 为全局最小值点。代入第一段，
\[f(\frac12)=\frac43\cdot\frac18-\frac14+\frac13=\frac16-\frac14+\frac13=\frac14.\]

\textbf{结论与检查}\quad 导函数为上述分段式，最小值为 \(\boxed{\frac14}\)。两段函数在 \(x=1\) 的值均为 \(\frac23\)，且两侧导数均为 \(2\)，因此接合点也已正确覆盖。
\end{QuestionSolution}

% Q 2018 18
\QuestionPage{2.6 函数性态的研究}{解答题}{2018 年第 18 题}{3}{构造函数证明不等式}{125 / 234}{***}
已知常数$k \ge \ln 2 - 1$，证明：$(x - 1)\left(x - \ln^2 x + 2k\ln x - 1 \right) \ge 0$．
\QuestionEnd

\begin{QuestionSolution}{2018 年第 18 题}
\textbf{已知与目标}\quad 已知 \(k\ge\ln2-1\)，要证明含 \(\ln x\) 的乘积非负，定义域为 \(x>0\)。

\textbf{方法与条件}\quad 若括号中的函数 \(h\) 单调非减且 \(h(1)=0\)，则它与 \(x-1\) 同号或为零，乘积非负。要保证 \(h'\ge0\)，可先求 \(xh'\) 中不含参数部分的最小值；由于 \(x>0\)，乘 \(x\) 不改变导数符号。

\textbf{推导}\quad 令 \(h=x-(\ln x)^2+2k\ln x-1\)，则 \(h(1)=0\)，且
\[h'=1-\frac{2\ln x}{x}+\frac{2k}{x},\qquad xh'=x-2\ln x+2k.\]
令 \(p=x-2\ln x\)，有 \(p'=1-\frac2x\)，在 \((0,2)\) 为负、\((2,+\infty)\) 为正，所以全局最小值为 \(p(2)=2-2\ln2\)。于是
\[
xh'(x)\ge 2-2\ln2+2k\ge 0.
\]
最后一个不等式由 \(k\ge\ln2-1\) 得到，故 \(h'\ge0\)。在 \(0<x<1\) 时有 \(h(x)\le h(1)=0\)，两因子均非正；在 \(x>1\) 时有 \(h(x)\ge0\)，两因子均非负。

\textbf{结论与检查}\quad \(\boxed{(x-1)h(x)\ge0\quad(x>0)}\)。\(x=1\) 时乘积为零；当参数取边界 \(k=\ln2-1\) 时仍有 \(h'\ge0\)，所以无需排除边界。
\end{QuestionSolution}

% Q 2019 15
\QuestionPage{2.6 函数性态的研究}{解答题}{2019 年第 15 题}{2}{求导与函数极值}{126 / 234}{***}
已知函数$f(x) = \begin{cases}
     x^{2x}, & x > 0, \\
  x\mathrm{e}^x + 1, & x \le 0.
\end{cases}$求$f'(x)$，并求函数$f(x)$的极值．
\QuestionEnd

\begin{QuestionSolution}{2019 年第 15 题}
\textbf{已知与目标}\quad 函数在零点两侧采用不同表达式，要求导函数和全部极值，接合点也需检查。

\textbf{方法与条件}\quad 正半轴的幂指函数先写成指数形式求导，负半轴用乘积法则。分段点处只能用差商定义。找极值时，除了导数为零的点，还要检查函数连续但不可导的接合点。

\textbf{推导}\quad 当 \(x<0\) 时，\(f'=(x+1)\mathrm e^x\)。当 \(x>0\) 时，\(x^{2x}=\mathrm e^{2x\ln x}\)，故 \(f'=2(\ln x+1)x^{2x}\)。\(f(0)=1\)，左侧函数值趋于 \(1\)，右侧因 \(x\ln x\to0\) 也趋于 \(1\)，所以零点连续。零点左差商为 \(\mathrm e^x\to1\)；右差商为
\[
\frac{x^{2x}-1}{x}
=\frac{\mathrm{e}^{2x\ln x}-1}{2x\ln x}\,2\ln x\longrightarrow-\infty,
\]
其中第一因子趋于 \(1\)，\(2\ln x\to-\infty\)，故无有限右导数，\(f'(0)\) 不存在。负半轴上导数与 \(x+1\) 同号，正半轴上与 \(\ln x+1\) 同号，故函数依次在 \((-\infty,-1)\) 递减、\((-1,0)\) 递增、\((0,\mathrm e^{-1})\) 递减、\((\mathrm e^{-1},+\infty)\) 递增。

\textbf{结论与检查}\quad 导函数为 \(\boxed{f'=(x+1)\mathrm e^x\ (x<0),\quad f'=2(\ln x+1)x^{2x}\ (x>0)}\)，零点不可导。极小值为 \(\boxed{f(-1)=1-\mathrm e^{-1}}\)、\(\boxed{f(\mathrm e^{-1})=\mathrm e^{-\frac2{\mathrm e}}}\)，极大值为 \(\boxed{f(0)=1}\)。零点两侧均趋于并小于 \(1\)，所以不可导并不妨碍此处取得极大值。
\end{QuestionSolution}

% Q 2021 18
\QuestionPage{2.6 函数性态的研究}{解答题}{2021 年第 18 题}{3}{凹凸区间与渐近线}{127 / 234}{***}
已知
\[f(x)=\frac{x|x|}{1+x},
\]
求 \(f(x)\) 的凹凸性及渐近线。
\QuestionEnd

\begin{QuestionSolution}{2021 年第 18 题}
\textbf{已知与目标}\quad 要求含绝对值的有理函数的凹凸区间和全部渐近线。定义域排除 \(x=-1\)，绝对值还在 \(0\) 处分段。

\textbf{方法与条件}\quad 二阶导数正时图像为凹（下凸），负时为凸（上凸）。渐近线需检查函数与直线的差是否趋零；有限点处若函数趋于无穷，则相应竖直直线为铅直渐近线。

\textbf{推导}\quad 因 \(x^2=(x+1)(x-1)+1\)，令 \(r=x-1+\frac1{1+x}\)，则 \(f=r\)（\(x\ge0\)），\(f=-r\)（\(x<0,x\ne-1\)）。求导有 \(r'=1-\frac1{(1+x)^2}\)、\(r''=\frac2{(1+x)^3}\)，所以
\[
f''(x)=\begin{cases}-\frac2{(1+x)^3},&x<0,\ x\ne-1,\\[2pt]\frac2{(1+x)^3},&x>0.\end{cases}
\]
逐段判断，\((-\infty,-1)\)、\((0,+\infty)\) 上为正，\((-1,0)\) 上为负。零点连续且两侧凹凸性改变，所以 \((0,0)\) 是拐点；\(-1\) 不在曲线上，不能作为拐点。

在 \(-1\) 附近用负半轴表达式 \(f=-x+1-\frac1{1+x}\)，当 \(x\to-1^-\) 时趋于 \(+\infty\)，当 \(x\to-1^+\) 时趋于 \(-\infty\)。在无穷远处，\(f-(x-1)=\frac1{1+x}\to0\)（\(x\to+\infty\)），\(f-(-x+1)=-\frac1{1+x}\to0\)（\(x\to-\infty\)）。

\textbf{结论与检查}\quad 凹区间为 \(\boxed{(-\infty,-1),(0,+\infty)}\)，凸区间为 \(\boxed{(-1,0)}\)。渐近线为
\[\boxed{x=-1,\quad y=x-1,\quad y=-x+1}.\]
两个方向函数均无界，故没有水平渐近线，其他有限点函数连续，无额外铅直渐近线。
\end{QuestionSolution}

% Q 2026 5
\QuestionPage{2.6 函数性态的研究}{选择题}{2026 年第 5 题}{1}{单调性、极值与凹凸性}{116 / 234}{****}
设函数 \(f(x)\) 在区间 \([-1,1]\) 上有定义，则
\choices{当 \(f(x)\) 在 \((-1,0)\) 单调递减、在 \((0,1)\) 单调递增时，\(f(0)\) 是极小值}{当 \(f(0)\) 是极小值时，\(f(x)\) 在 \((-1,0)\) 单调递减、在 \((0,1)\) 单调递增}{当 \(f(x)\) 的图形在 \([-1,1]\) 上是凹的时，\(\dfrac{f(x)-f(1)}{x-1}\) 在 \([-1,1)\) 上单调递增}{当 \(\dfrac{f(x)-f(1)}{x-1}\) 在 \([-1,1)\) 上单调递增时，\(f(x)\) 的图形在 \([-1,1]\) 上是凹的}
\QuestionEnd

\begin{QuestionSolution}{2026 年第 5 题}
\textbf{已知与目标}\quad 只知道 \(f\) 在区间内有定义，没有连续或可导条件，要判断关于极值、单调性和凹性的四个推断。

\textbf{方法与条件}\quad 此处凹图像采用教材的弦线约定：两点之间的曲线不高于两点连线。差商 \(\frac{f(x)-f(1)}{x-1}\) 是固定右端点 \((1,f(1))\) 的割线斜率，可直接用弦线不等式比较，不需要导数。

\textbf{推导}\quad 对 C，取 \(-1\le x<y<1\)。因 \(y=\frac{1-y}{1-x}x+\frac{y-x}{1-x}\cdot1\)，凹性给出 \(f(y)\le\frac{1-y}{1-x}f(x)+\frac{y-x}{1-x}f(1)\)，即
\[
f(1)-f(y)\ge \frac{1-y}{1-x}\,[f(1)-f(x)].
\]
除以 \(1-y>0\)，得 \(\frac{f(y)-f(1)}{y-1}\ge\frac{f(x)-f(1)}{x-1}\)，即差商单调递增。

对 A，取 \(f(x)=|x|\)（\(x\ne0\)）、\(f(0)=2\)，两侧单调性成立，但零点不是极小值，说明缺少连续性时不能连接两侧。对 B，取 \(f=x^2(1-x^2)\)，零点附近 \(f>0=f(0)\)，但 \(f'=2x(1-2x^2)\) 在 \((\frac1{\sqrt2},1)\) 为负，故并非在整个右半区间递增。对 D，取 \(f=(x-1)x^3\)，固定端点差商为 \(x^3\)，严格递增；但 \(f''=12x^2-6x<0\) 在 \((0,\frac12)\) 成立，不满足所要求的凹性。

\textbf{结论与检查}\quad 选 \(\boxed{\mathrm C}\)。C 的证明仅用函数值和弦线关系，符合题设没有可导性的范围；三个反例均在整个 \([-1,1]\) 有定义。
\end{QuestionSolution}

% Q 2025 21
\QuestionPage{2.6 函数性态的研究}{解答题}{2025 年第 21 题}{3}{导函数单调与函数凹凸性}{129 / 234}{****}
（本题满分 12 分）
设函数 \(f\left( x \right)\) 在区间 \(\left( a, b \right)\) 内可导，证明：导函数 \(f^{\prime}\left( x \right)\) 在 \(\left( a, b \right)\) 内严格单调递增的充分必要条件是：对 \(\left( a, b \right)\) 内任意的 \(x_{1}\)，\(x_{2}\)，\(x_{3}\)，当 \(x_{1} < x_{2} < x_{3}\) 时，有

\(\frac{f\left( x_{2} \right) - f\left( x_{1} \right)}{x_{2} - x_{1}} < \frac{f\left( x_{3} \right) - f\left( x_{2} \right)}{x_{3} - x_{2}}\).
\QuestionEnd

\begin{QuestionSolution}{2025 年第 21 题}
\textbf{已知与目标}\quad \(f\) 在开区间可导，要证明 \(f'\) 严格增加与任意三点的相邻割线斜率严格增加互相推出。

\textbf{方法与条件}\quad 割线斜率是两个函数值之差除以横坐标之差。充分性用 Lagrange 中值定理将割线斜率写为导数。必要性先令短割线趋于切线，得到导数非减，再排除导数在一段区间取相同值的情况；取极限一般只能保留非严格不等式，因此最后一步必不可少。

\textbf{推导}\quad 充分性：任取 \(x_1<x_2<x_3\)，可导保证在两段闭区间上连续、内部可导。分别应用 Lagrange 中值定理，两个割线斜率为 \(f'(\xi_1),f'(\xi_2)\)，且 \(x_1<\xi_1<x_2<\xi_2<x_3\)。由 \(f'\) 严格增加，前者严格小于后者。

必要性：固定任意 \(a<u<v<b\)，取充分小的 \(\varepsilon>0\)，使 \(u-\varepsilon,v+\varepsilon\) 仍在原区间内。分别把条件用在 \(u-\varepsilon<u<v\)、\(u<v<v+\varepsilon\)，得到
\[\frac{f(u)-f(u-\varepsilon)}\varepsilon<\frac{f(v)-f(u)}{v-u}<\frac{f(v+\varepsilon)-f(v)}\varepsilon.\]
由可导性令 \(\varepsilon\to0^+\)，得
\[
f'(u)\le\frac{f(v)-f(u)}{v-u}\le f'(v).
\]
所以 \(f'\) 非减。若某对 \(u<v\) 满足 \(f'(u)=f'(v)=c\)，则任意 \(t\in[u,v]\) 有 \(c=f'(u)\le f'(t)\le f'(v)=c\)，故 \(f'\) 在此恒为 \(c\)。取区间内三个不同点，每段由中值定理的割线斜率都等于 \(c\)，与题设严格不等式矛盾。因此任意 \(u<v\) 都有 \(f'(u)<f'(v)\)。

\textbf{结论与检查}\quad 两个方向成立，充要性得证。必要性明确排除了导数的常值区间，最终得到严格增加；证明没有使用题设未提供的二阶导数或导函数连续性。
\end{QuestionSolution}

% Q 2022 21
\QuestionPage{2.6 函数性态的研究}{解答题}{2022 年第 21 题}{4}{二阶导数与凹凸性}{128 / 234}{*****}
设函数 \(f\left( x \right)\) 在 \(\left( -\infty,+\infty \right)\) 上有二阶连续导数，证明：
\(f^{\prime\prime}\left( x \right) \ge 0\)
的充分必要条件是对任意不同的实数 \(a,b\)，都有
\[f\left( \frac{a+b}{2} \right) \le \frac{1}{b-a}\int_{a}^{b} f\left( x \right) \,\mathrm{d}x\]
成立.
\QuestionEnd

\begin{QuestionSolution}{2022 年第 21 题}
\textbf{已知与目标}\quad \(f\) 二阶连续可导，要证明二阶导数处处非负与任意区间的中点值不超过积分平均值互相推出。

\textbf{方法与条件}\quad 充分性将区间内关于中点对称的函数值配对，再积分。必要性则把满足条件的区间缩到任意一点，利用二阶 Taylor 展开读取积分平均误差中的二阶导数；须控制余项积分，不能直接忽略它。

\textbf{推导}\quad 先说明交换 \(a,b\) 时积分和分母同时变号，商不变，所以可用中点 \(m=\frac{a+b}{2}\) 和半长 \(h=\frac{|b-a|}{2}>0\) 表示区间。

充分性：若 \(f''\ge0\)，由中值定理 \(f'\) 非减。对 \(0\le t\le h\)，微积分基本定理给出
\[
f(m+t)+f(m-t)-2f(m)
=\int_0^t[f'(m+s)-f'(m-s)]\,\mathrm{d}s\ge0.
\]
积分中 \(m+s\ge m-s\)，所以被积函数非负。再对 \(t\) 从 \(0\) 积分至 \(h\)，得到 \(\int_{m-h}^{m+h}f(x)\,\mathrm dx\ge2hf(m)\)，除以 \(2h>0\)，就是题目所要求的不等式。

必要性：假设所有不同端点都满足不等式。固定任意实数 \(m\)，取端点 \(m-h,m+h\)。二阶 Taylor 公式为
\[f(m+t)=f(m)+f'(m)t+\frac12f''(m)t^2+r(t),\qquad r(t)=o(t^2).\]
这里 \(r(t)=o(t^2)\) 表示 \(\frac{r(t)}{t^2}\to0\)。在 \([-h,h]\) 积分时，一次项因奇性为零；二次项的平均值为 \(\frac1{2h}\frac12f''(m)\frac{2h^3}{3}=\frac{h^2}{6}f''(m)\)。因此
\[
\frac{1}{2h}\int_{-h}^h f(m+t)\,\mathrm{d}t-f(m)
=\frac{h^2}{6}f''(m)+o(h^2).
\]
余项的控制如下：任给 \(\varepsilon>0\)，充分小的 \(h\) 保证所有 \(|t|\le h\) 都有 \(|r(t)|\le\varepsilon t^2\)，所以
\[\left|\frac1{2h}\int_{-h}^h r(t)\,\mathrm dt\right|\le\frac\varepsilon{2h}\frac{2h^3}{3}=\frac\varepsilon3h^2,\]
这确实是 \(o(h^2)\)。由假设，平均误差非负；除以 \(h^2>0\) 后令 \(h\to0^+\)，得到 \(\frac16f''(m)\ge0\)。

\textbf{结论与检查}\quad \(m\) 任意，所以 \(f''\ge0\) 处处成立，两个方向均已证明。积分比较覆盖任意长区间，缩区间证明覆盖任意中心；\(h>0\) 来自 \(a\ne b\)，所有平均值除法均合法。
\end{QuestionSolution}

\SectionHeading{3.1}{一元函数积分学及其应用}{定积分的概念、存在条件与性质}{定理与要点\quad 教材 PDF 第 189 页起\hfill 本节 11 题}
\pdfbookmark[1]{一元函数积分学及其应用}{chapter-3.1}
\pdfbookmark[2]{3.1 定积分的概念、存在条件与性质}{section-3.1}
\Theory{定理 1.1}{195}{可积的充要条件}{有界函数可积，当且仅当分割足够细时，各小区间的振幅乘区间长度之和可任意小。}
\Theory{定理 1.2}{196}{连续函数可积}{闭区间上的连续函数 Riemann 可积。}
\Theory{定理 1.3}{196}{有限间断点及单调情形}{有界函数若只有有限个第一类间断点，或在闭区间上单调，则可积。}
\Theory{性质 1.1--1.5}{197--199}{定积分的性质}{积分具有线性、保序、绝对值估计和区间可加性。可积函数的乘积仍可积。}
\Theory{性质 1.6}{199}{积分中值定理}{若 \(f\) 在 \([a,b]\) 上连续，\(g\) 在 \([a,b]\) 上可积且不变号，则存在 \(\xi\in[a,b]\)，使 \(\int_a^b f(x)g(x)\,\mathrm{d}x=f(\xi)\int_a^b g(x)\,\mathrm{d}x\)。}
\Theory{推论 1.2}{200}{积分平均值}{取权函数恒为 \(1\)，则连续函数在闭区间上的积分平均值等于该区间某一点的函数值。}
\Theory{方法}{192--195}{从和式识别定积分}{把和式写成取样函数值乘小区间长度，再找分割的区间和取样点。}
% Q 2021 7
\QuestionPage{3.1 定积分的概念、存在条件与性质}{选择题}{2021 年第 7 题}{1}{定积分与 Riemann 和}{132 / 234}{*}
设函数 \(f(x)\) 在区间 \(\left[0,1\right]\) 上连续，则
\[\int_{0}^{1} f(x)\,\mathrm{d}x=
\]
\choices{\(\lim_{n \to \infty}\sum_{k=1}^{n} f\left(\dfrac{2k-1}{2n}\right)\dfrac{1}{2n}\)}{\(\lim_{n \to \infty}\sum_{k=1}^{n} f\left(\dfrac{2k-1}{2n}\right)\dfrac{1}{n}\)}{\(\lim_{n \to \infty}\sum_{k=1}^{2n} f\left(\dfrac{k-1}{2n}\right)\dfrac{1}{n}\)}{\(\lim_{n \to \infty}\sum_{k=1}^{2n} f\left(\dfrac{k}{2n}\right)\dfrac{2}{n}\)}
\QuestionEnd

\begin{QuestionSolution}{2021 年第 7 题}
\textbf{已知与目标}\quad 已知 \(f\) 在 \([0,1]\) 连续，要从四个和式中找出该区间的定积分。

\textbf{方法与条件}\quad Riemann 和是“每段取一个函数值，乘该段长度，再相加”。连续函数可积，分割的最大长度趋于零时，这种和的极限等于积分。

\textbf{推导}\quad 将 \([0,1]\) 等分成 \(n\) 段，第 \(k\) 段为 \([\frac{k-1}{n},\frac{k}{n}]\)，长度为 \(\frac{1}{n}\)，中点为 \(\frac{2k-1}{2n}\)。相应的中点和为
\[
\sum_{k=1}^n f\!\left(\frac{2k-1}{2n}\right)\frac1n
\]
其极限就是题设积分，故选 B。

\textbf{结论与检查}\quad 正确选项为 \(\boxed{\mathrm{B}}\)。取连续函数 \(f(x)\equiv1\) 检查，A、C、D 分别得到 \(\frac{1}{2},2,4\)，均不等于 \(\int_0^1 1\,\mathrm{d}x=1\)。
\end{QuestionSolution}

% Q 2016 10
\QuestionPage{3.1 定积分的概念、存在条件与性质}{填空题}{2016 年第 10 题}{2}{三角和式极限}{134 / 234}{**}
极限$\lim_{n \to \infty} \frac{1}{n^2}\left(\sin\frac{1}{n} + 2\sin\frac{2}{n}
+ \cdots + n\sin\frac{n}{n}\right) =$ \blank{}．
\QuestionEnd

\begin{QuestionSolution}{2016 年第 10 题}
\textbf{已知与目标}\quad 求带权和式的极限，权重 \(k\) 与正弦自变量 \(\frac{k}{n}\) 都依赖项数 \(n\)。

\textbf{方法与条件}\quad Riemann 和指函数值乘小区间宽度的总和。把 \(\frac{k}{n^2}\) 拆为 \(\frac{1}{n}\frac{k}{n}\)，可识别连续函数 \(h(x)=x\sin x\) 在 \([0,1]\) 上的积分。

\textbf{推导}\quad 原和式写为
\[
\sum_{k=1}^n\frac1n\,\frac{k}{n}\sin\frac{k}{n}.
\]
取样点 \(\frac{k}{n}\) 是第 \(k\) 个小区间的右端点，宽度为 \(\frac{1}{n}\)。由连续性，其极限为 \(\int_0^1x\sin x\,\mathrm{d}x\)。分部积分取 \(u=x\)、\(\mathrm{d}v=\sin x\,\mathrm{d}x\)，得到
\[
\int_0^1 x\sin x\,\mathrm{d}x
=\bigl[-x\cos x+\sin x\bigr]_0^1
=\boxed{\sin1-\cos1}.
\]
\textbf{结论与检查}\quad 极限为 \(\boxed{\sin1-\cos1}\)。由于 \(0<x<1\) 时 \(x\sin x>0\)，结果应为正，\(\sin1-\cos1\approx0.3012\) 符合这一检查。
\end{QuestionSolution}

% Q 2024 15
\QuestionPage{3.1 定积分的概念、存在条件与性质}{填空题}{2024 年第 15 题}{2}{速度的积分平均值}{135 / 234}{**}
某物体以速度
\[v(t)=t+k\sin\pi t
\]
做直线运动。若它从 \(t=0\) 到 \(t=3\) 的平均速度是 \(\dfrac52\)，则 \(k=\) \blank{}
\QuestionEnd

\begin{QuestionSolution}{2024 年第 15 题}
\textbf{已知与目标}\quad 速度为 \(v(t)=t+k\sin\pi t\)，\([0,3]\) 上的平均速度为 \(\frac{5}{2}\)，要求常数 \(k\)。

\textbf{方法与条件}\quad 平均速度等于位移除以时间长度；连续速度的积分给出位移，所以平均速度为 \(\frac{1}{3}\int_0^3v(t)\,\mathrm{d}t\)。

\textbf{推导}\quad 先分别积分：\(\int_0^3t\,\mathrm{d}t=\frac{9}{2}\)，且
\[
\int_0^3\sin\pi t\,\mathrm{d}t
=\left[-\frac{\cos\pi t}{\pi}\right]_0^3
=\frac{-\cos3\pi+\cos0}{\pi}=\frac{2}{\pi}.
\]
因此
\[
\frac52=\frac13\int_0^3(t+k\sin\pi t)\,\mathrm{d}t
=\frac13\left(\frac92+\frac{2k}{\pi}\right).
\]
移项得 \(\frac{2k}{\pi}=3\)，故 \(k=\boxed{\frac{3\pi}{2}}\)。
\textbf{结论与检查}\quad \(\boxed{k=\frac{3\pi}{2}}\)。代回后位移为 \(\frac{9}{2}+3=\frac{15}{2}\)，除以 \(3\) 正好得到平均速度 \(\frac{5}{2}\)。
\end{QuestionSolution}

% Q 2026 15
\QuestionPage{3.1 定积分的概念、存在条件与性质}{填空题}{2026 年第 15 题}{2}{函数的积分平均值}{137 / 234}{**}
函数
\[f(x)=\ln(2+x)
\]
在区间 \([0,2]\) 上的平均值为 \blank{}.
\QuestionEnd

\begin{QuestionSolution}{2026 年第 15 题}
\textbf{已知与目标}\quad 求连续函数 \(\ln(2+x)\) 在 \([0,2]\) 上的平均值。

\textbf{方法与条件}\quad 积分平均值等于积分除以区间长度 \(2-0=2\)。换元 \(u=2+x\) 后，\(\mathrm{d}u=\mathrm{d}x\)，上下限由 \(0,2\) 变为 \(2,4\)。

\textbf{推导}\quad 因 \(\int\ln u\,\mathrm{d}u=u\ln u-u+C\)，所求值为
\[
\frac12\int_0^2\ln(2+x)\,\mathrm{d}x
=\frac12\bigl[u\ln u-u\bigr]_2^4
=\boxed{3\ln2-1}.
\]
具体代入为 \(\frac{1}{2}[(4\ln4-4)-(2\ln2-2)]=3\ln2-1\)。

\textbf{结论与检查}\quad 平均值为 \(\boxed{3\ln2-1}\)。函数递增，平均值应在 \(\ln2\) 与 \(\ln4\) 之间；所得值约 \(1.0794\)，符合范围。
\end{QuestionSolution}

% Q 2012 4
\QuestionPage{3.1 定积分的概念、存在条件与性质}{选择题}{2012 年第 4 题}{1}{定积分比较}{130 / 234}{***}
设$I_k = \int_0^{k\pi} \mathrm{e}^{x^{2}} \sin x\,\mathrm{d} x$（$k = 1,2,3$），则有\pickout{}
\begin{abcd}
\item $I_1 < I_2 < I_3$．
\item $I_3 < I_2 < I_1$．
\item $I_2 < I_3 < I_1$．
\item $I_2 < I_1 < I_3$．
\end{abcd}
\QuestionEnd

\begin{QuestionSolution}{2012 年第 4 题}
\textbf{已知与目标}\quad 三个积分的被积函数相同，上限依次为 \(\pi,2\pi,3\pi\)。要比较它们，需同时考虑正弦变号和指数因子增大。

\textbf{方法与条件}\quad 将每个长度为 \(\pi\) 的区间平移到 \([0,\pi]\)。连续函数的积分具有保序性：被积函数在区间内严格较大，其积分也严格较大。

\textbf{推导}\quad 令 \(J_j=\int_0^\pi \mathrm{e}^{(t+j\pi)^2}\sin t\,\mathrm{d}t\)，其中 \(j=0,1,2\)。在 \(0<t<\pi\) 上，\(\sin t>0\)，且 \(t<t+\pi<t+2\pi\)，平方和取指数均保持严格大小关系，因此 \(0<J_0<J_1<J_2\)。分别在后两段令 \(x=t+\pi\)、\(x=t+2\pi\)，由 \(\sin(t+j\pi)=(-1)^j\sin t\)，得
\[
I_1=J_0,\qquad I_2=J_0-J_1,\qquad I_3=J_0-J_1+J_2.
\]
所以 \(I_1-I_2=J_1>0\)，\(I_3-I_1=J_2-J_1>0\)。

\textbf{结论与检查}\quad \(\boxed{I_2<I_1<I_3}\)，选 D。还可检查 \(I_2<0<I_1\)，与得到的次序一致。
\end{QuestionSolution}

% Q 2018 5
\QuestionPage{3.1 定积分的概念、存在条件与性质}{选择题}{2018 年第 5 题}{1}{定积分估计与比较}{131 / 234}{***}
设$M = \int_{-\frac{\pi}{2}}^{\frac{\pi}{2}} \frac{\left(1 + x \right)^2}{1 + x^2}\,\mathrm{d} x$,
  $N = \int_{-\frac{\pi}{2}}^{\frac{\pi}{2}} \frac{1 + x}{\mathrm{e}^x}\,\mathrm{d} x$,
  $K = \int_{-\frac{\pi}{2}}^{\frac{\pi}{2}} \left(1 + \sqrt{\cos x}\right)\,\mathrm{d} x$，
则\pickout{}
\begin{abcd}
\item $M > N > K$．
\item $M > K > N$．
\item $K > M > N$．
\item $K > N > M$．
\end{abcd}
\QuestionEnd

\begin{QuestionSolution}{2018 年第 5 题}
\textbf{已知与目标}\quad 三个积分的区间都是 \([-\frac{\pi}{2},\frac{\pi}{2}]\)，长度为 \(\pi\)。目标是比较积分值，无需求出全部原函数。

\textbf{方法与条件}\quad 奇函数满足 \(h(-x)=-h(x)\)，在对称区间上的积分为零；再用积分的保序性，以常数函数 \(1\) 的积分 \(\pi\) 为基准。

\textbf{推导}\quad 展开分子得到
\[
M=\int_{-\frac{\pi}{2}}^{\frac{\pi}{2}}\left(1+\frac{2x}{1+x^2}\right)\mathrm{d}x=\pi,
\]
因为第二项为奇函数。在区间内部 \(\cos x>0\)，故 \(\sqrt{\cos x}>0\)，于是 \(K>\pi\)。另一方面，令 \(h(x)=\mathrm{e}^x-1-x\)，则 \(h'(x)=\mathrm{e}^x-1\)，在零点左侧为负、右侧为正，所以 \(h\) 在零点取唯一最小值 \(0\)。因 \(\mathrm{e}^x>0\)，可除得 \((1+x)\mathrm{e}^{-x}\le1\)，并且除零点外严格小于 \(1\)。积分便有 \(N<\pi\)。

\textbf{结论与检查}\quad \(\boxed{K>M>N}\)，选 C。指数不等式对负数 \(x\) 仍成立，所以整个积分区间均已覆盖。
\end{QuestionSolution}

% Q 2022 7
\QuestionPage{3.1 定积分的概念、存在条件与性质}{选择题}{2022 年第 7 题}{2}{比较定积分大小}{133 / 234}{***}
若
\(I_{1} = \int_{0}^{1} \frac{x}{2\left( 1 + \cos x \right)} \,\mathrm{d}x,\quad I_{2} = \int_{0}^{1} \frac{\ln\left( 1 + x \right)}{1 + \cos x} \,\mathrm{d}x,\quad I_{3} = \int_{0}^{1} \frac{2x}{1 + \sin x} \,\mathrm{d}x\)
则 \(\left( \ \right)\)
\choices{\(I_{1} < I_{2} < I_{3}\)}{\(I_{2} < I_{1} < I_{3}\)}{\(I_{1} < I_{3} < I_{2}\)}{\(I_{3} < I_{2} < I_{1}\)}
\QuestionEnd

\begin{QuestionSolution}{2022 年第 7 题}
\textbf{已知与目标}\quad 三个积分区间相同，目标是确定 \(I_1,I_2,I_3\) 的严格大小关系。

\textbf{方法与条件}\quad 先比较被积函数，再积分。比较分式时要确认分母为正；对数可写成 \(\ln(1+x)=\int_0^x\frac{1}{1+t}\,\mathrm{d}t\) 来估计。

\textbf{推导}\quad 对 \(0<x<1\) 及 \(0<t<x\)，有 \(\frac{1}{2}<\frac{1}{1+t}<1\)，两边从 \(0\) 积分到 \(x\)，得 \(\frac{x}{2}<\ln(1+x)<x\)。除以正数 \(1+\cos x\)，可得第一项被积函数小于第二项，故 \(I_1<I_2\)。又因为 \(0<x<1<\frac{\pi}{2}\)，有 \(1+\cos x>1\)、\(1+\sin x<2\)，所以
\[
\frac{\ln(1+x)}{1+\cos x}<x<\frac{2x}{1+\sin x}.
\]
对这一严格不等式积分，便得 \(I_2<I_3\)。

\textbf{结论与检查}\quad \(\boxed{I_1<I_2<I_3}\)，选 A。端点 \(x=0\) 的等号不影响内部严格不等式导出的严格积分关系。
\end{QuestionSolution}

% Q 2025 13
\QuestionPage{3.1 定积分的概念、存在条件与性质}{填空题}{2025 年第 13 题}{2}{和式极限与定积分定义}{136 / 234}{***}
\(\lim_{n \to \infty} \frac{1}{n^{2}} \left[ \ln \frac{1}{n} + 2 \ln \frac{2}{n} + \cdots + \left( n - 1 \right) \ln \frac{n - 1}{n} \right] =\) \blank{}.
\QuestionEnd

\begin{QuestionSolution}{2025 年第 13 题}
\textbf{已知与目标}\quad 求 \(\frac{1}{n^2}\sum_{k=1}^{n-1}k\ln\frac{k}{n}\) 的极限。难点在于 \(\ln x\) 在零点无定义，而真正出现的函数是 \(x\ln x\)。

\textbf{方法与条件}\quad 先作连续延拓，即在缺失的端点补上极限值，再应用 Riemann 和的定义。因为令 \(x=\mathrm{e}^{-s}\) 后 \(x\ln x=-s\mathrm{e}^{-s}\to0\)，可令 \(h(0)=0\)、\(h(x)=x\ln x\)（\(x>0\)），使 \(h\) 在 \([0,1]\) 连续。

\textbf{推导}\quad 加上 \(k=n\) 项不会改变和式，因为 \(\ln1=0\)。于是原式为 \(\sum_{k=1}^n h(\frac{k}{n})\frac{1}{n}\)，其中宽度为 \(\frac{1}{n}\)，取样点为右端点。其极限为 \(\int_0^1x\ln x\,\mathrm{d}x\)。在 \([\varepsilon,1]\) 上分部积分，取 \(u=\ln x\)、\(\mathrm{d}v=x\,\mathrm{d}x\)，再令 \(\varepsilon\to0^+\)，得
\[
\int_0^1 x\ln x\,\mathrm{d}x
=\left[\frac{x^2}{2}\ln x-\frac{x^2}{4}\right]_{0^+}^1
=\boxed{-\frac14}.
\]
这里 \(x^2\ln x\to0\)，所以零点边界项确实消失。

\textbf{结论与检查}\quad 极限为 \(\boxed{-\frac{1}{4}}\)。原和式中 \(\frac{k}{n}<1\)，所有非零项均为负，结果的符号正确。
\end{QuestionSolution}

% Q 2010 16
\QuestionPage{3.1 定积分的概念、存在条件与性质}{解答题}{2010 年第 16 题}{2}{积分比较与数列极限}{138 / 234}{***}
\begin{enumerate}
\item
比较$\int_0^1 |\ln t|[\ln(1+t)]^n\,\mathrm{d} t$
与$\int_0^1 t^n |\ln t|\,\mathrm{d} t$（$n=1,2,\cdots$）的大小，说明理由；
\item
记$u_n = \int_0^1 |\ln t|[\ln(1+t)]^n\,\mathrm{d} t$（$n=1,2,\cdots$），求极限$\lim_{n\to\infty} u_n$．
\end{enumerate}
\QuestionEnd

\begin{QuestionSolution}{2010 年第 16 题}
\textbf{已知与目标}\quad 对每个正整数 \(n\)，先比较两个含对数的积分，再求较小积分 \(u_n\) 在 \(n\to\infty\) 时的极限。

\textbf{方法与条件}\quad 对正数取正整数次幂保持大小关系，乘非负函数也保持大小关系。夹逼定理指一个数列被两个同趋某值的数列夹住时，它也趋于该值。先验证积分存在，再用它估计 \(u_n\)。

\textbf{推导}\quad （1）对 \(0<t\le1\)，\(\ln(1+t)=\int_0^t\frac{1}{1+s}\,\mathrm{d}s<t\)。在 \(0<t<1\) 上 \(|\ln t|=-\ln t>0\)，所以
\[
0<|\ln t|[\ln(1+t)]^n<t^n|\ln t|.
\]
因为 \(-t^n\ln t\to0\)（\(t\to0^+\)），两个被积函数都可在零点补值 \(0\) 后连续，故积分存在。连续函数在整个内部严格不等，积分也严格不等。

（2）分部积分取 \(u=-\ln t\)、\(\mathrm{d}v=t^n\,\mathrm{d}t\)，得到
\[
\int_0^1 t^n|\ln t|\,\mathrm{d}t
=\left[-\frac{t^{n+1}\ln t}{n+1}\right]_{0^+}^{1}
+\frac{1}{n+1}\int_0^1t^n\,\mathrm{d}t
=\frac{1}{(n+1)^2}.
\]
边界项在 \(1\) 处为零，在 \(0^+\) 处也因 \(t^{n+1}\ln t\to0\) 为零。于是 \(0<u_n<\frac{1}{(n+1)^2}\to0\)。

\textbf{结论与检查}\quad 第一问的前一个积分严格小于后一个积分，第二问 \(\boxed{\lim_{n\to\infty}u_n=0}\)。估计对每个 \(n\ge1\) 成立，无需交换极限与积分。
\end{QuestionSolution}

% Q 2020 20
\QuestionPage{3.1 定积分的概念、存在条件与性质}{解答题}{2020 年第 20 题}{3}{积分中值定理}{140 / 234}{***}
设函数 \(f\left( x \right)=\int_{1}^{x}{\mathrm{e}^{t^{2}} \,\mathrm{d} t}\)。
\begin{enumerate}
\item 证明：存在 \(\xi \in \left( 1,2 \right)\)，使得 \(f\left( \xi \right)=\left( 2-\xi \right)\mathrm{e}^{\xi^{2}}\)；
\item 证明：存在 \(\eta \in \left( 1,2 \right)\)，使得 \(f\left( 2 \right)=\ln 2 \cdot \eta \mathrm{e}^{\eta^{2}}\)。
\end{enumerate}
\QuestionEnd

\begin{QuestionSolution}{2020 年第 20 题}
\textbf{已知与目标}\quad \(f(x)=\int_1^x\mathrm{e}^{t^2}\,\mathrm{d}t\)，要证明两个等式分别在某个 \((1,2)\) 内的点成立，两问中的点可以不同。

\textbf{方法与条件}\quad 第一问用零点存在定理：连续函数在两端异号，内部便有零点。第二问用加权积分中值定理：连续函数乘不变号权重的积分，等于某点函数值乘权重积分；还需说明该点能取在开区间内。

\textbf{推导}\quad （1）令 \(H(x)=f(x)-(2-x)\mathrm{e}^{x^2}\)，它在 \([1,2]\) 连续。代入端点，\(f(1)=0\)，所以 \(H(1)=-\mathrm{e}<0\)；又 \(H(2)=f(2)=\int_1^2\mathrm{e}^{t^2}\,\mathrm{d}t>0\)。零点存在定理给出 \(\xi\in(1,2)\)，满足 \(H(\xi)=0\)，即第一问。

（2）将 \(f(2)\) 写为
\[
\int_1^2 (t \mathrm{e}^{t^2})\frac{1}{t}\,\mathrm{d}t.
\]
令 \(p(t)=t\mathrm{e}^{t^2}\)，有 \(p'(t)=(1+2t^2)\mathrm{e}^{t^2}>0\)。在内部 \(p(1)<p(t)<p(2)\)，乘正权重 \(\frac{1}{t}\) 并积分，得
\[
p(1)<\frac{f(2)}{\int_1^2\frac{1}{t}\,\mathrm{d}t}<p(2).
\]
由 \(p\) 连续，存在 \(\eta\in(1,2)\) 使中间值等于 \(p(\eta)\)。再用 \(\int_1^2\frac{1}{t}\,\mathrm{d}t=\ln2\)，便得 \(f(2)=\ln2\cdot\eta\mathrm{e}^{\eta^2}\)。

\textbf{结论与检查}\quad 两问要求的内部点均存在。第一问的严格异号、第二问的严格端点界分别排除了点只能位于端点的可能。
\end{QuestionSolution}

% Q 2014 19
\QuestionPage{3.1 定积分的概念、存在条件与性质}{解答题}{2014 年第 19 题}{3}{定积分保序与比较}{139 / 234}{****}
设函数$f(x),g(x)$在区间$[a,b]$上连续，且$f(x)$单调增加，$0 \le g(x) \le 1$，证明：
\begin{enumerate}
\item $0 \le \int_a^x g(t)\,\mathrm{d} t \le x - a$，$x \in [a,b]$；
\item $\int_a^{a + \int_a^b g(t)\,\mathrm{d} t} f(x)\,\mathrm{d} x \le \int_a^b f(x)g(x)\,\mathrm{d} x$．
\end{enumerate}
\QuestionEnd

\begin{QuestionSolution}{2014 年第 19 题}
\textbf{已知与目标}\quad \(f,g\) 连续，\(f\) 单调增加，\(0\le g\le1\)。先控制 \(g\) 的积分范围，再证明截取左端一段所得的 \(f\) 积分不超过加权积分。

\textbf{方法与条件}\quad 设 \(c=a+\int_a^b g(t)\,\mathrm{d}t\)，比较右边减左边。权重指乘在函数旁边的因子，这里 \(g\) 与 \(1-g\) 都非负；用单调性与分界值 \(f(c)\) 比较时可保持不等号方向。

\textbf{推导}\quad （1）对 \(x\in[a,b]\)，将 \(0\le g(t)\le1\) 从 \(a\) 积分到 \(x\)，得
\[
0\le\int_a^xg(t)\,\mathrm{d}t\le\int_a^x1\,\mathrm{d}t=x-a.
\]
（2）取 \(x=b\)，可知 \(a\le c\le b\)，故分割积分合法。由 \(c-a=\int_a^b g\) 得
\[
\int_a^c(1-g(x))\,\mathrm{d}x
=(c-a)-\int_a^cg(x)\,\mathrm{d}x=\int_c^bg(x)\,\mathrm{d}x.
\]
这说明减去或加上常数 \(f(c)\) 所产生的两项正好抵消。因此所证不等式右边减左边为
\[
\begin{aligned}
\int_c^b f g-\int_a^c f(1-g)
={}&\int_c^b [f(x)-f(c)]g(x)\,\mathrm{d}x\\
&+\int_a^c [f(c)-f(x)](1-g(x))\,\mathrm{d}x\ge0 .
\end{aligned}
\]
当 \(x\in[c,b]\) 时 \(f(x)-f(c)\ge0\)；当 \(x\in[a,c]\) 时 \(f(c)-f(x)\ge0\)，故上式两项均非负。

\textbf{结论与检查}\quad 两个不等式均成立。证明没有假定 \(f\ge0\)，所以 \(f\) 取负值时仍适用；若 \(c=a\) 或 \(c=b\)，相应空区间积分为零，推导仍成立。
\end{QuestionSolution}

\SectionHeading{3.2}{一元函数积分学及其应用}{微积分基本公式与基本定理}{定理与要点\quad 教材 PDF 第 203 页起\hfill 本节 11 题}
\pdfbookmark[2]{3.2 微积分基本公式与基本定理}{section-3.2}
\Theory{定理 2.1}{204}{Newton-Leibniz 公式}{若可积函数 \(f\) 在 \([a,b]\) 上有原函数 \(F\)，则 \(\int_a^b f(x)\,\mathrm{d}x=F(b)-F(a)\)。}
\Theory{定理 2.2}{206}{微积分第一基本定理}{连续函数的变上限积分可导，其导数等于被积函数在上限的值。}
\Theory{定理 2.3}{207}{微积分第二基本定理}{若 \(F\) 是 \(f\) 的一个原函数，则 \(f\) 的全部原函数可写成 \(F+C\)，其中 \(C\) 为任意常数。}
\Theory{推论 2.1}{206}{连续函数的原函数}{连续函数有原函数，固定下限的变上限积分就是其中一个。}
\Theory{公式}{206--207}{上下限均变动}{上限贡献取正，下限贡献取负，各自用被积函数值乘对应端点的导数。}
% Q 2013 3
\QuestionPage{3.2 微积分基本公式与基本定理}{选择题}{2013 年第 3 题}{1}{变上限积分的连续性与可导性}{142 / 234}{**}
设函数$f(x) = \begin{cases}
  \sin x, & 0 \le x < \pi , \\
       2, & \pi \le x \le 2\pi ,
\end{cases}$ $F(x) = \int_0^x f(t) \,\mathrm{d} t$，则\pickout{}
\begin{abcd}
\item $x = \pi$是函数$F(x)$的跳跃间断点．
\item $x = \pi$是函数$F(x)$的可去间断点．
\item $F(x)$在$x = \pi$处连续但不可导．
\item $F(x)$在$x = \pi$处可导．
\end{abcd}
\QuestionEnd

\begin{QuestionSolution}{2013 年第 3 题}
\textbf{已知与目标}\quad \(f\) 在 \(\pi\) 处由正弦段接到常数段，要判断其积分函数 \(F\) 在接点的连续性与可导性。

\textbf{方法与条件}\quad 连续要求两侧极限等于函数值；可导要求左右差商极限相同。由于 \(f\) 在接点不连续，不能直接套用 \(F'(\pi)=f(\pi)\)，应先把积分算成分段式。

\textbf{推导}\quad 当 \(0\le x<\pi\)，\(F(x)=\int_0^x\sin t\,\mathrm{d}t=[-\cos t]_0^x=1-\cos x\)。当 \(\pi\le x\le2\pi\)，
\[
F(x)=\int_0^\pi\sin t\,\mathrm{d}t+\int_\pi^x2\,\mathrm{d}t=2+2(x-\pi).
\]
代入接点，左极限为 \(1-\cos\pi=2\)，右极限与 \(F(\pi)\) 也为 \(2\)，所以连续。左侧表达式的导数在接点趋于 \(\sin\pi=0\)，右侧线性表达式的斜率为 \(2\)，左右导数不同。

\textbf{结论与检查}\quad \(F\) 在 \(\pi\) 处连续但不可导，选 \(\boxed{\mathrm{C}}\)。积分函数连续并不要求被积函数在该点连续，这与本题接点的行为一致。
\end{QuestionSolution}

% Q 2016 2
\QuestionPage{3.2 微积分基本公式与基本定理}{选择题}{2016 年第 2 题}{1}{分段函数的原函数}{143 / 234}{**}
已知函数$f(x) = \begin{cases}
  2(x - 1), & x < 1, \\
     \ln x, & x \ge 1,
\end{cases}$，则$f(x)$的一个原函数是\pickout{}
\begin{abcd}
\item $F(x) = \begin{cases}
                 (x - 1)^2, & x < 1 \\
  x\left(\ln x - 1 \right), & x \ge 1
\end{cases}$．
\item $F(x) = \begin{cases}
                     (x - 1)^2, & x < 1 \\
  x\left(\ln x + 1 \right) - 1, & x \ge 1
\end{cases}$．
\item $F(x) = \begin{cases}
                     (x - 1)^2, & x < 1 \\
  x\left(\ln x + 1 \right) + 1, & x \ge 1
\end{cases}$．
\item $F(x) = \begin{cases}
                     (x - 1)^2, & x < 1 \\
  x\left(\ln x - 1 \right) + 1, & x \ge 1
\end{cases}$．
\end{abcd}
\QuestionEnd

\begin{QuestionSolution}{2016 年第 2 题}
\textbf{已知与目标}\quad \(f\) 在 \(1\) 的两侧表达式不同，要找在整个实轴上导数都等于 \(f\) 的函数 \(F\)，这种 \(F\) 称为 \(f\) 的原函数。

\textbf{方法与条件}\quad 先检查各开区间的导数，再检查接点。原函数必须可导，因此也必须连续，分段积分常数要使两段接上。

\textbf{推导}\quad 左段 \([(x-1)^2]'=2(x-1)\)。右段中
\[
[x(\ln x-1)]'=\ln x,\qquad [x(\ln x+1)]'=\ln x+2,
\]
所以 B、C 的右段导数不符合。A、D 的右段分别在 \(1\) 处取 \(-1\)、\(0\)，而左极限为 \(0\)，故只有 D 连续。对 D，其左右导数分别为 \(2(1-1)=0\) 与 \(\ln1=0\)，相等且等于 \(f(1)\)。

\textbf{结论与检查}\quad 选 \(\boxed{\mathrm{D}}\)。两侧区间和接点均已核对，故该函数确实在全定义域上满足 \(F'=f\)。
\end{QuestionSolution}

% Q 2015 11
\QuestionPage{3.2 微积分基本公式与基本定理}{填空题}{2015 年第 11 题}{2}{变上限积分求导}{145 / 234}{**}
设$f(x)$连续，$\varphi (x) = \int_0^{x^2} x f(t)\,\mathrm{d} t$，
若$\varphi (1) = 1$，$\varphi'(1) = 5$，则$f(1) =$ \blank{}．
\QuestionEnd

\begin{QuestionSolution}{2015 年第 11 题}
\textbf{已知与目标}\quad \(f\) 连续，已知 \(\varphi(1)=1\)、\(\varphi'(1)=5\)，求 \(f(1)\)。注意积分变量是 \(t\)，式中的 \(x\) 对积分而言是常数。

\textbf{方法与条件}\quad 把 \(x\) 提到积分外，先用乘积求导，再用变上限积分求导。连续性保证 \(\frac{\mathrm{d}}{\mathrm{d}x}\int_0^{x^2}f(t)\,\mathrm{d}t=f(x^2)\cdot2x\)。

\textbf{推导}\quad 写成 \(\varphi(x)=x\int_0^{x^2}f(t)\,\mathrm{d}t\)，则
\[
\varphi'(x)=\int_0^{x^2}f(t)\,\mathrm{d}t+2x^2 f(x^2).
\]
由 \(\varphi(1)=1\cdot\int_0^1f(t)\,\mathrm{d}t=1\)，把 \(x=1\) 代入导数式，得到 \(5=1+2f(1)\)，解得 \(f(1)=2\)。

\textbf{结论与检查}\quad \(\boxed{f(1)=2}\)。代回导数条件为 \(1+2\times2=5\)，与已知一致，外部因子和上限的两项贡献都已计入。
\end{QuestionSolution}

% Q 2023 15
\QuestionPage{3.2 微积分基本公式与基本定理}{填空题}{2023 年第 15 题}{2}{积分关系求函数值}{147 / 234}{**}
设连续函数 \(f(x)\) 满足
\[f(x+2)-f(x)=x,\qquad \int_0^2 f(x)\,\mathrm{d}x=0,
\]
则
\[\int_1^3 f(x)\,\mathrm{d}x=\blank{}.
\]
\QuestionEnd

\begin{QuestionSolution}{2023 年第 15 题}
\textbf{已知与目标}\quad 已知跨越 \(2\) 个单位的函数差 \(f(x+2)-f(x)=x\)，以及 \([0,2]\) 上积分为零，要求 \([1,3]\) 上积分。

\textbf{方法与条件}\quad 利用积分区间的可加性，比较两区间新增的 \([2,3]\) 与移出的 \([0,1]\)。两段平移 \(2\) 后可以用已知函数差相减，无需求出 \(f\)。

\textbf{推导}\quad 先写 \(\int_1^3f=\int_0^2f-\int_0^1f+\int_2^3f\)，第一项为零。在最后一项令原变量为 \(x+2\)，得
\[
\begin{aligned}
\int_1^3f(x)\,\mathrm{d}x
&=\int_2^3f(x)\,\mathrm{d}x-\int_0^1f(x)\,\mathrm{d}x\\
&=\int_0^1[f(x+2)-f(x)]\,\mathrm{d}x
=\int_0^1 x\,\mathrm{d}x=\boxed{\frac12}.
\end{aligned}
\]
\textbf{结论与检查}\quad \(\boxed{\int_1^3f(x)\,\mathrm{d}x=\frac{1}{2}}\)。新增段减移出段使用的是 \(f(x+2)-f(x)\)，与题设符号相同。
\end{QuestionSolution}

% Q 2009 6
\QuestionPage{3.2 微积分基本公式与基本定理}{选择题}{2009 年第 6 题}{1、2}{变上限积分的图像}{141 / 234}{***}
设函数$y = f(x)$在区间$\left[- 1,3 \right]$上的图形为
\par\centerline{\begin{tikzpicture}[scale=\DiagramScale,x=1pt,y=1pt,xscale=0.439724,yscale=-0.439724]
\path[use as bounding box] (0,0) rectangle (440,367);
\path[fill=black!100,even odd rule,shift={(133,1)}] svg {M0 0 C0.7 0 1.3 0 2 0 C3.6 3.8 5.1 7.6 6.7 11.4 C7.2 12.7 7.8 14 8.3 15.3 C9.1 17.2 9.8 19 10.6 20.9 C10.8 21.5 10.8 21.5 12 24.3 C13 27 13 27 13 29 C9.9 27.8 6.9 26.5 4 25 C3.2 27.9 2.9 30.4 2.9 33.4 C2.9 34.3 2.9 35.1 2.9 36 C2.9 36.9 2.9 37.8 2.9 38.8 C2.9 39.7 2.9 40.7 2.9 41.7 C2.9 44.7 2.9 47.7 2.9 50.8 C2.9 52.8 2.9 54.9 3 56.9 C3 61.9 3 67 3 72 C5 72 7 72 9 72 C9.3 73 9.7 74 10 75 C7.7 75 5.4 75 3 75 C3 98.1 3 121.2 3 145 C4.7 145.7 6.3 146.3 8 147 C8.7 147.3 9.3 147.7 10 148 C9.7 149 9.3 150 9 151 C8.3 151 7.7 151 7 151 C6.3 152.3 5.7 153.6 5 155 C4.3 155 3.7 155 3 155 C3 177.4 3 199.9 3 223 C26.8 223 50.5 223 75 223 C75 221.3 75 219.7 75 218 C75.7 216.1 76.4 214.2 77.1 212.3 C78.8 207.7 80 202.9 81.2 198.1 C108.7 94.3 108.7 94.3 144.8 73.3 C147.2 72.1 149.6 71 152 70 C152.7 72 153.3 74 154 76 C153.5 76.2 153.5 76.2 151 77 C151 76.3 151 75.7 151 75 C114.9 94.2 101.1 144.6 90 181 C86.7 192.3 83.8 203.6 81 215 C80.7 216.3 80.7 216.3 79 223 C101.1 223 123.2 223 146 223 C146.3 221.7 146.7 220.4 147 219 C147.5 218.7 147.5 218.7 150 217 C150.3 216 150.7 215 151 214 C152 214 153 214 154 214 C154 215.3 154 216.6 154 218 C157.6 220.2 160.5 220.3 164.7 220.3 C165.4 220.3 166 220.3 166.7 220.3 C168.9 220.4 171.1 220.4 173.3 220.4 C174.8 220.4 176.3 220.5 177.8 220.5 C181.8 220.5 185.8 220.6 189.8 220.6 C193.9 220.6 197.9 220.7 202 220.7 C210 220.8 218 220.9 226 221 C226 220 226 219 226 218 C227 218 228 218 229 218 C229.3 219.3 229.7 220.6 230 222 C234.5 222.2 234.5 222.2 257 223 C255.4 219.7 253.7 216.4 252 213 C256.2 213.6 259.7 214.6 263.6 216.3 C264.6 216.7 265.7 217.2 266.7 217.6 C267.7 218.1 268.8 218.6 269.9 219.1 C270.9 219.5 272 220 273.1 220.5 C275.7 221.7 278.4 222.8 281 224 C281 224.7 281 225.3 281 226 C277.2 227.6 273.4 229.2 269.6 230.8 C268.3 231.3 267 231.8 265.7 232.4 C263.8 233.2 262 233.9 260.1 234.7 C259 235.2 257.9 235.6 256.7 236.1 C254 237 254 237 252 236 C253.6 233 255.3 230.1 257 227 C244.2 226.9 231.3 226.9 218.5 226.8 C212.5 226.8 206.6 226.8 200.6 226.8 C194.8 226.7 189.1 226.7 183.3 226.7 C181.1 226.7 179 226.7 176.8 226.7 C173.7 226.7 170.6 226.7 167.5 226.7 C166.6 226.6 165.7 226.6 164.8 226.6 C160.1 226.6 157.5 226.7 154 230 C151 229.8 148.8 229.5 146.2 228 C143.2 226.7 140.8 226.8 137.5 226.8 C137.2 226.8 137.2 226.8 135.5 226.8 C133.4 226.8 131.3 226.8 129.2 226.8 C127.7 226.8 126.2 226.8 124.7 226.8 C120.9 226.8 117 226.8 113.1 226.9 C109.2 226.9 105.2 226.9 101.3 226.9 C93.5 226.9 85.8 227 78 227 C77.9 227.6 77.7 228.2 77.6 228.8 C77.4 229.6 77.3 230.4 77.1 231.3 C76.9 232 76.8 232.8 76.6 233.6 C75.7 237.3 74.1 240.5 72.4 243.9 C72.1 244.6 71.7 245.3 71.4 246 C59.2 269 35.7 282.7 13.2 294.1 C12.4 294.5 11.6 294.9 10.8 295.3 C5.3 298 5.3 298 3 298 C3 320.4 3 342.9 3 366 C1.7 365.7 0.4 365.3 -1 365 C-1 319.5 -1 273.9 -1 227 C-44.2 227 -87.5 227 -132 227 C-132 225.7 -132 224.4 -132 223 C-113.5 223 -95 223 -76 223 C-76 220 -76 217.1 -76 214 C-75 214 -74 214 -73 214 C-73 217 -73 219.9 -73 223 C-49.2 223 -25.5 223 -1 223 C-1 200.6 -1 178.1 -1 155 C-2 154.7 -3 154.3 -4 154 C-4.3 153.3 -4.7 152.7 -5 152 C-28.4 152 -51.9 152 -76 152 C-75.7 150.7 -75.3 149.4 -75 148 C-67.3 146.9 -59.9 146.9 -52.1 147 C-50.1 147 -48.1 147 -46.1 147 C-41.2 147.1 -36.3 147.1 -31.4 147.2 C-28.7 147.2 -25.9 147.2 -23.2 147.2 C-20.6 147.3 -18 147.3 -15.4 147.3 C-14.2 147.3 -13 147.3 -11.8 147.3 C-11.3 147.3 -11.3 147.3 -8.7 147.4 C-5.5 146.9 -4.3 146.2 -2 144 C-1.5 141.8 -1.5 141.8 -1.5 139.1 C-1.5 138.1 -1.5 137.1 -1.4 136.1 C-1.4 135.5 -1.4 135.5 -1.5 132.7 C-1.4 132.2 -1.4 132.2 -1.4 129.2 C-1.4 125.4 -1.4 121.5 -1.4 117.7 C-1.4 115 -1.4 112.3 -1.3 109.7 C-1.3 103.4 -1.3 97.1 -1.3 90.8 C-1.3 83.6 -1.2 76.4 -1.2 69.3 C-1.1 54.5 -1 39.8 -1 25 C-1.3 25.1 -1.3 25.1 -2.7 25.9 C-3.4 26.2 -4.1 26.6 -4.9 27 C-5.6 27.4 -6.3 27.7 -7.1 28.1 C-9 29 -9 29 -11 29 C-10.4 25.2 -9.2 21.8 -7.7 18.3 C-7.3 17.2 -6.8 16.1 -6.4 15 C-5.9 13.9 -5.4 12.8 -4.9 11.7 C-4.5 10.6 -4 9.4 -3.5 8.3 C-2.3 5.5 -1.2 2.8 0 0 Z M0 147 C-0.7 148.3 -1.3 149.6 -2 151 C-1.5 151.3 -1.5 151.3 1 153 C2 152.3 3 151.7 4 151 C3.3 149.7 2.7 148.4 2 147 C1.3 147 0.7 147 0 147 Z M150 221 C149.7 222 149.3 223 149 224 C149.3 224.7 149.7 225.3 150 226 C151.6 225.7 153.3 225.3 155 225 C154.7 223.7 154.3 222.4 154 221 C152.7 221 151.4 221 150 221 Z M3 227 C2.4 240.7 1.8 254.3 1.9 268 C1.9 270.7 1.9 273.4 1.9 276.1 C1.9 277.9 1.9 279.7 1.9 281.4 C1.9 282.2 1.9 283 1.9 283.9 C1.9 287.4 2 290.6 3 294 C9 292.2 14.5 290 20.1 287 C20.8 286.6 21.5 286.2 22.3 285.8 C34.2 279.3 45.4 271.6 55 262 C55.8 261.2 56.6 260.4 57.4 259.6 C65.6 250.9 74 239.5 74 227 C50.6 227 27.1 227 3 227 Z };
\path[fill=black!100,even odd rule,shift={(284.875,235.625)}] svg {M0 0 C2.9 0.1 4.1 0.3 6.7 1.9 C8.4 4.9 8.4 7 7.7 10.3 C6.8 13.6 5.3 15.7 3.1 18.4 C2.5 18.4 1.8 18.4 1.1 18.4 C0.8 19 0.5 19.7 0.2 20.3 C-0.9 22.4 -0.9 22.4 -1.9 23.4 C1.8 23.4 5.4 23.4 9.1 23.4 C8.8 24.2 8.8 24.2 7.1 28.4 C2.2 28.4 -2.8 28.4 -7.9 28.4 C-6.1 24.9 -4.5 22.3 -1.9 19.4 C-1.2 19.4 -0.6 19.4 0.1 19.4 C0.4 18.6 0.7 17.8 1.1 16.9 C2.1 14.4 2.1 14.4 3.1 13.4 C3.4 6.8 3.4 6.8 2.1 4.4 C-4.6 5.2 -4.6 5.2 -6.9 6.4 C-5.9 1.7 -5.1 0.5 0 0 Z };
\path[fill=black!100,even odd rule,shift={(114.875,56.625)}] svg {M0 0 C3 0.1 4 0.3 6.6 2 C8.7 5.3 8.8 6.6 8.1 10.4 C5.5 15.3 2.3 19.7 -1.9 23.4 C1.8 23.4 5.4 23.4 9.1 23.4 C8.8 24.7 8.5 26 8.1 27.4 C2.8 28.1 -2.5 28.5 -7.9 28.4 C-6.1 24.4 -3.5 21.6 -0.5 18.4 C2.8 14.4 3.5 11.5 3.1 6.4 C2.8 5.7 2.5 5.1 2.1 4.4 C-3.5 4.1 -3.5 4.1 -6.9 6.4 C-5.9 1.7 -5.1 0.5 0 0 Z };
\path[fill=black!100,even odd rule,shift={(365,236)}] svg {M0 0 C1.4 2.2 1.4 2.2 2 5 C1.6 6 1.3 7 0.9 8.1 C0.6 9 0.3 10 0 11 C1 12.3 2 13.7 3 15 C2.9 19 2.6 21.2 0.1 24.4 C-3.7 27.3 -6.3 27.3 -11 27 C-11.7 26 -12.3 25 -13 24 C-9.3 22.8 -6.8 23.4 -3 24 C-1 21 -1 21 -1.1 18.6 C-2.3 15.2 -4.1 14 -7 12 C-6.7 11.7 -6.7 11.7 -5 10.3 C-3 8 -3 8 -2.8 5.2 C-2.8 4.9 -2.8 4.9 -3 3 C-5.8 2 -6.9 2 -9.8 2.9 C-10.1 3.1 -10.1 3.1 -12 4 C-9.6 -0.7 -9.6 -0.7 -7.1 -1.9 C-4.2 -2 -2.5 -1.5 0 0 Z };
\path[fill=black!100,even odd rule,shift={(171,0)}] svg {M0 0 C3 0 6 0 9 0 C9 3 9 3 7.6 4.8 C3.5 10.5 3.6 19.2 4.6 25.9 C5.6 30.4 7.1 33.5 10 37 C5.9 35.5 3.3 32.6 1 29 C-1.7 21.7 -1.6 14.3 1.6 7.1 C2 6.1 2.5 5 3 4 C3 3 3 2 3 1 C2 0.7 1 0.3 0 0 Z };
\path[fill=black!100,even odd rule,shift={(169,1)}] svg {M0 0 C0.7 0.3 1.3 0.7 2 1 C0.7 4 -0.6 6.9 -2 10 C-1.3 10.3 -0.7 10.7 0 11 C-0.3 11.2 -0.3 11.2 -2 12 C-2.8 14.5 -3.6 17 -4.3 19.6 C-6.1 25.7 -7.8 31 -12 36 C-15.3 37 -15.3 37 -18 37 C-18 35.7 -18 34.4 -18 33 C-17.3 33 -16.7 33 -16 33 C-16 33.7 -16 34.3 -16 35 C-15.5 34.8 -15.5 34.8 -13 34 C-9.9 26.7 -8.3 18.8 -7 11 C-7.7 10.7 -8.3 10.3 -9 10 C-8.3 9.3 -7.5 8.6 -6.8 7.8 C-4.3 5.3 -2.1 2.8 0 0 Z };
\path[fill=black!100,even odd rule,shift={(121,295)}] svg {M0 0 C0 8.6 0 17.2 0 26 C0.7 26 1.3 26 2 26 C2 26.7 2 27.3 2 28 C-0.3 28 -2.6 28 -5 28 C-5.3 27.3 -5.7 26.7 -6 26 C-5.3 26 -4.7 26 -4 26 C-4 18.4 -4 10.8 -4 3 C-5 2.7 -6 2.3 -7 2 C-4.4 0.5 -3.1 0 0 0 Z };
\path[fill=black!100,even odd rule,shift={(115,234)}] svg {M0 0 C7.4 -0.5 7.4 -0.5 10.4 1.9 C12.8 5.1 13.5 7.5 13.4 11.6 C12.5 17 9.3 22.1 5.1 25.6 C3.2 26.7 3.2 26.7 0 28 C-0.7 27.7 -1.3 27.3 -2 27 C-0.6 25.5 -0.6 25.5 1 24 C1.7 24 2.3 24 3 24 C3.3 23.3 3.5 22.6 3.8 21.9 C4.7 19.8 5.6 17.7 6.6 15.6 C8.1 11.7 8.5 9.1 8 5 C6.9 3.1 6.9 3.1 5 2 C1.8 1.8 1.8 1.8 -1 2 C-0.7 1.3 -0.3 0.7 0 0 Z };
\path[fill=black!100,even odd rule,shift={(114,235)}] svg {M0 0 C0 0.7 0 1.3 0 2 C-0.8 2.7 -1.5 3.4 -2.3 4.1 C-6.5 8.6 -7.8 12.8 -8.3 18.8 C-8 22 -8 22 -6 25 C-3.4 25.7 -3.4 25.7 -1 26 C-1 26.3 -1 26.7 -1 27 C-7.3 27.4 -7.3 27.4 -10.4 25.7 C-13.4 22.5 -13.5 18.7 -13.4 14.5 C-12.5 9.1 -9 4.7 -5 1.2 C-3 0 -3 0 0 0 Z };
\path[fill=black!100,even odd rule,shift={(203,1)}] svg {M0 0 C4.7 1.9 7.1 5.1 9.2 9.6 C10.9 17.1 10.9 23.6 7.2 30.3 C5.9 32.1 4.6 33.5 3 35 C2.3 34.7 1.7 34.3 1 34 C1.5 33.3 1.9 32.5 2.4 31.8 C6.3 25 6.1 15.9 4.4 8.4 C3.3 5.3 1.8 2.7 0 0 Z };
\path[fill=black!100,even odd rule,shift={(155,137)}] svg {M0 0 C0 8.9 0 17.8 0 27 C1 27.3 2 27.7 3 28 C-2.6 29.1 -2.6 29.1 -6 28 C-5.3 27.3 -4.7 26.7 -4 26 C-3.8 23.8 -3.8 23.8 -3.8 21 C-3.8 20 -3.8 19 -3.8 18 C-3.8 17 -3.9 15.9 -3.9 14.9 C-3.9 13.8 -3.9 12.8 -3.9 11.7 C-3.9 9.1 -4 6.6 -4 4 C-4.7 3.7 -5.3 3.3 -6 3 C-1.1 0 -1.1 0 0 0 Z };
\path[fill=black!100,even odd rule,shift={(69,234)}] svg {M0 0 C1.3 9 1.1 17.9 1 27 C1.7 27.3 2.3 27.7 3 28 C0 28 -2.9 28 -6 28 C-5.3 27.3 -4.7 26.7 -4 26 C-3.7 23.8 -3.7 23.8 -3.7 21.1 C-3.7 20.1 -3.7 19.1 -3.7 18.1 C-3.7 17.1 -3.7 16.1 -3.8 15 C-3.7 14 -3.7 13 -3.7 11.9 C-3.7 11.4 -3.7 11.4 -3.7 8.9 C-3.7 8 -3.7 7.1 -3.7 6.2 C-4 4 -4 4 -6 2 C-2.2 0 -2.2 0 0 0 Z };
\path[fill=black!100,even odd rule,shift={(217,238)}] svg {M0 0 C0 0.7 0 1.5 0.1 2.3 C0.2 5.6 0.3 9 0.4 12.4 C0.4 13.6 0.4 14.7 0.5 15.9 C0.5 17.1 0.5 18.2 0.6 19.3 C0.6 20.4 0.6 21.4 0.7 22.5 C1 25 1 25 3 27 C0 27 -2.9 27 -6 27 C-5.3 26.3 -4.7 25.7 -4 25 C-3.8 22.8 -3.8 22.8 -3.8 20 C-3.8 19 -3.8 18 -3.8 17 C-3.8 16 -3.9 14.9 -3.9 13.9 C-3.9 12.8 -3.9 11.8 -3.9 10.7 C-3.9 8.1 -4 5.6 -4 3 C-5 2.7 -6 2.3 -7 2 C-4 0 -4 0 0 0 Z };
\path[fill=black!100,even odd rule,shift={(428,217)}] svg {M0 0 C2.2 1.9 2.2 1.9 4 4 C4.8 3.3 5.6 2.7 6.4 2 C9 0 9 0 11 0 C11 1 11 2 11 3 C10.6 3.1 10.6 3.1 8.6 3.6 C6 5 6 5 5.1 7.9 C5 10.6 5.1 12.5 6 15 C7 15 8 15 9 15 C8.3 16.3 7.7 17.6 7 19 C5.7 18.7 4.4 18.3 3 18 C2.8 17.2 2.8 17.2 2 13 C1 13.3 0 13.7 -1 14 C-0.8 12.9 -0.6 11.9 -0.4 10.8 C-0 7.1 0.1 5.9 -2 3 C-2.7 2.3 -3.3 1.7 -4 1 C-2.2 0.2 -2.2 0.2 0 0 Z };
\path[fill=black!100,even odd rule,shift={(185,11)}] svg {M0 0 C2 0 4 0 6 0 C6.3 1.3 6.7 2.6 7 4 C7.7 4 8.3 4 9 4 C9.3 7.6 9.7 11.3 10 15 C11.3 14.3 12.6 13.7 14 13 C13.2 14.9 13.2 14.9 12 17 C9.9 17.8 9.9 17.8 8 18 C6 15 6 15 6 12 C5.3 12 4.7 12 4 12 C4 11.4 3.9 10.8 3.9 10.1 C3.8 9.3 3.8 8.5 3.7 7.7 C3.6 6.9 3.6 6.1 3.5 5.3 C2.9 2.6 2 1.7 0 0 Z };
\path[fill=black!100,even odd rule,shift={(38,246)}] svg {M0 0 C5.6 0 11.2 0 17 0 C17 0.7 17 1.3 17 2 C11.4 2 5.8 2 0 2 C0 1.3 0 0.7 0 0 Z };
\path[fill=black!100,even odd rule,shift={(89,307)}] svg {M0 0 C5.9 0 11.9 0 18 0 C18 0.3 18 0.7 18 1 C12.1 1 6.1 1 0 1 C0 0.7 0 0.3 0 0 Z };
\path[fill=black!100,even odd rule,shift={(188,23)}] svg {M0 0 C0 3 0 3 -1.2 4.7 C-3 6 -3 6 -7 6 C-6.7 5 -6.3 4 -6 3 C-4.1 2.4 -4.1 2.4 -2 2 C-1.3 1.3 -0.7 0.7 0 0 Z };
\path[fill=black!100,even odd rule,shift={(426,231)}] svg {M0 0 C0 1 0 2 0 3 C-0.7 3.7 -1.3 4.3 -2 5 C-4.1 4.6 -4.1 4.6 -6 4 C-6 3.3 -6 2.7 -6 2 C-5.6 1.9 -5.6 1.9 -3.6 1.6 C-1 1 -1 1 0 0 Z };
\path[fill=black!100,even odd rule,shift={(57,190)}] svg {M0 0 C1 0 2 0 3 0 C3 1.7 3 3.3 3 5 C2 5 1 5 0 5 C0 3.3 0 1.7 0 0 Z };
\path[fill=black!100,even odd rule,shift={(284,156)}] svg {M0 0 C1 0.3 2 0.7 3 1 C2.7 2.3 2.3 3.6 2 5 C1 4.7 0 4.3 -1 4 C-0.7 2.7 -0.3 1.4 0 0 Z };
\path[fill=black!100,even odd rule,shift={(196,11)}] svg {M0 0 C1.7 0 3.3 0 5 0 C4.7 1 4.3 2 4 3 C2 3.3 0 3.7 -2 4 C-1.3 2.7 -0.7 1.4 0 0 Z };
\path[fill=black!100,even odd rule,shift={(283,124)}] svg {M0 0 C1.3 0 2.6 0 4 0 C3.7 1.3 3.3 2.6 3 4 C2 3.7 1 3.3 0 3 C0 2 0 1 0 0 Z };
\path[fill=black!100,even odd rule,shift={(284,115)}] svg {M0 0 C1 0.3 2 0.7 3 1 C2.7 2.3 2.3 3.6 2 5 C1 4.3 0 3.7 -1 3 C-0.7 2 -0.3 1 0 0 Z };
\path[fill=black!100,even odd rule,shift={(284,82)}] svg {M0 0 C1 0 2 0 3 0 C3 1.3 3 2.6 3 4 C1.7 3.7 0.4 3.3 -1 3 C-0.7 2 -0.3 1 0 0 Z };
\path[fill=black!100,even odd rule,shift={(284,206)}] svg {M0 0 C1 0.3 2 0.7 3 1 C2.7 2.3 2.3 3.6 2 5 C1.3 5 0.7 5 0 5 C-0.6 3.1 -0.6 3.1 -1 1 C-0.7 0.7 -0.3 0.3 0 0 Z };
\path[fill=black!100,even odd rule,shift={(284,198)}] svg {M0 0 C1 0 2 0 3 0 C3 1.3 3 2.6 3 4 C2 4 1 4 0 4 C-0.3 3 -0.7 2 -1 1 C-0.7 0.7 -0.3 0.3 0 0 Z };
\path[fill=black!100,even odd rule,shift={(284,165)}] svg {M0 0 C1 0 2 0 3 0 C3 1.3 3 2.6 3 4 C1.7 3.7 0.4 3.3 -1 3 C-0.7 2 -0.3 1 0 0 Z };
\path[fill=black!100,even odd rule,shift={(57,207)}] svg {M0 0 C1 0 2 0 3 0 C3 1.3 3 2.6 3 4 C2 4 1 4 0 4 C0 2.7 0 1.4 0 0 Z };
\path[fill=black!100,even odd rule,shift={(57,199)}] svg {M0 0 C1 0 2 0 3 0 C3 1.3 3 2.6 3 4 C2 4 1 4 0 4 C0 2.7 0 1.4 0 0 Z };
\path[fill=black!100,even odd rule,shift={(284,190)}] svg {M0 0 C1 0 2 0 3 0 C2.7 1.3 2.3 2.6 2 4 C1 3.7 0 3.3 -1 3 C-0.7 2 -0.3 1 0 0 Z };
\path[fill=black!100,even odd rule,shift={(283,182)}] svg {M0 0 C1.3 0 2.6 0 4 0 C3.7 1.3 3.3 2.6 3 4 C2.3 4 1.7 4 1 4 C0.7 2.7 0.3 1.4 0 0 Z };
\path[fill=black!100,even odd rule,shift={(57,182)}] svg {M0 0 C1 0 2 0 3 0 C3 1.3 3 2.6 3 4 C2 4 1 4 0 4 C0 2.7 0 1.4 0 0 Z };
\path[fill=black!100,even odd rule,shift={(57,174)}] svg {M0 0 C1 0 2 0 3 0 C3 1.3 3 2.6 3 4 C2 4 1 4 0 4 C0 2.7 0 1.4 0 0 Z };
\path[fill=black!100,even odd rule,shift={(284,173)}] svg {M0 0 C1 0.3 2 0.7 3 1 C3 2 3 3 3 4 C1.7 3.7 0.4 3.3 -1 3 C-0.7 2 -0.3 1 0 0 Z };
\path[fill=black!100,even odd rule,shift={(58,165)}] svg {M0 0 C0.7 0.3 1.3 0.7 2 1 C2 2.3 2 3.6 2 5 C1 4.7 0 4.3 -1 4 C-0.7 2.7 -0.3 1.4 0 0 Z };
\path[fill=black!100,even odd rule,shift={(57,157)}] svg {M0 0 C1 0 2 0 3 0 C3 1.3 3 2.6 3 4 C2 4 1 4 0 4 C0 2.7 0 1.4 0 0 Z };
\path[fill=black!100,even odd rule,shift={(284,107)}] svg {M0 0 C1 0 2 0 3 0 C3 1.3 3 2.6 3 4 C2 4 1 4 0 4 C0 2.7 0 1.4 0 0 Z };
\path[fill=black!100,even odd rule,shift={(284,99)}] svg {M0 0 C1 0 2 0 3 0 C3 1.3 3 2.6 3 4 C2 4 1 4 0 4 C0 2.7 0 1.4 0 0 Z };
\path[fill=black!100,even odd rule,shift={(284,90)}] svg {M0 0 C1 0.3 2 0.7 3 1 C3 2 3 3 3 4 C2 4.3 1 4.7 0 5 C0 3.3 0 1.7 0 0 Z };
\path[fill=black!100,even odd rule,shift={(284,148)}] svg {M0 0 C1 0.3 2 0.7 3 1 C3 2 3 3 3 4 C2 4 1 4 0 4 C0 2.7 0 1.4 0 0 Z };
\path[fill=black!100,even odd rule,shift={(284,140)}] svg {M0 0 C1 0.3 2 0.7 3 1 C3 2 3 3 3 4 C2 4 1 4 0 4 C0 2.7 0 1.4 0 0 Z };
\path[fill=black!100,even odd rule,shift={(284,132)}] svg {M0 0 C1 0.3 2 0.7 3 1 C2.7 2 2.3 3 2 4 C1 3.7 0 3.3 -1 3 C-0.7 2 -0.3 1 0 0 Z };
\path[fill=black!100,even odd rule,shift={(306,94)}] svg {M0 0 C0.7 0.3 0.7 0.3 4 2 C2.7 1.7 1.4 1.3 0 1 C0 0.7 0 0.3 0 0 Z };
\end{tikzpicture}}
则函数$F(x) = \int_0^x f(t)\,\mathrm{d} t$的图形为\pickout{}

\noindent
\graphchoice{A}{%
\begin{tikzpicture}[scale=\DiagramScale,x=1pt,y=1pt,xscale=0.290994,yscale=-0.290994]
\path[use as bounding box] (0,0) rectangle (440,266);
\path[fill=black!100,even odd rule,shift={(166,1)}] svg {M0 0 C0.7 0 1.3 0 2 0 C5 7.3 7.9 14.5 11 22 C10.5 22.2 10.5 22.2 8 23 C6 22 4 21 2 20 C2 47.4 2 74.8 2 103 C4 103 6 103 8 103 C8 104 8 105 8 106 C6.3 106.3 4.7 106.7 3 107 C2.7 126.8 2.3 146.6 2 167 C22.5 167 42.9 167 64 167 C64 165.3 64 163.7 64 162 C65 162.3 66 162.7 67 163 C67 164.3 67 165.6 67 167 C81.2 167 95.4 167 110 167 C110.7 165.3 111.3 163.7 112 162 C112.6 160.8 113.3 159.5 113.9 158.3 C114.3 157.6 114.7 156.9 115 156.2 C115.4 155.5 115.8 154.8 116.2 154.1 C120.3 146.4 124.1 138.7 128 131 C149.4 131 170.9 131 193 131 C193 132 193 133 193 134 C172.2 134.3 151.4 134.7 130 135 C128.2 138.6 128.2 138.6 119 157 C118.2 158.5 118.2 158.5 114 166 C118.6 166 123.2 166 128 166 C128.3 164.7 128.7 163.4 129 162 C130 162.3 131 162.7 132 163 C132 164.3 132 165.6 132 167 C152.1 166.7 172.3 166.3 193 166 C193.3 164.7 193.7 163.4 194 162 C194.7 162 195.3 162 196 162 C196.3 163.3 196.7 164.6 197 166 C208.6 166 220.1 166 232 166 C231 164.3 230 162.7 229 161 C229 160.3 229 159.7 229 159 C234 159.8 238.6 161.6 243.2 163.4 C244.1 163.7 244.9 164 245.8 164.3 C246.2 164.4 246.2 164.4 248.1 165.2 C248.5 165.3 248.5 165.3 250.3 166 C252 167 252 167 253 170 C252.2 170.3 251.4 170.6 250.5 171 C243.6 173.6 236.8 176.2 230 179 C229.7 178.3 229.3 177.7 229 177 C230.1 174.5 231.6 172.3 233 170 C217.3 169.7 201.6 169.5 185.9 169.4 C178.6 169.4 171.3 169.3 164 169.1 C157.7 169 151.3 168.9 145 168.9 C141.6 168.9 138.2 168.9 134.9 168.8 C114.4 168.2 114.4 168.2 108.7 172.9 C104.9 177.4 102.3 182.4 100 187.8 C97.9 192.5 94.3 196.2 91 200 C90.3 200.9 89.6 201.7 88.8 202.6 C83 209.5 76.2 216.2 67 218 C61.6 217.6 57.3 216.1 52.6 213.7 C51.9 213.3 51.2 213 50.5 212.6 C33.8 203.9 19.7 190.3 5.9 177.6 C4 176 4 176 2 175 C2 194.1 2 213.3 2 233 C4 233 6 233 8 233 C8 234 8 235 8 236 C6 236 4 236 2 236 C2 242.3 2 248.5 2 255 C1 255 0 255 -1 255 C-1 226.9 -1 198.9 -1 170 C-29.7 170 -58.4 170 -88 170 C-88 169 -88 168 -88 167 C-80.7 166.7 -73.5 166.3 -66 166 C-65.7 164.7 -65.3 163.4 -65 162 C-64.3 162 -63.7 162 -63 162 C-62.8 162.7 -62.8 162.7 -62 166 C-43.8 166 -25.7 166 -7 166 C-9.9 162.4 -12.7 159.2 -16 156 C-16.8 155.2 -17.7 154.3 -18.6 153.4 C-19.5 152.6 -20.4 151.7 -21.3 150.8 C-22.8 149.3 -24.3 147.7 -25.8 146.2 C-26.9 145.2 -27.9 144.2 -28.9 143.2 C-34.3 137.8 -39.7 132.4 -44.6 126.6 C-46 125 -47.4 123.6 -48.9 122.1 C-51 120 -51 120 -51 118 C-51.7 118 -52.3 118 -53 118 C-54.6 116.3 -56 114.5 -57.5 112.6 C-61.9 107.1 -61.9 107.1 -64 105 C-64 104 -64 103 -64 102 C-60.1 103.6 -57.9 106 -55.2 109.1 C-54.7 109.6 -54.7 109.6 -52.5 112.1 C-51.7 113.1 -50.9 114 -50 115 C-49.3 115.8 -48.7 116.5 -48 117.3 C-45.4 120.2 -42.8 123.2 -40.4 126.2 C-38 129 -35.6 131.4 -33 134 C-31.9 135.1 -30.7 136.3 -29.6 137.4 C-29.3 137.7 -29.3 137.7 -27.7 139.3 C-19.1 147.9 -10.6 156.4 -2 165 C-1.9 150.7 -1.7 136.3 -1.6 122 C-1.5 115.3 -1.5 108.6 -1.4 102 C-0.9 44.5 -0.9 44.5 -1 21 C-1.8 21.3 -2.5 21.7 -3.3 22 C-6 23 -6 23 -9 23 C-6.5 15.1 -3.3 7.6 0 0 Z M4 170 C5.6 173.3 6.4 174.6 8.9 176.9 C9.4 177.4 10 178 10.6 178.5 C11.2 179.1 11.8 179.6 12.4 180.2 C13.1 180.8 13.7 181.3 14.3 181.9 C50.8 215 50.8 215 66.8 215 C78.9 212.9 87.9 200.5 94.6 191.1 C96.5 188.4 98.2 185.7 100 183 C100.7 181.9 101.4 180.9 102.1 179.8 C107 172.3 107 172.3 107 170 C73 170 39 170 4 170 Z };
\path[fill=black!100,even odd rule,shift={(188,2)}] svg {M0 0 C5.6 0 11.2 0 17 0 C16.7 1.7 16.3 3.3 16 5 C15.3 5 14.7 5 14 5 C14 4 14 3 14 2 C11.4 2 8.7 2 6 2 C5 4.6 4 7.3 3 10 C5.3 9.3 7.6 8.7 10 8 C9.7 10 9.3 12 9 14 C6.7 13.7 4.4 13.3 2 13 C1.8 14.5 1.7 16 1.5 17.4 C1.4 18.3 1.3 19.1 1.2 19.9 C1 22 1 22 1 23 C-1.6 23 -4.3 23 -7 23 C-6.6 22.5 -6.1 21.9 -5.6 21.4 C-3.7 18.6 -2.8 15.9 -1.8 12.6 C-1.5 11.5 -1.1 10.4 -0.8 9.3 C0 6 0 6 0 0 Z };
\path[fill=black!100,even odd rule,shift={(28,114)}] svg {M0 0 C0.7 0 1.3 0 2 0 C2.2 0.7 2.4 1.3 2.6 2 C3.5 5 4.5 7.9 5.4 10.9 C5.5 11.4 5.5 11.4 6.3 14 C6.5 14.5 6.5 14.5 7.3 17 C7.6 17.9 7.8 18.8 8.1 19.8 C9 22 9 22 11 24 C8.4 24 5.7 24 3 24 C3.3 23 3.7 22 4 21 C3.1 18.3 3.1 18.3 2 16 C0.4 16 -1.3 16 -3 16 C-3.7 18 -4.4 20 -5 22 C-4.7 22.7 -4.3 23.3 -4 24 C-6 24 -8 24 -10 24 C-9 23.7 -8 23.3 -7 23 C-6.5 21 -6.1 18.9 -5.7 16.9 C-4.3 11.1 -2.1 5.6 0 0 Z M-1 7 C-1.7 9.3 -2.3 11.6 -3 14 C-1.4 14 0.3 14 2 14 C1.6 11.1 1.6 11.1 1 8 C0.3 7.7 -0.3 7.3 -1 7 Z };
\path[fill=black!100,even odd rule,shift={(299,179)}] svg {M0 0 C1.9 1.1 1.9 1.1 3 3 C3.6 6.9 3.2 8.7 1.2 12.1 C0.5 13.1 -0.3 14 -1 15 C-1.6 15.8 -2.2 16.6 -2.8 17.4 C-3.2 17.9 -3.6 18.4 -4 19 C-1.4 19 1.3 19 4 19 C3.7 20.3 3.3 21.6 3 23 C-1.3 23 -5.6 23 -10 23 C-4 15 -4 15 -2.3 12.9 C-0.5 10.2 -0.7 8.2 -1 5 C-2 4 -2 4 -4.4 3.9 C-7 4 -7 4 -9 5 C-8.6 3.1 -8.6 3.1 -8 1 C-5.2 -0.4 -3.1 -0.3 0 0 Z };
\path[fill=black!100,even odd rule,shift={(363,178)}] svg {M0 0 C1.9 2.1 1.9 2.1 3 5 C2.7 7 2.4 9 2 11 C2.7 12.3 3.3 13.7 4 15 C3.8 17.7 3.3 19.4 1.9 21.8 C-1.2 23.8 -4.3 23.7 -8 24 C-8.3 23 -8.7 22 -9 21 C-6.4 21 -3.7 21 -1 21 C-0.3 18.7 -0.3 18.7 0 16 C-1.3 14.3 -1.3 14.3 -3 13 C-3.7 12.3 -4.3 11.7 -5 11 C-4 10.3 -3 9.7 -2 9 C-1.8 6.4 -1.8 6.4 -2 4 C-4 4 -6 4 -8 4 C-7.7 3 -7.3 2 -7 1 C-4.4 -0.3 -2.9 -0.4 0 0 Z };
\path[fill=black!100,even odd rule,shift={(150,177)}] svg {M0 0 C3.1 -0.2 5.4 -0.3 8.3 1 C10.9 4.1 11.3 7.1 11 11 C9.4 16.3 6.8 20.1 2 23 C-0.8 23.8 -0.8 23.8 -3 24 C-1.3 22.3 0.3 20.7 2.1 19.2 C6.5 14.1 6.3 8.4 6 2 C3.7 2 1.4 2 -1 2 C-0.7 1.3 -0.3 0.7 0 0 Z };
\path[fill=black!100,even odd rule,shift={(149,178)}] svg {M0 0 C-1.5 3.1 -3.2 6 -5 9 C-6.6 12.2 -6.5 15.5 -6 19 C-4 21.5 -4 21.5 -2 23 C-5.6 22.8 -7.4 22.6 -9.9 19.9 C-11.3 16 -11.2 12.9 -10 9 C-7.9 5.1 -5 0 0 0 Z };
\path[fill=black!100,even odd rule,shift={(234,181)}] svg {M0 0 C0.3 7.6 0.7 15.2 1 23 C-1 23 -3 23 -5 23 C-4.7 22.3 -4.3 21.7 -4 21 C-3.9 19.6 -3.9 18.3 -3.9 16.9 C-3.9 16.1 -3.9 15.3 -3.9 14.5 C-3.9 14.1 -3.9 14.1 -3.9 12 C-3.9 11.2 -3.9 10.3 -3.9 9.5 C-3.9 8.7 -3.9 7.9 -3.9 7.1 C-3.9 6.3 -3.9 5.6 -3.9 4.8 C-4 3 -4 3 -5 1 C-3 0 -3 0 0 0 Z };
\path[fill=black!100,even odd rule,shift={(231,1)}] svg {M0 0 C4.1 3.1 6.4 4.9 8 10 C8.9 17.1 8.4 22.3 4 28 C1.8 29.4 1.8 29.4 0 30 C0.5 29.3 0.9 28.5 1.4 27.8 C4.6 22.2 4.5 16.3 4 10 C3.2 7.4 2.3 5.4 1 3 C0.7 2 0.3 1 0 0 Z };
\path[fill=black!100,even odd rule,shift={(156,94)}] svg {M0 0 C0 0.6 0.1 1.3 0.1 1.9 C0.2 4.7 0.3 7.5 0.4 10.4 C0.5 11.4 0.5 12.3 0.6 13.4 C0.6 14.3 0.6 15.2 0.7 16.2 C0.7 17.1 0.8 17.9 0.8 18.8 C1 21 1 21 2 23 C-0.3 23 -2.6 23 -5 23 C-4.3 22.7 -3.7 22.3 -3 22 C-3 15.7 -3 9.5 -3 3 C-3.7 2.7 -4.3 2.3 -5 2 C-2 0 -2 0 0 0 Z };
\path[fill=black!100,even odd rule,shift={(209,3)}] svg {M0 0 C0.7 0.3 1.3 0.7 2 1 C1.7 1.6 1.4 2.1 1 2.7 C-2.6 10.8 -0.7 19 2 27 C-1.4 25.9 -1.6 25.5 -3.2 22.5 C-5.3 17.9 -5.9 13.9 -5 9 C-3.8 5.8 -2.4 2.4 0 0 Z };
\path[fill=black!100,even odd rule,shift={(107,181)}] svg {M0 0 C0.3 7.6 0.7 15.2 1 23 C-1 23 -3 23 -5 23 C-4.3 22.7 -3.7 22.3 -3 22 C-3 15.7 -3 9.5 -3 3 C-3.7 2.7 -4.3 2.3 -5 2 C-2 0 -2 0 0 0 Z };
\path[fill=black!100,even odd rule,shift={(426,163)}] svg {M0 0 C1.6 0 3.3 0 5 0 C5.3 1 5.7 2 6 3 C6.6 2.5 7.2 2 7.8 1.5 C10 0 10 0 14 0 C14 0.7 14 1.3 14 2 C13.6 2.1 13.6 2.1 11.6 2.8 C10.7 3.2 9.9 3.6 9 4 C7.8 7.6 8.5 8.6 10 12 C10.7 11.7 11.3 11.3 12 11 C11.3 12.3 10.7 13.6 10 15 C8.1 14.7 8.1 14.7 6 14 C4 11 4 11 4 8 C2.4 2.4 2.4 2.4 0 0 Z };
\path[fill=black!100,even odd rule,shift={(216,10)}] svg {M0 0 C2.4 -0.1 2.4 -0.1 5 0 C5.7 0.7 6.3 1.3 7 2 C7.7 1.3 8.3 0.7 9 0 C10.3 0 11.6 0 13 0 C13 0.7 13 1.3 13 2 C11.3 2.7 9.7 3.3 8 4 C8 6.6 8 9.3 8 12 C9 12.3 10 12.7 11 13 C10.3 13.3 10.3 13.3 7 15 C6.3 13.7 5.7 12.4 5 11 C4 11.3 3 11.7 2 12 C2.2 10.9 2.4 9.8 2.6 8.7 C3 5 3 5 2 2 C1.3 1.3 0.7 0.7 0 0 Z };
\path[fill=black!100,even odd rule,shift={(7,113)}] svg {M0 0 C0 3 0 3 -1.4 5.8 C-4.1 11.4 -4.4 16.1 -2.7 22 C-1.5 25.3 0.2 28.1 2 31 C-2.2 29.5 -3.8 26.8 -6 23 C-8.1 17.7 -7.4 12.4 -5.2 7.2 C-3.8 4.3 -2.4 2.2 0 0 Z };
\path[fill=black!100,even odd rule,shift={(46,113)}] svg {M0 0 C3.6 1.6 5.4 3.2 7.2 6.7 C9 12 9.1 18.1 7.2 23.3 C5.8 25.8 4.6 27.7 2 29 C2.3 28.2 2.6 27.5 2.9 26.7 C5.5 20.1 6.6 13 3.9 6.2 C2.7 4.1 1.4 2 0 0 Z };
\path[fill=black!100,even odd rule,shift={(82,191)}] svg {M0 0 C4.3 0 8.6 0 13 0 C13 0.3 13 0.7 13 1 C8.7 1 4.4 1 0 1 C0 0.7 0 0.3 0 0 Z };
\path[fill=black!100,even odd rule,shift={(428,174)}] svg {M0 0 C0 1 0 2 0 3 C-1.3 3.3 -2.6 3.7 -4 4 C-4 3 -4 2 -4 1 C-2.7 0.7 -1.4 0.3 0 0 Z };
\path[fill=black!100,even odd rule,shift={(214,22)}] svg {M0 0 C1.3 0 2.6 0 4 0 C4 0.7 4 1.3 4 2 C2.3 2.3 0.7 2.7 -1 3 C-0.7 2 -0.3 1 0 0 Z };
\path[fill=black!100,even odd rule,shift={(212,1)}] svg {M0 0 C0.3 0.7 0.7 1.3 1 2 C0.3 1.7 -0.3 1.3 -1 1 C-0.7 0.7 -0.3 0.3 0 0 Z };
\end{tikzpicture}
}
\hfill
\graphchoice{B}{%
\begin{tikzpicture}[scale=\DiagramScale,x=1pt,y=1pt,xscale=0.290994,yscale=-0.290994]
\path[use as bounding box] (0,0) rectangle (440,264);
\path[fill=black!100,even odd rule,shift={(166,0)}] svg {M0 0 C3.3 3.3 4.5 7.7 6.2 11.9 C6.4 12.4 6.4 12.4 7.3 14.7 C7.7 15.6 8.1 16.5 8.4 17.4 C8.6 17.8 8.6 17.8 9.4 19.9 C10 22 10 22 9 25 C6.7 23.7 4.4 22.4 2 21 C2 47.7 2 74.5 2 102 C3.7 102 5.3 102 7 102 C7 103 7 104 7 105 C5.3 105 3.7 105 2 105 C2 125.1 2 145.3 2 166 C21.8 166 41.6 166 62 166 C62.3 165 62.7 164 63 163 C64 163 65 163 66 163 C66.3 164 66.7 165 67 166 C72 166.2 77 166.4 82 166.5 C83.7 166.6 85.4 166.7 87.1 166.8 C101.1 167.5 101.1 167.5 105.5 164.1 C107.4 161.5 108.7 158.9 110 156 C111 154.2 112.1 152.5 113.1 150.7 C114.1 149 115.1 147.2 116.1 145.4 C116.6 144.6 117.1 143.7 117.6 142.8 C121.6 135.7 121.6 135.7 123.4 132.4 C123.8 131.6 124.2 130.7 124.7 129.9 C125.1 129.2 125.4 128.5 125.8 127.7 C127 126 127 126 130 125 C129.1 129.7 129.1 129.7 128 132 C128.7 132 129.3 132 130 132 C130 133.3 130 134.6 130 136 C129.3 136 128.7 136 128 136 C128 135.3 128 134.7 128 134 C125.6 136.2 124.1 138.4 122.5 141.3 C122.1 142.2 121.6 143 121.1 143.9 C120.7 144.8 120.2 145.6 119.7 146.6 C118.7 148.4 117.7 150.2 116.7 152 C116.3 152.9 115.8 153.7 115.3 154.6 C113 158.8 110.5 162.9 108 167 C110.6 166.9 113.2 166.8 115.8 166.7 C116.6 166.7 117.3 166.6 118 166.6 C123.5 166.3 123.5 166.3 126 164.6 C127.2 162.8 127.6 161.1 128 159 C128.7 159 129.3 159 130 159 C130.2 159.8 130.5 160.6 130.7 161.4 C131.1 162.3 131.6 163.1 132 164 C135.5 165 139 165.1 142.6 165.2 C143.7 165.2 144.8 165.2 146 165.3 C147.1 165.3 148.3 165.3 149.5 165.3 C150.1 165.3 150.1 165.3 153.2 165.4 C157.1 165.4 160.9 165.5 164.8 165.6 C167.4 165.6 170.1 165.6 172.7 165.7 C179.1 165.8 185.6 165.9 192 166 C192.3 165 192.7 164 193 163 C193.7 163 194.3 163 195 163 C195 164.3 195 165.6 195 167 C206.9 167 218.8 167 231 167 C230.3 165.8 230.3 165.8 227 160 C227.3 159.3 227.7 158.7 228 158 C231.6 159.5 231.6 159.5 250 167 C250 168 250 169 250 170 C246.7 171.4 243.5 172.7 240.2 174.1 C239.3 174.5 238.3 174.8 237.4 175.2 C236.5 175.6 235.6 176 234.7 176.3 C233.8 176.7 233 177 232.2 177.4 C230 178 230 178 227 177 C228.3 174.7 229.6 172.4 231 170 C230.2 170 229.5 170.1 228.7 170.1 C224.8 170.2 220.9 170.3 216.9 170.4 C216.6 170.4 216.6 170.4 214.7 170.5 C202.7 170.9 190.7 171.1 178.8 171.3 C178.1 171.3 177.3 171.3 176.6 171.3 C169.5 171.4 162.3 171.5 155.1 171.6 C153.1 171.6 151 171.6 148.9 171.7 C146.4 171.7 143.8 171.7 141.3 171.8 C140.1 171.8 139 171.8 137.8 171.8 C136.8 171.8 135.8 171.8 134.8 171.8 C132 172 132 172 129.3 172.6 C128.5 172.8 127.8 172.9 127 173 C126.3 172.3 125.7 171.7 125 171 C118.7 169.4 111.2 169 105 171 C100.1 174.7 96.9 180.2 94 185.5 C93.7 186 93.3 186.5 93 187 C92.3 187 91.7 187 91 187 C90.9 187.3 90.9 187.3 90.2 188.7 C86 196.5 77.9 202.4 70 206 C56.8 209.8 42.7 202.3 31.3 196 C23.3 191.4 16.1 186.2 9 180.4 C6 178 6 178 1 175 C1.3 192.8 1.7 210.6 2 229 C3.7 229.3 5.3 229.7 7 230 C7 231 7 232 7 233 C5 233 3 233 1 233 C1 239.9 1 246.9 1 254 C0 254 -1 254 -2 254 C-2 227.6 -2 201.2 -2 174 C-6.4 177.5 -10.4 180.9 -13.8 185.4 C-16.3 188.4 -19.1 190.8 -22.1 193.3 C-24.4 195.3 -26.1 197.6 -28 200 C-29.7 201.6 -31.5 203 -33.3 204.5 C-35.7 206.6 -37.8 208.8 -40 211 C-41.3 212.3 -42.7 213.7 -44 215 C-44.6 215.6 -45.3 216.3 -45.9 216.9 C-47.4 218.4 -48.9 219.9 -50.4 221.4 C-52.2 223.2 -53.9 224.9 -55.6 226.7 C-56.5 227.5 -57.3 228.4 -58.2 229.3 C-59 230.2 -59.9 231.1 -60.8 232 C-63 234 -63 234 -65 234 C-64.6 229.5 -61.5 227.1 -58.5 224.1 C-57.9 223.6 -57.3 223 -56.7 222.4 C-54.8 220.5 -52.9 218.6 -51 216.7 C-49.7 215.4 -48.4 214.1 -47 212.8 C-43.6 209.3 -40.1 205.8 -36.6 202.4 C-33 198.8 -29.5 195.3 -25.9 191.8 C-19 184.8 -12 177.9 -5 171 C-25.8 170.4 -46.5 169.8 -67.3 169.9 C-69.4 169.9 -71.4 169.9 -73.5 170 C-78.3 170 -83.2 170 -88 170 C-88 168.7 -88 167.4 -88 166 C-80.7 166 -73.5 166 -66 166 C-65.7 164.3 -65.3 162.7 -65 161 C-64.3 161.3 -63.7 161.7 -63 162 C-63 163.3 -63 164.6 -63 166 C-61.9 166 -60.9 166 -59.8 166 C-55.7 166 -51.6 166 -47.5 165.9 C-45.8 165.9 -44.1 165.9 -42.3 165.9 C-28.8 165.9 -15.5 166.4 -2 167 C-2 119.2 -2 71.3 -2 22 C-3.7 22.7 -5.3 23.3 -7 24 C-7.5 24.2 -7.5 24.2 -10 25 C-8.2 19.2 -6.1 13.6 -3.9 8 C-3.6 7.3 -3.3 6.7 -3.1 6 C-1.1 1.1 -1.1 1.1 0 0 Z M128 167 C128 168 128 169 128 170 C128.7 170 129.3 170 130 170 C130 169 130 168 130 167 C129.3 167 128.7 167 128 167 Z M3 170 C8.6 176.8 8.6 176.8 11.5 178.9 C12.1 179.4 12.8 179.9 13.4 180.4 C14.1 180.9 14.8 181.4 15.4 181.9 C16.1 182.4 16.8 182.9 17.6 183.5 C31.2 193.4 47.6 204 65 204 C80 201.3 90.8 185.7 99 174 C99.7 172.7 100.3 171.4 101 170 C68.7 170 36.3 170 3 170 Z };
\path[fill=black!100,even odd rule,shift={(35.9375,114.25)}] svg {M0 0 C1.5 2.4 1.5 3.9 1.1 6.8 C-0.4 9.1 -0.4 9.1 -1.9 10.8 C-0.9 11.1 0 11.4 1.1 11.8 C1.8 14.5 1.8 14.5 2.1 17.8 C0.5 19.9 0.5 19.9 -1.9 21.8 C-5.4 22.1 -5.4 22.1 -9.2 22 C-10.5 22 -11.7 21.9 -13 21.9 C-13.5 21.9 -13.5 21.9 -15.9 21.8 C-15.3 21.4 -14.6 21.1 -13.9 20.8 C-13.9 13.8 -13.9 6.9 -13.9 -0.2 C-14.6 -0.6 -15.3 -0.9 -15.9 -1.2 C-10.5 -2.3 -4.9 -3.2 0 0 Z M-9.9 -0.2 C-9.9 2.7 -9.9 5.7 -9.9 8.8 C-8 8.8 -6 8.8 -3.9 8.8 C-3.3 5.9 -3.3 5.9 -2.9 2.8 C-3.6 1.8 -4.3 0.8 -4.9 -0.2 C-7.6 -0.4 -7.6 -0.4 -9.9 -0.2 Z M-9.9 10.8 C-9.9 14.1 -9.9 17.3 -9.9 20.8 C-8 20.8 -6 20.8 -3.9 20.8 C-1.7 15.3 -1.7 15.3 -2.8 12.4 C-3.2 11.8 -3.6 11.3 -3.9 10.8 C-5.9 10.8 -7.9 10.8 -9.9 10.8 Z };
\path[fill=black!100,even odd rule,shift={(186,1)}] svg {M0 0 C5.6 0 11.2 0 17 0 C16.7 1.7 16.3 3.3 16 5 C15.3 4.3 14.7 3.7 14 3 C11.8 2.6 11.8 2.6 9.4 2.4 C8.6 2.3 7.7 2.2 6.9 2.1 C6.3 2.1 5.6 2.1 5 2 C4.7 4.3 4.3 6.6 4 9 C8.8 8.1 8.8 8.1 11 7 C10.7 9.3 10.3 11.6 10 14 C9.3 14 8.7 14 8 14 C8 13.3 8 12.7 8 12 C6 12 4 12 2 12 C1.3 15 0.7 17.9 0 21 C0.7 21 1.3 21 2 21 C2 21.7 2 22.3 2 23 C-0.3 23 -2.6 23 -5 23 C-4.6 21.8 -4.2 20.6 -3.8 19.3 C-1.7 13.2 0.1 7.1 2 1 C1.3 0.7 0.7 0.3 0 0 Z };
\path[fill=black!100,even odd rule,shift={(295.625,177)}] svg {M0 0 C4 1.7 4 1.7 5.4 4 C5.7 11 1.5 15.6 -2.6 21 C4.1 20.1 4.1 20.1 6.4 19 C6 20.7 5.7 22.3 5.4 24 C1.1 24 -3.2 24 -7.6 24 C-6 20.7 -5 19.3 -2.6 16.8 C1 12.9 1.4 9.2 1.4 4 C0.9 4.1 0.9 4.1 -1.4 4.5 C-4.6 5 -4.6 5 -7.6 5 C-3.4 0 -3.4 0 0 0 Z };
\path[fill=black!100,even odd rule,shift={(362,177)}] svg {M0 0 C2 2 2 2 2.1 4.4 C2.1 5.2 2 6.1 2 7 C2 11 2 11 2.8 14.4 C3 18 3 18 1 20.9 C-2.6 23.4 -4.7 23.4 -9 23 C-9 22 -9 21 -9 20 C-6.4 20 -3.7 20 -1 20 C-1 18 -1 16 -1 14 C-1.8 13.7 -1.8 13.7 -6 12 C-5.3 11.4 -4.7 10.8 -4 10.2 C-2 8 -2 8 -2 4 C-5.1 3 -6 3.1 -9 4 C-8.7 3 -8.3 2 -8 1 C-5.2 -0.4 -3.1 -0.2 0 0 Z };
\path[fill=black!100,even odd rule,shift={(147,138)}] svg {M0 0 C3.9 -0.5 5.6 -0.3 8.9 1.9 C11.5 5.7 11.4 7.5 10.6 11.8 C9.2 16.6 6.4 20.6 2 23 C1 22.7 0 22.3 -1 22 C-0.2 21.1 0.6 20.2 1.4 19.2 C5.7 14 6.2 10.8 6 4 C3.7 3.3 1.4 2.7 -1 2 C-0.7 1.3 -0.3 0.7 0 0 Z };
\path[fill=black!100,even odd rule,shift={(146,139)}] svg {M0 0 C-1.1 2.2 -2.1 4.2 -3.6 6.2 C-5.5 10 -6 13.7 -6 18 C-5.7 18.7 -5.3 19.3 -5 20 C-2.4 20.6 -2.4 20.6 0 21 C0 21.7 0 22.3 0 23 C-7.4 22.6 -7.4 22.6 -10.1 19.9 C-11.3 16.1 -11.1 12.9 -10 9 C-7.8 5 -5 0 0 0 Z };
\path[fill=black!100,even odd rule,shift={(110,180)}] svg {M0 0 C0 0.6 0.1 1.3 0.1 1.9 C0.2 4.7 0.3 7.5 0.4 10.4 C0.5 11.4 0.5 12.3 0.6 13.4 C0.6 14.3 0.6 15.2 0.7 16.2 C0.7 17.1 0.8 17.9 0.8 18.8 C1 21 1 21 2 23 C-0.3 23 -2.6 23 -5 23 C-4.7 22 -4.3 21 -4 20 C-3.9 17.3 -3.9 14.5 -3.9 11.8 C-3.9 11.1 -3.9 10.3 -4 9.5 C-4 7.7 -4 5.8 -4 4 C-4.7 3.7 -5.3 3.3 -6 3 C-6 2.3 -6 1.7 -6 1 C-3 0 -3 0 0 0 Z };
\path[fill=black!100,even odd rule,shift={(233,176)}] svg {M0 0 C0 6.9 0 13.9 0 21 C0.7 21 1.3 21 2 21 C2 21.7 2 22.3 2 23 C-0.3 23 -2.6 23 -5 23 C-5 22.3 -5 21.7 -5 21 C-4.3 21 -3.7 21 -3 21 C-3 15.1 -3 9.1 -3 3 C-4 2.7 -5 2.3 -6 2 C-3 0 -3 0 0 0 Z };
\path[fill=black!100,even odd rule,shift={(229,0)}] svg {M0 0 C4.5 2.2 6.2 4.3 8 9 C9.5 15.3 9.2 20.3 6 26 C3.8 28.5 3.8 28.5 2 30 C1.3 29.7 0.7 29.3 0 29 C0.5 28.6 0.9 28.1 1.4 27.7 C4.7 24.2 4.2 20.1 4.2 15.6 C4.3 14.3 4.3 13 4.3 11.7 C4 7.8 3.2 6.1 1 3 C0.7 2 0.3 1 0 0 Z };
\path[fill=black!100,even odd rule,shift={(210,0)}] svg {M0 0 C0.7 0.3 1.3 0.7 2 1 C1.4 1.8 0.7 2.5 0.1 3.3 C-4 8.6 -3.5 14.6 -3 21 C-2.3 23.6 -1.3 25.6 0 28 C0.3 28.7 0.7 29.3 1 30 C-2.8 28.8 -4.4 26.5 -6.2 23 C-8.2 18.1 -8.5 13.9 -6.9 8.8 C-5.3 5 -3.1 2.7 0 0 Z };
\path[fill=black!100,even odd rule,shift={(155,92)}] svg {M0 0 C0 0.6 0.1 1.3 0.1 1.9 C0.2 4.7 0.3 7.5 0.4 10.4 C0.5 11.4 0.5 12.3 0.6 13.4 C0.6 14.3 0.6 15.2 0.7 16.2 C0.7 17.1 0.8 17.9 0.8 18.8 C1 21 1 21 2 23 C-0.6 23 -3.3 23 -6 23 C-5.3 22.3 -4.7 21.7 -4 21 C-3.7 19.1 -3.7 19.1 -3.6 16.9 C-3.5 16.1 -3.5 15.3 -3.5 14.4 C-3.4 13.6 -3.4 12.7 -3.4 11.9 C-3.3 11 -3.3 10.2 -3.3 9.3 C-3.2 7.2 -3.1 5.1 -3 3 C-4 2.7 -5 2.3 -6 2 C-2.2 0 -2.2 0 0 0 Z };
\path[fill=black!100,even odd rule,shift={(217,8)}] svg {M0 0 C1 0.7 2 1.3 3 2 C3 2.7 3 3.3 3 4 C3.7 4 4.3 4 5 4 C5.2 6.1 5.4 8.1 5.6 10.2 C5.7 10.8 5.9 11.4 6 12 C6.7 12.3 7.3 12.7 8 13 C8 13.7 8 14.3 8 15 C6.7 15.3 5.4 15.7 4 16 C3 14.7 2 13.4 1 12 C0 13.3 -1 14.6 -2 16 C-3.3 15.7 -4.6 15.3 -6 15 C-5.7 14.3 -5.3 13.7 -5 13 C-3.7 13 -2.4 13 -1 13 C-0.8 11.6 -0.5 10.1 -0.3 8.7 C-0.2 7.9 -0.1 7.1 0.1 6.3 C0 4 0 4 -1.5 2.2 C-2 1.8 -2.5 1.4 -3 1 C-2 0.7 -1 0.3 0 0 Z };
\path[fill=black!100,even odd rule,shift={(8,111)}] svg {M0 0 C0.7 0.3 1.3 0.7 2 1 C1.8 1.2 1.8 1.2 0.9 2.3 C-3.7 8.5 -4.1 14.4 -3 22 C-2 24.8 -0.8 26.6 1 29 C1.3 29.7 1.7 30.3 2 31 C-2.2 29.5 -3.8 26.8 -6 23 C-7.9 18.1 -7.4 13.2 -5.8 8.4 C-4.1 5.2 -2.3 2.7 0 0 Z };
\path[fill=black!100,even odd rule,shift={(46,111)}] svg {M0 0 C4.5 3.1 6.4 5.8 8 11 C9 17.5 7.7 22.6 4 28 C1.8 30 1.8 30 0 31 C0.3 30.4 0.5 29.7 0.8 29.1 C1.2 28.1 1.6 27 2 26 C2.3 25.3 2.6 24.6 2.8 23.9 C4.9 18.4 5.4 14.8 4 9 C2.8 6.6 1.4 4.3 0 2 C0 1.3 0 0.7 0 0 Z };
\path[fill=black!100,even odd rule,shift={(425,161)}] svg {M0 0 C0.8 0.2 0.8 0.2 5 1 C5 1.7 5 2.3 5 3 C5.6 2.5 6.2 2 6.8 1.5 C9 0 9 0 13 0 C13 0.7 13 1.3 13 2 C12.6 2.1 12.6 2.1 10.6 2.8 C9.7 3.2 8.9 3.6 8 4 C6.8 7.6 7.5 8.6 9 12 C9.7 11.7 10.3 11.3 11 11 C10.3 12.3 9.7 13.6 9 15 C7.1 14.9 7.1 14.9 5 14 C3.9 11.3 3.9 11.3 3 8 C1.1 2.2 1.1 2.2 0 0 Z };
\path[fill=black!100,even odd rule,shift={(96,189)}] svg {M0 0 C0.7 0.7 1.3 1.3 2 2 C-2.6 2 -7.2 2 -12 2 C-7.2 -0.4 -5.2 -0.3 0 0 Z };
\path[fill=black!100,even odd rule,shift={(427,171)}] svg {M0 0 C0 1 0 2 0 3 C-0.7 3.7 -1.3 4.3 -2 5 C-3 5 -4 5 -5 5 C-4 2 -4 2 -1.9 0.8 C-1.3 0.5 -0.7 0.3 0 0 Z };
\path[fill=black!100,even odd rule,shift={(225,8)}] svg {M0 0 C0.7 0.3 1.3 0.7 2 1 C2 1.7 2 2.3 2 3 C0.3 3.3 -1.3 3.7 -3 4 C-2.7 3 -2.3 2 -2 1 C-1.3 0.7 -0.7 0.3 0 0 Z };
\path[fill=black!100,even odd rule,shift={(294,139)}] svg {M0 0 C0.7 0 1.3 0 2 0 C2 1.3 2 2.6 2 4 C1.3 4 0.7 4 0 4 C-0.3 3 -0.7 2 -1 1 C-0.7 0.7 -0.3 0.3 0 0 Z };
\path[fill=black!100,even odd rule,shift={(294,153)}] svg {M0 0 C0.7 0 1.3 0 2 0 C2 1 2 2 2 3 C1.3 3 0.7 3 0 3 C0 2 0 1 0 0 Z };
\path[fill=black!100,even odd rule,shift={(294,146)}] svg {M0 0 C0.7 0 1.3 0 2 0 C2 1 2 2 2 3 C1.3 3 0.7 3 0 3 C0 2 0 1 0 0 Z };
\end{tikzpicture}
}
\hfill
\graphchoice{C}{%
\begin{tikzpicture}[scale=\DiagramScale,x=1pt,y=1pt,xscale=0.290994,yscale=-0.290994]
\path[use as bounding box] (0,0) rectangle (440,267);
\path[fill=black!100,even odd rule,shift={(164,11)}] svg {M0 0 C0.7 0 1.3 0 2 0 C3.5 3.6 5 7.2 6.5 10.9 C6.7 11.4 6.7 11.4 7.8 14 C8.2 15 8.6 16 9 17 C9.4 17.9 9.8 18.8 10.2 19.8 C11 22 11 22 11 24 C7.9 22.8 5 21.5 2 20 C2 49.7 2 79.4 2 110 C4 110 6 110 8 110 C8 111 8 112 8 113 C7.3 113 6.7 113 6 113 C18.6 128.7 39.4 149.1 60 153 C67.9 153.8 72.8 151.4 78.9 146.4 C94.6 131.8 106.1 111.1 116.4 92.6 C116.8 91.8 117.2 91.1 117.6 90.4 C119.5 86.8 121.4 83.3 123 79.6 C123.3 78.9 123.7 78.2 124 77.5 C124.6 76.2 125.2 74.8 125.7 73.5 C126.5 71.7 126.5 71.7 129 69 C136.6 67.3 145 67.9 152.8 67.9 C154 67.9 155.2 67.9 156.4 67.9 C158.9 67.9 161.3 67.9 163.8 67.9 C167.6 67.9 171.4 67.9 175.2 67.9 C177.6 67.9 180 67.9 182.5 67.9 C183.6 67.9 184.7 67.9 185.9 67.9 C192 67.9 197.9 68.2 204 69 C204 69.7 204 70.3 204 71 C180 71.6 156 72.2 132 72 C131.7 72.5 131.4 73.1 131.1 73.7 C90.7 149.5 90.7 149.5 69 157 C56.5 159 44.6 151.4 34.9 144.5 C24.5 136.8 14.8 128.1 5.9 118.8 C4 117 4 117 2 116 C2.3 135.1 2.7 154.3 3 174 C22.8 174 42.6 174 63 174 C63.3 173 63.3 173 65 168 C65.7 168.3 66.3 168.7 67 169 C67 171 67 173 67 175 C87.1 174.7 107.3 174.3 128 174 C128.3 172.3 128.7 170.7 129 169 C130 169.3 131 169.7 132 170 C132 171.7 132 173.3 132 175 C152.1 174.7 172.3 174.3 193 174 C193.3 172.3 193.7 170.7 194 169 C194.7 169 195.3 169 196 169 C196.3 171 196.7 173 197 175 C209.5 174.7 222.1 174.3 235 174 C234 172 233 170 232 168 C232.3 167.3 232.7 166.7 233 166 C236.6 167.5 236.6 167.5 255 175 C255 176 255 177 255 178 C251.7 179.4 248.5 180.7 245.2 182.1 C244.3 182.5 243.3 182.8 242.4 183.2 C241.5 183.6 240.6 184 239.7 184.3 C238.8 184.7 238 185 237.2 185.4 C235 186 235 186 232 185 C233.3 182.7 234.6 180.4 236 178 C158.8 178 81.6 178 2 178 C2 197.5 2 216.9 2 237 C4.3 237 6.6 237 9 237 C9 237.7 9 238.3 9 239 C7 239.3 5 239.7 3 240 C2.7 244.9 2.3 249.9 2 255 C1 255 0 255 -1 255 C-1 229.6 -1 204.2 -1 178 C-28.7 178 -56.4 178 -85 178 C-85.3 177 -85.7 176 -86 175 C-82.7 174.8 -82.7 174.8 -66 174 C-65.8 173 -65.8 173 -65 168 C-64.3 168 -63.7 168 -63 168 C-63 169 -63 170 -63 171 C-59.6 169.6 -57.8 167.8 -55.4 164.9 C-52.3 161.1 -52.3 161.1 -49 160 C-48.7 159.4 -48.4 158.8 -48.1 158.2 C-46.6 155.2 -44.4 153.3 -42 151 C-41.3 150.2 -40.6 149.5 -39.8 148.7 C-38 147 -38 147 -36 147 C-35.7 146.4 -35.5 145.8 -35.2 145.2 C-33.9 142.9 -32.9 141.8 -30.9 140.2 C-30.2 139.7 -29.6 139.2 -28.9 138.7 C-28.6 138.4 -28.6 138.4 -26.8 137 C-1.4 115.4 -1.4 115.4 0.3 102 C0.9 94.8 0.3 87.6 -0.3 80.4 C-1.4 69.4 -1.1 58.3 -1 47.1 C-0.9 41.4 -0.9 35.7 -1 30 C-1 29 -1 28 -1 26.9 C-1 24.6 -1 22.3 -1 20 C-2.3 20.7 -2.3 20.7 -9 24 C-6.7 15.6 -3.5 7.9 0 0 Z M-6.9 119.6 C-7.7 120.4 -8.5 121.2 -9.3 122 C-10.2 122.9 -11.1 123.8 -11.9 124.7 C-12.8 125.6 -13.7 126.4 -14.6 127.4 C-16.5 129.2 -18.4 131.1 -20.2 133 C-22.6 135.4 -25 137.8 -27.5 140.2 C-29.8 142.5 -32 144.8 -34.3 147.1 C-35.2 147.9 -36.1 148.8 -37 149.7 C-37.7 150.5 -38.5 151.3 -39.4 152.1 C-40.1 152.8 -40.8 153.5 -41.5 154.2 C-43 156 -43 156 -43 158 C-43.6 158.3 -44.2 158.5 -44.8 158.8 C-47.2 160.1 -48.8 161.6 -50.8 163.5 C-51.4 164.1 -52 164.7 -52.6 165.3 C-54 167 -54 167 -54 169 C-54.6 169.2 -55.3 169.5 -55.9 169.8 C-58 171 -58 171 -58.6 173.1 C-58.8 173.7 -58.9 174.3 -59 175 C-51.6 175 -44.2 175 -36.8 175.1 C-33.3 175.1 -29.9 175.1 -26.4 175.1 C-22.5 175.1 -18.5 175.1 -14.6 175.1 C-13.3 175.1 -12.1 175.1 -10.8 175.1 C-10.2 175.1 -10.2 175.1 -7.3 175.1 C-6.3 175.1 -5.3 175.1 -4.3 175.1 C-2 175 -2 175 -1 174 C-0.9 172.2 -0.9 170.4 -0.9 168.6 C-0.9 167.4 -0.9 166.2 -0.9 165 C-0.9 164.4 -0.9 164.4 -0.9 161.2 C-0.9 159.9 -0.9 158.6 -0.9 157.3 C-0.9 153.9 -0.9 150.5 -0.9 147 C-0.9 143.6 -0.9 140.1 -1 136.6 C-1 129.7 -1 122.9 -1 116 C-4 116 -4.8 117.6 -6.9 119.6 Z };
\path[fill=black!100,even odd rule,shift={(185,12)}] svg {M0 0 C5.3 0 10.6 0 16 0 C15.7 1.7 15.3 3.3 15 5 C14.7 4 14.3 3 14 2 C11 2 8.1 2 5 2 C4.3 4.6 3.7 7.3 3 10 C5.3 9.3 7.6 8.7 10 8 C9.7 10 9.3 12 9 14 C8 13.3 7 12.7 6 12 C3.8 12.4 3.8 12.4 2 13 C-0.1 19.8 -0.1 19.8 1 22 C-1.3 22 -3.6 22 -6 22 C-5.3 21.7 -4.7 21.3 -4 21 C-3.2 18.4 -2.5 15.9 -1.9 13.2 C-1.7 12.5 -1.5 11.8 -1.3 11.1 C-0.4 7.1 0 4.2 0 0 Z };
\path[fill=black!100,even odd rule,shift={(355,192)}] svg {M0 0 C2.8 -0.3 2.8 -0.3 6 0 C7.9 2.1 7.9 2.1 9 5 C8.7 7 8.4 9 8 11 C8.3 11.5 8.3 11.5 10 14 C9.8 17 9.5 19.2 7.9 21.8 C4.8 23.7 1.6 23.7 -2 24 C-2.3 23 -2.7 22 -3 21 C-1.7 21 -0.4 21 0.9 21.1 C2.3 21.1 3.6 21.1 5 21 C6 20 6 20 6.4 17.7 C6 15 6 15 3.5 12.7 C2.7 12.1 1.9 11.6 1 11 C1.7 10.4 2.3 9.8 3 9.2 C5 7 5 7 5 3 C2.7 3.3 0.4 3.7 -2 4 C-1.3 2.7 -0.7 1.4 0 0 Z };
\path[fill=black!100,even odd rule,shift={(24,125)}] svg {M0 0 C0.7 0.3 1.3 0.7 2 1 C1.7 1.5 1.7 1.5 0 4 C-0.2 5.9 -0.2 5.9 -0.2 8 C-0.2 8.8 -0.2 9.5 -0.2 10.3 C-0.2 11.1 -0.1 11.9 -0.1 12.7 C-0.1 13.1 -0.1 13.1 -0.1 15.1 C-0.1 17.1 -0 19 0 21 C3 21 5.9 21 9 21 C9.7 19.7 10.3 18.4 11 17 C11.3 17 11.7 17 12 17 C12 18.7 12 20.3 12 22 C7.9 23.4 4.3 23.8 0 23 C-3.6 20 -5.1 17.7 -6.1 13.1 C-6.4 8.5 -4.6 4.7 -2 1 C-1.3 0.7 -0.7 0.3 0 0 Z };
\path[fill=black!100,even odd rule,shift={(148,193)}] svg {M0 0 C7.4 0.4 7.4 0.4 10 3.2 C11.6 7.7 11.5 10.6 9.8 15.1 C7.3 19.2 5 22.4 0.2 23.8 C-0.5 23.9 -1.3 23.9 -2 24 C-1.6 23.6 -1.2 23.1 -0.8 22.6 C4.8 16.4 6.3 12.3 6 4 C3.7 3.3 1.4 2.7 -1 2 C-0.7 1.3 -0.3 0.7 0 0 Z };
\path[fill=black!100,even odd rule,shift={(146,194)}] svg {M0 0 C0 4 -1.8 5.7 -4 9 C-5 11.9 -5.1 13.7 -5.1 16.7 C-5.1 17.5 -5 18.3 -5 19.1 C-5 19.8 -5 20.4 -5 21 C-3.3 21.3 -1.7 21.7 0 22 C0 22.3 0 22.7 0 23 C-5.4 23.4 -5.4 23.4 -7.8 21.9 C-9.7 18.9 -9.9 16.6 -10 13 C-8.2 7.1 -4.9 3.6 0 0 Z };
\path[fill=black!100,even odd rule,shift={(105,195)}] svg {M0 0 C0 7.3 0 14.5 0 22 C0.7 22.3 1.3 22.7 2 23 C-0.3 23 -2.6 23 -5 23 C-4.7 22.3 -4.3 21.7 -4 21 C-3.8 19.6 -3.8 18.2 -3.7 16.9 C-3.6 16.1 -3.6 15.3 -3.6 14.4 C-3.5 13.6 -3.5 12.7 -3.4 11.9 C-3.4 11 -3.4 10.2 -3.3 9.3 C-3.2 7.2 -3.1 5.1 -3 3 C-3.7 2.7 -4.3 2.3 -5 2 C-2 0 -2 0 0 0 Z };
\path[fill=black!100,even odd rule,shift={(154,100)}] svg {M0 0 C0 0.6 0.1 1.2 0.1 1.8 C0.2 4.5 0.3 7.1 0.4 9.8 C0.5 10.7 0.5 11.7 0.6 12.6 C0.6 13.5 0.6 14.4 0.7 15.3 C0.7 16.2 0.8 17 0.8 17.8 C1 20 1 20 2 23 C0 23 -2 23 -4 23 C-4 15.7 -4 8.5 -4 1 C-2 0 -2 0 0 0 Z };
\path[fill=black!100,even odd rule,shift={(232,194)}] svg {M0 0 C0.3 7.6 0.7 15.2 1 23 C-1 23 -3 23 -5 23 C-4.3 22.7 -3.7 22.3 -3 22 C-3.1 20.7 -3.1 19.5 -3.2 18.2 C-3.3 16.5 -3.4 14.8 -3.5 13.2 C-3.5 12.8 -3.5 12.8 -3.6 10.7 C-3.7 10.3 -3.7 10.3 -3.8 8.2 C-3.8 7.5 -3.9 6.8 -3.9 6 C-4 4.3 -4 2.7 -4 1 C-2 0 -2 0 0 0 Z };
\path[fill=black!100,even odd rule,shift={(208,11)}] svg {M0 0 C0.7 0.3 1.3 0.7 2 1 C1.5 1.4 1.1 1.7 0.6 2.1 C-4 7.6 -3 14.1 -2.5 20.9 C-2 23.9 -1.2 26.3 0 29 C-4.8 24.5 -6.7 21.7 -7.4 15.2 C-7.2 9.8 -5 5.1 -1.4 1.3 C-1.1 1.1 -1.1 1.1 0 0 Z };
\path[fill=black!100,even odd rule,shift={(227,11)}] svg {M0 0 C3.8 1.5 5.7 3.7 8 7 C9.7 12 10 17.4 8.1 22.4 C7.1 24.3 6.1 26.2 5 28 C4 28 3 28 2 28 C2.7 27 3.3 26 4 25 C6.3 17.4 5.7 9.9 2 3 C1.3 2 0.7 1 0 0 Z };
\path[fill=black!100,even odd rule,shift={(215,19)}] svg {M0 0 C3 1 3 1 5 3 C6.7 2 8.3 1 10 0 C10.3 1 10.7 2 11 3 C9.3 3.7 7.7 4.3 6 5 C6 7.6 6 10.3 6 13 C7 12.7 8 12.3 9 12 C8.3 13.3 7.7 14.6 7 16 C6 16 5 16 4 16 C3.8 15.2 3.8 15.2 3 11 C2.3 11 1.7 11 1 11 C1 10.5 1 10.5 1.2 8.2 C1 4.7 0.4 3.6 -2 1 C-1.3 0.7 -0.7 0.3 0 0 Z };
\path[fill=black!100,even odd rule,shift={(426,180)}] svg {M0 0 C1.6 0 3.3 0 5 0 C5.3 1 5.7 2 6 3 C6.6 2.5 7.2 2 7.8 1.5 C10 0 10 0 13 0 C13 0.7 13 1.3 13 2 C10.5 3.6 10.5 3.6 8 5 C8.4 7.9 8.4 7.9 9 11 C9.7 11.3 10.3 11.7 11 12 C10.7 13 10.3 14 10 15 C7.8 14.6 6.2 14.1 4.2 12.9 C2.3 9.9 3 6.6 3 3 C1.5 1 1.5 1 0 0 Z };
\path[fill=black!100,even odd rule,shift={(7,123)}] svg {M0 0 C0 3 0 3 -1.4 5.6 C-4.2 11.5 -4.7 16.6 -2.7 22.8 C-1.6 25.6 -0.4 28.3 1 31 C-3.5 27.9 -5.4 25.2 -7 20 C-8 13.5 -6.7 8.4 -3 3 C-2 2 -1 1 0 0 Z };
\path[fill=black!100,even odd rule,shift={(290,193)}] svg {M0 0 C6.4 -0.4 6.4 -0.4 8.9 1.1 C10.5 4 10.4 5.8 10 9 C8 12.7 6.5 14.7 3 17 C4 15 5 13 6 11 C6 8.7 6 6.3 6 4 C0.2 3.9 0.2 3.9 -2 5 C-1.3 3.3 -0.7 1.7 0 0 Z };
\path[fill=black!100,even odd rule,shift={(44,123)}] svg {M0 0 C3.9 1.2 5.3 3.6 7.2 7.1 C9.7 12.8 9.2 18.3 7 24 C5.5 26.1 4 27.4 2 29 C3.9 23.2 3.9 23.2 5 21 C5.8 13.6 4.9 8.3 1 2 C0.7 1.3 0.3 0.7 0 0 Z };
\path[fill=black!100,even odd rule,shift={(291,211)}] svg {M0 0 C0.3 0.7 0.7 1.3 1 2 C2.5 1.8 2.5 1.8 10 1 C9.7 2.3 9.3 3.6 9 5 C4.7 5 0.4 5 -4 5 C-4 4.3 -4 3.7 -4 3 C-2 1.4 -2 1.4 0 0 Z };
\path[fill=black!100,even odd rule,shift={(26,124)}] svg {M0 0 C3.3 0.3 6.6 0.7 10 1 C9.7 2.7 9.3 4.3 9 6 C8.9 5.7 8.9 5.7 8.4 4.1 C7 2 7 2 3.4 1.2 C2.3 1.2 1.1 1.1 0 1 C0 0.7 0 0.3 0 0 Z };
\path[fill=black!100,even odd rule,shift={(80,205)}] svg {M0 0 C4.3 0 8.6 0 13 0 C13 0.3 13 0.7 13 1 C8.7 1 4.4 1 0 1 C0 0.7 0 0.3 0 0 Z };
\path[fill=black!100,even odd rule,shift={(215,30)}] svg {M0 0 C0 1 0 2 0 3 C-0.7 3.3 -0.7 3.3 -4 5 C-4.3 4 -4.7 3 -5 2 C-4.4 1.9 -3.7 1.7 -3.1 1.6 C-1 1 -1 1 0 0 Z };
\path[fill=black!100,even odd rule,shift={(428,191)}] svg {M0 0 C0 1 0 2 0 3 C-1.6 3.3 -3.3 3.7 -5 4 C-5 3.3 -5 2.7 -5 2 C-3.4 1.3 -1.7 0.7 0 0 Z };
\path[fill=black!100,even odd rule,shift={(208,40)}] svg {M0 0 C2 1 2 1 2 1 Z };
\end{tikzpicture}
}
\hfill
\graphchoice{D}{%
\begin{tikzpicture}[scale=\DiagramScale,x=1pt,y=1pt,xscale=0.290994,yscale=-0.290994]
\path[use as bounding box] (0,0) rectangle (440,265);
\path[fill=black!100,even odd rule,shift={(166,11)}] svg {M0 0 C0.7 0 1.3 0 2 0 C3.5 3.4 5.1 6.9 6.6 10.3 C6.8 10.8 6.8 10.8 7.9 13.3 C8.3 14.2 8.7 15.1 9.1 16.1 C9.5 17 9.9 17.8 10.3 18.7 C11 21 11 21 10 24 C8 23 6 22 4 21 C3.7 49.4 3.3 77.8 3 107 C5 107.7 7 108.3 9 109 C8.7 109.7 8.3 110.3 8 111 C6.3 111.3 4.7 111.7 3 112 C3 130.5 3 149 3 168 C23.1 168 43.3 168 64 168 C64.7 166.7 65.3 165.4 66 164 C66.7 165.3 67.3 166.6 68 168 C75.3 168.2 75.3 168.2 112 169 C112.7 167.3 113.3 165.7 114 164 C115.1 161.7 116.3 159.5 117.4 157.2 C117.8 156.6 118.1 155.9 118.4 155.3 C119.1 153.9 119.8 152.6 120.5 151.3 C121.6 149.2 122.6 147.2 123.7 145.2 C124.3 143.9 125 142.6 125.7 141.3 C126 140.7 126.3 140 126.7 139.4 C128.9 135.1 128.9 135.1 130 134 C132 133.9 134.1 133.9 136.1 133.9 C136.7 133.9 136.7 133.9 140 133.9 C141.5 133.9 142.9 133.9 144.3 133.9 C145.8 133.9 147.2 133.9 148.7 133.9 C152.5 133.9 156.4 133.9 160.2 133.9 C164.1 133.9 168 133.9 171.9 134 C179.6 134 187.3 134 195 134 C194.7 135 194.3 136 194 137 C173.2 137 152.4 137 131 137 C130.7 138.3 130.3 139.6 130 141 C129.3 142.5 128.6 144 127.9 145.5 C127.4 146.4 127 147.2 126.6 148.1 C126.1 149 125.7 149.9 125.2 150.8 C124.7 151.7 124.3 152.6 123.9 153.5 C121.4 158.4 118.8 163.2 116 168 C118.9 169 120.7 169.1 123.7 169.1 C124.5 169.1 125.3 169 126.1 169 C126.8 169 127.4 169 128 169 C128.3 167.7 128.7 166.4 129 165 C129.7 165 130.3 165 131 165 C131 166.3 131 167.6 131 169 C151.1 169 171.3 169 192 169 C192.3 168 192.7 167 193 166 C194 166 195 166 196 166 C196.3 167 196.7 168 197 169 C202.9 169.2 202.9 169.2 233 170 C232 167 231 164.1 230 161 C238.3 163.8 246.2 166.8 254 171 C251.8 173.2 250.3 173.8 247.5 174.9 C246.6 175.3 245.7 175.7 244.8 176 C243.9 176.4 243 176.7 242 177.1 C241.1 177.5 240.1 177.9 239.2 178.3 C232.3 181 232.3 181 230 181 C231.9 175.2 231.9 175.2 233 173 C193.1 173 153.1 173 112 173 C111.3 174.7 110.7 176.3 110 178 C101.6 194.3 87.5 216.6 69.6 223.4 C61.7 225.2 54.9 222.7 48 219 C47.6 218.8 47.6 218.8 45.7 217.8 C34.9 211.9 23.8 202.5 16 193 C16 192.3 16 191.7 16 191 C15.3 191 14.7 191 14 191 C12.4 189.2 11 187.4 9.5 185.6 C5.2 180.1 5.2 180.1 3 179 C3 196.2 3 213.3 3 231 C4.7 231.3 6.3 231.7 8 232 C7.7 233 7.3 234 7 235 C5.3 235 3.7 235 2 235 C2 241.3 2 247.5 2 254 C1 253.7 0 253.3 -1 253 C-1 252.2 -1 251.3 -1 250.5 C-1 242.6 -1 234.7 -0.9 226.8 C-0.9 222.7 -0.9 218.6 -0.9 214.6 C-0.9 210.7 -0.9 206.7 -0.9 202.8 C-0.9 201.3 -0.9 199.8 -0.9 198.3 C-0.9 196.2 -0.9 194.1 -0.9 192 C-0.9 190.8 -0.9 189.7 -0.9 188.4 C-1 184.9 -1.4 181.5 -2 178 C-2.4 178.4 -2.9 178.9 -3.3 179.3 C-6.6 182.6 -9.9 185.9 -13.2 189.2 C-13.6 189.6 -13.6 189.6 -15.2 191.2 C-16.7 192.7 -18.2 194.2 -19.7 195.7 C-21.2 197.2 -22.7 198.7 -24.2 200.2 C-24.8 200.9 -25.5 201.6 -26.2 202.3 C-26.5 202.7 -26.5 202.7 -28.2 204.4 C-30 206 -30 206 -32 206 C-32.1 206.3 -32.1 206.3 -32.8 207.9 C-34.2 210.4 -35.4 211 -38 212 C-38 212.7 -38 213.3 -38 214 C-39.5 215.4 -39.5 215.4 -41.4 217 C-44.9 219.8 -44.9 219.8 -46 222 C-46.7 222 -47.3 222 -48 222 C-48 222.7 -48 223.3 -48 224 C-48.7 224 -49.3 224 -50 224 C-50 224.7 -50 225.3 -50 226 C-53.6 230 -57.7 233.8 -62 237 C-62.7 237 -63.3 237 -64 237 C-63.6 232.6 -60.5 230.1 -57.5 227.2 C-56.9 226.6 -56.3 226 -55.7 225.4 C-53.8 223.5 -51.9 221.7 -50 219.8 C-48.7 218.5 -47.4 217.1 -46.1 215.8 C-42.6 212.4 -39.1 208.9 -35.6 205.5 C-30.1 200 -24.5 194.4 -18.9 188.9 C-17 187 -15 185 -13.1 183.1 C-11.9 181.9 -10.7 180.8 -9.5 179.6 C-9 179 -8.4 178.5 -7.9 177.9 C-4.1 174.2 -4.1 174.2 -3 172 C-30.4 172 -57.8 172 -86 172 C-86 171 -86 170 -86 169 C-79.1 169 -72.1 169 -65 169 C-65 167.7 -65 166.4 -65 165 C-64 165 -63 165 -62 165 C-62 166.3 -62 167.6 -62 169 C-41.9 169 -21.7 169 -1 169 C-1 160.7 -1 152.5 -1.1 143.9 C-1.1 116.2 -0.8 88.4 -0.5 60.7 C-0.5 59.7 -0.5 58.7 -0.5 57.7 C-0.3 45.5 -0.2 33.2 0 21 C-3 22 -5.9 23 -9 24 C-6 16 -3 8 0 0 Z M5 172 C5 176 6.5 177.2 9.1 180.2 C9.9 181.2 10.8 182.2 11.6 183.2 C12.4 184.1 13.2 185 14 186 C14.3 186.4 14.3 186.4 15.9 188.3 C17.2 189.9 18.6 191.5 20 193 C20.7 193 21.3 193 22 193 C22.3 194 22.7 195 23 196 C23.7 196 24.3 196 25 196 C25 196.7 25 197.3 25 198 C27.4 200.6 30.3 202.8 33 205 C34 205.8 35 206.7 36 207.5 C44.9 214.7 55 222.2 67 221 C84.3 216.4 96.5 194.9 105.2 180.7 C105.4 180.4 105.4 180.4 106.4 178.7 C107.4 177.2 108.2 175.6 109 174 C108.8 173.7 108.8 173.7 108 172 C74 172 40 172 5 172 Z };
\path[fill=black!100,even odd rule,shift={(188,10)}] svg {M0 0 C5.6 0 11.2 0 17 0 C16.7 1.7 16.3 3.3 16 5 C15.3 4.7 14.7 4.3 14 4 C14 3.3 14 2.7 14 2 C11 2 8.1 2 5 2 C4.9 2.9 4.7 3.9 4.6 4.9 C4 8 4 8 3 10 C5 10 7 10 9 10 C9 9.3 9 8.7 9 8 C9.7 8 10.3 8 11 8 C10.3 10.3 9.7 12.6 9 15 C8.3 14.3 7.7 13.7 7 13 C4.4 12.4 4.4 12.4 2 12 C1.8 13.5 1.6 14.9 1.4 16.4 C1.3 17.2 1.2 18 1.1 18.8 C1 21 1 21 2 23 C-0.6 23 -3.3 23 -6 23 C-5.3 22.7 -4.7 22.3 -4 22 C-1.6 15.1 0 7.3 0 0 Z };
\path[fill=black!100,even odd rule,shift={(297.1875,188)}] svg {M0 0 C2.8 1 2.8 1 4.8 4 C5 11.1 1 15.6 -3.2 21 C-0.2 20.7 2.8 20.3 5.8 20 C5.5 21.3 5.2 22.6 4.8 24 C0.5 24 -3.8 24 -8.2 24 C-6.5 20.6 -5 18.5 -2.6 15.7 C-0.2 12.7 1.3 10.8 1.2 6.9 C1.1 6 1 5 0.8 4 C-1.5 4 -3.8 4 -6.2 4 C-6.2 4.7 -6.2 5.3 -6.2 6 C-6.8 5.7 -7.5 5.3 -8.2 5 C-4 0 -4 0 0 0 Z };
\path[fill=black!100,even odd rule,shift={(361.625,187.625)}] svg {M0 0 C2.4 1.4 2.4 1.4 3.2 3.8 C3.4 6.4 3.4 6.4 2.8 8.6 C2.7 8.9 2.7 8.9 2.4 10.4 C3 11 3.7 11.7 4.4 12.4 C4.9 16 4.6 18 2.9 21.2 C0.1 23.6 -0.5 23.7 -3.9 23.7 C-5.1 23.6 -6.4 23.5 -7.6 23.4 C-7.6 22.4 -7.6 21.4 -7.6 20.4 C-5 20.4 -2.3 20.4 0.4 20.4 C0.4 18.4 0.4 16.4 0.4 14.4 C-0.9 13.7 -2.3 13.1 -3.6 12.4 C-3.6 11.7 -3.6 11.1 -3.6 10.4 C-2.6 9.7 -1.6 9.1 -0.6 8.4 C-0.8 7.6 -0.8 7.6 -1.6 3.4 C-3.6 3.7 -5.6 4 -7.6 4.4 C-6.8 2.4 -6.8 2.4 -5.6 0.4 C-2.6 -0.6 -2.6 -0.6 0 0 Z };
\path[fill=black!100,even odd rule,shift={(154,155)}] svg {M0 0 C-1.8 1.2 -3.7 2.4 -5.5 3.6 C-8.7 6.7 -10.1 9.7 -11 14 C-10.7 17.5 -10.7 17.5 -10 20 C-8 20.7 -6 21.3 -4 22 C-4.3 22.7 -4.7 23.3 -5 24 C-11.4 23.5 -11.4 23.5 -13.9 21.4 C-15.8 17.3 -15.8 13.3 -14.3 9.1 C-11.3 3.3 -6.8 -0.8 0 0 Z };
\path[fill=black!100,even odd rule,shift={(20,123)}] svg {M0 0 C4 0 7.9 0 12 0 C12 0.3 12 0.7 12 1 C10 1.3 8 1.7 6 2 C5.7 8.6 5.3 15.2 5 22 C7 22 9 22 11 22 C11 22.3 11 22.7 11 23 C7.4 23 3.7 23 0 23 C0.3 22.3 0.7 21.7 1 21 C1.1 19.6 1.1 18.1 1.1 16.7 C1.1 15.9 1.1 15 1.1 14.2 C1.1 13.3 1.1 12.4 1.1 11.5 C1.1 10.6 1.1 9.7 1.1 8.8 C1.1 8 1.1 7.1 1.1 6.3 C1.1 5.5 1.1 4.7 1.1 3.9 C1 2 1 2 0 0 Z };
\path[fill=black!100,even odd rule,shift={(211,10)}] svg {M0 0 C0.7 0.3 1.3 0.7 2 1 C1.5 1.6 1.1 2.2 0.6 2.8 C-3.2 8.1 -2.5 14.8 -2 21 C-0.5 24.6 -0.5 24.6 1 27 C1 28 1 29 1 30 C-6 23.6 -6 23.6 -7 18 C-7.1 11.5 -6.2 7 -2 2 C-1.3 1.3 -0.7 0.7 0 0 Z };
\path[fill=black!100,even odd rule,shift={(231,10)}] svg {M0 0 C4.3 1.5 5.9 4.1 8 8 C9.5 13.4 9.3 18 7 23.2 C5.4 26 3.4 27.9 1 30 C1 27 1 27 2.5 24.6 C4.9 18.9 5 12.6 3.4 6.7 C2.6 4.9 2.6 4.9 1 2 C0.7 1.3 0.3 0.7 0 0 Z };
\path[fill=black!100,even odd rule,shift={(154,155)}] svg {M0 0 C2.4 0.3 2.4 0.3 5 1 C7.7 5 7.7 8.2 7 13 C4.7 17.4 2 20.1 -2 23 C-2.7 22.7 -3.3 22.3 -4 22 C-3.6 21.4 -3.2 20.9 -2.7 20.3 C1 15.1 3.1 11.4 3 5 C1.6 3.1 1.6 3.1 0 2 C0 1.3 0 0.7 0 0 Z };
\path[fill=black!100,even odd rule,shift={(111,191)}] svg {M0 0 C0 0.6 0.1 1.3 0.1 1.9 C0.2 4.7 0.3 7.5 0.4 10.4 C0.5 11.4 0.5 12.3 0.6 13.4 C0.6 14.3 0.6 15.2 0.7 16.2 C0.7 17.1 0.8 17.9 0.8 18.8 C1 21 1 21 2 23 C-0.3 23 -2.6 23 -5 23 C-5 22.3 -5 21.7 -5 21 C-4.3 21 -3.7 21 -3 21 C-3.1 18 -3.2 15.1 -3.4 12.1 C-3.4 11.7 -3.4 11.7 -3.5 9.6 C-3.5 8.8 -3.5 8 -3.6 7.1 C-3.6 6.4 -3.6 5.6 -3.7 4.9 C-4 3 -4 3 -6 1 C-4 0 -4 0 0 0 Z };
\path[fill=black!100,even odd rule,shift={(155,109)}] svg {M0 0 C1 0.3 2 0.7 3 1 C3 1.6 3 2.2 3 2.8 C2.9 5.5 2.9 8.2 2.9 10.9 C2.9 11.8 2.8 12.7 2.8 13.7 C2.8 14.6 2.8 15.5 2.8 16.4 C2.8 17.3 2.8 18.1 2.8 18.9 C3 21 3 21 5 23 C2.4 23 -0.3 23 -3 23 C-2.3 22.7 -1.7 22.3 -1 22 C-1 15.7 -1 9.5 -1 3 C-1.7 2.7 -2.3 2.3 -3 2 C-2 1.3 -1 0.7 0 0 Z };
\path[fill=black!100,even odd rule,shift={(231,189)}] svg {M0 0 C1 0.3 2 0.7 3 1 C3 7.9 3 14.9 3 22 C3.7 22.3 4.3 22.7 5 23 C2.7 23 0.4 23 -2 23 C-1.7 22 -1.3 21 -1 20 C-0.9 17.3 -0.8 14.6 -0.8 11.9 C-0.8 11.2 -0.8 10.5 -0.8 9.7 C-0.8 7.8 -0.9 5.9 -1 4 C-1.7 3.3 -2.3 2.7 -3 2 C-2 1.3 -1 0.7 0 0 Z };
\path[fill=black!100,even odd rule,shift={(32,123)}] svg {M0 0 C3.3 1.8 4.9 2.8 7 6 C7.6 10.5 7.5 14.7 5.8 19 C3.7 21.3 2 22.1 -1 23 C-0.5 22.3 -0 21.6 0.5 20.9 C2.3 17.5 2.4 15.3 2.3 11.5 C2.3 10.4 2.3 9.3 2.3 8.2 C2 5.1 1.3 2.8 0 0 Z };
\path[fill=black!100,even odd rule,shift={(437,175)}] svg {M0 0 C1 0.3 2 0.7 3 1 C3 1.7 3 2.3 3 3 C1.4 3.7 -0.3 4.3 -2 5 C-2 7.3 -2 9.6 -2 12 C-1 12 -0 12 1 12 C0.3 13.3 -0.3 14.6 -1 16 C-2.9 15.9 -2.9 15.9 -5 15 C-6.9 10.9 -7.2 7.4 -7 3 C-8 2.7 -9 2.3 -10 2 C-10 1.7 -10 1.3 -10 1 C-7.6 0.9 -7.6 0.9 -5 1 C-4.3 1.7 -3.7 2.3 -3 3 C-2 2 -1 1 0 0 Z };
\path[fill=black!100,even odd rule,shift={(46,121)}] svg {M0 0 C4.1 1.8 5.7 4.1 8 8 C9.9 13.6 9.8 19 7.2 24.4 C5.6 27.3 4 29.3 1 31 C1.3 30.3 1.5 29.7 1.8 29 C2.2 28 2.6 27 3 26 C3.3 25.4 3.6 24.8 3.8 24.2 C6.1 18.9 6.6 14.6 5 9 C3.5 5.9 1.9 2.9 0 0 Z };
\path[fill=black!100,even odd rule,shift={(217,18)}] svg {M0 0 C1.3 0 2.6 0 4 0 C4.3 0.7 4.7 1.3 5 2 C5.6 1.7 6.2 1.3 6.8 1 C9 0 9 0 12 0 C11 3 11 3 8.9 4.2 C8.3 4.5 7.7 4.7 7 5 C7.3 7.3 7.7 9.6 8 12 C8.7 12 9.3 12 10 12 C9.7 13 9.3 14 9 15 C5.7 14.4 4.1 13.6 2 11 C1.7 8.7 1.7 8.7 1.8 6.3 C1.8 5.5 1.9 4.7 1.9 3.9 C1.9 3.2 2 2.6 2 2 C1.3 2 0.7 2 0 2 C0 1.3 0 0.7 0 0 Z };
\path[fill=black!100,even odd rule,shift={(7,122)}] svg {M0 0 C0 3 0 3 -1.5 6.1 C-3.8 11.3 -4.1 16.9 -2.1 22.3 C-1.2 24.3 -0.1 26.1 1 28 C1.3 28.7 1.7 29.3 2 30 C-2.3 28.2 -3.3 26.1 -5.2 22.1 C-7.4 16.4 -7 12.7 -5 7.1 C-3.7 4.4 -2.1 2.2 0 0 Z };
\path[fill=black!100,even odd rule,shift={(85,200)}] svg {M0 0 C4.6 0 9.2 0 14 0 C14 0.7 14 1.3 14 2 C9.4 2 4.8 2 0 2 C0 1.3 0 0.7 0 0 Z };
\path[fill=black!100,even odd rule,shift={(429,186)}] svg {M0 0 C0 1 0 2 0 3 C-0.7 3.7 -1.3 4.3 -2 5 C-3 5 -4 5 -5 5 C-5 4 -5 3 -5 2 C-4.4 1.9 -3.7 1.7 -3.1 1.6 C-1 1 -1 1 0 0 Z };
\path[fill=black!100,even odd rule,shift={(218,29)}] svg {M0 0 C0 1 0 2 0 3 C-1.7 3.3 -3.3 3.7 -5 4 C-4.7 3 -4.3 2 -4 1 C-1.9 0.3 -1.9 0.3 0 0 Z };
\end{tikzpicture}
}
\QuestionEnd

\begin{QuestionSolution}{2009 年第 6 题}
\textbf{已知与目标}\quad 已知 \(f\) 的分段图形，要找出从零点累积积分 \(F(x)=\int_0^xf(t)\,\mathrm{d}t\) 的图形。先确定起点和各段的增减，再检查接合。

\textbf{方法与条件}\quad 在 \(f\) 连续的点，微积分基本定理给出 \(F'(x)=f(x)\)：\(f>0\) 时 \(F\) 增加，\(f<0\) 时 \(F\) 减少。积分还满足 \(F(0)=0\)，反向积分须变号。

\textbf{推导}\quad 由图中左段 \(f(x)=1\)，对 \(-1\le x<0\)，
\[
F(x)=-\int_x^0 1\,\mathrm{d}t=-(0-x)=x,
\]
故图形经过 \((-1,-1)\) 与 \((0,0)\)。在 \((0,1)\) 上 \(f<0\)，所以 \(F\) 递减；在 \((1,2)\) 上 \(f>0\)，所以 \(F\) 递增。于是 \(x=1\) 为极小值点，且切线斜率 \(F'(1)=f(1)=0\)。在 \((2,3)\) 上 \(f=0\)，累积积分不再变化，故图形为水平线段。设 \(|f|\le M\)，则 \(|F(x)-F(y)|=|\int_y^xf|\le M|x-y|\)，从而接合点也连续。

\textbf{结论与检查}\quad 这些条件对应 \(\boxed{\mathrm{D}}\)。零点处的折角可以来自 \(f\) 的跳跃，但 \(F\) 本身不能跳跃；这也是筛选图形的检查条件。
\end{QuestionSolution}

% Q 2023 2
\QuestionPage{3.2 微积分基本公式与基本定理}{选择题}{2023 年第 2 题}{1}{分段函数的原函数}{144 / 234}{***}
函数
\[f(x)=
\begin{cases}
\dfrac{1}{\sqrt{1+x^2}}, & x\le 0,\\[0.6em]
(x+1)\cos x, & x>0
\end{cases}
\]
的一个原函数为（ ）。
\twochoices{\(F(x)=
\begin{cases}
\ln\left(\sqrt{1+x^2}-x\right), & x\le 0,\\
(x+1)\cos x-\sin x, & x>0
\end{cases}
\)}{\(F(x)=
\begin{cases}
\ln\left(\sqrt{1+x^2}-x\right)+1, & x\le 0,\\
(x+1)\cos x-\sin x, & x>0
\end{cases}
\)}{\(F(x)=
\begin{cases}
\ln\left(\sqrt{1+x^2}+x\right), & x\le 0,\\
(x+1)\sin x+\cos x, & x>0
\end{cases}
\)}{\(F(x)=
\begin{cases}
\ln\left(\sqrt{1+x^2}+x\right)+1, & x\le 0,\\
(x+1)\sin x+\cos x, & x>0
\end{cases}
\)}
\QuestionEnd

\begin{QuestionSolution}{2023 年第 2 题}
\textbf{已知与目标}\quad 原函数的含义是处处满足 \(F'(x)=f(x)\)。需要同时检查两侧表达式及零点处是否可导。

\textbf{方法与条件}\quad 对候选式直接求导，利用可导必连续确定常数。\(\sqrt{1+x^2}>|x|\)，所以两个对数的真数都为正，求导合法。

\textbf{推导}\quad 对左段的加号形式，链式法则给出
\[
\frac{\mathrm{d}}{\mathrm{d}x}\ln(\sqrt{1+x^2}+x)
=\frac{\frac{x}{\sqrt{1+x^2}}+1}{\sqrt{1+x^2}+x}
=\frac{1}{\sqrt{1+x^2}}.
\]
减号形式同样求导得到 \(-\frac{1}{\sqrt{1+x^2}}\)，故 A、B 不符合。C、D 的右段导数为
\[
[(x+1)\sin x+\cos x]'=\sin x+(x+1)\cos x-\sin x=(x+1)\cos x.
\]
右极限为 \(1\)，而左段对数在零点为 \(0\)，因此需要加常数 \(1\)，只有 D 连续。此时左、右导数都为 \(1=f(0)\)。

\textbf{结论与检查}\quad 选 \(\boxed{\mathrm{D}}\)。除检查各段导数外，接点的连续与相同单侧导数保证了零点也满足原函数定义。
\end{QuestionSolution}

% Q 2016 12
\QuestionPage{3.2 微积分基本公式与基本定理}{填空题}{2016 年第 12 题}{2}{变上限积分与高阶导数}{146 / 234}{***}
已知函数$f(x)$在$(-\infty,+\infty)$上连续，且$f(x) = (x + 1)^2 + 2\int_0^x f(t) \,\mathrm{d} t$，
则当$n \ge 2$时，$f^{(n)}(0) =$ \blank{}．
\QuestionEnd

\begin{QuestionSolution}{2016 年第 12 题}
\textbf{已知与目标}\quad \(f\) 起初只假定连续，但满足给定积分关系。要求任意整数 \(n\ge2\) 时零点处的 \(n\) 阶导数。

\textbf{方法与条件}\quad 连续被积函数的变上限积分可导，原等式因此先保证 \(f\) 可导；所得新等式继续保证更高阶可导。找出相邻阶导数之间的递推关系后，可直接计算任意阶。

\textbf{推导}\quad 原式令 \(x=0\)，得 \(f(0)=1\)。求一次导数，
\[
f'(x)=2(x+1)+2f(x),\qquad f'(0)=2+2\times1=4.
\]
右边可导，继续得到
\[
f''(x)=2+2f'(x),\qquad f''(0)=2+2\times4=10.
\]
再求导时常数项消失，此后每增加一阶均多乘 \(2\)：\(f^{(n)}(x)=2f^{(n-1)}(x)\)（\(n\ge3\)）。从二阶开始重复 \(n-2\) 次，得到 \(f^{(n)}(0)=10\cdot2^{n-2}\)。

\textbf{结论与检查}\quad \(\boxed{f^{(n)}(0)=10\cdot2^{n-2},\ n\ge2}\)。代入 \(n=2\) 得 \(10\)，代入 \(n=3\) 得 \(20=2f''(0)\)，符合初值与递推。
\end{QuestionSolution}

% Q 2026 18
\QuestionPage{3.2 微积分基本公式与基本定理}{解答题}{2026 年第 18 题}{2}{变上限积分的导数与连续性}{151 / 234}{***}
（本题满分 12 分）
已知函数 \(g(x)\) 连续.设
\[f(x)=\int_0^{x^2} g(xt)\,\mathrm{d}t,
\]
求 \(f'(x)\) 的表达式，并判断 \(f'(x)\) 在 \(x=0\) 处的连续性.
\QuestionEnd

\begin{QuestionSolution}{2026 年第 18 题}
\textbf{已知与目标}\quad 只知 \(g\) 连续，求 \(f(x)=\int_0^{x^2}g(xt)\,\mathrm{d}t\) 的导数，并判断导数在零点是否连续。

\textbf{方法与条件}\quad 不能直接对 \(g(xt)\) 求导，因为题设没有保证 \(g'\) 存在。令 \(u=xt\) 把 \(x\) 移到积分上下限和外部因子中；零点单独用差商处理。有向积分指交换上下限时积分变号，因此此换元对负 \(x\) 同样有效。

\textbf{推导}\quad 当 \(x\ne0\)，\(\mathrm{d}t=\frac{1}{x}\mathrm{d}u\)，上限 \(t=x^2\) 变为 \(u=x^3\)，故 \(f(x)=\frac{1}{x}\int_0^{x^3}g(u)\,\mathrm{d}u\)。由连续性与乘积求导，
\[
f'(x)=3xg(x^3)-\frac1{x^2}\int_0^{x^3} g(u)\,\mathrm{d}u.
\]
这里上限的导数为 \(3x^2\)，与外面的 \(\frac{1}{x}\) 相乘产生 \(3x\)。由于 \(g\) 连续，可写 \(g(u)=g(0)+o(1)\)，其中 \(o(1)\) 表示趋于零的误差。积分误差的绝对值不超过 \(|x|^3\sup_{|u|\le|x|^3}|g(u)-g(0)|\)，故
\[
\int_0^{x^3}g(u)\,\mathrm{d}u=x^3[g(0)+o(1)],\qquad
f(x)=x^2[g(0)+o(1)].
\]
因 \(f(0)=0\)，\(f'(0)=\lim_{x\to0}\frac{f(x)}{x}=0\)。非零点导数式中的两项分别为 \(3x[g(0)+o(1)]\)、\(-x[g(0)+o(1)]\)，均趋于零。

\textbf{结论与检查}\quad \(\boxed{f'(x)=3xg(x^3)-\frac{1}{x^2}\int_0^{x^3}g(u)\,\mathrm{d}u\ (x\ne0),\quad f'(0)=0}\)，且 \(f'\) 在零点连续。若 \(g\equiv c\)，原式为 \(cx^2\)，导数式给出 \(3cx-cx=2cx\)，核对正确。
\end{QuestionSolution}

% Q 2016 21
\QuestionPage{3.2 微积分基本公式与基本定理}{解答题}{2016 年第 21 题}{4}{原函数、平均值与零点}{148 / 234}{****}
已知$f(x)$在$\left[0,\frac{3\pi}{2} \right]$上连续，
在$\left(0,\frac{3\pi}{2} \right)$内是函数$\frac{\cos x}{2x - 3\pi}$的一个原函数，且$f(0) = 0$．
\begin{enumerate}
\item 求$f(x)$在区间$\left[0,\frac{3\pi}{2} \right]$上的平均值；
\item 证明$f(x)$在区间$\left(0,\frac{3\pi}{2} \right)$内存在唯一零点．
\end{enumerate}
\QuestionEnd

\begin{QuestionSolution}{2016 年第 21 题}
\textbf{已知与目标}\quad 令 \(L=\frac{3\pi}{2}\)。已知 \(f\) 在 \([0,L]\) 连续，在内部满足 \(f'(x)=\frac{\cos x}{2(x-L)}\)，且 \(f(0)=0\)。先求积分平均值，再证明内部零点存在且唯一。

\textbf{方法与条件}\quad 平均值为 \(\frac{1}{L}\int_0^Lf\)。分部积分时取 \(x-L\) 为 \(1\) 的原函数，可消去 \(f'\) 的分母。导数条件只在开区间成立，先在 \([\varepsilon,b]\subset(0,L)\) 积分，再利用 \(f\) 连续取端点极限。

\textbf{推导}\quad （1）分部积分后，\(\varepsilon\to0^+\)、\(b\to L^-\)，得
\[
\begin{aligned}
\int_0^L f(x)\,\mathrm{d}x
&=\bigl[(x-L)f(x)\bigr]_0^L
-\lim_{b\to L^-}\int_0^b (x-L)f'(x)\,\mathrm{d}x\\
&=-\frac12\int_0^L\cos x\,\mathrm{d}x=\frac12 .
\end{aligned}
\]
边界项在零点因 \(f(0)=0\) 消失，在 \(L\) 因 \(f\) 有限且 \(x-L\to0\) 消失。最后使用 \(\sin L=-1\)、\(\sin0=0\)。平均值为 \(\frac{1}{L}\times\frac{1}{2}=\frac{1}{3\pi}\)。

（2）内部的分母恒负，\(\cos x\) 在 \((0,\frac{\pi}{2})\) 为正，在 \((\frac{\pi}{2},L)\) 为负，故 \(f\) 在前段严格递减、后段严格递增。由 \(f(0)=0\)，前段及 \(\frac{\pi}{2}\) 处均为负。又积分 \(\frac{1}{2}>0\)，若内部处处非正，则由连续性积分必非正，矛盾。因此存在内点 \(d\) 满足 \(f(d)>0\)，且必有 \(d>\frac{\pi}{2}\)。连续性保证 \((\frac{\pi}{2},d)\) 内有零点；后段严格递增保证至多一个，前段没有零点。

\textbf{结论与检查}\quad 平均值为 \(\boxed{\frac{1}{3\pi}}\)，且 \((0,\frac{3\pi}{2})\) 内恰有一个零点。端点 \(f(0)=0\) 未计入内部零点数，积分为正补足了仅凭单调性无法得到的存在性。
\end{QuestionSolution}

% Q 2018 16
\QuestionPage{3.2 微积分基本公式与基本定理}{解答题}{2018 年第 16 题}{2}{变上限积分与平均值}{149 / 234}{****}
已知连续函数$f(x)$满足$\int_0^x f(t)\,\mathrm{d} t + \int_0^x tf(x - t)\,\mathrm{d} t = ax^2$．
\begin{enumerate}
\item 求$f(x)$；
\item 若$f(x)$在区间$[0,1]$上的平均值为$1$，求$a$的值．
\end{enumerate}
\QuestionEnd

\begin{QuestionSolution}{2018 年第 16 题}
\textbf{已知与目标}\quad 连续函数 \(f\) 满足含 \(f(x-t)\) 的积分关系。第一问求 \(f\)，第二问利用 \([0,1]\) 平均值为 \(1\) 确定参数 \(a\)。

\textbf{方法与条件}\quad 连续性允许对变上限积分求导，却没有直接给出 \(f'\)。先把第二项换元，再以 \(F(x)=\int_0^xf(t)\,\mathrm{d}t\) 为未知函数求一阶方程。积分因子是使等式左边成为乘积导数的乘数。

\textbf{推导}\quad （1）第二项令 \(u=x-t\)，上下限由 \(0,x\) 变为 \(x,0\)，且 \(\mathrm{d}t=-\mathrm{d}u\)，故它等于 \(\int_0^x(x-u)f(u)\,\mathrm{d}u\)。写为 \(xF(x)-\int_0^xu f(u)\,\mathrm{d}u\)，求导得到 \(F(x)+xf(x)-xf(x)=F(x)\)。于是
\[
F'(x)+F(x)=2ax,\qquad F(0)=0.
\]
乘 \(\mathrm{e}^x\)，左边为 \((\mathrm{e}^xF)'\)。从 \(0\) 积分至 \(x\)，利用 \(\int t\mathrm{e}^t\,\mathrm{d}t=(t-1)\mathrm{e}^t\)，得
\[
\mathrm{e}^xF(x)=2a\bigl[(x-1)\mathrm{e}^x+1\bigr].
\]
所以 \(F(x)=2a(x-1+\mathrm{e}^{-x})\)，求导得 \(f(x)=2a(1-\mathrm{e}^{-x})\)。

（2）区间长度为 \(1\)，平均值为
\[
\int_0^1 2a(1-\mathrm{e}^{-x})\,\mathrm{d}x
=2a[x+\mathrm{e}^{-x}]_0^1
=2a(1+\mathrm{e}^{-1}-1)=\frac{2a}{\mathrm{e}}.
\]
令其等于 \(1\)，得 \(a=\frac{\mathrm{e}}{2}\)。

\textbf{结论与检查}\quad \(\boxed{f(x)=2a(1-\mathrm{e}^{-x})}\)，平均值条件给出 \(\boxed{a=\frac{\mathrm{e}}{2}}\)。所得函数使原式两边导数相等，且两边在零点同为零，故原积分等式也成立，没有因求导引入多余解。
\end{QuestionSolution}

% Q 2020 16
\QuestionPage{3.2 微积分基本公式与基本定理}{解答题}{2020 年第 16 题}{2}{变上限积分求导与连续性}{150 / 234}{****}
已知函数 \(f\left( x \right)\) 连续且 \(\lim_{x\to 0}{\dfrac{f\left( x \right)}{x}} = 1\)，\(g\left( x \right)=\int_{0}^{1}{f\left( xt \right) \,\mathrm{d} t}\)，求 \(g^{\prime}\left( x \right)\)，并证明 \(g^{\prime}\left( x \right)\) 在 \(x=0\) 处连续。
\QuestionEnd

\begin{QuestionSolution}{2020 年第 16 题}
\textbf{已知与目标}\quad \(f\) 连续且 \(\frac{f(x)}{x}\to1\)。求 \(g\) 在所有点的导数，并证明 \(g'\) 在零点的极限等于 \(g'(0)\)。

\textbf{方法与条件}\quad 不假定 \(f\) 可导，先令 \(u=xt\) 化为变上限积分；\(x=0\) 时换元公式不能除以 \(x\)，须单独用导数定义。记号 \(o(x)\) 表示除以 \(x\) 后趋于零的误差。

\textbf{推导}\quad 极限条件使 \(f(x)\to0\)，由连续性 \(f(0)=0\)，所以 \(g(0)=0\)。当 \(x\ne0\)，换元及商的求导法则给出
\[
g(x)=\frac1x\int_0^x f(u)\,\mathrm{d}u,\qquad
g'(x)=\frac{xf(x)-\int_0^x f(u)\,\mathrm{d}u}{x^2}.
\]
为处理零点，令 \(r(u)=f(u)-u\)，则 \(r(u)=o(u)\)。任给 \(\varepsilon>0\)，充分靠近零点时 \(|r(u)|\le\varepsilon|u|\)，故对正、负 \(x\) 都有
\[
\left|\int_0^xr(u)\,\mathrm{d}u\right|\le\frac{\varepsilon x^2}{2}.
\]
因此 \(\int_0^xf(u)\,\mathrm{d}u=\frac{x^2}{2}+o(x^2)\)。于是
\[
g'(0)=\lim_{x\to0}\frac{g(x)-g(0)}{x}=\frac12 .
\]
再把 \(xf(x)=x^2+o(x^2)\) 和积分展开代入非零点的公式，得 \(g'(x)=\frac{1}{2}+o(1)\to\frac{1}{2}\)。

\textbf{结论与检查}\quad \(\boxed{g'(x)=\frac{xf(x)-\int_0^xf(u)\,\mathrm{d}u}{x^2}\ (x\ne0),\quad g'(0)=\frac{1}{2}}\)，且 \(g'\) 在零点连续。误差估计对两侧都成立，因而证明了双侧连续性。
\end{QuestionSolution}

\SectionHeading{3.3}{一元函数积分学及其应用}{两种基本积分法}{定理与要点\quad 教材 PDF 第 214 页起\hfill 本节 12 题}
\pdfbookmark[2]{3.3 两种基本积分法}{section-3.3}
\Theory{定理 3.1}{214}{不定积分换元法则 I}{辨认被积式中的 \(f(\varphi(x))\varphi'(x)\) 后，令 \(u=\varphi(x)\)，对 \(f(u)\) 积分后代回 \(u=\varphi(x)\)。}
\Theory{定理 3.2}{219}{不定积分换元法则 II}{令 \(x=\varphi(t)\)，在 \(\varphi'\) 定号且可反解的区间内，将被积式化为 \(f(\varphi(t))\varphi'(t)\,\mathrm{d}t\)。}
\Theory{定理 3.3}{222}{定积分换元}{连续被积函数和连续可导换元满足教材的区间条件时，上下限随代换改变，微分乘以导数。}
\Theory{公式}{225、229}{分部积分}{\(\int u\,\mathrm{d}v=uv-\int v\,\mathrm{d}u\)。定积分中的乘积 \(uv\) 取上、下限处的差。}
% Q 2011 6
\QuestionPage{3.3 两种基本积分法}{选择题}{2011 年第 6 题}{1}{定积分的对称换元}{152 / 234}{**}
设$I = \int_{0}^{\frac{\pi}{4}}\ln\sin x\,\mathrm{d} x$，$J = \int_{0}^{\frac{\pi}{4}}\ln\cot x\,\mathrm{d} x$，
$K = \int_{0}^{\frac{\pi}{4}}\ln\cos x\,\mathrm{d} x$，则$I,\;J,\;K$的大小关系是\pickout{}
\begin{abcd}
\item $I < J < K$．
\item $I < K < J$．
\item $J < I < K$．
\item $K < J < I$．
\end{abcd}
\QuestionEnd

\begin{QuestionSolution}{2011 年第 6 题}
\textbf{已知与目标}\quad 三个积分区间相同，比较 \(I,J,K\) 时可先观察三个被积函数的符号和大小。

\textbf{方法与条件}\quad 对数在正数上严格递增，故可先比较真数。\(\ln\sin x\) 和 \(\ln\cot x\) 在零点无界，须确认反常积分收敛，即截去端点后取极限得到有限值，才能比较积分。

\textbf{推导}\quad 在 \(0<x<\frac{\pi}{4}\) 上，\(0<\sin x<\cos x<1\)，\(\cot x>1\)，所以
\[
\ln\sin x<\ln\cos x<0<\ln\cot x.
\]
又 \(\ln\sin x=\ln x+\ln\frac{\sin x}{x}\)，第二项趋于 \(0\)，而 \(\int_0^a\ln x\,\mathrm{d}x=[x\ln x-x]_{0^+}^a\) 有限，故 \(I\) 收敛。\(K\) 的被积函数连续，\(\ln\cot x=\ln\cos x-\ln\sin x\)，因此 \(J\) 也收敛。对严格不等式积分得到 \(I<K<J\)。

\textbf{结论与检查}\quad 选 \(\boxed{\mathrm{B}}\)。\(I,K\) 均为负数，\(J\) 为正数，与次序一致。
\end{QuestionSolution}

% Q 2020 3
\QuestionPage{3.3 两种基本积分法}{选择题}{2020 年第 3 题}{1}{反三角函数的换元积分}{153 / 234}{**}
\(\int_{0}^{1}{\dfrac{\arcsin \sqrt{x}}{\sqrt{x\left( 1-x \right)}} \,\mathrm{d} x} =\)
\choices{\(\dfrac{\pi^{2}}{4}\)}{\(\dfrac{\pi^{2}}{8}\)}{\(\dfrac{\pi}{4}\)}{\(\dfrac{\pi}{8}\)}
\QuestionEnd

\begin{QuestionSolution}{2020 年第 3 题}
\textbf{已知与目标}\quad 计算含 \(\arcsin\sqrt{x}\) 和根式 \(\sqrt{x(1-x)}\) 的积分，积分端点处分母为零。

\textbf{方法与条件}\quad 令 \(u=\arcsin\sqrt{x}\)，即 \(x=\sin^2u\)。在 \(0<u<\frac{\pi}{2}\) 内正弦、余弦均为正，根式可直接写成乘积。先在内部换元，再取两端极限。

\textbf{推导}\quad \(x\) 由 \(0\) 增至 \(1\)，\(u\) 由 \(0\) 增至 \(\frac{\pi}{2}\)。又
\[
\mathrm{d}x=2\sin u\cos u\,\mathrm{d}u,\qquad
\sqrt{x(1-x)}=\sin u\cos u.
\]
约去相同因子后原积分变为
\[
2\int_0^{\frac{\pi}{2}}u\,\mathrm{d}u
=[u^2]_0^{\frac{\pi}{2}}=\frac{\pi^2}{4}.
\]
\textbf{结论与检查}\quad 积分收敛且等于 \(\boxed{\frac{\pi^2}{4}}\)，选 A。换元后的函数 \(2u\) 在闭区间连续，说明原积分端点虽有根式分母，极限仍有限。
\end{QuestionSolution}

% Q 2017 11
\QuestionPage{3.3 两种基本积分法}{填空题}{2017 年第 11 题}{2}{定积分计算}{155 / 234}{**}
$\int_0^{+\infty} \frac{\ln(1 + x)}{(1 + x)^2}\,\mathrm{d} x =$ \blank{}．
\QuestionEnd

\begin{QuestionSolution}{2017 年第 11 题}
\textbf{已知与目标}\quad 积分区间无上界，要计算截断积分在 \(R\to+\infty\) 时的极限。

\textbf{方法与条件}\quad 对 \(\ln(1+x)\) 分部积分。取 \(u=\ln(1+x)\)、\(\mathrm{d}v=\frac{1}{(1+x)^2}\,\mathrm{d}x\)，则 \(\mathrm{d}u=\frac{1}{1+x}\,\mathrm{d}x\)、\(v=-\frac{1}{1+x}\)。

\textbf{推导}\quad 先在有限区间上计算，
\[
\int_0^R\frac{\ln(1+x)}{(1+x)^2}\,\mathrm{d}x
=\left[-\frac{\ln(1+x)}{1+x}\right]_0^R
+\int_0^R\frac{\mathrm{d}x}{(1+x)^2}.
\]
右边等于 \(-\frac{\ln(1+R)}{1+R}+1-\frac{1}{1+R}\)。当 \(R\to+\infty\)，\(\frac{\ln(1+R)}{1+R}\to0\)，\(\frac{1}{1+R}\to0\)，所以极限为 \(1\)。

\textbf{结论与检查}\quad 反常积分收敛，值为 \(\boxed{1}\)。被积函数在 \(x>0\) 为正，结果也为正；下限处 \(\ln1=0\)，边界项没有遗漏常数。
\end{QuestionSolution}

% Q 2018 11
\QuestionPage{3.3 两种基本积分法}{填空题}{2018 年第 11 题}{2}{有理函数的积分}{157 / 234}{**}
$\int_5^{+\infty} \frac{1}{x^2 - 4x + 3}\,\mathrm{d} x =$ \blank{}．
\QuestionEnd

\begin{QuestionSolution}{2018 年第 11 题}
\textbf{已知与目标}\quad 求 \([5,+\infty)\) 上的反常积分，分母在该区间内无零点。

\textbf{方法与条件}\quad 因式分解后拆成部分分式，即简单分式的和。先在 \([5,R]\) 上积分，并将对数合并后取极限，避免把两个分别发散的量作差。

\textbf{推导}\quad \(x^2-4x+3=(x-1)(x-3)\)，并且
\[
\frac{1}{(x-1)(x-3)}=\frac{1}{2}\left(\frac{1}{x-3}-\frac{1}{x-1}\right).
\]
因 \(x\ge5\) 时两因子都为正，积分得
\[
\int_5^R \frac{\mathrm{d}x}{x^2-4x+3}
=\frac12\left[\ln\frac{x-3}{x-1}\right]_5^R .
\]
当 \(R\to+\infty\)，\(\frac{R-3}{R-1}\to1\)，上端对数趋于 \(0\)；下端为 \(\ln\frac{5-3}{5-1}=\ln\frac{1}{2}=-\ln2\)。所以结果为 \(\frac{1}{2}[0-(-\ln2)]\)。

\textbf{结论与检查}\quad 积分收敛且为 \(\boxed{\frac{1}{2}\ln2}\)。原函数的导数为 \(\frac{1}{2}(\frac{1}{x-3}-\frac{1}{x-1})\)，合并后正好回到原被积函数。
\end{QuestionSolution}

% Q 2022 13
\QuestionPage{3.3 两种基本积分法}{填空题}{2022 年第 13 题}{2}{有理函数的定积分}{158 / 234}{**}
\(\int_{0}^{1} \frac{2x + 3}{x^{2} - x + 1} \,\mathrm{d}x = \blank{}\)
\QuestionEnd

\begin{QuestionSolution}{2022 年第 13 题}
\textbf{已知与目标}\quad 计算分母为二次式的有理函数积分。分子与分母导数只差一个常数。

\textbf{方法与条件}\quad 分母 \(q(x)=x^2-x+1=(x-\frac{1}{2})^2+\frac{3}{4}>0\)，所以对数与反正切积分均适用。将分子写成 \(q'(x)+4\)。

\textbf{推导}\quad \(2x+3=(2x-1)+4\)，第一部分为 \(\int_0^1\frac{q'}{q}\,\mathrm{d}x=[\ln q]_0^1=\ln1-\ln1=0\)。对剩余部分用配方，得到
\[
4\int_0^1\frac{\mathrm{d}x}{x^2-x+1}
=\frac8{\sqrt3}\left[\arctan\frac{2x-1}{\sqrt3}\right]_0^1
=\boxed{\frac{8\pi}{3\sqrt3}}.
\]
具体代入反正切，\(\arctan\frac{1}{\sqrt3}=\frac{\pi}{6}\)，\(\arctan\frac{-1}{\sqrt3}=-\frac{\pi}{6}\)，差为 \(\frac{\pi}{3}\)。

\textbf{结论与检查}\quad 积分为 \(\boxed{\frac{8\pi}{3\sqrt3}}\)。\(\frac{2}{\sqrt3}\arctan\frac{2x-1}{\sqrt3}\) 的导数为 \(\frac{1}{q(x)}\)，因此配方后的系数正确。
\end{QuestionSolution}

% Q 2009 11
\QuestionPage{3.3 两种基本积分法}{填空题}{2009 年第 11 题}{2}{含参数积分与分部积分}{154 / 234}{***}
$\lim_{n \to \infty} \int_{0}^{1} \mathrm{e}^{- x}\sin nx\,\mathrm{d} x =$ \blank{}．
\QuestionEnd

\begin{QuestionSolution}{2009 年第 11 题}
\textbf{已知与目标}\quad 求 \(I_n=\int_0^1\mathrm{e}^{-x}\sin(nx)\,\mathrm{d}x\) 在 \(n\to\infty\) 时的极限。\(\sin(nx)\) 本身没有逐点极限，需估计整个积分。

\textbf{方法与条件}\quad 分部积分取 \(u=\mathrm{e}^{-x}\)、\(\mathrm{d}v=\sin(nx)\,\mathrm{d}x\)，从正弦的原函数中提出 \(\frac{1}{n}\)，再用 \(|\cos(nx)|\le1\) 给出与 \(n\) 无关的上界。

\textbf{推导}\quad 有 \(\mathrm{d}u=-\mathrm{e}^{-x}\,\mathrm{d}x\)、\(v=-\frac{\cos(nx)}{n}\)，所以
\[
I_n=\frac{1-\mathrm{e}^{-1}\cos n}{n}
-\frac1n\int_0^1 \mathrm{e}^{-x}\cos(nx)\,\mathrm{d}x.
\]
边界项的绝对值不超过 \(\frac{1+\mathrm{e}^{-1}}{n}\)；余下积分的绝对值不超过 \(\int_0^1\mathrm{e}^{-x}\,\mathrm{d}x=1-\mathrm{e}^{-1}\)。因此
\[
|I_n|\le\frac{1+\mathrm{e}^{-1}}{n}+\frac{1-\mathrm{e}^{-1}}{n}=\frac{2}{n}\longrightarrow0.
\]
\textbf{结论与检查}\quad 所求极限为 \(\boxed{0}\)。估计控制的是绝对值，因此同时排除了正、负方向的振荡，无需交换积分与极限。
\end{QuestionSolution}

% Q 2017 13
\QuestionPage{3.3 两种基本积分法}{填空题}{2017 年第 13 题}{2}{三角函数的定积分}{156 / 234}{***}
$\int_0^1 \,\mathrm{d} y \int_y^1 \frac{\tan x}{x}\,\mathrm{d} x =$ \blank{}．
\QuestionEnd

\begin{QuestionSolution}{2017 年第 13 题}
\textbf{已知与目标}\quad 计算先对 \(x\)、再对 \(y\) 的累次积分。积分范围是 \(0\le y\le x\le1\)。

\textbf{方法与条件}\quad 改变积分次序可先积掉不出现在被积函数中的 \(\!y\)。因 \(\lim_{x\to0}\frac{\tan x}{x}=1\)，在零点补值 \(1\) 后，被积函数在积分区域连续，允许交换积分次序。

\textbf{推导}\quad 固定 \(x\in[0,1]\) 后，\(y\) 从 \(0\) 到 \(x\)，因此
\[
\int_0^1\mathrm{d}y\int_y^1\frac{\tan x}{x}\,\mathrm{d}x
=\int_0^1\frac{\tan x}{x}\left(\int_0^x\mathrm{d}y\right)\mathrm{d}x
=\int_0^1\tan x\,\mathrm{d}x.
\]
内层积分 \(\int_0^x1\,\mathrm{d}y=x\) 正好约去分母。由于 \(0\le x\le1<\frac{\pi}{2}\)，\(\cos x>0\)，可用 \(\tan x=-(\ln\cos x)'\)，得到 \([-\ln\cos x]_0^1=-\ln\cos1\)。

\textbf{结论与检查}\quad \(\boxed{-\ln\cos1}\)。\(0<\cos1<1\)，所以答案为正，与原被积函数为正一致。
\end{QuestionSolution}

% Q 2017 17
\QuestionPage{3.3 两种基本积分法}{解答题}{2017 年第 17 题}{3}{和式极限与定积分}{160 / 234}{***}
求$\lim_{n \to \infty} \sum_{k = 1}^n \frac{k}{n^2} \ln \left(1 + \frac{k}{n} \right)$．
\QuestionEnd

\begin{QuestionSolution}{2017 年第 17 题}
\textbf{已知与目标}\quad 求含 \(\frac{k}{n}\) 的加权对数和式极限，可把 \(\frac{k}{n^2}\) 拆成区间宽度与取样点的乘积。

\textbf{方法与条件}\quad \(h(x)=x\ln(1+x)\) 在 \([0,1]\) 连续，其 Riemann 和，即函数值乘小区间宽度的和，趋于定积分。计算该积分时对对数分部积分。

\textbf{推导}\quad 原和式为 \(\sum_{k=1}^n\frac{1}{n}\frac{k}{n}\ln(1+\frac{k}{n})\)，故极限为 \(\int_0^1x\ln(1+x)\,\mathrm{d}x\)。取 \(u=\ln(1+x)\)、\(\mathrm{d}v=x\,\mathrm{d}x\)，并利用 \(\frac{x^2}{1+x}=x-1+\frac{1}{1+x}\)，得
\[
\begin{aligned}
\int_0^1 x\ln(1+x)\,\mathrm{d}x
&=\frac12\ln2-\frac12\int_0^1\frac{x^2}{1+x}\,\mathrm{d}x\\
&=\frac12\ln2-\frac12\int_0^1\left(x-1+\frac1{1+x}\right)\mathrm{d}x
=\boxed{\frac14}.
\end{aligned}
\]
其中最后一步使用 \(\int_0^1(x-1+\frac{1}{1+x})\,\mathrm{d}x=\frac{1}{2}-1+\ln2\)，对数项相消。

\textbf{结论与检查}\quad 极限为 \(\boxed{\frac{1}{4}}\)。被积函数非负，答案为正；宽度 \(\frac{1}{n}\) 乘 \(n\) 项覆盖的总区间长度正好为 \(1\)。
\end{QuestionSolution}

% Q 2018 15
\QuestionPage{3.3 两种基本积分法}{解答题}{2018 年第 15 题}{2}{换元与分部积分}{161 / 234}{***}
求不定积分$\int \mathrm{e}^{2x}\arctan \sqrt{\mathrm{e}^x - 1} \,\mathrm{d} x$．
\QuestionEnd

\begin{QuestionSolution}{2018 年第 15 题}
\textbf{已知与目标}\quad 求 \(\mathrm{e}^{2x}\arctan\sqrt{\mathrm{e}^x-1}\) 的原函数。实数定义域为 \(x\ge0\)，通常在开区间 \(x>0\) 上写原函数。

\textbf{方法与条件}\quad 反正切求导能消去部分指数因子，故先分部积分；剩余根式积分令 \(u=\sqrt{\mathrm{e}^x-1}\) 消根。

\textbf{推导}\quad 对 \(x>0\)，链式求导得
\[
\left(\arctan\sqrt{\mathrm{e}^x-1}\right)'
=\frac{1}{1+(\mathrm{e}^x-1)}\frac{\mathrm{e}^x}{2\sqrt{\mathrm{e}^x-1}}
=\frac{1}{2\sqrt{\mathrm{e}^x-1}}.
\]
以 \(\frac{1}{2}\mathrm{e}^{2x}\) 为指数项的原函数，分部积分得
\[
\int \mathrm{e}^{2x}\arctan\sqrt{\mathrm{e}^x-1}\,\mathrm{d}x
=\frac{\mathrm{e}^{2x}}2\arctan\sqrt{\mathrm{e}^x-1}
-\frac14\int\frac{\mathrm{e}^{2x}}{\sqrt{\mathrm{e}^x-1}}\,\mathrm{d}x .
\]
令 \(u=\sqrt{\mathrm{e}^x-1}\)，则 \(\mathrm{e}^x=1+u^2\)、\(\mathrm{e}^x\mathrm{d}x=2u\,\mathrm{d}u\)，所以末积分为 \(2\int(1+u^2)\,\mathrm{d}u=2u+\frac{2}{3}u^3\)。代回得到
\[
\frac{\mathrm{e}^{2x}}2\arctan u-\frac u2-\frac{u^3}{6}+C,\qquad
u=\sqrt{\mathrm{e}^x-1},\quad x>0.
\]
\textbf{结论与检查}\quad 上式给出所求原函数，\(C\) 任意。求导时首项多出的 \(\frac{\mathrm{e}^{2x}}{4u}\)，与 \(-\frac{u}{2}-\frac{u^3}{6}\) 的导数 \(-\frac{1+u^2}{2}u'=-\frac{\mathrm{e}^{2x}}{4u}\) 相消，回到原被积函数。
\end{QuestionSolution}

% Q 2019 16
\QuestionPage{3.3 两种基本积分法}{解答题}{2019 年第 16 题}{2}{不定积分计算}{162 / 234}{***}
求不定积分$\int \frac{3x + 6}{(x - 1)^2(x^2 + x + 1)}\,\mathrm{d} x$．
\QuestionEnd

\begin{QuestionSolution}{2019 年第 16 题}
\textbf{已知与目标}\quad 求有理函数的原函数，分母含二重因子 \((x-1)^2\) 和二次因子 \(q(x)=x^2+x+1\)，定义域排除 \(x=1\)。

\textbf{方法与条件}\quad 部分分式分解要同时保留 \(\frac{1}{x-1}\) 和 \(\frac{1}{(x-1)^2}\) 两项。又 \(q(x)=(x+\frac{1}{2})^2+\frac{3}{4}>0\)，所以 \(\ln q\) 无需绝对值。

\textbf{推导}\quad 设分解为 \(\frac{A}{(x-1)^2}+\frac{B}{x-1}+\frac{Cx+D}{q(x)}\)，清分母后
\[
3x+6=Aq(x)+B(x-1)q(x)+(Cx+D)(x-1)^2.
\]
令 \(x=1\)，得 \(9=3A\)，故 \(A=3\)。两侧求导后再令 \(x=1\)，得 \(3=3A+3B\)，故 \(B=-2\)。比较三次项，\(B+C=0\)，得 \(C=2\)。令 \(x=0\)，\(6=A-B+D\)，得 \(D=1\)。于是
\[
\frac{3x+6}{(x-1)^2(x^2+x+1)}
=\frac3{(x-1)^2}-\frac2{x-1}+\frac{2x+1}{x^2+x+1}.
\]
逐项积分：第一项为 \(-\frac{3}{x-1}\)，第二项为 \(-2\ln|x-1|\)，第三项因分子等于 \(q'\)，积分为 \(\ln q\)。

\textbf{结论与检查}\quad \(\boxed{-\frac{3}{x-1}-2\ln|x-1|+\ln(x^2+x+1)+C}\)，适用于 \(x<1\) 或 \(x>1\)，两个区间的常数可分别选择。求导便得到上面的部分分式，通分可复原原式。
\end{QuestionSolution}

% Q 2025 17
\QuestionPage{3.3 两种基本积分法}{解答题}{2025 年第 17 题}{3}{定积分计算}{163 / 234}{***}
（本题满分 10 分）
计算 \(\int_{0}^{1} \frac{1}{\left( x + 1 \right)\left( x^{2} - 2x + 2 \right)} \,\mathrm{d} x\).
\QuestionEnd

\begin{QuestionSolution}{2025 年第 17 题}
\textbf{已知与目标}\quad 计算 \([0,1]\) 上的有理函数积分。两个分母因子 \(x+1\) 与 \(q(x)=x^2-2x+2=(x-1)^2+1\) 均为正。

\textbf{方法与条件}\quad 先拆部分分式，再把二次项的分子分解为分母导数和常数，分别产生对数与反正切。

\textbf{推导}\quad 设 \(\frac{1}{(x+1)q(x)}=\frac{A}{x+1}+\frac{Bx+C}{q(x)}\)。清分母比较系数，得到 \(A+B=0\)、\(-2A+B+C=0\)、\(2A+C=1\)，解得 \(A=\frac{1}{5}\)、\(B=-\frac{1}{5}\)、\(C=\frac{3}{5}\)。所以
\[
\frac1{(x+1)q(x)}
=\frac15\left(\frac1{x+1}+\frac{-x+3}{q(x)}\right).
\]
因 \(-x+3=-\frac{1}{2}q'(x)+2\)，原积分为
\[
\frac15\left[\ln(x+1)-\frac12\ln q(x)+2\arctan(x-1)\right]_0^1
=\boxed{\frac{3\ln2+\pi}{10}}.
\]
上限代入得 \(\ln2\)，下限代入得 \(-\frac{1}{2}\ln2-\frac{\pi}{2}\)，相减后乘 \(\frac{1}{5}\)，即上式。

\textbf{结论与检查}\quad \(\boxed{\frac{3\ln2+\pi}{10}}\)。对括号内原函数求导得到 \(\frac{1}{x+1}+\frac{-x+3}{q(x)}\)，与部分分式一致；积分值为正，符合被积函数的符号。
\end{QuestionSolution}

% Q 2009 16
\QuestionPage{3.3 两种基本积分法}{解答题}{2009 年第 16 题}{3}{对数函数的不定积分}{159 / 234}{****}
计算不定积分$\int \ln \left(1 + \sqrt{\frac{1 + x}{x}} \right)\,\mathrm{d} x$（$x > 0$）．
\QuestionEnd

\begin{QuestionSolution}{2009 年第 16 题}
\textbf{已知与目标}\quad 在 \(x>0\) 上求 \(L(x)=\ln(1+\sqrt{\frac{1+x}{x}})\) 的全部原函数。

\textbf{方法与条件}\quad 把根式通分后拆对数，再分部积分。因为 \(x>0\)，平方根和对数真数都为正，可写 \(L(x)=A(x)-\frac{1}{2}\ln x\)，其中 \(A(x)=\ln(\sqrt{x+1}+\sqrt{x})\)。

\textbf{推导}\quad 由链式法则，
\[
A'(x)=\frac{\frac{1}{2\sqrt{x+1}}+\frac{1}{2\sqrt{x}}}{\sqrt{x+1}+\sqrt{x}}
=\frac{1}{2\sqrt{x(x+1)}},\quad
L'(x)=\frac{1}{2\sqrt{x(x+1)}}-\frac{1}{2x}.
\]
取 \(u=L(x)\)、\(\mathrm{d}v=\mathrm{d}x\) 分部积分，得
\[
\int L(x)\,\mathrm{d}x
=xL(x)+\frac{x}{2}-\frac12\int\frac{x}{\sqrt{x(x+1)}}\,\mathrm{d}x .
\]
令 \(S(x)=\sqrt{x(x+1)}\)，则 \(S'(x)=\frac{2x+1}{2S(x)}\)，因此
\[
\frac{x}{S(x)}=S'(x)-\frac{1}{2S(x)}=S'(x)-A'(x).
\]
故剩余积分等于 \(S(x)-A(x)+C\)。代回得
\[
xL(x)+\frac{x}{2}-\frac12\sqrt{x(x+1)}
+\frac12\ln(\sqrt{x+1}+\sqrt{x})+C,\qquad x>0.
\]
\textbf{结论与检查}\quad 上式为 \(x>0\) 上的全部原函数，其中 \(L(x)\) 是原被积函数，\(C\) 为任意常数。求导时除 \(L\) 外的项合为 \(xL'+\frac{1}{2}-\frac{1}{2}(S'-A')=0\)，所以导数确为 \(L\)。
\end{QuestionSolution}

\SectionHeading{3.4}{一元函数积分学及其应用}{定积分的应用}{定理与要点\quad 教材 PDF 第 233 页起\hfill 本节 25 题}
\pdfbookmark[2]{3.4 定积分的应用}{section-3.4}
\Theory{方法}{234--235}{微元法}{选择积分变量，把一个小区间的贡献写成主部 \(\mathrm{d}Q\)，再对整个区间积分。}
\Theory{公式}{235--237}{平面面积}{直角坐标下，对上、下边界之差积分。极坐标下，扇形面积微元为 \(\mathrm{d}S=\frac{\rho^2}{2}\,\mathrm{d}\theta\)。}
\Theory{公式}{237--239}{旋转体积}{截面积沿轴积分。圆柱壳法用周长乘壳高再乘厚度。}
\Theory{方法}{239--242}{物理应用}{以薄层或小段建立功、压力等贡献的微元，并保留单位和力的作用方向。}
% Q 2010 12
\QuestionPage{3.4 定积分的应用}{填空题}{2010 年第 12 题}{2}{极坐标曲线弧长}{167 / 234}{*}
当$0 \le \theta \le \pi$时，对数螺线$r = \mathrm{e}^{\theta}$的弧长为 \blank{}．
\QuestionEnd

\begin{QuestionSolution}{2010 年第 12 题}
\textbf{已知与目标}\quad 求 \(r=\mathrm{e}^\theta\) 在 \(0\le\theta\le\pi\) 上的曲线长度。

\textbf{方法与条件}\quad 极坐标曲线可写成 \(x=r\cos\theta,y=r\sin\theta\)，由平面弧长公式得到微小弧长 \(\mathrm{d}s=\sqrt{r^2+(r')^2}\,\mathrm{d}\theta\)。本题 \(r\) 连续可导，可以直接使用。

\textbf{推导}\quad \(r'=\mathrm{e}^\theta\)，且指数函数为正，所以根式为 \(\sqrt2\mathrm{e}^\theta\)。于是
\[
s=\int_0^\pi\sqrt{\mathrm{e}^{2\theta}+\mathrm{e}^{2\theta}}\,\mathrm{d}\theta
=\sqrt2\int_0^\pi \mathrm{e}^\theta\,\mathrm{d}\theta
=\boxed{\sqrt2(\mathrm{e}^\pi-1)}.
\]
\textbf{结论与检查}\quad 弧长为 \(\boxed{\sqrt2(\mathrm{e}^\pi-1)}\)。弧长大于两端半径之差 \(\mathrm{e}^\pi-1\)，与曲线同时沿径向和角向变化一致。
\end{QuestionSolution}

% Q 2017 6
\QuestionPage{3.4 定积分的应用}{选择题}{2017 年第 6 题}{1}{速度图像与追及距离}{164 / 234}{**}
甲、乙两人赛跑，计时开始时，甲在乙前方$10$（单位：m）处，
图中实线表示甲的速度曲线$v = v_1(t)$（单位：$\mathrm{m}\,\mathrm{s}^{-1}$），虚线表示乙的速度$v = v_2(t)$，
三块阴影部分面积的数值依次为$10$、$20$、$3$，计时开始后乙追上甲的时刻记为$t_0$（单位：s），
则\pickout{}
\par\centerline{\begin{tikzpicture}[scale=\DiagramScale,x=1pt,y=1pt,xscale=0.439724,yscale=-0.439724]
\path[use as bounding box] (0,0) rectangle (440,275);
\path[fill=black!78,even odd rule,shift={(36,19)}] svg {M0 0 C2.3 2.1 3.5 3.8 4.6 6.7 C4.8 7.4 5.1 8.1 5.3 8.8 C5.6 9.5 5.9 10.3 6.1 11 C6.3 11.4 6.3 11.4 6.9 13.2 C8 16.1 9 19 10 22 C9.6 21.9 9.6 21.9 7.6 21.5 C5 21 5 21 3 21 C3 44.1 3 67.2 3 91 C5.6 88.4 8 85.9 10.4 83.1 C13.4 79.6 16.6 76.3 20 73 C20.3 73.2 20.3 73.2 22 74 C18.1 79.5 13.8 84.5 9.1 89.3 C3.4 96.6 2.4 102.7 2.5 111.8 C2.5 113 2.5 114.2 2.5 115.4 C2.5 119.2 2.6 123 2.6 126.8 C2.6 129.5 2.6 132.1 2.6 134.8 C2.6 141 2.7 147.3 2.7 153.6 C2.8 160.7 2.8 167.8 2.8 175 C2.8 189.7 2.9 204.3 3 219 C3.8 218.9 4.6 218.9 5.5 218.8 C12.4 218.2 19.3 217.9 26.2 217.9 C27.1 217.9 28 217.9 28.9 217.9 C31.7 217.9 34.5 217.9 37.3 217.9 C39.2 217.9 41.1 217.9 43.1 218 C47.7 218 52.4 218 57 218 C57 216 57 214 57 212 C58 212 59 212 60 212 C60 214 60 216 60 218 C60.6 218 61.2 218 61.8 218 C68.2 217.8 74.5 217.7 80.8 217.6 C83.1 217.5 85.5 217.5 87.8 217.4 C91.2 217.4 94.6 217.3 98 217.2 C98.5 217.2 98.5 217.2 101.2 217.2 C102.2 217.1 103.2 217.1 104.2 217.1 C104.6 217.1 104.6 217.1 106.8 217 C110.1 217 113.5 217 116.8 217.1 C117.3 217.1 117.3 217.1 119.8 217.2 C120.9 217.2 121.9 217.2 123 217.2 C124.1 217.2 125.2 217.3 126.3 217.3 C129.8 217.4 133.2 217.4 136.7 217.5 C139 217.5 141.4 217.6 143.7 217.6 C149.5 217.8 155.2 217.9 161 218 C161 216 161 214 161 212 C162 212 163 212 164 212 C164 214 164 216 164 218 C180.2 218 196.3 218 213 218 C213 216 213 214 213 212 C214 212 215 212 216 212 C216 214 216 216 216 218 C232.2 218 248.3 218 265 218 C265 217.3 265 216.7 265 216 C266 216 267 216 268 216 C268 216.7 268 217.3 268 218 C284.2 218 300.3 218 317 218 C316.9 213.6 316.9 209.2 316.8 204.7 C316.7 198.7 316.6 192.8 316.5 186.8 C316.5 186.4 316.5 186.4 316.5 184.5 C316.1 158 315.9 131.5 316 105 C317 105.3 318 105.7 319 106 C319 143 319 179.9 319 218 C326.6 218 334.2 218 342 218 C341 216.3 340 214.7 339 213 C339.3 212 339.7 211 340 210 C346.3 212.6 352.5 215.3 359 218 C359 219 359 220 359 221 C356.2 222.2 353.5 223.3 350.7 224.5 C349.9 224.8 349.1 225.2 348.3 225.5 C347.9 225.7 347.9 225.7 346 226.5 C345.3 226.8 344.6 227.1 343.9 227.4 C342 228 342 228 339 228 C339.7 226 340.3 224 341 222 C229.5 222 117.9 222 3 222 C3 232.2 3 242.5 3 253 C1.7 253 0.4 253 -1 253 C-1 242.8 -1 232.5 -1 222 C-5 222.3 -8.9 222.7 -13 223 C-15.1 223 -17.1 223 -19.2 222.9 C-20.2 222.8 -21.2 222.8 -22.2 222.7 C-23.2 222.7 -24.2 222.6 -25.2 222.6 C-26.3 222.5 -27.3 222.5 -28.4 222.4 C-30.9 222.3 -33.5 222.1 -36 222 C-36 221 -36 220 -36 219 C-31.3 218.8 -26.7 218.7 -22 218.5 C-20.4 218.5 -18.8 218.4 -17.2 218.4 C-15 218.3 -12.7 218.2 -10.4 218.1 C-10.1 218.1 -10.1 218.1 -8.3 218 C-3.2 217.9 -3.2 217.9 -1 219 C-1.3 153.3 -1.7 87.7 -2 20 C-4 20.7 -6 21.3 -8 22 C-7.4 16.7 -5.6 12.4 -3.4 7.6 C-3.1 6.9 -2.8 6.2 -2.4 5.4 C-1.6 3.6 -0.8 1.8 0 0 Z };
\path[fill=black!76,even odd rule,shift={(95,80)}] svg {M0 0 C13.5 -1.3 26.1 7 36.3 15.1 C37.9 16.4 39.4 17.7 41 19 C41.9 19.7 42.7 20.3 43.6 21 C45.7 22.7 47.3 24.4 49 26.4 C52.9 31.2 57.1 35.2 62 39 C62 38.3 62 37.7 62 37 C63 36.7 64 36.3 65 36 C64.8 36.4 64.8 36.4 64.1 38.5 C64 43.3 66 45 69.2 48.4 C71 50 71 50 73 51 C74 50 75 49 76 48 C77 51 77 51 76 54 C77.3 55.3 78.6 56.7 79.9 58 C80.2 58.4 80.2 58.4 82.1 60.2 C84 62 84 62 86 62 C86.4 61.2 86.8 60.5 87.2 59.7 C87.7 58.7 88.3 57.7 88.8 56.6 C89.3 55.6 89.9 54.6 90.4 53.6 C92 51 92 51 95 49 C94.4 52.7 93.4 55.4 91.4 58.6 C89 62.7 89 62.7 89 66 C90.1 66.6 91.3 67.2 92.4 67.8 C95.9 69.6 95.9 69.6 97 73 C87.8 71.9 77.8 62.7 72.1 55.6 C69.6 52.5 66.7 50.1 63.6 47.6 C63.3 47.3 63.3 47.3 62 46 C62 45.3 62 44.7 62 44 C61.3 44 60.7 44 60 44 C58.5 42.5 58.5 42.5 56.9 40.5 C56.3 39.9 55.8 39.2 55.2 38.5 C54.8 38 54.4 37.5 54 37 C53.7 38 53.3 39 53 40 C52 39.8 52 39.8 47 39 C46.6 39.8 46.1 40.5 45.7 41.3 C44.3 43.6 43 45.2 41 47 C40 46.7 39 46.3 38 46 C38.4 45.3 38.8 44.7 39.2 44 C47 31.2 47 31.2 47 28 C46.3 28 45.7 28 45 28 C43.3 26.3 41.7 24.7 40 23 C39 22.7 38 22.3 37 22 C34.6 26.4 32.2 30.9 29.8 35.3 C29 36.9 28.2 38.4 27.4 39.9 C26.2 42.1 25.1 44.2 23.9 46.4 C23.7 46.7 23.7 46.7 22.8 48.4 C21.2 51.4 19.9 54.2 18.9 57.4 C18.6 58.3 18.3 59.1 18 60 C15.9 60.9 15.9 60.9 14 61 C14.4 54.7 17.5 50 20.6 44.6 C21.7 42.7 22.7 40.7 23.8 38.8 C24.4 37.8 24.9 36.8 25.5 35.8 C28.8 29.6 31.9 23.3 35 17 C34.4 16.7 33.8 16.4 33.2 16.1 C31 15 31 15 28.4 13.4 C26 12 26 12 24 12 C21.4 16.6 18.7 21.2 16 26 C15.3 26 14.7 26 14 26 C14.6 21.8 16.1 18.5 18.1 14.7 C18.6 13.6 19.2 12.6 19.7 11.4 C19.9 11 19.9 11 21 9 C19.4 8.3 17.8 7.7 16.1 7 C15.2 6.6 14.3 6.3 13.4 5.9 C11 5 11 5 9 5 C6 10.6 3.1 16.2 0 22 C-1 22 -2 22 -3 22 C-1.6 17.4 0.3 13.4 2.6 9.1 C3.3 8 3.9 6.8 4.5 5.6 C5 4.8 5.5 3.9 6 3 C5 2.8 5 2.8 0 2 C0 1.3 0 0.7 0 0 Z M49 33 C48.7 33.7 48.3 34.3 48 35 C49 35.3 50 35.7 51 36 C50.7 35 50.3 34 50 33 C49.7 33 49.3 33 49 33 Z };
\path[fill=black!71,even odd rule,shift={(91.25,78.875)}] svg {M0 0 C0.7 0 1.5 0 2.3 0 C4.1 0.1 5.9 0.1 7.8 0.1 C6.4 0.5 5.1 0.8 3.8 1.1 C3.8 1.8 3.8 2.4 3.8 3.1 C0.8 3.1 -2.2 3.1 -5.2 3.1 C-5.5 4 -5.8 4.9 -6.1 5.8 C-7.2 8.9 -8.5 11.7 -10.1 14.7 C-10.7 15.7 -11.2 16.8 -11.8 17.8 C-12.4 18.9 -13 20 -13.6 21.2 C-14.2 22.3 -14.9 23.5 -15.5 24.6 C-21.3 35.2 -27.2 45.7 -33.2 56.1 C-34.6 56.1 -35.9 56.1 -37.2 56.1 C-37.2 55.1 -37.2 54.1 -37.2 53.1 C-36.3 53.1 -35.3 53.1 -34.2 53.1 C-34 52.4 -33.7 51.8 -33.5 51.1 C-32 47.6 -30.3 44.3 -28.5 41 C-28.2 40.3 -27.8 39.7 -27.5 39 C-26.3 36.8 -25.2 34.7 -24 32.5 C-23.2 31.1 -22.5 29.6 -21.7 28.2 C-17.6 20.5 -13.5 12.8 -9.2 5.1 C-10.2 4.1 -11.2 3.1 -12.2 2.1 C-8.4 -0.4 -4.4 -0.1 0 0 Z };
\path[fill=black!74,even odd rule,shift={(281,60)}] svg {M0 0 C1 0.3 2 0.7 3 1 C3 1.7 3 2.3 3 3 C2.3 3 1.7 3 1 3 C0.9 3.4 0.9 3.4 0.2 5.1 C-1 8.1 -2.4 10.9 -3.9 13.8 C-4.5 14.9 -5.1 16 -5.6 17.1 C-6.3 18.3 -6.9 19.5 -7.5 20.7 C-7.8 21.3 -8.1 21.9 -8.4 22.5 C-13.4 32.1 -18.5 41.5 -23.8 50.9 C-25.1 53.1 -26.3 55.4 -27.5 57.6 C-27.9 58.4 -28.3 59.1 -28.7 59.9 C-29.5 61.3 -30.3 62.7 -31 64.1 C-33.7 68.9 -33.7 68.9 -37 70 C-35.7 65.8 -34.1 62.1 -32.1 58.3 C-31.5 57.2 -30.9 56.1 -30.3 54.9 C-29.6 53.7 -29 52.6 -28.3 51.4 C-27.6 50.1 -27 48.9 -26.3 47.7 C-25.3 45.7 -24.2 43.8 -23.2 41.9 C-20.7 37.4 -18.3 32.8 -15.8 28.2 C-15.2 26.9 -14.5 25.7 -13.8 24.4 C-12.6 22 -11.3 19.6 -10 17.2 C-6.9 11.3 -3.6 5.6 0 0 Z };
\path[fill=black!76,even odd rule,shift={(102,12)}] svg {M0 0 C0.7 0 1.3 0 2 0 C1.4 5.1 0.3 9.5 -1.4 14.4 C-1.7 15 -1.7 15 -2.7 18.1 C-4 21 -4 21 -6 22 C-8.4 21.9 -10.8 21.8 -13.1 21.6 C-14.4 21.5 -15.7 21.4 -17 21.3 C-18 21.2 -19 21.1 -20 21 C-20 16.7 -20 12.4 -20 8 C-18.7 8 -17.4 8 -16 8 C-13.7 7.7 -11.3 7.4 -9 7 C-7.9 10.7 -6.9 14.3 -6 18 C-4 12.1 -2 6.1 0 0 Z M-14 10 C-15 11 -15 11 -15.1 12.8 C-15.1 13.6 -15.1 14.3 -15.1 15.1 C-15 16.4 -15 17.7 -15 19 C-13.7 19 -12.4 19 -11 19 C-11 16 -11 13.1 -11 10 C-12 10 -13 10 -14 10 Z };
\path[fill=black!77,even odd rule,shift={(232,116)}] svg {M0 0 C1 0 2 0 3 0 C1.2 5.1 -1.5 9.7 -4.2 14.4 C-8.7 22.5 -13 30.7 -17 39 C-14.7 39 -12.4 39 -10 39 C-10 39.7 -10 40.3 -10 41 C-17.1 42.8 -25.2 44.4 -32 41 C-32 40.3 -32 39.7 -32 39 C-28.4 39 -24.7 39 -21 39 C-20.5 38 -19.9 37 -19.4 36 C-17.4 32.2 -15.4 28.5 -13.5 24.7 C-12.6 23.1 -11.8 21.5 -10.9 19.9 C-9.7 17.6 -8.5 15.3 -7.2 12.9 C-6.9 12.2 -6.5 11.5 -6.1 10.8 C-5.1 8.8 -4 6.9 -3 5 C-2.5 4 -2 3 -1.4 2 C-1 1.3 -0.5 0.7 0 0 Z };
\path[fill=black!75,even odd rule,shift={(224,36)}] svg {M0 0 C0 0.7 0 1.3 0 2 C-0.6 2.5 -1.2 3.1 -1.7 3.7 C-4.6 6.6 -6.3 9.9 -8.1 13.4 C-8.3 13.8 -8.3 13.8 -9.3 15.6 C-10.5 17.9 -11.7 20.3 -12.9 22.6 C-14.1 24.8 -15.3 27.1 -16.5 29.4 C-17.2 30.9 -18 32.4 -18.8 33.8 C-21.8 39.6 -24.9 45.3 -28 51 C-29 51 -30 51 -31 51 C-29.7 46.8 -28.1 43.1 -26 39.3 C-25.4 38.2 -24.9 37.1 -24.3 36 C-23.6 34.9 -23 33.7 -22.4 32.6 C-21.1 30.2 -19.8 27.9 -18.6 25.6 C-18 24.4 -17.4 23.3 -16.7 22.2 C-14.5 18.1 -12.3 14 -10.1 9.9 C-9.5 8.7 -8.9 7.6 -8.3 6.4 C-8.1 6 -8.1 6 -7 4 C-8 3.7 -9 3.3 -10 3 C-4 0 -4 0 0 0 Z };
\path[fill=black!74,even odd rule,shift={(237,35)}] svg {M0 0 C1.3 0.3 2.6 0.7 4 1 C2.6 4.5 1.1 7.8 -0.6 11.2 C-0.9 11.7 -0.9 11.7 -2.2 14.2 C-2.7 15.3 -3.3 16.3 -3.9 17.4 C-4.2 18 -4.2 18 -5.6 20.8 C-9.6 28.3 -13.6 35.9 -17.7 43.4 C-18.2 44.3 -18.7 45.2 -19.2 46.2 C-21.1 49.5 -22.9 52.8 -25 56 C-25.7 56 -26.3 56 -27 56 C-26.4 51.5 -24.6 48 -22.5 44 C-22.1 43.4 -21.8 42.7 -21.4 42 C-20.2 39.9 -19.1 37.8 -17.9 35.6 C-17.2 34.2 -16.4 32.8 -15.7 31.4 C-12.7 25.9 -9.7 20.3 -6.7 14.8 C-6 13.4 -5.2 12.1 -4.5 10.7 C-4.2 10.1 -3.8 9.5 -3.5 8.8 C-1.9 5.9 -0.6 3.3 0 0 Z };
\path[fill=black!76,even odd rule,shift={(417,225)}] svg {M0 0 C0 0.7 0 1.3 0 2 C-0.5 2.4 -1 2.9 -1.4 3.3 C-4.4 6.5 -4 10.9 -4 15 C-3.2 19 -2 22.4 0 26 C-2.9 24.7 -4.9 23.4 -7 21 C-7.8 18.8 -7.8 18.8 -8.2 16.5 C-8.3 16.1 -8.3 16.1 -8.8 14.2 C-9 12 -9 12 -8 9 C-8.3 9.2 -8.3 9.2 -10 10 C-10.7 13.1 -11 15.9 -11 19 C-11.7 19.3 -11.7 19.3 -15 21 C-16 20 -16 20 -16 17.2 C-16 16 -15.9 14.8 -15.8 13.6 C-15.8 12.5 -15.7 11.3 -15.6 10.1 C-15 7 -15 7 -13.1 5.1 C-11 4 -11 4 -9 4 C-8.7 5 -8.3 6 -8 7 C-7.7 6.5 -7.7 6.5 -6.1 4.1 C-3.3 0 -3.3 0 0 0 Z };
\path[fill=black!73,even odd rule,shift={(263,244)}] svg {M0 0 C2.7 2.9 3.4 5 3.4 8.9 C3.2 13 2.9 15.9 0 19 C-2.4 20 -2.4 20 -5 20 C-8.1 18.3 -8.9 17.4 -10 14 C-10.5 8.6 -10.6 4.4 -7 0 C-4.2 -1.4 -2.9 -0.9 0 0 Z M-4 1 C-6.3 4.5 -6.3 6 -6.3 10.1 C-6.3 11.1 -6.3 12.2 -6.4 13.3 C-6.2 14.2 -6.1 15.1 -6 16 C-5.5 16.3 -5.5 16.3 -3 18 C-0.2 15.2 -0.7 13.4 -0.7 9.5 C-0.7 8.3 -0.7 7.2 -0.6 6 C-0.8 5 -0.9 4 -1 3 C-2 2.3 -3 1.7 -4 1 Z };
\path[fill=black!77,even odd rule,shift={(232,92)}] svg {M0 0 C2.9 3.1 4.3 5.8 4.5 10 C4.3 14.2 2.9 16.9 0 20 C-3.1 20.6 -3.1 20.6 -6 20 C-9.2 17.4 -9.7 14.1 -10.1 10.1 C-9.9 6.5 -8.7 4.1 -7 1 C-4.3 -0.3 -2.9 -0.5 0 0 Z M-4 2 C-6.8 6.9 -5.8 11.6 -5 17 C-3.7 17 -2.4 17 -1 17 C0.5 12.5 0.8 7.9 -1 3.5 C-1.3 3 -1.7 2.5 -2 2 C-2.7 2 -3.3 2 -4 2 Z };
\path[fill=black!74,even odd rule,shift={(253,40)}] svg {M0 0 C1 0.3 2 0.7 3 1 C2.4 6.1 -0 10 -2.6 14.3 C-6 20.3 -9.2 26.3 -12.2 32.5 C-18.9 45.9 -18.9 45.9 -22 50 C-22.5 50.2 -22.5 50.2 -25 51 C-24.3 47 -22.3 43.6 -20.4 40 C-20 39.3 -19.7 38.6 -19.3 37.9 C-18.1 35.8 -17 33.6 -15.8 31.4 C-15 30 -14.3 28.6 -13.5 27.1 C-9.8 20.1 -6 13.1 -2.1 6.2 C-1.1 4.1 -0.5 2.3 0 0 Z };
\path[fill=black!78,even odd rule,shift={(366,243)}] svg {M0 0 C2.8 2.3 3.3 4.9 3.9 8.4 C4.2 12.7 3.6 15.4 1 19 C-1.7 20.4 -1.7 20.4 -4 21 C-7.2 18.3 -8.8 16.7 -9.9 12.6 C-10 8.6 -9.5 5.7 -8 2 C-4.8 -0.1 -3.7 -0.6 0 0 Z M-4 2 C-6.3 7.3 -5.9 11.4 -5 17 C-3.7 17 -2.4 17 -1 17 C0.2 14.7 0.2 13.2 0.2 10.6 C0.2 9.9 0.2 9.1 0.2 8.3 C-0 5.8 -0.7 4.2 -2 2 C-2.7 2 -3.3 2 -4 2 Z };
\path[fill=black!75,even odd rule,shift={(296,68)}] svg {M0 0 C1 0.3 2 0.7 3 1 C3 1.7 3 2.3 3 3 C2 3 1 3 0 3 C-0.3 3.9 -0.5 4.7 -0.8 5.6 C-2 8.9 -3.4 11.9 -5 15 C-4.6 14.3 -4.1 13.6 -3.6 12.9 C-3 12 -2.4 11.1 -1.8 10.2 C-1.2 9.3 -0.6 8.4 -0 7.5 C2.1 4.9 4 3.5 7 2 C7.3 3.3 7.7 4.6 8 6 C7.3 6.4 6.5 6.7 5.7 7.1 C2.5 9.3 1.1 11.5 -0.9 14.9 C-4.7 20.9 -8.1 24.7 -15 27 C-13 21.7 -10.3 16.8 -7.5 11.8 C-7 10.9 -6.5 10 -6 9 C-5.5 8.2 -5 7.3 -4.5 6.4 C-4.1 5.6 -3.7 4.8 -3.2 4 C-2 2 -2 2 0 0 Z M-7 16 C-6 18 -6 18 -6 18 Z };
\path[fill=black!73,even odd rule,shift={(214,116)}] svg {M0 0 C0.7 0 1.3 0 2 0 C1.4 3 0.4 5.4 -1 8 C-1.4 8.8 -1.8 9.6 -2.2 10.3 C-2.7 11.2 -3.1 12 -3.6 12.8 C-3.8 13.2 -3.8 13.2 -4.9 15.3 C-6.3 17.9 -7.6 20.4 -9 23 C-9.8 24.5 -10.6 26.1 -11.4 27.6 C-13.3 31.1 -15.1 34.5 -17 38 C-15.7 38.3 -14.4 38.7 -13 39 C-13.3 39.7 -13.7 40.3 -14 41 C-17 40 -19.9 39 -23 38 C-23 37 -23 36 -23 35 C-21.7 35.3 -20.4 35.7 -19 36 C-18.7 35.1 -18.5 34.3 -18.2 33.4 C-17 30 -15.6 27 -13.8 23.9 C-13.2 22.7 -12.6 21.6 -12 20.4 C-11.7 19.8 -11.4 19.3 -11 18.7 C-10.1 16.9 -9.1 15.1 -8.1 13.3 C-2.5 2.9 -2.5 2.9 0 0 Z };
\path[fill=black!82,even odd rule,shift={(353,48)}] svg {M0 0 C0.3 0 0.7 0 1 0 C1.8 4.9 2.1 9.5 2.1 14.5 C2.1 15.1 2.1 15.8 2.1 16.5 C2.1 18.7 2.1 20.8 2.1 22.9 C2.1 24.4 2.1 25.9 2 27.3 C2 30.9 2 34.4 2 38 C1 38 0 38 -1 38 C-1.7 36.3 -2.3 34.7 -3 33 C-2.3 33 -1.7 33 -1 33 C-1 30 -1 27.1 -1 24 C-1.3 24.6 -1.3 24.6 -3 27.8 C-3.4 28.5 -3.7 29.2 -4.1 30 C-5 32 -5 32 -5 35 C-6.3 34.7 -7.6 34.3 -9 34 C-8.1 29.4 -5.7 25.7 -3.4 21.7 C-1.1 17.3 -1 14 -1 9.1 C-1 6 -0.5 3.1 0 0 Z };
\path[fill=black!68,even odd rule,shift={(242,243)}] svg {M0 0 C3.6 0.4 3.6 0.4 7 1 C8.2 3.6 8.2 3.6 9 7 C7.7 10.5 5.7 13.3 3 16 C5.3 16 7.6 16 10 16 C9.7 16.8 9.7 16.8 8 21 C4.4 21 0.7 21 -3 21 C-3 16.6 -0.5 15.1 2.5 12.1 C4.3 9.6 4.3 8.1 4 5 C1.7 5.3 -0.6 5.7 -3 6 C-2.4 3.1 -2.1 2.1 0 0 Z };
\path[fill=black!73,even odd rule,shift={(183,72)}] svg {M0 0 C1 3 1 3 0.1 5.3 C-0.2 5.7 -0.2 5.7 -1.4 7.9 C-1.7 8.4 -1.7 8.4 -3 10.9 C-3.6 11.9 -4.2 13 -4.8 14.1 C-5.4 15.1 -6 16.2 -6.6 17.3 C-8.4 20.5 -10.2 23.8 -12 27 C-13 28.9 -14.1 30.7 -15.1 32.6 C-17.4 36.7 -19.7 40.9 -22 45 C-23 45 -24 45 -25 45 C-23.7 41.5 -22.2 38.3 -20.5 35 C-20 34 -19.5 33.1 -19 32.1 C-18.4 31.1 -17.9 30.1 -17.4 29.1 C-16.8 28 -16.3 27 -15.8 26 C-3.9 3.5 -3.9 3.5 0 0 Z };
\path[fill=black!72,even odd rule,shift={(273,111)}] svg {M0 0 C0.2 0.5 0.2 0.5 1 3 C-0.1 4.7 -0.1 4.7 -1.9 6.5 C-4.9 9.7 -7.6 13 -10.4 16.4 C-16.5 24 -22.9 30.6 -31 36 C-31.2 35.5 -31.2 35.5 -32 33 C-31.2 32.3 -30.4 31.6 -29.6 30.8 C-28.5 29.8 -27.4 28.9 -26.3 27.9 C-25.8 27.4 -25.3 26.9 -24.7 26.4 C-22 24 -19.5 21.5 -17.2 18.6 C-16.5 17.8 -15.8 16.9 -15 16 C-14.3 16 -13.7 16 -13 16 C-12.7 15.4 -12.5 14.9 -12.2 14.3 C-10.9 11.7 -9.2 9.6 -7.4 7.4 C-4.1 3.2 -4.1 3.2 -3 1 C-2 0.7 -1 0.3 0 0 Z };
\path[fill=black!71,even odd rule,shift={(298,243)}] svg {M0 0 C1.9 1.1 1.9 1.1 3 3 C3 9.1 1 11.8 -3 16 C-0.7 16 1.6 16 4 16 C3.1 19.9 3.1 19.9 2 21 C0.3 21.1 -1.4 21.1 -3.1 21.1 C-4 21.1 -4.9 21 -5.8 21 C-6.6 21 -7.3 21 -8 21 C-8 16.6 -7.9 16.4 -5.1 13.5 C-2.7 11 -1.6 9.4 -1 6 C-2 5.3 -3 4.7 -4 4 C-6.7 4.9 -6.7 4.9 -9 6 C-8.7 4.1 -8.7 4.1 -8 2 C-4.8 -0.2 -3.7 -0.5 0 0 Z };
\path[fill=black!71,even odd rule,shift={(192,87)}] svg {M0 0 C1 0.3 2 0.7 3 1 C-12.1 30.3 -12.1 30.3 -19 42 C-20 42 -21 42 -22 42 C-20.4 37.7 -18.3 33.7 -16.2 29.7 C-15.8 29 -15.5 28.3 -15.1 27.6 C-14.3 26.1 -13.5 24.7 -12.8 23.2 C-11.6 21 -10.4 18.8 -9.2 16.6 C-8.5 15.2 -7.7 13.7 -6.9 12.3 C-6.6 11.7 -6.2 11 -5.9 10.3 C-4 6.8 -2.1 3.4 0 0 Z };
\path[fill=black!75,even odd rule,shift={(329,45)}] svg {M0 0 C-1.5 3.1 -3.2 6.2 -4.8 9.2 C-5.9 11.3 -7 13.4 -8.1 15.4 C-8.7 16.5 -9.3 17.5 -9.8 18.6 C-10.4 19.6 -10.9 20.6 -11.4 21.6 C-11.9 22.5 -12.4 23.4 -12.9 24.4 C-14.1 27.1 -14.2 29.1 -14 32 C-15.3 32 -16.6 32 -18 32 C-17.5 25.8 -15.4 21.4 -12.4 16.1 C-12 15.3 -11.6 14.5 -11.2 13.7 C-10.1 11.8 -9.1 9.9 -8 8 C-8.8 8.5 -8.8 8.5 -13 11 C-11.8 6 -8.9 4 -5 1.1 C-3 0 -3 0 0 0 Z };
\path[fill=black!77,even odd rule,shift={(214,91.5)}] svg {M0 0 C3 0.5 3 0.5 4.9 1.6 C6.7 4.8 5.9 7 5 10.5 C3.8 12.7 2.5 14.5 1 16.5 C3 16.2 5 15.8 7 15.5 C6.7 17.2 6.3 18.8 6 20.5 C2 20.5 -1.9 20.5 -6 20.5 C-4 16.5 -4 16.5 -1.4 13.6 C1.1 10.4 1.7 8.6 2 4.5 C1.6 4.6 1.6 4.6 -0.3 5 C-3 5.5 -3 5.5 -6 5.5 C-3.8 0.6 -3.8 0.6 0 0 Z };
\path[fill=black!77,even odd rule,shift={(91,244)}] svg {M0 0 C2.6 0 5.3 0 8 0 C7.7 1.3 7.3 2.6 7 4 C5.3 4 3.7 4 2 4 C2.5 4.3 2.5 4.3 5 6.1 C7.7 8.3 8 8.8 8.6 12.6 C8 16 8 16 6.4 18.4 C3.2 20.5 0.8 20.3 -3 20 C-3.3 19 -3.7 18 -4 17 C-0.9 16 -0 16.1 3 17 C3.6 14.6 3.6 14.6 4 12 C1.8 9.8 -0.1 9.1 -3 8 C-2 5.4 -1 2.7 0 0 Z };
\path[fill=black!75,even odd rule,shift={(204,244)}] svg {M0 0 C1.9 -0.3 1.9 -0.3 4.1 -0.2 C5.4 -0.1 6.7 -0.1 8 0 C7.7 1 7.3 2 7 3 C5 3 3 3 1 3 C1 3.7 1 4.3 1 5 C1.8 5.2 2.6 5.5 3.4 5.8 C6 7 6 7 7.4 8.8 C8.2 11.9 7.9 13.9 7 17 C4.6 19.7 3.4 20 -0.2 20.2 C-0.7 20.2 -0.7 20.2 -3 20 C-3.3 18.7 -3.7 17.4 -4 16 C-1.7 16 0.6 16 3 16 C3.3 15 3.7 14 4 13 C1.8 9.7 0.7 9.2 -3 8 C-2.5 4.6 -1.9 2.9 0 0 Z };
\path[fill=black!72,even odd rule,shift={(308,243)}] svg {M0 0 C2.3 0 4.6 0 7 0 C7.3 0.7 7.7 1.3 8 2 C7.7 2.7 7.3 3.3 7 4 C3.9 4.6 3.9 4.6 1 5 C2 5.5 3 6 4 6.5 C7 9 7 9 7.6 10.9 C7.8 14.8 7.5 16.5 4.8 19.5 C1.9 21.1 0.3 21.4 -3 21 C-3.3 19.7 -3.7 18.4 -4 17 C-1.7 17 0.6 17 3 17 C3.3 16 3.7 15 4 14 C1.6 11.5 0.3 10.1 -3 9 C-2.4 5.7 -1.6 3 0 0 Z };
\path[fill=black!72,even odd rule,shift={(243,130)}] svg {M0 0 C1 0.3 2 0.7 3 1 C0.7 5.6 -1.6 10.2 -4 15 C-3.3 15 -2.7 15 -2 15 C-2.4 17.4 -2.6 18.7 -4.6 20.2 C-5.2 20.5 -5.9 20.8 -6.5 21.1 C-7.2 21.5 -7.9 21.8 -8.7 22.2 C-9.4 22.5 -10.2 22.9 -10.9 23.2 C-11.7 23.6 -12.4 24 -13.2 24.4 C-18.7 27 -18.7 27 -21 27 C-21.3 26 -21.7 25 -22 24 C-20.4 23.3 -18.8 22.7 -17.1 22 C-16.2 21.6 -15.3 21.3 -14.4 20.9 C-12 20 -12 20 -10 20 C-9.8 19.3 -9.5 18.7 -9.3 18 C-7.7 14.3 -5.8 10.7 -3.9 7.2 C-3.5 6.5 -3.1 5.8 -2.7 5.1 C-1.8 3.4 -0.9 1.7 0 0 Z };
\path[fill=black!78,even odd rule,shift={(339,49)}] svg {M0 0 C0.4 2 0.7 4 1 6 C1.1 6.8 1.3 7.6 1.4 8.5 C2.4 15.4 2.4 15.4 0.3 18.6 C-1.9 20 -3.4 20.7 -6 21 C-7 20.7 -8 20.3 -9 20 C-9 19 -9 18 -9 17 C-6.4 17 -3.7 17 -1 17 C-1.8 14.6 -1.8 14.6 -3 12 C-4 11.7 -5 11.3 -6 11 C-5.7 10 -5.3 9 -5 8 C-4.3 8 -3.7 8 -3 8 C-3 6.4 -3 4.7 -3 3 C-5 3.7 -7 4.3 -9 5 C-8.8 3.1 -8.8 3.1 -8 1 C-5.1 -1.1 -3.3 -1.2 0 0 Z };
\path[fill=black!75,even odd rule,shift={(348.4375,242.8125)}] svg {M0 0 C2.6 0.2 2.6 0.2 4.6 2.2 C4.6 4.5 4.6 6.9 4.6 9.2 C4.9 9.5 5.2 9.8 5.6 10.2 C5.8 13.4 5.9 15.6 4.4 18.5 C1.9 20.8 -0.1 21 -3.4 21.2 C-4.1 20.9 -4.8 20.5 -5.4 20.2 C-5.4 19.2 -5.4 18.2 -5.4 17.2 C-3.1 17.5 -0.8 17.8 1.6 18.2 C1.6 16.5 1.6 14.9 1.6 13.2 C0.2 12.5 -1.1 11.9 -2.4 11.2 C-2.4 10.5 -2.4 9.9 -2.4 9.2 C-1.4 8.5 -0.5 7.9 0.6 7.2 C0.6 6.2 0.6 5.2 0.6 4.2 C-1.4 4.5 -3.4 4.8 -5.4 5.2 C-3 0.2 -3 0.2 0 0 Z };
\path[fill=black!71,even odd rule,shift={(57,99)}] svg {M0 0 C1 3 1 3 -0.1 5.5 C-0.6 6.4 -1.2 7.3 -1.7 8.3 C-2.3 9.4 -2.9 10.4 -3.5 11.4 C-4.2 12.5 -4.8 13.6 -5.4 14.7 C-6.1 15.8 -6.7 16.9 -7.3 18 C-8.9 20.7 -10.4 23.3 -12 26 C-11.3 26.7 -10.7 27.3 -10 28 C-10 29 -10 30 -10 31 C-10.6 30.7 -11.1 30.4 -11.7 30.1 C-14.1 29 -16.5 28 -19 27 C-18.7 26 -18.3 25 -18 24 C-16.7 24.3 -15.4 24.7 -14 25 C-13.8 24.2 -13.5 23.4 -13.3 22.6 C-11.6 18 -9.4 13.9 -7 9.7 C-6.8 9.3 -6.8 9.3 -5.6 7.2 C-5.2 6.4 -4.8 5.6 -4.3 4.8 C-3.9 4.1 -3.5 3.4 -3.1 2.7 C-2 1 -2 1 0 0 Z };
\path[fill=black!71,even odd rule,shift={(352.625,43.5)}] svg {M0 0 C2.7 2.8 2.7 3.5 2.8 7.2 C2.7 8.7 2.5 10.1 2.4 11.5 C2 11.5 1.7 11.5 1.4 11.5 C1.4 9.2 1.4 6.9 1.4 4.5 C0.6 4.2 -0.2 3.9 -1 3.6 C-3.6 2.5 -3.6 2.5 -5.6 1.5 C-6.3 3.1 -6.9 4.8 -7.6 6.5 C-8.6 6.2 -9.6 5.8 -10.6 5.5 C-10 3.9 -9.3 2.2 -8.6 0.5 C-11.1 1 -11.1 1 -23.6 3.5 C-24 2.5 -24.3 1.5 -24.6 0.5 C-15.8 -2.5 -8.7 -3.6 0 0 Z };
\path[fill=black!80,even odd rule,shift={(20,244)}] svg {M0 0 C7.4 -0.4 7.4 -0.4 9.9 1.2 C12 4.7 12.1 6.4 11.3 10.4 C10 14.1 8.7 16.3 6 19 C3.8 19.8 3.8 19.8 2 20 C2.5 17.5 3.1 15.8 4.5 13.6 C6 10.9 6.7 9 7 6 C6.7 5 6.3 4 6 3 C3.7 3 1.4 3 -1 3 C-0.7 2 -0.3 1 0 0 Z };
\path[fill=black!69,even odd rule,shift={(426,231)}] svg {M0 0 C0.3 1.7 0.7 3.3 1 5 C0.3 5 -0.3 5 -1 5 C-1.3 4.3 -1.7 3.7 -2 3 C-2.3 3.2 -2.3 3.2 -4 4 C-3.2 4.4 -2.4 4.8 -1.6 5.2 C1 7 1 7 2 9.5 C2 12 2 12 0.9 13.8 C-1.8 15.5 -3.9 15.3 -7 15 C-8 14 -8 14 -8.1 11.4 C-8 10.6 -8 9.8 -8 9 C-7.3 9 -6.7 9 -6 9 C-5.7 9.5 -5.7 9.5 -4 12 C-3.3 12 -2.7 12 -2 12 C-1.7 11.3 -1.3 10.7 -1 10 C-1.4 9.9 -1.4 9.9 -3.4 9.1 C-6 8 -6 8 -8 6 C-7.8 3.4 -7.8 3.4 -7 1 C-4.5 -0.2 -2.7 -0.1 0 0 Z };
\path[fill=black!78,even odd rule,shift={(80,81)}] svg {M0 0 C-0.3 1 -0.7 2 -1 3 C-3.1 4 -3.1 4 -5.7 4.9 C-13 7.7 -17.1 10.2 -20.3 17.4 C-20.6 17.9 -20.8 18.5 -21 19 C-22 19 -23 19 -24 19 C-24.8 16.9 -24.8 16.9 -25 14 C-20.3 6.1 -9.1 0 0 0 Z };
\path[fill=black!74,even odd rule,shift={(107,19)}] svg {M0 0 C3.2 -0.3 3.2 -0.3 6 0 C6.3 1.6 6.7 3.3 7 5 C6 4.7 5 4.3 4 4 C4 3.3 4 2.7 4 2 C3 2 2 2 1 2 C1 3 1 4 1 5 C1.6 5.1 2.2 5.2 2.9 5.4 C5 6 5 6 7 8 C7.3 10.5 7.3 10.5 7 13 C4 15 4 15 0.8 14.6 C-0.1 14.4 -1 14.2 -2 14 C-2 12.4 -2 10.7 -2 9 C-1.3 9 -0.7 9 0 9 C0.3 9.7 0.7 10.3 1 11 C2 11.3 3 11.7 4 12 C3.8 11.8 3.8 11.8 2.8 10.8 C0.8 8.8 -1.1 6.9 -3 5 C-1.9 2.5 -1.9 2.5 0 0 Z };
\path[fill=black!66,even odd rule,shift={(83,101)}] svg {M0 0 C1 0.3 2 0.7 3 1 C3 6.9 3 12.9 3 19 C3.7 19.3 4.3 19.7 5 20 C2.3 21.4 -0 21.1 -3 21 C-3 20.3 -3 19.7 -3 19 C-2.3 19 -1.7 19 -1 19 C-1 14 -1 9.1 -1 4 C-1.7 4 -2.3 4 -3 4 C-3 3.3 -3 2.7 -3 2 C-2 1.3 -1 0.7 0 0 Z };
\path[fill=black!65,even odd rule,shift={(192,243)}] svg {M0 0 C0.7 0.3 1.3 0.7 2 1 C2 6.9 2 12.9 2 19 C2.7 19.3 3.3 19.7 4 20 C3 21 3 21 -0.1 21.1 C-1 21 -2 21 -3 21 C-3.3 20.3 -3.7 19.7 -4 19 C-3.3 19 -2.7 19 -2 19 C-2 14.1 -2 9.1 -2 4 C-2.7 4 -3.3 4 -4 4 C-4 3.3 -4 2.7 -4 2 C-2.1 0.9 -2.1 0.9 0 0 Z };
\path[fill=black!68,even odd rule,shift={(205,98)}] svg {M0 0 C0.3 0 0.7 0 1 0 C0.6 6.3 -1.6 10.6 -4.8 16 C-5.3 16.8 -5.7 17.5 -6.1 18.3 C-7.4 20.6 -8.7 22.8 -10 25 C-11.3 27.3 -12.7 29.7 -14 32 C-15 32 -16 32 -17 32 C-15.5 28.1 -13.7 24.5 -11.8 20.9 C-11.2 19.7 -10.5 18.6 -9.9 17.4 C-9.6 16.8 -9.3 16.2 -8.9 15.6 C-7.9 13.7 -7 11.9 -6 10.1 C-1.1 1.1 -1.1 1.1 0 0 Z };
\path[fill=black!75,even odd rule,shift={(115,13)}] svg {M0 0 C3.1 0.3 4 1 6.2 3.4 C9.1 7.6 9.4 10.9 9 16 C7.6 20.3 6.3 23.3 2.2 25.5 C1.9 25.6 1.9 25.6 0 26 C1.2 23.5 2.5 21.3 4 19 C4.8 12.7 4.5 7.4 1 2 C0.7 1.3 0.3 0.7 0 0 Z };
\path[fill=black!84,even odd rule,shift={(288,90)}] svg {M0 0 C0.3 0.7 0.7 1.3 1 2 C-1.2 5.2 -3.4 8.4 -5.7 11.6 C-7 13.6 -8.3 15.5 -9.6 17.4 C-10.7 19 -11.8 20.5 -13 22 C-14 22 -15 22 -16 22 C-14.5 18.2 -13 14.5 -11.2 10.9 C-10.8 9.8 -10.3 8.8 -9.8 7.7 C-7.6 4.4 -5.6 3.6 -2 2 C-1.3 1.3 -0.7 0.7 0 0 Z };
\path[fill=black!67,even odd rule,shift={(429,225)}] svg {M0 0 C4.3 1.5 6 4 8 8 C8.8 13.1 8.6 16.9 6.2 21.5 C4.2 23.8 2.9 25 0 26 C1 24 2 22 3 20 C3.1 17.9 3.2 15.8 3.2 13.7 C3.3 12.6 3.3 11.4 3.3 10.3 C3 6.7 2 4.9 0 2 C0 1.3 0 0.7 0 0 Z };
\path[fill=black!77,even odd rule,shift={(72,12)}] svg {M0 0 C0.3 1 0.7 2 1 3 C0.3 3 -0.3 3 -1 3 C-2.1 6.5 -2.2 9.8 -2.2 13.4 C-2.3 14.5 -2.3 15.7 -2.3 16.8 C-2 20.3 -1 22.1 1 25 C1.3 25.7 1.7 26.3 2 27 C-1.1 26.7 -2 26 -4.2 23.6 C-7.1 19.4 -7.3 16.1 -7 11 C-5.4 6.3 -3.7 3.2 0 0 Z };
\path[fill=black!73,even odd rule,shift={(267,50)}] svg {M0 0 C1.3 0.3 2.6 0.7 4 1 C4 1.7 4 2.3 4 3 C3.6 3.4 3.1 3.9 2.7 4.3 C0.7 6.3 -0.4 8.4 -1.7 10.9 C-2.2 11.8 -2.7 12.8 -3.2 13.7 C-3.4 14.2 -3.4 14.2 -4.7 16.7 C-5.2 17.7 -5.7 18.7 -6.2 19.7 C-7.5 22.1 -8.8 24.6 -10 27 C-11 27 -12 27 -13 27 C-11.3 20.9 -8.3 15.6 -5.4 10.1 C-5.1 9.6 -5.1 9.6 -3.8 7.2 C-2.5 4.8 -1.3 2.4 0 0 Z };
\path[fill=black!78,even odd rule,shift={(17,246)}] svg {M0 0 C1 0.3 2 0.7 3 1 C2.5 2 2 3 1.5 4.1 C-0.1 7.8 -0.2 10.9 0 15 C1.6 15.3 3.3 15.7 5 16 C4.7 17 4.3 18 4 19 C1.1 18.7 -1.3 18.4 -3.9 16.9 C-5.7 13.8 -5.6 10.5 -5 7 C-3.5 4.3 -2.2 2.2 0 0 Z };
\path[fill=black!71,even odd rule,shift={(253,77)}] svg {M0 0 C1 0.3 2 0.7 3 1 C1.1 4.8 -0.9 8.7 -2.9 12.5 C-3.5 13.8 -4.2 15.1 -4.9 16.4 C-5.8 18.3 -6.8 20.1 -7.7 22 C-8 22.6 -8 22.6 -9.5 25.4 C-11 28 -11 28 -13 29 C-12.7 21.9 -9.4 16.6 -5.9 10.6 C-5.4 9.5 -4.8 8.5 -4.2 7.5 C-2.8 5 -1.4 2.5 0 0 Z };
\path[fill=black!73,even odd rule,shift={(62,19)}] svg {M0 0 C1.3 0.7 2.6 1.3 4 2 C2.6 5 1 7.6 -0.9 10.2 C-1.4 11 -1.9 11.7 -2.5 12.4 C-3 12.9 -3.5 13.5 -4 14 C-5 14 -6 14 -7 14 C-7.3 10.4 -7.7 6.7 -8 3 C-8.7 3 -9.3 3 -10 3 C-10 2.3 -10 1.7 -10 1 C-8.4 1 -6.7 1 -5 1 C-4.3 3.6 -3.7 6.3 -3 9 C-0.8 5.7 -0.9 4.7 -1 1 C-0.7 0.7 -0.3 0.3 0 0 Z };
\path[fill=black!99,even odd rule,shift={(353,86)}] svg {M0 0 C0.7 0 1.3 0 2 0 C2 12.9 2 25.7 2 39 C1.3 39 0.7 39 0 39 C-0.2 34.1 -0.3 29.2 -0.5 24.3 C-0.5 22.6 -0.6 20.9 -0.6 19.3 C-0.7 16.9 -0.8 14.4 -0.9 12 C-0.9 11.3 -0.9 10.6 -1 9.8 C-1.1 6.4 -1.1 3.3 0 0 Z };
\path[fill=black!69,even odd rule,shift={(99,123)}] svg {M0 0 C0.7 0 1.3 0 2 0 C0.3 6.2 -2.9 11.5 -6 17 C-5 17.3 -4 17.7 -3 18 C-3.3 19 -3.7 20 -4 21 C-4.6 20.8 -5.2 20.6 -5.8 20.4 C-8.2 20 -9.6 20.2 -12 21 C-12.3 20 -12.7 19 -13 18 C-12 17.7 -11 17.3 -10 17 C-8.8 15.4 -8.8 15.4 -7.8 13.4 C-7.4 12.7 -7 12 -6.6 11.3 C-6.2 10.5 -5.8 9.8 -5.4 9 C-4.6 7.5 -3.8 6 -3 4.6 C-2.6 3.9 -2.3 3.2 -1.9 2.6 C-1 1 -1 1 0 0 Z };
\path[fill=black!68,even odd rule,shift={(195,53)}] svg {M0 0 C-1.7 4.4 -3.7 8.6 -5.9 12.8 C-6.2 13.4 -6.6 14.1 -6.9 14.8 C-9.5 19.8 -9.5 19.8 -12.1 20.7 C-12.7 20.8 -13.4 20.9 -14 21 C-12.8 17.9 -11.5 14.9 -10 12 C-10.7 11.3 -11.3 10.7 -12 10 C-11.7 9.3 -11.3 8.7 -11 8 C-9.7 8.3 -8.4 8.7 -7 9 C-7.3 8.4 -7.5 7.9 -7.8 7.3 C-8 5 -8 5 -5.8 2.2 C-3 0 -3 0 0 0 Z };
\path[fill=black!81,even odd rule,shift={(157,244)}] svg {M0 0 C4 2 4 2 5 4 C5.5 9.1 5.8 13.6 3.3 18.2 C1 20 1 20 -1.8 20.3 C-2.5 20.2 -3.2 20.1 -4 20 C-2.7 19.3 -1.4 18.7 0 18 C0 12.1 0 6.1 0 0 Z };
\path[fill=black!88,even odd rule,shift={(143,244)}] svg {M0 0 C0.1 1.2 0.1 2.4 0.2 3.6 C0.3 5.2 0.4 6.7 0.4 8.3 C0.5 9.1 0.5 9.9 0.6 10.7 C0.6 11.5 0.6 12.2 0.7 13 C0.7 13.7 0.8 14.4 0.8 15.1 C1 17 1 17 2 20 C-0.3 20 -2.6 20 -5 20 C-4.7 19 -4.3 18 -4 17 C-3.8 14.6 -3.6 12.2 -3.4 9.8 C-3.4 8.5 -3.3 7.3 -3.2 6 C-3.1 5 -3.1 4 -3 3 C-4 2.7 -5 2.3 -6 2 C-3 0 -3 0 0 0 Z };
\path[fill=black!76,even odd rule,shift={(157,243)}] svg {M0 0 C0 1 0 2 0 3 C-1 3 -2 3 -3 3 C-3 5.3 -3 7.5 -3.1 9.8 C-3.1 11.1 -3.1 12.3 -3.1 13.6 C-3 17 -3 17 -2 20 C-3.9 19.8 -3.9 19.8 -6 19 C-8.6 14.9 -8.7 10.7 -8 6 C-5.9 2.6 -4.2 0 0 0 Z };
\path[fill=black!78,even odd rule,shift={(96,101)}] svg {M0 0 C3.7 0.5 5.6 1 8 4 C9.1 8 9.3 11.8 8.2 15.9 C7 18 5.9 19.5 4 21 C3.3 21 2.7 21 2 21 C2.3 20.2 2.7 19.3 3 18.5 C4.7 12.5 4.2 7.1 3 1 C2 0.7 1 0.3 0 0 Z };
\path[fill=black!82,even odd rule,shift={(95,102)}] svg {M0 0 C1.3 0.3 2.6 0.7 4 1 C3.3 1.7 2.7 2.3 2 3 C0.6 9.1 0.9 14.2 3 20 C-0.3 19 -1 17.9 -2.8 14.8 C-4.1 10.6 -4.1 8.3 -3 4 C-1.4 1.6 -1.4 1.6 0 0 Z };
\path[fill=black!67,even odd rule,shift={(316,52)}] svg {M0 0 C0.7 0.3 1.3 0.7 2 1 C-0.7 5.2 -3.6 9 -6.9 12.7 C-7.8 13.6 -8.6 14.6 -9.5 15.5 C-11.9 17.9 -14.1 19.4 -17 21 C-18 20.3 -19 19.7 -20 19 C-19.3 18.8 -18.5 18.6 -17.8 18.3 C-14.3 16.7 -12.9 14.8 -10.6 11.8 C-7.2 7.5 -3.9 3.7 0 0 Z };
\path[fill=black!82,even odd rule,shift={(83,19)}] svg {M0 0 C0 1 0 2 0 3 C-1 3.3 -2 3.7 -3 4 C-3 5.3 -3.1 6.6 -3.1 7.9 C-3.1 8.6 -3.2 9.3 -3.2 10.1 C-3 12 -3 12 -1 14 C-3.6 14 -6.3 14 -9 14 C-9 13.7 -9 13.7 -8.7 11.9 C-8.6 11 -8.5 10.1 -8.4 9.2 C-8.3 8.3 -8.2 7.4 -8.1 6.5 C-8 4 -8 4 -9 1 C-6 -0.5 -3.3 -0.1 0 0 Z };
\path[fill=black!70,even odd rule,shift={(80,123)}] svg {M0 0 C0.7 0.3 1.3 0.7 2 1 C0.6 4.7 -1 8.2 -2.9 11.8 C-3.4 12.7 -3.9 13.7 -4.4 14.7 C-6 17 -6 17 -8.2 17.8 C-8.8 17.9 -9.4 17.9 -10 18 C-10 15 -10 15 -8.7 13.7 C-8.2 13.1 -7.6 12.6 -7 12 C-5.7 10 -4.5 7.9 -3.3 5.8 C-2.7 4.7 -2.1 3.6 -1.4 2.5 C-1 1.7 -0.5 0.9 0 0 Z };
\path[fill=black!71,even odd rule,shift={(330,72)}] svg {M0 0 C0.7 0 1.3 0 2 0 C1.1 4.8 1.1 4.8 0 7 C0.7 7 1.3 7 2 7 C2 7.7 2 8.3 2 9 C0.4 8.7 -1.3 8.3 -3 8 C-2.5 4.6 -1.9 2.9 0 0 Z };
\path[fill=black!58,even odd rule,shift={(301,81)}] svg {M0 0 C1 0.3 2 0.7 3 1 C2.7 2.3 2.3 3.6 2 5 C1.3 5 0.7 5 0 5 C-0.6 3.1 -0.6 3.1 -1 1 C-0.7 0.7 -0.3 0.3 0 0 Z };
\path[fill=black!55,even odd rule,shift={(301,161)}] svg {M0 0 C1 0 2 0 3 0 C3 1.3 3 2.6 3 4 C1.7 4 0.4 4 -1 4 C-0.7 2.7 -0.3 1.4 0 0 Z };
\path[fill=black!60,even odd rule,shift={(301,98)}] svg {M0 0 C1 0.3 2 0.7 3 1 C2.7 2.3 2.3 3.6 2 5 C1 4.3 0 3.7 -1 3 C-0.7 2 -0.3 1 0 0 Z };
\path[fill=black!63,even odd rule,shift={(301,172)}] svg {M0 0 C1 0.3 2 0.7 3 1 C2.7 2.3 2.3 3.6 2 5 C1 4 0 3 -1 2 C-0.7 1.3 -0.3 0.7 0 0 Z };
\path[fill=black!63,even odd rule,shift={(301,155)}] svg {M0 0 C1 0.3 2 0.7 3 1 C2.7 2.3 2.3 3.6 2 5 C1.3 5 0.7 5 0 5 C0 3.3 0 1.7 0 0 Z };
\path[fill=black!63,even odd rule,shift={(301,133)}] svg {M0 0 C1 0 2 0 3 0 C2.7 1.3 2.3 2.6 2 4 C1 3.7 0 3.3 -1 3 C-0.7 2 -0.3 1 0 0 Z };
\path[fill=black!57,even odd rule,shift={(301,104)}] svg {M0 0 C1 0 2 0 3 0 C3 1.3 3 2.6 3 4 C1.7 3.7 0.4 3.3 -1 3 C-0.7 2 -0.3 1 0 0 Z };
\path[fill=black!60,even odd rule,shift={(145,230)}] svg {M0 0 C1 0 2 0 3 0 C3 1.3 3 2.6 3 4 C2 4 1 4 0 4 C0 2.7 0 1.4 0 0 Z };
\path[fill=black!65,even odd rule,shift={(301,229)}] svg {M0 0 C1 0.3 2 0.7 3 1 C2.7 2.3 2.3 3.6 2 5 C1.3 4.7 0.7 4.3 0 4 C0 2.7 0 1.4 0 0 Z };
\path[fill=black!60,even odd rule,shift={(301,224)}] svg {M0 0 C1 0 2 0 3 0 C3 1.3 3 2.6 3 4 C2 4 1 4 0 4 C0 2.7 0 1.4 0 0 Z };
\path[fill=black!60,even odd rule,shift={(301,218)}] svg {M0 0 C1 0 2 0 3 0 C3 1.3 3 2.6 3 4 C2 4 1 4 0 4 C0 2.7 0 1.4 0 0 Z };
\path[fill=black!71,even odd rule,shift={(145,213)}] svg {M0 0 C1 0 2 0 3 0 C3 1.3 3 2.6 3 4 C2 4 1 4 0 4 C0 2.7 0 1.4 0 0 Z };
\path[fill=black!65,even odd rule,shift={(145,207)}] svg {M0 0 C1 0 2 0 3 0 C3 1.3 3 2.6 3 4 C2 4 1 4 0 4 C0 2.7 0 1.4 0 0 Z };
\path[fill=black!60,even odd rule,shift={(301,201)}] svg {M0 0 C1 0 2 0 3 0 C3 1.3 3 2.6 3 4 C2 4 1 4 0 4 C0 2.7 0 1.4 0 0 Z };
\path[fill=black!65,even odd rule,shift={(145,201)}] svg {M0 0 C1 0.3 2 0.7 3 1 C2.7 2.3 2.3 3.6 2 5 C1.3 4.7 0.7 4.3 0 4 C0 2.7 0 1.4 0 0 Z };
\path[fill=black!71,even odd rule,shift={(145,196)}] svg {M0 0 C1 0 2 0 3 0 C3 1.3 3 2.6 3 4 C2 4 1 4 0 4 C0 2.7 0 1.4 0 0 Z };
\path[fill=black!71,even odd rule,shift={(145,190)}] svg {M0 0 C1 0 2 0 3 0 C3 1.3 3 2.6 3 4 C2 4 1 4 0 4 C0 2.7 0 1.4 0 0 Z };
\path[fill=black!60,even odd rule,shift={(301,184)}] svg {M0 0 C1 0 2 0 3 0 C3 1.3 3 2.6 3 4 C2 4 1 4 0 4 C0 2.7 0 1.4 0 0 Z };
\path[fill=black!71,even odd rule,shift={(145,184)}] svg {M0 0 C1 0 2 0 3 0 C3 1.3 3 2.6 3 4 C2 4 1 4 0 4 C0 2.7 0 1.4 0 0 Z };
\path[fill=black!71,even odd rule,shift={(145,179)}] svg {M0 0 C1 0 2 0 3 0 C3 1.3 3 2.6 3 4 C2 4 1 4 0 4 C0 2.7 0 1.4 0 0 Z };
\path[fill=black!65,even odd rule,shift={(301,178)}] svg {M0 0 C1 0.3 2 0.7 3 1 C2.7 2.3 2.3 3.6 2 5 C1.3 4.7 0.7 4.3 0 4 C0 2.7 0 1.4 0 0 Z };
\path[fill=black!71,even odd rule,shift={(145,173)}] svg {M0 0 C1 0 2 0 3 0 C3 1.3 3 2.6 3 4 C2 4 1 4 0 4 C0 2.7 0 1.4 0 0 Z };
\path[fill=black!71,even odd rule,shift={(145,167)}] svg {M0 0 C1 0 2 0 3 0 C3 1.3 3 2.6 3 4 C2 4 1 4 0 4 C0 2.7 0 1.4 0 0 Z };
\path[fill=black!83,even odd rule,shift={(145,156)}] svg {M0 0 C1 0 2 0 3 0 C3 1.3 3 2.6 3 4 C2 4 1 4 0 4 C0 2.7 0 1.4 0 0 Z };
\path[fill=black!77,even odd rule,shift={(145,150)}] svg {M0 0 C1 0 2 0 3 0 C3 1.3 3 2.6 3 4 C2 4 1 4 0 4 C0 2.7 0 1.4 0 0 Z };
\path[fill=black!54,even odd rule,shift={(301,144)}] svg {M0 0 C1 0 2 0 3 0 C2.7 1.3 2.3 2.6 2 4 C1.3 4 0.7 4 0 4 C-0.3 3 -0.7 2 -1 1 C-0.7 0.7 -0.3 0.3 0 0 Z };
\path[fill=black!71,even odd rule,shift={(145,139)}] svg {M0 0 C1 0 2 0 3 0 C3 1.3 3 2.6 3 4 C2 4 1 4 0 4 C0 2.7 0 1.4 0 0 Z };
\path[fill=black!65,even odd rule,shift={(301,138)}] svg {M0 0 C1 0.3 2 0.7 3 1 C3 2 3 3 3 4 C1.7 3.7 0.4 3.3 -1 3 C-0.7 2 -0.3 1 0 0 Z };
\path[fill=black!71,even odd rule,shift={(145,133)}] svg {M0 0 C1 0 2 0 3 0 C3 1.3 3 2.6 3 4 C2 4 1 4 0 4 C0 2.7 0 1.4 0 0 Z };
\path[fill=black!65,even odd rule,shift={(301,127)}] svg {M0 0 C1 0 2 0 3 0 C2.7 1.3 2.3 2.6 2 4 C1.3 4 0.7 4 0 4 C-0.3 3 -0.7 2 -1 1 C-0.7 0.7 -0.3 0.3 0 0 Z };
\path[fill=black!77,even odd rule,shift={(145,127)}] svg {M0 0 C1 0 2 0 3 0 C3 1.3 3 2.6 3 4 C2 4 1 4 0 4 C0 2.7 0 1.4 0 0 Z };
\path[fill=black!71,even odd rule,shift={(146,121)}] svg {M0 0 C0.7 0.3 1.3 0.7 2 1 C1.7 2.3 1.3 3.6 1 5 C0.3 5 -0.3 5 -1 5 C-0.7 3.3 -0.3 1.7 0 0 Z };
\path[fill=black!60,even odd rule,shift={(301,121)}] svg {M0 0 C1 0 2 0 3 0 C3 1.3 3 2.6 3 4 C2 4 1 4 0 4 C0 2.7 0 1.4 0 0 Z };
\path[fill=black!60,even odd rule,shift={(301,110)}] svg {M0 0 C1 0 2 0 3 0 C2.7 1.3 2.3 2.6 2 4 C1.3 4 0.7 4 0 4 C-0.3 3 -0.7 2 -1 1 C-0.7 0.7 -0.3 0.3 0 0 Z };
\path[fill=black!65,even odd rule,shift={(301,93)}] svg {M0 0 C1 0 2 0 3 0 C2.7 1.3 2.3 2.6 2 4 C1 3.7 0 3.3 -1 3 C-0.7 2 -0.3 1 0 0 Z };
\path[fill=black!65,even odd rule,shift={(301,75)}] svg {M0 0 C1 0.3 2 0.7 3 1 C2.7 2.3 2.3 3.6 2 5 C1.3 5 0.7 5 0 5 C0 3.3 0 1.7 0 0 Z };
\path[fill=black!75,even odd rule,shift={(146,224)}] svg {M0 0 C0.7 0 1.3 0 2 0 C2 1.3 2 2.6 2 4 C1 4 0 4 -1 4 C-0.7 2.7 -0.3 1.4 0 0 Z };
\path[fill=black!69,even odd rule,shift={(146,218)}] svg {M0 0 C0.7 0.3 1.3 0.7 2 1 C1.7 2.3 1.3 3.6 1 5 C0.3 4.7 -0.3 4.3 -1 4 C-0.7 2.7 -0.3 1.4 0 0 Z };
\path[fill=black!62,even odd rule,shift={(301,207)}] svg {M0 0 C1 0 2 0 3 0 C2.7 1.3 2.3 2.6 2 4 C1.3 4 0.7 4 0 4 C0 2.7 0 1.4 0 0 Z };
\path[fill=black!69,even odd rule,shift={(301,195)}] svg {M0 0 C1 0.3 2 0.7 3 1 C3 2 3 3 3 4 C2 4 1 4 0 4 C0 2.7 0 1.4 0 0 Z };
\path[fill=black!62,even odd rule,shift={(301,190)}] svg {M0 0 C1 0 2 0 3 0 C2.7 1.3 2.3 2.6 2 4 C1.3 4 0.7 4 0 4 C0 2.7 0 1.4 0 0 Z };
\path[fill=black!69,even odd rule,shift={(301,167)}] svg {M0 0 C1 0 2 0 3 0 C2.7 1.3 2.3 2.6 2 4 C1.3 4 0.7 4 0 4 C0 2.7 0 1.4 0 0 Z };
\path[fill=black!75,even odd rule,shift={(145,161)}] svg {M0 0 C1 0.3 2 0.7 3 1 C3 2 3 3 3 4 C2 4 1 4 0 4 C0 2.7 0 1.4 0 0 Z };
\path[fill=black!69,even odd rule,shift={(301,150)}] svg {M0 0 C1 0 2 0 3 0 C2.7 1.3 2.3 2.6 2 4 C1.3 4 0.7 4 0 4 C0 2.7 0 1.4 0 0 Z };
\path[fill=black!69,even odd rule,shift={(146,144)}] svg {M0 0 C0.7 0 1.3 0 2 0 C2 1.3 2 2.6 2 4 C1 4 0 4 -1 4 C-0.7 2.7 -0.3 1.4 0 0 Z };
\path[fill=black!56,even odd rule,shift={(116,135)}] svg {M0 0 C0.7 0.3 1.3 0.7 2 1 C2 1.7 2 2.3 2 3 C0.7 3.3 -0.6 3.7 -2 4 C-2 3 -2 2 -2 1 C-1.3 0.7 -0.7 0.3 0 0 Z };
\path[fill=black!49,even odd rule,shift={(125,130)}] svg {M0 0 C1 0 2 0 3 0 C2.3 1.3 1.7 2.6 1 4 C0.3 3.7 -0.3 3.3 -1 3 C-0.7 2 -0.3 1 0 0 Z };
\path[fill=black!62,even odd rule,shift={(49,129)}] svg {M0 0 C1 0.3 2 0.7 3 1 C3 2 3 3 3 4 C1.7 3.3 0.4 2.7 -1 2 C-0.7 1.3 -0.3 0.7 0 0 Z };
\path[fill=black!69,even odd rule,shift={(301,115)}] svg {M0 0 C1 0.3 2 0.7 3 1 C3 2 3 3 3 4 C2 4 1 4 0 4 C0 2.7 0 1.4 0 0 Z };
\path[fill=black!62,even odd rule,shift={(301,87)}] svg {M0 0 C1 0.3 2 0.7 3 1 C3 2 3 3 3 4 C2 4 1 4 0 4 C0 2.7 0 1.4 0 0 Z };
\path[fill=black!43,even odd rule,shift={(339,80)}] svg {M0 0 C1.3 0 2.6 0 4 0 C4 1 4 2 4 3 C2.7 2.7 1.4 2.3 0 2 C0 1.3 0 0.7 0 0 Z };
\path[fill=black!49,even odd rule,shift={(317,76)}] svg {M0 0 C1.3 0.3 2.6 0.7 4 1 C4 1.7 4 2.3 4 3 C2.7 3 1.4 3 0 3 C0 2 0 1 0 0 Z };
\path[fill=black!49,even odd rule,shift={(277,56)}] svg {M0 0 C0.5 0.3 0.5 0.3 3 2 C3 2.7 3 3.3 3 4 C1.7 3.7 0.4 3.3 -1 3 C-0.7 2 -0.3 1 0 0 Z };
\path[fill=black!72,even odd rule,shift={(302,212)}] svg {M0 0 C0.7 0.3 1.3 0.7 2 1 C2 2 2 3 2 4 C1 4 0 4 -1 4 C-0.7 2.7 -0.3 1.4 0 0 Z };
\path[fill=black!58,even odd rule,shift={(107,139)}] svg {M0 0 C0 1 0 2 0 3 C-1.3 3 -2.6 3 -4 3 C-4 2.3 -4 1.7 -4 1 C-2 0 -2 0 0 0 Z };
\path[fill=black!65,even odd rule,shift={(64,136)}] svg {M0 0 C1.3 0.3 2.6 0.7 4 1 C4 1.7 4 2.3 4 3 C2.7 3 1.4 3 0 3 C0 2 0 1 0 0 Z };
\path[fill=black!65,even odd rule,shift={(59,134)}] svg {M0 0 C1.3 0.7 2.6 1.3 4 2 C2.7 2.3 1.4 2.7 0 3 C0 2 0 1 0 0 Z };
\path[fill=black!51,even odd rule,shift={(120,133)}] svg {M0 0 C1 0 2 0 3 0 C2.7 1 2.3 2 2 3 C1 3 0 3 -1 3 C-0.7 2 -0.3 1 0 0 Z };
\path[fill=black!51,even odd rule,shift={(151,107)}] svg {M0 0 C0.7 0 1.3 0 2 0 C1.7 1.3 1.3 2.6 1 4 C0.3 4 -0.3 4 -1 4 C-0.7 2.7 -0.3 1.4 0 0 Z };
\path[fill=black!65,even odd rule,shift={(158,98)}] svg {M0 0 C0.7 0 1.3 0 2 0 C1.7 1.3 1.3 2.6 1 4 C0 3.7 -1 3.3 -2 3 C-1.3 2 -0.7 1 0 0 Z };
\path[fill=black!58,even odd rule,shift={(161,93)}] svg {M0 0 C0.7 0 1.3 0 2 0 C1.7 1.3 1.3 2.6 1 4 C0.3 4 -0.3 4 -1 4 C-0.7 2.7 -0.3 1.4 0 0 Z };
\path[fill=black!51,even odd rule,shift={(334,79)}] svg {M0 0 C1 0 2 0 3 0 C3 1 3 2 3 3 C1.7 2.7 0.4 2.3 -1 2 C-0.7 1.3 -0.3 0.7 0 0 Z };
\path[fill=black!65,even odd rule,shift={(174,74)}] svg {M0 0 C0.7 0 1.3 0 2 0 C1.7 1.3 1.3 2.6 1 4 C0.3 4 -0.3 4 -1 4 C-0.7 2.7 -0.3 1.4 0 0 Z };
\path[fill=black!51,even odd rule,shift={(306,73)}] svg {M0 0 C1.3 0.3 2.6 0.7 4 1 C4 1.7 4 2.3 4 3 C2.7 2.7 1.4 2.3 0 2 C0 1.3 0 0.7 0 0 Z };
\path[fill=black!58,even odd rule,shift={(177,70)}] svg {M0 0 C1 0.3 2 0.7 3 1 C2 2 1 3 0 4 C0 2.7 0 1.4 0 0 Z };
\path[fill=black!58,even odd rule,shift={(181,65)}] svg {M0 0 C0.7 0 1.3 0 2 0 C1.7 1.3 1.3 2.6 1 4 C0.3 4 -0.3 4 -1 4 C-0.7 2.7 -0.3 1.4 0 0 Z };
\path[fill=black!58,even odd rule,shift={(263,46)}] svg {M0 0 C1 0.3 2 0.7 3 1 C3 2 3 3 3 4 C2 3.7 1 3.3 0 3 C0 2 0 1 0 0 Z };
\path[fill=black!65,even odd rule,shift={(206,42)}] svg {M0 0 C0.7 0 1.3 0 2 0 C2 1 2 2 2 3 C1 3.3 0 3.7 -1 4 C-1 1 -1 1 0 0 Z };
\path[fill=black!51,even odd rule,shift={(242,36)}] svg {M0 0 C1.3 0.3 2.6 0.7 4 1 C4 1.7 4 2.3 4 3 C2.7 2.7 1.4 2.3 0 2 C0 1.3 0 0.7 0 0 Z };
\path[fill=black!58,even odd rule,shift={(232,34)}] svg {M0 0 C1 0.3 2 0.7 3 1 C3 1.7 3 2.3 3 3 C1.7 3 0.4 3 -1 3 C-0.7 2 -0.3 1 0 0 Z };
\path[fill=black!62,even odd rule,shift={(98,140)}] svg {M0 0 C1 0 2 0 3 0 C3 1 3 2 3 3 C2 3 1 3 0 3 C0 2 0 1 0 0 Z };
\path[fill=black!62,even odd rule,shift={(130,127)}] svg {M0 0 C1 0.3 2 0.7 3 1 C2.7 1.7 2.3 2.3 2 3 C1 3 0 3 -1 3 C-0.7 2 -0.3 1 0 0 Z };
\path[fill=black!62,even odd rule,shift={(154,103)}] svg {M0 0 C1 0.3 2 0.7 3 1 C2.7 1.7 2.3 2.3 2 3 C1 3 0 3 -1 3 C-0.7 2 -0.3 1 0 0 Z };
\path[fill=black!69,even odd rule,shift={(164,88)}] svg {M0 0 C0.7 0.3 1.3 0.7 2 1 C1.7 2 1.3 3 1 4 C0.3 4 -0.3 4 -1 4 C-0.7 2.7 -0.3 1.4 0 0 Z };
\path[fill=black!62,even odd rule,shift={(167,84)}] svg {M0 0 C0.7 0 1.3 0 2 0 C1.7 1.3 1.3 2.6 1 4 C0.3 3.7 -0.3 3.3 -1 3 C-0.7 2 -0.3 1 0 0 Z };
\path[fill=black!54,even odd rule,shift={(171,79)}] svg {M0 0 C0.7 0.7 1.3 1.3 2 2 C1.3 2.3 1.3 2.3 -2 4 C-1.3 2.7 -0.7 1.4 0 0 Z };
\path[fill=black!54,even odd rule,shift={(323,77)}] svg {M0 0 C1 0.3 2 0.7 3 1 C3 1.7 3 2.3 3 3 C1.7 3 0.4 3 -1 3 C-0.7 2 -0.3 1 0 0 Z };
\path[fill=black!62,even odd rule,shift={(290,66)}] svg {M0 0 C1.3 0.3 2.6 0.7 4 1 C4 1.7 4 2.3 4 3 C3 3 2 3 1 3 C0.7 2 0.3 1 0 0 Z };
\path[fill=black!54,even odd rule,shift={(286,63)}] svg {M0 0 C1 0.3 2 0.7 3 1 C3 1.7 3 2.3 3 3 C1.7 2.7 0.4 2.3 -1 2 C-0.7 1.3 -0.3 0.7 0 0 Z };
\path[fill=black!62,even odd rule,shift={(197,49)}] svg {M0 0 C0.7 0 1.3 0 2 0 C1.3 1.3 0.7 2.6 0 4 C-0.7 3.3 -1.3 2.7 -2 2 C-1.3 1.3 -0.7 0.7 0 0 Z };
\path[fill=black!54,even odd rule,shift={(248,38)}] svg {M0 0 C1.3 0.3 2.6 0.7 4 1 C3.7 1.7 3.3 2.3 3 3 C2 2.7 1 2.3 0 2 C0 1.3 0 0.7 0 0 Z };
\path[fill=black!65,even odd rule,shift={(75,140)}] svg {M0 0 C1.3 0 2.6 0 4 0 C4 0.7 4 1.3 4 2 C2.7 2 1.4 2 0 2 C0 1.3 0 0.7 0 0 Z };
\path[fill=black!74,even odd rule,shift={(202,45)}] svg {M0 0 C0.3 1 0.7 2 1 3 C0 3.3 -1 3.7 -2 4 C-1.3 2.7 -0.7 1.4 0 0 Z };
\path[fill=black!74,even odd rule,shift={(258,43)}] svg {M0 0 C1 0.3 2 0.7 3 1 C3 1.7 3 2.3 3 3 C2 3 1 3 0 3 C0 2 0 1 0 0 Z };
\path[fill=black!56,even odd rule,shift={(210,40)}] svg {M0 0 C1 0 2 0 3 0 C3 0.7 3 1.3 3 2 C2 2.3 1 2.7 0 3 C0 2 0 1 0 0 Z };
\path[fill=black!90,even odd rule,shift={(81,140)}] svg {M0 0 C1 0.3 2 0.7 3 1 C3 1.7 3 2.3 3 3 C2 3 1 3 0 3 C0 2 0 1 0 0 Z };
\path[fill=black!80,even odd rule,shift={(272,53)}] svg {M0 0 C1 0.3 2 0.7 3 1 C3 1.7 3 2.3 3 3 C2 2.7 1 2.3 0 2 C0 1.3 0 0.7 0 0 Z };
\path[fill=black!80,even odd rule,shift={(226,35)}] svg {M0 0 C1 0 2 0 3 0 C3 0.7 3 1.3 3 2 C1.7 1.7 0.4 1.3 -1 1 C-0.7 0.7 -0.3 0.3 0 0 Z };
\path[fill=black!30,even odd rule,shift={(62,1)}] svg {M0 0 C0.7 0.3 1.3 0.7 2 1 C1.7 1.7 1.3 2.3 1 3 C0.3 2.7 -0.3 2.3 -1 2 C-0.7 1.3 -0.3 0.7 0 0 Z };
\path[fill=black!30,even odd rule,shift={(58,4)}] svg {M0 0 C0.7 0.7 1.3 1.3 2 2 C1 1.7 0 1.3 -1 1 C-0.7 0.7 -0.3 0.3 0 0 Z };
\path[fill=black!30,even odd rule,shift={(42,23)}] svg {M0 0 C2 1 2 1 2 1 Z };
\end{tikzpicture}}
\begin{abcd}
\item $t_0 = 10$．
\item $15 < t_0 < 20$．
\item $t_0 = 25$．
\item $t_0 > 25$．
\end{abcd}
\QuestionEnd

\begin{QuestionSolution}{2017 年第 6 题}
\textbf{已知与目标}\quad 初始甲领先 \(10\,\mathrm{m}\)，两条曲线给出双方速度。要求乙首次追上甲的时刻，即双方位置差首次变成零的时刻。

\textbf{方法与条件}\quad 速度积分给出位移，速度差的积分给出位置差的变化。有向面积指曲线在横轴上方记正、下方记负；这里对速度差 \(v_1-v_2\) 积分。

\textbf{推导}\quad 设 \(d(t)\) 为甲的位置减乙的位置，则
\[
d(t)=10+\int_0^t[v_1(s)-v_2(s)]\,\mathrm{d}s.
\]
在 \(0\) 到 \(10\) 秒内实线较高，第一块面积 \(10\) 增加甲的领先距离，所以 \(d(10)=10+10=20\)。在 \(10\) 到 \(25\) 秒内虚线较高，\(d\) 严格减少，第二块面积 \(20\) 恰好抵消此前领先距离，故 \(d(25)=20-20=0\)。该段末端以前只抵消了不足 \(20\) 米，因而 \(d(t)>0\)。

\textbf{结论与检查}\quad \(\boxed{t_0=25\,\mathrm{s}}\)，选 C。必须计入初始的 \(10\) 米；第三块面积发生在首次追上之后，不影响本问。
\end{QuestionSolution}

% Q 2011 11
\QuestionPage{3.4 定积分的应用}{填空题}{2011 年第 11 题}{1}{平面曲线弧长}{168 / 234}{**}
曲线$y = \int_0^x \tan t\,\mathrm{d} t$（$0 \le x \le \frac{\pi}{4}$）的弧长$s =$ \blank{}．
\QuestionEnd

\begin{QuestionSolution}{2011 年第 11 题}
\textbf{已知与目标}\quad 曲线由变上限积分给出，要计算它在 \([0,\frac{\pi}{4}]\) 上的弧长，无需求出 \(y\) 的表达式。

\textbf{方法与条件}\quad 连续被积函数使 \(y'=\tan x\)。弧长公式为 \(s=\int\sqrt{1+(y')^2}\,\mathrm{d}x\)；记 \(\sec x=\frac{1}{\cos x}\)，用 \(1+\tan^2x=\sec^2x\) 化简。

\textbf{推导}\quad 本区间 \(\cos x>0\)，所以 \(\sqrt{\sec^2x}=\sec x\)。又 \([\ln(\sec x+\tan x)]'=\sec x\)，因此
\[
s=\int_0^{\frac{\pi}{4}}\sqrt{1+\tan^2x}\,\mathrm{d}x
=\bigl[\ln(\sec x+\tan x)\bigr]_0^{\frac{\pi}{4}}
=\boxed{\ln(\sqrt2+1)} .
\]
上限 \(\sec\frac{\pi}{4}=\sqrt2\)、\(\tan\frac{\pi}{4}=1\)，下限的对数为 \(\ln1=0\)。

\textbf{结论与检查}\quad \(\boxed{s=\ln(\sqrt2+1)}\)。它约为 \(0.8814\)，大于水平距离 \(\frac{\pi}{4}\approx0.7854\)，符合弧长的几何范围。
\end{QuestionSolution}

% Q 2013 11
\QuestionPage{3.4 定积分的应用}{填空题}{2013 年第 11 题}{2}{极坐标曲线围成的面积}{169 / 234}{**}
设封闭曲线$L$的极坐标方程为$r = \cos 3\theta$（$- \frac{\pi}{6} \le \theta \le \frac{\pi}{6}$），
则$L$所围成的平面图形的面积为 \blank{}．
\QuestionEnd

\begin{QuestionSolution}{2013 年第 11 题}
\textbf{已知与目标}\quad 在给定角区间内 \(r=\cos3\theta\)，要求这段封闭曲线所围面积。

\textbf{方法与条件}\quad 极坐标下面积微元是小扇形面积 \(\frac{1}{2}r^2\,\mathrm{d}\theta\)。本区间 \(3\theta\in[-\frac{\pi}{2},\frac{\pi}{2}]\)，半径非负，两端为零，中间为正，恰好描出一瓣且不重复。

\textbf{推导}\quad \(\cos^2(3\theta)\) 是偶函数，可把对称区间积分化为正半段的两倍。再用平方降幂公式，
\[
S=\frac12\int_{-\frac{\pi}{6}}^{\frac{\pi}{6}}\cos^2(3\theta)\,\mathrm{d}\theta
=\int_0^{\frac{\pi}{6}}\frac{1+\cos 6\theta}{2}\,\mathrm{d}\theta
=\boxed{\frac{\pi}{12}}.
\]
最后的原函数为 \(\frac{\theta}{2}+\frac{\sin6\theta}{12}\)，在 \(\frac{\pi}{6}\) 处正弦为零，在零点也为零。

\textbf{结论与检查}\quad \(\boxed{S=\frac{\pi}{12}}\)。积分只计算题设角区间的一瓣，没有乘整条玫瑰线的瓣数。
\end{QuestionSolution}

% Q 2014 13
\QuestionPage{3.4 定积分的应用}{填空题}{2014 年第 13 题}{2}{细棒的质心}{170 / 234}{**}
一根长为$1$的细棒位于$x$轴的区间$[0,1]$上，若其线密度$\rho (x) = - x^2 + 2x + 1$，
则该细棒的质心坐标$\overline{x} =$ \blank{}．
\QuestionEnd

\begin{QuestionSolution}{2014 年第 13 题}
\textbf{已知与目标}\quad 细棒位于 \([0,1]\)，线密度 \(\rho(x)=-x^2+2x+1\)，求它的质心横坐标。

\textbf{方法与条件}\quad 线密度是单位长度的质量；一阶矩指位置乘质量的总和。质心是按质量加权的位置平均值，故 \(\overline{x}=\frac{\int_0^1x\rho(x)\,\mathrm{d}x}{\int_0^1\rho(x)\,\mathrm{d}x}\)。本区间 \(\rho(x)=2-(x-1)^2\ge1\)，总质量为正。

\textbf{推导}\quad 总质量为
\[
M=\left[-\frac{x^3}{3}+x^2+x\right]_0^1=-\frac{1}{3}+1+1=\frac{5}{3}.
\]
一阶矩为
\[
\int_0^1x(-x^2+2x+1)\,\mathrm{d}x=-\frac14+\frac23+\frac12=\frac{11}{12}.
\]
相除得 \(\overline{x}=\frac{11}{12}\div\frac{5}{3}=\frac{11}{20}\)。

\textbf{结论与检查}\quad \(\boxed{\overline{x}=\frac{11}{20}}\)。它位于棒内部并略大于中点 \(\frac{1}{2}\)，符合密度向右增大的分布。
\end{QuestionSolution}

% Q 2019 12
\QuestionPage{3.4 定积分的应用}{填空题}{2019 年第 12 题}{2}{对数曲线的弧长}{171 / 234}{**}
曲线$y = \ln \cos x$（$0 \le x \le \frac{\pi}{6}$）的弧长为 \blank{}．
\QuestionEnd

\begin{QuestionSolution}{2019 年第 12 题}
\textbf{已知与目标}\quad 求 \(y=\ln\cos x\) 在 \([0,\frac{\pi}{6}]\) 上的长度；该区间 \(\cos x>0\)，函数连续可导。

\textbf{方法与条件}\quad 弧长是 \(\int\sqrt{1+(y')^2}\,\mathrm{d}x\)。\(y'=-\tan x\)，由 \(1+\tan^2x=\sec^2x\)，其中 \(\sec x=\frac{1}{\cos x}>0\)，根式可化成 \(\sec x\)。

\textbf{推导}\quad 利用 \(\int\sec x\,\mathrm{d}x=\ln(\sec x+\tan x)+C\)，
\[
s=\int_0^{\frac{\pi}{6}}\sec x\,\mathrm{d}x
=\bigl[\ln(\sec x+\tan x)\bigr]_0^{\frac{\pi}{6}}
=\ln\left(\frac2{\sqrt3}+\frac1{\sqrt3}\right)
=\boxed{\frac{1}{2}\ln3}.
\]
这里 \(\sec\frac{\pi}{6}=\frac{2}{\sqrt3}\)、\(\tan\frac{\pi}{6}=\frac{1}{\sqrt3}\)，相加为 \(\sqrt3\)；下限为 \(\ln1=0\)。

\textbf{结论与检查}\quad 弧长为 \(\boxed{\frac{1}{2}\ln3}\)，约 \(0.5493\)，大于水平距离 \(\frac{\pi}{6}\approx0.5236\)，符合几何关系。
\end{QuestionSolution}

% Q 2022 15
\QuestionPage{3.4 定积分的应用}{填空题}{2022 年第 15 题}{2}{极坐标曲线围成的面积}{173 / 234}{**}
已知曲线 \(L\) 的极坐标方程为
\(r = \sin 3\theta \quad \left( 0 \le \theta \le \frac{\pi}{3} \right)\)
则 \(L\) 围成的有界区域的面积为 \(\blank{}\)
\QuestionEnd

\begin{QuestionSolution}{2022 年第 15 题}
\textbf{已知与目标}\quad 给定 \(r=\sin3\theta\)、\(0\le\theta\le\frac{\pi}{3}\)，求对应封闭区域的面积。

\textbf{方法与条件}\quad 该区间 \(3\theta\in[0,\pi]\)，半径两端为零，内部为正，描出一瓣且不重复。极坐标面积用小扇形面积 \(\frac{1}{2}r^2\,\mathrm{d}\theta\) 积分。

\textbf{推导}\quad 代入半径，再用 \(\sin^2u=\frac{1-\cos2u}{2}\)，
\[
S=\frac12\int_0^{\frac{\pi}{3}}\sin^2(3\theta)\,\mathrm{d}\theta
=\frac14\int_0^{\frac{\pi}{3}}(1-\cos6\theta)\,\mathrm{d}\theta
=\boxed{\frac{\pi}{12}}.
\]
原函数为 \(\frac{\theta}{4}-\frac{\sin6\theta}{24}\)，代入上限 \(\frac{\pi}{3}\) 与下限 \(0\)，正弦项均为零，留下 \(\frac{\pi}{12}\)。

\textbf{结论与检查}\quad \(\boxed{S=\frac{\pi}{12}}\)。本题只给出一个花瓣的角度范围，所以积分中无需额外乘瓣数。
\end{QuestionSolution}

% Q 2013 16
\QuestionPage{3.4 定积分的应用}{解答题}{2013 年第 16 题}{2}{两个旋转体的体积关系}{178 / 234}{**}
设$D$是由曲线$y = x^{\frac{1}{3}}$，直线$x = a$（$a > 0$）及$x$轴所围成的平面图形，
$V_x$，$V_y$分别是$D$绕$x$轴，$y$轴旋转一周所得旋转体的体积，若$V_y = 10V_x$，求$a$的值．
\QuestionEnd

\begin{QuestionSolution}{2013 年第 16 题}
\textbf{已知与目标}\quad 区域为 \(0\le x\le a\)、\(0\le y\le x^{\frac{1}{3}}\)，\(a>0\)。求使绕两坐标轴的体积满足 \(V_y=10V_x\) 的 \(a\)。

\textbf{方法与条件}\quad 竖条绕 \(x\) 轴形成圆盘，体积微元为 \(\pi y^2\,\mathrm{d}x\)；绕 \(y\) 轴形成圆柱壳，体积微元为周长 \(2\pi x\) 乘高度 \(y\) 再乘厚度 \(\mathrm{d}x\)。

\textbf{推导}\quad 两种体积分别为
\[
V_x=\pi\int_0^a x^{\frac{2}{3}}\,\mathrm{d}x=\frac{3\pi}{5}a^{\frac{5}{3}},\qquad
V_y=2\pi\int_0^a x^{\frac{4}{3}}\,\mathrm{d}x=\frac{6\pi}{7}a^{\frac{7}{3}}.
\]
用幂函数积分公式 \(\int_0^ax^p\,\mathrm{d}x=\frac{a^{p+1}}{p+1}\) 得到上述值。因 \(a>0\)，可除以 \(V_x>0\)，得 \(\frac{V_y}{V_x}=\frac{10}{7}a^{\frac{2}{3}}=10\)，所以 \(a^{\frac{2}{3}}=7\)，\(a=7^{\frac{3}{2}}=7\sqrt7\)。

\textbf{结论与检查}\quad \(\boxed{a=7\sqrt7}\)。代回体积比为 \(\frac{10}{7}\times7=10\)，且 \(a>0\) 满足题设。
\end{QuestionSolution}

% Q 2015 16
\QuestionPage{3.4 定积分的应用}{解答题}{2015 年第 16 题}{2}{旋转体积的比值}{181 / 234}{**}
设$A > 0$，$D$是由曲线段$y = A\sin x$（$0 \le x \le \frac{\pi}{2}$）
及直线$y = 0$，$x = \frac{\pi}{2}$所围成的平面区域，
$V_1$，$V_2$分别表示$D$绕$x$轴与绕$y$轴旋转成旋转体的体积，若$V_1 = V_2$，求$A$的值．
\QuestionEnd

\begin{QuestionSolution}{2015 年第 16 题}
\textbf{已知与目标}\quad 区域为正弦曲线下方 \(0\le x\le\frac{\pi}{2}\) 的部分，参数 \(A>0\)。要求绕两坐标轴所得体积相等时的 \(A\)。

\textbf{方法与条件}\quad 竖条绕 \(x\) 轴形成圆盘，半径为 \(A\sin x\)；绕 \(y\) 轴形成圆柱壳，其周长为 \(2\pi x\)、高度为 \(A\sin x\)。因区域在第一象限，半径和高度都非负。

\textbf{推导}\quad 绕 \(x\) 轴时，用 \(\sin^2x=\frac{1-\cos2x}{2}\)，得
\[
V_1=\pi A^2\int_0^{\frac{\pi}{2}}\sin^2x\,\mathrm{d}x=\frac{\pi^2 A^2}{4}.
\]
绕 \(y\) 轴时，壳半径为 \(x\)、高为 \(A\sin x\)，故
\[
V_2=2\pi A\int_0^{\frac{\pi}{2}}x\sin x\,\mathrm{d}x
=2\pi A[-x\cos x+\sin x]_0^{\frac{\pi}{2}}=2\pi A.
\]
令 \(\frac{\pi^2A^2}{4}=2\pi A\)，由于 \(A>0\)，可除以 \(\pi A\)，得到 \(\frac{\pi A}{4}=2\)，故 \(A=\frac{8}{\pi}\)。

\textbf{结论与检查}\quad \(\boxed{A=\frac{8}{\pi}}\)。此时 \(V_1=\frac{\pi^2}{4}\frac{64}{\pi^2}=16\)，\(V_2=2\pi\frac{8}{\pi}=16\)，满足条件。
\end{QuestionSolution}

% Q 2025 6
\QuestionPage{3.4 定积分的应用}{选择题}{2025 年第 6 题}{1}{克服引力做功}{165 / 234}{***}
设单位质点 \(P\)，\(Q\) 分别位于点 \(\left( 0, 0 \right)\) 和 \(\left( 0, 1 \right)\) 处，\(P\) 从点 \(\left( 0, 0 \right)\) 出发沿 \(x\) 轴正向移动，记 \(G\) 为引力常量，则当质点 \(P\) 移动到点 \(\left( 1, 0 \right)\) 时，克服质点 \(Q\) 的引力所做的功为
\choices{\(\int_{0}^{1} \frac{G}{x^{2} + 1} \,\mathrm{d} x\)}{\(\int_{0}^{1} \frac{Gx}{\left( x^{2} + 1 \right)^{\frac{3}{2}}} \,\mathrm{d} x\)}{\(\int_{0}^{1} \frac{G}{\left( x^{2} + 1 \right)^{\frac{3}{2}}} \,\mathrm{d} x\)}{\(\int_{0}^{1} \frac{G\left( x + 1 \right)}{\left( x^{2} + 1 \right)^{\frac{3}{2}}} \,\mathrm{d} x\)}
\QuestionEnd

\begin{QuestionSolution}{2025 年第 6 题}
\textbf{已知与目标}\quad 两质点的质量均为 \(1\)，\(P\) 沿 \(x\) 轴从零点移到 \((1,0)\)。求克服固定质点 \(Q=(0,1)\) 的引力所做的功。

\textbf{方法与条件}\quad 引力大小为 \(\frac{G}{r^2}\)，其中 \(r\) 为两点距离。功等于沿位移方向的力分量乘微小位移后积分，竖直分力因没有竖直位移而不做功。

\textbf{推导}\quad 在 \(P=(x,0)\) 处，\(r=\sqrt{x^2+1}\)，从 \(P\) 指向 \(Q\) 的单位向量为 \(\frac{1}{r}(-x,1)\)。因此引力水平分量为 \(-\frac{Gx}{r^3}\)，克服它的外力水平分量为正的 \(\frac{Gx}{r^3}\)。功为
\[
W=\int_0^1\frac{Gx}{(1+x^2)^{\frac{3}{2}}}\,\mathrm{d}x
=G\left(1-\frac1{\sqrt2}\right),
\]
换元 \(u=1+x^2\)，\(\mathrm{d}u=2x\,\mathrm{d}x\)，便得 \(W=[-G(1+x^2)^{-\frac{1}{2}}]_0^1\)。

\textbf{结论与检查}\quad 选 \(\boxed{\mathrm{B}}\)，功为 \(G(1-\frac{1}{\sqrt2})\)。运动中两点距离增大，需要正功，结果为正；在 \(x=0\) 引力完全竖直，水平分量确为零。
\end{QuestionSolution}

% Q 2026 4
\QuestionPage{3.4 定积分的应用}{选择题}{2026 年第 4 题}{1}{连续分布质点的引力}{166 / 234}{***}
设线密度为 \(1\) 的细直棒两端点分别位于 \((-1,0)\) 和 \((1,0)\)，质量为 \(m\) 的质点位于 \((0,1)\)，\(G\) 为引力常量，则该细直棒对该质点的引力大小为
\choices{\(\displaystyle \int_0^1 \frac{2Gmx}{(x^2+1)^{\frac{1}{2}}}\,\mathrm{d}x\)}{\(\displaystyle \int_0^1 \frac{2Gm}{(x^2+1)^{\frac{1}{2}}}\,\mathrm{d}x\)}{\(\displaystyle \int_0^1 \frac{2Gmx}{(x^2+1)^{\frac{3}{2}}}\,\mathrm{d}x\)}{\(\displaystyle \int_0^1 \frac{2Gm}{(x^2+1)^{\frac{3}{2}}}\,\mathrm{d}x\)}
\QuestionEnd

\begin{QuestionSolution}{2026 年第 4 题}
\textbf{已知与目标}\quad 线密度为 \(1\) 的棒位于 \([-1,1]\) 上，质点在 \((0,1)\)。求整个棒产生的合引力大小，而非各小段引力大小的直接相加。

\textbf{方法与条件}\quad 线密度是单位长度的质量，所以长度 \(\mathrm{d}x\) 的棒元质量为 \(\mathrm{d}x\)。利用左右对称使水平分力抵消，再积分竖直分力。方向余弦是力方向与指定轴夹角的余弦，用它求该轴上的投影。

\textbf{推导}\quad 取 \(x\in[0,1]\) 的棒元，到质点的距离为 \(r=\sqrt{1+x^2}\)，引力大小为 \(\frac{Gm}{1+x^2}\mathrm{d}x\)。竖直距离为 \(1\)，故向下分量需再乘 \(\frac{1}{r}\)。与 \(-x\) 处配对，两个水平分量相反、竖直分量相同，所以
\[
2Gm\int_0^1\frac{\mathrm{d}x}{(x^2+1)^{\frac{3}{2}}}
=\sqrt2\,Gm .
\]
计算时可用原函数 \(\frac{x}{\sqrt{1+x^2}}\)，故 \(F=2Gm[\frac{x}{\sqrt{1+x^2}}]_0^1=\sqrt2Gm\)。

\textbf{结论与检查}\quad 积分式选 \(\boxed{\mathrm{D}}\)，合力大小为 \(\sqrt2Gm\)，方向竖直向下。被积函数中没有多余的 \(x\) 因子，因为求的是竖直分量，其竖直距离恒为 \(1\)。
\end{QuestionSolution}

% Q 2020 12
\QuestionPage{3.4 定积分的应用}{填空题}{2020 年第 12 题}{2}{水压力的微元法}{172 / 234}{***}
斜边长为 \(2a\) 的等腰直角三角形平板铅直地沉没在水中，且斜边与水面相齐。记重力加速度为 \(g\)，水的密度为 \(\rho\)，则该平板一侧所受的水压力为
\blank{}
\QuestionEnd

\begin{QuestionSolution}{2020 年第 12 题}
\textbf{已知与目标}\quad 等腰直角三角形平板斜边长 \(2a\)，斜边在水面，直角顶点在下方，求水对平板一侧的总压力。

\textbf{方法与条件}\quad 静水在深度 \(h\) 的压强为 \(\rho gh\)。取水平窄条，条内压强可用同一深度的值作微元主部，压力微元为压强乘面积；条宽由相似三角形确定。

\textbf{推导}\quad 等腰直角三角形斜边中点到直角顶点的距离为斜边一半，即 \(a\)。深度 \(h\in[0,a]\) 处，条宽与剩余高度 \(a-h\) 成比例，故宽为 \(2a\frac{a-h}{a}=2(a-h)\)。面积微元为 \(2(a-h)\,\mathrm{d}h\)，所以
\[
F=\int_0^a \rho gh\,2(a-h)\,\mathrm{d}h
=2\rho g\left(\frac{a^3}{2}-\frac{a^3}{3}\right)
=\boxed{\frac{\rho ga^3}{3}}.
\]
\textbf{结论与检查}\quad 单侧总压力为 \(\boxed{\frac{\rho ga^3}{3}}\)。也可用面积 \(a^2\) 乘形心深度 \(\frac{a}{3}\) 处的压强 \(\frac{\rho ga}{3}\) 核对，结果相同。
\end{QuestionSolution}

% Q 2023 12
\QuestionPage{3.4 定积分的应用}{填空题}{2023 年第 12 题}{2}{由积分表示的曲线弧长}{174 / 234}{***}
曲线
\[y=\int_{-\sqrt3}^{x}\sqrt{3-t^2}\,\mathrm{d}t
\]
的弧长为
\[\blank{}.
\]
\QuestionEnd

\begin{QuestionSolution}{2023 年第 12 题}
\textbf{已知与目标}\quad 曲线由根式积分给出，题目未另给横坐标范围，需先确定整个实数定义域，再求弧长。

\textbf{方法与条件}\quad 被积函数要求 \(3-t^2\ge0\)，故变上限只能在 \([-\sqrt3,\sqrt3]\) 内。连续性保证 \(y'=\sqrt{3-x^2}\)，用弧长公式 \(s=\int\sqrt{1+(y')^2}\,\mathrm{d}x\)。

\textbf{推导}\quad 弧长被积函数为 \(\sqrt{1+3-x^2}=\sqrt{4-x^2}\)，它是偶函数，所以
\[
s=2\int_0^{\sqrt3}\sqrt{4-x^2}\,\mathrm{d}x
=\left[x\sqrt{4-x^2}+4\arcsin(\frac{x}{2})\right]_0^{\sqrt3}
=\boxed{\sqrt3+\frac{4\pi}{3}} .
\]
该原函数可由 \(x=2\sin u\) 得到。代入 \(x=\sqrt3\)，\(\sqrt{4-x^2}=1\)、\(\arcsin\frac{\sqrt3}{2}=\frac{\pi}{3}\)；下限为零。

\textbf{结论与检查}\quad \(\boxed{s=\sqrt3+\frac{4\pi}{3}}\)。弧长在 \(2\sqrt3\) 与 \(4\sqrt3\) 之间，因为弧长被积函数在 \([1,2]\) 内；所得值约 \(5.9208\)，符合该范围。
\end{QuestionSolution}

% Q 2010 18
\QuestionPage{3.4 定积分的应用}{解答题}{2010 年第 18 题}{3}{液体质量的积分模型}{175 / 234}{***}
一个高为$l$的柱体形贮油罐，底面是长轴为$2a$，短轴为$2b$的椭圆．
现将贮油罐平放，当油罐中油面高度为$\frac{3}{2}b$时（如图），计算油的质量．
\par\centerline{\begin{tikzpicture}[scale=\DiagramScale,x=1pt,y=1pt,xscale=0.439724,yscale=-0.439724]
\path[use as bounding box] (0,0) rectangle (440,280);
\path[fill=black!80,even odd rule,shift={(204.0625,46.8125)}] svg {M0 0 C0.6 -0 0.6 -0 3.4 -0 C6.5 -0 9.7 0 12.8 0 C13.3 0 13.3 0 15.8 0 C16.2 0 16.2 0 18.4 0 C18.8 0 18.8 0 20.8 0.1 C22.9 0.2 22.9 0.2 25.1 0.7 C27.4 1.3 28.7 0.8 30.9 0.2 C31.9 0.5 32.9 0.8 33.9 1.2 C34.6 1 35.3 0.8 35.9 0.6 C36.6 0.5 37.3 0.3 37.9 0.2 C38.9 0.8 39.9 1.5 40.9 2.2 C41.3 1.5 41.6 0.9 41.9 0.2 C43.6 0.5 45.2 0.8 46.9 1.2 C46.9 1.8 46.9 2.5 46.9 3.2 C47.3 3.4 47.3 3.4 48.9 4.2 C48.9 2.9 48.9 1.5 48.9 0.2 C50.6 0.5 52.2 0.8 53.9 1.2 C53.6 2.2 53.3 3.2 52.9 4.2 C53.2 4.2 53.2 4.2 54.7 4.6 C83.8 10.6 113.6 16.9 138.9 33.2 C139.8 33.7 140.6 34.2 141.4 34.7 C144.2 36.4 146.9 38.2 149.5 40.1 C150.4 40.8 151.4 41.5 152.3 42.2 C154.7 44 156.8 46 158.9 48.2 C159.3 47.5 159.6 46.9 159.9 46.2 C161.6 46.2 163.2 46.2 164.9 46.2 C164.9 47.2 164.9 48.2 164.9 49.2 C164.4 49.4 164.4 49.4 161.9 50.2 C162.5 50.9 163.1 51.7 163.7 52.5 C172.9 64.3 178.3 74.7 178.2 90 C178.2 91 178.2 91.9 178.2 92.9 C178.2 93.3 178.2 93.3 178.2 95.6 C178.2 96.5 178.2 97.3 178.2 98.1 C177.9 100.2 177.9 100.2 175.9 102.2 C176.6 102.8 177.3 103.5 177.9 104.2 C177.9 105.5 177.9 106.8 177.9 108.2 C177.6 108.1 177.6 108.1 175.6 107.9 C172.9 108.2 172.9 108.2 171.4 109.8 C171 110.6 170.6 111.3 170.1 112.1 C150.5 144 107.5 157.5 73.2 165.7 C68.1 166.8 63.1 167.7 58 168.6 C54.9 169.2 54.9 169.2 52.9 170.2 C52.3 172.3 52.3 172.3 51.9 174.2 C50.6 173.9 49.3 173.5 47.9 173.2 C47.9 172.2 47.9 171.2 47.9 170.2 C47.3 170.2 46.6 170.2 45.9 170.2 C45.6 171.5 45.3 172.8 44.9 174.2 C44.3 174 43.7 173.9 43.1 173.7 C40.6 173.1 38.5 173 36 173.7 C33.9 174.2 33.9 174.2 29.9 174.2 C27.8 174.2 25.7 174.2 23.5 174.3 C-20.2 174.7 -62.2 174.9 -105.1 164.2 C-105.9 164 -106.8 163.8 -107.7 163.5 C-118.5 160.8 -128.8 157.4 -139.1 153.2 C-139.8 152.9 -140.5 152.6 -141.2 152.3 C-157.1 145.9 -174.6 137.3 -186.1 124.2 C-186.1 123.5 -186.1 122.9 -186.1 122.2 C-186.7 122.2 -187.4 122.2 -188.1 122.2 C-189.5 120.4 -190.8 118.6 -192.1 116.8 C-192.8 115.8 -193.6 114.8 -194.3 113.8 C-196.1 111.2 -196.1 111.2 -197.1 109.2 C-197.7 109.4 -197.7 109.4 -201.1 110.2 C-201.4 109.2 -201.7 108.2 -202.1 107.2 C-201.4 106.2 -200.7 105.2 -200.1 104.2 C-200.7 103.2 -201.4 102.2 -202.1 101.2 C-203.5 90.6 -203.2 78.7 -198.1 69.2 C-197.6 68.4 -197.2 67.5 -196.8 66.7 C-193.2 60.1 -188.5 54.3 -183.1 49.2 C-182.4 49.2 -181.7 49.2 -181.1 49.2 C-181.1 48.5 -181.1 47.9 -181.1 47.2 C-171.3 38.2 -159.1 31.5 -147.1 26.2 C-146.4 25.9 -145.7 25.6 -145 25.3 C-116.4 12.7 -86.1 5.6 -55.1 2.2 C-54 2.1 -52.9 2 -51.8 1.8 C-34.6 0.1 -17.3 0 0 0 Z M-171.1 45.2 C-170.7 45.8 -170.4 46.5 -170.1 47.2 C-166.8 47.2 -163.6 47.2 -160.3 47.2 C-159.9 47.3 -159.9 47.3 -157.5 47.3 C-156.6 47.3 -155.7 47.3 -154.8 47.3 C-154.4 47.3 -154.4 47.3 -152.4 47.3 C-150.1 47.2 -150.1 47.2 -147.8 46.6 C-147.2 46.5 -146.7 46.3 -146.1 46.2 C-145.9 46.4 -145.9 46.4 -145.1 47.2 C-143.4 47.3 -141.7 47.3 -140.1 47.3 C-139 47.3 -138 47.3 -136.9 47.3 C-136.3 47.3 -136.3 47.3 -133.3 47.3 C-132.7 47.3 -132.7 47.3 -129.7 47.3 C-125.8 47.2 -121.9 47.2 -118.1 47.2 C-116.9 47.2 -115.7 47.2 -114.5 47.2 C-107.2 47.1 -99.9 47.1 -92.6 47 C-88.5 47 -84.4 46.9 -80.3 46.9 C-79.6 46.9 -78.9 46.9 -78.2 46.9 C-76.7 46.9 -75.2 46.9 -73.8 46.8 C-71.6 46.8 -69.4 46.8 -67.1 46.8 C14.7 46.1 14.7 46.1 49.8 46.2 C50.4 46.2 50.4 46.2 53.2 46.2 C63.7 46.2 74.2 46.2 84.7 46.2 C89.3 46.2 93.9 46.2 98.5 46.2 C99.4 46.3 100.2 46.3 101.1 46.3 C116.4 46.3 131.7 45.9 146.9 45.2 C124.1 21.7 81.8 13.8 50.8 8.9 C49.9 8.8 48.9 8.6 47.9 8.5 C29 5.7 10 4.8 -9.1 4.9 C-9.6 4.9 -9.6 4.9 -12.3 4.9 C-65.1 5 -127.4 12.4 -171.1 45.2 Z M-27.4 50.7 C-44.3 51.1 -61.1 51.3 -77.9 51.2 C-82.1 51.1 -86.2 51.1 -90.3 51.1 C-101.9 51.1 -113.6 51 -125.2 50.9 C-131.8 50.9 -138.3 50.9 -144.8 50.9 C-148.2 50.8 -151.7 50.8 -155.1 50.8 C-158.9 50.8 -162.6 50.8 -166.4 50.7 C-167.5 50.7 -168.6 50.7 -169.8 50.7 C-177.3 50.7 -177.3 50.7 -181.2 53.1 C-183.3 55.5 -185.4 57.9 -187.4 60.4 C-188.1 61.3 -188.8 62.2 -189.5 63.1 C-190 63.8 -190.5 64.5 -191.1 65.2 C-190.1 66.8 -189.1 68.5 -188.1 70.2 C-189.7 70.2 -191.4 70.2 -193.1 70.2 C-196.1 75.5 -197.9 80 -198.1 86.2 C-197.4 87.2 -196.7 88.2 -196.1 89.2 C-196.4 89.4 -196.4 89.4 -198.1 90.2 C-196.4 102.1 -196.4 102.1 -192.1 107.2 C-191.4 107.2 -190.7 107.2 -190.1 107.2 C-190.3 107.7 -190.5 108.3 -190.8 108.9 C-191.1 111.2 -191.1 111.2 -189.8 113.5 C-184.8 119.6 -179.4 124.5 -173.1 129.2 C-172.7 128.5 -172.4 127.9 -172.1 127.2 C-170.4 127.2 -168.8 127.2 -167.1 127.2 C-167.1 128.2 -167.1 129.2 -167.1 130.2 C-167.7 130.4 -167.7 130.4 -171.1 131.2 C-169.8 132 -168.5 132.9 -167.2 133.7 C-166.5 134.2 -165.8 134.6 -165.1 135.1 C-163.1 136.2 -163.1 136.2 -160.1 136.2 C-160.1 136.8 -160.1 137.5 -160.1 138.2 C-157.5 139.5 -155 140.8 -152.4 142.1 C-151.7 142.4 -151 142.8 -150.3 143.2 C-139.5 148.6 -139.5 148.6 -135.2 147.3 C-134.5 146.9 -133.8 146.6 -133.1 146.2 C-132.7 147.8 -132.4 149.5 -132.1 151.2 C-131.5 151.4 -130.9 151.5 -130.3 151.7 C-129.9 151.8 -129.9 151.8 -127.9 152.4 C-127.4 152.6 -127.4 152.6 -125.1 153.3 C-123.1 153.9 -121.1 154.5 -119.1 155.1 C-97.7 161.9 -75.6 167.2 -53.1 167.2 C-50.7 167.5 -48.4 167.8 -46.1 168.2 C-46.1 167.5 -46.1 166.9 -46.1 166.2 C-45.5 166.4 -45.5 166.4 -42.9 167.3 C-38.9 168.2 -37.6 168.1 -34.1 166.2 C-33.7 167.2 -33.4 168.2 -33.1 169.2 C-32.4 168.5 -31.7 167.9 -31.1 167.2 C-28.9 166.6 -28.9 166.6 -27.1 166.2 C-26.4 167.5 -25.7 168.8 -25.1 170.2 C-24.7 168.9 -24.4 167.5 -24.1 166.2 C-22.7 166.5 -21.4 166.8 -20.1 167.2 C-20.1 168.2 -20.1 169.2 -20.1 170.2 C-19.1 170.2 -18.1 170.2 -17.1 170.2 C-17.1 168.9 -17.1 167.5 -17.1 166.2 C-15.4 166.2 -13.8 166.2 -12.1 166.2 C-11.7 167.2 -11.4 168.2 -11.1 169.2 C-11.1 168.5 -11.1 167.9 -11.1 167.2 C-9.1 167.2 -7.1 167.2 -5.1 167.2 C-5.1 168.2 -5.1 169.2 -5.1 170.2 C-4.1 169.9 -3.1 169.5 -2.1 169.2 C-2.1 168.2 -2.1 167.2 -2.1 166.2 C-0.4 166.2 1.2 166.2 2.9 166.2 C2.9 167.2 2.9 168.2 2.9 169.2 C3.6 169.2 4.3 169.2 4.9 169.2 C4.9 168.2 4.9 167.2 4.9 166.2 C6.6 166.2 8.2 166.2 9.9 166.2 C9.9 167.2 9.9 168.2 9.9 169.2 C10.6 169.2 11.3 169.2 11.9 169.2 C12.3 168.5 12.6 167.9 12.9 167.2 C15 166.6 15 166.6 16.9 166.2 C16.9 166.8 16.9 167.5 16.9 168.2 C17.9 168.2 18.9 168.2 19.9 168.2 C19.9 167.5 19.9 166.9 19.9 166.2 C21.6 166.2 23.2 166.2 24.9 166.2 C25.3 166.8 25.6 167.5 25.9 168.2 C26.2 168 26.2 168 27.6 167.1 C30.5 165.9 31.9 166.6 34.9 167.2 C37.3 167.1 37.3 167.1 39.6 166.8 C40.6 166.7 41.5 166.6 42.4 166.5 C43.4 166.3 44.3 166.2 45.3 166.1 C45.8 166 45.8 166 48.4 165.6 C64.2 163.4 79.5 159.8 94.8 155 C95.9 154.7 97 154.3 98.1 154 C99.1 153.6 100.1 153.3 101.2 153 C102.1 152.7 103 152.4 103.9 152.1 C105.9 151.2 105.9 151.2 106.9 149.2 C105.9 149.2 105 149.2 103.9 149.2 C104.3 148.2 104.6 147.2 104.9 146.2 C105.9 146.2 106.9 146.2 107.9 146.2 C108.3 147.2 108.6 148.2 108.9 149.2 C109.6 149.2 110.3 149.2 110.9 149.2 C111.3 148.2 111.6 147.2 111.9 146.2 C113.2 146.1 114.5 146.1 115.8 146 C121.8 145.3 126.4 142.8 131.6 139.9 C132.6 139.3 133.5 138.8 134.5 138.3 C142.9 133.5 150 127.9 156.9 121.2 C157.7 120.5 158.5 119.8 159.3 119.1 C162.5 116 164.7 113.6 165.9 109.2 C164.9 108.9 164 108.5 162.9 108.2 C163.3 107.5 163.6 106.9 163.9 106.2 C165.6 106.2 167.2 106.2 168.9 106.2 C171.8 100.4 173.5 95.6 173.9 89.2 C174 88.1 174.1 87.1 174.2 86 C173.8 80.2 172.4 74.5 169.9 69.2 C167.8 68.3 167.8 68.3 165.9 68.2 C166 67.3 166.1 66.5 166.2 65.6 C165.8 59.9 162.2 56.5 158.3 52.7 C152.8 48.1 146.1 48.6 139.3 48.8 C137.7 48.8 136.2 48.8 134.6 48.8 C130.4 48.8 126.2 48.9 122.1 48.9 C117.6 49 113.2 49 108.7 49 C100.3 49.1 91.9 49.2 83.4 49.3 C74 49.4 64.5 49.4 55.1 49.5 C27.6 49.6 0.1 49.9 -27.4 50.7 Z };
\path[fill=black!72,even odd rule,shift={(26,245)}] svg {M0 0 C1 4 1 4 0 6 C36.6 6 73.3 6 111 6 C111 7 111 8 111 9 C108.1 9.8 105.5 10.1 102.5 10.2 C101.6 10.2 100.6 10.2 99.7 10.2 C98.7 10.2 97.7 10.2 96.7 10.2 C95.7 10.2 94.6 10.2 93.6 10.2 C90.1 10.3 86.7 10.3 83.2 10.3 C80.8 10.3 78.4 10.4 76 10.4 C69.7 10.4 63.4 10.5 57.1 10.5 C50.7 10.6 44.3 10.6 37.8 10.7 C25.2 10.8 12.6 10.9 0 11 C0.3 12.6 0.7 14.3 1 16 C-5.4 14.3 -11.1 12.1 -17 9 C-17 8.3 -17 7.7 -17 7 C-11.4 4.7 -5.8 2.4 0 0 Z };
\path[fill=black!82,even odd rule,shift={(361,243)}] svg {M0 0 C6 2 12 4 18 6 C17.9 7.8 17.9 7.8 17 10 C14.7 11.3 12.6 12.3 10.2 13.3 C9.6 13.6 8.9 13.9 8.3 14.2 C5.5 15.3 3.5 16.1 0.6 15.7 C0 15.4 -0.5 15.2 -1 15 C-0.3 13.7 0.3 12.4 1 11 C-40.2 10.7 -81.5 10.3 -124 10 C-123.7 9 -123.3 8 -123 7 C-96.2 6.1 -69.4 5.9 -42.5 5.9 C-37.6 5.9 -32.7 5.9 -27.7 6 C-18.1 6 -8.6 6 1 6 C0.7 5.2 0.7 5.2 -1 1 C-0.7 0.7 -0.3 0.3 0 0 Z };
\path[fill=black!71,even odd rule,shift={(416,188)}] svg {M0 0 C1 0.3 2 0.7 3 1 C3.2 2.5 3.2 2.5 4 10 C5.3 9.3 6.6 8.7 8 8 C9.1 11.4 8.9 11.9 7.6 15 C7.2 15.8 6.9 16.5 6.6 17.3 C6.3 18.1 5.9 18.9 5.6 19.7 C5.2 20.5 4.9 21.3 4.5 22.1 C3.7 24.1 2.9 26 2 28 C6.6 28 11.2 28 16 28 C16 28.3 16 28.7 16 29 C0.8 29 -14.4 29 -30 29 C-30 28.7 -30 28.3 -30 28 C-25.1 27.8 -25.1 27.8 0 27 C-3 20.3 -3 20.3 -3.9 18.1 C-5.4 14.8 -6 12.7 -6 9 C-4.7 9.3 -3.4 9.7 -2 10 C-0.9 6.7 -0.5 3.5 0 0 Z };
\path[fill=black!81,even odd rule,shift={(415,2)}] svg {M0 0 C1 0 2 0 3 0 C3 7.6 3 15.2 3 23 C5 22.3 7 21.7 9 21 C10 23 10 23 9.5 24.8 C9.2 25.5 8.9 26.2 8.6 26.9 C8.2 27.7 7.9 28.4 7.6 29.2 C7.3 30 6.9 30.8 6.6 31.6 C6.2 32.4 5.9 33.2 5.5 34.1 C4.7 36 3.9 38 3 40 C2 39.7 1 39.3 0 39 C-1.2 36.5 -2.2 34.2 -3.2 31.6 C-3.5 30.9 -3.7 30.3 -4 29.5 C-5.2 26.6 -6 24.2 -6 21 C-4 21.7 -2 22.3 0 23 C0 15.4 0 7.8 0 0 Z };
\path[fill=black!77,even odd rule,shift={(209,247)}] svg {M0 0 C1.2 3.6 0.7 4.7 -0.4 8.2 C-0.7 9.1 -1 10 -1.3 10.9 C-1.5 11.6 -1.8 12.3 -2 13 C-1.3 13 -0.7 13 0 13 C-0.9 16.9 -0.9 16.9 -2 18 C-3.7 18 -5.3 18 -7 18 C-7 17.3 -7 16.7 -7 16 C-7.7 16.5 -8.5 17 -9.2 17.6 C-12 19 -12 19 -14.2 18.8 C-16 18 -16 18 -17 16 C-17.4 12.2 -17.2 10.3 -15.3 6.9 C-10.4 0.8 -7.7 -0.5 0 0 Z M-8 3 C-10.4 6 -11.9 8 -12.2 11.9 C-12.2 12.6 -12.1 13.3 -12 14 C-11 14 -10 14 -9 14 C-6.6 11 -5.1 9 -4.8 5.1 C-4.8 4.4 -4.9 3.7 -5 3 C-6 3 -7 3 -8 3 Z };
\path[fill=black!80,even odd rule,shift={(416,98)}] svg {M0 0 C3.5 2.7 4.7 6.7 6.2 10.7 C6.5 11.4 6.7 12.1 7 12.8 C7.7 14.5 8.3 16.3 9 18 C8 18.3 7 18.7 6 19 C5.3 18.7 4.7 18.3 4 18 C3.7 21 3.3 23.9 3 27 C1.7 27 0.4 27 -1 27 C-1.3 24 -1.7 21.1 -2 18 C-3.6 18.3 -5.3 18.7 -7 19 C-5.2 12.4 -2.9 6.2 0 0 Z };
\path[fill=black!76,even odd rule,shift={(377,216)}] svg {M0 0 C18.5 0.3 37 0.7 56 1 C55.7 2 55.3 3 55 4 C47.5 4.2 40 4.3 32.5 4.5 C30 4.5 27.4 4.6 24.9 4.6 C21.2 4.7 17.5 4.8 13.9 4.9 C12.7 4.9 11.6 4.9 10.4 5 C9.4 5 8.3 5 7.2 5 C6.7 5 6.7 5 4.4 5.1 C2 5 2 5 -1 4 C-0.7 2.7 -0.3 1.4 0 0 Z };
\path[fill=black!78,even odd rule,shift={(415,88)}] svg {M0 0 C0.7 0 1.3 0 2 0 C2 1 2 2 2 3 C3 3 4 3 5 3 C5 2.3 5 1.7 5 1 C6 1 7 1 8 1 C8 2 8 3 8 4 C10.6 4 13.3 4 16 4 C16.2 4.7 16.2 4.7 17 8 C4.5 8 -8.1 8 -21 8 C-21 7 -21 6 -21 5 C-20.6 5 -20.6 5 -18.3 4.8 C-17.2 4.7 -16 4.6 -14.8 4.6 C-13.7 4.5 -12.5 4.4 -11.3 4.3 C-3.2 3.5 -3.2 3.5 0 0 Z };
\path[fill=black!85,even odd rule,shift={(380,223)}] svg {M0 0 C0.7 0.3 1.3 0.7 2 1 C2 18.2 2 35.3 2 53 C1.5 53.2 1.5 53.2 -1 54 C-1 46.5 -1 38.9 -1.1 31.4 C-1.1 28.8 -1.1 26.3 -1.1 23.7 C-1.1 20 -1.1 16.4 -1.1 12.7 C-1.1 11.5 -1.1 10.4 -1.1 9.2 C-1.1 8.1 -1.1 7.1 -1.1 6 C-1.1 5 -1.1 4.1 -1.1 3.1 C-1 1 -1 1 0 0 Z };
\path[fill=black!91,even odd rule,shift={(432,47)}] svg {M0 0 C0 1 0 2 0 3 C-0.8 3 -1.6 3 -2.4 3 C-12.3 3.5 -12.3 3.5 -15.2 6.6 C-15.8 7.4 -16.4 8.2 -17 9 C-17.3 7 -17.7 5 -18 3 C-25.6 3 -33.2 3 -41 3 C-41 2 -41 1 -41 0 C-36.6 -0.2 -32.2 -0.3 -27.8 -0.5 C-26.3 -0.5 -24.8 -0.6 -23.4 -0.6 C-15.5 -1 -7.8 -1.2 0 0 Z };
\path[fill=black!91,even odd rule,shift={(3,232)}] svg {M0 0 C1 0.3 2 0.7 3 1 C3 15.8 3 30.7 3 46 C2 46 1 46 0 46 C0 30.8 0 15.6 0 0 Z };
\path[fill=black!78,even odd rule,shift={(181.375,238.6875)}] svg {M0 0 C3.8 0.3 5.2 0.5 7.6 3.3 C8.9 6.4 9.1 8.3 7.7 11.4 C5.8 14.8 4.4 16.9 0.8 18.8 C0.5 18.9 0.5 18.9 -1.4 19.3 C0.6 15.3 0.6 15.3 1.8 13.4 C2.9 10.7 2.8 8.3 2.6 5.3 C-0.3 5.7 -0.3 5.7 -3.4 6.3 C-3.7 7 -4 7.6 -4.4 8.3 C-5 8.3 -5.7 8.3 -6.4 8.3 C-5.4 0.5 -5.4 0.5 0 0 Z };
\path[fill=black!69,even odd rule,shift={(431,142)}] svg {M0 0 C0.7 0.3 1.3 0.7 2 1 C1.7 2.7 1.3 4.3 1 6 C2.6 6.3 4.3 6.7 6 7 C7.2 9.9 7 11.1 5.9 14.1 C4 17 4 17 1 18.4 C-2 19 -2 19 -5 18 C-4.9 16.9 -4.7 15.7 -4.6 14.6 C-4 10 -3.5 5.5 -3 1 C-2 0.7 -1 0.3 0 0 Z M1 9 C0 11.3 -1 13.6 -2 16 C1.9 15.1 1.9 15.1 3 14 C3 12.3 3 10.7 3 9 C2.3 9 1.7 9 1 9 Z };
\path[fill=black!71,even odd rule,shift={(417.4375,156.5625)}] svg {M0 0 C2.6 0.4 2.6 0.4 4.6 2.6 C5.6 5.4 5.6 5.4 4.4 8.6 C3.1 10.5 1.9 12.5 0.6 14.4 C2.2 14.1 3.9 13.8 5.6 13.4 C5.2 14.8 4.9 16.1 4.6 17.4 C0.9 18.4 -1.7 18.4 -5.4 17.4 C-4.3 14.9 -3.1 13.1 -1.2 10.9 C0.6 8.4 0.6 8.4 0.3 6.1 C0.1 5.6 -0.2 5 -0.4 4.4 C-1.1 4.8 -1.1 4.8 -4.4 6.4 C-4.4 2.1 -4.4 0.8 0 0 Z };
\path[fill=black!69,even odd rule,shift={(418,131.3125)}] svg {M0 0 C2 0.7 2 0.7 4 2.7 C4 4.7 4 6.7 4 8.7 C4.3 9.7 4.7 10.7 5 11.7 C4.2 14.3 4.2 14.3 3 16.7 C0.3 18 -2 17.8 -5 17.7 C-5.3 16.7 -5.7 15.7 -6 14.7 C-3.7 14.7 -1.4 14.7 1 14.7 C-0.8 10.8 -0.8 10.8 -3 9.7 C-3 9 -3 8.4 -3 7.7 C-2.3 7.4 -1.7 7 -1 6.7 C-0.7 5.7 -0.3 4.7 0 3.7 C-1.6 4 -3.3 4.3 -5 4.7 C-4.6 2.8 -4.6 2.8 -4 0.7 C-2 -0.3 -2 -0.3 0 0 Z };
\path[fill=black!71,even odd rule,shift={(417,54)}] svg {M0 0 C0.3 0.7 0.7 1.3 1 2 C3.6 2.6 3.6 2.6 6 3 C7 6 7 6 5.8 9.1 C4 12 4 12 1 14 C-1.6 14.2 -1.6 14.2 -4 14 C-5 13 -5 13 -5 11.1 C-4.5 4.6 -4.5 4.6 -1.9 1.4 C-1.3 1 -0.7 0.5 0 0 Z M0 5 C-0.3 7.3 -0.7 9.6 -1 12 C1 10.8 1 10.8 3 9 C3 7.7 3 6.4 3 5 C2 5 1 5 0 5 Z };
\path[fill=black!86,even odd rule,shift={(182,256)}] svg {M0 0 C0.3 0.7 0.7 1.3 1 2 C0.7 2.7 0.3 3.3 0 4 C1.5 3.8 1.5 3.8 9 3 C8.7 5.4 8.7 5.4 8 8 C4.3 10.5 2.7 10.1 -1.7 9.6 C-2.9 9.5 -4 9.4 -5.3 9.3 C-6.2 9.2 -7.1 9.1 -8 9 C-6.1 4.7 -3.6 2.8 0 0 Z };
\path[fill=black!72,even odd rule,shift={(422,77)}] svg {M0 0 C0.5 5.7 -0.5 8.5 -4 13 C-3.3 13.3 -2.7 13.7 -2 14 C-3 14 -4 14 -5 14 C-6.5 10.2 -5.4 7.7 -4 4 C-4 3.3 -4 2.7 -4 2 C-6 2.7 -8 3.3 -10 4 C-9.2 1.6 -9.2 1.6 -8 -1 C-4.4 -2.2 -3.4 -1.5 0 0 Z };
\path[fill=black!71,even odd rule,shift={(410,151)}] svg {M0 0 C4.6 0 9.2 0 14 0 C14.3 1 14.7 2 15 3 C9.7 3 4.4 3 -1 3 C-0.7 2 -0.3 1 0 0 Z };
\path[fill=black!51,even odd rule,shift={(409,70)}] svg {M0 0 C4.9 0 9.9 0 15 0 C15 1 15 2 15 3 C10.1 3 5.1 3 0 3 C0 2 0 1 0 0 Z };
\path[fill=black!58,even odd rule,shift={(81,193)}] svg {M0 0 C1.7 0 3.3 0 5 0 C5 1.3 5 2.6 5 4 C3.3 4 1.7 4 0 4 C0 2.7 0 1.4 0 0 Z };
\path[fill=black!58,even odd rule,shift={(84,173)}] svg {M0 0 C1.7 0.3 3.3 0.7 5 1 C4.7 2 4.3 3 4 4 C2.3 3.7 0.7 3.3 -1 3 C-0.7 2 -0.3 1 0 0 Z };
\path[fill=black!58,even odd rule,shift={(297,132)}] svg {M0 0 C1.6 0.3 3.3 0.7 5 1 C4.7 2 4.3 3 4 4 C2.4 3.7 0.7 3.3 -1 3 C-0.7 2 -0.3 1 0 0 Z };
\path[fill=black!58,even odd rule,shift={(239,112)}] svg {M0 0 C1.7 0.3 3.3 0.7 5 1 C4.7 2 4.3 3 4 4 C2.3 3.7 0.7 3.3 -1 3 C-0.7 2 -0.3 1 0 0 Z };
\path[fill=black!60,even odd rule,shift={(106,173)}] svg {M0 0 C1.7 0.3 3.3 0.7 5 1 C4.7 2 4.3 3 4 4 C2.3 3.7 0.7 3.3 -1 3 C-0.7 2 -0.3 1 0 0 Z };
\path[fill=black!60,even odd rule,shift={(67,174)}] svg {M0 0 C-0.3 1 -0.7 2 -1 3 C-2.7 2.7 -4.3 2.3 -6 2 C-6 1.3 -6 0.7 -6 0 C-3.5 -1.2 -2.6 -0.8 0 0 Z };
\path[fill=black!60,even odd rule,shift={(155,153)}] svg {M0 0 C1.7 0 3.3 0 5 0 C5.3 0.7 5.7 1.3 6 2 C5.7 2.7 5.3 3.3 5 4 C3.3 3.7 1.7 3.3 0 3 C0 2 0 1 0 0 Z };
\path[fill=black!60,even odd rule,shift={(119,153)}] svg {M0 0 C1.7 0 3.3 0 5 0 C4.7 1.3 4.3 2.6 4 4 C2.7 4 1.4 4 0 4 C0 2.7 0 1.4 0 0 Z };
\path[fill=black!60,even odd rule,shift={(90,153)}] svg {M0 0 C1.7 0.3 3.3 0.7 5 1 C4.7 2 4.3 3 4 4 C2.3 3.7 0.7 3.3 -1 3 C-0.7 2 -0.3 1 0 0 Z };
\path[fill=black!60,even odd rule,shift={(324,152)}] svg {M0 0 C1.7 0.3 3.3 0.6 5 1 C4.7 2 4.3 3 4 4 C2.7 4 1.4 4 0 4 C-0.3 3 -0.7 2 -1 1 C-0.7 0.7 -0.3 0.3 0 0 Z };
\path[fill=black!60,even odd rule,shift={(317,152)}] svg {M0 0 C1.6 0.3 3.3 0.7 5 1 C5 1.7 5 2.3 5 3 C3.4 3.3 1.7 3.7 0 4 C0 2.7 0 1.4 0 0 Z };
\path[fill=black!60,even odd rule,shift={(295,152)}] svg {M0 0 C1.6 0.3 3.3 0.7 5 1 C4.7 2 4.3 3 4 4 C2.4 3.7 0.7 3.3 -1 3 C-0.7 2 -0.3 1 0 0 Z };
\path[fill=black!60,even odd rule,shift={(92,133)}] svg {M0 0 C1.7 0.3 3.3 0.7 5 1 C4.7 2 4.3 3 4 4 C2.3 3.7 0.7 3.3 -1 3 C-0.7 2 -0.3 1 0 0 Z };
\path[fill=black!60,even odd rule,shift={(48,133)}] svg {M0 0 C1.3 0 2.6 0 4 0 C4.3 1 4.7 2 5 3 C3.4 3.3 1.7 3.7 0 4 C0 2.7 0 1.4 0 0 Z };
\path[fill=black!60,even odd rule,shift={(251,133)}] svg {M0 0 C-0.3 1 -0.7 2 -1 3 C-2.7 2.7 -4.3 2.3 -6 2 C-6 1.3 -6 0.7 -6 0 C-3.5 -1.2 -2.6 -0.8 0 0 Z };
\path[fill=black!67,even odd rule,shift={(115,112)}] svg {M0 0 C0.7 0.2 0.7 0.2 4 1 C3.7 2.3 3.3 3.6 3 5 C1.7 4.7 0.4 4.3 -1 4 C-0.7 2.7 -0.3 1.4 0 0 Z };
\path[fill=black!60,even odd rule,shift={(41,113)}] svg {M0 0 C1.6 0.3 3.3 0.7 5 1 C4.7 2 4.3 3 4 4 C2.4 3.7 0.7 3.3 -1 3 C-0.7 2 -0.3 1 0 0 Z };
\path[fill=black!60,even odd rule,shift={(290,112)}] svg {M0 0 C1.7 0.3 3.3 0.6 5 1 C4.7 2 4.3 3 4 4 C2.7 4 1.4 4 0 4 C-0.3 3 -0.7 2 -1 1 C-0.7 0.7 -0.3 0.3 0 0 Z };
\path[fill=black!56,even odd rule,shift={(283,112)}] svg {M0 0 C1.6 0.3 3.3 0.7 5 1 C4.7 2 4.3 3 4 4 C2.4 3.7 0.7 3.3 -1 3 C-0.7 2 -0.3 1 0 0 Z };
\path[fill=black!58,even odd rule,shift={(274,217)}] svg {M0 0 C1.6 0 3.3 0 5 0 C5 1 5 2 5 3 C1.2 4.1 1.2 4.1 -1 3 C-0.7 2 -0.3 1 0 0 Z };
\path[fill=black!54,even odd rule,shift={(280,192)}] svg {M0 0 C0.7 0.2 0.7 0.2 4 1 C4 2 4 3 4 4 C2.4 4 0.7 4 -1 4 C-0.7 2.7 -0.3 1.4 0 0 Z };
\path[fill=black!58,even odd rule,shift={(147,193)}] svg {M0 0 C1.7 0 3.3 0 5 0 C4.7 1.3 4.3 2.6 4 4 C2.7 3.7 1.4 3.3 0 3 C0 2 0 1 0 0 Z };
\path[fill=black!58,even odd rule,shift={(140,193)}] svg {M0 0 C0.8 0.2 0.8 0.2 5 1 C4.7 2 4.3 3 4 4 C2.7 4 1.4 4 0 4 C0 2.7 0 1.4 0 0 Z };
\path[fill=black!62,even odd rule,shift={(133,193)}] svg {M0 0 C1.3 0 2.6 0 4 0 C4 1.3 4 2.6 4 4 C2.3 3.7 0.7 3.3 -1 3 C-0.7 2 -0.3 1 0 0 Z };
\path[fill=black!62,even odd rule,shift={(111,193)}] svg {M0 0 C1.7 0.3 3.3 0.7 5 1 C4.7 2 4.3 3 4 4 C2.3 3.7 0.7 3.3 -1 3 C-0.7 2 -0.3 1 0 0 Z };
\path[fill=black!62,even odd rule,shift={(75,193)}] svg {M0 0 C0.7 0.2 0.7 0.2 4 1 C4 2 4 3 4 4 C2.3 4 0.7 4 -1 4 C-0.7 2.7 -0.3 1.4 0 0 Z };
\path[fill=black!54,even odd rule,shift={(142,173)}] svg {M0 0 C1.7 0 3.3 0 5 0 C4.7 1.3 4.3 2.6 4 4 C2.7 3.7 1.4 3.3 0 3 C0 2 0 1 0 0 Z };
\path[fill=black!58,even odd rule,shift={(120,173)}] svg {M0 0 C1.7 0 3.3 0 5 0 C4.7 1.3 4.3 2.6 4 4 C2.7 3.7 1.4 3.3 0 3 C0 2 0 1 0 0 Z };
\path[fill=black!62,even odd rule,shift={(113,173)}] svg {M0 0 C0.8 0.2 0.8 0.2 5 1 C4.7 2 4.3 3 4 4 C2.7 4 1.4 4 0 4 C0 2.7 0 1.4 0 0 Z };
\path[fill=black!58,even odd rule,shift={(77,173)}] svg {M0 0 C0.7 0.2 0.7 0.2 4 1 C4 2 4 3 4 4 C2.3 4 0.7 4 -1 4 C-0.7 2.7 -0.3 1.4 0 0 Z };
\path[fill=black!58,even odd rule,shift={(347,172)}] svg {M0 0 C1.6 0 3.3 0 5 0 C4.7 1.3 4.3 2.6 4 4 C2.7 3.7 1.4 3.3 0 3 C0 2 0 1 0 0 Z };
\path[fill=black!54,even odd rule,shift={(333,172)}] svg {M0 0 C1.3 0 2.6 0 4 0 C4 1.3 4 2.6 4 4 C2.4 3.7 0.7 3.3 -1 3 C-0.7 2 -0.3 1 0 0 Z };
\path[fill=black!62,even odd rule,shift={(318,172)}] svg {M0 0 C0.8 0.2 0.8 0.2 5 1 C4.7 2 4.3 3 4 4 C2.7 4 1.4 4 0 4 C0 2.7 0 1.4 0 0 Z };
\path[fill=black!62,even odd rule,shift={(3,166)}] svg {M0 0 C1 0.3 2 0.7 3 1 C2.7 2.7 2.3 4.3 2 6 C1 5.3 0 4.7 -1 4 C-0.7 2.7 -0.3 1.4 0 0 Z };
\path[fill=black!54,even odd rule,shift={(39,153)}] svg {M0 0 C0.7 0.2 0.7 0.2 4 1 C4 2 4 3 4 4 C2.4 4 0.7 4 -1 4 C-0.7 2.7 -0.3 1.4 0 0 Z };
\path[fill=black!58,even odd rule,shift={(126,153)}] svg {M0 0 C1.7 0 3.3 0 5 0 C4.7 1.3 4.3 2.6 4 4 C2.7 3.7 1.4 3.3 0 3 C0 2 0 1 0 0 Z };
\path[fill=black!62,even odd rule,shift={(97,153)}] svg {M0 0 C0.8 0.2 0.8 0.2 5 1 C4.7 2 4.3 3 4 4 C2.7 4 1.4 4 0 4 C0 2.7 0 1.4 0 0 Z };
\path[fill=black!58,even odd rule,shift={(83,153)}] svg {M0 0 C1.3 0 2.6 0 4 0 C4 1.3 4 2.6 4 4 C2.3 3.7 0.7 3.3 -1 3 C-0.7 2 -0.3 1 0 0 Z };
\path[fill=black!62,even odd rule,shift={(61,153)}] svg {M0 0 C0.7 0.2 0.7 0.2 4 1 C4 2 4 3 4 4 C2.4 4 0.7 4 -1 4 C-0.7 2.7 -0.3 1.4 0 0 Z };
\path[fill=black!62,even odd rule,shift={(46,153)}] svg {M0 0 C1.3 0 2.6 0 4 0 C4 1.3 4 2.6 4 4 C2.4 3.7 0.7 3.3 -1 3 C-0.7 2 -0.3 1 0 0 Z };
\path[fill=black!62,even odd rule,shift={(346,152)}] svg {M0 0 C1.6 0.3 3.3 0.7 5 1 C4.7 1.5 4.7 1.5 3 4 C1.7 3.7 0.4 3.3 -1 3 C-0.7 2 -0.3 1 0 0 Z };
\path[fill=black!58,even odd rule,shift={(288,152)}] svg {M0 0 C1.3 0 2.6 0 4 0 C4 1.3 4 2.6 4 4 C2.4 3.7 0.7 3.3 -1 3 C-0.7 2 -0.3 1 0 0 Z };
\path[fill=black!54,even odd rule,shift={(283,132)}] svg {M0 0 C0.7 0.2 0.7 0.2 4 1 C4 2 4 3 4 4 C2.4 4 0.7 4 -1 4 C-0.7 2.7 -0.3 1.4 0 0 Z };
\path[fill=black!62,even odd rule,shift={(136,133)}] svg {M0 0 C1.7 0.3 3.3 0.7 5 1 C4.7 1.5 4.7 1.5 3 4 C1.7 3.7 0.4 3.3 -1 3 C-0.7 2 -0.3 1 0 0 Z };
\path[fill=black!62,even odd rule,shift={(77,133)}] svg {M0 0 C0.8 0.2 0.8 0.2 5 1 C4.7 2 4.3 3 4 4 C2.7 4 1.4 4 0 4 C0 2.7 0 1.4 0 0 Z };
\path[fill=black!62,even odd rule,shift={(71,133)}] svg {M0 0 C1.3 0.3 2.6 0.7 4 1 C3.7 2 3.3 3 3 4 C1.3 3.7 -0.3 3.3 -2 3 C-1.3 2 -0.7 1 0 0 Z };
\path[fill=black!58,even odd rule,shift={(41,133)}] svg {M0 0 C0.7 0.2 0.7 0.2 4 1 C4 2 4 3 4 4 C2.4 4 0.7 4 -1 4 C-0.7 2.7 -0.3 1.4 0 0 Z };
\path[fill=black!62,even odd rule,shift={(33,133)}] svg {M0 0 C0.8 0.2 0.8 0.2 5 1 C5 2 5 3 5 4 C3.4 4 1.7 4 0 4 C0 2.7 0 1.4 0 0 Z };
\path[fill=black!62,even odd rule,shift={(26,133)}] svg {M0 0 C0.7 0.2 0.7 0.2 4 1 C4 2 4 3 4 4 C2.4 3.7 0.7 3.3 -1 3 C-0.7 2 -0.3 1 0 0 Z };
\path[fill=black!58,even odd rule,shift={(19,133)}] svg {M0 0 C0.7 0.2 0.7 0.2 4 1 C4 2 4 3 4 4 C2.4 4 0.7 4 -1 4 C-0.7 2.7 -0.3 1.4 0 0 Z };
\path[fill=black!58,even odd rule,shift={(362,132)}] svg {M0 0 C1.6 0 3.3 0 5 0 C4.7 1.3 4.3 2.6 4 4 C2.7 3.7 1.4 3.3 0 3 C0 2 0 1 0 0 Z };
\path[fill=black!58,even odd rule,shift={(326,132)}] svg {M0 0 C0.8 0.2 0.8 0.2 5 1 C4.7 2 4.3 3 4 4 C2.7 4 1.4 4 0 4 C0 2.7 0 1.4 0 0 Z };
\path[fill=black!62,even odd rule,shift={(319,132)}] svg {M0 0 C1.3 0 2.6 0 4 0 C4 1.3 4 2.6 4 4 C2.4 3.7 0.7 3.3 -1 3 C-0.7 2 -0.3 1 0 0 Z };
\path[fill=black!62,even odd rule,shift={(268,132)}] svg {M0 0 C1.3 0 2.6 0 4 0 C4 1.3 4 2.6 4 4 C2.4 3.7 0.7 3.3 -1 3 C-0.7 2 -0.3 1 0 0 Z };
\path[fill=black!58,even odd rule,shift={(70,113)}] svg {M0 0 C0.8 0.2 0.8 0.2 5 1 C4.7 2 4.3 3 4 4 C2.7 4 1.4 4 0 4 C0 2.7 0 1.4 0 0 Z };
\path[fill=black!62,even odd rule,shift={(63,113)}] svg {M0 0 C1.7 0.3 3.3 0.6 5 1 C4.7 2 4.3 3 4 4 C2.7 4 1.4 4 0 4 C-0.3 3 -0.7 2 -1 1 C-0.7 0.7 -0.3 0.3 0 0 Z };
\path[fill=black!62,even odd rule,shift={(56,113)}] svg {M0 0 C1.3 0 2.6 0 4 0 C4 1.3 4 2.6 4 4 C2.4 3.7 0.7 3.3 -1 3 C-0.7 2 -0.3 1 0 0 Z };
\path[fill=black!58,even odd rule,shift={(34,113)}] svg {M0 0 C0.7 0.2 0.7 0.2 4 1 C4 2 4 3 4 4 C2.4 4 0.7 4 -1 4 C-0.7 2.7 -0.3 1.4 0 0 Z };
\path[fill=black!58,even odd rule,shift={(319,112)}] svg {M0 0 C1.6 0 3.3 0 5 0 C4.7 1.3 4.3 2.6 4 4 C2.7 3.7 1.4 3.3 0 3 C0 2 0 1 0 0 Z };
\path[fill=black!62,even odd rule,shift={(268,112)}] svg {M0 0 C0.8 0.2 0.8 0.2 5 1 C4.7 2 4.3 3 4 4 C2.7 4 1.4 4 0 4 C0 2.7 0 1.4 0 0 Z };
\path[fill=black!62,even odd rule,shift={(261,112)}] svg {M0 0 C1.3 0 2.6 0 4 0 C4 1.3 4 2.6 4 4 C2.4 3.7 0.7 3.3 -1 3 C-0.7 2 -0.3 1 0 0 Z };
\path[fill=black!54,even odd rule,shift={(261,47)}] svg {M0 0 C0.7 0.2 0.7 0.2 4 1 C4 2 4 3 4 4 C2.4 4 0.7 4 -1 4 C-0.7 2.7 -0.3 1.4 0 0 Z };
\path[fill=black!54,even odd rule,shift={(355,47)}] svg {M0 0 C1.6 0 3.3 0 5 0 C4.7 1.3 4.3 2.6 4 4 C2.7 3.7 1.4 3.3 0 3 C0 2 0 1 0 0 Z };
\path[fill=black!58,even odd rule,shift={(326,47)}] svg {M0 0 C0.8 0.2 0.8 0.2 5 1 C4.7 2 4.3 3 4 4 C2.7 4 1.4 4 0 4 C0 2.7 0 1.4 0 0 Z };
\path[fill=black!62,even odd rule,shift={(319,47)}] svg {M0 0 C0.8 0.2 0.8 0.2 5 1 C4.7 2 4.3 3 4 4 C2.7 4 1.4 4 0 4 C0 2.7 0 1.4 0 0 Z };
\path[fill=black!54,even odd rule,shift={(312,47)}] svg {M0 0 C1.3 0 2.6 0 4 0 C4 1.3 4 2.6 4 4 C2.4 3.7 0.7 3.3 -1 3 C-0.7 2 -0.3 1 0 0 Z };
\path[fill=black!71,even odd rule,shift={(267,217)}] svg {M0 0 C1.7 0.3 3.3 0.6 5 1 C5 1.7 5 2.3 5 3 C3.7 3.3 2.4 3.7 1 4 C0.3 3 -0.3 2 -1 1 C-0.7 0.7 -0.3 0.3 0 0 Z };
\path[fill=black!67,even odd rule,shift={(356,216)}] svg {M0 0 C1 0.3 2 0.7 3 1 C3 2 3 3 3 4 C1.4 4 -0.3 4 -2 4 C-2 3 -2 2 -2 1 C-1.3 0.7 -0.7 0.3 0 0 Z };
\path[fill=black!63,even odd rule,shift={(126,193)}] svg {M0 0 C1.3 0 2.6 0 4 0 C3.7 1.3 3.3 2.6 3 4 C1.7 3.7 0.4 3.3 -1 3 C-0.7 2 -0.3 1 0 0 Z };
\path[fill=black!59,even odd rule,shift={(118,193)}] svg {M0 0 C0.8 0.2 0.8 0.2 5 1 C4.7 2 4.3 3 4 4 C2.7 3.7 1.4 3.3 0 3 C0 2 0 1 0 0 Z };
\path[fill=black!63,even odd rule,shift={(104,193)}] svg {M0 0 C1.3 0 2.6 0 4 0 C3.7 1.3 3.3 2.6 3 4 C1.7 3.7 0.4 3.3 -1 3 C-0.7 2 -0.3 1 0 0 Z };
\path[fill=black!63,even odd rule,shift={(97,193)}] svg {M0 0 C0.7 0.2 0.7 0.2 4 1 C3.7 2 3.3 3 3 4 C1.7 4 0.4 4 -1 4 C-0.7 2.7 -0.3 1.4 0 0 Z };
\path[fill=black!63,even odd rule,shift={(89,193)}] svg {M0 0 C0.7 0.2 0.7 0.2 4 1 C4 2 4 3 4 4 C2.3 3.7 0.7 3.3 -1 3 C-0.7 2 -0.3 1 0 0 Z };
\path[fill=black!59,even odd rule,shift={(287,192)}] svg {M0 0 C0.7 0.2 0.7 0.2 4 1 C4 2 4 3 4 4 C2.4 3.7 0.7 3.3 -1 3 C-0.7 2 -0.3 1 0 0 Z };
\path[fill=black!59,even odd rule,shift={(3,181)}] svg {M0 0 C1 0 2 0 3 0 C3 1.7 3 3.3 3 5 C1.7 4.7 0.4 4.3 -1 4 C-0.7 2.7 -0.3 1.4 0 0 Z };
\path[fill=black!67,even odd rule,shift={(216,173)}] svg {M0 0 C1.7 0 3.3 0 5 0 C5 1 5 2 5 3 C3 3 1 3 -1 3 C-0.7 2 -0.3 1 0 0 Z };
\path[fill=black!59,even odd rule,shift={(157,173)}] svg {M0 0 C0.8 0.2 0.8 0.2 5 1 C4.7 2 4.3 3 4 4 C2.7 3.7 1.4 3.3 0 3 C0 2 0 1 0 0 Z };
\path[fill=black!63,even odd rule,shift={(128,173)}] svg {M0 0 C1.3 0 2.6 0 4 0 C3.7 1.3 3.3 2.6 3 4 C1.7 3.7 0.4 3.3 -1 3 C-0.7 2 -0.3 1 0 0 Z };
\path[fill=black!63,even odd rule,shift={(92,173)}] svg {M0 0 C0.7 0.2 0.7 0.2 4 1 C3.7 2 3.3 3 3 4 C1.7 4 0.4 4 -1 4 C-0.7 2.7 -0.3 1.4 0 0 Z };
\path[fill=black!59,even odd rule,shift={(55,173)}] svg {M0 0 C0.7 0.2 0.7 0.2 4 1 C4 2 4 3 4 4 C2.4 3.7 0.7 3.3 -1 3 C-0.7 2 -0.3 1 0 0 Z };
\path[fill=black!63,even odd rule,shift={(41,173)}] svg {M0 0 C1 0.3 2 0.7 3 1 C3 2 3 3 3 4 C1.4 4 -0.3 4 -2 4 C-2 3 -2 2 -2 1 C-1.3 0.7 -0.7 0.3 0 0 Z };
\path[fill=black!63,even odd rule,shift={(311,172)}] svg {M0 0 C1.3 0 2.6 0 4 0 C3.7 1.3 3.3 2.6 3 4 C1.7 3.7 0.4 3.3 -1 3 C-0.7 2 -0.3 1 0 0 Z };
\path[fill=black!55,even odd rule,shift={(304,172)}] svg {M0 0 C0.7 0.2 0.7 0.2 4 1 C3.7 2 3.3 3 3 4 C1.7 4 0.4 4 -1 4 C-0.7 2.7 -0.3 1.4 0 0 Z };
\path[fill=black!55,even odd rule,shift={(297,172)}] svg {M0 0 C0.7 0.2 0.7 0.2 4 1 C3.7 2 3.3 3 3 4 C1.7 4 0.4 4 -1 4 C-0.7 2.7 -0.3 1.4 0 0 Z };
\path[fill=black!59,even odd rule,shift={(282,172)}] svg {M0 0 C0.7 0.2 0.7 0.2 4 1 C4 2 4 3 4 4 C2.4 3.7 0.7 3.3 -1 3 C-0.7 2 -0.3 1 0 0 Z };
\path[fill=black!59,even odd rule,shift={(379,165)}] svg {M0 0 C1 0 2 0 3 0 C2.7 1.7 2.3 3.3 2 5 C1 4.7 0 4.3 -1 4 C-0.7 2.7 -0.3 1.4 0 0 Z };
\path[fill=black!59,even odd rule,shift={(3,159)}] svg {M0 0 C1 0.3 2 0.7 3 1 C3 2.3 3 3.6 3 5 C1.7 4.7 0.4 4.3 -1 4 C-0.7 2.7 -0.3 1.4 0 0 Z };
\path[fill=black!59,even odd rule,shift={(105,153)}] svg {M0 0 C1.3 0 2.6 0 4 0 C3.7 1.3 3.3 2.6 3 4 C1.7 3.7 0.4 3.3 -1 3 C-0.7 2 -0.3 1 0 0 Z };
\path[fill=black!55,even odd rule,shift={(76,153)}] svg {M0 0 C0.7 0.2 0.7 0.2 4 1 C3.7 2 3.3 3 3 4 C1.7 4 0.4 4 -1 4 C-0.7 2.7 -0.3 1.4 0 0 Z };
\path[fill=black!63,even odd rule,shift={(68,153)}] svg {M0 0 C0.7 0.2 0.7 0.2 4 1 C4 2 4 3 4 4 C2.3 3.7 0.7 3.3 -1 3 C-0.7 2 -0.3 1 0 0 Z };
\path[fill=black!59,even odd rule,shift={(54,153)}] svg {M0 0 C0.7 0.2 0.7 0.2 4 1 C3.7 2 3.3 3 3 4 C1.7 4 0.4 4 -1 4 C-0.7 2.7 -0.3 1.4 0 0 Z };
\path[fill=black!59,even odd rule,shift={(32,153)}] svg {M0 0 C0.7 0.2 0.7 0.2 4 1 C3.7 2 3.3 3 3 4 C1.7 4 0.4 4 -1 4 C-0.7 2.7 -0.3 1.4 0 0 Z };
\path[fill=black!59,even odd rule,shift={(18,153)}] svg {M0 0 C1 0.3 2 0.7 3 1 C3 2 3 3 3 4 C1.4 4 -0.3 4 -2 4 C-2 3 -2 2 -2 1 C-1.3 0.7 -0.7 0.3 0 0 Z };
\path[fill=black!55,even odd rule,shift={(353,152)}] svg {M0 0 C1.6 0 3.3 0 5 0 C5 1 5 2 5 3 C3.7 3.3 2.4 3.7 1 4 C0.7 2.7 0.3 1.4 0 0 Z };
\path[fill=black!55,even odd rule,shift={(339,152)}] svg {M0 0 C1.3 0 2.6 0 4 0 C3.7 1.3 3.3 2.6 3 4 C1.7 3.7 0.4 3.3 -1 3 C-0.7 2 -0.3 1 0 0 Z };
\path[fill=black!59,even odd rule,shift={(303,152)}] svg {M0 0 C0.7 0.2 0.7 0.2 4 1 C3.7 2 3.3 3 3 4 C1.7 4 0.4 4 -1 4 C-0.7 2.7 -0.3 1.4 0 0 Z };
\path[fill=black!55,even odd rule,shift={(261,132)}] svg {M0 0 C0.7 0.2 0.7 0.2 4 1 C4 2 4 3 4 4 C2.4 4 0.7 4 -1 4 C-0.7 2.7 -0.3 1.4 0 0 Z };
\path[fill=black!63,even odd rule,shift={(150,133)}] svg {M0 0 C1.7 0 3.3 0 5 0 C5 1 5 2 5 3 C3.7 3.3 2.4 3.7 1 4 C0.7 2.7 0.3 1.4 0 0 Z };
\path[fill=black!59,even odd rule,shift={(121,133)}] svg {M0 0 C0.8 0.2 0.8 0.2 5 1 C4.7 2 4.3 3 4 4 C2.7 3.7 1.4 3.3 0 3 C0 2 0 1 0 0 Z };
\path[fill=black!63,even odd rule,shift={(114,133)}] svg {M0 0 C1.3 0 2.6 0 4 0 C4 1.3 4 2.6 4 4 C2.7 4 1.4 4 0 4 C-0.3 3 -0.7 2 -1 1 C-0.7 0.7 -0.3 0.3 0 0 Z };
\path[fill=black!55,even odd rule,shift={(107,133)}] svg {M0 0 C1.3 0 2.6 0 4 0 C3.7 1.3 3.3 2.6 3 4 C1.7 3.7 0.4 3.3 -1 3 C-0.7 2 -0.3 1 0 0 Z };
\path[fill=black!63,even odd rule,shift={(99,133)}] svg {M0 0 C0.8 0.2 0.8 0.2 5 1 C4.7 2 4.3 3 4 4 C2.7 3.7 1.4 3.3 0 3 C0 2 0 1 0 0 Z };
\path[fill=black!55,even odd rule,shift={(85,133)}] svg {M0 0 C0.7 0.2 0.7 0.2 4 1 C3.7 2 3.3 3 3 4 C1.7 4 0.4 4 -1 4 C-0.7 2.7 -0.3 1.4 0 0 Z };
\path[fill=black!55,even odd rule,shift={(63,133)}] svg {M0 0 C0.7 0.2 0.7 0.2 4 1 C4 2 4 3 4 4 C2.3 3.7 0.7 3.3 -1 3 C-0.7 2 -0.3 1 0 0 Z };
\path[fill=black!63,even odd rule,shift={(56,133)}] svg {M0 0 C0.7 0.2 0.7 0.2 4 1 C3.7 2 3.3 3 3 4 C1.7 4 0.4 4 -1 4 C-0.7 2.7 -0.3 1.4 0 0 Z };
\path[fill=black!59,even odd rule,shift={(341,132)}] svg {M0 0 C1.3 0 2.6 0 4 0 C3.7 1.3 3.3 2.6 3 4 C1.7 3.7 0.4 3.3 -1 3 C-0.7 2 -0.3 1 0 0 Z };
\path[fill=black!59,even odd rule,shift={(334,132)}] svg {M0 0 C1.3 0 2.6 0 4 0 C3.7 1.3 3.3 2.6 3 4 C1.7 3.7 0.4 3.3 -1 3 C-0.7 2 -0.3 1 0 0 Z };
\path[fill=black!55,even odd rule,shift={(312,132)}] svg {M0 0 C1.3 0 2.6 0 4 0 C3.7 1.3 3.3 2.6 3 4 C1.7 3.7 0.4 3.3 -1 3 C-0.7 2 -0.3 1 0 0 Z };
\path[fill=black!59,even odd rule,shift={(290,132)}] svg {M0 0 C1.3 0 2.6 0 4 0 C3.7 1.3 3.3 2.6 3 4 C1.7 3.7 0.4 3.3 -1 3 C-0.7 2 -0.3 1 0 0 Z };
\path[fill=black!63,even odd rule,shift={(276,132)}] svg {M0 0 C0.7 0.2 0.7 0.2 4 1 C3.7 2 3.3 3 3 4 C1.7 4 0.4 4 -1 4 C-0.7 2.7 -0.3 1.4 0 0 Z };
\path[fill=black!59,even odd rule,shift={(254,132)}] svg {M0 0 C0.7 0.2 0.7 0.2 4 1 C3.7 2 3.3 3 3 4 C1.7 4 0.4 4 -1 4 C-0.7 2.7 -0.3 1.4 0 0 Z };
\path[fill=black!63,even odd rule,shift={(188,112)}] svg {M0 0 C0.7 0.2 0.7 0.2 4 1 C4 2 4 3 4 4 C2.3 4 0.7 4 -1 4 C-0.7 2.7 -0.3 1.4 0 0 Z };
\path[fill=black!63,even odd rule,shift={(121,113)}] svg {M0 0 C1.7 0 3.3 0 5 0 C4.7 1.3 4.3 2.6 4 4 C2.7 3.7 1.4 3.3 0 3 C0 2 0 1 0 0 Z };
\path[fill=black!55,even odd rule,shift={(93,113)}] svg {M0 0 C1.3 0 2.6 0 4 0 C3.7 1.3 3.3 2.6 3 4 C1.7 3.7 0.4 3.3 -1 3 C-0.7 2 -0.3 1 0 0 Z };
\path[fill=black!55,even odd rule,shift={(78,113)}] svg {M0 0 C1.3 0 2.6 0 4 0 C3.7 1.3 3.3 2.6 3 4 C1.7 3.7 0.4 3.3 -1 3 C-0.7 2 -0.3 1 0 0 Z };
\path[fill=black!55,even odd rule,shift={(326,112)}] svg {M0 0 C0.8 0.2 0.8 0.2 5 1 C4.7 2 4.3 3 4 4 C2.7 3.7 1.4 3.3 0 3 C0 2 0 1 0 0 Z };
\path[fill=black!59,even odd rule,shift={(305,112)}] svg {M0 0 C1.3 0 2.6 0 4 0 C3.7 1.3 3.3 2.6 3 4 C1.7 3.7 0.4 3.3 -1 3 C-0.7 2 -0.3 1 0 0 Z };
\path[fill=black!55,even odd rule,shift={(298,112)}] svg {M0 0 C1.3 0 2.6 0 4 0 C3.7 1.3 3.3 2.6 3 4 C1.7 3.7 0.4 3.3 -1 3 C-0.7 2 -0.3 1 0 0 Z };
\path[fill=black!55,even odd rule,shift={(254,112)}] svg {M0 0 C0.7 0.2 0.7 0.2 4 1 C4 2 4 3 4 4 C2.4 3.7 0.7 3.3 -1 3 C-0.7 2 -0.3 1 0 0 Z };
\path[fill=black!55,even odd rule,shift={(247,112)}] svg {M0 0 C0.7 0.2 0.7 0.2 4 1 C3.7 2 3.3 3 3 4 C1.7 4 0.4 4 -1 4 C-0.7 2.7 -0.3 1.4 0 0 Z };
\path[fill=black!59,even odd rule,shift={(226,112)}] svg {M0 0 C1 0.3 2 0.7 3 1 C3 2 3 3 3 4 C1.3 4 -0.3 4 -2 4 C-2 3 -2 2 -2 1 C-1.3 0.7 -0.7 0.3 0 0 Z };
\path[fill=black!71,even odd rule,shift={(159,112)}] svg {M0 0 C0.7 0.2 0.7 0.2 4 1 C4 2 4 3 4 4 C2.3 4 0.7 4 -1 4 C-0.7 2.7 -0.3 1.4 0 0 Z };
\path[fill=black!55,even odd rule,shift={(334,47)}] svg {M0 0 C1.3 0 2.6 0 4 0 C3.7 1.3 3.3 2.6 3 4 C1.7 3.7 0.4 3.3 -1 3 C-0.7 2 -0.3 1 0 0 Z };
\path[fill=black!55,even odd rule,shift={(283,47)}] svg {M0 0 C0.7 0.2 0.7 0.2 4 1 C4 2 4 3 4 4 C2.4 3.7 0.7 3.3 -1 3 C-0.7 2 -0.3 1 0 0 Z };
\path[fill=black!55,even odd rule,shift={(276,47)}] svg {M0 0 C0.7 0.2 0.7 0.2 4 1 C3.7 2 3.3 3 3 4 C1.7 4 0.4 4 -1 4 C-0.7 2.7 -0.3 1.4 0 0 Z };
\path[fill=black!61,even odd rule,shift={(260,217)}] svg {M0 0 C1.3 0 2.6 0 4 0 C3.7 1.3 3.3 2.6 3 4 C1.7 3.7 0.4 3.3 -1 3 C-0.7 2 -0.3 1 0 0 Z };
\path[fill=black!61,even odd rule,shift={(371,216)}] svg {M0 0 C1 0.3 2 0.7 3 1 C2.7 2 2.3 3 2 4 C0.4 3.7 -1.3 3.3 -3 3 C-2 2 -1 1 0 0 Z };
\path[fill=black!65,even odd rule,shift={(4,195)}] svg {M0 0 C0.7 0.3 1.3 0.7 2 1 C2 2.7 2 4.3 2 6 C1 6 0 6 -1 6 C-1 4.3 -1 2.7 -1 1 C-0.7 0.7 -0.3 0.3 0 0 Z };
\path[fill=black!56,even odd rule,shift={(259,192)}] svg {M0 0 C1 0.3 2 0.7 3 1 C3 2 3 3 3 4 C1.4 4 -0.3 4 -2 4 C-2 3 -2 2 -2 1 C-1.3 0.7 -0.7 0.3 0 0 Z };
\path[fill=black!61,even odd rule,shift={(302,192)}] svg {M0 0 C0.7 0.2 0.7 0.2 4 1 C3.7 2 3.3 3 3 4 C1.7 3.7 0.4 3.3 -1 3 C-0.7 2 -0.3 1 0 0 Z };
\path[fill=black!56,even odd rule,shift={(232,172)}] svg {M0 0 C1 0.3 2 0.7 3 1 C3 2 3 3 3 4 C1.3 4 -0.3 4 -2 4 C-2 3 -2 2 -2 1 C-1.3 0.7 -0.7 0.3 0 0 Z };
\path[fill=black!61,even odd rule,shift={(48,173)}] svg {M0 0 C0.7 0.2 0.7 0.2 4 1 C4 2 4 3 4 4 C2.4 4 0.7 4 -1 4 C-0.7 2.7 -0.3 1.4 0 0 Z };
\path[fill=black!56,even odd rule,shift={(255,172)}] svg {M0 0 C0.7 0.3 1.3 0.7 2 1 C2 2 2 3 2 4 C0.3 4 -1.3 4 -3 4 C-3 3 -3 2 -3 1 C-2 0.7 -1 0.3 0 0 Z };
\path[fill=black!65,even odd rule,shift={(226,172)}] svg {M0 0 C0.7 0.3 1.3 0.7 2 1 C2 2 2 3 2 4 C0.3 4 -1.3 4 -3 4 C-3 3 -3 2 -3 1 C-2 0.7 -1 0.3 0 0 Z };
\path[fill=black!65,even odd rule,shift={(172,173)}] svg {M0 0 C1.7 0.3 3.3 0.7 5 1 C5 1.7 5 2.3 5 3 C3 3 1 3 -1 3 C-0.7 2 -0.3 1 0 0 Z };
\path[fill=black!65,even odd rule,shift={(150,173)}] svg {M0 0 C1.3 0 2.6 0 4 0 C3.7 1.3 3.3 2.6 3 4 C1.7 3.7 0.4 3.3 -1 3 C-0.7 2 -0.3 1 0 0 Z };
\path[fill=black!61,even odd rule,shift={(135,173)}] svg {M0 0 C0.8 0.2 0.8 0.2 5 1 C4.7 2 4.3 3 4 4 C2.7 3.7 1.4 3.3 0 3 C0 2 0 1 0 0 Z };
\path[fill=black!56,even odd rule,shift={(99,173)}] svg {M0 0 C0.7 0.2 0.7 0.2 4 1 C3.7 2 3.3 3 3 4 C1.7 3.7 0.4 3.3 -1 3 C-0.7 2 -0.3 1 0 0 Z };
\path[fill=black!65,even odd rule,shift={(70,173)}] svg {M0 0 C0.7 0.2 0.7 0.2 4 1 C3.7 2 3.3 3 3 4 C1.7 4 0.4 4 -1 4 C-0.7 2.7 -0.3 1.4 0 0 Z };
\path[fill=black!61,even odd rule,shift={(379,172)}] svg {M0 0 C1 0 2 0 3 0 C2.7 1.7 2.3 3.3 2 5 C1.3 5 0.7 5 0 5 C-0.6 3.1 -0.6 3.1 -1 1 C-0.7 0.7 -0.3 0.3 0 0 Z };
\path[fill=black!56,even odd rule,shift={(326,172)}] svg {M0 0 C0.7 0.2 0.7 0.2 4 1 C3.7 2 3.3 3 3 4 C1.7 3.7 0.4 3.3 -1 3 C-0.7 2 -0.3 1 0 0 Z };
\path[fill=black!56,even odd rule,shift={(275,172)}] svg {M0 0 C0.7 0.2 0.7 0.2 4 1 C3.7 2 3.3 3 3 4 C1.7 3.7 0.4 3.3 -1 3 C-0.7 2 -0.3 1 0 0 Z };
\path[fill=black!65,even odd rule,shift={(239,172)}] svg {M0 0 C1 0.3 2 0.7 3 1 C3 2 3 3 3 4 C1.3 4 -0.3 4 -2 4 C-1.3 2.7 -0.7 1.4 0 0 Z };
\path[fill=black!65,even odd rule,shift={(380,157)}] svg {M0 0 C0.7 0.3 1.3 0.7 2 1 C2 2.3 2 3.6 2 5 C0.7 5 -0.6 5 -2 5 C-1.1 1.1 -1.1 1.1 0 0 Z };
\path[fill=black!61,even odd rule,shift={(148,153)}] svg {M0 0 C1.7 0 3.3 0 5 0 C5 1 5 2 5 3 C4 3.3 3 3.7 2 4 C1.3 3.7 0.7 3.3 0 3 C0 2 0 1 0 0 Z };
\path[fill=black!56,even odd rule,shift={(361,152)}] svg {M0 0 C1.3 0 2.6 0 4 0 C4 1 4 2 4 3 C2.7 3.3 1.4 3.7 0 4 C0 2.7 0 1.4 0 0 Z };
\path[fill=black!56,even odd rule,shift={(332,152)}] svg {M0 0 C0.7 0.2 0.7 0.2 4 1 C3.7 2 3.3 3 3 4 C1.7 3.7 0.4 3.3 -1 3 C-0.7 2 -0.3 1 0 0 Z };
\path[fill=black!56,even odd rule,shift={(310,152)}] svg {M0 0 C0.7 0.2 0.7 0.2 4 1 C3.7 2 3.3 3 3 4 C1.7 3.7 0.4 3.3 -1 3 C-0.7 2 -0.3 1 0 0 Z };
\path[fill=black!56,even odd rule,shift={(281,152)}] svg {M0 0 C0.7 0.2 0.7 0.2 4 1 C3.7 2 3.3 3 3 4 C1.7 3.7 0.4 3.3 -1 3 C-0.7 2 -0.3 1 0 0 Z };
\path[fill=black!61,even odd rule,shift={(274,152)}] svg {M0 0 C0.7 0.2 0.7 0.2 4 1 C3.7 2 3.3 3 3 4 C1.7 4 0.4 4 -1 4 C-0.7 2.7 -0.3 1.4 0 0 Z };
\path[fill=black!61,even odd rule,shift={(253,152)}] svg {M0 0 C1 0.3 2 0.7 3 1 C2.7 2 2.3 3 2 4 C0.7 4 -0.6 4 -2 4 C-2 3 -2 2 -2 1 C-1.3 0.7 -0.7 0.3 0 0 Z };
\path[fill=black!65,even odd rule,shift={(143,133)}] svg {M0 0 C0.8 0.2 0.8 0.2 5 1 C5 1.7 5 2.3 5 3 C3.7 3.3 2.4 3.7 1 4 C0.7 2.7 0.3 1.4 0 0 Z };
\path[fill=black!61,even odd rule,shift={(129,133)}] svg {M0 0 C1.3 0 2.6 0 4 0 C3.7 1.3 3.3 2.6 3 4 C1.7 3.7 0.4 3.3 -1 3 C-0.7 2 -0.3 1 0 0 Z };
\path[fill=black!61,even odd rule,shift={(12,133)}] svg {M0 0 C0.7 0.2 0.7 0.2 4 1 C3.7 2 3.3 3 3 4 C1.7 4 0.4 4 -1 4 C-0.7 2.7 -0.3 1.4 0 0 Z };
\path[fill=black!61,even odd rule,shift={(348,132)}] svg {M0 0 C1.7 0.3 3.3 0.6 5 1 C5 1.7 5 2.3 5 3 C3.7 3.3 2.4 3.7 1 4 C0.3 3 -0.3 2 -1 1 C-0.7 0.7 -0.3 0.3 0 0 Z };
\path[fill=black!56,even odd rule,shift={(305,132)}] svg {M0 0 C0.7 0.2 0.7 0.2 4 1 C3.7 2 3.3 3 3 4 C1.7 3.7 0.4 3.3 -1 3 C-0.7 2 -0.3 1 0 0 Z };
\path[fill=black!65,even odd rule,shift={(240,132)}] svg {M0 0 C1 0.3 2 0.7 3 1 C3 2 3 3 3 4 C1.3 3.7 -0.3 3.3 -2 3 C-2 2.3 -2 1.7 -2 1 C-1.3 0.7 -0.7 0.3 0 0 Z };
\path[fill=black!70,even odd rule,shift={(188,132)}] svg {M0 0 C1.3 0.7 2.6 1.3 4 2 C4 2.7 4 3.3 4 4 C2.3 4 0.7 4 -1 4 C-0.7 2.7 -0.3 1.4 0 0 Z };
\path[fill=black!65,even odd rule,shift={(27,113)}] svg {M0 0 C0.7 0.2 0.7 0.2 4 1 C3.7 2 3.3 3 3 4 C1.7 4 0.4 4 -1 4 C-0.7 2.7 -0.3 1.4 0 0 Z };
\path[fill=black!61,even odd rule,shift={(183,112)}] svg {M0 0 C0.7 0.3 1.3 0.7 2 1 C2 2 2 3 2 4 C0.3 4 -1.3 4 -3 4 C-3 3 -3 2 -3 1 C-2 0.7 -1 0.3 0 0 Z };
\path[fill=black!65,even odd rule,shift={(49,113)}] svg {M0 0 C0.7 0.2 0.7 0.2 4 1 C3.7 2 3.3 3 3 4 C1.7 3.7 0.4 3.3 -1 3 C-0.7 2 -0.3 1 0 0 Z };
\path[fill=black!61,even odd rule,shift={(334,112)}] svg {M0 0 C1.3 0 2.6 0 4 0 C3.7 1.3 3.3 2.6 3 4 C1.7 3.7 0.4 3.3 -1 3 C-0.7 2 -0.3 1 0 0 Z };
\path[fill=black!56,even odd rule,shift={(276,112)}] svg {M0 0 C0.7 0.2 0.7 0.2 4 1 C3.7 2 3.3 3 3 4 C1.7 3.7 0.4 3.3 -1 3 C-0.7 2 -0.3 1 0 0 Z };
\path[fill=black!61,even odd rule,shift={(233,112)}] svg {M0 0 C1 0.3 2 0.7 3 1 C3 2 3 3 3 4 C1.3 3.7 -0.3 3.3 -2 3 C-2 2.3 -2 1.7 -2 1 C-1.3 0.7 -0.7 0.3 0 0 Z };
\path[fill=black!61,even odd rule,shift={(348,47)}] svg {M0 0 C0.8 0.2 0.8 0.2 5 1 C4.7 2 4.3 3 4 4 C2.7 3.7 1.4 3.3 0 3 C0 2 0 1 0 0 Z };
\path[fill=black!56,even odd rule,shift={(305,47)}] svg {M0 0 C0.7 0.2 0.7 0.2 4 1 C3.7 2 3.3 3 3 4 C1.7 3.7 0.4 3.3 -1 3 C-0.7 2 -0.3 1 0 0 Z };
\path[fill=black!56,even odd rule,shift={(298,47)}] svg {M0 0 C0.7 0.2 0.7 0.2 4 1 C3.7 2 3.3 3 3 4 C1.7 3.7 0.4 3.3 -1 3 C-0.7 2 -0.3 1 0 0 Z };
\path[fill=black!65,even odd rule,shift={(290,47)}] svg {M0 0 C0.7 0.2 0.7 0.2 4 1 C3.7 2 3.3 3 3 4 C1.7 3.7 0.4 3.3 -1 3 C-0.7 2 -0.3 1 0 0 Z };
\path[fill=black!63,even odd rule,shift={(3,225)}] svg {M0 0 C1 0 2 0 3 0 C3 1.7 3 3.3 3 5 C2 5 1 5 0 5 C0 3.3 0 1.7 0 0 Z };
\path[fill=black!63,even odd rule,shift={(347,217)}] svg {M0 0 C1.6 0 3.3 0 5 0 C5 1 5 2 5 3 C3.4 3 1.7 3 0 3 C0 2 0 1 0 0 Z };
\path[fill=black!63,even odd rule,shift={(339,217)}] svg {M0 0 C1.6 0 3.3 0 5 0 C5 1 5 2 5 3 C3.4 2.7 1.7 2.3 0 2 C0 1.3 0 0.7 0 0 Z };
\path[fill=black!72,even odd rule,shift={(332,217)}] svg {M0 0 C1.6 0 3.3 0 5 0 C5 1 5 2 5 3 C3.4 3 1.7 3 0 3 C0 2 0 1 0 0 Z };
\path[fill=black!72,even odd rule,shift={(325,217)}] svg {M0 0 C1.6 0 3.3 0 5 0 C5 1 5 2 5 3 C3.4 3 1.7 3 0 3 C0 2 0 1 0 0 Z };
\path[fill=black!67,even odd rule,shift={(310,217)}] svg {M0 0 C1.6 0 3.3 0 5 0 C5 1 5 2 5 3 C3.4 3 1.7 3 0 3 C0 2 0 1 0 0 Z };
\path[fill=black!67,even odd rule,shift={(303,217)}] svg {M0 0 C1.6 0 3.3 0 5 0 C5 1 5 2 5 3 C3.4 3 1.7 3 0 3 C0 2 0 1 0 0 Z };
\path[fill=black!67,even odd rule,shift={(296,217)}] svg {M0 0 C1.6 0 3.3 0 5 0 C5 1 5 2 5 3 C3.4 3 1.7 3 0 3 C0 2 0 1 0 0 Z };
\path[fill=black!67,even odd rule,shift={(282,217)}] svg {M0 0 C1.3 0 2.6 0 4 0 C4 1 4 2 4 3 C2.7 3.3 1.4 3.7 0 4 C0 2.7 0 1.4 0 0 Z };
\path[fill=black!63,even odd rule,shift={(363,216)}] svg {M0 0 C1 0.3 2 0.7 3 1 C3 2 3 3 3 4 C1.4 3.7 -0.3 3.3 -2 3 C-1.3 2 -0.7 1 0 0 Z };
\path[fill=black!63,even odd rule,shift={(380,208)}] svg {M0 0 C0.7 0.3 1.3 0.7 2 1 C1.7 2.7 1.3 4.3 1 6 C0.3 6 -0.3 6 -1 6 C-1 4.3 -1 2.7 -1 1 C-0.7 0.7 -0.3 0.3 0 0 Z };
\path[fill=black!58,even odd rule,shift={(3,203)}] svg {M0 0 C1 0 2 0 3 0 C3 1.7 3 3.3 3 5 C2 5 1 5 0 5 C0 3.3 0 1.7 0 0 Z };
\path[fill=black!63,even odd rule,shift={(379,201)}] svg {M0 0 C1 0 2 0 3 0 C3 1.7 3 3.3 3 5 C2 5 1 5 0 5 C0 3.3 0 1.7 0 0 Z };
\path[fill=black!58,even odd rule,shift={(379,194)}] svg {M0 0 C1 0 2 0 3 0 C2.7 1.7 2.3 3.3 2 5 C1.3 5 0.7 5 0 5 C-0.6 3.1 -0.6 3.1 -1 1 C-0.7 0.7 -0.3 0.3 0 0 Z };
\path[fill=black!63,even odd rule,shift={(235,193)}] svg {M0 0 C1.7 0 3.3 0 5 0 C5 1 5 2 5 3 C3.3 3 1.7 3 0 3 C0 2 0 1 0 0 Z };
\path[fill=black!77,even odd rule,shift={(213,193)}] svg {M0 0 C1.7 0 3.3 0 5 0 C5 1 5 2 5 3 C3.3 3 1.7 3 0 3 C0 2 0 1 0 0 Z };
\path[fill=black!82,even odd rule,shift={(191,193)}] svg {M0 0 C1.7 0 3.3 0 5 0 C5 1 5 2 5 3 C3.3 3 1.7 3 0 3 C0 2 0 1 0 0 Z };
\path[fill=black!77,even odd rule,shift={(184,193)}] svg {M0 0 C1.7 0 3.3 0 5 0 C5 1 5 2 5 3 C3.3 3 1.7 3 0 3 C0 2 0 1 0 0 Z };
\path[fill=black!77,even odd rule,shift={(169,193)}] svg {M0 0 C1.7 0 3.3 0 5 0 C5 1 5 2 5 3 C3.3 3 1.7 3 0 3 C0 2 0 1 0 0 Z };
\path[fill=black!63,even odd rule,shift={(162,193)}] svg {M0 0 C0.8 0.2 0.8 0.2 5 1 C5 1.7 5 2.3 5 3 C4 3.3 3 3.7 2 4 C1.3 3.7 0.7 3.3 0 3 C0 2 0 1 0 0 Z };
\path[fill=black!67,even odd rule,shift={(295,192)}] svg {M0 0 C1.3 0.7 2.6 1.3 4 2 C3.7 2.7 3.3 3.3 3 4 C1.7 4 0.4 4 -1 4 C-0.7 2.7 -0.3 1.4 0 0 Z };
\path[fill=black!63,even odd rule,shift={(252,192)}] svg {M0 0 C1 0.3 2 0.7 3 1 C2.7 2 2.3 3 2 4 C0.7 4 -0.6 4 -2 4 C-2 3 -2 2 -2 1 C-1.3 0.7 -0.7 0.3 0 0 Z };
\path[fill=black!63,even odd rule,shift={(4,188)}] svg {M0 0 C0.7 0.3 1.3 0.7 2 1 C1.7 2.7 1.3 4.3 1 6 C0 5 -1 4 -2 3 C-1.3 2 -0.7 1 0 0 Z };
\path[fill=black!63,even odd rule,shift={(379,187)}] svg {M0 0 C1 0 2 0 3 0 C2.7 1.7 2.3 3.3 2 5 C1 4.7 0 4.3 -1 4 C-0.7 2.7 -0.3 1.4 0 0 Z };
\path[fill=black!63,even odd rule,shift={(3,174)}] svg {M0 0 C1 0.3 2 0.7 3 1 C2.7 2.3 2.3 3.6 2 5 C1 4.7 0 4.3 -1 4 C-0.7 2.7 -0.3 1.4 0 0 Z };
\path[fill=black!63,even odd rule,shift={(262,172)}] svg {M0 0 C0.7 0.3 1.3 0.7 2 1 C2 2 2 3 2 4 C0.4 3.7 -1.3 3.3 -3 3 C-3 2.3 -3 1.7 -3 1 C-2 0.7 -1 0.3 0 0 Z };
\path[fill=black!77,even odd rule,shift={(208,173)}] svg {M0 0 C1.7 0 3.3 0 5 0 C5 1 5 2 5 3 C3.3 3 1.7 3 0 3 C0 2 0 1 0 0 Z };
\path[fill=black!67,even odd rule,shift={(193,173)}] svg {M0 0 C1.7 0 3.3 0 5 0 C5 1 5 2 5 3 C3.3 3 1.7 3 0 3 C0 2 0 1 0 0 Z };
\path[fill=black!77,even odd rule,shift={(186,173)}] svg {M0 0 C1.7 0 3.3 0 5 0 C5 1 5 2 5 3 C3.3 3 1.7 3 0 3 C0 2 0 1 0 0 Z };
\path[fill=black!63,even odd rule,shift={(165,173)}] svg {M0 0 C1.3 0 2.6 0 4 0 C3.7 1.3 3.3 2.6 3 4 C1.7 3.7 0.4 3.3 -1 3 C-0.7 2 -0.3 1 0 0 Z };
\path[fill=black!67,even odd rule,shift={(340,172)}] svg {M0 0 C1.3 0 2.6 0 4 0 C3.7 1.3 3.3 2.6 3 4 C2 4 1 4 0 4 C0 2.7 0 1.4 0 0 Z };
\path[fill=black!63,even odd rule,shift={(246,172)}] svg {M0 0 C1.3 0.7 2.6 1.3 4 2 C3.7 2.7 3.3 3.3 3 4 C1.7 4 0.4 4 -1 4 C-0.7 2.7 -0.3 1.4 0 0 Z };
\path[fill=black!63,even odd rule,shift={(258,153)}] svg {M0 0 C1.6 0 3.3 0 5 0 C5 1 5 2 5 3 C3.4 3 1.7 3 0 3 C0 2 0 1 0 0 Z };
\path[fill=black!58,even odd rule,shift={(236,153)}] svg {M0 0 C1.7 0 3.3 0 5 0 C5 1 5 2 5 3 C3.3 3 1.7 3 0 3 C0 2 0 1 0 0 Z };
\path[fill=black!63,even odd rule,shift={(230,152)}] svg {M0 0 C0.7 0.2 0.7 0.2 4 1 C3.7 2 3.3 3 3 4 C1.7 4 0.4 4 -1 4 C-0.7 2.7 -0.3 1.4 0 0 Z };
\path[fill=black!72,even odd rule,shift={(214,153)}] svg {M0 0 C1.7 0 3.3 0 5 0 C5 1 5 2 5 3 C3.3 3 1.7 3 0 3 C0 2 0 1 0 0 Z };
\path[fill=black!77,even odd rule,shift={(207,153)}] svg {M0 0 C1.7 0 3.3 0 5 0 C5 1 5 2 5 3 C3.3 3 1.7 3 0 3 C0 2 0 1 0 0 Z };
\path[fill=black!72,even odd rule,shift={(192,153)}] svg {M0 0 C1.7 0 3.3 0 5 0 C5 1 5 2 5 3 C3.3 3 1.7 3 0 3 C0 2 0 1 0 0 Z };
\path[fill=black!77,even odd rule,shift={(185,153)}] svg {M0 0 C1.7 0 3.3 0 5 0 C5 1 5 2 5 3 C3.3 3 1.7 3 0 3 C0 2 0 1 0 0 Z };
\path[fill=black!72,even odd rule,shift={(178,153)}] svg {M0 0 C1.7 0.3 3.3 0.7 5 1 C4.7 1.7 4.3 2.3 4 3 C2.3 3 0.7 3 -1 3 C-0.7 2 -0.3 1 0 0 Z };
\path[fill=black!63,even odd rule,shift={(170,153)}] svg {M0 0 C1.7 0 3.3 0 5 0 C5 1 5 2 5 3 C3.3 3 1.7 3 0 3 C0 2 0 1 0 0 Z };
\path[fill=black!72,even odd rule,shift={(163,153)}] svg {M0 0 C1.7 0 3.3 0 5 0 C5 1 5 2 5 3 C3.3 3 1.7 3 0 3 C0 2 0 1 0 0 Z };
\path[fill=black!67,even odd rule,shift={(134,153)}] svg {M0 0 C1.3 0 2.6 0 4 0 C3.7 1.3 3.3 2.6 3 4 C2 4 1 4 0 4 C0 2.7 0 1.4 0 0 Z };
\path[fill=black!67,even odd rule,shift={(112,153)}] svg {M0 0 C0.8 0.2 0.8 0.2 5 1 C4.7 1.5 4.7 1.5 3 4 C2 3.7 1 3.3 0 3 C0 2 0 1 0 0 Z };
\path[fill=black!72,even odd rule,shift={(25,153)}] svg {M0 0 C1 0.3 2 0.7 3 1 C3 2 3 3 3 4 C1.4 3.7 -0.3 3.3 -2 3 C-2 2.3 -2 1.7 -2 1 C-1.3 0.7 -0.7 0.3 0 0 Z };
\path[fill=black!63,even odd rule,shift={(267,152)}] svg {M0 0 C1 0.3 2 0.7 3 1 C3 2 3 3 3 4 C1.4 3.7 -0.3 3.3 -2 3 C-2 2.3 -2 1.7 -2 1 C-1.3 0.7 -0.7 0.3 0 0 Z };
\path[fill=black!67,even odd rule,shift={(202,133)}] svg {M0 0 C1.7 0.3 3.3 0.7 5 1 C4.7 1.7 4.3 2.3 4 3 C2.3 3 0.7 3 -1 3 C-0.7 2 -0.3 1 0 0 Z };
\path[fill=black!77,even odd rule,shift={(194,133)}] svg {M0 0 C1.7 0 3.3 0 5 0 C5 1 5 2 5 3 C3.3 3 1.7 3 0 3 C0 2 0 1 0 0 Z };
\path[fill=black!67,even odd rule,shift={(172,133)}] svg {M0 0 C1.7 0 3.3 0 5 0 C5 1 5 2 5 3 C3.3 3 1.7 3 0 3 C0 2 0 1 0 0 Z };
\path[fill=black!63,even odd rule,shift={(158,133)}] svg {M0 0 C1.7 0.3 3.3 0.7 5 1 C4.7 1.7 4.3 2.3 4 3 C2.3 3 0.7 3 -1 3 C-0.7 2 -0.3 1 0 0 Z };
\path[fill=black!72,even odd rule,shift={(370,132)}] svg {M0 0 C1.3 0 2.6 0 4 0 C4 1 4 2 4 3 C2.7 3.3 1.4 3.7 0 4 C0 2.7 0 1.4 0 0 Z };
\path[fill=black!67,even odd rule,shift={(355,132)}] svg {M0 0 C1.6 0 3.3 0 5 0 C5 1 5 2 5 3 C3.4 3 1.7 3 0 3 C0 2 0 1 0 0 Z };
\path[fill=black!72,even odd rule,shift={(211,132)}] svg {M0 0 C1 0.3 2 0.7 3 1 C2.7 2 2.3 3 2 4 C0.7 4 -0.6 4 -2 4 C-2 3 -2 2 -2 1 C-1.3 0.7 -0.7 0.3 0 0 Z };
\path[fill=black!58,even odd rule,shift={(202,113)}] svg {M0 0 C1.7 0 3.3 0 5 0 C5 1 5 2 5 3 C3.3 3 1.7 3 0 3 C0 2 0 1 0 0 Z };
\path[fill=black!72,even odd rule,shift={(165,113)}] svg {M0 0 C1.7 0 3.3 0 5 0 C5 1 5 2 5 3 C3.3 3 1.7 3 0 3 C0 2 0 1 0 0 Z };
\path[fill=black!72,even odd rule,shift={(143,113)}] svg {M0 0 C1.7 0 3.3 0 5 0 C5 1 5 2 5 3 C3.3 3 1.7 3 0 3 C0 2 0 1 0 0 Z };
\path[fill=black!72,even odd rule,shift={(136,113)}] svg {M0 0 C1.7 0 3.3 0 5 0 C5 1 5 2 5 3 C3.3 3 1.7 3 0 3 C0 2 0 1 0 0 Z };
\path[fill=black!67,even odd rule,shift={(85,113)}] svg {M0 0 C1.3 0 2.6 0 4 0 C3.7 1.3 3.3 2.6 3 4 C2 4 1 4 0 4 C0 2.7 0 1.4 0 0 Z };
\path[fill=black!67,even odd rule,shift={(20,113)}] svg {M0 0 C1 0.3 2 0.7 3 1 C3 2 3 3 3 4 C1.4 3.7 -0.3 3.3 -2 3 C-1.3 2 -0.7 1 0 0 Z };
\path[fill=black!67,even odd rule,shift={(312,112)}] svg {M0 0 C1.3 0 2.6 0 4 0 C4 1.3 4 2.6 4 4 C2.7 3.7 1.4 3.3 0 3 C0 2 0 1 0 0 Z };
\path[fill=black!67,even odd rule,shift={(197,112)}] svg {M0 0 C1 1 2 2 3 3 C1.3 3.3 -0.3 3.7 -2 4 C-2 3 -2 2 -2 1 C-1.3 0.7 -0.7 0.3 0 0 Z };
\path[fill=black!63,even odd rule,shift={(386,93)}] svg {M0 0 C1.6 0 3.3 0 5 0 C5 1 5 2 5 3 C3.4 3 1.7 3 0 3 C0 2 0 1 0 0 Z };
\path[fill=black!58,even odd rule,shift={(379,93)}] svg {M0 0 C1.6 0 3.3 0 5 0 C5 1 5 2 5 3 C3.4 3 1.7 3 0 3 C0 2 0 1 0 0 Z };
\path[fill=black!72,even odd rule,shift={(384,47)}] svg {M0 0 C1.6 0 3.3 0 5 0 C5 1 5 2 5 3 C3.4 3 1.7 3 0 3 C0 2 0 1 0 0 Z };
\path[fill=black!63,even odd rule,shift={(377,47)}] svg {M0 0 C1.6 0 3.3 0 5 0 C5 1 5 2 5 3 C3.4 3 1.7 3 0 3 C0 2 0 1 0 0 Z };
\path[fill=black!67,even odd rule,shift={(341,47)}] svg {M0 0 C1.3 0 2.6 0 4 0 C4 1 4 2 4 3 C2.7 3.3 1.4 3.7 0 4 C-0.3 3.3 -0.7 2.7 -1 2 C-0.7 1.3 -0.3 0.7 0 0 Z };
\path[fill=black!65,even odd rule,shift={(3,218)}] svg {M0 0 C1 0 2 0 3 0 C2.7 1.7 2.3 3.3 2 5 C1.3 5 0.7 5 0 5 C0 3.3 0 1.7 0 0 Z };
\path[fill=black!70,even odd rule,shift={(319,216)}] svg {M0 0 C1.3 0.7 2.6 1.3 4 2 C3.7 2.7 3.3 3.3 3 4 C1.7 4 0.4 4 -1 4 C-0.7 2.7 -0.3 1.4 0 0 Z };
\path[fill=black!65,even odd rule,shift={(264,193)}] svg {M0 0 C1.6 0 3.3 0 5 0 C5 1 5 2 5 3 C3.4 2.7 1.7 2.3 0 2 C0 1.3 0 0.7 0 0 Z };
\path[fill=black!70,even odd rule,shift={(228,193)}] svg {M0 0 C1.7 0 3.3 0 5 0 C4.7 1 4.3 2 4 3 C2.7 3 1.4 3 0 3 C0 2 0 1 0 0 Z };
\path[fill=black!75,even odd rule,shift={(220,193)}] svg {M0 0 C1.7 0 3.3 0 5 0 C5 1 5 2 5 3 C3.3 2.7 1.7 2.3 0 2 C0 1.3 0 0.7 0 0 Z };
\path[fill=black!80,even odd rule,shift={(206,193)}] svg {M0 0 C1.7 0 3.3 0 5 0 C4.7 1 4.3 2 4 3 C2.7 3 1.4 3 0 3 C0 2 0 1 0 0 Z };
\path[fill=black!70,even odd rule,shift={(379,179)}] svg {M0 0 C1 0.3 2 0.7 3 1 C3 2.3 3 3.6 3 5 C2 5 1 5 0 5 C0 3.3 0 1.7 0 0 Z };
\path[fill=black!75,even odd rule,shift={(201,173)}] svg {M0 0 C0.8 0.2 0.8 0.2 5 1 C5 1.7 5 2.3 5 3 C3.3 3 1.7 3 0 3 C0 2 0 1 0 0 Z };
\path[fill=black!75,even odd rule,shift={(179,173)}] svg {M0 0 C0.8 0.2 0.8 0.2 5 1 C5 1.7 5 2.3 5 3 C3.3 3 1.7 3 0 3 C0 2 0 1 0 0 Z };
\path[fill=black!70,even odd rule,shift={(290,172)}] svg {M0 0 C1 0.3 2 0.7 3 1 C3 2 3 3 3 4 C1.7 4 0.4 4 -1 4 C-0.7 2.7 -0.3 1.4 0 0 Z };
\path[fill=black!65,even odd rule,shift={(268,172)}] svg {M0 0 C1 0.3 2 0.7 3 1 C3 2 3 3 3 4 C1.7 4 0.4 4 -1 4 C-0.7 2.7 -0.3 1.4 0 0 Z };
\path[fill=black!65,even odd rule,shift={(243,153)}] svg {M0 0 C1.7 0 3.3 0 5 0 C5 1 5 2 5 3 C3.3 2.7 1.7 2.3 0 2 C0 1.3 0 0.7 0 0 Z };
\path[fill=black!70,even odd rule,shift={(200,153)}] svg {M0 0 C1.7 0.3 3.3 0.7 5 1 C4.7 1.7 4.3 2.3 4 3 C2.3 2.7 0.7 2.3 -1 2 C-0.7 1.3 -0.3 0.7 0 0 Z };
\path[fill=black!65,even odd rule,shift={(141,153)}] svg {M0 0 C0.8 0.2 0.8 0.2 5 1 C5 1.7 5 2.3 5 3 C3.3 3 1.7 3 0 3 C0 2 0 1 0 0 Z };
\path[fill=black!70,even odd rule,shift={(231,133)}] svg {M0 0 C1.7 0 3.3 0 5 0 C4.7 1 4.3 2 4 3 C2.7 3 1.4 3 0 3 C0 2 0 1 0 0 Z };
\path[fill=black!65,even odd rule,shift={(224,133)}] svg {M0 0 C1.7 0.3 3.3 0.7 5 1 C4.7 1.7 4.3 2.3 4 3 C2.3 2.7 0.7 2.3 -1 2 C-0.7 1.3 -0.3 0.7 0 0 Z };
\path[fill=black!70,even odd rule,shift={(216,133)}] svg {M0 0 C1.7 0 3.3 0 5 0 C5 1 5 2 5 3 C3.3 2.7 1.7 2.3 0 2 C0 1.3 0 0.7 0 0 Z };
\path[fill=black!75,even odd rule,shift={(180,133)}] svg {M0 0 C1.7 0.3 3.3 0.7 5 1 C4.7 1.7 4.3 2.3 4 3 C2.3 2.7 0.7 2.3 -1 2 C-0.7 1.3 -0.3 0.7 0 0 Z };
\path[fill=black!75,even odd rule,shift={(165,133)}] svg {M0 0 C0.8 0.2 0.8 0.2 5 1 C5 1.7 5 2.3 5 3 C3.3 3 1.7 3 0 3 C0 2 0 1 0 0 Z };
\path[fill=black!65,even odd rule,shift={(209,113)}] svg {M0 0 C1.7 0 3.3 0 5 0 C5 1 5 2 5 3 C3.3 2.7 1.7 2.3 0 2 C0 1.3 0 0.7 0 0 Z };
\path[fill=black!65,even odd rule,shift={(100,113)}] svg {M0 0 C1.3 0 2.6 0 4 0 C4 1 4 2 4 3 C2.3 3 0.7 3 -1 3 C-0.7 2 -0.3 1 0 0 Z };
\path[fill=black!65,even odd rule,shift={(363,112)}] svg {M0 0 C1.6 0.3 3.3 0.7 5 1 C4.7 1.7 4.3 2.3 4 3 C2.4 2.7 0.7 2.3 -1 2 C-0.7 1.3 -0.3 0.7 0 0 Z };
\path[fill=black!60,even odd rule,shift={(356,112)}] svg {M0 0 C1.3 0 2.6 0 4 0 C4 1 4 2 4 3 C2.4 3 0.7 3 -1 3 C-0.7 2 -0.3 1 0 0 Z };
\path[fill=black!60,even odd rule,shift={(348,112)}] svg {M0 0 C0.8 0.2 0.8 0.2 5 1 C5 1.7 5 2.3 5 3 C3.4 3 1.7 3 0 3 C0 2 0 1 0 0 Z };
\path[fill=black!70,even odd rule,shift={(341,112)}] svg {M0 0 C0.8 0.2 0.8 0.2 5 1 C5 1.7 5 2.3 5 3 C3.4 3 1.7 3 0 3 C0 2 0 1 0 0 Z };
\path[fill=black!70,even odd rule,shift={(219,112)}] svg {M0 0 C0.7 0.3 1.3 0.7 2 1 C2 2 2 3 2 4 C0.3 3.7 -1.3 3.3 -3 3 C-2 2 -1 1 0 0 Z };
\path[fill=black!60,even odd rule,shift={(372,93)}] svg {M0 0 C1.6 0.3 3.3 0.7 5 1 C4.7 1.7 4.3 2.3 4 3 C2.4 2.7 0.7 2.3 -1 2 C-0.7 1.3 -0.3 0.7 0 0 Z };
\path[fill=black!65,even odd rule,shift={(370,47)}] svg {M0 0 C1.3 0 2.6 0 4 0 C4 1 4 2 4 3 C2.4 3 0.7 3 -1 3 C-0.7 2 -0.3 1 0 0 Z };
\path[fill=black!68,even odd rule,shift={(4,210)}] svg {M0 0 C0.7 0.3 1.3 0.7 2 1 C2 2.3 2 3.6 2 5 C1 5 0 5 -1 5 C-0.7 3.3 -0.3 1.7 0 0 Z };
\path[fill=black!68,even odd rule,shift={(272,193)}] svg {M0 0 C0.8 0.2 0.8 0.2 5 1 C4.7 1.7 4.3 2.3 4 3 C2.7 3 1.4 3 0 3 C0 2 0 1 0 0 Z };
\path[fill=black!68,even odd rule,shift={(243,193)}] svg {M0 0 C1.3 0 2.6 0 4 0 C4 1 4 2 4 3 C2.3 2.7 0.7 2.3 -1 2 C-0.7 1.3 -0.3 0.7 0 0 Z };
\path[fill=black!79,even odd rule,shift={(199,193)}] svg {M0 0 C1.3 0 2.6 0 4 0 C4 1 4 2 4 3 C2.3 2.7 0.7 2.3 -1 2 C-0.7 1.3 -0.3 0.7 0 0 Z };
\path[fill=black!84,even odd rule,shift={(177,193)}] svg {M0 0 C1.3 0 2.6 0 4 0 C4 1 4 2 4 3 C2.3 2.7 0.7 2.3 -1 2 C-0.7 1.3 -0.3 0.7 0 0 Z };
\path[fill=black!73,even odd rule,shift={(222,153)}] svg {M0 0 C1.3 0 2.6 0 4 0 C4 1 4 2 4 3 C2.3 2.7 0.7 2.3 -1 2 C-0.7 1.3 -0.3 0.7 0 0 Z };
\path[fill=black!73,even odd rule,shift={(173,113)}] svg {M0 0 C1.3 0 2.6 0 4 0 C4 1 4 2 4 3 C2.3 2.7 0.7 2.3 -1 2 C-0.7 1.3 -0.3 0.7 0 0 Z };
\path[fill=black!63,even odd rule,shift={(129,113)}] svg {M0 0 C1.3 0 2.6 0 4 0 C4 1 4 2 4 3 C2.3 2.7 0.7 2.3 -1 2 C-0.7 1.3 -0.3 0.7 0 0 Z };
\path[fill=black!63,even odd rule,shift={(363,47)}] svg {M0 0 C1.3 0 2.6 0 4 0 C4 1 4 2 4 3 C2.4 2.7 0.7 2.3 -1 2 C-0.7 1.3 -0.3 0.7 0 0 Z };
\path[fill=black!65,even odd rule,shift={(289,217)}] svg {M0 0 C1.3 0 2.6 0 4 0 C4 1 4 2 4 3 C2.7 3 1.4 3 0 3 C0 2 0 1 0 0 Z };
\path[fill=black!77,even odd rule,shift={(155,193)}] svg {M0 0 C1.3 0 2.6 0 4 0 C4 1 4 2 4 3 C2.7 3 1.4 3 0 3 C0 2 0 1 0 0 Z };
\path[fill=black!77,even odd rule,shift={(151,113)}] svg {M0 0 C1.3 0 2.6 0 4 0 C4 1 4 2 4 3 C2.7 3 1.4 3 0 3 C0 2 0 1 0 0 Z };
\path[fill=black!77,even odd rule,shift={(107,113)}] svg {M0 0 C1.3 0 2.6 0 4 0 C4 1 4 2 4 3 C2.7 3 1.4 3 0 3 C0 2 0 1 0 0 Z };
\path[fill=black!77,even odd rule,shift={(268,48)}] svg {M0 0 C1.3 0 2.6 0 4 0 C4 1 4 2 4 3 C2.7 3 1.4 3 0 3 C0 2 0 1 0 0 Z };
\end{tikzpicture}}
（长度单位为$\mathrm{m}$，质量单位为$\mathrm{kg}$，油的密度为常数$\rho\,\mathrm{kg}\,\mathrm{m}^{-3}$）
\QuestionEnd

\begin{QuestionSolution}{2010 年第 18 题}
\textbf{已知与目标}\quad 椭圆柱形油罐轴长为 \(l\)，横截面长、短半轴为 \(a,b\)，油深为 \(\frac{3b}{2}\)。密度为常数 \(\rho\)，质量等于油体积乘密度。

\textbf{方法与条件}\quad 平放后各垂直于罐轴的油截面相同，故体积为该截面积乘 \(l\)。取椭圆中心为原点、\(y\) 轴竖直向上，用水平窄条宽度积分求面积。

\textbf{推导}\quad 椭圆满足 \(\frac{x^2}{a^2}+\frac{y^2}{b^2}=1\)，底端为 \(y=-b\)，油面为 \(-b+\frac{3b}{2}=\frac{b}{2}\)。固定高度 \(y\) 时两侧横坐标为 \(\pm a\sqrt{1-\frac{y^2}{b^2}}\)，故条宽为其两倍。令 \(u=\frac{y}{b}\)、\(\mathrm{d}y=b\,\mathrm{d}u\)，得
\[
S=2ab\int_{-1}^{\frac{1}{2}}\sqrt{1-u^2}\,\mathrm{d}u
=ab\left[u\sqrt{1-u^2}+\arcsin u\right]_{-1}^{\frac{1}{2}}
=ab\left(\frac{\sqrt3}{4}+\frac{2\pi}{3}\right).
\]
上限代入为 \(\frac{\sqrt3}{4}+\frac{\pi}{6}\)，下限为 \(-\frac{\pi}{2}\)，相减便得上式。乘轴长 \(l\) 和密度 \(\rho\) 即得质量。

\textbf{结论与检查}\quad \(\boxed{m=\rho lab(\frac{\sqrt3}{4}+\frac{2\pi}{3})\,\mathrm{kg}}\)。因油深大于半罐而小于全罐，质量应在 \(\frac{1}{2}\rho l\pi ab\) 与 \(\rho l\pi ab\) 之间，所得系数满足这一范围。
\end{QuestionSolution}

% Q 2013 21
\QuestionPage{3.4 定积分的应用}{解答题}{2013 年第 21 题}{4}{弧长与平面图形的形心}{179 / 234}{***}
设曲线$L$的方程为$y = \frac{1}{4}x^2 - \frac{1}{2}\ln x$（$1 \le x \le \mathrm{e}$）．
\begin{enumerate}
\item 求$L$的弧长；
\item 设$D$是由曲线$L$，直线$x = 1,x = \mathrm{e}$及$x$轴所围平面图形，求$D$的形心的横坐标．
\end{enumerate}
\QuestionEnd

\begin{QuestionSolution}{2013 年第 21 题}
\textbf{已知与目标}\quad 曲线为 \(f(x)=\frac{x^2}{4}-\frac{1}{2}\ln x\)，\(1\le x\le\mathrm{e}\)。先求曲线长度，再求曲线下方区域形心的横坐标。形心是按面积均匀加权的位置平均值。

\textbf{方法与条件}\quad 弧长用 \(\int\sqrt{1+(f')^2}\,\mathrm{d}x\)。形心横坐标为 \(\overline{x}=\frac{M_y}{S}\)，其中 \(S\) 为面积，\(M_y=\int x\,\mathrm{d}S\) 是对 \(y\) 轴的一阶矩；先确认竖条高度为正。

\textbf{推导}\quad （1）\(f'(x)=\frac{1}{2}(x-\frac{1}{x})\)，因此 \(1+(f')^2=\frac{1}{4}(x+\frac{1}{x})^2\)。因 \(x>0\)，开根号得
\[
\sqrt{1+[f'(x)]^2}=\frac{x+x^{-1}}{2},
\qquad
s=\left[\frac{x^2}{4}+\frac12\ln x\right]_1^{\mathrm{e}}
=\boxed{\frac{\mathrm{e}^2+1}{4}}.
\]
弧长代入端点为 \(\frac{\mathrm{e}^2}{4}+\frac{1}{2}-\frac{1}{4}=\frac{\mathrm{e}^2+1}{4}\)。

（2）\(f(1)=\frac{1}{4}>0\)，且 \(f'(x)\ge0\)，所以 \(f\) 始终为正，竖条面积为 \(f(x)\,\mathrm{d}x\)。对对数项分部积分，面积与一阶矩分别为
\[
\begin{aligned}
S&=\int_1^{\mathrm{e}} f(x)\,\mathrm{d}x
=\left[\frac{x^3}{12}-\frac{x\ln x-x}{2}\right]_1^{\mathrm{e}}
=\frac{\mathrm{e}^3-7}{12},\\
M_y&=\int_1^{\mathrm{e}} xf(x)\,\mathrm{d}x
=\left[\frac{x^4}{16}-\frac{x^2}{4}\ln x+\frac{x^2}{8}\right]_1^{\mathrm{e}}
=\frac{\mathrm{e}^4-2\mathrm{e}^2-3}{16}.
\end{aligned}
\]
因此 \(\overline{x}=\frac{M_y}{S}=\boxed{\frac{3(\mathrm{e}^4-2\mathrm{e}^2-3)}{4(\mathrm{e}^3-7)}}\)。
\textbf{结论与检查}\quad \(\boxed{s=\frac{\mathrm{e}^2+1}{4}}\)，形心横坐标如上，约为 \(2.1103\)，在 \((1,\mathrm{e})\) 内。形心使用的是区域的面积权重，不含弧长权重。
\end{QuestionSolution}

% Q 2014 20
\QuestionPage{3.4 定积分的应用}{解答题}{2014 年第 20 题}{3}{递推函数围成的面积}{180 / 234}{***}
设函数$f(x) = \frac{x}{1 + x}$，$x \in [0,1]$．定义函数列
\[
f_1(x) = f(x),\quad f_2(x) = f(f_1(x)),\quad \cdots,\quad f_n(x) = f(f_{n - 1}(x)), \cdots
\]
设$S_n$是曲线$y = f_n(x)$，直线$x = 1$, $y = 0$所围平面图形的面积．求极限$\lim_{n \to \infty} nS_n$．
\QuestionEnd

\begin{QuestionSolution}{2014 年第 20 题}
\textbf{已知与目标}\quad \(f_n\) 是 \(f(x)=\frac{x}{1+x}\) 反复复合 \(n\) 次的函数，\(S_n\) 是其图像下方的面积，要求 \(nS_n\) 的极限。

\textbf{方法与条件}\quad 对 \(x>0\) 先取倒数，把复合递推化成加法递推；\(x=0\) 单独代入。得到明确函数式后直接积分，避免未经证明地交换极限和积分。

\textbf{推导}\quad 对 \(x>0\)，\(\frac{1}{f_n(x)}=1+\frac{1}{f_{n-1}(x)}\)。重复展开共加上 \(n\) 个 \(1\)，得 \(\frac{1}{f_n(x)}=n+\frac{1}{x}\)，所以 \(f_n(x)=\frac{x}{1+nx}\)。在零点原递推始终给出 \(0\)，也符合该式。因 \(f_n\ge0\)，面积为其积分，于是
\[
nS_n=\int_0^1\frac{nx}{1+nx}\,\mathrm{d}x
=1-\int_0^1\frac{\mathrm{d}x}{1+nx}
=1-\frac{\ln(1+n)}{n}.
\]
这里把 \(\frac{nx}{1+nx}\) 写为 \(1-\frac{1}{1+nx}\)，而 \(\int_0^1\frac{1}{1+nx}\,\mathrm{d}x=\frac{\ln(1+n)}{n}\to0\)。

\textbf{结论与检查}\quad \(\boxed{\lim_{n\to\infty}nS_n=1}\)。对每个 \(n\)，\(0<nS_n<1\)，所得极限与该范围一致。
\end{QuestionSolution}

% Q 2019 19
\QuestionPage{3.4 定积分的应用}{解答题}{2019 年第 19 题}{3}{分段面积与极限}{183 / 234}{***}
设$n$是正整数，记$S_n$为曲线$y = \mathrm{e}^{- x}\sin x$（$0 \le x \le n\pi$）
与$x$轴所围成图形的面积，求$S_n$，并求$\lim_{n \to \infty} S_n$．
\QuestionEnd

\begin{QuestionSolution}{2019 年第 19 题}
\textbf{已知与目标}\quad 求 \(0\le x\le n\pi\) 内曲线与横轴之间的总面积，再求 \(n\to\infty\) 的极限。面积必须将横轴下方的部分也按正值计算。

\textbf{方法与条件}\quad 用 \(S_n=\int_0^{n\pi}\mathrm{e}^{-x}|\sin x|\,\mathrm{d}x\)，按零点 \(k\pi\) 分段。平移后正弦形状相同，指数因子使各段面积构成等比数列，即后一项是前一项的固定倍数。

\textbf{推导}\quad 令 \(q=\mathrm{e}^{-\pi}\in(0,1)\)。在第 \(k+1\) 段令 \(x=k\pi+t\)，则 \(\mathrm{e}^{-x}=q^k\mathrm{e}^{-t}\)，\(|\sin(k\pi+t)|=\sin t\)。因此该段面积为 \(q^kA\)，其中
\[
A=\int_0^\pi \mathrm{e}^{-t}\sin t\,\mathrm{d}t
=\left[-\frac12 \mathrm{e}^{-t}(\sin t+\cos t)\right]_0^\pi
=\frac{1+q}{2}.
\]
故
\[
S_n=A\sum_{k=0}^{n-1}q^k
=\boxed{\frac{(1+q)(1-q^n)}{2(1-q)}},
\qquad
\lim_{n\to\infty}S_n=\boxed{\frac{1+q}{2(1-q)}}.
\]
等比和公式由 \((1-q)\sum_{k=0}^{n-1}q^k=1-q^n\) 得到，且 \(0<q<1\) 保证 \(q^n\to0\)。

\textbf{结论与检查}\quad \(\boxed{S_n=\frac{(1+\mathrm{e}^{-\pi})(1-\mathrm{e}^{-n\pi})}{2(1-\mathrm{e}^{-\pi})}}\)，极限为 \(\boxed{\frac{1+\mathrm{e}^{-\pi}}{2(1-\mathrm{e}^{-\pi})}}\)。取 \(n=1\) 恢复首段面积 \(A\)，且 \(S_n\) 随 \(n\) 增加，符合总面积性质。
\end{QuestionSolution}

% Q 2021 19
\QuestionPage{3.4 定积分的应用}{解答题}{2021 年第 19 题}{3}{弧长与旋转曲面面积}{185 / 234}{***}
\(f(x)\) 满足
\[\int \frac{f(x)}{\sqrt{x}} \,\mathrm{d}x=\frac{1}{6}x^{2}-x+C,
\]
\(L\) 为曲线 \(y=f(x)\)，其中 \(4 \le x \le 9\)。\(L\) 的弧长为 \(s\)，\(L\) 绕 \(x\) 轴旋转一周所形成的曲面的面积为 \(A\)，求 \(s\) 和 \(A\)。
\QuestionEnd

\begin{QuestionSolution}{2021 年第 19 题}
\textbf{已知与目标}\quad 已知 \(\frac{f(x)}{\sqrt{x}}\) 的一个原函数，求曲线在 \([4,9]\) 上的长度和绕 \(x\) 轴形成的曲面面积。

\textbf{方法与条件}\quad 对不定积分等式求导恢复被积函数。弧长微元为 \(\mathrm{d}s=\sqrt{1+(f')^2}\,\mathrm{d}x\)，旋转曲面面积微元为 \(2\pi|f(x)|\,\mathrm{d}s\)，其中绝对值表示到轴的距离。

\textbf{推导}\quad 求导得 \(\frac{f(x)}{\sqrt{x}}=\frac{x}{3}-1\)，故 \(f(x)=\frac{1}{3}x^{\frac{3}{2}}-x^{\frac{1}{2}}\)。在 \([4,9]\) 上为正，并且
\[
f'(x)=\frac{x-1}{2\sqrt{x}},\qquad
1+(f')^2=\frac{4x+(x-1)^2}{4x}=\frac{(x+1)^2}{4x}.
\]
因 \(x>0\)，根式等于 \(\frac{x+1}{2\sqrt{x}}\)。因此弧长为
\[
s=\int_4^9\frac{x+1}{2\sqrt{x}}\,\mathrm{d}x
=\left[\frac{x^{\frac{3}{2}}}3+\sqrt{x}\right]_4^9
=\boxed{\frac{22}{3}}.
\]
以 \(2\pi f(x)\) 乘弧长微元，得曲面面积
\[
\begin{aligned}
A&=\pi\int_4^9(\frac{x}{3}-1)(x+1)\,\mathrm{d}x\\
&=\pi\left[\frac{x^3}{9}-\frac{x^2}{3}-x\right]_4^9
=\boxed{\frac{425\pi}{9}}.
\end{aligned}
\]
代入弧长端点为 \((9+3)-(\frac{8}{3}+2)=\frac{22}{3}\)；面积原函数在 \(9\) 处为 \(45\)，在 \(4\) 处为 \(-\frac{20}{9}\)，相减得 \(\frac{425}{9}\)。

\textbf{结论与检查}\quad \(\boxed{s=\frac{22}{3}}\)，\(\boxed{A=\frac{425\pi}{9}}\)。\(s>9-4=5\)，且曲面面积为正，符合长度与面积的基本要求。
\end{QuestionSolution}

% Q 2023 19
\QuestionPage{3.4 定积分的应用}{解答题}{2023 年第 19 题}{3}{无界区域的面积与旋转体积}{186 / 234}{***}
已知平面区域
\[D=\left\{(x,y)\mid 0\le y\le \frac{1}{x\sqrt{1+x^2}},\ x\ge 1\right\}.
\]
\begin{enumerate}
\item 求 \(D\) 的面积；
\item 求 \(D\) 绕 \(x\) 轴旋转所成旋转体的体积。
\end{enumerate}
\QuestionEnd

\begin{QuestionSolution}{2023 年第 19 题}
\textbf{已知与目标}\quad 区域无界且位于非负函数 \(y=\frac{1}{x\sqrt{1+x^2}}\) 下方，\(x\ge1\)。分别求面积和绕 \(x\) 轴的体积，需验证两个反常积分有限。

\textbf{方法与条件}\quad 面积为高度积分，体积为圆盘截面积积分。先取上限 \(R\)，计算后令 \(R\to+\infty\)。面积用倒数换元，体积把平方后的分母拆开。

\textbf{推导}\quad （1）令 \(u=\frac{1}{x}>0\)，\(\mathrm{d}x=-\frac{1}{u^2}\,\mathrm{d}u\)，且 \(x\sqrt{1+x^2}=\frac{\sqrt{1+u^2}}{u^2}\)。上下限从 \(1,+\infty\) 变为 \(1,0\)，负号反转后，
\[
S=\int_1^\infty\frac{\mathrm{d}x}{x\sqrt{1+x^2}}
=\int_0^1\frac{\mathrm{d}u}{\sqrt{1+u^2}}
=\left[\ln(u+\sqrt{1+u^2})\right]_0^1
=\boxed{\ln(1+\sqrt2)}.
\]
（2）圆盘半径为原函数值，平方后 \(\frac{1}{x^2(1+x^2)}=\frac{1}{x^2}-\frac{1}{1+x^2}\)，故
\[
V=\pi\int_1^\infty\left(\frac1{x^2}-\frac1{1+x^2}\right)\mathrm{d}x
=\pi\left[-\frac1x-\arctan x\right]_1^\infty
=\boxed{\pi\left(1-\frac\pi4\right)} .
\]
上端原函数极限为 \(-\frac{\pi}{2}\)，下端为 \(-1-\frac{\pi}{4}\)，相减为 \(1-\frac{\pi}{4}\)。

\textbf{结论与检查}\quad \(\boxed{S=\ln(1+\sqrt2)}\)，\(\boxed{V=\pi(1-\frac{\pi}{4})}\)。两个极限都有限，说明无界区域仍可有有限面积与体积；两结果均为正。
\end{QuestionSolution}

% Q 2024 19
\QuestionPage{3.4 定积分的应用}{解答题}{2024 年第 19 题}{3}{旋转体积及其最大值}{187 / 234}{***}
设 \(t>0\)，曲线
\[y=\sqrt{x}\mathrm{e}^{-x}
\]
与直线 \(x=t,\ x=2t\) 及 \(x\) 轴所围平面图形绕 \(x\) 轴旋转所得的旋转体体积为 \(V(t)\)，求 \(V(t)\) 的最大值。
\QuestionEnd

\begin{QuestionSolution}{2024 年第 19 题}
\textbf{已知与目标}\quad \(t>0\)，积分区间从 \(t\) 到 \(2t\) 随参数变化，要求旋转体积 \(V(t)\) 在整个正半轴上的最大值。

\textbf{方法与条件}\quad 圆盘法以函数值平方为截面积。对上下限都变化的积分求导时，上限项乘上限导数，下限项乘下限导数并取负号。通过全定义域上的导数符号确认全局最大值。

\textbf{推导}\quad \(y^2=x\mathrm{e}^{-2x}\)，故 \(V(t)=\pi\int_t^{2t}x\mathrm{e}^{-2x}\,\mathrm{d}x\)。被积函数连续，求导得
\[
V'(t)=\pi(4t\mathrm{e}^{-4t}-t\mathrm{e}^{-2t})
=\pi t\mathrm{e}^{-4t}(4-\mathrm{e}^{2t}).
\]
故 \(V\) 在 \(0<t<\ln2\) 上递增，在 \(t>\ln2\) 上递减，最大值在 \(t=\ln2\) 取得。分部积分得
\[
V(t)=\frac{\pi}{4}\bigl[(2t+1)\mathrm{e}^{-2t}-(4t+1)\mathrm{e}^{-4t}\bigr].
\]
在 \(t=\ln2\) 时，\(\mathrm{e}^{-2t}=\frac{1}{4}\)、\(\mathrm{e}^{-4t}=\frac{1}{16}\)，所以
\[
V(\ln2)=\frac{\pi}{64}\bigl[4(2\ln2+1)-(4\ln2+1)\bigr]
=\frac{\pi(4\ln2+3)}{64}.
\]
\textbf{结论与检查}\quad \(\boxed{V_{\max}=\frac{\pi(4\ln2+3)}{64}}\)，仅在 \(t=\ln2\) 取得。\(V(t)\) 在 \(t\to0^+\) 和 \(t\to+\infty\) 时都趋零，与先增后减的结论一致。
\end{QuestionSolution}

% Q 2011 20
\QuestionPage{3.4 定积分的应用}{解答题}{2011 年第 20 题}{3}{旋转体容积与抽水做功}{176 / 234}{****}
一容器的内侧是由图中曲线绕$y$轴旋转一周而成的曲面，
该曲线由$x^2 + y^2 = 2y$（$y \ge \frac{1}{2}$）与$x^2 + y^2 = 1$（$y \le \frac{1}{2}$）连接而成．
\par\centerline{\begin{tikzpicture}[scale=\DiagramScale,x=1pt,y=1pt,xscale=0.439724,yscale=-0.439724]
\path[use as bounding box] (0,0) rectangle (440,381);
\path[fill=black!100,even odd rule,shift={(178,6)}] svg {M0 0 C2.7 1.3 2.8 2.6 3.9 5.4 C4.3 6.3 4.7 7.2 5 8.1 C5.2 8.6 5.2 8.6 6.1 10.9 C6.5 11.9 6.9 12.8 7.3 13.8 C10 20.7 10 20.7 10 23 C5.3 21.1 5.3 21.1 3 20 C2.7 33.2 2.3 46.4 2 60 C4.3 60.2 4.3 60.2 16 61 C31.4 64.2 46.4 71.4 58 82 C58.8 82.7 59.7 83.5 60.5 84.3 C67.1 90.4 71.6 97.2 76 105 C76.4 105.7 76.8 106.3 77.2 107 C79.8 111.6 81.2 116.5 82.7 121.5 C82.9 122.2 83.1 122.9 83.3 123.7 C89 143.3 86.3 163.3 78.3 181.8 C75.9 187.7 75.8 190.4 78.3 196.2 C78.5 196.6 78.5 196.6 79.3 198.6 C83.6 209 85 217.6 85 229 C128.9 229 172.8 229 218 229 C217 226.7 216 224.4 215 222 C220.4 222.6 224.9 224.4 229.9 226.6 C230.7 226.9 231.4 227.2 232.2 227.6 C234.2 228.4 236.1 229.2 238 230 C238 230.7 238 231.3 238 232 C234.7 233.4 231.5 234.7 228.2 236.1 C227.3 236.5 226.3 236.8 225.4 237.2 C224.5 237.6 223.6 238 222.7 238.3 C221.8 238.7 221 239 220.2 239.4 C218 240 218 240 215 239 C216 237 217 235 218 233 C174.4 233 130.9 233 86 233 C85 238.9 84 244.9 83 251 C78.5 265.5 72.2 277.7 62 289 C61.4 289.7 60.8 290.3 60.2 291 C45.2 306.9 24.2 313.8 3 316 C1.8 319.8 1.8 323.5 1.7 327.4 C1.7 328.1 1.6 328.8 1.6 329.5 C1.6 331.8 1.5 334 1.4 336.2 C1.4 337.8 1.4 339.3 1.3 340.8 C1.2 344.5 1.1 348.3 1 352 C0.3 352 -0.3 352 -1 352 C-1.3 340.1 -1.7 328.2 -2 316 C-7.3 315.3 -12.6 314.7 -18 314 C-32.5 310.7 -44.7 304.6 -56 295 C-57 294.2 -58 293.5 -59.1 292.7 C-75.1 278 -83.2 258.2 -86 237 C-86 235.3 -86.1 233.7 -86 232 C-88 232 -90 232 -92 232 C-92 238.9 -92 245.9 -92 253 C-91.3 253 -90.7 253 -90 253 C-90 253.7 -90 254.3 -90 255 C-92.3 255 -94.6 255 -97 255 C-97 254.3 -97 253.7 -97 253 C-96.3 253 -95.7 253 -95 253 C-95 247.4 -95 241.8 -95 236 C-96 235.7 -97 235.3 -98 235 C-97 234.7 -96 234.3 -95 234 C-95 233.3 -95 232.7 -95 232 C-116.5 232 -137.9 232 -160 232 C-159.7 231 -159.3 230 -159 229 C-134.9 229 -110.8 229 -86 229 C-85.3 222.7 -84.7 216.5 -84 210 C-82.2 204.3 -80.4 198.9 -77.8 193.5 C-75.4 187.8 -76.3 184 -78.6 178.4 C-78.9 177.8 -79.2 177.1 -79.5 176.4 C-87.7 156.6 -87 133.9 -79.2 114.1 C-76.7 108.1 -73.6 102.4 -70 97 C-69.6 96.3 -69.1 95.7 -68.7 95 C-66.1 91.3 -63.3 88.1 -60 85 C-59.7 84.7 -59.7 84.7 -58 83 C-41.8 67.3 -23.3 62.5 -1 60 C-1.3 46.8 -1.7 33.6 -2 20 C-3.7 21 -5.3 22 -7 23 C-7.7 23 -8.3 23 -9 23 C-6.9 14.9 -3.6 7.5 0 0 Z M3 63 C1 65 1.8 68.7 1.8 71.4 C1.8 72.2 1.8 73 1.8 73.8 C1.8 76.3 1.8 78.9 1.8 81.5 C1.8 83.3 1.8 85.1 1.8 86.8 C1.8 90.6 1.8 94.3 1.8 98.1 C1.8 102.9 1.7 107.7 1.7 112.5 C1.7 116.2 1.7 119.9 1.7 123.6 C1.7 125.4 1.7 127.1 1.7 128.9 C1.6 131.4 1.6 133.9 1.6 136.3 C1.6 137.1 1.6 137.8 1.6 138.6 C1.7 143.6 1.7 143.6 3.5 146 C4 146.3 4.5 146.7 5 147 C4 147 3 147 2 147 C2 152.5 1.9 158 1.9 163.5 C1.9 165.4 1.9 167.2 1.8 169.1 C1.8 171.8 1.8 174.5 1.8 177.2 C1.8 178 1.8 178.9 1.8 179.7 C1.8 185.7 1.8 185.7 4 189 C3.3 189 2.7 189 2 189 C1.8 193.7 1.7 198.5 1.5 203.2 C1.5 204.8 1.4 206.4 1.4 208 C1.3 210.3 1.2 212.7 1.1 215 C1.1 215.7 1.1 216.4 1 217.1 C0.9 221.3 1.3 224.8 2 229 C6.9 229 11.9 229 17 229 C17.3 228.3 17.7 227.7 18 227 C18.7 227 19.3 227 20 227 C20.2 227.3 20.2 227.3 21 229 C28.8 229.2 36.6 229.3 44.4 229.4 C48 229.5 51.7 229.6 55.3 229.6 C58.8 229.7 62.3 229.8 65.8 229.8 C67.8 229.9 69.8 229.9 71.7 230 C72.9 230 74.1 230 75.4 230 C76.4 230 77.5 230.1 78.6 230.1 C81 230 81 230 82 229 C83.4 215.9 79.5 203.3 73 192 C72 192.7 71 193.3 70 194 C70.2 192.9 70.4 191.8 70.6 190.6 C71 187 71 187 70 185 C70.6 184.9 71.2 184.8 71.9 184.7 C74 184 74 184 76 180.8 C84.8 161.5 86.3 138 79 118 C70.5 96.8 55.6 79.2 34.5 70 C24.2 65.8 14.1 63.3 3 63 Z M-50 80 C-50.7 80.5 -51.4 81 -52.1 81.6 C-53.8 82.9 -53.8 82.9 -56 85 C-56 85.7 -56 86.3 -56 87 C-56.6 87.2 -57.2 87.5 -57.8 87.7 C-70.6 95.1 -76.5 114.3 -80.8 127.5 C-85.1 144.4 -82.7 167 -74 182 C-73 182.3 -72 182.7 -71 183 C-70.8 182.7 -70.8 182.7 -70 181 C-69.3 181.3 -68.7 181.7 -68 182 C-68.7 183.3 -69.3 184.6 -70 186 C-69 186.7 -68 187.3 -67 188 C-67.3 188.2 -67.3 188.2 -69 189 C-69 190 -69 191 -69 192 C-70.3 191.7 -71.6 191.3 -73 191 C-77.1 197.9 -79.7 204.5 -81.4 212.2 C-81.5 213 -81.7 213.7 -81.9 214.5 C-82.9 219.6 -83 223.5 -83 229 C-62.5 229 -42.1 229 -21 229 C-21 228 -21 227 -21 226 C-20.8 226.2 -20.8 226.2 -20 227.5 C-16.9 229.8 -13.6 229.4 -9.9 229.2 C-9.5 229.2 -9.5 229.2 -7.6 229.2 C-2.2 229.2 -2.2 229.2 -1 228 C-0.9 225.5 -0.9 223 -0.9 220.5 C-0.9 219.7 -0.9 218.9 -0.9 218.1 C-0.9 215.4 -0.9 212.8 -0.9 210.2 C-0.9 208.4 -0.9 206.6 -0.9 204.7 C-0.9 200 -0.9 195.2 -0.9 190.4 C-0.9 185.5 -0.9 180.6 -1 175.7 C-1 166.2 -1 156.6 -1 147 C-2.3 147 -3.6 147 -5 147 C-4.3 146.3 -3.7 145.7 -3 145 C-2.3 145 -1.7 145 -1 145 C-1 117.9 -1 90.9 -1 63 C-18.8 63 -35.9 69.3 -50 80 Z M-71 187 C-71.3 187.7 -71.7 188.3 -72 189 C-71.3 189 -70.7 189 -70 189 C-70.3 188.3 -70.7 187.7 -71 187 Z M-83 233 C-82.9 251 -75.8 267.8 -65 282 C-64.5 282.6 -64 283.3 -63.5 284 C-52.8 297.3 -36.7 306.5 -20.2 310.5 C-19.9 310.6 -19.9 310.6 -18.2 311 C-14.2 311.8 -10.4 312.1 -6.2 312 C-5.7 312 -5.7 312 -3.3 312 C-2.5 312 -1.8 312 -1 312 C-1 301.7 -1 291.4 -0.9 281.2 C-0.9 276.4 -0.9 271.6 -0.9 266.8 C-0.9 262.2 -0.9 257.6 -0.9 253 C-0.9 251.3 -0.9 249.5 -0.9 247.7 C-0.9 245.3 -0.9 242.8 -0.9 240.4 C-0.9 240 -0.9 240 -0.9 238.1 C-0.9 233.1 -0.9 233.1 -2 232 C-4.5 231.9 -7 231.9 -9.4 231.9 C-10.2 231.9 -11 231.9 -11.8 231.9 C-14.4 231.9 -17 231.9 -19.6 231.9 C-21.4 231.9 -23.1 231.9 -24.9 231.9 C-28.7 231.9 -32.4 231.9 -36.2 231.9 C-41 231.9 -45.9 231.9 -50.7 231.9 C-54.4 231.9 -58.1 231.9 -61.8 231.9 C-63.5 231.9 -65.3 231.9 -67.1 231.9 C-69.6 231.9 -72.1 231.9 -74.6 231.9 C-75.3 231.9 -76 231.9 -76.8 231.9 C-81.9 231.9 -81.9 231.9 -83 233 Z M2 232 C0.9 282.7 0.9 282.7 0.9 299.6 C0.9 300.6 0.9 301.7 0.9 302.8 C0.9 303.8 0.9 304.7 0.9 305.7 C0.9 306.5 0.9 307.3 0.9 308.1 C1 310 1 310 2 312 C21.3 313.7 40.5 304 55 292 C69.6 279.5 80.6 259.2 82.3 239.9 C82.3 237.5 82.3 235.3 82 233 C80.3 231.3 78.1 231.9 75.9 231.9 C75.5 231.9 75.5 231.9 73.6 231.9 C73.2 231.9 73.2 231.9 71.3 231.9 C68.7 231.9 66.2 231.9 63.6 231.9 C61.9 231.9 60.1 231.9 58.3 231.9 C53.6 231.9 49 231.9 44.3 231.9 C39.6 231.9 34.8 231.9 30 232 C20.7 232 11.3 232 2 232 Z };
\path[fill=black!100,even odd rule,shift={(308,70)}] svg {M0 0 C0.8 0.2 0.8 0.2 5 1 C4.2 7.6 4.2 7.6 2 11 C3.3 10.7 4.6 10.3 6 10 C5.7 11 5.3 12 5 13 C2.4 13 -0.3 13 -3 13 C-3.2 13.7 -3.4 14.4 -3.6 15.1 C-5.4 18.9 -8 22 -11 25 C-11.3 25 -11.7 25 -12 25 C-12.3 20.4 -12.7 15.8 -13 11 C-14 10.7 -15 10.3 -16 10 C-14.4 9.7 -12.7 9.3 -11 9 C-9.4 12.2 -9.9 15.4 -10 19 C-7.4 16.8 -7.4 16.8 -5 14 C-5.2 11.7 -5.2 11.7 -6 10 C-5 9.7 -4 9.3 -3 9 C-3 9.7 -3 10.3 -3 11 C-2 11 -1 11 0 11 C1 8 2 5.1 3 2 C1.4 2.3 -0.3 2.7 -2 3 C-1.3 2 -0.7 1 0 0 Z };
\path[fill=black!100,even odd rule,shift={(341.5625,71.1875)}] svg {M0 0 C2.4 0.8 2.4 0.8 4.4 3.8 C5 7.1 4.9 9.1 3.1 11.9 C2.1 13.3 1 14.7 -0.2 16.1 C-1.6 17.8 -1.6 17.8 -2.6 19.8 C-1.2 19.6 -1.2 19.6 5.4 18.8 C4.8 20.1 4.1 21.5 3.4 22.8 C-0.2 22.8 -3.8 22.8 -7.6 22.8 C-6.6 20.8 -5.6 18.8 -4.6 16.8 C-3.9 16.8 -3.2 16.8 -2.6 16.8 C-0.2 13 -0 9.2 0.4 4.8 C-0.6 4.6 -0.6 4.6 -5.6 3.8 C-5.6 4.5 -5.6 5.1 -5.6 5.8 C-6.2 5.8 -6.9 5.8 -7.6 5.8 C-5.9 2.1 -4.5 -0.3 0 0 Z };
\path[fill=black!100,even odd rule,shift={(190.0625,41.1875)}] svg {M0 0 C2.9 0.8 2.9 0.8 4.9 3.8 C4.9 10.4 2 14 -2.1 18.8 C-0.7 18.6 -0.7 18.6 5.9 17.8 C5.6 19.1 5.3 20.5 4.9 21.8 C0.6 21.8 -3.6 21.8 -8.1 21.8 C-7.1 20.8 -6.2 19.7 -5.2 18.6 C-0.4 12.9 -0.4 12.9 0.9 9.8 C0.6 6.5 0.6 6.5 -0.1 3.8 C-4.8 4.7 -4.8 4.7 -7.1 5.8 C-6.8 3.9 -6.8 3.9 -6.1 1.8 C-3.1 -0.2 -3.1 -0.2 0 0 Z };
\path[fill=black!100,even odd rule,shift={(163,240)}] svg {M0 0 C3.5 -0.3 5.4 -0.4 8.4 1.5 C10.7 5 10.8 8 10 12 C8.1 16.6 6.3 19.4 1.7 21.5 C1.2 21.6 1.2 21.6 -1 22 C-0.2 21 0.6 20.1 1.4 19.1 C4.3 15 5.6 10.9 6 6 C5.7 5 5.3 4 5 3 C4 2.8 4 2.8 -1 2 C-0.7 1.3 -0.3 0.7 0 0 Z };
\path[fill=black!100,even odd rule,shift={(162,241)}] svg {M0 0 C-1.1 2.2 -2.1 4.1 -3.6 6.1 C-5.6 10.3 -5.7 14.4 -6 19 C-4 19.3 -2 19.7 0 20 C0 20.7 0 21.3 0 22 C-6.4 21.5 -6.4 21.5 -8.9 19.4 C-10.5 15.9 -10.7 12.7 -10 9 C-7.9 5 -5 0 0 0 Z };
\path[fill=black!100,even odd rule,shift={(379.625,195.125)}] svg {M0 0 C2.4 0.9 2.4 0.9 4.4 2.9 C5 6.2 4.8 8.2 3 11 C1.5 13 -0 15 -1.6 16.9 C0.7 16.9 3 16.9 5.4 16.9 C5 17.9 4.7 18.9 4.4 19.9 C0.7 19.9 -2.9 19.9 -6.6 19.9 C-6.2 19.3 -6.2 19.3 -3.8 16.2 C-1.2 12.9 0.5 10 1.4 5.9 C1 5.2 0.7 4.6 0.4 3.9 C-0.9 3.9 -2.3 3.9 -3.6 3.9 C-4.3 4.5 -4.9 5.2 -5.6 5.9 C-6 5.2 -6.3 4.6 -6.6 3.9 C-3.3 -0.2 -3.3 -0.2 0 0 Z };
\path[fill=black!100,even odd rule,shift={(351,79)}] svg {M0 0 C2.1 0.4 2.1 0.4 4 1 C4.2 3.1 4.4 5.1 4.6 7.2 C4.7 7.8 4.9 8.4 5 9 C5.7 9.3 6.3 9.7 7 10 C9 3.4 9 3.4 9 0 C10 0.3 11 0.7 12 1 C11.5 5.4 9.8 8 7.2 11.6 C6.5 12.5 5.8 13.4 5.1 14.4 C3.2 16.8 1.4 19 -1 21 C-2.3 21 -3.6 21 -5 21 C-4.7 20 -4.3 19 -4 18 C-2 17.4 -2 17.4 0 17 C1 16 1 16 1.1 12.8 C1.1 11.6 1.1 10.3 1.1 8.9 C1.1 7.6 1 6.4 1 5 C1 4 1 3 1 2 C0.3 1.7 -0.3 1.3 -1 1 C-0.7 0.7 -0.3 0.3 0 0 Z };
\path[fill=black!100,even odd rule,shift={(201,2)}] svg {M0 0 C0.7 3.3 1.3 6.6 2 10 C4.1 7.7 4.1 7.7 6 5 C5.8 2.7 5.8 2.7 5 1 C6 0.7 7 0.3 8 0 C9 3 9 3 8.2 5 C5.2 10.4 1.9 16 -3 20 C-5.7 20.7 -5.7 20.7 -8 21 C-8 20 -8 19 -8 18 C-7.2 17.9 -6.4 17.8 -5.6 17.6 C-4.7 17.4 -3.9 17.2 -3 17 C-0.9 12.8 -1.7 7.4 -3 3 C-3.7 2.3 -4.3 1.7 -5 1 C-3 0 -3 0 0 0 Z };
\path[fill=black!100,even odd rule,shift={(252,79)}] svg {M0 0 C4 0 4 0 5.8 1.5 C6 1.7 6 1.7 7 3 C7 2.3 7 1.7 7 1 C8.6 0.7 10.3 0.3 12 0 C11 3 11 3 8.9 4.2 C8.3 4.5 7.7 4.7 7 5 C7 7.3 7 9.6 7 12 C8 11.7 9 11.3 10 11 C9.3 12.3 8.7 13.6 8 15 C5 14 5 14 3 11 C2.8 11.3 2.8 11.3 1.7 13 C1.1 13.7 0.6 14.3 0 15 C-1 15 -2 15 -3 15 C-3 14 -3 13 -3 12 C-1.7 12 -0.4 12 1 12 C1.2 10.7 1.3 9.4 1.5 8.1 C1.6 7.4 1.7 6.7 1.8 5.9 C2 4 2 4 2 2 C1.3 2 0.7 2 0 2 C0 1.3 0 0.7 0 0 Z };
\path[fill=black!100,even odd rule,shift={(212,324)}] svg {M0 0 C1.7 0.3 3.3 0.7 5 1 C5 1.7 5 2.3 5 3 C5.8 2.5 6.6 2 7.4 1.5 C10 0 10 0 12 0 C12 1 12 2 12 3 C10.7 3.3 9.4 3.7 8 4 C8 6.6 8 9.3 8 12 C8.7 12 9.3 12 10 12 C10 12.7 10 13.3 10 14 C8.7 14.3 7.4 14.7 6 15 C5.7 13.7 5.3 12.4 5 11 C4.2 11.7 3.5 12.4 2.7 13.1 C1.8 13.7 0.9 14.3 0 15 C-1 14.7 -2 14.3 -3 14 C-3 13.3 -3 12.7 -3 12 C-2.2 11.8 -1.4 11.5 -0.6 11.3 C2 10 2 10 3 7.6 C3 5 3 5 1.6 3.3 C1 2.9 0.5 2.4 0 2 C0 1.3 0 0.7 0 0 Z };
\path[fill=black!100,even odd rule,shift={(170,323)}] svg {M0 0 C0 0.6 0 1.2 0.1 1.8 C0.2 4.5 0.3 7.2 0.4 9.9 C0.4 10.8 0.4 11.7 0.5 12.7 C0.5 13.6 0.5 14.5 0.6 15.4 C0.6 16.3 0.6 17.1 0.7 17.9 C1 20 1 20 3 22 C0.4 22 -2.3 22 -5 22 C-4.7 21.3 -4.3 20.7 -4 20 C-3.8 17.1 -3.6 14.3 -3.4 11.4 C-3.4 10.6 -3.4 9.8 -3.3 8.9 C-3.2 7 -3.1 5 -3 3 C-3.7 2.7 -4.3 2.3 -5 2 C-2 0 -2 0 0 0 Z };
\path[fill=black!100,even odd rule,shift={(191,124)}] svg {M0 0 C0 0.6 0.1 1.2 0.1 1.8 C0.2 4.5 0.3 7.2 0.4 9.9 C0.5 10.8 0.5 11.7 0.6 12.7 C0.6 13.6 0.6 14.5 0.7 15.4 C0.7 16.3 0.8 17.1 0.8 17.9 C1 20 1 20 2 22 C-0.3 22 -2.6 22 -5 22 C-4.7 21.3 -4.3 20.7 -4 20 C-3.8 17.1 -3.6 14.3 -3.4 11.4 C-3.4 10.6 -3.4 9.8 -3.3 8.9 C-3.2 7 -3.1 5 -3 3 C-4 2.7 -5 2.3 -6 2 C-2.2 0 -2.2 0 0 0 Z };
\path[fill=black!100,even odd rule,shift={(300,317)}] svg {M0 0 C1 0 2 0 3 0 C3 0.6 3.1 1.2 3.1 1.8 C3.2 4.5 3.3 7.2 3.4 9.9 C3.5 10.8 3.5 11.7 3.6 12.7 C3.6 13.6 3.6 14.5 3.7 15.4 C3.7 16.3 3.8 17.1 3.8 17.9 C4 20 4 20 5 22 C2.7 22 0.4 22 -2 22 C-1.7 21.3 -1.3 20.7 -1 20 C-0.8 17.1 -0.6 14.3 -0.4 11.4 C-0.4 10.6 -0.4 9.8 -0.3 8.9 C-0.2 7 -0.1 5 0 3 C-0.7 2.7 -1.3 2.3 -2 2 C-1.3 1.3 -0.7 0.7 0 0 Z };
\path[fill=black!100,even odd rule,shift={(272,242)}] svg {M0 0 C0.7 0.3 1.3 0.7 2 1 C2 7.6 2 14.2 2 21 C2.7 21.3 3.3 21.7 4 22 C1.7 22 -0.6 22 -3 22 C-2.7 21.3 -2.3 20.7 -2 20 C-1.9 17.3 -1.9 14.7 -1.9 12.1 C-1.9 11.7 -1.9 11.7 -2 9.9 C-2 7.9 -2 6 -2 4 C-2 3 -2 2 -2 1 C-1.3 0.7 -0.7 0.3 0 0 Z };
\path[fill=black!100,even odd rule,shift={(427,228)}] svg {M0 0 C1 0.3 2 0.7 3 1 C3 1.7 3 2.3 3 3 C5 2 7 1 9 0 C9.3 1 9.7 2 10 3 C8.7 3.3 7.4 3.7 6 4 C6 6.3 6 8.6 6 11 C7 11 8 11 9 11 C8.7 12 8.3 13 8 14 C5.9 14.7 5.9 14.7 4 15 C3.3 13.7 2.7 12.4 2 11 C1.7 11.2 1.7 11.2 0 12 C0.2 10.9 0.5 9.9 0.7 8.8 C1 5 1 5 -0.4 2.6 C-1 2 -1.5 1.5 -2 1 C-1.3 0.7 -0.7 0.3 0 0 Z };
\path[fill=black!100,even odd rule,shift={(268,315)}] svg {M0 0 C1.7 -0 3.3 -0 5 0 C5.3 0.3 5.7 0.7 6 1 C5.6 5.1 5.3 7.5 3 11 C4.3 11.3 5.6 11.7 7 12 C6 13 6 13 2.1 13.4 C-0.7 13.7 -1.7 13.8 -4 15.6 C-4.3 16 -4.7 16.5 -5 17 C-5 14.7 -5 12.4 -5 10 C-4 10 -3 10 -2 10 C0.1 9.4 0.1 9.4 2 8 C3.2 4.9 3.2 4.9 4 2 C2.4 2.3 0.7 2.7 -1 3 C-0.7 2 -0.3 1 0 0 Z };
\path[fill=black!100,even odd rule,shift={(381,166)}] svg {M0 0 C0 5.9 0 11.9 0 18 C0.7 18.3 1.3 18.7 2 19 C-0.3 19 -2.6 19 -5 19 C-4.3 18.7 -3.7 18.3 -3 18 C-3 12.7 -3 7.4 -3 2 C-3.7 1.7 -4.3 1.3 -5 1 C-3 0 -3 0 0 0 Z };
\path[fill=black!100,even odd rule,shift={(242,320)}] svg {M0 0 C0.7 0 1.3 0 2 0 C2 2.3 2 4.6 2 7 C4.3 7 6.6 7 9 7 C9 7.7 9 8.3 9 9 C6.7 9 4.4 9 2 9 C2 11.3 2 13.6 2 16 C1.3 16 0.7 16 0 16 C0 13.7 0 11.4 0 9 C-2.3 9 -4.6 9 -7 9 C-7 8.3 -7 7.7 -7 7 C-4.7 7 -2.4 7 0 7 C0 4.7 0 2.4 0 0 Z };
\path[fill=black!100,even odd rule,shift={(336,185)}] svg {M0 0 C2.1 3.2 2.5 5.3 3 9 C4 8.3 5 7.7 6 7 C4.5 11.5 2.3 14.6 -1 18 C-3.8 19.4 -3.8 19.4 -6 20 C-6 19 -6 18 -6 17 C-5.4 16.9 -4.7 16.8 -4.1 16.6 C-3.4 16.4 -2.7 16.2 -2 16 C-0.6 13.2 -0.9 11 -0.9 7.9 C-0.9 6.8 -1 5.7 -1 4.6 C-1 3.7 -1 2.9 -1 2 C-1.7 1.7 -2.3 1.3 -3 1 C-2 0.7 -1 0.3 0 0 Z };
\path[fill=black!100,even odd rule,shift={(255,324)}] svg {M0 0 C1 0.3 2 0.7 3 1 C3.7 3.1 3.7 3.1 4.2 5.6 C4.3 6.4 4.5 7.2 4.7 8.1 C4.8 8.7 4.9 9.3 5 10 C6 9.3 7 8.7 8 8 C6.4 11.8 4 14.2 1 17 C1 16 0.9 15 0.9 14 C0.8 12.7 0.7 11.4 0.6 10.1 C0.6 8.8 0.5 7.6 0.4 6.3 C0 3 0 3 -2 1 C-1.3 0.7 -0.7 0.3 0 0 Z };
\path[fill=black!100,even odd rule,shift={(282,74)}] svg {M0 0 C0.3 0 0.7 0 1 0 C1.3 2.3 1.7 4.6 2 7 C4.3 7.3 6.6 7.7 9 8 C6 10 6 10 2 10 C1.7 12.3 1.3 14.6 1 17 C0.7 17 0.3 17 0 17 C-0.3 14.7 -0.7 12.4 -1 10 C-3 10 -5 10 -7 10 C-7.3 9.3 -7.7 8.7 -8 8 C-6.8 7.8 -6.8 7.8 -1 7 C-0.7 4.7 -0.3 2.4 0 0 Z };
\path[fill=black!100,even odd rule,shift={(266,70)}] svg {M0 0 C1.6 0 3.3 0 5 0 C6.1 3.6 6.1 4.8 4.6 8.3 C4.3 8.8 4.3 8.8 3 11 C4 10.7 5 10.3 6 10 C6 11 6 12 6 13 C3.7 13 1.4 13 -1 13 C-0.7 12.5 -0.5 12 -0.2 11.4 C1.3 8.3 2.7 5.2 4 2 C2.4 2.3 0.7 2.7 -1 3 C-0.7 2 -0.3 1 0 0 Z };
\path[fill=black!100,even odd rule,shift={(370,190)}] svg {M0 0 C6.3 0 12.5 0 19 0 C19 0.7 19 1.3 19 2 C12.7 2 6.5 2 0 2 C0 1.3 0 0.7 0 0 Z };
\path[fill=black!100,even odd rule,shift={(316,84)}] svg {M0 0 C5.3 0 10.6 0 16 0 C16 0.7 16 1.3 16 2 C10.4 2 4.8 2 -1 2 C-0.7 1.3 -0.3 0.7 0 0 Z };
\path[fill=black!100,even odd rule,shift={(316,79)}] svg {M0 0 C5.3 0 10.6 0 16 0 C16 0.7 16 1.3 16 2 C10.4 2 4.8 2 -1 2 C-0.7 1.3 -0.3 0.7 0 0 Z };
\path[fill=black!100,even odd rule,shift={(277,324)}] svg {M0 0 C4.9 0 9.9 0 15 0 C15.3 0.7 15.7 1.3 16 2 C10.4 2 4.8 2 -1 2 C-0.7 1.3 -0.3 0.7 0 0 Z };
\path[fill=black!100,even odd rule,shift={(277,329)}] svg {M0 0 C4.9 0 9.9 0 15 0 C15.3 0.7 15.7 1.3 16 2 C10.4 2 4.8 2 -1 2 C-0.7 1.3 -0.3 0.7 0 0 Z };
\path[fill=black!100,even odd rule,shift={(348,190)}] svg {M0 0 C4.6 0 9.2 0 14 0 C14 0.7 14 1.3 14 2 C9.4 2 4.8 2 0 2 C0 1.3 0 0.7 0 0 Z };
\path[fill=black!100,even odd rule,shift={(146,332)}] svg {M0 0 C4.3 0 8.6 0 13 0 C13 0.7 13 1.3 13 2 C8.7 2 4.4 2 0 2 C0 1.3 0 0.7 0 0 Z };
\path[fill=black!100,even odd rule,shift={(227,315)}] svg {M0 0 C1.7 0.3 3.3 0.7 5 1 C5.8 2.8 5.8 2.8 6 5 C4.1 7.2 4.1 7.2 2 9 C2.3 7 2.7 5 3 3 C1.3 3 -0.3 3 -2 3 C-1.3 2 -0.7 1 0 0 Z };
\path[fill=black!100,even odd rule,shift={(352.53125,184.734375)}] svg {M0 0 C1 0 2 0 3 0 C3.9 0 4.9 0 5.9 0 C8.5 0.3 8.5 0.3 10.5 2.3 C5.5 2.3 0.6 2.3 -4.5 2.3 C-2.5 0.3 -2.5 0.3 0 0 Z };
\path[fill=black!100,even odd rule,shift={(228,324)}] svg {M0 0 C0.3 0.7 0.7 1.3 1 2 C2.3 1.7 3.6 1.3 5 1 C5 2 5 3 5 4 C2.4 4 -0.3 4 -3 4 C-3 3.3 -3 2.7 -3 2 C-2 1.3 -1 0.7 0 0 Z };
\path[fill=black!100,even odd rule,shift={(61,248)}] svg {M0 0 C4.6 0 9.2 0 14 0 C14 0.3 14 0.7 14 1 C11.9 1.2 11.9 1.2 1 2 C0.7 1.3 0.3 0.7 0 0 Z };
\path[fill=black!100,even odd rule,shift={(342,185)}] svg {M0 0 C0.7 0 1.3 0 2 0 C2.7 1.8 2.7 1.8 3 4 C1.6 5.8 1.6 5.8 0 7 C0 4.7 0 2.4 0 0 Z };
\path[fill=black!100,even odd rule,shift={(294,96)}] svg {M0 0 C0 1 0 2 0 3 C-1.6 3.3 -3.3 3.7 -5 4 C-4.7 3 -4.3 2 -4 1 C-1.9 0.3 -1.9 0.3 0 0 Z };
\path[fill=black!100,even odd rule,shift={(202,231)}] svg {M0 0 C1 0 2 0 3 0 C3 1 3 2 3 3 C1.7 3 0.4 3 -1 3 C-0.7 2 -0.3 1 0 0 Z };
\path[fill=black!100,even odd rule,shift={(134,222)}] svg {M0 0 C1 0.3 2 0.7 3 1 C2.7 2 2.3 3 2 4 C1 3.7 0 3.3 -1 3 C-0.7 2 -0.3 1 0 0 Z };
\path[fill=black!100,even odd rule,shift={(153,192)}] svg {M0 0 C1.3 0.3 2.6 0.7 4 1 C4 1.7 4 2.3 4 3 C2.7 3 1.4 3 0 3 C0 2 0 1 0 0 Z };
\path[fill=black!100,even odd rule,shift={(112,181)}] svg {M0 0 C1 0 2 0 3 0 C2.3 1.3 1.7 2.6 1 4 C0.3 3.7 -0.3 3.3 -1 3 C-0.7 2 -0.3 1 0 0 Z };
\path[fill=black!100,even odd rule,shift={(130,164)}] svg {M0 0 C1.3 0 2.6 0 4 0 C3.7 1 3.3 2 3 3 C2 3 1 3 0 3 C0 2 0 1 0 0 Z };
\path[fill=black!100,even odd rule,shift={(423,240)}] svg {M0 0 C1.3 0 2.6 0 4 0 C4 0.7 4 1.3 4 2 C2.7 2.3 1.4 2.7 0 3 C0 2 0 1 0 0 Z };
\path[fill=black!100,even odd rule,shift={(236,211)}] svg {M0 0 C0.7 0.3 1.3 0.7 2 1 C1.7 2 1.3 3 1 4 C0 3.7 -1 3.3 -2 3 C-1.3 2 -0.7 1 0 0 Z };
\path[fill=black!100,even odd rule,shift={(119,210)}] svg {M0 0 C1 1 2 2 3 3 C1.7 3 0.4 3 -1 3 C-0.7 2 -0.3 1 0 0 Z };
\path[fill=black!100,even odd rule,shift={(244,202)}] svg {M0 0 C1 0.3 2 0.7 3 1 C2 2 1 3 0 4 C0 2.7 0 1.4 0 0 Z };
\path[fill=black!100,even odd rule,shift={(301,192)}] svg {M0 0 C1 0.3 2 0.7 3 1 C3 1.7 3 2.3 3 3 C1.7 3 0.4 3 -1 3 C-0.7 2 -0.3 1 0 0 Z };
\path[fill=black!100,even odd rule,shift={(295,192)}] svg {M0 0 C1 0.3 2 0.7 3 1 C3 1.7 3 2.3 3 3 C1.7 3 0.4 3 -1 3 C-0.7 2 -0.3 1 0 0 Z };
\path[fill=black!100,even odd rule,shift={(263,192)}] svg {M0 0 C1 0.3 2 0.7 3 1 C3 1.7 3 2.3 3 3 C1.7 3 0.4 3 -1 3 C-0.7 2 -0.3 1 0 0 Z };
\path[fill=black!100,even odd rule,shift={(173,192)}] svg {M0 0 C1 0.3 2 0.7 3 1 C3 1.7 3 2.3 3 3 C1.7 3 0.4 3 -1 3 C-0.7 2 -0.3 1 0 0 Z };
\path[fill=black!100,even odd rule,shift={(160,192)}] svg {M0 0 C1 0 2 0 3 0 C3 1 3 2 3 3 C1.7 2.7 0.4 2.3 -1 2 C-0.7 1.3 -0.3 0.7 0 0 Z };
\path[fill=black!100,even odd rule,shift={(141,192)}] svg {M0 0 C1 0.3 2 0.7 3 1 C3 1.7 3 2.3 3 3 C1.7 3 0.4 3 -1 3 C-0.7 2 -0.3 1 0 0 Z };
\path[fill=black!100,even odd rule,shift={(128,192)}] svg {M0 0 C1 0.3 2 0.7 3 1 C3 1.7 3 2.3 3 3 C1.7 3 0.4 3 -1 3 C-0.7 2 -0.3 1 0 0 Z };
\path[fill=black!100,even odd rule,shift={(96,192)}] svg {M0 0 C1 0.3 2 0.7 3 1 C3 1.7 3 2.3 3 3 C1.7 3 0.4 3 -1 3 C-0.7 2 -0.3 1 0 0 Z };
\path[fill=black!100,even odd rule,shift={(83,192)}] svg {M0 0 C1 0.3 2 0.7 3 1 C3 1.7 3 2.3 3 3 C1.7 3 0.4 3 -1 3 C-0.7 2 -0.3 1 0 0 Z };
\path[fill=black!100,even odd rule,shift={(70,192)}] svg {M0 0 C1 0.3 2 0.7 3 1 C3 1.7 3 2.3 3 3 C1.7 3 0.4 3 -1 3 C-0.7 2 -0.3 1 0 0 Z };
\path[fill=black!100,even odd rule,shift={(51,192)}] svg {M0 0 C1 0.3 2 0.7 3 1 C3 1.7 3 2.3 3 3 C1.7 3 0.4 3 -1 3 C-0.7 2 -0.3 1 0 0 Z };
\path[fill=black!100,even odd rule,shift={(38,192)}] svg {M0 0 C1 0.3 2 0.7 3 1 C3 1.7 3 2.3 3 3 C1.7 3 0.4 3 -1 3 C-0.7 2 -0.3 1 0 0 Z };
\path[fill=black!100,even odd rule,shift={(25,192)}] svg {M0 0 C1 0.3 2 0.7 3 1 C3 1.7 3 2.3 3 3 C1.7 3 0.4 3 -1 3 C-0.7 2 -0.3 1 0 0 Z };
\path[fill=black!100,even odd rule,shift={(241,180)}] svg {M0 0 C1 0.3 2 0.7 3 1 C3 2 3 3 3 4 C2 3.7 1 3.3 0 3 C0 2 0 1 0 0 Z };
\path[fill=black!100,even odd rule,shift={(121,171)}] svg {M0 0 C0.7 0 1.3 0 2 0 C1.7 1.3 1.3 2.6 1 4 C0.3 4 -0.3 4 -1 4 C-0.7 2.7 -0.3 1.4 0 0 Z };
\path[fill=black!100,even odd rule,shift={(227,167)}] svg {M0 0 C1 0.3 2 0.7 3 1 C3 2 3 3 3 4 C2 3.7 1 3.3 0 3 C0 2 0 1 0 0 Z };
\path[fill=black!100,even odd rule,shift={(127,167)}] svg {M0 0 C0.7 0.7 1.3 1.3 2 2 C1.7 2.7 1.3 3.3 1 4 C0 3.7 -1 3.3 -2 3 C-2 2.3 -2 1.7 -2 1 C-1.3 0.7 -0.7 0.3 0 0 Z };
\path[fill=black!100,even odd rule,shift={(222,163)}] svg {M0 0 C1 0.3 2 0.7 3 1 C2.7 2 2.3 3 2 4 C1 3.7 0 3.3 -1 3 C-0.7 2 -0.3 1 0 0 Z };
\path[fill=black!100,even odd rule,shift={(137,160)}] svg {M0 0 C0.7 0 1.3 0 2 0 C1.7 1.3 1.3 2.6 1 4 C0.3 4 -0.3 4 -1 4 C-0.7 2.7 -0.3 1.4 0 0 Z };
\path[fill=black!100,even odd rule,shift={(151,155)}] svg {M0 0 C0 1 0 2 0 3 C-1 3.3 -2 3.7 -3 4 C-3.3 3 -3.7 2 -4 1 C-2 0 -2 0 0 0 Z };
\path[fill=black!100,even odd rule,shift={(186,151)}] svg {M0 0 C1 0 2 0 3 0 C3 1 3 2 3 3 C1.7 2.7 0.4 2.3 -1 2 C-0.7 1.3 -0.3 0.7 0 0 Z };
\path[fill=black!100,even odd rule,shift={(167,151)}] svg {M0 0 C1 0 2 0 3 0 C2.7 1 2.3 2 2 3 C1 3 0 3 -1 3 C-0.7 2 -0.3 1 0 0 Z };
\path[fill=black!100,even odd rule,shift={(250,342)}] svg {M0 0 C0.7 0.3 0.7 0.3 4 2 C2.7 2.3 1.4 2.7 0 3 C0 2 0 1 0 0 Z };
\path[fill=black!100,even odd rule,shift={(208,229)}] svg {M0 0 C1 0 2 0 3 0 C2.7 1 2.3 2 2 3 C1 2.7 0 2.3 -1 2 C-0.7 1.3 -0.3 0.7 0 0 Z };
\path[fill=black!100,even odd rule,shift={(215,226)}] svg {M0 0 C0.7 0.7 1.3 1.3 2 2 C1 2.7 0 3.3 -1 4 C-1.3 3 -1.7 2 -2 1 C-1.3 0.7 -0.7 0.3 0 0 Z };
\path[fill=black!100,even odd rule,shift={(140,225)}] svg {M0 0 C0.7 0.3 1.3 0.7 2 1 C2 2 2 3 2 4 C1 3.7 0 3.3 -1 3 C-0.7 2 -0.3 1 0 0 Z };
\path[fill=black!100,even odd rule,shift={(128,219)}] svg {M0 0 C1.3 0.3 2.6 0.7 4 1 C3.7 1.7 3.3 2.3 3 3 C2 2.7 1 2.3 0 2 C0 1.3 0 0.7 0 0 Z };
\path[fill=black!100,even odd rule,shift={(114,206)}] svg {M0 0 C1 0 2 0 3 0 C3 1 3 2 3 3 C2 3 1 3 0 3 C0 2 0 1 0 0 Z };
\path[fill=black!100,even odd rule,shift={(218,192)}] svg {M0 0 C1 0.3 2 0.7 3 1 C3 1.7 3 2.3 3 3 C1.7 3 0.4 3 -1 3 C-0.7 2 -0.3 1 0 0 Z };
\path[fill=black!100,even odd rule,shift={(283,192)}] svg {M0 0 C0.7 0.3 1.3 0.7 2 1 C2 1.7 2 2.3 2 3 C0.7 3 -0.6 3 -2 3 C-2 2.3 -2 1.7 -2 1 C-1.3 0.7 -0.7 0.3 0 0 Z };
\path[fill=black!100,even odd rule,shift={(231,192)}] svg {M0 0 C1 0.3 2 0.7 3 1 C2.7 1.7 2.3 2.3 2 3 C1 3 0 3 -1 3 C-0.7 2 -0.3 1 0 0 Z };
\path[fill=black!100,even odd rule,shift={(224,192)}] svg {M0 0 C1 0 2 0 3 0 C3 1 3 2 3 3 C2 3 1 3 0 3 C0 2 0 1 0 0 Z };
\path[fill=black!100,even odd rule,shift={(199,192)}] svg {M0 0 C0.7 0 1.3 0 2 0 C2.3 1 2.7 2 3 3 C1.7 3 0.4 3 -1 3 C-0.7 2 -0.3 1 0 0 Z };
\path[fill=black!100,even odd rule,shift={(192,192)}] svg {M0 0 C1 0 2 0 3 0 C3 1 3 2 3 3 C2 3 1 3 0 3 C0 2 0 1 0 0 Z };
\path[fill=black!100,even odd rule,shift={(187,192)}] svg {M0 0 C0.7 0.3 1.3 0.7 2 1 C2 1.7 2 2.3 2 3 C0.7 3 -0.6 3 -2 3 C-2 2.3 -2 1.7 -2 1 C-1.3 0.7 -0.7 0.3 0 0 Z };
\path[fill=black!100,even odd rule,shift={(147,192)}] svg {M0 0 C1 0 2 0 3 0 C3 1 3 2 3 3 C2 3 1 3 0 3 C0 2 0 1 0 0 Z };
\path[fill=black!100,even odd rule,shift={(134,192)}] svg {M0 0 C1 0 2 0 3 0 C3 1 3 2 3 3 C2 3 1 3 0 3 C0 2 0 1 0 0 Z };
\path[fill=black!100,even odd rule,shift={(115,192)}] svg {M0 0 C1 0 2 0 3 0 C3 1 3 2 3 3 C2 3 1 3 0 3 C0 2 0 1 0 0 Z };
\path[fill=black!100,even odd rule,shift={(89,192)}] svg {M0 0 C1 0 2 0 3 0 C3 1 3 2 3 3 C2 3 1 3 0 3 C0 2 0 1 0 0 Z };
\path[fill=black!100,even odd rule,shift={(76,192)}] svg {M0 0 C1 0 2 0 3 0 C3 1 3 2 3 3 C2 3 1 3 0 3 C0 2 0 1 0 0 Z };
\path[fill=black!100,even odd rule,shift={(64,192)}] svg {M0 0 C1 0.3 2 0.7 3 1 C2.7 1.7 2.3 2.3 2 3 C1 3 0 3 -1 3 C-0.7 2 -0.3 1 0 0 Z };
\path[fill=black!100,even odd rule,shift={(57,192)}] svg {M0 0 C1 0.3 2 0.7 3 1 C3 1.7 3 2.3 3 3 C1.7 2.7 0.4 2.3 -1 2 C-0.7 1.3 -0.3 0.7 0 0 Z };
\path[fill=black!100,even odd rule,shift={(44,192)}] svg {M0 0 C1 0 2 0 3 0 C3 1 3 2 3 3 C2 3 1 3 0 3 C0 2 0 1 0 0 Z };
\path[fill=black!100,even odd rule,shift={(18,192)}] svg {M0 0 C1 0 2 0 3 0 C3 1 3 2 3 3 C2 3 1 3 0 3 C0 2 0 1 0 0 Z };
\path[fill=black!100,even odd rule,shift={(246,185)}] svg {M0 0 C0.7 0.3 1.3 0.7 2 1 C2 2 2 3 2 4 C1 3.7 0 3.3 -1 3 C-0.7 2 -0.3 1 0 0 Z };
\path[fill=black!100,even odd rule,shift={(237,175)}] svg {M0 0 C1 1 2 2 3 3 C1.7 3 0.4 3 -1 3 C-0.7 2 -0.3 1 0 0 Z };
\path[fill=black!100,even odd rule,shift={(232,171)}] svg {M0 0 C1 0.3 2 0.7 3 1 C2.7 2 2.3 3 2 4 C1.3 3.7 0.7 3.3 0 3 C0 2 0 1 0 0 Z };
\path[fill=black!100,even odd rule,shift={(216,160)}] svg {M0 0 C1 0.3 2 0.7 3 1 C2.7 2 2.3 3 2 4 C1.3 3.7 0.7 3.3 0 3 C0 2 0 1 0 0 Z };
\path[fill=black!100,even odd rule,shift={(204,155)}] svg {M0 0 C0.7 0.3 0.7 0.3 4 2 C2.7 2.3 1.4 2.7 0 3 C0 2 0 1 0 0 Z };
\path[fill=black!100,even odd rule,shift={(192,152)}] svg {M0 0 C1 0 2 0 3 0 C3 1 3 2 3 3 C2 3 1 3 0 3 C0 2 0 1 0 0 Z };
\path[fill=black!100,even odd rule,shift={(160,152)}] svg {M0 0 C1 0 2 0 3 0 C3 1 3 2 3 3 C2 3 1 3 0 3 C0 2 0 1 0 0 Z };
\path[fill=black!100,even odd rule,shift={(152,230)}] svg {M0 0 C0.7 0.3 1.3 0.7 2 1 C1.7 2 1.3 3 1 4 C0.3 3.7 -0.3 3.3 -1 3 C-0.7 2 -0.3 1 0 0 Z };
\path[fill=black!100,even odd rule,shift={(145,228)}] svg {M0 0 C1 0.3 2 0.7 3 1 C3 1.7 3 2.3 3 3 C2 3 1 3 0 3 C0 2 0 1 0 0 Z };
\path[fill=black!100,even odd rule,shift={(220,223)}] svg {M0 0 C0.7 0 1.3 0 2 0 C2 1 2 2 2 3 C1 3 0 3 -1 3 C-0.7 2 -0.3 1 0 0 Z };
\path[fill=black!100,even odd rule,shift={(225,220)}] svg {M0 0 C1 0 2 0 3 0 C2.7 1 2.3 2 2 3 C1.3 3 0.7 3 0 3 C0 2 0 1 0 0 Z };
\path[fill=black!100,even odd rule,shift={(230,216)}] svg {M0 0 C1 0 2 0 3 0 C2.7 1 2.3 2 2 3 C1.3 3 0.7 3 0 3 C0 2 0 1 0 0 Z };
\path[fill=black!100,even odd rule,shift={(240,207)}] svg {M0 0 C0.7 0 1.3 0 2 0 C2 1 2 2 2 3 C1 3 0 3 -1 3 C-0.7 2 -0.3 1 0 0 Z };
\path[fill=black!100,even odd rule,shift={(110,201)}] svg {M0 0 C1 0.3 2 0.7 3 1 C3 1.7 3 2.3 3 3 C2 3 1 3 0 3 C0 2 0 1 0 0 Z };
\path[fill=black!100,even odd rule,shift={(313,193)}] svg {M0 0 C1.3 0 2.6 0 4 0 C4 0.7 4 1.3 4 2 C2.7 2 1.4 2 0 2 C0 1.3 0 0.7 0 0 Z };
\path[fill=black!100,even odd rule,shift={(307,192)}] svg {M0 0 C1 0.3 2 0.7 3 1 C3 1.7 3 2.3 3 3 C2 3 1 3 0 3 C0 2 0 1 0 0 Z };
\path[fill=black!100,even odd rule,shift={(276,192)}] svg {M0 0 C0.7 0 1.3 0 2 0 C2 1 2 2 2 3 C1 3 0 3 -1 3 C-0.7 2 -0.3 1 0 0 Z };
\path[fill=black!100,even odd rule,shift={(269,192)}] svg {M0 0 C1 0.3 2 0.7 3 1 C3 1.7 3 2.3 3 3 C2 3 1 3 0 3 C0 2 0 1 0 0 Z };
\path[fill=black!100,even odd rule,shift={(256,192)}] svg {M0 0 C1 0.3 2 0.7 3 1 C3 1.7 3 2.3 3 3 C2 3 1 3 0 3 C0 2 0 1 0 0 Z };
\path[fill=black!100,even odd rule,shift={(244,192)}] svg {M0 0 C0.7 0 1.3 0 2 0 C2 1 2 2 2 3 C1 3 0 3 -1 3 C-0.7 2 -0.3 1 0 0 Z };
\path[fill=black!100,even odd rule,shift={(237,192)}] svg {M0 0 C1 0.3 2 0.7 3 1 C3 1.7 3 2.3 3 3 C2 3 1 3 0 3 C0 2 0 1 0 0 Z };
\path[fill=black!100,even odd rule,shift={(212,192)}] svg {M0 0 C0.7 0 1.3 0 2 0 C2 1 2 2 2 3 C1 3 0 3 -1 3 C-0.7 2 -0.3 1 0 0 Z };
\path[fill=black!100,even odd rule,shift={(205,192)}] svg {M0 0 C1 0.3 2 0.7 3 1 C3 1.7 3 2.3 3 3 C2 3 1 3 0 3 C0 2 0 1 0 0 Z };
\path[fill=black!100,even odd rule,shift={(167,192)}] svg {M0 0 C0.7 0 1.3 0 2 0 C2 1 2 2 2 3 C1 3 0 3 -1 3 C-0.7 2 -0.3 1 0 0 Z };
\path[fill=black!100,even odd rule,shift={(122,192)}] svg {M0 0 C0.7 0 1.3 0 2 0 C2 1 2 2 2 3 C1 3 0 3 -1 3 C-0.7 2 -0.3 1 0 0 Z };
\path[fill=black!100,even odd rule,shift={(32,192)}] svg {M0 0 C0.7 0 1.3 0 2 0 C2 1 2 2 2 3 C1 3 0 3 -1 3 C-0.7 2 -0.3 1 0 0 Z };
\path[fill=black!100,even odd rule,shift={(116,176)}] svg {M0 0 C1 0 2 0 3 0 C2.7 1 2.3 2 2 3 C1.3 3 0.7 3 0 3 C0 2 0 1 0 0 Z };
\path[fill=black!100,even odd rule,shift={(210,158)}] svg {M0 0 C1.3 0.3 2.6 0.7 4 1 C3.3 1.7 2.7 2.3 2 3 C1.3 2.7 0.7 2.3 0 2 C0 1.3 0 0.7 0 0 Z };
\path[fill=black!100,even odd rule,shift={(142,158)}] svg {M0 0 C1 0 2 0 3 0 C2.7 1 2.3 2 2 3 C1.3 3 0.7 3 0 3 C0 2 0 1 0 0 Z };
\path[fill=black!100,even odd rule,shift={(199,153)}] svg {M0 0 C0.5 0.3 0.5 0.3 3 2 C1.7 2.3 0.4 2.7 -1 3 C-0.7 2 -0.3 1 0 0 Z };
\path[fill=black!100,even odd rule,shift={(123,215)}] svg {M0 0 C1 0.3 2 0.7 3 1 C3 1.7 3 2.3 3 3 C2 2.7 1 2.3 0 2 C0 1.3 0 0.7 0 0 Z };
\path[fill=black!100,even odd rule,shift={(321,192)}] svg {M0 0 C0.7 0.3 1.3 0.7 2 1 C2 1.7 2 2.3 2 3 C1 3 0 3 -1 3 C-0.7 2 -0.3 1 0 0 Z };
\path[fill=black!100,even odd rule,shift={(290,192)}] svg {M0 0 C0.3 1 0.7 2 1 3 C0 3 -1 3 -2 3 C-2 2.3 -2 1.7 -2 1 C-1.3 0.7 -0.7 0.3 0 0 Z };
\path[fill=black!100,even odd rule,shift={(154,154)}] svg {M0 0 C1 0 2 0 3 0 C3 0.7 3 1.3 3 2 C2 2 1 2 0 2 C0 1.3 0 0.7 0 0 Z };
\path[fill=black!100,even odd rule,shift={(254,341)}] svg {M0 0 C0.3 0.2 0.3 0.2 2 1 C1.7 1.2 1.7 1.2 0 2 C0 1.3 0 0.7 0 0 Z };
\end{tikzpicture}}
\begin{enumerate}
\item 求容器的容积；
\item 若将容器内盛满的水从容器顶部全部抽出，至少需要做多少功？
\end{enumerate}
（长度单位：$\,\mathrm{m}$，重力加速度为$g\,\mathrm{m}\,\mathrm{s}^{-2}$，水的密度为$10^3\,\mathrm{kg}\,\mathrm{m}^{-3}$．）
\QuestionEnd

\begin{QuestionSolution}{2011 年第 20 题}
\textbf{已知与目标}\quad 容器下部来自单位圆，上部来自圆 \(x^2+(y-1)^2=1\)，在 \(y=\frac{1}{2}\) 相接。求容积及把满罐水提升到顶部的最小功。

\textbf{方法与条件}\quad 水平截面绕 \(y\) 轴形成圆盘，面积为 \(\pi x^2\)。体积微元为 \(A(y)\,\mathrm{d}y\)，功微元为这一层水的重量乘提升距离。最小功只计克服重力，出口高度为最高点 \(y=2\)。

\textbf{推导}\quad （1）底部在 \(y=-1\)，顶部在 \(y=2\)，截面积为
\[
A(y)=\begin{cases}\pi(1-y^2),&-1\le y\le\frac{1}{2},\\
\pi(2y-y^2),&\frac{1}{2}\le y\le2.
\end{cases}
\]
上段中以 \(1-y\) 代入下段，\(2(1-y)-(1-y)^2=1-y^2\)，可见两段关于 \(y=\frac{1}{2}\) 对称。故
\[
V=2\pi\int_{-1}^{\frac{1}{2}}(1-y^2)\,\mathrm{d}y
=2\pi\left[y-\frac{y^3}{3}\right]_{-1}^{\frac{1}{2}}
=\boxed{\frac{9\pi}{4}\,\mathrm{m}^3}.
\]
具体下段积分为 \((\frac{1}{2}-\frac{1}{24})-(-1+\frac{1}{3})=\frac{9}{8}\)。

（2）高度 \(y\) 处的水层重量为 \(10^3gA(y)\,\mathrm{d}y\)，提升距离为 \(2-y\)。令 \(J=\int_{-1}^2yA(y)\,\mathrm{d}y\)，换元 \(z=1-y\) 并用对称性，得 \(J=\int_{-1}^2(1-z)A(z)\,\mathrm{d}z=V-J\)，故 \(J=\frac{V}{2}\)。于是
\[
W=10^3g\int_{-1}^2(2-y)A(y)\,\mathrm{d}y
=\frac32\cdot10^3gV
=\boxed{\frac{27\pi}{8}\,10^3g\,\mathrm{J}} .
\]
\textbf{结论与检查}\quad \(\boxed{V=\frac{9\pi}{4}\,\mathrm{m}^3}\)，\(\boxed{W=\frac{27\pi}{8}\,10^3g\,\mathrm{J}}\)。水的质心在 \(y=\frac{1}{2}\)，平均提升距离为 \(2-\frac{1}{2}=\frac{3}{2}\)，也给出相同的功。
\end{QuestionSolution}

% Q 2012 17
\QuestionPage{3.4 定积分的应用}{解答题}{2012 年第 17 题}{2}{切线围成的面积与旋转体积}{177 / 234}{****}
过$(0,1)$点作曲线$L:\;y = \ln x$的切线，切点为$A$，又$L$与$x$轴交于$B$点，区域$D$由$L$与直线$AB$围成，
求区域$D$的面积及$D$绕$x$轴旋转一周所得旋转体的体积．
\QuestionEnd

\begin{QuestionSolution}{2012 年第 17 题}
\textbf{已知与目标}\quad 先由过 \((0,1)\) 的切线确定切点 \(A\)，再求对数曲线与弦 \(AB\) 围成的面积和绕 \(x\) 轴旋转的体积。区域边界使用的是弦 \(AB\)。

\textbf{方法与条件}\quad 切线斜率等于导数；面积用上边界减下边界。两边界都非负时，绕 \(x\) 轴的截面为圆环，其面积是 \(\pi\) 乘外、内半径平方之差。

\textbf{推导}\quad 设切点横坐标为 \(a>0\)，切线为 \(y=\ln a+\frac{x-a}{a}\)。代入 \((0,1)\)，\(1=\ln a-1\)，得 \(a=\mathrm{e}^2\)，故 \(A=(\mathrm{e}^2,2)\)。由 \(\ln x=0\) 得 \(B=(1,0)\)，于是弦为 \(h(x)=\frac{2(x-1)}{\mathrm{e}^2-1}\)。\((\ln x)''=-\frac{1}{x^2}<0\) 表示曲线向下弯，位于两点连线之上，故在 \([1,\mathrm{e}^2]\) 有 \(0\le h(x)\le\ln x\)。

面积为
\[
S=\int_1^{\mathrm{e}^2}\ln x\,\mathrm{d}x
-\frac2{\mathrm{e}^2-1}\int_1^{\mathrm{e}^2}(x-1)\,\mathrm{d}x
=(\mathrm{e}^2+1)-(\mathrm{e}^2-1)=\boxed{2}.
\]
绕 \(x\) 轴旋转时，垫片面积为外圆面积减去内圆面积。由分部积分得 \(\int (\ln x)^2\,\mathrm{d}x=x[(\ln x)^2-2\ln x+2]+C\)。所以
\[
\begin{aligned}
V&=\pi\int_1^{\mathrm{e}^2}\left[(\ln x)^2-\frac{4(x-1)^2}{(\mathrm{e}^2-1)^2}\right]\mathrm{d}x\\
&=\pi\left[2(\mathrm{e}^2-1)-\frac43(\mathrm{e}^2-1)\right]
=\boxed{\frac{2\pi}{3}(\mathrm{e}^2-1)}.
\end{aligned}
\]
其中 \(\int_1^{\mathrm{e}^2}(x-1)^2\,\mathrm{d}x=\frac{(\mathrm{e}^2-1)^3}{3}\)，而对数平方积分的两端代入值之差为 \(2\mathrm{e}^2-2\)。

\textbf{结论与检查}\quad \(\boxed{S=2}\)，\(\boxed{V=\frac{2\pi}{3}(\mathrm{e}^2-1)}\)。上下边界在两个端点相等，内部严格分开，所以面积和体积均为正，与结果一致。
\end{QuestionSolution}

% Q 2016 20
\QuestionPage{3.4 定积分的应用}{解答题}{2016 年第 20 题}{3}{旋转体积和表面积}{182 / 234}{****}
设$D$是由曲线$y = \sqrt{1 - x^2}$（$0 \le x \le 1$）与$\begin{cases}
  x = \cos^3 t \\
  y = \sin^3 t
\end{cases}$（$0 \le t \le \frac{\pi}{2}$）围成的平面区域，求$D$绕$x$轴旋转一周所得旋转体的体积和表面积．
\QuestionEnd

\begin{QuestionSolution}{2016 年第 20 题}
\textbf{已知与目标}\quad 区域位于第一象限的单位圆弧与参数曲线之间，求绕 \(x\) 轴后的体积及全部边界表面积。

\textbf{方法与条件}\quad 圆环截面体积用外圆盘减内圆盘；旋转曲面面积用 \(2\pi y\,\mathrm{d}s\)，即旋转周长乘弧长微元。表面积应把内、外两壁相加，还需检查有无端面。

\textbf{推导}\quad 参数曲线通常称为星形线。它满足 \(x^2+y^2=\cos^6t+\sin^6t\le(\cos^2t+\sin^2t)^3=1\)，故位于圆弧下方。外半球体积为 \(\pi\int_0^1(1-x^2)\,\mathrm{d}x=\frac{2\pi}{3}\)。计算内体积时，\(t\) 从 \(0\) 到 \(\frac{\pi}{2}\)，\(x\) 从 \(1\) 减到 \(0\)，须反转积分方向，故
\[
\begin{aligned}
V_{\text{内}}
&=\pi\int_0^{\frac{\pi}{2}}\sin^6 t\,3\cos^2t\sin t\,\mathrm{d}t\\
&=3\pi\int_0^1 u^2(1-u^2)^3\,\mathrm{d}u
=\frac{16\pi}{105},
\end{aligned}
\]
其中令 \(u=\cos t\)，展开多项式后积分为 \(3\pi(\frac{1}{3}-\frac{3}{5}+\frac{3}{7}-\frac{1}{9})=\frac{16\pi}{105}\)。因此 \(V=\frac{2\pi}{3}-\frac{16\pi}{105}=\frac{18\pi}{35}\)。

外壁为单位半球曲面，面积 \(2\pi\)。内壁的弧长微元为
\(\sqrt{(x'(t))^2+(y'(t))^2}\,\mathrm{d}t=3\sin t\cos t\,\mathrm{d}t\)，故内壁面积为
\[
2\pi\int_0^{\frac{\pi}{2}}\sin^3t\cdot3\sin t\cos t\,\mathrm{d}t
=\frac{6\pi}{5}.
\]
这里内壁面积积分可令 \(z=\sin t\)，得到 \(6\pi\int_0^1z^4\,\mathrm{d}z=\frac{6\pi}{5}\)。在 \(x=0\) 处两半径都为 \(1\)，在 \(x=1\) 都为零，没有额外端面面积。

\textbf{结论与检查}\quad \(\boxed{V=\frac{18\pi}{35}}\)，\(\boxed{S=2\pi+\frac{6\pi}{5}=\frac{16\pi}{5}}\)。体积小于外半球体积，表面积大于外半球面积，分别符合挖去内部体积及增加内壁的几何关系。
\end{QuestionSolution}

% Q 2020 18
\QuestionPage{3.4 定积分的应用}{解答题}{2020 年第 18 题}{3}{函数关系与旋转体积}{184 / 234}{****}
设函数 \(f\left( x \right)\) 的定义域为 \(\left( 0,+\infty \right)\) 且满足
\[2f\left( x \right)+x^{2}f\left( \dfrac{1}{x} \right)=\dfrac{x^{2}+2x}{\sqrt{1+x^{2}}}
\]
求 \(f\left( x \right)\)，并求曲线 \(y=f\left( x \right)\)，\(y=\dfrac{1}{2}\)，\(y=\dfrac{\sqrt{3}}{2}\) 及 \(y\) 轴所围图形绕 \(x\) 轴旋转所成旋转体的体积。
\QuestionEnd

\begin{QuestionSolution}{2020 年第 18 题}
\textbf{已知与目标}\quad 通过同时含 \(f(x)\) 与 \(f(\frac{1}{x})\) 的关系求函数，再求两水平线之间、曲线左侧区域绕 \(x\) 轴的体积。

\textbf{方法与条件}\quad \(x>0\) 时 \(\frac{1}{x}\) 仍在定义域内，可用替换自变量得到第二个方程。区域的上下边界以 \(y\) 给出，用水平条形成圆柱壳，体积微元为周长乘条长再乘厚度。

\textbf{推导}\quad 设 \(F=f(x)\)、\(H=x^2f(\frac{1}{x})\)，原式给出
\[
2F+H=\frac{x^2+2x}{\sqrt{1+x^2}}.
\]
以 \(\frac{1}{x}\) 代入原式，再乘 \(x^2\)，因 \(\sqrt{1+\frac{1}{x^2}}=\frac{\sqrt{1+x^2}}{x}\)，得到 \(F+2H=\frac{x+2x^2}{\sqrt{1+x^2}}\)。第一式乘 \(2\) 减去第二式，得 \(3F=\frac{3x}{\sqrt{1+x^2}}\)，所以 \(f(x)=\frac{x}{\sqrt{1+x^2}}\)。

对 \(\frac{1}{2}\le y\le\frac{\sqrt3}{2}\)，右边界反解为 \(x=\frac{y}{\sqrt{1-y^2}}\)，左边界为 \(x=0\)。壳半径为 \(y\)，长度为右边界的横坐标，故
\[
V=2\pi\int_{\frac{1}{2}}^{\frac{\sqrt3}{2}}\frac{y^2}{\sqrt{1-y^2}}\,\mathrm{d}y
=2\pi\int_{\frac{\pi}{6}}^{\frac{\pi}{3}}\sin^2\theta\,\mathrm{d}\theta
=\boxed{\frac{\pi^2}{6}}.
\]
第二步令 \(y=\sin\theta\)，\(\mathrm{d}y=\cos\theta\,\mathrm{d}\theta\)，\(\sqrt{1-y^2}=\cos\theta>0\)。于是 \(V=\pi[\theta-\frac{1}{2}\sin2\theta]_{\frac{\pi}{6}}^{\frac{\pi}{3}}\)，两端的 \(\sin2\theta\) 都为 \(\frac{\sqrt3}{2}\)，相减消去。

\textbf{结论与检查}\quad \(\boxed{f(x)=\frac{x}{\sqrt{1+x^2}},\ x>0}\)，\(\boxed{V=\frac{\pi^2}{6}}\)。代回函数关系，左边为 \(\frac{2x+x^2}{\sqrt{1+x^2}}\)，符合已知；反解中取正根也符合第一象限区域。
\end{QuestionSolution}

% Q 2026 20
\QuestionPage{3.4 定积分的应用}{解答题}{2026 年第 20 题}{2}{拐点与无界旋转体积}{188 / 234}{****}
（本题满分 12 分）
已知 \(M(x_0,y_0)\) 是曲线
\[y=\frac{1}{1+x^2}\qquad (x\ge 0)
\]
的拐点，\(O\) 为原点.记 \(D\) 是第一象限中以曲线
\[y=\frac{1}{1+x^2}\qquad (x\ge x_0),
\]
线段 \(OM\) 及 \(x\) 正半轴为边界的无界区域，求 \(D\) 绕 \(x\) 轴旋转所成旋转体的体积.
\QuestionEnd

\begin{QuestionSolution}{2026 年第 20 题}
\textbf{已知与目标}\quad 区域左侧上边界为线段 \(OM\)，右侧为给定曲线，分界点 \(M\) 是拐点。求整个无界区域绕 \(x\) 轴的体积。

\textbf{方法与条件}\quad 拐点表示曲线弯曲方向改变，需检查二阶导数变号，不能只求其零点。先定位拐点，再把体积分为圆锥和右侧圆盘积分；右侧无界，须取上限极限。

\textbf{推导}\quad \(f'(x)=-\frac{2x}{(1+x^2)^2}\)，\(f''(x)=\frac{6x^2-2}{(1+x^2)^3}\)。分母为正，分子在 \(x_0=\frac{1}{\sqrt3}\) 由负变正，所以这就是唯一内部拐点，\(y_0=\frac{1}{1+\frac{1}{3}}=\frac{3}{4}\)。线段 \(OM\) 旋转成高 \(x_0\)、底半径 \(y_0\) 的圆锥，
\[
V_{\text{左}}=\frac{\pi}{3}\left(\frac{3}{4}\right)^2\frac{1}{\sqrt3}
=\frac{\pi\sqrt3}{16}.
\]
右侧的圆盘半径为 \(f(x)\)，所以
\[
\begin{aligned}
V_{\text{右}}
&=\pi\int_{\frac{1}{\sqrt3}}^\infty\frac{\mathrm{d}x}{(1+x^2)^2}\\
&=\pi\left[\frac{x}{2(1+x^2)}+\frac12\arctan x\right]_{\frac{1}{\sqrt3}}^\infty
=\pi\left(\frac\pi6-\frac{\sqrt3}{8}\right).
\end{aligned}
\]
原函数可令 \(x=\tan\theta\) 得到：\(\frac{\mathrm{d}x}{(1+x^2)^2}=\cos^2\theta\,\mathrm{d}\theta\)。代回后，上端为 \(\frac{\pi}{4}\)，下端为 \(\frac{\sqrt3}{8}+\frac{\pi}{12}\)，相减为 \(\frac{\pi}{6}-\frac{\sqrt3}{8}\)。

\textbf{结论与检查}\quad 两部分相加，\(\boxed{V=\pi(\frac{\pi}{6}-\frac{\sqrt3}{16})}\)。右侧被积函数按 \(\frac{1}{x^4}\) 的速度衰减，积分收敛；两个分区没有交叠，公共截面不增加体积。
\end{QuestionSolution}

\SectionHeading{3.5}{一元函数积分学及其应用}{反常积分}{定理与要点\quad 教材 PDF 第 244 页起\hfill 本节 13 题}
\pdfbookmark[2]{3.5 反常积分}{section-3.5}
\Theory{定理 5.1}{250}{无穷区间比较准则 I}{当 \(x\) 充分大时若有 \(0\le f(x)\le g(x)\)，则 \(g\) 的反常积分收敛可推出 \(f\) 的反常积分收敛，\(f\) 的反常积分发散可推出 \(g\) 的反常积分发散。}
\Theory{定理 5.2}{251}{无穷区间比较准则 II}{非负被积函数之比趋于有限正数，则同敛散。比值趋零或趋正无穷时按相应方向比较。}
\Theory{定理 5.3}{252}{无穷区间绝对收敛准则}{绝对值的反常积分收敛，则原反常积分收敛。}
\Theory{定理 5.4}{253}{无界函数比较准则 I}{在同一奇点附近，非负被积函数可按大小关系比较反常积分的敛散。}
\Theory{定理 5.5}{253}{无界函数比较准则 II}{靠近同一奇点时，非负被积函数比值趋有限正数，则同敛散。零及正无穷的情形按方向比较。}
\Theory{定理 5.6}{254}{无界函数绝对收敛准则}{在奇点处，绝对值的反常积分收敛，则原反常积分收敛。}
\Theory{要点}{245--249}{\(p\) 型积分与多个端点}{\(x^{-p}\) 的积分在无穷远处收敛的条件是 \(p>1\)，在零点附近收敛的条件是 \(p<1\)。有多个反常点时须分段检查。}
% Q 2015 1
\QuestionPage{3.5 反常积分}{选择题}{2015 年第 1 题}{1}{反常积分的收敛判别}{191 / 234}{**}
下列反常积分收敛的是\pickout{}
\begin{abcd}
\item $\int_2^{+\infty} \frac{1}{\sqrt{x}}\,\mathrm{d} x$．
\item $\int_2^{+\infty} \frac{\ln x}{x}\,\mathrm{d} x$．
\item $\int_2^{+\infty} \frac{1}{x\ln x}\,\mathrm{d} x$．
\item $\int_2^{+\infty} \frac{x}{\mathrm{e}^x}\,\mathrm{d} x$．
\end{abcd}
\QuestionEnd

\begin{QuestionSolution}{2015 年第 1 题}
\textbf{已知与目标}\quad 四个积分在有限下限 \(2\) 处都没有奇点，只需找出上限趋于无穷时有有限积分值的一项。

\textbf{方法与条件}\quad 反常积分收敛是指 \(\lim_{R\to\infty}\int_2^Rh(x)\,\mathrm{d}x\) 存在且有限。逐项求原函数即可直接判断。

\textbf{推导}\quad A 的截断积分为 \(2\sqrt R-2\sqrt2\to+\infty\)。B 为 \(\frac{1}{2}[(\ln R)^2-(\ln2)^2]\to+\infty\)。C 为 \(\ln\ln R-\ln\ln2\to+\infty\)。D 分部积分得原函数 \(-(x+1)\mathrm{e}^{-x}\)，所以
\[
\int_2^Rx\mathrm{e}^{-x}\,\mathrm{d}x
=-(R+1)\mathrm{e}^{-R}+3\mathrm{e}^{-2}\longrightarrow3\mathrm{e}^{-2}.
\]
\textbf{结论与检查}\quad 只有 D 收敛，选 \(\boxed{\mathrm{D}}\)。指数衰减使 \((R+1)\mathrm{e}^{-R}\to0\)，计算所得正值也符合被积函数为正的性质。
\end{QuestionSolution}

% Q 2016 3
\QuestionPage{3.5 反常积分}{选择题}{2016 年第 3 题}{1}{反常积分敛散性}{192 / 234}{**}
反常积分\circnum{1}$\int_{-\infty}^0 \frac{1}{x^2}\mathrm{e}^{\frac{1}{x}} \,\mathrm{d} x$，
\par\circnum{2}$\int_0^{\text{+}\infty} \frac{1}{x^2}\mathrm{e}^{\frac{1}{x}} \,\mathrm{d} x$的敛散性为\pickout{}
\begin{abcd}
\item \circnum{1}收敛，\circnum{2}收敛．
\item \circnum{1}收敛，\circnum{2}发散．
\item \circnum{1}发散，\circnum{2}收敛．
\item \circnum{1}发散，\circnum{2}发散．
\end{abcd}
\QuestionEnd

\begin{QuestionSolution}{2016 年第 3 题}
\textbf{已知与目标}\quad 两积分的被积函数相同，但分别靠近零点左、右两侧，指数行为不同。需要检查各自的无穷端和零点端。

\textbf{方法与条件}\quad \([-\mathrm{e}^{\frac{1}{x}}]'=\frac{1}{x^2}\mathrm{e}^{\frac{1}{x}}\)。有两个反常端点时须分段分别取极限，不能让两端的发散相互抵消。

\textbf{推导}\quad 第一项在 \(-1\) 处分段。\(x\to-\infty\) 时 \(\mathrm{e}^{\frac{1}{x}}\to1\)，\(x\to0^-\) 时趋于 \(0\)，故
\[
\int_{-\infty}^{-1}\frac{\mathrm{e}^{\frac{1}{x}}}{x^2}\,\mathrm{d}x=1-\mathrm{e}^{-1},\qquad
\int_{-1}^{0}\frac{\mathrm{e}^{\frac{1}{x}}}{x^2}\,\mathrm{d}x=\mathrm{e}^{-1}.
\]
两段均有限，总值为 \(1\)。第二项靠近零点时，\(\int_\varepsilon^1\frac{\mathrm{e}^{\frac{1}{x}}}{x^2}\,\mathrm{d}x=\mathrm{e}^{\frac{1}{\varepsilon}}-\mathrm{e}\to+\infty\)。所以即使无穷端积分有限，第二项仍发散。

\textbf{结论与检查}\quad 第一项收敛，第二项发散，选 \(\boxed{\mathrm{B}}\)。零点两侧指数分别趋于零和无穷，不能把两侧混用。
\end{QuestionSolution}

% Q 2019 3
\QuestionPage{3.5 反常积分}{选择题}{2019 年第 3 题}{1}{反常积分审敛}{193 / 234}{**}
下列反常积分发散的是\pickout{}
\begin{abcd}
\item $\int_0^{+\infty} x\mathrm{e}^{- x} \,\mathrm{d} x$．
\item $\int_0^{+\infty} x\mathrm{e}^{- x^2} \,\mathrm{d} x$．
\item $\int_0^{+\infty} \frac{\arctan x}{1 + x^2}\,\mathrm{d} x$．
\item $\int_0^{+\infty} \frac{x}{1 + x^2}\,\mathrm{d} x$．
\end{abcd}
\QuestionEnd

\begin{QuestionSolution}{2019 年第 3 题}
\textbf{已知与目标}\quad 从四个积分中找出发散的一项。所有被积函数在零点连续，因此只需检查无穷远端。

\textbf{方法与条件}\quad 分别用分部积分或换元求值。有限结果表示收敛，截断积分趋于 \(+\infty\) 表示发散。

\textbf{推导}\quad A 的原函数为 \(-(x+1)\mathrm{e}^{-x}\)，取两端极限得到 \(1\)。B 令 \(u=x^2\)、\(\mathrm{d}u=2x\,\mathrm{d}x\)，得 \(\frac{1}{2}\int_0^\infty\mathrm{e}^{-u}\,\mathrm{d}u=\frac{1}{2}\)。C 令 \(v=\arctan x\)、\(\mathrm{d}v=\frac{1}{1+x^2}\mathrm{d}x\)，得 \(\int_0^{\frac{\pi}{2}}v\,\mathrm{d}v=\frac{\pi^2}{8}\)。D 在上限 \(R\) 处的积分为
\[
\int_0^R\frac{x}{1+x^2}\,\mathrm{d}x
=\left[\frac{1}{2}\ln(1+x^2)\right]_0^R
=\frac{1}{2}\ln(1+R^2)\longrightarrow+\infty.
\]
\textbf{结论与检查}\quad 发散项是 \(\boxed{\mathrm{D}}\)。其被积函数在无穷远与 \(\frac{1}{x}\) 等价，也能用对数发散核对。
\end{QuestionSolution}

% Q 2009 10
\QuestionPage{3.5 反常积分}{填空题}{2009 年第 10 题}{2}{反常积分求参数}{196 / 234}{**}
已知$\int_{-\infty}^{+\infty}\mathrm{e}^{k|x|}\,\mathrm{d} x = 1$，则$k=$ \blank{}．
\QuestionEnd

\begin{QuestionSolution}{2009 年第 10 题}
\textbf{已知与目标}\quad 已知整个实轴上 \(\mathrm{e}^{k|x|}\) 的积分为 \(1\)，求实参数 \(k\)，并保证反常积分确实收敛。

\textbf{方法与条件}\quad 被积函数是偶函数，即正、负 \(x\) 对应的值相同。两半轴积分相等，分别收敛后才可写成正半轴的两倍。

\textbf{推导}\quad 若 \(k=0\)，被积函数恒为 \(1\)，积分发散；若 \(k>0\)，正半轴指数增长，也发散。因此必须 \(k<0\)。此时两半轴均收敛，且
\[
\int_{-\infty}^{\infty}\mathrm{e}^{k|x|}\,\mathrm{d}x
=2\int_0^\infty\mathrm{e}^{kx}\,\mathrm{d}x
=2\left[\frac{\mathrm{e}^{kx}}{k}\right]_0^\infty=-\frac{2}{k}.
\]
令 \(-\frac{2}{k}=1\)，解得 \(k=-2\)。

\textbf{结论与检查}\quad \(\boxed{k=-2}\)。它满足必要条件 \(k<0\)，代回积分为 \(-\frac{2}{-2}=1\)。
\end{QuestionSolution}

% Q 2011 12
\QuestionPage{3.5 反常积分}{填空题}{2011 年第 12 题}{2}{无穷区间上的积分}{197 / 234}{**}
设函数$f(x) = \begin{cases}
  \lambda \mathrm{e}^{- \lambda x}, & x > 0, \\
                         0, & x \le 0,
\end{cases}$ $\lambda > 0$，则$\int_{-\infty}^{+\infty} xf(x)\,\mathrm{d} x =$ \blank{}．
\QuestionEnd

\begin{QuestionSolution}{2011 年第 12 题}
\textbf{已知与目标}\quad \(\lambda>0\)，负半轴 \(f=0\)，正半轴为 \(\lambda\mathrm{e}^{-\lambda x}\)，求 \(xf(x)\) 的全轴积分。

\textbf{方法与条件}\quad 负半轴没有贡献，正半轴用分部积分。\(\lambda>0\) 使指数衰减，保证边界项 \(R\mathrm{e}^{-\lambda R}\to0\)。

\textbf{推导}\quad 先在 \([0,R]\) 积分，取 \(u=x\)、\(\mathrm{d}v=\lambda\mathrm{e}^{-\lambda x}\,\mathrm{d}x\)，得
\[
\int_0^R x\lambda\mathrm{e}^{-\lambda x}\,\mathrm{d}x
=[-x\mathrm{e}^{-\lambda x}]_0^R+\int_0^R\mathrm{e}^{-\lambda x}\,\mathrm{d}x
=-R\mathrm{e}^{-\lambda R}+\frac{1-\mathrm{e}^{-\lambda R}}{\lambda}.
\]
令 \(R\to\infty\)，于是
\[
\int_{-\infty}^\infty xf(x)\,\mathrm{d}x
=\int_0^\infty x\lambda \mathrm{e}^{-\lambda x}\,\mathrm{d}x
=\int_0^\infty \mathrm{e}^{-\lambda x}\,\mathrm{d}x
=\boxed{\frac1\lambda}.
\]
\textbf{结论与检查}\quad \(\boxed{\frac{1}{\lambda}}\)，条件为 \(\lambda>0\)。令 \(z=\lambda x\) 也可将积分化为 \(\frac{1}{\lambda}\int_0^\infty z\mathrm{e}^{-z}\,\mathrm{d}z\)，与结果一致。
\end{QuestionSolution}

% Q 2014 9
\QuestionPage{3.5 反常积分}{填空题}{2014 年第 9 题}{2}{无穷区间上的反常积分}{198 / 234}{**}
$\int_{-\infty}^{1} \frac{1}{x^2 + 2x + 5}\,\mathrm{d} x =$ \blank{}．
\QuestionEnd

\begin{QuestionSolution}{2014 年第 9 题}
\textbf{已知与目标}\quad 求下限为负无穷的有理函数积分，分母 \(x^2+2x+5\) 始终为正，没有有限奇点。

\textbf{方法与条件}\quad 配方 \(x^2+2x+5=(x+1)^2+4\)，换元化成 \(\frac{1}{1+u^2}\) 的积分，其原函数为 \(\arctan u\)。下端必须使用 \(\arctan(-\infty)=-\frac{\pi}{2}\)。

\textbf{推导}\quad 令 \(u=\frac{x+1}{2}\)，\(\mathrm{d}x=2\,\mathrm{d}u\)，分母为 \(4(1+u^2)\)，所以系数为 \(\frac{1}{2}\)。当 \(x=1\)，\(u=1\)；当 \(x\to-\infty\)，\(u\to-\infty\)。故
\[
\int_{-\infty}^1\frac{\mathrm{d}x}{x^2+2x+5}
=\frac12[\arctan u]_{-\infty}^1
=\frac12\left(\frac\pi4+\frac\pi2\right)
=\boxed{\frac{3\pi}{8}}.
\]
\textbf{结论与检查}\quad \(\boxed{\frac{3\pi}{8}}\)。全实轴上的同一积分为 \(\frac{\pi}{2}\)，本题是其中一部分，\(0<\frac{3\pi}{8}<\frac{\pi}{2}\) 符合范围。
\end{QuestionSolution}

% Q 2021 11
\QuestionPage{3.5 反常积分}{填空题}{2021 年第 11 题}{2}{无穷区间上的积分}{199 / 234}{**}
\[\int_{-\infty}^{+\infty} |x|3^{-x^2}\,\mathrm{d}x= \blank{}
\]
\QuestionEnd

\begin{QuestionSolution}{2021 年第 11 题}
\textbf{已知与目标}\quad 求含 \(|x|\) 与 \(x^2\) 的全轴积分。两者均关于零点对称，可只算正半轴。

\textbf{方法与条件}\quad 偶函数的两半轴积分相等；先计算正半轴并确认收敛。令 \(u=x^2\)，微分 \(\mathrm{d}u=2x\,\mathrm{d}x\) 与外部两倍因子配合。

\textbf{推导}\quad 在正半轴 \(|x|=x\)，故原积分为 \(2\int_0^\infty x3^{-x^2}\,\mathrm{d}x\)。换元后 \(u\) 仍从零到无穷，得到
\[
\int_0^\infty 3^{-u}\,\mathrm{d}u
=\left[-\frac{3^{-u}}{\ln3}\right]_0^\infty
=\boxed{\frac1{\ln3}}.
\]
由于 \(3^{-u}=\mathrm{e}^{-u\ln3}\)，且 \(\ln3>0\)，无穷端的指数项确实趋于零。

\textbf{结论与检查}\quad \(\boxed{\frac{1}{\ln3}}\)。正半轴积分有限，负半轴由对称性也收敛，因而全轴积分成立；外部的 \(2\) 已被换元的 \(2x\) 吸收。
\end{QuestionSolution}

% Q 2010 4
\QuestionPage{3.5 反常积分}{选择题}{2010 年第 4 题}{1}{反常积分审敛}{189 / 234}{***}
设$m,n$是正整数，则反常积分$\int_0^1 \frac{\sqrt[m]{\ln^2(1 - x)}}{\sqrt[n]{x}}\,\mathrm{d} x$的收敛性\pickout{}
\begin{abcd}
\item 仅与$m$的取值有关．
\item 仅与$n$的取值有关．
\item 与$m,n$取值都有关．
\item 与$m,n$取值都无关．
\end{abcd}
\QuestionEnd

\begin{QuestionSolution}{2010 年第 4 题}
\textbf{已知与目标}\quad \(m,n\) 是正整数，被积函数非负，在 \(0\) 和 \(1\) 都可能出现无界现象。要判断收敛性是否依赖参数。

\textbf{方法与条件}\quad 多个反常点须分开检查。等价比较指两个非负函数之比趋于有限正数时，积分同敛散；在零点 \(\int_0^a x^q\,\mathrm{d}x\) 收敛当且仅当 \(q>-1\)。

\textbf{推导}\quad 当 \(x\to0^+\)，\(|\ln(1-x)|\sim x\)，故被积函数等价于 \(x^{\frac{2}{m}-\frac{1}{n}}\)。因 \(\frac{1}{n}\le1\) 且 \(\frac{2}{m}>0\)，指数严格大于 \(-1\)，左端总收敛。在 \([\frac{1}{2},1)\) 上，\(x^{-\frac{1}{n}}\le2^{\frac{1}{n}}\)，只需检查对数因子。令 \(u=-\ln(1-x)\)，\(\mathrm{d}x=\mathrm{e}^{-u}\,\mathrm{d}u\)，得
\[
\int_{\frac{1}{2}}^1|\ln(1-x)|^{\frac{2}{m}}\,\mathrm{d}x
=\int_{\ln2}^\infty u^{\frac{2}{m}}\mathrm{e}^{-u}\,\mathrm{d}u<\infty .
\]
最后一个积分有限，因为对固定 \(m\)，当 \(u\) 足够大时 \(u^{\frac{2}{m}}\le\mathrm{e}^{\frac{u}{2}}\)，其被积函数不超过可积的 \(\mathrm{e}^{-\frac{u}{2}}\)。

\textbf{结论与检查}\quad 对所有正整数 \(m,n\) 都收敛，选 \(\boxed{\mathrm{D}}\)。即使 \(n=1\)，左端指数也是 \(-1+\frac{2}{m}>-1\)，临界情形仍已覆盖。
\end{QuestionSolution}

% Q 2013 4
\QuestionPage{3.5 反常积分}{选择题}{2013 年第 4 题}{1}{分段函数的反常积分}{190 / 234}{***}
设函数$f(x) =\begin{cases}
  \frac{1}{(x - 1)^{\alpha - 1}}, & 1 < x < \mathrm{e}, \\
   \frac{1}{x\ln^{\alpha + 1} x}, & x \ge \mathrm{e},
\end{cases}$若反常积分$\int_1^{+\infty} f(x) \,\mathrm{d} x$收敛，则\pickout{}
\begin{abcd}
\item $\alpha < - 2$．
\item $\alpha > 2$．
\item $- 2 < \alpha < 0$．
\item $0 < \alpha < 2$．
\end{abcd}
\QuestionEnd

\begin{QuestionSolution}{2013 年第 4 题}
\textbf{已知与目标}\quad 被积函数在 \(\mathrm{e}\) 处分段，需使 \(x=1\) 附近和 \(+\infty\) 附近的积分都收敛，求 \(\alpha\) 的条件。

\textbf{方法与条件}\quad 零点处的幂积分 \(\int_0^a u^q\,\mathrm{d}u\) 要求 \(q>-1\)；无穷远的 \(\int_1^\infty v^{-p}\,\mathrm{d}v\) 要求 \(p>1\)。两端条件须同时满足。

\textbf{推导}\quad 左段令 \(u=x-1\)，得到 \(\int_0^{\mathrm{e}-1}u^{1-\alpha}\,\mathrm{d}u\)。要求 \(1-\alpha>-1\)，即 \(\alpha<2\)。右段令 \(v=\ln x\)，则 \(\mathrm{d}v=\frac{1}{x}\mathrm{d}x\)，上、下限变为 \(+\infty,1\)，积分为 \(\int_1^\infty v^{-\alpha-1}\,\mathrm{d}v\)。要求 \(\alpha+1>1\)，即 \(\alpha>0\)。取交集便得 \(0<\alpha<2\)。

\textbf{结论与检查}\quad 选 \(\boxed{\mathrm{D}}\)。\(\alpha=0\) 时右段为 \(\int\frac{1}{v}\,\mathrm{d}v\) 发散，\(\alpha=2\) 时左段为 \(\int\frac{1}{u}\,\mathrm{d}u\) 发散，因此两个端点均排除。
\end{QuestionSolution}

% Q 2022 5
\QuestionPage{3.5 反常积分}{选择题}{2022 年第 5 题}{1}{反常积分的收敛参数}{194 / 234}{***}
设 \(p\) 为常数，若反常积分
\(\int_{0}^{1} \frac{\ln x}{x^{p}\left( 1-x \right)^{1-p}} \,\mathrm{d}x\)
收敛，则 \(p\) 的取值范围是 \(\left( \ \right)\)
\choices{\(\left( -1,1 \right)\)}{\(\left( -1,2 \right)\)}{\(\left( -\infty,1 \right)\)}{\(\left( -\infty,2 \right)\)}
\QuestionEnd

\begin{QuestionSolution}{2022 年第 5 题}
\textbf{已知与目标}\quad 被积函数在 \((0,1)\) 恒负，两个端点都可能影响收敛，求参数 \(p\) 的完整范围。

\textbf{方法与条件}\quad 恒负函数的积分收敛等价于其绝对值积分收敛。两非负函数之比趋于 \(1\) 时可用等价比较；零点幂积分 \(\int_0^a t^q\,\mathrm{d}t\) 的条件为 \(q>-1\)。

\textbf{推导}\quad 在 \(x\to0^+\)，\((1-x)^{1-p}\to1\)，绝对值等价于 \(x^{-p}(-\ln x)\)。令 \(u=-\ln x\)，\(\mathrm{d}x=-\mathrm{e}^{-u}\,\mathrm{d}u\)，其积分成为 \(\int_{u_0}^\infty u\mathrm{e}^{(p-1)u}\,\mathrm{d}u\)：\(p<1\) 时指数衰减而收敛，\(p=1\) 时为 \(\int u\,\mathrm{d}u\) 发散，\(p>1\) 时指数增长也发散。

在 \(x\to1^-\)，\(-\ln x\sim1-x\)，\(x^p\to1\)，因此绝对值等价于 \((1-x)^p\)，收敛恰需 \(p>-1\)。两条件同时成立得到 \(-1<p<1\)。

\textbf{结论与检查}\quad \(\boxed{p\in(-1,1)}\)，选 A。\(p=-1\) 在右端产生 \(\frac{1}{1-x}\) 型发散，\(p=1\) 在左端发散，两端点均不能加入。
\end{QuestionSolution}

% Q 2025 11
\QuestionPage{3.5 反常积分}{填空题}{2025 年第 11 题}{2}{反常积分求参数}{200 / 234}{***}
设 \(\int_{1}^{+\infty} \frac{a}{x\left( 2x + a \right)} \,\mathrm{d} x = \ln 2\)，则 \(a =\) \blank{}.
\QuestionEnd

\begin{QuestionSolution}{2025 年第 11 题}
\textbf{已知与目标}\quad 已知积分值为 \(\ln2\)，求参数 \(a\)。必须先保证积分路径上没有导致发散的分母零点。

\textbf{方法与条件}\quad 当 \(a\le-2\)，零点 \(x=-\frac{a}{2}\ge1\) 位于积分区间内或端点，分子非零，出现 \(\frac{1}{x-c}\) 型不可积奇点，故排除。对 \(a>-2\)，分母在 \([1,\infty)\) 为正，可拆部分分式并把对数合并取极限。

\textbf{推导}\quad 对允许的参数，
\[
\frac{a}{x(2x+a)}=\frac1x-\frac{2}{2x+a},\qquad
\int_1^\infty \frac{a}{x(2x+a)}\,\mathrm{d}x
=\left[\ln\frac{x}{2x+a}\right]_1^\infty
=\ln\frac{2+a}{2}.
\]
上端对数趋于 \(\ln\frac{1}{2}\)，下端为 \(\ln\frac{1}{2+a}\)，相减为 \(\ln\frac{2+a}{2}\)。由 \(\ln\frac{2+a}{2}=\ln2\) 及对数严格递增，得 \(\frac{2+a}{2}=2\)，所以 \(a=2\)。

\textbf{结论与检查}\quad \(\boxed{a=2}\)。它满足 \(a>-2\)，代回积分为 \(\ln\frac{4}{2}=\ln2\)，且不存在积分区间内的极点。
\end{QuestionSolution}

% Q 2026 11
\QuestionPage{3.5 反常积分}{填空题}{2026 年第 11 题}{2}{反常积分的收敛参数}{201 / 234}{***}
设 \(p\) 为常数，若反常积分
\[\int_0^{+\infty}\frac{\arctan x}{x^p(x+1)}\,\mathrm{d}x
\]
收敛，则 \(p\) 的取值范围是 \blank{}.
\QuestionEnd

\begin{QuestionSolution}{2026 年第 11 题}
\textbf{已知与目标}\quad 对任意实数 \(p\)，被积函数在 \(x>0\) 非负连续。零点和无穷远都要检查，目标是两端同时收敛的参数范围。

\textbf{方法与条件}\quad 等价比较指非负函数之比趋于有限正数时同敛散。\(\int_0^1x^q\,\mathrm{d}x\) 要求 \(q>-1\)，\(\int_1^\infty x^{-r}\,\mathrm{d}x\) 要求 \(r>1\)，两端的判定方向不同。

\textbf{推导}\quad 在 \(x\to0^+\)，\(\arctan x\sim x\)、\(x+1\sim1\)，所以被积函数等价于 \(x^{1-p}\)。要求 \(1-p>-1\)，得到 \(p<2\)。在 \(x\to\infty\)，\(\arctan x\to\frac{\pi}{2}\)、\(x+1\sim x\)，所以被积函数等价于 \(\frac{\pi}{2}x^{-p-1}\)。要求 \(p+1>1\)，得到 \(p>0\)。两条件取交集得 \(0<p<2\)。

\textbf{结论与检查}\quad \(\boxed{p\in(0,2)}\)。\(p=0\) 时无穷端是 \(\frac{1}{x}\) 型发散，\(p=2\) 时零点端是 \(\frac{1}{x}\) 型发散，故两个边界均排除。
\end{QuestionSolution}

% Q 2024 7
\QuestionPage{3.5 反常积分}{选择题}{2024 年第 7 题}{1}{非负函数的反常积分审敛}{195 / 234}{****}
设非负函数 \(f(x)\) 在 \([0,+\infty)\) 上连续，给出三个命题：
\begin{enumerate}
\item 若 \(\int_0^{+\infty}f^2(x)\,\mathrm{d}x\) 收敛，则 \(\int_0^{+\infty}f(x)\,\mathrm{d}x\) 收敛；
\item 若存在 \(p>1\)，使极限 \(\lim_{x\to+\infty}x^p f(x)\) 存在，则 \(\int_0^{+\infty}f(x)\,\mathrm{d}x\) 收敛；
\item 若 \(\int_0^{+\infty}f(x)\,\mathrm{d}x\) 收敛，则存在 \(p>1\)，使极限 \(\lim_{x\to+\infty}x^p f(x)\) 存在。
\end{enumerate}
其中正确命题的个数是（ ）
\choices{\(0\)}{\(1\)}{\(2\)}{\(3\)}
\QuestionEnd

\begin{QuestionSolution}{2024 年第 7 题}
\textbf{已知与目标}\quad \(f\) 在非负半轴连续且非负，需逐一判断平方可积、幂次衰减和可积之间的三个命题，再数正确命题。

\textbf{方法与条件}\quad “极限存在”按有限实数极限理解。证明一个普遍命题需对全部合条件函数成立；否定它只需给出一个合条件反例。比较判别允许用可积的上界控制非负函数。

\textbf{推导}\quad （1）取 \(f(x)=\frac{1}{1+x}\)，它连续非负，\(\int_0^\infty f^2=\left[-\frac{1}{1+x}\right]_0^\infty=1\)，但 \(\int_0^R f=\ln(1+R)\to\infty\)，所以命题一为假。

（2）若 \(x^pf(x)\) 有有限极限，便在充分大的 \(x\) 上有上界 \(C>0\)，因此 \(0\le f(x)\le\frac{C}{x^p}\)。因 \(p>1\)，右边的无穷积分收敛；剩余有限区间由连续性可积。命题二为真。

（3）取 \(f(x)=\frac{1}{(x+\mathrm{e})\ln^2(x+\mathrm{e})}\)，它在 \([0,\infty)\) 连续非负。令 \(u=\ln(x+\mathrm{e})\)，得 \(\int_0^\infty f=\int_1^\infty\frac{1}{u^2}\,\mathrm{d}u=1\)。但对任意 \(p>1\)，\(x^pf(x)\sim\frac{x^{p-1}}{(\ln x)^2}\to+\infty\)，因为任意正次幂都比对数平方增长更快，有限极限不存在。因此命题三为假。

\textbf{结论与检查}\quad 仅命题二正确，共 \(\boxed{1}\) 个，选 B。两个反例均在整个指定区间连续，且满足非负条件。
\end{QuestionSolution}

\SectionHeading{4.1}{常微分方程}{几类简单的微分方程}{定理与要点\quad 教材 PDF 第 262 页起\hfill 本节 17 题}
\pdfbookmark[1]{常微分方程}{chapter-4.1}
\pdfbookmark[2]{4.1 几类简单的微分方程}{section-4.1}
\Theory{方法}{263--265}{方程、初值与通解}{先识别阶数、未知函数及初值。通解中的任意常数由给定条件确定。}
\Theory{方法}{266--267}{分离变量}{将同一变量的项移到等式同侧后积分。除以含未知函数的因子前，检查其为零时的常值解。}
\Theory{方法}{267--270}{一阶线性方程}{对 \(y'+p(x)y=q(x)\)，先求对应的齐次方程通解，再用常数变易法求非齐次方程通解。}
\Theory{方法}{270--274}{变量代换}{齐次型方程可令 \(u=\frac{y}{x}\)。Bernoulli 方程可作幂函数代换化为线性方程。}
\Theory{方法}{274--277}{高阶方程降阶}{方程不显含 \(y\) 时可令 \(p=y'\)。方程不显含 \(x\) 时可将 \(p\) 视为 \(y\) 的函数，写成 \(y''=p\frac{\mathrm{d}p}{\mathrm{d}y}\)。}
% Q 2010 2
\QuestionPage{4.1 几类简单的微分方程}{选择题}{2010 年第 2 题}{1}{一阶线性方程解的关系}{202 / 234}{**}
设$y_1,y_2$是一阶线性非齐次微分方程$y' + p(x)y = q(x)$的两个特解，
若常数$\lambda ,\mu$使$\lambda y_1 + \mu y_2$是该方程的解，
$\lambda y_1 - \mu y_2$是该方程对应的齐次方程的解，则\pickout{}
\begin{abcd}
\item $\lambda = \frac{1}{2},\mu = \frac{1}{2}$．
\item $\lambda = - \frac{1}{2},\mu = - \frac{1}{2}$．
\item $\lambda = \frac{2}{3},\mu = \frac{1}{3}$．
\item $\lambda = \frac{2}{3},\mu = \frac{2}{3}$．
\end{abcd}
\QuestionEnd

\begin{QuestionSolution}{2010 年第 2 题}
\textbf{已知与目标}\quad \(y_1,y_2\) 都满足同一个非齐次方程，要确定两种组合分别仍满足原方程、对应齐次方程时的常数 \(\lambda,\mu\)。齐次方程指右端改为零的方程。

\textbf{方法与条件}\quad 导数对常数倍及加减满足分配规律，因此可以直接代入组合，不必先求出 \(y_1,y_2\)。非齐次表示 \(q(x)\) 不恒等于零。

\textbf{推导}\quad 已知 \(y_1'+py_1=y_2'+py_2=q\)。把两种组合分别代入左端，得到
\[
(\lambda+\mu)q(x)=q(x),\qquad(\lambda-\mu)q(x)=0.
\]
由于至少存在一点使 \(q(x)\ne0\)，而两式要在整个区间成立，必须有 \(\lambda+\mu=1\)、\(\lambda-\mu=0\)。两式相加得 \(2\lambda=1\)，继而 \(\mu=\lambda=\frac{1}{2}\)。

\textbf{结论与检查}\quad 选 A。代回系数关系，和为 \(1\)、差为 \(0\)，所以两种组合的右端确为 \(q\) 和 \(0\)，符合题意。
\end{QuestionSolution}

% Q 2011 10
\QuestionPage{4.1 几类简单的微分方程}{填空题}{2011 年第 10 题}{1}{一阶线性方程初值问题}{203 / 234}{**}
微分方程$y' + y = \mathrm{e}^{- x}\cos x$满足条件$y(0) = 0$的解为$y =$ \blank{}．
\QuestionEnd

\begin{QuestionSolution}{2011 年第 10 题}
\textbf{已知与目标}\quad 已知一阶方程 \(y'+y=\mathrm{e}^{-x}\cos x\) 及初值 \(y(0)=0\)，求满足两者的函数。

\textbf{方法与条件}\quad 用积分因子法：寻找一个乘数，使左边成为一个乘积的导数。本题乘数为 \(\mathrm{e}^x\)，它处处非零，乘上或除去它均不改变解。

\textbf{推导}\quad 由乘积求导公式，\((\mathrm{e}^xy)'=\mathrm{e}^xy'+\mathrm{e}^xy\)。原方程两边乘 \(\mathrm{e}^x\)，得
\[
(\mathrm{e}^xy)'=\cos x,
\qquad \mathrm{e}^xy=\int\cos x\,\mathrm{d}x=\sin x+C.
\]
代入 \(x=0,y=0\) 得 \(C=0\)，故 \(\boxed{y=\mathrm{e}^{-x}\sin x}\)。

\textbf{结论与检查}\quad 求导得 \(y'=\mathrm{e}^{-x}(\cos x-\sin x)\)，所以 \(y'+y=\mathrm{e}^{-x}\cos x\)；同时 \(y(0)=0\)。答案在全部实数范围成立。
\end{QuestionSolution}

% Q 2012 12
\QuestionPage{4.1 几类简单的微分方程}{填空题}{2012 年第 12 题}{2}{一阶微分方程初值问题}{204 / 234}{**}
微分方程$y\,\mathrm{d} x + (x - 3y^2)\,\mathrm{d} y = 0$满足初始条件$y|_{x = 1}$=1的解为$y =$ \blank{}．
\QuestionEnd

\begin{QuestionSolution}{2012 年第 12 题}
\textbf{已知与目标}\quad 已知微分关系及经过点 \((1,1)\) 的条件，求以 \(x\) 为自变量的解 \(y(x)\)。

\textbf{方法与条件}\quad 乘积微分公式给出 \(\mathrm{d}(xy)=y\,\mathrm{d}x+x\,\mathrm{d}y\)，又有 \(\mathrm{d}(y^3)=3y^2\,\mathrm{d}y\)，所以可以整体积分，避免先除以可能为零的 \(y\)。

\textbf{推导}\quad 原式为 \(\mathrm{d}(xy-y^3)=0\)，因此 \(xy-y^3=C\)。用初值计算 \(C=1\cdot1-1^3=0\)，得到
\[
y(x-y^2)=0.
\]
初值处 \(y=1>0\)，连续性保证其附近 \(y\ne0\)，故该分支必须满足 \(y^2=x\)。两个平方根中，只有正根在 \(x=1\) 时等于 \(1\)，于是 \(\boxed{y=\sqrt{x},\ x>0}\)。常值分支 \(y=0\) 不满足初值。

\textbf{结论与检查}\quad \(y'=\frac{1}{2\sqrt{x}}\)，代入得 \(y+(x-3y^2)y'=\sqrt{x}-\frac{2x}{2\sqrt{x}}=0\)。\(y(1)=1\)，且 \(x>0\) 时导数有限；不能把 \(x=0\) 纳入这个可导解的区间。
\end{QuestionSolution}

% Q 2016 11
\QuestionPage{4.1 几类简单的微分方程}{填空题}{2016 年第 11 题}{2}{由特解确定一阶线性方程}{205 / 234}{**}
以$y = x^2 - \mathrm{e}^x$和$y = x^2$为特解的一阶非齐次线性微分方程为 \blank{}．
\QuestionEnd

\begin{QuestionSolution}{2016 年第 11 题}
\textbf{已知与目标}\quad 两个指定函数是同一个一阶非齐次线性方程的解，求这个方程。

\textbf{方法与条件}\quad 把方程写成 \(y'+p(x)y=q(x)\)。将两个解满足的等式相减，右端消去，就能先求 \(p(x)\)。这说明两个非齐次解之差满足右端为零的齐次方程。

\textbf{推导}\quad 两解之差为 \(-\mathrm{e}^x\)，其导数也是 \(-\mathrm{e}^x\)，故
\[
-\mathrm{e}^x-p(x)\mathrm{e}^x=0.
\]
因为 \(\mathrm{e}^x>0\)，得 \(p(x)=-1\)。再将较简单的解 \(y=x^2\) 代入，求得 \(q(x)=2x-x^2\)。所以可写为 \(\boxed{y'-y=2x-x^2}\)。

\textbf{结论与检查}\quad 对另一解，\((x^2-\mathrm{e}^x)'-(x^2-\mathrm{e}^x)=2x-x^2\)，同样成立。右端不恒为零，符合非齐次要求。
\end{QuestionSolution}

% Q 2024 13
\QuestionPage{4.1 几类简单的微分方程}{填空题}{2024 年第 13 题}{2}{一阶微分方程}{206 / 234}{**}
微分方程
\[y'=\frac{1}{(x+y)^2}
\]
满足初始条件 \(y(1)=0\) 的解为 \blank{}
\QuestionEnd

\begin{QuestionSolution}{2024 年第 13 题}
\textbf{已知与目标}\quad 求 \(y'=\frac{1}{(x+y)^2}\) 经过 \((1,0)\) 的解；分母要求 \(x+y\ne0\)。

\textbf{方法与条件}\quad 右端只含组合 \(x+y\)，令 \(u=x+y\)。由 \(u'=1+y'\) 可将方程变成只含 \(u\) 的关系。\(u'>0\) 保证可以局部反解，把 \(x\) 看成 \(u\) 的函数。

\textbf{推导}\quad \(u(1)=1\)，初值所在分支有 \(u>0\)，且 \(u'=1+\frac{1}{u^2}\)。取倒数并分解被积式，得
\[
\frac{\mathrm{d}x}{\mathrm{d}u}=\frac{u^2}{1+u^2},
\qquad x=\int\left(1-\frac{1}{1+u^2}\right)\mathrm{d}u=u-\arctan u+C.
\]
代入 \(x=u=1\)，得 \(1=1-\frac{\pi}{4}+C\)，所以 \(C=\frac{\pi}{4}\)。还原 \(u=x+y\)，得隐式解（用 \(x,y\) 共同满足的等式表示的解）
\[
\boxed{y=\arctan(x+y)-\frac{\pi}{4}},\qquad x+y>0.
\]

\textbf{结论与检查}\quad 对答案求导，\(y'=\frac{1+y'}{1+(x+y)^2}\)，移项得 \((x+y)^2y'=1\)，正是原式。\((1,0)\) 满足答案。由 \(u-\arctan u\) 在 \(u>0\) 上严格增加、值域为 \((0,\infty)\)，此分支的 \(x\) 区间是 \((\frac{\pi}{4},\infty)\)。
\end{QuestionSolution}

% Q 2025 15
\QuestionPage{4.1 几类简单的微分方程}{填空题}{2025 年第 15 题}{2}{一阶微分方程初值问题}{207 / 234}{***}
微分方程 \(\left( 2y - 3x \right) \mathrm{d} x + \left( 2x - 5y \right) \mathrm{d} y = 0\) 满足条件 \(y\left( 1 \right) = 1\) 的解为 \blank{}.
\QuestionEnd

\begin{QuestionSolution}{2025 年第 15 题}
\textbf{已知与目标}\quad 已知微分关系及 \(y(1)=1\)，求经过该点的函数分支。

\textbf{方法与条件}\quad \(2y\,\mathrm{d}x+2x\,\mathrm{d}y=\mathrm{d}(2xy)\)，其余两项也是平方函数的微分，所以先整体积分，再由初值选择二次方程的根。

\textbf{推导}\quad 原式可写为 \(\mathrm{d}(2xy-\frac{3}{2}x^2-\frac{5}{2}y^2)=0\)，故括号等于常数。代入 \((1,1)\)，得常数为 \(2-\frac{3}{2}-\frac{5}{2}=-2\)。两边乘 \(-2\)，整理为
\[
5y^2-4xy+3x^2-4=0.
\]
把它看作关于 \(y\) 的二次方程，判别式为 \(16x^2-20(3x^2-4)=4(20-11x^2)\)，所以
\[
y=\frac{4x\pm2\sqrt{20-11x^2}}{10}
=\frac{2x\pm\sqrt{20-11x^2}}{5}.
\]
在 \(x=1\) 时，正根为 \(1\)，负根为 \(-\frac{1}{5}\)，因此
\[
\boxed{y=\frac{2x+\sqrt{20-11x^2}}{5}},\qquad |x|<\sqrt{\frac{20}{11}}.
\]

\textbf{结论与检查}\quad 对隐式关系求导得 \((6x-4y)+(10y-4x)y'=0\)，除以 \(-2\) 就是原方程。所取分支有 \(2x-5y=-\sqrt{20-11x^2}\ne0\)，故在上述开区间内可导；初值亦满足。
\end{QuestionSolution}

% Q 2009 18
\QuestionPage{4.1 几类简单的微分方程}{解答题}{2009 年第 18 题}{3}{降阶方程与旋转体积}{208 / 234}{***}
设非负函数$y = y(x)$（$x \ge 0$）满足微分方程$xy'' - y' + 2 = 0$，
当曲线$y = y(x)$过原点时，其与直线$x = 1$及$y = 0$围成平面区域$D$的面积为$2$，
求$D$绕$y$轴旋转所得旋转体体积．
\QuestionEnd

\begin{QuestionSolution}{2009 年第 18 题}
\textbf{已知与目标}\quad 曲线非负、过原点、满足二阶方程，且在 \([0,1]\) 下方区域面积为 \(2\)。求该区域绕 \(y\) 轴的旋转体积。

\textbf{方法与条件}\quad 方程不直接含 \(y\)，令 \(p=y'\) 可降低一阶。确定曲线后，用圆柱壳法：宽为 \(\mathrm{d}x\) 的竖条绕 \(y\) 轴旋转，薄壳体积为周长 \(2\pi x\) 乘高 \(y(x)\) 乘厚度 \(\mathrm{d}x\)。

\textbf{推导}\quad 在 \(x>0\) 时，\(xp'-p+2=0\)。因为
\[
\left(\frac{p-2}{x}\right)'=\frac{xp'-p+2}{x^2}=0,
\]
所以 \(p=2+Cx\)。再积分得 \(y=2x+\frac{C}{2}x^2+D\)。用 \(y(0)=0\) 得 \(D=0\)，把 \(\frac{C}{2}\) 记为 \(a\)。由非负性，面积直接为函数积分：
\[
2=\int_0^1(2x+ax^2)\,\mathrm{d}x
=\left[x^2+\frac{a}{3}x^3\right]_0^1=1+\frac{a}{3},
\]
得 \(a=3\)，即 \(y=2x+3x^2\)。于是
\[
\begin{aligned}
V&=2\pi\int_0^1x(2x+3x^2)\,\mathrm{d}x\\
&=2\pi\left[\frac{2}{3}x^3+\frac{3}{4}x^4\right]_0^1
=\boxed{\frac{17\pi}{6}}.
\end{aligned}
\]

\textbf{结论与检查}\quad \(y'=2+6x,y''=6\)，所以 \(xy''-y'+2=0\)，曲线在 \(x\ge0\) 非负且过原点。体积为正；题目未给具体长度单位，结果以相应的立方单位计。
\end{QuestionSolution}

% Q 2011 18
\QuestionPage{4.1 几类简单的微分方程}{解答题}{2011 年第 18 题}{3}{由切线倾角建立微分方程}{209 / 234}{***}
设函数$y(x)$具有二阶导数，且曲线$l:y = y(x)$与直线$y = x$相切于原点，
记$\alpha$为曲线$l$在点$(x,y)$处切线的倾角，若$\frac{\mathrm{d}\alpha}{\mathrm{d} x} = \frac{\mathrm{d} y}{\mathrm{d} x}$，求$y(x)$的表达式．
\QuestionEnd

\begin{QuestionSolution}{2011 年第 18 题}
\textbf{已知与目标}\quad 曲线在原点与 \(y=x\) 相切，故 \(y(0)=0,y'(0)=1\)。还已知切线倾角的变化率 \(\alpha'\) 等于 \(y'\)，要求曲线表达式。

\textbf{方法与条件}\quad 倾角是切线与 \(x\) 轴正方向所成的角，斜率为 \(\tan\alpha\)。在初值点附近取连续的角度分支；\(y'(0)=1\) 对应 \(\alpha(0)=\frac{\pi}{4}\)。用正弦的导数可避免直接积分正切的倒数。

\textbf{推导}\quad 由 \(\alpha'=y'\) 积分，得 \(\alpha-y=C\)，用初值得 \(C=\frac{\pi}{4}\)。又有 \(\alpha'=y'=\tan\alpha\)，所以
\[
\frac{\mathrm{d}}{\mathrm{d}x}\sin\alpha
=\cos\alpha\,\alpha'=\sin\alpha.
\]
该方程的解为 \(\sin\alpha=K\mathrm{e}^x\)。在 \(x=0\) 时 \(K=\sin\frac{\pi}{4}=\frac{1}{\sqrt{2}}\)。初值附近 \(0<\alpha<\frac{\pi}{2}\)，因此取反正弦的主值，得到
\[
\boxed{y(x)=\arcsin\left(\frac{\mathrm{e}^x}{\sqrt{2}}\right)-\frac{\pi}{4}},
\qquad x<\frac{1}{2}\ln2.
\]

\textbf{结论与检查}\quad \(y'=\frac{\mathrm{e}^x}{\sqrt{2-\mathrm{e}^{2x}}}\)，所以 \(y(0)=0,y'(0)=1\)。令 \(\alpha=y+\frac{\pi}{4}\)，则 \(\tan\alpha=y'\)、\(\alpha'=y'\) 均成立。区间右端的根号变为零，导数不再有限，故必须取严格不等号。
\end{QuestionSolution}

% Q 2014 16
\QuestionPage{4.1 几类简单的微分方程}{解答题}{2014 年第 16 题}{2}{可分离变量方程与极值}{210 / 234}{***}
已知函数$y = y(x)$满足微分方程$x^2 + y^2 y' = 1 - y'$，且$y(2) = 0$，求$y(x)$的极大值和极小值．
\QuestionEnd

\begin{QuestionSolution}{2014 年第 16 题}
\textbf{已知与目标}\quad 已知微分方程及 \(y(2)=0\)，求解函数的极大值和极小值，即在某个邻域中比附近值大或小的函数值。

\textbf{方法与条件}\quad 先用导数符号判断极值位置，再求这些位置的函数值。因为 \(1+y^2>0\)，导数符号完全由 \(1-x^2\) 决定，不必显式解出三次方程。

\textbf{推导}\quad 整理原式得 \(y'=\frac{1-x^2}{1+y^2}\)。因此 \(x<-1\) 时导数为负，\(-1<x<1\) 时为正，\(x>1\) 时为负。故 \(x=-1\) 为极小值点，\(x=1\) 为极大值点。

分离变量 \((1+y^2)\,\mathrm{d}y=(1-x^2)\,\mathrm{d}x\)，积分得
\[
y+\frac{y^3}{3}=x-\frac{x^3}{3}+C .
\]
代入初值得 \(0=2-\frac{8}{3}+C\)，所以 \(C=\frac{2}{3}\)。在 \(x=-1\) 时右端为 \(-1+\frac{1}{3}+\frac{2}{3}=0\)，故 \(y=0\)；在 \(x=1\) 时右端为 \(\frac{4}{3}\)，而 \(1+\frac{1}{3}=\frac{4}{3}\)，故 \(y=1\)。

\textbf{结论与检查}\quad 极大值为 \(\boxed{1}\)，极小值为 \(\boxed{0}\)。函数 \(H(y)=y+\frac{y^3}{3}\) 满足 \(H'(y)=1+y^2>0\)，且值域为全部实数，因此上述两个 \(y\) 值均唯一，没有遗漏其他根。导数确实在两点发生相应变号。
\end{QuestionSolution}

% Q 2015 20
\QuestionPage{4.1 几类简单的微分方程}{解答题}{2015 年第 20 题}{3}{Newton 冷却模型}{211 / 234}{***}
已知高温物体置于低温介质中，任一时刻该物体温度对时间的变化率与该时刻物体和介质的温差成正比，
现将一初始温度为$120^{\circ}C$的物体在$20^{\circ}C$的恒温介质中冷却，$30$min后该物体降至$30^{\circ}C$，
若要将该物体的温度继续降至$21^{\circ}C$，还需冷却多长时间？
\QuestionEnd

\begin{QuestionSolution}{2015 年第 20 题}
\textbf{已知与目标}\quad 环境温度恒为 \(20^\circ\mathrm{C}\)，物体在 \(30\) 分钟内由 \(120^\circ\mathrm{C}\) 降到 \(30^\circ\mathrm{C}\)。求从 \(30^\circ\mathrm{C}\) 继续降到 \(21^\circ\mathrm{C}\) 所需时间。

\textbf{方法与条件}\quad 按题给的冷却规律，变化率正比于温差。设 \(t\) 的单位为分钟，温差 \(u(t)=T(t)-20\) 的单位为摄氏度。冷却时导数为负，故 \(u'=-ku\)，其中 \(k>0\)，单位为 \(\mathrm{min}^{-1}\)。

\textbf{推导}\quad 分离变量并积分：\(\frac{\mathrm{d}u}{u}=-k\,\mathrm{d}t\)，得 \(u=A\mathrm{e}^{-kt}\)。由 \(u(0)=100\)，得 \(u(t)=100\mathrm{e}^{-kt}\)。用 \(u(30)=10\)，得
\[
\mathrm{e}^{-30k}=\frac{1}{10},\qquad k=\frac{\ln10}{30}\ \mathrm{min}^{-1}.
\]
目标温差为 \(21-20=1^\circ\mathrm{C}\)。设还需 \(\Delta t\)，则 \(10\mathrm{e}^{-k\Delta t}=1\)，即温差也要缩小到此前的 \(\frac{1}{10}\)，所以 \(\Delta t=\frac{\ln10}{k}=30\) 分钟。

\textbf{结论与检查}\quad 还需 \(\boxed{30\text{ 分钟}}\)。总时长 \(60\) 分钟时，\(T(60)=20+100(\frac{1}{10})^2=21^\circ\mathrm{C}\)。温度始终高于环境温度且不断降低，与冷却过程一致。
\end{QuestionSolution}

% Q 2019 17
\QuestionPage{4.1 几类简单的微分方程}{解答题}{2019 年第 17 题}{2}{一阶线性方程与旋转体积}{213 / 234}{***}
设函数$y(x)$是微分方程$y'-xy=\frac{1}{2\sqrt{x}}\mathrm{e}^{\frac{x^2}{2}}$满足条件$y(1)=\sqrt{\mathrm{e}}$的特解．
\begin{enumerate}
\item 求$y(x)$的表达式；
\item 设平面区域$D = \{(x,y)|1 \le x \le 2,0 \le y \le y(x)\}$，
      求$D$绕$x$轴旋转一周所形成的旋转体的体积．
\end{enumerate}
\QuestionEnd

\begin{QuestionSolution}{2019 年第 17 题}
\textbf{已知与目标}\quad 先求一阶方程满足 \(y(1)=\sqrt{\mathrm{e}}\) 的解，再求其在 \([1,2]\) 下方区域绕 \(x\) 轴的旋转体积。第二问要使用第一问得到的函数。

\textbf{方法与条件}\quad （1）一阶线性方程用积分因子 \(\mathrm{e}^{\int -x\,\mathrm{d}x}=\mathrm{e}^{-\frac{x^2}{2}}\)，使左端成为乘积的导数。原式含 \(\sqrt{x}\) 作分母，故在 \(x>0\) 上求解。（2）绕 \(x\) 轴时垂直于该轴的截面为圆盘，面积为 \(\pi y^2\)。

\textbf{推导}\quad （1）两边乘积分因子，由乘积求导公式得
\[
(y\mathrm{e}^{-\frac{x^2}{2}})'=\frac{1}{2\sqrt{x}},
\qquad y\mathrm{e}^{-\frac{x^2}{2}}=\sqrt{x}+C.
\]
代入初值，左端为 \(\sqrt{\mathrm{e}}\,\mathrm{e}^{-\frac{1}{2}}=1\)，故 \(1=1+C\)，得 \(C=0\)。于是 \(\boxed{y=\sqrt{x}\,\mathrm{e}^{\frac{x^2}{2}},\ x>0}\)。

（2）\(y(x)>0\)，所以直接用圆盘面积积分。令 \(u=x^2\)，则 \(\mathrm{d}u=2x\,\mathrm{d}x\)，两端变为 \(1,4\)，因此
\[
V=\pi\int_1^2 y(x)^2\,\mathrm{d}x
=\pi\int_1^2 x\mathrm{e}^{x^2}\,\mathrm{d}x
=\frac{\pi}{2}\int_1^4\mathrm{e}^u\,\mathrm{d}u
=\boxed{\frac{\pi}{2}(\mathrm{e}^4-\mathrm{e})}.
\]

\textbf{结论与检查}\quad 对第一问的解求导，\(y'=\frac{\mathrm{e}^{\frac{x^2}{2}}}{2\sqrt{x}}+xy\)，与原方程一致，初值也成立。第二问体积为正；题目没有指定长度单位，按相应的立方单位计。
\end{QuestionSolution}

% Q 2021 20
\QuestionPage{4.1 几类简单的微分方程}{解答题}{2021 年第 20 题}{3}{一阶方程与距离最值}{215 / 234}{***}
设函数 \(y=y(x)\)（\(x>0\)）满足微分方程
\[xy'-6y=-6,
\]
且 \(y\left(\sqrt{3}\right)=10\)。
\begin{enumerate}
\item 求 \(y(x)\)；
\item \(P\) 为曲线 \(y=y(x)\) 上的一点，曲线在点 \(P\) 的法线在 \(y\) 轴上的截距为 \(I_y\)。为使 \(I_y\) 最小，求 \(P\) 的坐标。
\end{enumerate}
\QuestionEnd

\begin{QuestionSolution}{2021 年第 20 题}
\textbf{已知与目标}\quad 先由方程和 \(y(\sqrt{3})=10\) 求曲线，再确定法线纵截距最小时的点。截距是直线在 \(y\) 轴上的交点纵坐标。

\textbf{方法与条件}\quad 将 \(y-1\) 作为整体可化简方程；法线垂直于切线，切线斜率非零时，两斜率乘积为 \(-1\)。最后在题给的 \(x>0\) 上研究截距函数的增减。

\textbf{推导}\quad （1）原式为 \(xy'-6(y-1)=0\)。由于
\[
\bigl(x^{-6}(y-1)\bigr)'=x^{-7}\bigl(xy'-6(y-1)\bigr)=0,
\]
可得 \(y=1+Cx^6\)。初值给出 \(10=1+27C\)，故 \(C=\frac{1}{3}\)，即 \(\boxed{y=1+\frac{x^6}{3}}\)。

（2）设所求点横坐标为 \(x>0\)。切线斜率为 \(y'=2x^5\ne0\)，法线方程为 \(Y-y=-\frac{1}{2x^5}(X-x)\)。令 \(X=0\)，得
\[
I_y(x)=y+\frac{1}{2x^4}=1+\frac{x^6}{3}+\frac{1}{2x^4},
\qquad I_y'(x)=2x^5-\frac{2}{x^5}=\frac{2(x^{10}-1)}{x^5}.
\]
分母为正，故 \(0<x<1\) 时递减，\(x>1\) 时递增。全区间最小值在 \(x=1\) 取得，对应 \(\boxed{P=(1,\frac{4}{3})}\)。

\textbf{结论与检查}\quad 解满足 \(xy'-6y=-6\) 及给定初值。最优点的法线为 \(Y-\frac{4}{3}=-\frac{1}{2}(X-1)\)，纵截距为 \(\frac{11}{6}\)。两端 \(x\to0^+\) 和 \(x\to\infty\) 时截距均趋于无穷，内部最小值没有遗漏边界。
\end{QuestionSolution}

% Q 2022 18
\QuestionPage{4.1 几类简单的微分方程}{解答题}{2022 年第 18 题}{3}{一阶方程与曲线弧长}{216 / 234}{***}
设函数 \(y\left( x \right)\) 是微分方程
\(2x y^{\prime} - 4y = 2\ln x - 1\)
的满足条件
\(y\left( 1 \right) = \frac{1}{4}\)
的解，求曲线 \(y = y\left( x \right)\) 在 \(\left[ 1,\mathrm{e} \right]\) 上的弧长.
\QuestionEnd

\begin{QuestionSolution}{2022 年第 18 题}
\textbf{已知与目标}\quad 曲线满足一阶方程及 \(y(1)=\frac{1}{4}\)，求 \([1,\mathrm{e}]\) 上的弧长。必须先确定导数 \(y'\)。

\textbf{方法与条件}\quad 原式含 \(\ln x\)，在 \(x>0\) 上计算。用一个非齐次特解加上齐次通解求函数，再用弧长公式 \(s=\int_1^{\mathrm{e}}\sqrt{1+(y')^2}\,\mathrm{d}x\)。该公式适用于导数连续的曲线。

\textbf{推导}\quad 取 \(y_p=-\frac{1}{2}\ln x\)，其导数为 \(-\frac{1}{2x}\)，故 \(2xy_p'-4y_p=-1+2\ln x\)，恰为题给右端。对于齐次部分，\(2xy'-4y=0\) 等价于 \((\frac{y}{x^2})'=0\)，得 \(y_h=Cx^2\)。因此 \(y=Cx^2-\frac{1}{2}\ln x\)。由初值得 \(C=\frac{1}{4}\)，于是
\[
y(x)=\frac{x^2}{4}-\frac{1}{2}\ln x,\qquad
y'(x)=\frac{x}{2}-\frac{1}{2x}.
\]
注意 \(1+(y')^2=\frac{x^2+2+x^{-2}}{4}=(\frac{x+x^{-1}}{2})^2\)。由于 \(x>0\)，括号为正，开方后不用绝对值，得到
\[
s=\int_1^{\mathrm{e}}\left(\frac{x}{2}+\frac{1}{2x}\right)\mathrm{d}x
=\left[\frac{x^2}{4}+\frac{1}{2}\ln x\right]_1^{\mathrm{e}}
=\boxed{\frac{\mathrm{e}^2+1}{4}}.
\]

\textbf{结论与检查}\quad 代回可得 \(2xy'-4y=(x^2-1)-(x^2-2\ln x)=2\ln x-1\)，初值成立。弧长约为 \(2.097\)，大于水平跨度 \(\mathrm{e}-1\approx1.718\)，符合几何要求。未指定长度单位，按相应长度单位计。
\end{QuestionSolution}

% Q 2023 17
\QuestionPage{4.1 几类简单的微分方程}{解答题}{2023 年第 17 题}{3}{由切线截距建立一阶方程}{217 / 234}{***}
设曲线 \(\boldsymbol L:y=y(x)\)（\(x>\mathrm{e}\)）经过点 \((\mathrm{e}^2,0)\)，且 \(\boldsymbol L\) 上任一点 \(P(x,y)\) 到 \(y\) 轴的距离等于该点处切线在 \(y\) 轴上的截距。
\begin{enumerate}
\item 求 \(y(x)\)；
\item 在 \(\boldsymbol L\) 上求一点，使该点处的切线与两坐标轴所围三角形的面积最小，并求此最小面积。
\end{enumerate}
\QuestionEnd

\begin{QuestionSolution}{2023 年第 17 题}
\textbf{已知与目标}\quad 在 \(x>\mathrm{e}\) 上，曲线经过 \((\mathrm{e}^2,0)\)，且点到 \(y\) 轴的距离等于切线纵截距。先求曲线，再求切线围成三角形的最小面积及切点。

\textbf{方法与条件}\quad 将切线的截距写成 \(y,y'\) 的关系建立微分方程。两轴截距确定后，三角形面积为它们乘积绝对值的一半；先判断符号再使用公式。

\textbf{推导}\quad （1）切线为 \(Y-y=y'(X-x)\)，令 \(X=0\) 得纵截距 \(y-xy'\)。因 \(x>0\)，点到 \(y\) 轴的距离为 \(x\)，故 \(y-xy'=x\)。由商的求导公式，
\[
\left(\frac{y}{x}\right)'=\frac{xy'-y}{x^2}=-\frac{1}{x},
\qquad \frac{y}{x}=-\ln x+C.
\]
代入 \((\mathrm{e}^2,0)\) 得 \(C=2\)，即 \(\boxed{y=x(2-\ln x)}\)。

（2）\(y'=1-\ln x<0\)，切线纵截距为 \(x>0\)，横截距为 \(-\frac{x}{1-\ln x}=\frac{x}{\ln x-1}>0\)。所以
\[
S(x)=\frac{x^2}{2(\ln x-1)},\qquad
S'(x)=\frac{x(2\ln x-3)}{2(\ln x-1)^2}.
\]
分母为正，导数在 \(\mathrm{e}<x<\mathrm{e}^{\frac{3}{2}}\) 时为负，在 \(x>\mathrm{e}^{\frac{3}{2}}\) 时为正。因此所求点为
\[
\boxed{\left(\mathrm{e}^{\frac{3}{2}},\frac{\mathrm{e}^{\frac{3}{2}}}{2}\right)},
\qquad S_{\min}=\frac{\mathrm{e}^3}{2(\frac{3}{2}-1)}=\boxed{\mathrm{e}^3}.
\]

\textbf{结论与检查}\quad 曲线在 \(x=\mathrm{e}^2\) 时为零，且 \(y-xy'=x\) 恒成立。最优横坐标大于 \(\mathrm{e}\)；面积在左右端均趋于无穷，所得最小值是全区间最小值。
\end{QuestionSolution}

% Q 2017 21
\QuestionPage{4.1 几类简单的微分方程}{解答题}{2017 年第 21 题}{4}{由切线和法线条件建立方程}{212 / 234}{****}
设$y(x)$是区间$\left(0,\frac{3}{2} \right)$内的可导函数，且$y(1) = 0$．
点$P$是曲线$L:y = y(x)$上的任意一点．$L$在$P$处的切线与$y$轴相交于点$\left(0,Y_p\right)$，
法线与$x$轴相交于点$\left(X_p,0 \right)$，若$X_p = Y_p$，求$L$上点的坐标$(x,y)$满足的方程．
\QuestionEnd

\begin{QuestionSolution}{2017 年第 21 题}
\textbf{已知与目标}\quad 曲线在 \(0<x<\frac{3}{2}\) 可导，经过 \((1,0)\)。每点的切线纵截距等于法线横截距，求曲线的隐式方程。

\textbf{方法与条件}\quad 切线斜率为 \(y'\)，法线是与切线垂直的直线。用方向向量表示法线，可以同时处理 \(y'=0\) 时的竖直法线；得到方程后，因右端只含 \(\frac{y}{x}\)，令 \(u=\frac{y}{x}\) 化为可分离变量方程。

\textbf{推导}\quad 切线 \(Y-y=y'(X-x)\) 给出 \(Y_p=y-xy'\)。切线的方向向量是 \((1,y')\)，与之垂直的向量可取 \((-y',1)\)，故法线写为 \((X,Y)=(x,y)+s(-y',1)\)。令 \(Y=0\) 得 \(s=-y\)，进而 \(X_p=x+yy'\)。令两截距相等，得
\[
(x+y)y'=y-x.
\]
若 \(x+y=0\)，则右端为 \(-2x\ne0\)，产生矛盾，所以可以除以 \(x+y\)。由 \(y=xu\)、\(y'=u+xu'\)，得
\[
u+xu'=\frac{u-1}{u+1},\qquad
xu'=-\frac{1+u^2}{1+u}.
\]
移项积分（\(1+u^2\) 恒正）得
\[
\int\frac{1+u}{1+u^2}\,\mathrm{d}u=-\int\frac{1}{x}\,\mathrm{d}x,
\qquad \arctan u+\frac{1}{2}\ln(1+u^2)=-\ln x+C.
\]
初值 \((x,u)=(1,0)\) 给出 \(C=0\)。还原后，利用 \(x>0\) 合并对数，得
\[
\boxed{\arctan\left(\frac{y}{x}\right)+\frac{1}{2}\ln(x^2+y^2)=0},
\qquad 0<x<\frac{3}{2}.
\]

\textbf{结论与检查}\quad 取经过 \((1,0)\) 的连续分支，此分支有 \(x+y>0\)。对答案求导，左边导数为 \(\frac{xy'-y+x+yy'}{x^2+y^2}\)，令它为零便恢复截距关系；初值也使左边等于零。全过程未除以 \(y'\)，所以没有遗漏水平切线的情形。
\end{QuestionSolution}

% Q 2020 21
\QuestionPage{4.1 几类简单的微分方程}{解答题}{2020 年第 21 题}{4}{由面积比建立微分方程}{214 / 234}{****}
设函数 \(f\left( x \right)\) 可导，且 \(f^{\prime}\left( x \right)>0\)，曲线 \(y=f\left( x \right)\left( x\geq 0 \right)\) 经过坐标原点 \(O\)。\(M\) 为其上任意一点，点 \(M\) 处的切线与 \(x\) 轴交于点 \(T\)，又 \(MP\) 垂直 \(x\) 轴于点 \(P\)。已知由曲线 \(y=f\left( x \right)\)、直线 \(MP\) 以及 \(x\) 轴所围图形的面积与 \(\triangle MTP\) 的面积之比恒为 \(3:2\)，求满足上述条件的曲线方程。
\QuestionEnd

\begin{QuestionSolution}{2020 年第 21 题}
\textbf{已知与目标}\quad 曲线过原点、在正半轴严格上升，曲边图形与切线三角形的面积比为 \(3:2\)，求曲线。面积比只在 \(M\ne O\) 的非退化图形上有意义。

\textbf{方法与条件}\quad 设 \(F(x)=\int_0^x f(t)\,\mathrm{d}t\) 表示曲边面积。由 \(f(0)=0,f'>0\) 得 \(x>0\) 时 \(f(x)>0,F(x)>0\)，可安全作后面的除法。利用 \(F'=f\) 把面积关系化成可分离变量的方程。

\textbf{推导}\quad 在 \(M=(x,f(x))\) 的切线方程中令纵坐标为零，得 \(T\) 的横坐标 \(x-\frac{f(x)}{f'(x)}\)。因此底边 \(TP=\frac{f(x)}{f'(x)}\)，高为 \(f(x)\)，三角形面积为 \(\frac{f(x)^2}{2f'(x)}\)。由面积比，
\[
\frac{F(x)}{\frac{f(x)^2}{2f'(x)}}=\frac{3}{2},
\qquad F(x)f'(x)=\frac{3}{4}f(x)^2.
\]
这里不需要对 \(f'\) 再求导。直接计算
\[
\left(\frac{f}{F^{\frac{3}{4}}}\right)'
=F^{-\frac{7}{4}}\left(Ff'-\frac{3}{4}fF'\right)=0.
\]
故 \(f=kF^{\frac{3}{4}}\)，且 \(k>0\)。由 \(F'=f\)，得 \((F^{\frac{1}{4}})'=\frac{1}{4}F^{-\frac{3}{4}}F'=\frac{k}{4}\)。积分并令 \(x\to0^+\)，利用 \(F(0)=0\) 得 \(F^{\frac{1}{4}}=\frac{kx}{4}\)。因此
\[
F(x)=\left(\frac{kx}{4}\right)^4,
\qquad f(x)=\frac{k^4}{64}x^3=Cx^3,\quad C>0.
\]

\textbf{结论与检查}\quad 所求曲线为 \(\boxed{y=Cx^3,\ C>0,\ x\ge0}\)。曲边面积为 \(\frac{Cx^4}{4}\)，三角形面积为 \(\frac{Cx^4}{6}\)，其比确为 \(3:2\)；\(x>0\) 时导数为 \(3Cx^2>0\)。此解在原点的右导数为 \(0\)。因此这里将严格导数条件理解为针对 \(x>0\) 的非退化点；若还要求原点右导数严格大于零，则与面积比推得的结果矛盾，在该额外要求下无解。
\end{QuestionSolution}

% Q 2026 21
\QuestionPage{4.1 几类简单的微分方程}{解答题}{2026 年第 21 题}{2}{二阶微分方程降阶}{218 / 234}{****}
（本题满分 12 分）
求微分方程
\[x^2y''-2xy'-(y')^2=0\qquad (x>2)
\]
满足条件
\[y\big|_{x=3}=\frac12,\qquad y'\big|_{x=3}=-9
\]
的解.
\QuestionEnd

\begin{QuestionSolution}{2026 年第 21 题}
\textbf{已知与目标}\quad 在 \(x>2\) 上求二阶方程满足 \(y(3)=\frac{1}{2}\)、\(y'(3)=-9\) 的解。

\textbf{方法与条件}\quad 原式没有直接出现 \(y\)，先令 \(p=y'\) 降阶。所得式含 \(p^2\)，再令 \(q=\frac{1}{p}\) 化成一阶线性方程。初值处 \(p\ne0\)，这个代换在其附近合法；最后须检查所得解在整个 \(x>2\) 上均无零分母。

\textbf{推导}\quad 降阶得 \(x^2p'-2xp-p^2=0\)。由 \(q'=-\frac{p'}{p^2}\)，两边除以 \(p^2\) 并移项得
\[
q'+\frac{2}{x}q=-\frac{1}{x^2},\qquad
(x^2q)'=-1.
\]
第二式由两边乘积分因子 \(x^2\) 得到。积分得 \(x^2q=C-x\)；代入 \(q(3)=-\frac{1}{9}\)，有 \(-1=C-3\)，故 \(C=2\)。取倒数并作多项式除法：
\[
y'=p=\frac{x^2}{2-x}=-\frac{x^2}{x-2}=-x-2-\frac{4}{x-2}.
\]
这里用了 \(x^2=(x-2)(x+2)+4\)。逐项积分，且 \(x-2>0\)，得到 \(y=-\frac{x^2}{2}-2x-4\ln(x-2)+C_1\)。由 \(\frac{1}{2}=-\frac{9}{2}-6+C_1\)，得 \(C_1=11\)。因此
\[
\boxed{y=11-\frac{x^2}{2}-2x-4\ln(x-2)},\qquad x>2 .
\]

\textbf{结论与检查}\quad \(y'=-\frac{x^2}{x-2}\ne0\)，\(y''=-\frac{x(x-4)}{(x-2)^2}\)。代回方程的分子为 \(-x^3(x-4)+2x^3(x-2)-x^4=0\)。在 \(x=3\) 时分别得到 \(y=\frac{1}{2}\)、\(y'=-9\)。常值解虽满足原方程，但不满足导数初值，故除以 \(p^2\) 没有遗漏所求解。
\end{QuestionSolution}

\SectionHeading{4.2}{常微分方程}{高阶线性微分方程}{定理与要点\quad 教材 PDF 第 284 页起\hfill 本节 16 题}
\pdfbookmark[2]{4.2 高阶线性微分方程}{section-4.2}
\Theory{定理 2.1}{287}{齐次解的叠合性}{线性齐次方程的解作任意常系数线性组合，仍是解。}
\Theory{定理 2.2}{290}{基本解组与通解}{\(n\) 阶线性齐次方程若有 \(n\) 个线性无关的解，则它们的任意常系数组合构成通解。}
\Theory{定理 2.3}{290}{解的线性无关判别法}{\(n\) 个齐次解线性无关的充要条件是其 Wronski 行列式在区间内某点非零。}
\Theory{定理 2.4}{292}{线性齐次方程的通解结构}{任何一个解都能由 \(n\) 个线性无关的齐次解作常系数线性组合表示。}
\Theory{定理 2.5}{293}{线性非齐次方程解的结构}{任一非齐次特解加上对应齐次通解，即得非齐次通解。}
\Theory{定理 2.6}{294}{不同右端项的叠合}{右端项为 \(F_1\)、\(F_2\) 的两个线性方程特解相加，可得到右端项为 \(F_1+F_2\) 的特解。}
\Theory{方法}{294--307}{常系数方程}{齐次方程按特征根构造基本解。非齐次待定系数法遇到重合根时，按重数乘相应的自变量幂。}
% Q 2011 4
\QuestionPage{4.2 高阶线性微分方程}{选择题}{2011 年第 4 题}{1}{非齐次方程的特解形式}{219 / 234}{**}
微分方程$y'' - \lambda^2 y = \mathrm{e}^{\lambda x} + \mathrm{e}^{- \lambda x}$（$\lambda > 0$）的特解形式为\pickout{}
\begin{abcd}
\item $a(\mathrm{e}^{\lambda x} + \mathrm{e}^{- \lambda x})$．
\item $ax(\mathrm{e}^{\lambda x} + \mathrm{e}^{- \lambda x})$．
\item $x(a\mathrm{e}^{\lambda x} + b\mathrm{e}^{- \lambda x})$．
\item $x^2(a\mathrm{e}^{\lambda x} + b\mathrm{e}^{- \lambda x})$．
\end{abcd}
\QuestionEnd

\begin{QuestionSolution}{2011 年第 4 题}
\textbf{已知与目标}\quad 方程右端有两个指数项，\(\lambda>0\)。要求能包含一个特解的试探形式，系数可由代入确定。

\textbf{方法与条件}\quad 待定系数法先按右端函数的类型设特解。特征方程来自将 \(y=\mathrm{e}^{rx}\) 代入齐次方程：若右端指数的指数系数恰为一个单特征根，试探项须乘 \(x\)，以免代入后左端恒为零。

\textbf{推导}\quad 特征方程为 \(r^2-\lambda^2=0\)，根为 \(\lambda,-\lambda\)。因 \(\lambda>0\)，两根不同，均为单根。因此设
\[
y^*=ax\mathrm{e}^{\lambda x}+bx\mathrm{e}^{-\lambda x}.
\]
例如 \((x\mathrm{e}^{\lambda x})'=(1+\lambda x)\mathrm{e}^{\lambda x}\)，二阶导数为 \((2\lambda+\lambda^2x)\mathrm{e}^{\lambda x}\)，扣去 \(\lambda^2x\mathrm{e}^{\lambda x}\) 后仅余 \(2\lambda\mathrm{e}^{\lambda x}\)。同理负指数项给出 \(-2\lambda\mathrm{e}^{-\lambda x}\)。故代入得到
\[
2\lambda a\mathrm{e}^{\lambda x}-2\lambda b\mathrm{e}^{-\lambda x}
=\mathrm{e}^{\lambda x}+\mathrm{e}^{-\lambda x},
\quad a=\frac{1}{2\lambda},\ b=-\frac{1}{2\lambda}.
\]

\textbf{结论与检查}\quad 选 C。两个系数需分别确定，实际一正一负，B 中共用同一个系数无法满足两项。以上代入也直接验证了该形式可行。
\end{QuestionSolution}

% Q 2019 4
\QuestionPage{4.2 高阶线性微分方程}{选择题}{2019 年第 4 题}{1}{由通解确定方程系数}{221 / 234}{**}
已知微分方程$y'' + ay' + by = c\mathrm{e}^x$的通解为$y = (C_1 + C_2 x) \mathrm{e}^{- x} + \mathrm{e}^x$，则$a,b,c$依次为\pickout{}
\begin{abcd}
\item $1,0,1$．
\item $1,0,2$．
\item $2,1,3$．
\item $2,1,4$．
\end{abcd}
\QuestionEnd

\begin{QuestionSolution}{2019 年第 4 题}
\textbf{已知与目标}\quad 给出方程通解，求常数 \(a,b,c\)。通解指包含所有解、任意常数可自由取值的表达式。

\textbf{方法与条件}\quad 非齐次通解等于齐次通解加一个固定特解。含任意常数的部分可反推出齐次方程的特征根；再代入固定项求右端系数。

\textbf{推导}\quad 齐次部分为 \((C_1+C_2x)\mathrm{e}^{-x}\)。出现 \(\mathrm{e}^{-x}\) 和 \(x\mathrm{e}^{-x}\)，说明特征根 \(-1\) 重复两次，因此
\[
r^2+ar+b=(r+1)^2=r^2+2r+1.
\]
比较同次幂系数得 \(a=2,b=1\)。固定项 \(y_p=\mathrm{e}^x\) 的一、二阶导数均为 \(\mathrm{e}^x\)，所以 \(y_p''+2y_p'+y_p=4\mathrm{e}^x\)，得 \(c=4\)。

\textbf{结论与检查}\quad \(\boxed{a=2,b=1,c=4}\)，选 D。将两个齐次基解分别代入左端均得零，加上特解后恰得到 \(4\mathrm{e}^x\)，与通解结构一致。
\end{QuestionSolution}

% Q 2026 2
\QuestionPage{4.2 高阶线性微分方程}{选择题}{2026 年第 2 题}{1}{非齐次方程的解的结构}{224 / 234}{**}
设 \(y_1(x),y_2(x)\) 是某二阶非齐次线性微分方程的两个特解.若常数 \(\lambda,\mu\) 使得 \(2\lambda y_1(x)+\mu y_2(x)\) 是该方程的解，\(\lambda y_1(x)-2\mu y_2(x)\) 是该方程对应齐次方程的解，则
\choices{\(\lambda=\frac15,\ \mu=\frac25\)}{\(\lambda=\frac25,\ \mu=\frac15\)}{\(\lambda=\frac14,\ \mu=\frac12\)}{\(\lambda=\frac12,\ \mu=\frac14\)}
\QuestionEnd

\begin{QuestionSolution}{2026 年第 2 题}
\textbf{已知与目标}\quad \(y_1,y_2\) 都是同一非齐次方程的特解，要求两种指定组合分别为非齐次解和齐次解时的 \(\lambda,\mu\)。

\textbf{方法与条件}\quad 把左端的求导、乘系数及相加的运算整体记为 \(L[y]\)，称为线性微分算子。线性表示 \(L[cy+dz]=cL[y]+dL[z]\)，其中 \(c,d\) 为常数。非齐次右端 \(q\) 不恒为零。

\textbf{推导}\quad 已知 \(L[y_1]=L[y_2]=q\)，所以题设给出
\[
L[2\lambda y_1+\mu y_2]=(2\lambda+\mu)q=q,
\qquad L[\lambda y_1-2\mu y_2]=(\lambda-2\mu)q=0.
\]
由于 \(q\not\equiv0\)，可得 \(2\lambda+\mu=1\)、\(\lambda-2\mu=0\)。第二式给出 \(\lambda=2\mu\)，代入第一式得 \(5\mu=1\)，故 \(\lambda=\frac{2}{5},\mu=\frac{1}{5}\)。

\textbf{结论与检查}\quad 选 B。代回有 \(2\lambda+\mu=\frac{4}{5}+\frac{1}{5}=1\)，\(\lambda-2\mu=0\)，分别恢复题目要求的两个右端。
\end{QuestionSolution}

% Q 2010 9
\QuestionPage{4.2 高阶线性微分方程}{填空题}{2010 年第 9 题}{2}{三阶常系数齐次方程}{225 / 234}{**}
$3$阶常系数线性齐次微分方程$y''' - 2y'' + y' - 2y = 0$的通解为$y =$ \blank{}．
\QuestionEnd

\begin{QuestionSolution}{2010 年第 9 题}
\textbf{已知与目标}\quad 求三阶常系数齐次方程的通解，应包含三个独立的任意常数。

\textbf{方法与条件}\quad 令试探解 \(y=\mathrm{e}^{rx}\)，求导后每次多乘一个 \(r\)，便得到特征方程。实根 \(r\) 对应 \(\mathrm{e}^{rx}\)，共轭复根 \(\alpha\pm i\beta\) 对应两个实函数 \(\mathrm{e}^{\alpha x}\cos\beta x\)、\(\mathrm{e}^{\alpha x}\sin\beta x\)。

\textbf{推导}\quad 代入并约去非零的 \(\mathrm{e}^{rx}\)，得到
\[
r^3-2r^2+r-2=r^2(r-2)+(r-2)=(r-2)(r^2+1)=0.
\]
根为 \(2,i,-i\)，互不相同。因此
\[
\boxed{y=C_1\mathrm{e}^{2x}+C_2\cos x+C_3\sin x}.
\]

\textbf{结论与检查}\quad 三个函数分别代入原式均为零。例如余弦项给出 \(\sin x+2\cos x-\sin x-2\cos x=0\)。它们线性无关，且共有三个任意常数，构成三阶方程的全部解。
\end{QuestionSolution}

% Q 2015 12
\QuestionPage{4.2 高阶线性微分方程}{填空题}{2015 年第 12 题}{2}{常系数齐次方程与极值条件}{227 / 234}{**}
设函数$y = y(x)$是微分方程$y'' + y' - 2y = 0$的解，且在$x = 0$处$y(x)$取得极值$3$，则$y(x)$= \blank{}．
\QuestionEnd

\begin{QuestionSolution}{2015 年第 12 题}
\textbf{已知与目标}\quad 解在 \(x=0\) 处取得极值 \(3\)，求该解。这里不仅知道函数值，还知道极值点处的导数信息。

\textbf{方法与条件}\quad 可导函数在内部极值点处导数为零，因此条件转成 \(y(0)=3,y'(0)=0\)。用特征方程求通解，再确定两个常数。导数为零只是极值的必要条件，最后还要检查确有极值。

\textbf{推导}\quad 特征方程为 \(r^2+r-2=(r-1)(r+2)=0\)，得 \(y=C_1\mathrm{e}^x+C_2\mathrm{e}^{-2x}\)，\(y'=C_1\mathrm{e}^x-2C_2\mathrm{e}^{-2x}\)。代入条件：
\[
C_1+C_2=3,\qquad C_1-2C_2=0.
\]
由第二式 \(C_1=2C_2\)，代入第一式得 \(3C_2=3\)，故 \(C_1=2,C_2=1\)。因此 \(\boxed{y=2\mathrm{e}^x+\mathrm{e}^{-2x}}\)。

\textbf{结论与检查}\quad \(y''=2\mathrm{e}^x+4\mathrm{e}^{-2x}>0\)，所以在 \(x=0\) 处确为极小值，且值为 \(3\)。代入 \(y''+y'-2y\) 后两个指数项的系数均为零，原方程成立。
\end{QuestionSolution}

% Q 2021 15
\QuestionPage{4.2 高阶线性微分方程}{填空题}{2021 年第 15 题}{2}{常系数齐次方程的通解}{229 / 234}{**}
微分方程 \(y'''-y=0\) 的通解
\[y= \blank{}
\]
\QuestionEnd

\begin{QuestionSolution}{2021 年第 15 题}
\textbf{已知与目标}\quad 求 \(y'''-y=0\) 的实通解，需保留三个任意常数。

\textbf{方法与条件}\quad 设 \(y=\mathrm{e}^{rx}\) 得特征方程。共轭复根 \(\alpha\pm i\beta\) 成对出现时，可用实函数 \(\mathrm{e}^{\alpha x}\cos\beta x\)、\(\mathrm{e}^{\alpha x}\sin\beta x\) 表示它们对应的解，其中 \(i^2=-1\)。

\textbf{推导}\quad 特征方程为 \(r^3-1=0\)。利用立方差公式，
\[
r^3-1=(r-1)(r^2+r+1),\qquad
r^2+r+1=\left(r+\frac{1}{2}\right)^2+\frac{3}{4}.
\]
故一个实根为 \(1\)，另两根为 \(-\frac{1}{2}\pm\frac{i\sqrt{3}}{2}\)。按上述对应关系，
\[
\boxed{y=C_1\mathrm{e}^x+\mathrm{e}^{-\frac{x}{2}}
\left(C_2\cos\frac{\sqrt3x}{2}+C_3\sin\frac{\sqrt3x}{2}\right)}.
\]

\textbf{结论与检查}\quad 三个特征根不同，给出的三个实基解线性无关，共有三个任意常数。每个指数形式代入原式均因 \(r^3-1=0\) 而成立，其实部和虚部也满足实系数方程，故实通解正确。
\end{QuestionSolution}

% Q 2022 14
\QuestionPage{4.2 高阶线性微分方程}{填空题}{2022 年第 14 题}{2}{三阶常系数齐次方程}{230 / 234}{**}
微分方程
\(y^{\prime\prime\prime} - 2y^{\prime\prime} + 5y^{\prime} = 0\)
的通解为 \(y\left( x \right) = \blank{}\)
\QuestionEnd

\begin{QuestionSolution}{2022 年第 14 题}
\textbf{已知与目标}\quad 求三阶齐次方程的通解。原式不直接含 \(y\)，常数函数一定是其中一类解。

\textbf{方法与条件}\quad 用 \(\mathrm{e}^{rx}\) 代入得到特征方程。零特征根对应 \(\mathrm{e}^{0x}=1\)，复根的实部决定指数因子，虚部决定三角函数频率。

\textbf{推导}\quad 特征方程为
\[
r^3-2r^2+5r=r\bigl((r-1)^2+4\bigr)=0.
\]
所以根为 \(0,1+2i,1-2i\)，其中 \(i^2=-1\)。后两根对应 \(\mathrm{e}^x\cos2x\) 与 \(\mathrm{e}^x\sin2x\)。三个根均为单根，故
\[
\boxed{y=C_1+\mathrm{e}^x(C_2\cos2x+C_3\sin2x)}.
\]

\textbf{结论与检查}\quad 常数项的各阶导数均为零，确实满足原式。另两个基解分别来自特征根，代入亦为零；三个任意常数与三阶方程相符，未遗漏零根对应的常数解。
\end{QuestionSolution}

% Q 2017 4
\QuestionPage{4.2 高阶线性微分方程}{选择题}{2017 年第 4 题}{1}{常系数非齐次方程特解}{220 / 234}{***}
微分方程$y'' - 4y' + 8y = \mathrm{e}^{2x}(1 + \cos 2x)$的特解可设为$y^* =$\pickout{}
\begin{abcd}
\item $A\mathrm{e}^{2x} + \mathrm{e}^{2x}\left(B\cos 2x + C\sin 2x \right)$．
\item $Ax \mathrm{e}^{2x} + \mathrm{e}^{2x}\left(B\cos 2x + C\sin 2x \right)$．
\item $A\mathrm{e}^{2x} + x\mathrm{e}^{2x}\left(B\cos 2x + C\sin 2x \right)$．
\item $Ax \mathrm{e}^{2x} + x\mathrm{e}^{2x}\left(B\cos 2x + C\sin 2x \right)$．
\end{abcd}
\QuestionEnd

\begin{QuestionSolution}{2017 年第 4 题}
\textbf{已知与目标}\quad 右端为 \(\mathrm{e}^{2x}+\mathrm{e}^{2x}\cos2x\)，求两个部分合在一起时应设的特解形式。

\textbf{方法与条件}\quad 先令 \(y=\mathrm{e}^{2x}z\) 消去共同指数。若试探项本来就是齐次解，代入后会得到零，须乘 \(x\) 形成新的试探项。这种右端与齐次解重合的情况常称为共振。

\textbf{推导}\quad 由 \(y'=\mathrm{e}^{2x}(z'+2z)\)、\(y''=\mathrm{e}^{2x}(z''+4z'+4z)\)，代入并约去非零的指数得 \(z''+4z=1+\cos2x\)。常数右端可取 \(z_1=\frac{1}{4}\)，不需乘 \(x\)。齐次方程 \(z''+4z=0\) 的解含 \(\cos2x,\sin2x\)，所以余弦右端需要设 \(z_2=x(B\cos2x+C\sin2x)\)。例如
\[
(x\sin2x)'=\sin2x+2x\cos2x,
\qquad (x\sin2x)''=4\cos2x-4x\sin2x.
\]
因此可取 \(z_2=\frac{x\sin2x}{4}\)。乘回 \(\mathrm{e}^{2x}\)，总特解包含于 \(A\mathrm{e}^{2x}+x\mathrm{e}^{2x}(B\cos2x+C\sin2x)\)。

\textbf{结论与检查}\quad 选 C。取 \(A=\frac{1}{4},B=0,C=\frac{1}{4}\) 即能使右端为题给函数。只有余弦对应的部分需要乘 \(x\)，不能让两部分一律乘 \(x\)。
\end{QuestionSolution}

% Q 2023 4
\QuestionPage{4.2 高阶线性微分方程}{选择题}{2023 年第 4 题}{1}{常系数方程解的有界性}{222 / 234}{***}
若微分方程
\[y''+ay'+by=0
\]
的解在 \((-\infty,+\infty)\) 上有界，则（ ）。
\choices{\(a<0,\ b>0\)}{\(a>0,\ b>0\)}{\(a=0,\ b>0\)}{\(a=0,\ b<0\)}
\QuestionEnd

\begin{QuestionSolution}{2023 年第 4 题}
\textbf{已知与目标}\quad 方程的所有解在整个实轴上有界，即每个解的绝对值都可由某个有限常数控制。求 \(a,b\) 应满足的关系。

\textbf{方法与条件}\quad 特征根决定解中的指数和三角因子。题目要求在 \(x\to+\infty\) 和 \(x\to-\infty\) 两端都不增长，必须逐类排除可能的根。

\textbf{推导}\quad 特征方程为 \(r^2+ar+b=0\)。若有非零实根 \(r\)，解 \(\mathrm{e}^{rx}\) 在一端增长到无穷，故不允许。若两根都是零，则通解 \(C_1+C_2x\) 含无界的 \(x\)，也不允许。因此只能是共轭复根 \(\sigma\pm i\omega\)，其中 \(\omega>0\)，对应解
\[
y=\mathrm{e}^{\sigma x}(C_1\cos\omega x+C_2\sin\omega x).
\]
如果 \(\sigma\ne0\)，取余弦等于 \(1\) 的点列，指数仍在某一端无限增长，故必须 \(\sigma=0\)。由根的和为 \(-a\)、积为 \(b\)，得到 \(-a=0\)、\(b=\omega^2>0\)。

\textbf{结论与检查}\quad 选 C，即 \(\boxed{a=0,b>0}\)。反过来，此时 \(y=C_1\cos(\sqrt{b}x)+C_2\sin(\sqrt{b}x)\)，满足 \(|y|\le |C_1|+|C_2|\)，确实处处有界。边界 \(b=0\) 会出现线性增长解，故不能取等号。
\end{QuestionSolution}

% Q 2013 13
\QuestionPage{4.2 高阶线性微分方程}{填空题}{2013 年第 13 题}{2}{非齐次方程的特解与初值}{226 / 234}{***}
已知$y_1 = \mathrm{e}^{3x} - x\mathrm{e}^{2x}$，$y_2 = \mathrm{e}^x - x\mathrm{e}^{2x}$，$y_3 = - x\mathrm{e}^{2x}$
是某二阶常系数非齐次线性微分方程的$3$个解，
该方程满足条件$y\big|_{x = 0} = 0$，$y'\big|_{x = 0} = 1$的解为$y =$ \blank{}．
\QuestionEnd

\begin{QuestionSolution}{2013 年第 13 题}
\textbf{已知与目标}\quad 已知同一二阶常系数非齐次方程的三个解，要求满足 \(y(0)=0,y'(0)=1\) 的解，不需要先写出原方程。

\textbf{方法与条件}\quad 两个非齐次解相减会消去共同右端，得到齐次解。若两个齐次解无法互相表示成常数倍，就线性无关，可组成二阶齐次通解。再加一个已知非齐次解即可。

\textbf{推导}\quad 相减得 \(y_1-y_3=\mathrm{e}^{3x}\)、\(y_2-y_3=\mathrm{e}^x\)。两者比值为 \(\mathrm{e}^{2x}\)，不是常数，故线性无关。原方程通解为
\[
y=C_1\mathrm{e}^{3x}+C_2\mathrm{e}^x-x\mathrm{e}^{2x},
\quad y'=3C_1\mathrm{e}^{3x}+C_2\mathrm{e}^x-(1+2x)\mathrm{e}^{2x}.
\]
代入两个初值，得
\[
C_1+C_2=0,\qquad 3C_1+C_2-1=1 .
\]
第二式减第一式得 \(2C_1=2\)，故 \(C_1=1,C_2=-1\)。于是 \(\boxed{y=\mathrm{e}^{3x}-\mathrm{e}^x-x\mathrm{e}^{2x}}\)。

\textbf{结论与检查}\quad 答案就是 \(y_1-y_2+y_3\)，三项的系数和为 \(1\)，因此仍满足同一非齐次方程。\(y(0)=1-1=0\)，\(y'(0)=3-1-1=1\)，两初值均正确。
\end{QuestionSolution}

% Q 2020 13
\QuestionPage{4.2 高阶线性微分方程}{填空题}{2020 年第 13 题}{2}{齐次方程初值问题与积分}{228 / 234}{***}
设 \(y = y\left( x \right)\) 满足 \(y^{\prime\prime} + 2y^{\prime} + y = 0\)，且 \(y\left( 0 \right) = 0\)，\(y^{\prime}\left( 0 \right) = 1\)，则 \(\int_{0}^{+\infty}{y\left( x \right) \,\mathrm{d} x} =\)
\blank{}
\QuestionEnd

\begin{QuestionSolution}{2020 年第 13 题}
\textbf{已知与目标}\quad 已知二阶齐次方程和两个初值，求解函数在正半轴上的积分。

\textbf{方法与条件}\quad 先用特征方程及初值确定 \(y\)，再把无穷上限积分定义为有限上限积分的极限。计算采用分部积分，把 \(x\) 求导、指数积分，可以减少多项式次数。

\textbf{推导}\quad 特征方程 \(r^2+2r+1=(r+1)^2=0\)，故 \(y=(C_1+C_2x)\mathrm{e}^{-x}\)。其导数为 \((C_2-C_1-C_2x)\mathrm{e}^{-x}\)。初值给出 \(C_1=0,C_2-C_1=1\)，所以 \(y=x\mathrm{e}^{-x}\)。对 \(R>0\)，
\[
\int_0^R x\mathrm{e}^{-x}\,\mathrm{d}x
=[-x\mathrm{e}^{-x}]_0^R+\int_0^R\mathrm{e}^{-x}\,\mathrm{d}x
=1-(R+1)\mathrm{e}^{-R}.
\]
因为 \(\lim_{R\to\infty}\frac{R+1}{\mathrm{e}^R}=0\)（可用 L'Hospital 法则求得），所以所求为 \(\boxed{1}\)。

\textbf{结论与检查}\quad \(y'= (1-x)\mathrm{e}^{-x}\)、\(y''=(x-2)\mathrm{e}^{-x}\)，代回方程和初值均成立。被积函数非负且有指数衰减，所得积分为正并收敛。
\end{QuestionSolution}

% Q 2025 3
\QuestionPage{4.2 高阶线性微分方程}{选择题}{2025 年第 3 题}{1}{常系数方程与反常积分收敛}{223 / 234}{****}
如果对微分方程 \(y^{\prime\prime} - 2a y^{\prime} + \left( a + 2 \right) y = 0\) 的任一解 \(y\left( x \right)\)，反常积分 \(\int_{0}^{+\infty} y\left( x \right) \,\mathrm{d} x\) 均收敛，那么 \(a\) 的取值范围是
\choices{\(\left( -2, -1 \right]\)}{\(\left( -\infty, -1 \right]\)}{\(\left( -2, 0 \right)\)}{\(\left( -\infty, 0 \right)\)}
\QuestionEnd

\begin{QuestionSolution}{2025 年第 3 题}
\textbf{已知与目标}\quad 要求方程每一个解在 \([0,\infty)\) 上的反常积分都收敛，即 \(\lim_{R\to\infty}\int_0^R y(x)\,\mathrm{d}x\) 是有限实数，求参数 \(a\) 的范围。

\textbf{方法与条件}\quad 常系数方程的解由特征根构造。先判断每个基本解的积分能否收敛，再用二次方程根的和、积表达参数条件，可以避免逐一讨论判别式的复杂区间。

\textbf{推导}\quad 特征方程为 \(r^2-2ar+(a+2)=0\)。若所有根的实部都为负，解中的指数随 \(x\) 增大而衰减，\(\mathrm{e}^{rx}\)、重根时的 \(x\mathrm{e}^{rx}\) 以及带正弦、余弦的衰减项，绝对值积分均有限。

反过来，若有非负实根，\(\mathrm{e}^{rx}\) 的积分就发散。若有复根 \(\sigma\pm i\omega\)，\(\sigma\ge0,\omega>0\)，则解 \(\mathrm{e}^{\sigma x}\cos\omega x\) 的一个原函数为
\[
\frac{\mathrm{e}^{\sigma x}(\sigma\cos\omega x+\omega\sin\omega x)}
{\sigma^2+\omega^2},
\]
它在无穷远或振荡或增长，没有有限极限。因此所求条件确为两根实部均负。

设两根为 \(r_1,r_2\)，有 \(r_1+r_2=2a\)、\(r_1r_2=a+2\)。实根时，积正使两根同号，再由和负判定两根都负；复根时，实部为 \(a\)，积为模长平方。因此统一条件为
\[
2a<0,\qquad a+2>0,
\quad\Longleftrightarrow\quad -2<a<0.
\]

\textbf{结论与检查}\quad 选 C，\(\boxed{a\in(-2,0)}\)。在 \(a=-2\) 时出现零根，常数解积分发散；在 \(a=0\) 时出现纯振荡解，积分无极限，两端均不包括。
\end{QuestionSolution}

% Q 2009 20
\QuestionPage{4.2 高阶线性微分方程}{解答题}{2009 年第 20 题}{4}{分段求解微分方程}{231 / 234}{****}
设$y = y(x)$是区间$(- \pi ,\pi)$内过点
$\left(- \frac{\pi}{\sqrt{2}},\frac{\pi}{\sqrt{2}} \right)$的光滑曲线，
当$- \pi < x < 0$时，曲线上任一点处的法线都过原点，
当$0 \le x < \pi$时，函数$y(x)$满足$y'' + y + x = 0$．求$y(x)$的表达式．
\QuestionEnd

\begin{QuestionSolution}{2009 年第 20 题}
\textbf{已知与目标}\quad 左半区间由法线条件约束，右半区间满足二阶方程，整条曲线经过指定点且光滑相接。求两段函数，并确定它们在 \(x=0\) 处的连接条件。

\textbf{方法与条件}\quad 法线垂直于切线，所以法线经过原点可转成向量内积为零。先从已知点求左段，再由连接处函数值与一阶导数连续，为右段提供两个初值。右段采用齐次通解加特解。

\textbf{推导}\quad 对 \(-\pi<x<0\)，切向量为 \((1,y')\)，原点到曲线点的向量为 \((x,y)\)，法线条件给出 \((x,y)\cdot(1,y')=x+yy'=0\)。由于 \((x^2+y^2)'=2x+2yy'\)，积分得 \(x^2+y^2=C\)。代入已知点，得 \(C=\frac{\pi^2}{2}+\frac{\pi^2}{2}=\pi^2\)。在该开区间 \(y^2=\pi^2-x^2>0\)，由连续性及已知点纵坐标为正，整段应取正根：
\[
y=\sqrt{\pi^2-x^2},\qquad y'=-\frac{x}{\sqrt{\pi^2-x^2}}.
\]
令 \(x\to0^-\)，得到 \(y(0)=\pi,y'(0)=0\)，这是光滑连接要求的右段初值。

在右侧，方程写成 \(y''+y=-x\)。齐次特征根为 \(\pm i\)，所以齐次通解为 \(A\cos x+B\sin x\)；取特解 \(y_p=-x\)，它满足 \(y_p''+y_p=-x\)。因此 \(y=A\cos x+B\sin x-x\)，\(y'=-A\sin x+B\cos x-1\)。代入初值得 \(A=\pi,B-1=0\)。综上，
\[
\boxed{y(x)=
\begin{cases}
\sqrt{\pi^2-x^2},&-\pi<x<0,\\
\pi\cos x+\sin x-x,&0\le x<\pi.
\end{cases}}
\]

\textbf{结论与检查}\quad 左段满足 \(x+yy'=0\) 且经过指定点；两段在零点的函数值均为 \(\pi\)，斜率均为 \(0\)，相接时没有折角。右段求导两次后有 \(y''=-\pi\cos x-\sin x\)，故在右侧 \(y''+y+x=0\)。
\end{QuestionSolution}

% Q 2012 19
\QuestionPage{4.2 高阶线性微分方程}{解答题}{2012 年第 19 题}{3}{常系数方程与积分曲线拐点}{232 / 234}{****}
设函数$f(x)$满足方程$f''(x) + f'(x) - 2f(x) = 0$及$f'(x) + f(x) = 2\mathrm{e}^x$．
\begin{enumerate}
\item 求表达式$f(x)$；
\item 求曲线$y = f(x^2)\int_0^x f(- t^2)\,\mathrm{d} t$的拐点．
\end{enumerate}
\QuestionEnd

\begin{QuestionSolution}{2012 年第 19 题}
\textbf{已知与目标}\quad （1）求同时满足二阶方程及一阶关系的 \(f\)；（2）将它代入含积分的曲线，求拐点。拐点是曲线凹凸性发生改变的点，不能仅凭二阶导数等于零判定。

\textbf{方法与条件}\quad 先求常系数方程通解，再用附加关系确定常数。第二问的积分无需求原函数：用乘积求导与变上限积分求导，随后判断 \(y''\) 的符号。

\textbf{推导}\quad （1）特征方程 \(r^2+r-2=(r-1)(r+2)=0\)，得 \(f=A\mathrm{e}^x+B\mathrm{e}^{-2x}\)。代入附加关系，
\[
f'+f=2A\mathrm{e}^x-B\mathrm{e}^{-2x}=2\mathrm{e}^x.
\]
两种指数函数线性无关，比较系数得 \(A=1,B=0\)，所以 \(\boxed{f(x)=\mathrm{e}^x}\)。

（2）代入后 \(y=\mathrm{e}^{x^2}\int_0^x\mathrm{e}^{-t^2}\,\mathrm{d}t\)。被积函数连续，积分求导为 \(\mathrm{e}^{-x^2}\)，因此
\[
y'=2x\mathrm{e}^{x^2}\int_0^x\mathrm{e}^{-t^2}\,\mathrm{d}t
+\mathrm{e}^{x^2}\mathrm{e}^{-x^2}=2xy+1,
\]
再求一次导数并代入 \(y'\)，得到
\[
y''=2y+2xy'=2y+2x(2xy+1)=(2+4x^2)y+2x.
\]
由于 \(\mathrm{e}^{-t^2}>0\)，\(x>0\) 时积分为正，\(x<0\) 时因上下限反向积分为负，所以 \(y\) 与 \(x\) 同号。系数 \(2+4x^2\) 恒正，故 \(x<0\) 时 \(y''<0\)，\(x>0\) 时 \(y''>0\)。

\textbf{结论与检查}\quad \(f=\mathrm{e}^x\) 代回两个原关系均成立。第二问只有零点两侧的凹凸性改变，且 \(y(0)=0\)，所以唯一拐点为 \(\boxed{(0,0)}\)。两侧的严格符号同时排除了其他拐点。
\end{QuestionSolution}

% Q 2016 19
\QuestionPage{4.2 高阶线性微分方程}{解答题}{2016 年第 19 题}{3}{二阶变系数线性方程}{233 / 234}{****}
已知$y_1(x) = \mathrm{e}^x$，$y_2(x) = u(x) \mathrm{e}^x$是二阶微分方程$(2x - 1)y'' - (2x + 1)y' + 2y = 0$的解，
若$u(- 1) = \mathrm{e}$，$u(0) = - 1$，求$u(x)$，并写出该微分方程的通解．
\QuestionEnd

\begin{QuestionSolution}{2016 年第 19 题}
\textbf{已知与目标}\quad 已知 \(\mathrm{e}^x\) 和 \(u\mathrm{e}^x\) 均满足二阶方程，且 \(u(-1)=\mathrm{e},u(0)=-1\)。求 \(u\) 及原方程通解。

\textbf{方法与条件}\quad 代入 \(y=u\mathrm{e}^x\)，利用已知解消去不含 \(u\) 导数的项，再令 \(v=u'\) 降阶。最高阶系数 \(2x-1\) 在 \(x=\frac{1}{2}\) 为零，除法先在不含该点的区间进行；两个给定点都在其左侧。

\textbf{推导}\quad 乘积求导得 \(y'=\mathrm{e}^x(u'+u)\)、\(y''=\mathrm{e}^x(u''+2u'+u)\)。代入原方程，\(u\) 的系数为 \((2x-1)-(2x+1)+2=0\)，故
\[
(2x-1)u''+(2x-3)u'=0.
\]
令 \(v=u'\)，方程变为 \((2x-1)v'+(2x-3)v=0\)。在 \(x\ne\frac{1}{2}\) 的区间，直接求导可验证
\[
\left(\frac{\mathrm{e}^xv}{2x-1}\right)'
=\frac{\mathrm{e}^x\bigl((2x-1)v'+(2x-3)v\bigr)}{(2x-1)^2}=0.
\]
因此 \(v=C(2x-1)\mathrm{e}^{-x}\)。由于 \(((2x+1)\mathrm{e}^{-x})'=-(2x-1)\mathrm{e}^{-x}\)，积分并重命名常数后有 \(u=A+B(2x+1)\mathrm{e}^{-x}\)。代入给定条件，得
\[
A-B\mathrm{e}=\mathrm{e},\qquad A+B=-1.
\]
由 \(A=-1-B\) 代入第一式，得 \(-1-(1+\mathrm{e})B=\mathrm{e}\)，所以 \(B=-1,A=0\)。因此 \(\boxed{u=-(2x+1)\mathrm{e}^{-x}}\)，第二个原方程解为 \(-(2x+1)\)。它与 \(\mathrm{e}^x\) 无法互为常数倍，故在通常的非奇异区间上通解为
\[
\boxed{y=C_1\mathrm{e}^x+C_2(2x+1)}.
\]

\textbf{结论与检查}\quad \(u(-1)=\mathrm{e},u(0)=-1\)。对 \(y=2x+1\) 有 \(y'=2,y''=0\)，原方程左端为 \(-2(2x+1)+2(2x+1)=0\)；\(\mathrm{e}^x\) 亦满足原式。

若要求跨过 \(x=\frac{1}{2}\) 的二阶可导解，左右常数也不能任意拼接。两侧 \(u'\) 都趋于零，而 \(u''(\frac{1}{2})\) 的单侧差商为 \(-2B\mathrm{e}^{-\frac{1}{2}}\)，其存在迫使两侧 \(B\) 相同；再由 \(u\) 连续知 \(A\) 相同。因此上述表达式也给出这样的整体解。
\end{QuestionSolution}

% Q 2024 18
\QuestionPage{4.2 高阶线性微分方程}{解答题}{2024 年第 18 题}{3}{Euler 方程与定积分}{234 / 234}{****}
设 \(y=y(x)\) 满足方程
\[x^2y''+xy'-9y=0,
\]
且
\[y\big|_{x=1}=2,\qquad y'\big|_{x=1}=6.
\]
\begin{enumerate}
\item 利用变换 \(x=\mathrm{e}^t\) 化简方程，并求 \(y(x)\)；
\item 求
\[\int_1^2 y(x)\sqrt{4-x^2}\,\mathrm{d}x.
\]
\end{enumerate}
\QuestionEnd

\begin{QuestionSolution}{2024 年第 18 题}
\textbf{已知与目标}\quad （1）按指定代换 \(x=\mathrm{e}^t\) 解方程及两个初值；（2）用所得 \(y\) 计算定积分。第二问的根号在积分区间 \([1,2]\) 上为实数。

\textbf{方法与条件}\quad Euler 方程的导数系数是对应阶数的 \(x\) 幂，指数代换可把它变成常系数方程。该代换对应 \(x>0\)，涵盖初值及积分区间。第二问令 \(u=4-x^2\)，使根号变成幂函数。

\textbf{推导}\quad （1）记 \(Y(t)=y(\mathrm{e}^t)\)，区分对 \(t\) 与对 \(x\) 的求导。由链式法则，
\[
\frac{\mathrm{d}Y}{\mathrm{d}t}=xy',\qquad
\frac{\mathrm{d}^2Y}{\mathrm{d}t^2}=x\frac{\mathrm{d}}{\mathrm{d}x}(xy')=xy'+x^2y''.
\]
故原式变成 \(Y''-9Y=0\)，特征根为 \(3,-3\)，得 \(Y=C_1\mathrm{e}^{3t}+C_2\mathrm{e}^{-3t}\)。初值 \(x=1\) 对应 \(t=0\)，且 \(Y'(0)=1\cdot y'(1)=6\)，于是
\[
C_1+C_2=2,\qquad 3C_1-3C_2=6.
\]
两式相加、相减可得 \(C_1=2,C_2=0\)。由 \(\mathrm{e}^{3t}=x^3\)，还原为 \(\boxed{y(x)=2x^3,\ x>0}\)。

（2）令 \(u=4-x^2\)，则 \(2x\,\mathrm{d}x=-\mathrm{d}u\)，\(x^2=4-u\)。当 \(x=1,2\) 时，\(u=3,0\)，所以
\[
\begin{aligned}
\int_1^2 2x^3\sqrt{4-x^2}\,\mathrm{d}x
&=\int_0^3 (4-u)\sqrt{u}\,\mathrm{d}u\\
&=4\int_0^3u^{\frac{1}{2}}\,\mathrm{d}u-\int_0^3u^{\frac{3}{2}}\,\mathrm{d}u\\
&=\left[\frac{8}{3}u^{\frac{3}{2}}-\frac{2}{5}u^{\frac{5}{2}}\right]_0^3
=\boxed{\frac{22\sqrt3}{5}}.
\end{aligned}
\]
其中 \(3^{\frac{3}{2}}=3\sqrt{3}\)、\(3^{\frac{5}{2}}=9\sqrt{3}\)，故末式为 \((8-\frac{18}{5})\sqrt{3}\)。

\textbf{结论与检查}\quad \(y'=6x^2,y''=12x\)，代回得 \(12x^3+6x^3-18x^3=0\)，两个初值也成立。被积函数非负，结果 \(\frac{22\sqrt{3}}{5}>0\)；换元倒置积分上下限时已经同时消去负号。
\end{QuestionSolution}

\SectionHeading{4.3}{常微分方程}{线性微分方程组（教材星号节）}{定理与要点\quad 教材 PDF 第 310 页起\hfill 本节 0 题}
\pdfbookmark[2]{4.3 线性微分方程组（教材星号节）}{section-4.3}
\Theory{定理 3.1}{313}{齐次方程组解的相关性}{一组齐次解线性相关，等价于它们在区间内某一点取值所得向量组线性相关。}
\Theory{定理 3.2}{314}{齐次方程组的通解}{\(n\) 维线性齐次方程组有 \(n\) 个线性无关的解，通解为它们的常系数线性组合。}
\Theory{定理 3.3}{315}{Wronski 行列式判别}{\(n\) 个齐次解线性无关的充要条件是其 Wronski 行列式在区间内某点非零。}
\Theory{定理 3.4}{316}{非齐次方程组解的结构}{非齐次通解等于任一特解加对应齐次方程组的通解。}
\Theory{定理 3.5}{317}{基解矩阵与常数变易法}{取基解矩阵 \(X\)，写成 \(\boldsymbol{x}=X\boldsymbol{c}\)，对所得的 \(\boldsymbol{c}'\) 积分可构造零初值特解。}
\Theory{定理 3.6}{320}{特征向量法}{若系数矩阵有 \(n\) 个线性无关的特征向量，则用各特征向量乘对应的指数因子组成基解矩阵。}
\Theory{定理 3.7}{323}{重特征值对应的特解}{重数为 \(m\) 的特征值提供 \(m\) 个线性无关的解，每个解是指数函数乘次数至多为 \(m-1\) 的向量多项式。}
\Theory{定理 3.8}{324}{由各特征值构造基本解组}{合并各不同特征值对应的线性无关解，组成 \(n\) 维常系数齐次系统的基本解组。}
\Theory{定理 3.9}{329}{常系数非齐次系统的通解}{取满足 \(X(0)=E\) 的基解矩阵，非齐次项的贡献由 \(X(t-\tau)\) 乘右端项后积分得到。}
本节在所收原卷中没有以该节为主考点的题目。
\end{document}
